\clearpage
\thispagestyle{plain}
\markboth{예비사항}{예비사항}
\oldpage[IV-1]{5}
\phantomsection
\label{chapter:0IV1-fr}

\begin{center}
{\large 제0장 (계속)\footnotemark\par}
\vspace{1.5em}
{\Large\bfseries 예비사항\par}
\vspace{2em}
{\large\bfseries 목차\par}
\end{center}

\footnotetext{기호 $0_N$ (각각 $\operatorname{Err}_N$)는 제$N$장에
딸린 제0장의 부분 (각각 정오표)을 가리킴을 상기하자.}

\begin{center}
\begin{tabular}{@{}r l@{}}
\S~14. & 위상공간의 조합적 차원.\\
\S~15. & $M$-정칙열과 $\mathscr F$-정칙열.\\
\S~16. & 뇌터 국소환의 차원과 깊이.\\
\S~17. & 정칙환.\\
\S~18. & 대수의 확장에 관한 보충.\\
\S~19. & 형식적으로 매끄러운 대수와 Cohen 환.\\
\S~20. & 미분과 미분형식.\\
\S~21. & 표수가 $p$인 환에서의 미분형식.\\
\S~22. & 형식적 매끄러움과 정칙성의 미분 판정법.\\
\S~23. & 일본환.
\end{tabular}
\end{center}

앞선 절들의 거의 전부는 제IV장에서 사용할 가환대수학의 개념을
설명하는 데 바쳐져 있다. 이 개념들의 상당 부분은 이미 여러 문헌
([1], [12], [13], [17], [30])에 실려 있지만, 독자가 연속적이고 거의
독립적인 설명을 접하는 편이 더 편리하리라고 생각하였다. 이 절들은
제IV장의 \S\S~5, 6, 7 (여기서는 스킴의 언어를 사용한다)과 합쳐져,
우리 논고 안에서 제I장부터 제III장까지와 거의 독립적인 작은 전문
논고를 이룬다. 그 목적은 완비화나 정수적 폐포 같은 연산에 관하여
``좋은 거동을 보이는'' 환의 성질을 일관된 형태로 제시하고, 이
성질들을 몇 가지 일반적인 관점과 체계적으로 연결하는 데
있다\footnote{여기서 다루는 성질의 대부분은 슈발레, 자리스키,
나가타, 세르가 발견하였다. 여기서 사용하는 방법은 A.~그로텐디크가
1961년 가을 하버드 대학교에서 강의한 과정에서 처음 전개하였다.}.

\setcounter{section}{13}
\section{위상공간의 조합적 차원}
\label{section:0.14-fr}

\setcounter{subsection}{0}
\subsection{위상공간의 조합적 차원}
\label{subsection:0.14.1-fr}

\oldpage[0\textsubscript{IV-1}]{6}
\begin{env}[14.1.1]
\label{0.14.1.1-fr}
$I$를 순서집합이라 하자~; $I$의 원소들의 \emph{사슬}은 정의상
$I$의 원소들로 이루어진 유한한 강한 증가열
$i_0<i_1<\cdots<i_n$ ($n\geq 0$)이고~; 이 사슬의 \emph{길이}는
정의상 $n$이다. $X$가 위상공간이면, 그 \emph{기약 닫힌 부분집합}들의
집합은 포함관계로 순서가 정해지며, 따라서 $X$의 기약 닫힌
부분집합들의 \emph{사슬}이라는 개념을 얻는다.
\end{env}

\begin{definition}[14.1.2]
\label{0.14.1.2-fr}
$X$를 위상공간이라 하자. $X$의 기약 닫힌 부분집합들의 사슬의 길이들의
상한을 $X$의 \emph{조합적 차원} (혼동의 여지가 없으면 단순히 $X$의
\emph{차원})이라 하고, $\dim_c(X)$ 또는 단순히 $\dim(X)$로 나타낸다.
각 $x\in X$에 대하여, $X$에서 $x$의 열린 근방 전체를 $U$가 달릴 때의
수 $\inf_U(\dim(U))$를 \emph{점 $x$에서 $X$의 조합적 차원} (또는
단순히 \emph{점 $x$에서 $X$의 차원})이라 하고 $\dim_x(X)$로 나타낸다.
\end{definition}

이 정의로부터 다음을 얻는다.
\[
  \dim(\varnothing)=-\infty
\]
(공집합의 $\overline{\mathbb R}$에서의 상한은 $-\infty$이다). $X$의
기약성분들의 족을 $(X_\alpha)$라 하면,
\[
  \dim(X)=\sup_\alpha\dim(X_\alpha)
  \tag{14.1.2.1}\label{0.14.1.2.1-fr}
\]
이다. 실제로 $X$의 기약 닫힌 부분집합들의 모든 사슬은 정의상 하나의
기약성분에 포함되고, 거꾸로 이 성분들은 $X$에서 닫혀 있으므로 각
$X_\alpha$의 모든 기약 닫힌 부분집합은 $X$의 기약 닫힌 부분집합이다.

\begin{definition}[14.1.3]
\label{0.14.1.3-fr}
위상공간 $X$의 모든 기약성분이 같은 차원을 가지면 $X$가
\emph{등차원}이라고 한다 (이 공통 차원은 (14.1.2.1)에 따라
$\dim(X)$와 같다).
\end{definition}

\begin{proposition}[14.1.4]
\label{0.14.1.4-fr}
\begin{enumerate}
  \item[\rm(i)] 위상공간 $X$의 임의의 부분집합 $Y$에 대하여
  $\dim(Y)\leq\dim(X)$이다.
  \item[\rm(ii)] 위상공간 $X$가 유한 개의 닫힌 부분집합 $X_i$들의
  합집합이면, $\dim(X)=\sup_i\dim(X_i)$이다.
\end{enumerate}
\end{proposition}

\begin{proof}
$Y$의 임의의 기약 닫힌 부분집합 $Z$에 대하여, $X$에서 $Z$의 폐포
$\overline Z$는 기약이고 ($0_{\mathrm I}$, 2.1.2),
$\overline Z\cap Y=Z$이므로 (i)을 얻는다. 한편
$X=\bigcup_{i=1}^n X_i$이고 $X_i$들이 닫혀 있으면, $X$의 모든 기약
닫힌 부분집합은 어느 한 $X_i$에 포함된다 ($0_{\mathrm I}$, 2.1.1).
따라서 $X$의 기약 닫힌 부분집합들의 모든 사슬은 어느 한 $X_i$에
포함되며, 이로부터 (ii)를 얻는다.
\end{proof}

(14.1.4, (i))로부터 모든 $x\in X$에 대하여 다음과 같이 쓸 수도 있다.
\[
  \dim_x(X)=\lim_U\dim(U)
  \tag{14.1.4.1}\label{0.14.1.4.1-fr}
\]
여기서 극한은 $X$에서 $x$의 열린 근방들의 하향 여과집합을 따라
취한다.

\oldpage[0\textsubscript{IV-1}]{7}
\begin{corollary}[14.1.5]
\label{0.14.1.5-fr}
$X$를 위상공간, $x$를 $X$의 점, $U$를 $x$의 근방이라 하고,
$Y_i$ ($1\leq i\leq n$)를 $U$의 닫힌 부분집합들이라 하자. 모든 $i$에
대하여 $x\in Y_i$이고 $U$가 $Y_i$들의 합집합이면,
\[
  \dim_x(X)=\sup_i\bigl(\dim_x(Y_i)\bigr).
  \tag{14.1.5.1}\label{0.14.1.5.1-fr}
\]
\end{corollary}

\begin{proof}
(14.1.4, (ii))에 의하여, $V$가 $U$에 포함되는 $x$의 열린 근방
전체를 달릴 때
$\dim_x(X)=\inf_V(\sup_i(\dim(Y_i\cap V)))$이다. 마찬가지로 모든
$i$에 대하여 $\dim_x(Y_i)=\inf_V(\dim(Y_i\cap V))$이다. 따라서
\[
  \sup_i\bigl(\dim_x(Y_i)\bigr)=+\infty
\]
이면 따름정리는 자명하다. 그렇지 않으면, $x$의 열린 근방
$V_0\subset U$가 존재하여 $1\leq i\leq n$ 및 모든 $V\subset V_0$에
대하여 $\dim(Y_i\cap V)=\dim_x(Y_i)$가 되므로 결론을 얻는다.
\end{proof}

\begin{proposition}[14.1.6]
\label{0.14.1.6-fr}
모든 위상공간 $X$에 대하여
$\dim(X)=\sup_{x\in X}\dim_x(X)$이다.
\end{proposition}

\begin{proof}
정의 (14.1.2)와 명제 (14.1.4)에 의하여 모든 $x\in X$에 대하여
$\dim_x(X)\leq\dim(X)$이다. 한편
$Z_0\subset Z_1\subset\cdots\subset Z_n$을 $X$의 기약 닫힌
부분집합들의 사슬이라 하고 $x\in Z_0$라 하자. $x$를 포함하는 임의의
열린집합 $U\subset X$에 대하여 $U\cap Z_i$는 기약이고
($0_{\mathrm I}$, 2.1.6) $U$에서 닫혀 있다. 또한 $X$에서
$\overline{U\cap Z_i}=Z_i$이므로 $U\cap Z_i$들은 서로 다르다.
따라서 $\dim(U)\geq n$이고, 이로써 증명이 끝난다.
\end{proof}

\begin{corollary}[14.1.7]
\label{0.14.1.7-fr}
$(X_\alpha)$가 $X$의 열린 덮개이거나 국소 유한 닫힌 덮개이면,
$\dim(X)=\sup_\alpha(\dim(X_\alpha))$이다.
\end{corollary}

\begin{proof}
$X_\alpha$가 $x\in X$의 근방이면
$\dim_x(X)\leq\dim(X_\alpha)$이므로 첫 번째 주장을 얻는다. 한편
$X_\alpha$들이 닫혀 있고 $U$가 유한 개의 $X_\alpha$만 만나는
$x$의 근방이면, (14.1.4)에 의하여
\[
  \dim_x(X)\leq\dim(U)=\sup_\alpha\bigl(\dim(U\cap X_\alpha)\bigr)
  \leq\sup_\alpha\bigl(\dim(X_\alpha)\bigr)
\]
이므로 두 번째 주장을 얻는다.
\end{proof}

\begin{corollary}[14.1.8]
\label{0.14.1.8-fr}
$X$를 뇌터 콜모고로프 공간 ($0_{\mathrm I}$, 2.1.3)이라 하고,
$F$를 $X$의 닫힌점들의 집합이라 하자. 그러면
$\dim(X)=\sup_{x\in F}\dim_x(X)$이다.
\end{corollary}

\begin{proof}
(14.1.6)의 증명에서 사용한 기호에 따르면, $Z_0$ 안에 닫힌점이
존재함을 지적하면 충분하다 ($0_{\mathrm I}$, 2.1.3).
\end{proof}

\begin{proposition}[14.1.9]
\label{0.14.1.9-fr}
$X$를 공집합이 아닌 뇌터 콜모고로프 공간이라 하자.
$\dim(X)=0$일 필요충분조건은 $X$가 유한 이산공간인 것이다.
\end{proposition}

\begin{proof}
공간 $X$가 하우스도르프 공간이면 (특히 $X$가 이산공간이면), $X$의 기약 닫힌
부분집합은 한 점으로 이루어진 것뿐이므로 $\dim(X)=0$이다. 거꾸로,
$X$가 $\dim(X)=0$인 뇌터 콜모고로프 공간이라고 하자. $X$의 모든
기약성분은 닫힌점을 하나 포함하므로 ($0_{\mathrm I}$, 2.1.3), 그
성분은 이 점 하나로 이루어져야 한다. $X$의 기약성분은 유한 개뿐이므로
$X$는 유한 이산공간이다.
\end{proof}

\begin{corollary}[14.1.10]
\label{0.14.1.10-fr}
$X$를 뇌터 콜모고로프 공간이라 하자. 점 $x\in X$가 고립점일
필요충분조건은 $\dim_x(X)=0$인 것이다.
\end{corollary}

\begin{proof}
그 조건이 필요함은 자명하다 ($X$에 아무 가정도 하지 않아도 된다).
또한 그 조건은 충분하다.
\oldpage[0\textsubscript{IV-1}]{8}
실제로 이 조건으로부터 $x$의 어떤 열린 근방 $U$에 대하여
$\dim(U)=0$임을 얻으며, $U$는 뇌터 콜모고로프 공간이므로 유한
이산공간이다.
\end{proof}

\begin{proposition}[14.1.11]
\label{0.14.1.11-fr}
함수 $x\mapsto\dim_x(X)$는 $X$에서 상반연속이다.
\end{proposition}

\begin{proof}
이 함수의 값이 $+\infty$인 모든 점에서 상반연속임은 자명하다. 이제
$\dim_x(X)=n<+\infty$라 하자. 공식 (14.1.4.1)에 의하여 $x$의 열린
근방 $U_0$가 존재하여, $x$의 모든 열린 근방 $U\subset U_0$에 대하여
$\dim(U)=n$이다. 그러면 모든 $y\in U_0$와 $y$의 모든 열린 근방
$V\subset U_0$에 대하여
$\dim(V)\leq\dim(U_0)=n$이고 (14.1.4), 따라서 (14.1.4.1)로부터
$\dim_y(X)\leq n$을 얻는다.
\end{proof}

\begin{remark}[14.1.12]
\label{0.14.1.12-fr}
$X$, $Y$가 두 위상공간이고 $f:X\to Y$가 연속사상이더라도
$\dim(f(X))>\dim(X)$일 수 있다. 예를 들어, $X$는 두 원소 $a$, $b$로
이루어진 이산공간으로 하고, $Y$는 집합 $\{a,b\}$에 닫힌집합들이
$\varnothing$, $\{a\}$, $\{a,b\}$가 되도록 위상을 준 공간으로 하자.
$f:X\to Y$가 항등사상이면 $\dim(Y)=1$이고 $\dim(X)=0$이다. 여기서
$Y$는 이산 값매김환 $A$의 스펙트럼으로 볼 수 있고, $a$는 유일한
닫힌점이며 $b$는 일반점이다. $K$와 $k$를 각각 $A$의 분수체와
잉여체라 하면, $X$는 환 $k\times K$의 스펙트럼이고 $f$는 준동형
$(\varphi,\psi):A\to k\times K$에 대응하는 연속사상이다. 여기서
$\varphi:A\to k$와 $\psi:A\to K$는 표준 준동형이다 (IV, 5.4.3 참조).
\end{remark}

\subsection{닫힌 부분집합의 여차원}
\label{subsection:0.14.2-fr}

\begin{definition}[14.2.1]
\label{0.14.2.1-fr}
위상공간 $X$의 기약 닫힌 부분집합 $Y$가 주어졌다고 하자. $Y$가
최소 원소인 $X$의 기약 닫힌 부분집합들의 사슬 길이들의 상한을
$X$에서 $Y$의 \emph{조합적 여차원} (또는 단순히 \emph{여차원})이라
하고, $\operatorname{codim}(Y,X)$로 나타낸다. $Y$가 $X$의 임의의
닫힌 부분집합이면, $Y$의 기약성분들이 $X$에서 갖는 여차원들의 하한을
$X$에서 $Y$의 \emph{여차원}이라 하고, 역시
$\operatorname{codim}(Y,X)$로 나타낸다. $X$의 모든 극소 기약 닫힌
부분집합이 $X$에서 같은 여차원을 가지면 $X$가 \emph{등여차원}이라고
한다.
\end{definition}

이 정의로부터 $\operatorname{codim}(\varnothing,X)=+\infty$임을
얻는다. 실제로 $\overline{\mathbb R}$에서 공집합의 하한은
$+\infty$이다. $Y$가 $X$에서 닫혀 있고 $(X_\alpha)$ (각각
$(Y_\beta)$)가 $X$ (각각 $Y$)의 기약성분들의 족이면, 모든
$Y_\beta$는 어느 $X_\alpha$에 포함된다. 더 일반적으로,
$Y_\beta$가 최소 원소인 $X$의 기약 닫힌 부분집합들의 모든 사슬은
어느 $X_\alpha$의 부분집합들로 이루어진다. 따라서
\[
  \operatorname{codim}(Y,X)=
  \inf_\beta\left(\sup_\alpha
  \bigl(\operatorname{codim}(Y_\beta,X_\alpha)\bigr)\right)
  \tag{14.2.1.1}\label{0.14.2.1.1-fr}
\]
이고, 여기서 각 $\beta$에 대하여 $\alpha$는
$Y_\beta\subset X_\alpha$인 지표 전체를 달린다.

\begin{proposition}[14.2.2]
\label{0.14.2.2-fr}
$X$를 위상공간이라 하자.
\begin{enumerate}
  \item[\rm(i)] $\Phi$가 $X$의 기약 닫힌 부분집합 전체의 집합이면,
  \[
    \dim(X)=\sup_{Y\in\Phi}\bigl(\operatorname{codim}(Y,X)\bigr).
    \tag{14.2.2.1}\label{0.14.2.2.1-fr}
  \]
  \oldpage[0\textsubscript{IV-1}]{9}
  \item[\rm(ii)] $X$의 공집합이 아닌 모든 닫힌 부분집합 $Y$에 대하여
  \[
    \dim(Y)+\operatorname{codim}(Y,X)\leq\dim(X).
    \tag{14.2.2.2}\label{0.14.2.2.2-fr}
  \]
  \item[\rm(iii)] $Y\subset Z\subset T$인 $X$의 세 닫힌 부분집합
  $Y$, $Z$, $T$에 대하여
  \[
    \operatorname{codim}(Y,Z)+\operatorname{codim}(Z,T)
    \leq\operatorname{codim}(Y,T).
    \tag{14.2.2.3}\label{0.14.2.2.3-fr}
  \]
  \item[\rm(iv)] $X$의 닫힌 부분집합 $Y$에 대하여
  $\operatorname{codim}(Y,X)=0$일 필요충분조건은 $Y$가 $X$의
  기약성분 하나를 포함하는 것이다.
\end{enumerate}
\end{proposition}

\begin{proof}
(i)과 (iv)는 정의 (14.2.1)의 직접적인 결과이다. (ii)를 증명할 때는
$Y$가 기약인 경우만 생각해도 되고, 이 경우 공식은 정의 (14.1.1)과
(14.2.1)에서 따른다. 마지막으로 (iii)을 증명할 때는 정의 (14.2.1)에
의하여 먼저 $Y$가 기약인 경우만 생각해도 된다. 그러면 $Y$를 포함하는
$Z$의 기약성분들을 $Z_\alpha$라 할 때
$\operatorname{codim}(Y,Z)=
\sup_\alpha(\operatorname{codim}(Y,Z_\alpha))$이다. 분명히
$\operatorname{codim}(Y,T)\geq\operatorname{codim}(Y,Z)$이므로
$\operatorname{codim}(Y,Z)=+\infty$이면 부등식은 성립한다. 그렇지
않으면 어떤 $\alpha$에 대하여
$\operatorname{codim}(Y,Z)=\operatorname{codim}(Y,Z_\alpha)$이고,
(14.2.1)에 의하여 $Z$도 기약인 경우만 생각해도 된다. 이 경우
부등식 (14.2.2.3)은 정의 (14.2.1)의 직접적인 결과이다.
\end{proof}

\begin{proposition}[14.2.3]
\label{0.14.2.3-fr}
$X$를 위상공간, $Y$를 $X$의 닫힌 부분집합이라 하자. $X$의 모든
열린집합 $U$에 대하여
\[
  \operatorname{codim}(Y\cap U,U)\geq\operatorname{codim}(Y,X).
  \tag{14.2.3.1}\label{0.14.2.3.1-fr}
\]
또한 (14.2.3.1)의 양변이 같을 필요충분조건은, $U$와 만나는 $Y$의
기약성분들의 족을 $(Y_\alpha)$라 할 때
$\operatorname{codim}(Y,X)=
\inf_\alpha(\operatorname{codim}(Y_\alpha,X))$인 것이다.
\end{proposition}

\begin{proof}
$Z\mapsto\overline Z$가 $U$의 기약 닫힌 부분집합들의 집합에서 $U$와
만나는 $X$의 기약 닫힌 부분집합들의 집합으로 가는 전단사임이 알려져
있다 ($0_{\mathrm I}$, 2.1.6). 특히 이 전단사는 $Y\cap U$의
기약성분들을 $U$와 만나는 $Y$의 기약성분들에 대응시킨다. $Y_\alpha$가
후자의 기약성분 하나이면
$\operatorname{codim}(Y_\alpha,X)=
\operatorname{codim}(Y_\alpha\cap U,U)$이고, 따라서 명제는 정의
(14.2.1)에서 따른다.
\end{proof}

\begin{definition}[14.2.4]
\label{0.14.2.4-fr}
$X$를 위상공간, $Y$를 $X$의 닫힌 부분집합, $x$를 $X$의 점이라 하자.
$U$가 $X$에서 $x$의 열린 근방 전체를 달릴 때의 수
$\sup_U(\operatorname{codim}(Y\cap U,U))$를
\emph{점 $x$에서 $X$ 안의 $Y$의 여차원}이라 하고
$\operatorname{codim}_x(Y,X)$로 나타낸다.
\end{definition}

(14.2.3)에 의하여 다음과 같이 쓸 수도 있다.
\[
  \operatorname{codim}_x(Y,X)=
  \lim_U\bigl(\operatorname{codim}(Y\cap U,U)\bigr)
  \tag{14.2.4.1}\label{0.14.2.4.1-fr}
\]
여기서 극한은 $X$에서 $x$의 열린 근방들의 하향 여과집합을 따라
취한다. 또한 다음이 성립함에 주의하자.
\[
  \operatorname{codim}_x(Y,X)=+\infty\qquad(x\in X-Y\text{이면}).
\]

\oldpage[0\textsubscript{IV-1}]{10}
\begin{proposition}[14.2.5]
\label{0.14.2.5-fr}
$(Y_i)_{1\leq i\leq n}$을 위상공간 $X$의 닫힌 부분집합들로 이루어진
유한한 족이라 하고, $Y$를 이 족의 합집합이라 하자. 그러면
\[
  \operatorname{codim}(Y,X)=
  \inf_i\bigl(\operatorname{codim}(Y_i,X)\bigr).
  \tag{14.2.5.1}\label{0.14.2.5.1-fr}
\]
\end{proposition}

\begin{proof}
실제로 각 $Y_i$의 모든 기약성분은 $Y$의 어느 기약성분에 포함되고,
거꾸로 $Y$의 모든 기약성분은 어느 한 $Y_i$의 기약성분이기도 하다
($0_{\mathrm I}$, 2.1.1). 따라서 결론은 정의 (14.2.1)과 부등식
(14.2.2.3)에서 따른다.
\end{proof}

\begin{corollary}[14.2.6]
\label{0.14.2.6-fr}
$X$를 위상공간, $Y$를 $X$의 국소 뇌터인 닫힌 부분공간이라 하자.
\begin{enumerate}
  \item[\rm(i)] 모든 $x\in X$에 대하여 $x$를 포함하는 $Y$의
  기약성분 $Y_i$ ($1\leq i\leq n$)는 유한 개뿐이고,
  \[
    \operatorname{codim}_x(Y,X)=
    \inf_i\bigl(\operatorname{codim}(Y_i,X)\bigr).
  \]
  \item[\rm(ii)] 함수
  $x\mapsto\operatorname{codim}_x(Y,X)$는 $X$에서 하반연속이다.
\end{enumerate}
\end{corollary}

\begin{proof}
실제로 가정에 의하여 $X$에서 $x$의 열린 근방 $U_0$가 존재하여
$Y\cap U_0$가 뇌터이다. 따라서 그 기약성분은 유한 개뿐이며, 이들은
$Y$의 기약성분들을 $U_0$에 제한한 자취들이다. 특히 $x$를 포함하는
$Y$의 기약성분 $Y_i$ ($1\leq i\leq n$)는 유한 개뿐이다. $x$를
포함하지 않는 $Y_j$들과 만나지 않는 $x$의 열린 근방
$U\subset U_0$로 $U_0$를 바꾸면, $Y_i\cap U$들이 $Y\cap U$의
기약성분들이라고 가정할 수 있다. 그러면 $X$에서 $x$의 모든 열린 근방
$V\subset U$에 대하여 $Y_i\cap V$들이 $Y\cap V$의 기약성분들이고,
(14.2.3)에 의하여
$\operatorname{codim}(Y_i,X)=\operatorname{codim}(Y_i\cap V,V)$이다.
이로써 (i)이 증명된다. 또한 모든 $x'\in U$에 대하여 $x'$를 포함하는
$Y$의 기약성분들은 $Y_i$들 가운데 일부이므로
$\operatorname{codim}_{x'}(Y,X)\geq\operatorname{codim}_x(Y,X)$이고,
이로써 (ii)가 증명된다.
\end{proof}

\subsection{사슬 조건}
\label{subsection:0.14.3-fr}

\begin{env}[14.3.1]
\label{0.14.3.1-fr}
위상공간 $X$에서 기약 닫힌 부분집합들의 사슬
$Z_0\subset Z_1\subset\cdots\subset Z_n$에 대하여, 어떤 지표 $k$와
$Z_k\subset Z'\subset Z_{k+1}$을 만족하면서 모든 $Z_i$와 다른 기약
닫힌 부분집합 $Z'$가 존재하지 않으면 이 사슬을 \emph{포화 사슬}이라고
한다.
\end{env}

\begin{proposition}[14.3.2]
\label{0.14.3.2-fr}
$X$를 다음 성질을 갖는 위상공간이라 하자. $Y\subset Z$인 $X$의 임의의
두 기약 닫힌 부분집합 $Y$, $Z$에 대하여
$\operatorname{codim}(Y,Z)<+\infty$이다. 다음 두 조건은 동치이다.
\begin{enumerate}
  \item[\rm(a)] 양 끝점이 같은 기약 닫힌 부분집합들의 두 포화 사슬은
  길이가 같다.
  \item[\rm(b)] $Y\subset Z\subset T$인 $X$의 세 기약 닫힌 부분집합
  $Y$, $Z$, $T$에 대하여
  \[
    \operatorname{codim}(Y,T)=
    \operatorname{codim}(Y,Z)+\operatorname{codim}(Z,T).
    \tag{14.3.2.1}\label{0.14.3.2.1-fr}
  \]
\end{enumerate}
\end{proposition}

\begin{proof}
a)가 b)를 함의함은 즉시 알 수 있다. 거꾸로 b)가 성립한다고 하자.
양 끝점이 같은 두 포화 사슬의 길이가 $m$과 $n\leq m$이면 반드시
$m=n$임을 보이자. $n$에 관한 귀납법을 쓰면 $n=1$인 경우는 자명하다.
이제 $n>1$, $n<m$이라 하고,
$Z_0\subset Z_1\subset\cdots\subset Z_n$을 포화 사슬이라 하자.
\oldpage[0\textsubscript{IV-1}]{11}
양 끝점이 $Z_0$, $Z_n$이고 길이가 $m$인 다른 포화 사슬이 존재한다고
하자. $\operatorname{codim}(Z_0,Z_n)\geq m>n$이고
$\operatorname{codim}(Z_0,Z_1)=1$이므로, b)에 의하여
$\operatorname{codim}(Z_1,Z_n)=\operatorname{codim}(Z_0,Z_n)-1>n-1$이다.
이는 귀납가정에 모순이다.
\end{proof}

(14.3.2)의 조건들이 성립할 때 $X$가 \emph{사슬 조건}을 만족한다고,
또는 $X$가 \emph{카테너리 공간}이라고 한다. 카테너리 공간의 모든
닫힌 부분공간이 카테너리임은 분명하다.

\begin{proposition}[14.3.3]
\label{0.14.3.3-fr}
$X$를 유한 차원의 뇌터 콜모고로프 공간이라 하자. 다음 조건들은
동치이다.
\begin{enumerate}
  \item[\rm(a)] $X$의 기약 닫힌 부분집합들로 이루어진 두 극대 사슬은
  길이가 같다.
  \item[\rm(b)] $X$는 등차원이고 등여차원이며 카테너리이다.
  \item[\rm(c)] $X$는 등차원이고, $Y\subset Z$인 $X$의 두 기약 닫힌
  부분집합 $Y$, $Z$에 대하여
  \[
    \dim(Z)=\dim(Y)+\operatorname{codim}(Y,Z).
    \tag{14.3.3.1}\label{0.14.3.3.1-fr}
  \]
  \item[\rm(d)] $X$는 등여차원이고, $Y\subset Z$인 $X$의 두 기약
  닫힌 부분집합 $Y$, $Z$에 대하여
  \[
    \operatorname{codim}(Y,X)=
    \operatorname{codim}(Y,Z)+\operatorname{codim}(Z,X).
    \tag{14.3.3.2}\label{0.14.3.3.2-fr}
  \]
\end{enumerate}
\end{proposition}

\begin{proof}
$X$에 관한 가정에 의하여, $X$의 기약 닫힌 부분집합들로 이루어진 극대
사슬의 양 끝점은 반드시 닫힌점과 $X$의 기약성분이다
($0_{\mathrm I}$, 2.1.3). 또한 양 끝점이 $Y$, $Z$ ($Y\subset Z$)인
모든 포화 사슬은 극대 사슬에 포함되며, 이 극대 사슬에서 주어진 사슬에
속하지 않는 원소들은 $Y$에 포함되거나 $Z$를 포함한다. 이 관찰들로부터
a)와 b)의 동치가 즉시 따른다. 또한 a)가 성립하면 $X$의 모든 기약
닫힌 부분집합 $Y$에 대하여
\[
  \dim(Y)+\operatorname{codim}(Y,X)=\dim(X).
  \tag{14.3.3.3}\label{0.14.3.3.3-fr}
\]
공식 (14.3.2.1)이 성립하므로, (14.3.3.3)으로부터 (14.3.3.1)과
(14.3.3.2)를 즉시 얻는다. 거꾸로 (14.3.3.1)은 (14.3.2.1)을
함의하므로, (14.3.2)에 의하여 사슬 조건을 함의한다. 또한 $Y$가
$X$의 닫힌점 $x$ 하나로 이루어지고 $Z$가 $X$의 기약성분인 경우에
(14.3.3.1)을 적용하면
$\operatorname{codim}(\{x\},X)=\dim(Z)$를 얻는다. 따라서 c)는 b)를
함의한다. 마찬가지로 (14.3.3.2)는 (14.3.2.1)을, 따라서 사슬 조건을
함의한다. 위와 똑같이 $Y$와 $Z$를 택하면 (14.3.3.2)에서도
$\operatorname{codim}(\{x\},X)=\dim(Z)$를 얻는다. 그러므로 $X$의
모든 기약성분이 닫힌점을 포함한다는 사실 ($0_{\mathrm I}$, 2.1.3)에
의하여 d)는 b)를 함의한다.
\end{proof}

뇌터 콜모고로프 공간이 \emph{유한 차원}이고 (14.3.3)의 동치인 조건들을
만족하면 이 공간을 \emph{쌍등차원}이라고 한다.

\begin{corollary}[14.3.4]
\label{0.14.3.4-fr}
$X$를 쌍등차원인 뇌터 콜모고로프 공간이라 하자. 그러면 $X$의 모든
닫힌점 $x$와 모든 기약성분 $Z$에 대하여
\[
  \dim(X)=\dim(Z)=\operatorname{codim}(\{x\},X)=\dim_x(X).
  \tag{14.3.4.1}\label{0.14.3.4.1-fr}
\]
\end{corollary}

\oldpage[0\textsubscript{IV-1}]{12}
\begin{proof}
마지막 등식은 다음 사실에서 따른다. $Y_0=\{x\}\subset Y_1\subset
\cdots\subset Y_m$이 $X$의 기약 닫힌 부분집합들로 이루어진 극대
사슬이고 $U$가 $x$의 열린 근방이면, $U\cap Y_i$들은 $U$의 서로 다른
기약 닫힌 부분집합들이다 (실제로 $\overline{U\cap Y_i}=Y_i$이다).
따라서 (14.1.4)에 의하여 $\dim(U)=\dim(X)$이다.
\end{proof}

\begin{corollary}[14.3.5]
\label{0.14.3.5-fr}
$X$를 뇌터 콜모고로프 공간이라 하자. $X$가 쌍등차원이면 $X$의
기약성분들의 임의의 합집합과 $X$의 모든 기약 닫힌 부분집합도
쌍등차원이다. 또한 $X$의 모든 닫힌 부분집합 $Y$에 대하여
\[
  \dim(Y)+\operatorname{codim}(Y,X)=\dim(X).
  \tag{14.3.5.1}\label{0.14.3.5.1-fr}
\]
\end{corollary}

\begin{proof}
$X$의 기약 닫힌 부분집합들로 이루어진 모든 사슬은 $X$의 어느
기약성분에 포함되므로, 첫 번째 주장은 (14.3.3)에서 즉시 따른다. 또한
$X'$가 $X$의 기약 닫힌 부분집합이면 $X'$는 조건 (14.3.3, c))를
자명하게 만족하므로 두 번째 주장을 얻는다.

마지막으로 (14.3.5.1)을 증명하자. (14.3.3)의 증명에서 $Y$가 기약일
때 이 관계가 성립함을 보았다. $Y_i$ ($1\leq i\leq m$)들이 $Y$의
기약성분들이면, $\dim(Y_i)$가 가장 큰 $Y_i$는
$\operatorname{codim}(Y_i,X)$가 가장 작은 $Y_i$이기도 하다. 따라서
(14.3.5.1)은 $\dim(Y)$와 $\operatorname{codim}(Y,X)$의 정의에서
따른다.
\end{proof}

\begin{remark}[14.3.6]
\label{0.14.3.6-fr}
(14.3.2)의 증명에는 위상공간의 기약 닫힌 부분집합이라는 사실이 전혀
쓰이지 않았으므로, 그 증명은 임의의 순서집합에도 적용된다. (14.3.3)의
증명도 다음 성질을 갖는 순서집합 $E$에서 마찬가지로 적용된다. 모든
$x\in E$에 대하여 $E$의 \emph{극소} 원소인 $z\leq x$가 존재하고,
$E$의 원소들로 이루어진 사슬들의 길이가 유계이다.
\end{remark}

\section{\texorpdfstring{$M$}{M}-정칙열과
\texorpdfstring{$\mathscr F$}{F}-정칙열}
\label{section:0.15-fr}

\subsection{\texorpdfstring{$M$}{M}-정칙열과
\texorpdfstring{$M$}{M}-준정칙열}
\label{subsection:0.15.1-fr}
\label{0.15.1-fr}

\begin{env}[15.1.1]
\label{0.15.1.1-fr}
\oldpage[0\textsubscript{IV-1}]{12}
$A$를 환, $N$을 $A$-가군이라 하자. 이하에서는 간단히
$N[T_1,\ldots,T_n]$으로 등급 $A$-가군
$N\otimes_A A[T_1,\ldots,T_n]$ ($T_i$들은 부정원)을 나타내며, 여기서
$N$에는 자명한 등급을, $A[T_1,\ldots,T_n]$에는 전체 차수에 따른 등급을
준다.

$M$을 $A$-가군, $f_i$ ($1\leq i\leq n$)를 $A$의 원소들이라 하고,
$\mathfrak{J}=\sum_{i=1}^n f_iA$를 $f_i$들이 생성하는 $A$의
아이디얼이라 하자. $M$에 $\mathfrak J$-준아딕 여과를 주면, 다음과
같이 차수 $0$인 \emph{등급 $A$-가군의 전사 등급 준동형}
\[
  \varphi : (\gr_0(M))[T_1,\ldots,T_n]\longrightarrow\gr_\bullet(M)
  \tag{15.1.1.1}\label{0.15.1.1.1-fr}
\]
을 정의한다. $\xi\in\gr_0(M)$가 $x\in M$의
$\mathfrak{J}M$을 법으로 한 동치류이고
$P[T_1,\ldots,T_n]$이 차수 $k$인 동차 다항식이면, $\xi$와 $P$의
쌍에 $P(f_1,\ldots,f_n)x$의 $\mathfrak{J}^{k+1}M$을 법으로 한
동치류를 대응시킨다. 이 동치류가
\oldpage[0\textsubscript{IV-1}]{13}
$\xi$와 $P$에만 의존함은 즉시 알 수 있다. 또한 $S_k$가
$A[T_1,\ldots,T_n]$에서 차수 $k$인 동차 다항식들의 집합이면, 이로써
$A$-선형사상
\[
  \varphi_k:\gr_0(M)\otimes_A S_k\longrightarrow\gr_k(M)
\]
이 정의되며, 이는 $\mathfrak{J}^kM$의 정의에 따라 전사이다.

이하에서 $f\in A$가 $A$-가군 $M$에서 \emph{영인자가 아니다}라는 것은
$M$의 스칼라곱 자기준동형 $f_M:x\mapsto fx$가 \emph{단사}라는 뜻으로
사용한다.
\end{env}

\begin{lemma}[15.1.2]
\label{0.15.1.2-fr}
$M$을 $A$-가군, $f$를 $A$의 원소라 하고, $M$에 $(fA)$-준아딕
여과를 주자. $f$가 $M$에서 영인자가 아니면 (15.1.1)에서 정의한 표준
준동형
\[
  \varphi:(M/fM)[T]\longrightarrow\gr_\bullet(M)
  \tag{15.1.2.1}\label{0.15.1.2.1-fr}
\]
은 전단사이다. $M$이 $(fA)$-준아딕 위상에 관하여 분리이면 그 역도
성립한다.
\end{lemma}

\begin{proof}
여기서는 $\gr_k(M)=f^kM/f^{k+1}M$이다. $\varphi$가 단사라는 것은 모든
$x\in M$과 모든 $k$에 대하여 $f^kx\in f^{k+1}M$이면 $x\in fM$이라는
뜻이다. $f$가 $M$에서 영인자가 아니면 이 성질이 성립함은 명백하다.
거꾸로 스칼라곱 자기준동형 $f_M:x\mapsto fx$는 몫으로 이행하여
준동형 $\gr_k(M)\to\gr_{k+1}(M)$들을 정의하고, 따라서
$\gr_\bullet(M)$에서 자기 자신으로 가는 차수 $1$의 등급 준동형 $g$를
정의한다. 그런데 $\varphi$가 단사라는 것은 $g$가 단사라는 뜻이기도
하다. $M$의 여과가 분리라고 가정하면, 이로부터 $f_M$이 단사임이
따른다 (Bourbaki, \emph{Alg.~comm.}, chap.~III, \S~2, no~8, cor.~1 du
th.~1).
\end{proof}

\begin{corollary}[15.1.3]
\label{0.15.1.3-fr}
$A$를 뇌터 환, $M$을 유한형 $A$-가군이라 하자. $f$가 $M$에서
영인자가 아니기 위한 필요충분조건은 $\varphi$가 전단사인 것이다.
\end{corollary}

\begin{proof}
조건이 충분함만 증명하면 된다. $f_M$이 단사임을 보이려면 $A$의 모든
소아이디얼 $\mathfrak p$에 대하여 $A_{\mathfrak p}$-가군
$M_{\mathfrak p}=M'$에서 $f'=f/1\in A'=A_{\mathfrak p}$에 의한
스칼라곱 자기준동형이 단사임을 보이면 충분하다 (Bourbaki,
\emph{Alg.~comm.}, chap.~II, \S~3, no~4, th.~1). 그런데
$f'^kM'/f'^{k+1}M'=(f^kM/f^{k+1}M)\otimes_A A'$이고, 표준 준동형
$\varphi':(M'/f'M')[T]\to\gr_\bullet(M')$은
$\varphi\otimes 1_{A'}$와 같음이 즉시 따른다. 따라서 $\varphi$가
단사이면 $\varphi'$도 단사이다 ($0_{\mathrm I}$, 1.3.2). 그러므로
$A$가 \emph{국소} 뇌터 환인 경우로 환원되며, 그 극대 아이디얼을
$\mathfrak m$이라 하자. $f\notin\mathfrak m$이면 $f$는 가역원이므로
증명할 것이 없다. $f\in\mathfrak m$이면 $M$은 $(fA)$-준아딕 위상에
관하여 분리이고 ($0_{\mathrm I}$, 7.3.5), 따라서 (15.1.2)를 적용하면
된다.
\end{proof}

\begin{definition}[15.1.4]
\label{0.15.1.4-fr}
$A$를 환, $M$을 $A$-가군이라 하자. $f\in A$가 $M$에서 영인자가
아니면 $f$가 \emph{$M$-정칙}이라고 하고, (15.1.2.1)의 준동형이
전단사이면 \emph{$M$-준정칙}이라고 한다.
\end{definition}

따라서 방금 모든 $M$-정칙 원소는 $M$-준정칙이고, $A$가 뇌터 환이며
$M$이 유한형 $A$-가군이면 그 역도 성립함을 증명하였다. $M$이 유한형이라는
가정을 버리면 이 역은 성립하지 않을 수 있다. 예를 들어 $A$가 체가 아닌
주 아이디얼환이고 $K$가 그 분수체이며 $M$이 $A$-가군 $K/A$이면,
$A$의 모든 $f\neq0$에 대하여 $fM=M$이다. 따라서 (15.1.2.1)의 양변은
모두 $0$이지만, 이때 $f$는 $M$에서 영인자이다.

\begin{env}[15.1.5]
\label{0.15.1.5-fr}
\oldpage[0\textsubscript{IV-1}]{14}
이제 $M$에 $M_0=M$인 \emph{여과} $(M_k)$가 주어졌다고 하고,
$\gr_\bullet(M)=\bigoplus_{k\geq0}M_k/M_{k+1}$을 그 연관 등급
$A$-가군이라 하자. (15.1.1)의 표기를 유지하여 모든
$k\geq0$에 대하여
\[
  M'_k=M_k+\mathfrak{J}M_{k-1}+\cdots+
  \mathfrak{J}^{k-1}M_1+\mathfrak{J}^kM_0
  \tag{15.1.5.1}\label{0.15.1.5.1-fr}
\]
로 놓는다. $(M'_k)$가 $M$의 여과임은 명백하다. $\mathfrak Q$가
$A$의 아이디얼이고 $M_k=\mathfrak{Q}^kM$이면, 여과 $(M'_k)$는 바로
$(\mathfrak{J}+\mathfrak{Q})$-준아딕 여과이다. $(M'_k)$의 연관 등급
$A$-가군을 $\gr'_\bullet(M)$으로 나타내자. 또한
\[
  (\gr_\bullet(M)\otimes_A(A/\mathfrak{J}))[T_1,\ldots,T_n]
\]
이라는 기호는 등급 $A$-가군들의 텐서곱 ($\mathrm{II}$, 2.1.2)
\[
  (\gr_\bullet(M)\otimes_A(A/\mathfrak{J}))
  \otimes_A A[T_1,\ldots,T_n]
\]
을 뜻하는 것으로 한다. 다음과 같이 차수 $0$인
\emph{등급 $A$-가군의 전사 등급 준동형}
\[
  \psi:(\gr_\bullet(M)\otimes_A(A/\mathfrak{J}))[T_1,\ldots,T_n]
  \longrightarrow\gr'_\bullet(M)
  \tag{15.1.5.2}\label{0.15.1.5.2-fr}
\]
을 정의한다. $\alpha\in A/\mathfrak{J}$가 $a\in A$의 동치류,
$\xi\in M_k/M_{k+1}$가 $x\in M_k$의 동치류이고
$P[T_1,\ldots,T_n]$이 차수 $h$인 동차 다항식이면, 삼중항
$(\alpha,\xi,P)$에 $P(f_1,\ldots,f_n)ax$의 $M'_{h+k+1}$을 법으로 한
동치류를 대응시킨다. 이 원소가 $\xi$, $\alpha$, $P$에만 의존함도
즉시 확인되며, 이로써 등급 준동형 (15.1.5.2)가 정의된다. 이는
(15.1.5.1)의 정의에 따라 전사이다. $M_1=0$인 $M$의 여과를 택하면
준동형 (15.1.1.1)을 다시 얻는다.
\end{env}

\begin{lemma}[15.1.6]
\label{0.15.1.6-fr}
(15.1.5)의 표기를 유지하고 $n=1$, $f=f_1$이라 하자. $f$가
$\gr_\bullet(M)$-정칙이면 준동형 (15.1.5.2)는 전단사이다.
\end{lemma}

\begin{proof}
(15.1.5.2)의 왼쪽 변과 오른쪽 변에서 차수 $k$인 항들의 $A$-부분가군을
각각 $Q_k$와 $Q'_k$라 하자. $Q_k$에 다음 여과를 준다.
\[
  (Q_k)_i=\sum_{j\leq k-i}
  (\gr_{k-j}(M)\otimes_A(A/fA))T^j
  \tag{15.1.6.1}\label{0.15.1.6.1-fr}
\]
그러면 $(Q_k)_0=Q_k$, $(Q_k)_{k+1}=0$이고, 이 여과에 대하여
\[
  \gr_i(Q_k)=(\gr_i(M)\otimes_A(A/fA))T^{k-i}.
\]
$Q'_k$에는 $\psi((Q_k)_i)=(Q'_k)_i$로 이루어진 여과를 준다. 이
여과들은 \emph{유한}하므로, $\psi_k$에서 유도되는 준동형
$\psi_{ki}:\gr_i(Q_k)\to\gr_i(Q'_k)$이 단사임을 증명하면 충분하다
(Bourbaki, \emph{Alg.~comm.}, chap.~III, \S~3, no~8, cor.~1 du th.~1).
그런데
\[
  \gr_i(Q_k)=((M_i/M_{i+1})\otimes_A(A/fA))T^{k-i}
\]
이다. 한편 $(Q'_k)_{i+1}$은
\[
  M_k+fM_{k-1}+\cdots+f^{k-i-1}M_{i+1}
\]
의 $M'_k/M'_{k+1}$에서의 상이다. 따라서 $x\in M_i$에 대하여 관계
\oldpage[0\textsubscript{IV-1}]{15}
\[
  f^{k-i}x\in M_k+fM_{k-1}+\cdots+
  f^{k-i-1}M_{i+1}+M'_{k+1}
  \tag{15.1.6.2}\label{0.15.1.6.2-fr}
\]
가 $x\in fM_i+M_{i+1}$을 함의함을 보이는 문제로 환원된다. 그런데
(15.1.6.2)의 오른쪽 변은 $M_{i+1}+M'_{k+1}$에 포함되고,
(15.1.5.1)에 따라 $M'_{k+1}$은 $M_{i+1}+f^{k+1-i}M$에 포함된다.
$f$에 대한 가정으로부터 모든 $h$에 대하여 $f$는 $M/M_h$에서
영인자가 아니다 (Bourbaki, \emph{loc.~cit.}). 따라서
$f^{k-i}x\in f^{k+1-i}M+M_{i+1}$이라는 관계로부터
($M/M_{i+1}$에서 생각하여) $x-fy\in M_{i+1}$인 $y\in M$이 존재한다.
또한 $x\in M_i$이므로 $fy\in M_i$이고, 이제 $M/M_i$에서 생각하면
$y\in M_i$이다. 마침내 $x\in fM_i+M_{i+1}$을 얻는다. C.Q.F.D.
\end{proof}

\begin{definition}[15.1.7]
\label{0.15.1.7-fr}
$(f_i)_{1\leq i\leq n}$을 $A$의 원소들로 이루어진 열이라 하자. 모든
$1\leq i\leq n$에 대하여 $f_i$가
$M/(\sum_{1\leq j\leq i-1}f_jM)$-정칙이면 이 열을
\emph{$M$-정칙열}이라 한다. (15.1.1.1)에서 정의한 표준 준동형
$\varphi$가 전단사이면 열 $(f_i)$를 \emph{$M$-준정칙열}이라 한다.

$M=A$일 때에는 단순히 “정칙열”, “준정칙열”이라고 한다.
\end{definition}

\begin{proposition}[15.1.8]
\label{0.15.1.8-fr}
(15.1.5)의 표기를 유지하고, 열 $(f_i)_{1\leq i\leq n}$이
$\gr_\bullet(M)$-정칙이라고 가정하자. 그러면 표준 준동형 (15.1.5.2)는
전단사이다.
\end{proposition}

\begin{proof}
$n=1$일 때 이 명제는 바로 (15.1.6)이므로, $n$에 관하여 귀납하자.
다음과 같이 놓는다.
\[
  \mathfrak{J}''=\sum_{i=1}^{n-1}f_iA.
\]
그러면 $\mathfrak{J}=\mathfrak{J}''+f_nA$이다. 또한
\[
  M''_k=M_k+\mathfrak{J}''M_{k-1}+\cdots+
  \mathfrak{J}''^{k-1}M_1+\mathfrak{J}''^kM_0
\]
로 놓으면
\[
  M'_k=M''_k+f_nM''_{k-1}+\cdots+f_n^{k-1}M''_1+f_n^kM''_0
\]
로 쓸 수 있다. $M$의 여과 $(M''_k)$에 연관된 등급 $A$-가군을
$\gr''_\bullet(M)$으로 나타내자. 표준 동형들을 통하여
\[
  (\gr_\bullet(M)\otimes(A/\mathfrak{J}))\otimes_A A[T_1,\ldots,T_n]
  =
  (\gr_\bullet(M)\otimes(A/\mathfrak{J}''))[T_1,\ldots,T_{n-1}]
  \otimes(A/f_nA)\otimes_A A[T_n]
\]
로 쓸 수 있다. 따라서 (15.1.5.2)는
\[
  (\gr''_\bullet(M)\otimes(A/f_nA))[T_n]\longrightarrow\gr'_\bullet(M)
  \tag{15.1.8.1}\label{0.15.1.8.1-fr}
\]
과
\[
  ((\gr_\bullet(M)\otimes(A/\mathfrak{J}''))[T_1,\ldots,T_{n-1}])
  \otimes(A/f_nA)\otimes_A A[T_n]
  \xrightarrow{\psi'\otimes1}
  \gr''_\bullet(M)\otimes(A/f_nA)\otimes_A A[T_n]
\]
으로 분해된다. 여기서 $\psi'$은 $n$을 $n-1$로, $\mathfrak{J}$를
$\mathfrak{J}''$로, $\gr'_\bullet(M)$을 $\gr''_\bullet(M)$으로 바꾼
(15.1.5.2)의 사상이다. 귀납 가정에 따라 $\psi'$은 전단사이므로,
(15.1.8.1)도 전단사임을 보이는 일만 남는다. (15.1.6)을 적용하면
$f_n$이 $\gr''_\bullet(M)$-정칙임을 보이면 충분하다. 그런데
$\gr''_\bullet(M)$은
\[
  (\gr_\bullet(M)\otimes(A/\mathfrak{J}''))
  [T_1,\ldots,T_{n-1}]
\]
과 동일시된다. 가정에 따라
$\gr_\bullet(M)\otimes(A/\mathfrak{J}'')$에서 $f_n$에 의한
스칼라곱 자기준동형은 단사이다. $A[T_1,\ldots,T_{n-1}]$이 자유
$A$-가군이므로 $\gr''_\bullet(M)$에서 $f_n$에 의한 스칼라곱
자기준동형도 단사이다.
\end{proof}

\begin{proposition}[15.1.9]
\label{0.15.1.9-fr}
$M$을 $A$-가군, $(f_i)_{1\leq i\leq n}$을 $A$의 원소들로 이루어진
열이라 하자. 열 $(f_i)$가 $M$-정칙이면 $M$-준정칙이다. 더 나아가
$A$-가군들
\[
  M,\quad M/f_1M,\quad\ldots,\quad
  M/\left(\sum_{1\leq j\leq n-1}f_jM\right)
\]
이 $\mathfrak{J}$-준아딕 위상에 관하여 분리라고 가정하면 그 역도
성립한다. 여기서 $\mathfrak{J}=\sum_{i=1}^n f_iA$이다.
\end{proposition}

\begin{proof}
첫 번째 주장은 (15.1.8)에서 $M_1=0$으로 놓으면 따른다. 실제로 이때
준동형 (15.1.5.2)는 (15.1.1.1)로 환원된다. 역으로 (15.1.1.1)의
준동형 $\varphi$가 전단사라고 가정하자. 다음과 같이 놓고
\[
  \mathfrak{J}''=\sum_{i=1}^{n-1}f_iA,
\]
\oldpage[0\textsubscript{IV-1}]{16}
$M$에 $\mathfrak{J}''$-준아딕 여과를 주었을 때의 연관 등급
$A$-가군을 $\gr''_\bullet(M)$으로 나타내면, $\varphi$는
\[
  (\gr''_\bullet(M)\otimes(A/f_nA))[T_n]\longrightarrow\gr_\bullet(M)
  \tag{15.1.9.1}\label{0.15.1.9.1-fr}
\]
과 $\varphi'\otimes1$로 분해된다. 여기서 $\varphi'$은 표준 준동형
\[
  (\gr''_0(M))[T_1,\ldots,T_{n-1}]
  \longrightarrow\gr''_\bullet(M)
\]
이다. 이 두 준동형은 전사이므로, $\varphi$가 전단사라는 가정 아래
둘 다 전단사이다. 따라서 $\varphi'$은 단사이며, 이미 전사이므로
전단사이다. 이제 $n$에 관하여 귀납할 수 있다 ($n=1$인 경우는
(15.1.2)에서 따른다). 실제로
\[
  M,\quad M/f_1M,\quad\ldots,\quad
  M/\left(\sum_{j=1}^{n-2}f_jM\right)
\]
위의 $\mathfrak{J}''$-준아딕 위상은 $\mathfrak{J}$-준아딕 위상보다
더 섬세하므로 분리이다. 따라서 열 $(f_i)_{1\leq i\leq n-1}$이
$M$-정칙임을 알 수 있다. 한편 다음을 보이자.
가정에 따라 $y\in M/\mathfrak{J}''M$이고
\[
  f_n^ky\in f_n^{k+1}(M/\mathfrak{J}''M)
\]
이면 $y\in f_n(M/\mathfrak{J}''M)$이다. 실제로 $x$를 $y$의 한
대표원이라 하자. $f_n^kx\in\mathfrak{J}^{k+1}M$임을 보이면,
(15.1.9.1)이 단사라는 사실로부터 우리의 주장이 따른다. 가정에 따라
\[
  f_n^kx\in f_n^{k+1}M+\mathfrak{J}''M\subset\mathfrak{J}M
\]
이다. 따라서 $\varphi$가 단사이므로
$f_n^{k-1}x\in\mathfrak{J}^kM$이고, 귀납하면
$x\in\mathfrak{J}M$을 얻으며, 마침내
$f_n^kx\in\mathfrak{J}^{k+1}M$이다. (15.1.2)와 같이 논증을
마무리한다. $M/\mathfrak{J}''M$ 위에서 $\mathfrak{J}$-준아딕 여과는
$(f_nA)$-준아딕 여과와 같고, 가정에 따라 이 여과는 분리이다. 따라서
$M/\mathfrak{J}''M$에서 $f_n$에 의한 스칼라곱 자기준동형은 단사이다.
이로써 열 $(f_i)_{1\leq i\leq n}$이 $M$-정칙임이 증명되었다.
\end{proof}

\begin{corollary}[15.1.10]
\label{0.15.1.10-fr}
(15.1.9)의 표기를 유지하자. 모든 유한 생성 $A$-가군 $N$에 대하여
$M\otimes_A N$이 $\mathfrak{J}$-준아딕 위상에 관하여 분리라고
가정하자. 그러면 열 $(f_i)$가 $M$-정칙일 필요충분조건은 이 열이
$M$-준정칙인 것이다.
\end{corollary}

(15.1.10)의 가정으로부터 다음 사실이 따름을 지적해 두자. 열
$(f_i)_{1\leq i\leq n}$이 $M$-정칙이면, 구간 $[1,n]$의 임의의
순열 $\sigma$에 대하여 열 $(f_{\sigma(i)})$도 $M$-정칙이다.

\begin{corollary}[15.1.11]
\label{0.15.1.11-fr}
$A$를 뇌터 환, $M$을 유한 생성 $A$-가군,
$(f_i)_{1\leq i\leq n}$을 $A$의 근기에 포함되는 원소들의 유한열이라
하자. 그러면 열 $(f_i)$가 $M$-정칙일 필요충분조건은 이 열이
$M$-준정칙인 것이다.
\end{corollary}

\begin{proof}
이는 (15.1.10)의 특수한 경우이다. 실제로 이 경우 모든 유한 생성
$A$-가군은 $\mathfrak{J}$-준아딕 위상에 관하여 분리이다
($0_{\mathrm I}$, 7.3.5).
\end{proof}

\begin{remarks}[15.1.12]
\label{0.15.1.12-fr}
\mbox{}
\begin{enumerate}
  \item[(i)] $A$가 뇌터 국소환이고 $M$이 유한 생성 $A$-가군이며
  $f_i$들이 $A$의 극대 아이디얼에 포함될 때, $M$-정칙열의 개념은
  J.-P.~Serre가 $M$-\emph{열}이라는 이름으로 도입하였다 [17].
  \item[(ii)] $A$가 뇌터 환이고 $M=A$인 경우에도, $A$의 두 원소로
  이루어진 준정칙열 $(f,g)$가 반드시 정칙열인 것은 아니다. 예를 들어
  열 $(0,1)$은 준정칙이다. 실제로 이때
  $\mathfrak{J}=fA+gA=A$이므로 $\gr_\bullet(A)=0$이다. 그러나
  $A\neq0$이면 이 열은 정칙열이 아니다. 한편 이 경우 열 $(1,0)$은
  정칙열임을 지적해 두자. $1$은 $A$에서 영인자가 아니고, $A/A=0$이므로
  $0$도 이 가군에서 영인자가 아니다 ((15.1.1) 참조).
  \item[(iii)]
  \oldpage[0\textsubscript{IV-1}]{17}
  (15.1.6)에서 $f$가 $\gr_\bullet(M)$-정칙이라는 가정을 일반적으로
  $f$가 $\gr_\bullet(M)$-준정칙이라는 가정으로 바꿀 수 없다. 이를
  보기 위하여 (15.1.4)의 예로 돌아가자. 여기서 $A$는 분수체
  $K\neq A$를 갖는 주 아이디얼환이다. $M=K$로 놓고, $M$의 여과
  $(M_k)$를 $M_0=M=K$, $M_1=A$, $M_k=0$ ($k\geq2$)로 정의하면
  $\gr_0(M)=K/A$, $\gr_1(M)=A$, $\gr_k(M)=0$ ($k\geq2$)이다.
  $A$에서 $f\neq0$이면 앞에서 본 바와 같이 $f$는
  $\gr_\bullet(M)$-정칙이 아니다. 그러나
  \[
    f^k\gr_\bullet(M)/f^{k+1}\gr_\bullet(M)=f^kA/f^{k+1}A
  \]
  이므로, (15.1.2.1)에서 $M$을 $\gr_\bullet(M)$으로 바꾸어 얻는
  준동형이 단사임은 명백하다. 따라서 $f$는
  $\gr_\bullet(M)$-준정칙이다. 그러나 (15.1.6)의 표기에 따르면 이때
  모든 $k\geq0$에 대하여 $M'_k=K$이므로 $\gr'_\bullet(M)=0$인 반면,
  (15.1.5.2)의 왼쪽 변은 $0$이 아니다.
\end{enumerate}
\end{remarks}


\begin{proposition}[15.1.13]
\label{0.15.1.13-fr}
$A$를 환, $M$, $N$을 두 $A$-가군, $(f_i)_{1\leq i\leq n}$을 $A$의
원소들로 이루어진 열이라 하자. 열 $(f_i)$가 $M$-정칙이고 $N$이 평탄
$A$-가군이면, 열 $(f_i)$는 $M\otimes_A N$-정칙이다. $N$이 충실 평탄
$A$-가군이면 그 역도 성립한다.
\end{proposition}

\begin{proof}
실제로 ($N$에 아무 가정도 하지 않아도)
\[
  (M\otimes_A N)/
  \left(\sum_{j=1}^{i-1}f_j(M\otimes_A N)\right)
\]
은
\[
  \left(M/\left(\sum_{j=1}^{i-1}f_jM\right)\right)\otimes_A N
\]
과 동형이다. 이제 $M/(\sum_{j=1}^{i-1}f_jM)$에서 $f_i$에 의한
스칼라곱 자기준동형에 ($0_{\mathrm I}$, 6.1.1)과
($0_{\mathrm I}$, 6.4.1)을 적용하면 충분하다.
\end{proof}

\begin{corollary}[15.1.14]
\label{0.15.1.14-fr}
$A$를 환, $M$을 $A$-가군, $(f_i)_{1\leq i\leq n}$을 $A$의 원소들로
이루어진 유한열이라 하자. $\rho:A\to B$를 환 준동형이라 하고
$M'=M_{(B)}=M\otimes_A B$, $f'_i=\rho(f_i)$
($1\leq i\leq n$)로 놓자. $B$가 평탄 $A$-가군이고 열 $(f_i)$가
$M$-정칙이면 열 $(f'_i)$는 $M'$-정칙이다. $B$가 충실 평탄
$A$-가군이면 그 역도 성립한다.
\end{corollary}

\begin{proof}
실제로
\[
  M'/\left(\sum_{j=1}^{i-1}f'_jM'\right)
\]
은
\[
  \left(M/\left(\sum_{j=1}^{i-1}f_jM\right)\right)\otimes_A B
\]
와 동일시된다. 또한 $B$-가군
$M'/(\sum_{j=1}^{i-1}f'_jM')$에서 $f'_i$에 의한 스칼라곱
자기준동형은 $A$-가군
\[
  \left(M/\left(\sum_{j=1}^{i-1}f_jM\right)\right)\otimes_A B
\]
에서 $f_i$에 의한 스칼라곱 자기준동형과 동일시된다. 따라서 주장은
(15.1.13)에서 따른다.
\end{proof}

\begin{proposition}[15.1.15]
\label{0.15.1.15-fr}
$A$를 환, $B$를 가환 $A$-대수, $(f_i)_{1\leq i\leq n}$을 $B$의
원소들로 이루어진 유한열, $M$을 $B$-가군이라 하자. 열 $(f_i)$가
$M$-정칙이고 $A$-가군
\[
  M/\left(\sum_{j=1}^{i-1}f_jM\right)\quad(1\leq i\leq n+1)
\]
각각이 평탄이라고 가정하자. 그러면 임의의 환 준동형
$\rho:A\to A'$에 대하여
$B'=B\otimes_A A'$, $M'=M\otimes_A A'$, $f'_i=f_i\otimes1$
($1\leq i\leq n$)로 놓으면, 열 $(f'_i)$는 $M'$-정칙이고
$A'$-가군
\[
  M'/\left(\sum_{j=1}^{i-1}f'_jM'\right)\quad(1\leq i\leq n+1)
\]
각각은 평탄이다.
\end{proposition}

\begin{proof}
가정에 따라 완전열인 다음 열을 생각하자.
\[
  0\longrightarrow M\xrightarrow{f_1}M\longrightarrow M/f_1M
  \longrightarrow0
\]
여기서 화살표 $M\xrightarrow{f_1}M$은 $M$에서 $f_1$에 의한 스칼라곱
자기준동형이다. $M/f_1M$이 평탄 $A$-가군이라는 가정으로부터
\[
  0\longrightarrow M\otimes_A A'
  \xrightarrow{f_1\otimes1}M\otimes_A A'
  \longrightarrow(M/f_1M)\otimes_A A'\longrightarrow0
\]
\oldpage[0\textsubscript{IV-1}]{18}
이 완전열임이 따른다 ($0_{\mathrm I}$, 6.1.2). 따라서 $M'$에서
$f'_1$에 의한 스칼라곱 자기준동형은 단사이다.

이제 $1\leq i\leq n-1$에 대하여
\[
  M_i=M/\left(\sum_{j=1}^i f_jM\right),\qquad
  M'_i=M'/\left(\sum_{j=1}^i f'_jM'\right)
\]
로 놓자. 그러면 $M'_i=M_i\otimes_A A'$,
$M_i/f_{i+1}M_i=M_{i+1}$,
$M'_i/f'_{i+1}M'_i=M'_{i+1}$이다. 앞의 논증에서 $M$과 $f_1$을
$M_i$와 $f_{i+1}$로 바꾸면, $M_i/f_{i+1}M_i$의 평탄성에 따라
$M'_i$에서 $f'_{i+1}$에 의한 스칼라곱 자기준동형이 단사임을 얻는다.
마지막 주장은 ($0_{\mathrm I}$, 6.2.1)에서 따른다.
\end{proof}

\begin{proposition}[15.1.16]
\label{0.15.1.16-fr}
$A$, $B$를 두 뇌터 국소환, $k$를 $A$의 잉여체,
$\varphi:A\to B$를 국소 준동형, $M$을 유한 생성 $B$-가군이라 하자.
$(f_i)_{1\leq i\leq n}$을 $B$의 극대 아이디얼의 원소들로 이루어진
열이라 하고, 각 $i$에 대하여 $g_i$를 $B\otimes_A k$에서 $f_i$의
상이라 하자. 다음 조건들은 서로 동치이다.
\begin{enumerate}
  \item[\rm(a)] 열 $(f_i)$가 $M$-정칙이고, $1\leq i\leq n$에 대하여
  몫가군
  \[
    M_i=M/\left(\sum_{j=1}^i f_jM\right)
  \]
  는 평탄 $A$-가군이다.
  \item[\rm(b)] 열 $(f_i)$가 $M$-정칙이고 $A$-가군 $M_n$은 평탄이다.
  \item[\rm(c)] $A$-가군 $M$은 평탄이고, 열
  $(g_i)_{1\leq i\leq n}$은 $(M\otimes_A k)$-정칙이다.
  \item[\rm(d)] $A$-가군 $M$은 평탄이고, 임의의 환 준동형
  $\rho:A\to A'$에 대하여 $B'=B\otimes_A A'$,
  $M'=M\otimes_A A'$, $f'_i=f_i\otimes1$ ($1\leq i\leq n$)로
  놓으면 열 $(f'_i)$는 $M'$-정칙이다.
\end{enumerate}
\end{proposition}

\begin{proof}
a)가 b)를 함의함은 명백하다. b)가 d)를 함의함을 보이기 위해서는
(15.1.15)에 따라 b)가 $0\leq i\leq n-1$에 대하여 $M_i$가 평탄
$A$-가군임을 함의한다는 것을 보이면 충분하다 (여기서 $M_0=M$으로
놓는다). $M_{i+1}=M_i/f_{i+1}M_i$이므로, 하강 귀납법에 따라 다음을
증명하는 문제로 환원된다. $h$가 $B$의 극대 아이디얼의 원소이고
$M$에서 $h$에 의한 스칼라곱 자기준동형이 단사이며 $M/hM$이 평탄
$A$-가군이면 $M$도 평탄 $A$-가군이다. 그러나 이는 바로
($0_{\mathrm{III}}$, 10.2.7)이다. $A'=k$로 놓으면 c)는 d)의 특수한
경우임이 명백하다. 이제 c)가 a)를 함의함을 보이는 일만 남았다.
$M$은 평탄 $A$-가군이자 유한 생성 $B$-가군이므로, c)와
($0_{\mathrm{III}}$, 10.2.4)에서 $M$에서 $f_1$에 의한 스칼라곱
자기준동형이 단사이고 그 여핵 $M/f_1M=M_1$이 평탄 $A$-가군임이
따른다. $i$에 관하여 귀납하고 $M_i$가 평탄 $A$-가군임이 증명되었다고
가정하자. 가정에 따라 $M_i\otimes_A k$에서 $g_{i+1}$에 의한
스칼라곱 자기준동형이 단사이므로, 다시 ($0_{\mathrm{III}}$, 10.2.4)에
따라 $M_i$에서 $f_{i+1}$에 의한 스칼라곱 자기준동형이 단사이고
$M_i/f_{i+1}M_i=M_{i+1}$이 평탄 $A$-가군이다. 이로써 증명이
끝난다. c)가 a)를 함의한다는 사실의 증명에서는 $f_i$들이 $B$의
극대 아이디얼에 속한다는 사실을 사용하지 않았음을 지적해 두자.
\end{proof}

\begin{corollary}[15.1.17]
\label{0.15.1.17-fr}
$A$와 $B$에 관한 가정은 (15.1.16)과 같다고 하자. $M$을 $A$ 위에서
평탄인 유한 생성 $B$-가군, $R$을 $M$의 $B$-부분가군이라 하고,
$N=M/R$이 평탄 $A$-가군이라고 하자. $(f_i)_{1\leq i\leq n}$을
$B$의 극대 아이디얼의 원소들로 이루어진 열로서 다음 성질들을
만족한다고 하자.
\begin{enumerate}
  \item[(i)] $1\leq i\leq n$에 대하여 $f_iM\subset R$이다.
  \item[(ii)] $g_i$가 $B\otimes_A k$에서 $f_i$의 상이면, 열
  $(g_i)_{1\leq i\leq n}$은 $(M\otimes_A k)$-정칙이다.
  \item[(iii)] $M\otimes_A k$의 $(B\otimes_A k)$-부분가군
  $g_i(M\otimes_A k)$들의 합은 $M\otimes_A k$에서
  $R\otimes_A k$의 표준상과 같다. 이 표준상은 표준 준동형
  $M\otimes_A k\to N\otimes_A k$의 핵과 같다.
\end{enumerate}
그러면 열 $(f_i)$는 $M$-정칙이고
$R=\sum_{i=1}^n f_iM$이다.
\end{corollary}

\begin{proof}
\oldpage[0\textsubscript{IV-1}]{19}
$N$이 평탄 $A$-가군이므로, 완전열 $0\to R\to M\to N\to0$에서
\[
  0\longrightarrow R\otimes_A k\longrightarrow M\otimes_A k
  \longrightarrow N\otimes_A k\longrightarrow0
\]
이 완전열임이 따른다 ($0_{\mathrm I}$, 6.1.2). 따라서
$R\otimes_A k=\sum_i g_i(M\otimes_A k)$이다. $M$이 유한 생성
$B$-가군이고 $B$가 뇌터 환이므로 $R$도 유한 생성 $B$-가군이다.
$f_i$들이 $B$의 극대 아이디얼에 속하므로, 나카야마 보조정리에 따라
$R=\sum_i f_iM$이다.
\end{proof}

\begin{lemma}[15.1.18]
\label{0.15.1.18-fr}
$A$를 환, $M$을 분리이고 소진적인 여과 $(M_k)_{k\geq0}$가 주어진
$A$-가군, $f\in A$를 $\gr_\bullet(M)$-정칙 원소라 하자. 그러면
다음이 성립한다.
\begin{enumerate}
  \item[(i)] $f$는 $M$-정칙이다.
  \item[(ii)] 모든 $k$에 대하여 $M_k\cap fM=fM_k$이다. 따라서
  $M'=M/fM$에서 $M_k$의 표준상 $M'_k$는 $M_k/fM_k$와 표준적으로
  동형이다.
  \item[(iii)] 열
  \[
    0\longrightarrow M_{k+1}/fM_{k+1}\longrightarrow M_k/fM_k
    \longrightarrow\gr_k(M)/f\cdot\gr_k(M)\longrightarrow0
  \]
  은 완전하다. $M'$에 몫여과 $(M'_k)$를 주면
  $\gr_\bullet(M')$은
  $\gr_\bullet(M)/f\cdot\gr_\bullet(M)$과 표준적으로 동형이다.
\end{enumerate}
\end{lemma}

\begin{proof}
(i)은 $f$에 관한 가정과 Bourbaki, \emph{Alg.~comm.}, chap.~III,
\S~2, no~8, cor.~1 du th.~1에서 따른다.
$\gr_\bullet(M/M_k)$은 $\gr_\bullet(M)$의 직합인자이므로, 모든
$k$에 대하여 $f$는 $\gr_\bullet(M/M_k)$-정칙이고 따라서 (i)에
의하여 $(M/M_k)$-정칙이다. 이로써 (ii)가 증명된다. (iii)을
증명하려면 사상
$M_{k+1}/fM_{k+1}\to M_k/fM_k$이 단사임을 보이면 충분하다.
그런데 $x\in M_{k+1}$이고 $x\in fM_k$이면, (ii)에서
$x\in fM_{k+1}$이 따른다. 이로써 보조정리의 증명이 끝난다.
\end{proof}

\begin{proposition}[15.1.19]
\label{0.15.1.19-fr}
$A$를 환, $M$을 $M_0=M$인 여과 $(M_k)_{k\geq0}$가 주어진
$A$-가군, $(f_i)_{1\leq i\leq n}$을 $A$의 원소들로 이루어진
$\gr_\bullet(M)$-정칙열이라 하자. 또한 $0\leq i\leq n-1$에 대하여
\[
  M/\left(\sum_{j\leq i}f_jM\right)
\]
위의 $(M_k)$의 몫여과가 분리라고 가정하자. 그러면 열 $(f_i)$는
$M$-정칙이다. 또한
\[
  M'=M/\left(\sum_{i\leq n}f_iM\right)
\]
로 놓고 $M'$에 $(M_k)$의 몫여과를 주면, $\gr_\bullet(M')$은
\[
  \gr_\bullet(M)/
  \left(\sum_{i\leq n}f_i\gr_\bullet(M)\right)
\]
과 표준적으로 동형이다.
\end{proposition}

\begin{proof}
$n=1$일 때 명제는 보조정리 (15.1.18)에서 따른다. $n$에 관하여
귀납하고
\[
  N=M/\left(\sum_{i\leq n-1}f_iM\right)
\]
로 놓자. $N$ 위의 $(M_k)$의 몫여과를 $(N_k)$로 나타내면 이는 가정에
따라 분리이다. 또한 $\gr_\bullet(N)$은
\[
  \gr_\bullet(M)/
  \left(\sum_{i\leq n-1}f_i\gr_\bullet(M)\right)
\]
과 동형이다. 따라서 가정에 따라 $f_n$은
$\gr_\bullet(N)$-정칙이다. 그러므로 $N$에 보조정리 (15.1.18)을
적용하면 $f_n$은 $N$-정칙이고 $\gr_\bullet(M')$은
$\gr_\bullet(N)/f_n\gr_\bullet(N)$과 동형이다. 이로써 명제가
증명되었다.
\end{proof}

명제 (15.1.19)는 특히 다음 두 가정 가운데 하나를 만족하는 여과에
적용된다.
\begin{enumerate}
  \item[$1^\circ$] 여과 $(M_k)$가 유한이고 분리이다 (이 조건으로부터
  충분히 큰 $k$에 대하여 $M_k=0$임이 따르기 때문이다).
  \item[$2^\circ$] $A$가 뇌터 환이고, $\mathfrak{J}$가 $A$의
  근기에 포함되는 아이디얼이며, $M$이 유한 생성 $A$-가군이고,
  $(M_k)$가 $\mathfrak{J}$-준아딕 여과이다
  ($0_{\mathrm I}$, 7.3.5).
\end{enumerate}

\begin{corollary}[15.1.20]
\label{0.15.1.20-fr}
\oldpage[0\textsubscript{IV-1}]{20}
$A$를 환, $M$을 $A$-가군, $(f_i)_{1\leq i\leq n}$을 $A$의 원소들로
이루어진 $M$-정칙열이라 하자. 그러면 정수 $\nu_i\geq1$을 어떻게
택하든 열 $(f_i^{\nu_i})$은 $M$-정칙이고,
$M/(\sum_{i\leq n}f_i^{\nu_i}M)$은 그 몫들이 모두
$M/(\sum_{i\leq n}f_iM)$과 동형인 합성열을 갖는다.
\end{corollary}

\begin{proof}
$n$에 관하여 귀납하자. $n=1$일 때에는 가정에 따라 $M$에서 $f$에
의한 스칼라곱 자기준동형이 단사이고
$f^jM/f^{j+1}M$이 $M/fM$과 동형이므로 두 주장은 즉시 따른다. 이제
\[
  N=M/\left(\sum_{i\leq n-1}f_i^{\nu_i}M\right)
\]
로 놓자. 귀납 가정에 따라 $N$은 각 $\gr_k(N)$이
$M/(\sum_{i\leq n-1}f_iM)$과 동형인 유한 여과 $(N_k)$를 갖는다.
따라서 가정에 의해 $f_n$은 $\gr_\bullet(N)$-정칙이고, (15.1.19)에
따라 $f_n$은 $N$-정칙이다. 처음에 증명한 $n=1$인 경우를 적용하면
$f_n^{\nu_n}$은 $N$-정칙이고, $N/f_n^{\nu_n}N$은 그 몫들이
$N/f_nN$과 동형인 합성열을 갖는다. 그런데 $(N_k)$에서 유도되는
$N/f_nN$의 몫여과에 대하여, 모든 $k$에서 (15.1.19)에 따라
$\gr_k(N/f_nN)$은 $\gr_k(N)/f_n\gr_k(N)$과 동형이고, 따라서
$M/(\sum_{i\leq n}f_iM)$과 동형이다. 이로써 따름정리가 증명되었다.
\end{proof}

\begin{proposition}[15.1.21]
\label{0.15.1.21-fr}
$A$, $B$를 두 뇌터 국소환, $\varphi:A\to B$를 국소 준동형,
$M$을 유한 생성 $B$-가군, $(f_i)_{1\leq i\leq n}$을 $A$의 극대
아이디얼의 원소들로 이루어진 $A$-정칙열이라 하자. 다음 성질들은
서로 동치이다.
\begin{enumerate}
  \item[a)] $M$은 평탄 $A$-가군이다.
  \item[b)] $\mathfrak{J}=\sum_{i\leq n}f_iA$라 하면,
  $M/\mathfrak{J}M$은 평탄 $(A/\mathfrak{J})$-가군이고 열 $(f_i)$는
  $M$-정칙이다.
\end{enumerate}
\end{proposition}

\begin{proof}
a)가 b)를 함의함은 (15.1.13)과
($0_{\mathrm{III}}$, 10.2.2)에서 따른다. 역으로 열 $(f_i)$가
$M$-정칙이면 $M$-준정칙이므로 (15.1.9), 준동형 (15.1.1.1)은
전단사이다. 한편 열 $(f_i)_{1\leq i\leq n}$은 $A$-정칙이라고도
가정했으므로 $A$-준정칙이고 (15.1.9), 이에 대응하는 표준 준동형
$(A/\mathfrak{J})[T_1,\ldots,T_n]\to\gr_\bullet(A)$도 전단사이다.
따라서 표준 준동형 ($0_{\mathrm{III}}$, 10.1.1.2)
\[
  \gr_0(M)\otimes_{A/\mathfrak{J}}\gr_\bullet(A)
  \longrightarrow\gr_\bullet(M)
\]
이 전단사이다. 이제 $f_i$들이 $A$의 극대 아이디얼에 속하고
$\varphi$가 국소 준동형이므로, 결론은
($0_{\mathrm{III}}$, 10.2.2)에서 따른다.
\end{proof}

\subsection{\texorpdfstring{$\mathscr{F}$}{F}-정칙열}
\label{subsection:0.15.2-fr}

\begin{env}[15.2.1]
\label{0.15.2.1-fr}
$X$를 환 달린 공간, $\mathscr{F}$를 $\mathcal{O}_X$-가군,
$f_i$ ($1\leq i\leq n$)를 $X$ 위의 $\mathcal{O}_X$의 단면들로
이루어진 열, $\mathscr{J}$를 $f_i$들이 생성하는 $\mathcal{O}_X$의
아이디얼층이라 하자. 곧 $\mathscr{J}$는 $f_i$들이 표준적으로 정하는
준동형 $\mathcal{O}_X^n\to\mathcal{O}_X$의 상이다
($0_{\mathrm I}$, 5.1.1). 부분-$\mathcal{O}_X$-가군
$\mathscr{J}^k\mathscr{F}$들은 $\mathscr{F}$ 위의 여과를 정하며, 다시
$\gr_k(\mathscr{F})=\mathscr{J}^k\mathscr{F}/
\mathscr{J}^{k+1}\mathscr{F}$로 놓는다. $S_k$를
$\mathbb{Z}[T_1,\ldots,T_n]$에서 전체 차수가 $k$인 동차다항식들의
집합이라 하면, $S_k$를 $X$ 위의 상수층으로 볼 수 있고 (15.1.1)에서와
같이 표준 준동형
\[
\label{0.15.2.1.1-fr}
  \varphi_k:\gr_0(\mathscr{F})\otimes_{\mathbb Z}S_k
  \longrightarrow\gr_k(\mathscr{F}).
\tag{15.2.1.1}
\]
\oldpage[0\textsubscript{IV-1}]{21}
다음과 같이 정의한다. $X$의 임의의 열린집합 $U$ (또는 $X$의
위상에 대한 기저의 열린집합
$U$)에 대하여, $n$개의 단면
$f_i|U\in\Gamma(U,\mathcal{O}_X)$가 정하는 준동형 (15.1.1.1)
\[
  \varphi_{k,U}:\gr_0(\Gamma(U,\mathscr{F}))\otimes_{\mathbb Z}S_k
  \longrightarrow\gr_k(\Gamma(U,\mathscr{F}))
\]
을 생각한다. 여기서 $\Gamma(U,\mathcal{O}_X)$-가군
$\Gamma(U,\mathscr{F})$에는 부분가군
$(\Gamma(U,\mathscr{J}))^k\Gamma(U,\mathscr{F})$들이 이루는 여과를
준다. 각 $k$에 대하여 $\mathcal{O}_X$-가군
$\gr_k(\mathscr{F})$는 준층
$U\mapsto\gr_k(\Gamma(U,\mathscr{F}))$에 결부된 층이므로,
$\varphi_{k,U}$들은 $\mathcal{O}_X$-가군의 준동형 (15.2.1.1)을
정한다. 가군의 귀납극한에 대한 완전성 성질과 귀납극한이 텐서곱과
가환한다는 성질에 따라, 앞의 구성으로부터 모든 $x\in X$에 대하여
(15.2.1.1)의 두 변의 $x$에서의 줄기 사이 준동형
$(\varphi_k)_x$는 (15.1.1)의 준동형
\[
\label{0.15.2.1.2-fr}
  \gr_0(\mathscr{F}_x)\otimes_{\mathbb Z}S_k
  \longrightarrow\gr_k(\mathscr{F}_x)
\tag{15.2.1.2}
\]
과 동일시된다. 이 준동형은 $\mathcal{O}_x$의 원소들로 이루어진 열
$((f_i)_x)_{1\leq i\leq n}$이 정한다.
\end{env}

\begin{env}[15.2.2]
\label{0.15.2.2-fr}
표기는 (15.2.1)과 같이 하자. $f_j\mathscr{F}$로 $f_j$가 정하는
$\mathscr{F}$의 자기준동형의 상인 부분-$\mathcal{O}_X$-가군을
나타낼 때, 모든 $1\leq i\leq n$에 대하여 $f_i$가 정하는
$\mathscr{F}/(\sum_{j=1}^{i-1}f_j\mathscr{F})$의 자기준동형이
단사이면 열 $(f_i)_{1\leq i\leq n}$을
$\mathscr{F}$-정칙이라 한다. 준동형들 (15.2.1.1)이 단사이면 열
$(f_i)$를 $\mathscr{F}$-준정칙이라 한다. 이 정의들과 (15.2.1)의
설명에 따라, 열 $(f_i)$가 $\mathscr{F}$-정칙 (각각
$\mathscr{F}$-준정칙)이기 위한 필요충분조건은 모든 $x\in X$에
대하여 열 $((f_i)_x)_{1\leq i\leq n}$이
$\mathscr{F}_x$-정칙 (각각 $\mathscr{F}_x$-준정칙)인 것이다.
따라서 모든 $\mathscr{F}$-정칙열은 $\mathscr{F}$-준정칙열이다
(15.1.9). 역도 참인 것은 $\mathscr{F}$가 유한형
$\mathcal{O}_X$-가군이고, 모든 $x\in X$에 대하여
$\mathcal{O}_x$가 뇌터 환이며, 더 나아가 모든 $(f_i)_x$가
$\mathcal{O}_x$의 근기에 속하는 경우이다 (15.1.11).

$X=\operatorname{Spec}(A)$가 아핀 스킴이고
$\mathscr{F}=\widetilde{M}$이면, (I, 1.3)에 따라 열 $(f_i)$가
$\mathscr{F}$-정칙 (각각 $\mathscr{F}$-준정칙)이기 위한
필요충분조건은 이 열이 $M$-정칙 (각각 $M$-준정칙)인 것이다.

$\mathscr{F}=\mathcal{O}_X$일 때에는 앞의 정의들에서
$\mathscr{F}$를 언급하지 않기로 한다.
\end{env}

\begin{remark}[15.2.3]
\label{0.15.2.3-fr}
$X=\operatorname{Spec}(A)$이고
$\mathfrak{J}=\sum_{i=1}^n f_iA$라 하자. $V(\mathfrak{J})$와 만나지
않는 임의의 열린집합 $U$에서 열 $(f_j|U)$는 모든 $A$-가군 $M$에
대하여 $(\widetilde{M}|U)$-준정칙이다. 실제로 모든 $x\in U$에서
$\mathfrak{J}_x=\mathcal{O}_x$이고 (15.2.1.2)의 두 변은 모두
영이다. 따라서 $\widetilde{M}$-정칙열이라는 개념은
$V(\mathfrak{J})$의 점들의 근방에서만 의미가 있다.
\end{remark}

\begin{proposition}[15.2.4]
\label{0.15.2.4-fr}
$X$를 환 달린 공간, $\mathscr{F}$를 연접 $\mathcal{O}_X$-가군,
$x$를 $X$의 한 점, $(f_i)_{1\leq i\leq n}$을
$\Gamma(X,\mathcal{O}_X)$의 원소들로 이루어진 열이라 하자. 열
$((f_i)_x)_{1\leq i\leq n}$이 $\mathscr{F}_x$-정칙이면, $x$의 어떤
열린근방 $U$가 존재하여 열 $(f_i|U)_{1\leq i\leq n}$은
$(\mathscr{F}|U)$-정칙이다.
\end{proposition}

\begin{proof}
실제로 $1\leq i\leq n$에 대하여 $\mathcal{O}_X$-가군
$\mathscr{F}/(\sum_{j=1}^{i-1}f_j\mathscr{F})$는 연접이고
($0_{\mathrm I}$, 5.3.4), 따라서 $f_i$가 정하는 이
$\mathcal{O}_X$-가군의 자기준동형의 핵도 연접이다
($0_{\mathrm I}$, 5.3.4). 이 핵의 $x$에서의 줄기는 영이므로, $x$의
어떤 근방에 대한 그 제한도 영이다 ($0_{\mathrm I}$, 5.2.2).
\end{proof}

\begin{proposition}[15.2.5]
\label{0.15.2.5-fr}
\oldpage[0\textsubscript{IV-1}]{22}
$u=(\psi,\theta):X\to Y$를 환 달린 공간의 평탄 사상이라 하자
($0_{\mathrm I}$, 6.7.1). $\mathscr{F}$를
$\mathcal{O}_Y$-가군, $\mathscr{G}=u^*(\mathscr{F})$로 놓고,
$(f_i)_{1\leq i\leq n}$을 $\Gamma(Y,\mathcal{O}_Y)$의 원소들로
이루어진 열, $(g_i)_{1\leq i\leq n}$을 그 상들
$g_i=\theta(f_i)\in\Gamma(Y,u_*(\mathcal{O}_X))
=\Gamma(X,\mathcal{O}_X)$로 이루어진 열이라 하자. 열 $(f_i)$가
$\mathscr{F}$-정칙이면 열 $(g_i)$는
$\mathscr{G}$-정칙이다. $u$가 충실 평탄 사상이면 역도 참이다
($0_{\mathrm I}$, 6.7.8).
\end{proposition}

\begin{proof}
$x$를 $X$의 한 점, $y=u(x)=\psi(x)$라 하자. 다음이 성립한다.
\[
  \mathscr{G}_x=\psi_x^*(\mathscr{F}_y)
  \otimes_{\psi_x^*(\mathcal{O}_y)}\mathcal{O}_x,
  \qquad\text{그리고}\qquad
  (g_i)_x=\theta_x^\#(\psi_x^*((f_i)_y)).
\]
가정에 따라
$\theta_x^\#:\psi_x^*(\mathcal{O}_y)\to\mathcal{O}_x$는
$\mathcal{O}_x$를 평탄 $\psi_x^*(\mathcal{O}_y)$-가군으로 만들고
함자 $\psi^*$는 완전하므로, 첫째 주장은 (15.1.14)의 대응하는
주장에서 따른다. 둘째 주장도 마찬가지이다. 여기서는 $u$가 충실
평탄하다는 가정이 $\psi$가 전사이고 $\theta_x^\#$가
$\mathcal{O}_x$를 충실 평탄 $\psi_x^*(\mathcal{O}_y)$-가군으로
만든다는 뜻임을 함께 사용한다.
\end{proof}

\section{뇌터 국소환에서의 차원과 깊이}
\label{section:0.16-fr}

\subsection{환의 차원}
\label{subsection:0.16.1-fr}

\begin{env}[16.1.1]
\label{0.16.1.1-fr}
환 $A$의 스펙트럼 $\operatorname{Spec}(A)$의 (조합적) 차원
(14.2.1)을 $A$의 \emph{차원} (또는 \emph{크룰 차원})이라 하고
$\dim A$로 나타낸다. $\operatorname{Spec}(A)$의 기약 닫힌
부분집합은 $\mathfrak{p}$가 $A$의 소 아이디얼일 때의
$V(\mathfrak{p})$들이므로 (Bourbaki, \emph{Alg.~comm.}, chap.~II,
\S~4, no~3, prop.~14의 cor.~1), 이는 $\dim(A)$가 $A$의 소
아이디얼 사슬 (14.1.1)
\[
  \mathfrak{p}_0\subset\mathfrak{p}_1\subset\cdots\subset\mathfrak{p}_n
\]
의 길이들의 상한이라는 것과 같다. $A\neq0$이면 $\dim A\geq0$이다.
영이 아닌 아르틴 환 (특히 체)의 차원은 $0$이다 (Bourbaki,
\emph{Alg.~comm.}, chap.~IV, \S~2, no~5, prop.~7). 체가 아닌
데데킨트 환 (특히 $\mathbb{Z}$ 또는 체 $k$ 위의 한 부정원 다항식환
$k[T]$)의 차원은 $1$이다 (Bourbaki, \emph{Alg.~comm.}, chap.~VII,
\S~2). 뇌터 환의 차원이 반드시 유한한 것은 아니다 [30].

$(\mathfrak{p}_\alpha)$가 $A$의 극소 소 아이디얼들의 족이면
(14.1.2.1)에 따라
\[
\label{0.16.1.1.1-fr}
  \dim(A)=\sup_\alpha(\dim(A/\mathfrak{p}_\alpha)).
\tag{16.1.1.1}
\]
이다.
\end{env}

\begin{env}[16.1.2]
\label{0.16.1.2-fr}
$A$의 임의의 아이디얼 $\mathfrak{a}$에 대하여
$A/\mathfrak{a}$의 소 아이디얼들은 $\mathfrak{a}$를 포함하는
$A$의 소 아이디얼들과 전단사로 대응하므로
\[
\label{0.16.1.2.1-fr}
  \dim(A/\mathfrak{a})\leq\dim(A).
\tag{16.1.2.1}
\]
이다. 더 나아가 $\mathfrak{a}$가 $A$의 어느 극소 소 아이디얼
$\mathfrak{p}_\alpha$에도 포함되지 않으면, $\mathfrak{a}$를 포함하는
$A$의 모든 소 아이디얼 사슬에 $\mathfrak{p}_\alpha$ 가운데 하나를
덧붙일 수 있으므로
\[
\label{0.16.1.2.2-fr}
  1+\dim(A/\mathfrak{a})\leq\dim(A).
\tag{16.1.2.2}
\]
이다.
\end{env}

\begin{env}[16.1.3]
\label{0.16.1.3-fr}
\oldpage[0\textsubscript{IV-1}]{23}
$S$가 $A$의 곱셈적 부분집합이면 $S^{-1}A$의 소 아이디얼들은
$S$와 만나지 않는 $A$의 소 아이디얼들과 전단사로 대응하므로
\[
\label{0.16.1.3.1-fr}
  \dim(S^{-1}A)\leq\dim(A).
\tag{16.1.3.1}
\]
이다. $A$의 임의의 소 아이디얼 $\mathfrak{p}$에 대하여 정의에 따라
(14.2.1)
\[
\label{0.16.1.3.2-fr}
  \dim(A_{\mathfrak{p}})
  =\operatorname{codim}(V(\mathfrak{p}),\operatorname{Spec}(A)).
\tag{16.1.3.2}
\]
이다. 이 수를 소 아이디얼 $\mathfrak{p}$의 \emph{높이}라고도 하며
$\operatorname{ht}(\mathfrak{p})$로 나타낸다. 더 일반적으로 $A$의
임의의 아이디얼 $\mathfrak{a}$의 \emph{높이}를 수
\[
\label{0.16.1.3.3-fr}
  \operatorname{ht}(\mathfrak{a})
  =\operatorname{codim}(V(\mathfrak{a}),\operatorname{Spec}(A)).
\tag{16.1.3.3}
\]
로 정의한다. $\mathfrak{a}$를 포함하는 소 아이디얼들 가운데 극소인
것들의 족을 $(\mathfrak{p}_\lambda)$라 하면, 따라서 (14.2.1)에 의해
\[
\label{0.16.1.3.4-fr}
  \operatorname{ht}(\mathfrak{a})
  =\inf_\lambda(\operatorname{ht}(\mathfrak{p}_\lambda)).
\tag{16.1.3.4}
\]
이다. 명백히
\[
\label{0.16.1.3.5-fr}
  \dim(A)=\sup_{\mathfrak{m}}(\dim(A_{\mathfrak{m}}))
\tag{16.1.3.5}
\]
이고, 여기서 $\mathfrak{m}$은 $A$의 모든 극대 아이디얼을 달린다.
\end{env}

\begin{env}[16.1.4]
\label{0.16.1.4-fr}
$A$의 임의의 소 아이디얼 $\mathfrak{p}$에 대하여 (14.2.2.2)
\[
\label{0.16.1.4.1-fr}
  \dim(A_{\mathfrak{p}})+\dim(A/\mathfrak{p})\leq\dim A.
\tag{16.1.4.1}
\]
이다. $\operatorname{Spec}(A)$가 카테너리 (각각 등차원,
등여차원, 쌍등차원)이면 $A$를 카테너리 (각각 등차원, 등여차원,
쌍등차원)라고 한다. $A$가 뇌터이고 쌍등차원이면 모든 소 아이디얼
$\mathfrak{p}$에 대하여 (16.1.4.1)의 두 변은 같다 (14.3.5).

$A$의 임의의 아이디얼 $\mathfrak{a}$에 대하여
$A/\mathfrak{a}$의 소 아이디얼들은 $\mathfrak{a}$를 포함하는
$A$의 소 아이디얼들과 전단사로 대응하므로, $A$가 카테너리이면
$A/\mathfrak{a}$도 카테너리이다. 마찬가지로 $A$의 임의의 곱셈적
부분집합 $S$에 대하여 $S^{-1}A$의 소 아이디얼들은 $S$와 만나지 않는
$A$의 소 아이디얼들과 전단사로 대응하므로, $A$가 카테너리이면
$S^{-1}A$도 카테너리이다.

또한 소 아이디얼 $\mathfrak{p}$에 포함되는 $A$의 소 아이디얼들은
$A_{\mathfrak{p}}$의 소 아이디얼들과 전단사로 대응하므로, $A$가
카테너리이기 위한 필요충분조건은 모든 $\mathfrak{p}$에 대하여
$A_{\mathfrak{p}}$가 카테너리인 것이다. 따라서 $A$가 카테너리라는
것은 ((14.3.2)와 (16.1.3.2)) $A$의 모든 소 아이디얼 쌍
$\mathfrak{p}$, $\mathfrak{q}$에 대하여
$\mathfrak{q}\subset\mathfrak{p}$이면
\[
\label{0.16.1.4.2-fr}
  \dim(A_{\mathfrak{q}})
  +\dim(A_{\mathfrak{p}}/\mathfrak{q}A_{\mathfrak{p}})
  =\dim(A_{\mathfrak{p}}).
\tag{16.1.4.2}
\]
라는 뜻이다. $A$가 쌍등차원이면 모든 소 아이디얼
$\mathfrak{p}$에 대하여 $A_{\mathfrak{p}}$도 쌍등차원이다.
\end{env}

\begin{proposition}[16.1.5]
\label{0.16.1.5-fr}
$\rho:A\to B$를 $B$가 $A$ 위에서 정수적인 대수가 되게 하는 환
준동형이라 하자. 그러면 $\dim(B)\leq\dim(A)$이다. 더 나아가
$\rho$가 단사이면 $\dim(B)=\dim(A)$이다.
\end{proposition}

\begin{proof}
\oldpage[0\textsubscript{IV-1}]{24}
$\mathfrak{n}$이 $\rho$의 핵이면 $\rho(A)$는 $A/\mathfrak{n}$과
동형이므로 (16.1.2.1)에 따라 $\dim\rho(A)\leq\dim(A)$이고, $B$는
$\rho(A)$ 위에서 정수적이므로 $A\subset B$인 경우의 둘째 주장만
증명하면 충분하다. $B$의 소 아이디얼 사슬
$\mathfrak{p}'_0\subset\mathfrak{p}'_1\subset\cdots\subset
\mathfrak{p}'_n$이 주어지면, $\mathfrak{p}'_i\cap A$들은 $A$ 안에서
서로 다른 소 아이디얼들이므로 (Bourbaki, \emph{Alg.~comm.}, chap.~V,
\S~2, no~1, prop.~1의 cor.~1) 소 아이디얼 사슬을 이룬다. 역으로
$A$의 임의의 소 아이디얼 사슬
$\mathfrak{p}_0\subset\mathfrak{p}_1\subset\cdots\subset
\mathfrak{p}_m$에 대하여, 각 $i$마다
$\mathfrak{p}'_i\cap A=\mathfrak{p}_i$이고 $0\leq i\leq m-1$일 때
$\mathfrak{p}'_i\subset\mathfrak{p}'_{i+1}$인 $B$의 소 아이디얼
$\mathfrak{p}'_i$가 존재함이 알려져 있다 (\emph{loc.~cit.}, th.~1의
cor.~2). 이로부터 결론이 따른다.
\end{proof}

\begin{proposition}[16.1.6]
\label{0.16.1.6-fr}
$B$를 정역, $A$를 $B$의 국소 부분환, $\mathfrak{m}$을 그 극대
아이디얼이라 하자. $A$가 정수적으로 닫혀 있고 $B$가 $A$ 위에서
정수적이라고 가정하자. 그러면 $B$의 모든 극대 아이디얼
$\mathfrak{n}$에 대하여
$\dim(B_{\mathfrak{n}})=\dim(A)$이다.
\end{proposition}

\begin{proof}
(16.1.5)의 논증의 첫 부분으로부터
$\dim(B_{\mathfrak{n}})\leq\dim(A)$를 얻는다. 역으로 $A$의 소
아이디얼 사슬
$\mathfrak{p}_0\subset\mathfrak{p}_1\subset\cdots\subset
\mathfrak{p}_k=\mathfrak{m}$을 생각하자.
$\mathfrak{n}\cap A=\mathfrak{m}$임이 알려져 있고 (Bourbaki,
\emph{Alg.~comm.}, chap.~V, \S~2, no~1, prop.~1), 주어진 가정에 따라
제2 코언--자이덴베르크 정리를 적용할 수 있다 (Bourbaki,
\emph{Alg.~comm.}, chap.~V, \S~2, no~4, th.~3). 이 정리는 모든
$i$에 대하여 $\mathfrak{p}'_i\cap A=\mathfrak{p}_i$인 $B$의 소
아이디얼 사슬
$\mathfrak{p}'_0\subset\mathfrak{p}'_1\subset\cdots\subset
\mathfrak{p}'_k=\mathfrak{n}$의 존재를 보장한다. 따라서 곧바로
$\dim(B_{\mathfrak{n}})\geq\dim(A)$를 얻으며, 이것으로 명제가
증명되었다.
\end{proof}

\begin{env}[16.1.7]
\label{0.16.1.7-fr}
$M$을 $A$-가군이라 하자. 환 $A/\operatorname{Ann}(M)$의 차원을
$M$의 \emph{차원}이라 하고 $\dim_A(M)$ 또는 $\dim(M)$으로
나타낸다. 이 환은 $M$의 스칼라곱환 $A_M$과 동형이다. 따라서 모든
$k>0$에 대하여 $\dim(M^k)=\dim(M)$이다. $M$이 유한형이면
$\operatorname{Ann}(M)$을 포함하는 $A$의 소 아이디얼들은 정확히
$\operatorname{Supp}(M)$에 속하는 소 아이디얼들이고
($0_{\mathrm I}$, 1.7.4), 이 지지집합은 $\operatorname{Spec}(A)$에서
닫혀 있으므로
\[
\label{0.16.1.7.1-fr}
  \dim(M)=\dim(\operatorname{Supp}(M))\leq\dim(A).
\tag{16.1.7.1}
\]
이다. $M$, $N$이 $N\subset M$인 두 $A$-가군이면
\[
\label{0.16.1.7.2-fr}
  \dim(N)\leq\dim(M),\qquad \dim(M/N)\leq\dim(M).
\tag{16.1.7.2}
\]
이다. $S$가 $A$의 곱셈적 부분집합이고 $M$이 유한형이면
\[
\label{0.16.1.7.3-fr}
  \dim_{S^{-1}A}(S^{-1}M)\leq\dim_A(M).
\tag{16.1.7.3}
\]
이다. 실제로
$\operatorname{Ann}(S^{-1}M)=S^{-1}\operatorname{Ann}(M)$이고
(Bourbaki, \emph{Alg.~comm.}, chap.~II, \S~2, no~4),
$S^{-1}A/S^{-1}\operatorname{Ann}(M)
=S^{-1}(A/\operatorname{Ann}(M))$이다.

한편 $\operatorname{Ann}(M)$을 포함하는 모든 극대인 소 아이디얼 사슬은
어떤 극대 아이디얼 $\mathfrak{m}$에서 끝난다. 따라서 그 길이는
$A_{\mathfrak{m}}$의 대응하는 소 아이디얼 사슬의 길이와 같으며, 이
소 아이디얼들은
$\operatorname{Ann}(M_{\mathfrak{m}})
=(\operatorname{Ann}(M))_{\mathfrak{m}}$을 포함한다. 이 설명과
(16.1.7.3)으로부터
\[
\label{0.16.1.7.4-fr}
  \dim_A(M)=\sup_{\mathfrak{m}}
  (\dim_{A_{\mathfrak{m}}}(M_{\mathfrak{m}}))
\tag{16.1.7.4}
\]
를 얻고, 여기서 $\mathfrak{m}$은 $A$의 모든 극대 아이디얼을 달린다.

$A/\operatorname{Ann}(M)$이 카테너리 (각각 등차원, 등여차원,
쌍등차원)이면 $M$을 카테너리 (각각 등차원, 등여차원, 쌍등차원)라고
한다.
\end{env}

\begin{proposition}[16.1.8]
\label{0.16.1.8-fr}
\oldpage[0\textsubscript{IV-1}]{25}
$A$를 환, $M$을 유한형 $A$-가군, $\mathfrak{p}$를 $A$의 소
아이디얼이라 하자. 그러면
\[
\label{0.16.1.8.1-fr}
  \dim_{A_{\mathfrak{p}}}(M_{\mathfrak{p}})
  +\dim_A(M/\mathfrak{p}M)\leq\dim_A(M).
\tag{16.1.8.1}
\]
이다.
\end{proposition}

\begin{proof}
$\operatorname{Ann}(M_{\mathfrak{p}})
=(\operatorname{Ann}(M))_{\mathfrak{p}}$이므로,
$\operatorname{Ann}(M_{\mathfrak{p}})$을 포함하는
$A_{\mathfrak{p}}$의 소 아이디얼들은 $\operatorname{Ann}(M)$을
포함하고 $\mathfrak{p}$에 포함되는 $A$의 소 아이디얼들과 전단사로
대응한다. 한편 $\operatorname{Ann}(M)=\mathfrak{n}$이라 하면
$V(\operatorname{Ann}(M/\mathfrak{p}M))
=V(\mathfrak{p}+\mathfrak{n})$이고 ($0_{\mathrm I}$, 1.7.5), 따라서
$\operatorname{Ann}(M/\mathfrak{p}M)$을 포함하는 $A$의 소
아이디얼들은 $\mathfrak{p}+\mathfrak{n}$을 포함한다. 이로부터 결론이
따른다.

마지막 설명으로부터 더 나아가
\[
\label{0.16.1.8.2-fr}
  \dim_{A/\mathfrak{p}}(M/\mathfrak{p}M)=\dim_A(M/\mathfrak{p}M).
\tag{16.1.8.2}
\]
를 얻는다.
\end{proof}

\begin{proposition}[16.1.9]
\label{0.16.1.9-fr}
$A$를 뇌터 환, $\rho:A\to B$를 환 준동형, $M$을
$M_{[\rho]}$가 유한형 $A$-가군이 되는 $B$-가군이라 하자. 그러면
\[
\label{0.16.1.9.1-fr}
  \dim_A(M_{[\rho]})=\dim_B(M).
\tag{16.1.9.1}
\]
이다.
\end{proposition}

\begin{proof}
$B$-가군 $M$의 스칼라곱환 $B_M$은
$C=\operatorname{End}_A(M_{[\rho]})$의 부분환과 동일시되고, $C$는
유한형 $A$-가군이므로 (Bourbaki, \emph{Alg.~comm.}, chap.~III,
\S~3, no~1, lemme~2) $B_M$도 그러하다. 한편 $A$-가군
$M_{[\rho]}$의 스칼라곱환 $A_M$은 $B_M$의 부분환이며 $C$ 안에서
$A$의 표준상이다. 따라서 $B_M$은 $A_M$ 위에서 유한이고, 결론은
(16.1.5)에서 따른다.
\end{proof}

\begin{env}[16.1.10]
\label{0.16.1.10-fr}
$A$를 뇌터 환이라 하자. $M$이 유한 길이 $A$-가군이면
$\operatorname{Supp}(M)$이 유한 이산공간임이 알려져 있으므로
(Bourbaki, \emph{Alg.~comm.}, chap.~IV, \S~2, no~5, prop.~7과
prop.~7의 cor.~1), $M\neq0$일 때 $\dim(M)=0$이다. 역으로 $M$이
$\dim(M)=0$인 유한형 $A$-가군이면 $\operatorname{Supp}(M)$의 모든
점은 닫혀 있다. 다시 말해 모두 $A$의 극대 아이디얼이므로
(\emph{loc.~cit.}, prop.~7), $M$은 유한 길이이다.
\end{env}

\subsection[뇌터 반국소환의 차원]{뇌터 반국소환의
차원\protect\footnote{이 번호는 세르가 1958년 콜레주 드 프랑스
강좌에서 한 강의를 재현한 것이다.}}
\label{subsection:0.16.2-fr}

\begin{env}[16.2.1]
\label{0.16.2.1-fr}
$A$를 뇌터 반국소환, $\mathfrak{r}$을 그 근기라 하자. $A$의
\emph{정의 아이디얼} $\mathfrak{q}$란
$\mathfrak{q}\subset\mathfrak{r}$이고 $\mathfrak{r}$의 어떤
거듭제곱을 포함하는 아이디얼임을 상기하자. 이는 $\mathfrak{q}$를
포함하는 소 아이디얼들이 모두 극대라는 것, 또는
$A/\mathfrak{q}$가 아르틴 환이라는 것과 같다. $M$을 유한형
$A$-가군이라 하자. 모든 정수 $n>0$에 대하여
$M/\mathfrak{q}^nM$은 유한 길이 $A$-가군이고, 그 길이는
$\sum_{i=0}^{n-1}\operatorname{long}(\gr_i(M))$과 같다. 여기서
$\gr_\bullet(M)$은 $\mathfrak{q}$-준아딕 여과를 준 $M$에 결부된
$(A/\mathfrak{q})$-등급가군이다. 충분히 큰 $n$에 대하여
$\operatorname{long}(\gr_n(M))=Q(n)$이 되는 유리수 계수 다항식
$Q(n)$이 존재함이 알려져 있다 (III, 2.5.4). 따라서 충분히 큰
$n$에 대하여
\[
  \operatorname{long}(M/\mathfrak{q}^nM)=P_{\mathfrak{q}}(M,n)
\]
이 되는 $n$의 유리수 계수 다항식 $P_{\mathfrak{q}}(M,n)$이 존재한다.
이 다항식이 유일하고 $M\neq0$이면 그 최고차항 계수가 $>0$임은
명백하다. 이를 $\mathfrak{q}$-준아딕 여과에 대한 $M$의
\emph{힐베르트--사무엘 다항식}이라 한다.
\end{env}

\begin{lemma}[16.2.2.1]
\label{0.16.2.2.1-fr}
\oldpage[0\textsubscript{IV-1}]{26}
$(M_n)$을 $M$의 $\mathfrak{q}$-좋은 여과라 하자 (Bourbaki,
\emph{Alg.~comm.}, chap.~III, \S~3, no~1). 그러면 충분히 큰 $n$에
대하여 $\operatorname{long}(M/M_n)$은 $n$의 한 다항식과 같으며, 그
차수와 최고차항 계수는 $P_{\mathfrak{q}}(M,n)$의 차수와 최고차항
계수와 각각 같다.
\end{lemma}

\begin{proof}
실제로 어떤 $n_0$가 존재하여 $n\geq n_0$이면
$M_{n+1}=\mathfrak{q}M_n$이다. 따라서 충분히 큰 $n$에 대하여
$\mathfrak{q}^{n+n_0}M\subset M_{n+n_0}=\mathfrak{q}^nM_{n_0}
\subset\mathfrak{q}^nM\subset M_n$이고, 이로부터
\[
  P_{\mathfrak{q}}(M,n+n_0)\geq
  \operatorname{long}(M/M_{n+n_0})\geq
  P_{\mathfrak{q}}(M,n)\geq\operatorname{long}(M/M_n)
\]
를 얻으므로 결론이 따른다.
\end{proof}

$\mathfrak{q}'$이 두 번째 정의 아이디얼이면 어떤 정수 $m$에 대하여
$\mathfrak{q}'\supset\mathfrak{q}^m$이므로
$P_{\mathfrak{q}'}(M,n)\leq P_{\mathfrak{q}}(M,mn)$이다. 따라서
$P_{\mathfrak{q}}(M,n)$의 차수는 택한 정의 아이디얼
$\mathfrak{q}$에 의존하지 않는다. 이를 $d(M)$이라 쓴다.

\begin{lemma}[16.2.2.2]
\label{0.16.2.2.2-fr}
$0\to M'\to M\to M''\to0$을 유한형 $A$-가군들의 완전열이라 하자.
그러면 다항식
\[
  P_{\mathfrak{q}}(M,n)-P_{\mathfrak{q}}(M',n)
  -P_{\mathfrak{q}}(M'',n)
\]
의 차수는 $\leq d(M')-1\leq d(M)-1$이다.
\end{lemma}

\begin{proof}
실제로 $M'_n=M'\cap\mathfrak{q}^nM$으로 주어지는 여과는
$\mathfrak{q}$-좋은 여과이다 (Bourbaki, \emph{Alg.~comm.},
chap.~III, \S~3, no~1, prop.~1). 또한
\[
  \operatorname{long}(M/\mathfrak{q}^nM)
  =\operatorname{long}(M''/\mathfrak{q}^nM'')
  +\operatorname{long}(M'/M'_n)
\]
이므로, $M'$에 (16.2.2.1)을 적용하면 보조정리가 곧바로 따른다.
\end{proof}

\begin{theorem}[16.2.3]
\label{0.16.2.3-fr}
\textup{(크룰--슈발레--사무엘).}
$A$를 뇌터 반국소환, $\mathfrak{r}$을 그 근기, $M\neq0$을 유한형
$A$-가군이라 하자. $d(M)$을 $M$의 힐베르트--사무엘 다항식의
차수라 하되, $A$의 어느 정의 아이디얼을 택해도 좋다. 한편
$M/(x_1M+\cdots+x_nM)$이 유한 길이가 되는 원소
$x_1,\ldots,x_n\in\mathfrak{r}$가 존재하도록 하는 정수 $n$들의
하한을 $s(M)$이라 하자. 그러면
\[
  d(M)=s(M)=\dim(M).
\]
이다.
\end{theorem}

\begin{lemma}[16.2.3.1]
\label{0.16.2.3.1-fr}
각 $x\in\mathfrak{r}$에 대하여 $x$에 의해 소멸되는 $M$의 원소들로
이루어진 부분가군을 ${}_xM$이라 하자. 그러면 다음이 성립한다.
\begin{enumerate}
  \item[\rm{(i)}] $s(M)\leq s(M/xM)+1$이다.
  \item[\rm{(ii)}] $\mathfrak{p}_i$들을
  $\operatorname{Supp}(M)$에 속하고
  \[
    \dim(A/\mathfrak{p}_i)=\dim(M)\quad(1\leq i\leq m)
  \]
  을 만족하는 소 아이디얼들이라 하자.
\end{enumerate}
모든 $i$에 대하여 $x\notin\mathfrak{p}_i$이면
$\dim(M/xM)\leq\dim(M)-1$이다.
\begin{enumerate}
  \item[\rm{(iii)}] $\mathfrak{q}$가 $A$의 정의 아이디얼이면 다항식
  \[
    P_{\mathfrak{q}}({}_xM,n)-P_{\mathfrak{q}}(M/xM,n)
  \]
  의 차수는 $\leq d(M)-1$이다.
\end{enumerate}
\end{lemma}

\begin{proof}
(i)은 자명하다. 실제로 $N=M/xM$에 대하여
$(x_i)_{1\leq i\leq m}$이
$N/(x_1N+\cdots+x_mN)$을 유한 길이로 만드는 족이면, 이 가군이
$M/(xM+x_1M+\cdots+x_mM)$과 동형임을 관찰하면 된다.

(ii)를 증명하기 위해 $\mathfrak{a}$를 $M$의 소멸자로 두자.
$M/xM$의 소멸자를 포함하는 소 아이디얼들은
$\mathfrak{a}+Ax$를 포함하는 소 아이디얼들이다
($0_{\mathrm I}$, 1.7.5). 따라서 어느 $\mathfrak{p}_i$도
$\mathfrak{a}+Ax$를 포함할 수 없으며, $\dim(M)$의 정의 (16.1.7)에
따라 $\mathfrak{a}+Ax$를 포함하는 소 아이디얼의 모든 사슬은 길이가
$\leq\dim(M)-1$이다.

\oldpage[0\textsubscript{IV-1}]{27}

마지막으로 (iii)을 증명하기 위해 다음 두 완전열에 주목하자.
\[
  0\to{}_xM\to M\xrightarrow{x}xM\to0,\qquad
  0\to xM\to M\to M/xM\to0.
\]
그러면 (16.2.2.2)로부터 주장이 곧바로 따른다.
\end{proof}

(16.2.3)의 증명은 다음 세 단계로 이루어진다.

\begin{env}[16.2.3.2]
\label{0.16.2.3.2-fr}
$\dim(M)\leq d(M)$이다.
\end{env}

$d(M)$에 관해 귀납하자. $d(M)=0$인 경우는 자명하다. 이제
$d(M)\geq1$이라 하고,
$\dim(A/\mathfrak{p}_0)=\dim(M)$을 만족하는
$\mathfrak{p}_0\in\operatorname{Supp}(M)$을 택하자. 그러면
$\mathfrak{p}_0$는 $\operatorname{Supp}(M)$의 \emph{극소} 원소이고,
따라서 $M$의 연관 소 아이디얼이다 (Bourbaki,
\emph{Alg.~comm.}, chap.~IV, \S~1, no~4, th.~2). 그러므로 $M$은
$A/\mathfrak{p}_0$와 동형인 부분가군 $N$을 포함한다
(\emph{loc.~cit.}, \S~1, no~1). (16.2.2.2)에 의해
$d(M)\geq d(N)$이므로 $\dim(N)\leq d(N)$임을 증명하면 충분하다.
따라서
$\mathfrak{p}_0\subset\mathfrak{p}_1\subset\cdots
\subset\mathfrak{p}_n$을 $A$의 서로 다른 소 아이디얼들로 이루어진
사슬이라 하고
$n\leq d(N)$임을 보이자. $n=0$이면 자명하다. 그렇지 않으면
$x\notin\mathfrak{p}_0$인
$x\in\mathfrak{p}_1\cap\mathfrak{r}$가 존재한다. 실제로
$\mathfrak{p}_0\supset\mathfrak{p}_1\cap\mathfrak{r}$이면
$\mathfrak{p}_0$가 $\mathfrak{p}_1$을 포함하지 않으므로
$\mathfrak{p}_0\supset\mathfrak{r}$가 되어 $\mathfrak{p}_0$가
극대가 되는 모순이 생긴다. $i\geq1$이면
$\mathfrak{p}_i\in\operatorname{Supp}(N/xN)$이므로
$n-1\leq\dim(N/xN)$이다 (16.1.7). $x$의 선택에 의해
${}_xN=0$이므로 (16.2.3.1, (iii))에서
$d(N/xN)\leq d(N)-1$을 얻는다. 이제 귀납가정으로
$d(N/xN)=\dim(N/xN)\geq n-1$이므로
$n-1\leq d(N)-1$이다.

\begin{env}[16.2.3.3]
\label{0.16.2.3.3-fr}
$d(M)\leq s(M)$이다.
\end{env}

실제로 $(x_i)_{1\leq i\leq n}$을 $\mathfrak{r}$의 원소들로 이루어진
족이라 하고 $\mathfrak{a}=x_1A+\cdots+x_nA$라 두었을 때
$M/\mathfrak{a}M$이 유한 길이라고 하자. 그러면
$M/\mathfrak{a}M$의 연관 소 아이디얼들은 극대이다 (Bourbaki,
\emph{Alg.~comm.}, chap.~IV, \S~2, no~5, prop.~7). 또한 이들은
$\mathfrak{q}=\mathfrak{a}+
(\mathfrak{r}\cap\operatorname{Ann}(M))$을 포함하는 유일한 소
아이디얼들이므로 $A/\mathfrak{q}$는 아르틴 환이고
(\emph{loc.~cit.}, prop.~9), $\mathfrak{q}$는 $A$의 정의 아이디얼이다.
한편 $\mathfrak{q}^mM/\mathfrak{q}^{m+1}M
=\mathfrak{a}^mM/\mathfrak{a}^{m+1}M$이고, $M$이 $r$개의 원소
$z_j$로 생성되면
$\mathfrak{a}^mM/\mathfrak{a}^{m+1}M$은
$\sum_i\nu_i=m$인
$x_1^{\nu_1}x_2^{\nu_2}\cdots x_n^{\nu_n}z_j$들의 표준상들로
생성되는 $(A/\mathfrak{q})$-가군이다. 따라서 이 가군은
$r\binom{m+n-1}{n-1}$개의 원소로 생성되고, 그 길이는 어떤 상수
$a$에 대하여 $\leq ar\binom{m+n-1}{n-1}$이다. 이로부터
$P_{\mathfrak{q}}(M,n)$의 차수가 $\leq n$임이 곧바로 따르므로
$d(M)\leq s(M)$이다.

\begin{env}[16.2.3.4]
\label{0.16.2.3.4-fr}
$s(M)\leq\dim(M)$이다.
\end{env}

먼저 $n=\dim(M)$은 (16.2.3.2)에 의해 유한함에 주목하자. $n$에
관해 귀납한다. $n=0$이면 $M$은 유한 길이이므로 (16.1.10) 주장은
자명하다. 이제 $n\geq1$이라 하고, $\mathfrak{p}_i$
$(1\leq i\leq m)$를
$\dim(A/\mathfrak{p}_i)=n$을 만족하는
$\operatorname{Supp}(M)$의 소 아이디얼들이라 하자. 이 정의로부터
$\mathfrak{p}_i$들은 $\operatorname{Supp}(M)$의 극소 원소들이므로
유한 개뿐이다 (Bourbaki, \emph{Alg.~comm.}, chap.~IV, \S~1,
no~4, th.~2). $n\geq1$이므로 $\mathfrak{p}_i$들은 극대가 아니며,
따라서 모든 $i$에 대하여 $x\notin\mathfrak{p}_i$인
$x\in\mathfrak{r}$가 존재한다 (Bourbaki, \emph{Alg.~comm.},
chap.~II, \S~1, no~1, prop.~2). (16.2.3.1)에 의해
$s(M)\leq s(M/xM)+1$이고
$\dim(M)\geq\dim(M/xM)+1$이다. 귀납가정에서
$s(M/xM)\leq\dim(M/xM)$이므로 결론이 따른다.

\begin{corollary}[16.2.4]
\label{0.16.2.4-fr}
(16.2.3)의 가정 아래
\begin{equation}
\label{0.16.2.4.1-fr}
\tag{16.2.4.1}
  \dim_{\widehat{A}}(\widehat{M})=\dim_A(M).
\end{equation}
\end{corollary}

\oldpage[0\textsubscript{IV-1}]{28}

\begin{proof}
실제로 $M/\mathfrak{r}^nM$과
$\widehat{M}/\mathfrak{r}^n\widehat{M}$은 동형이므로
$d(\widehat{M})=d(M)$이다.
\end{proof}

\begin{corollary}[16.2.5]
\label{0.16.2.5-fr}
뇌터 반국소환 $A$의 차원은 유한하며, $A$의 정의 아이디얼을 생성하는
$A$의 근기 원소들의 최소 개수와 같다.
\end{corollary}

\begin{proof}
이는 $M=A$에 대한 등식 $s(M)=\dim(M)$이다.
\end{proof}

특히 다음을 얻는다.

\begin{corollary}[16.2.6]
\label{0.16.2.6-fr}
$A$를 뇌터 국소환, $\mathfrak{m}$을 그 극대 아이디얼,
$k=A/\mathfrak{m}$을 그 잉여체라 하자. 그러면
\begin{equation}
\label{0.16.2.6.1-fr}
\tag{16.2.6.1}
  \dim(A)\leq\operatorname{rg}_k(\mathfrak{m}/\mathfrak{m}^2).
\end{equation}
\end{corollary}

\begin{proof}
실제로 $x_i$ $(1\leq i\leq n)$가 $\mathfrak{m}$의 원소들이고 그
$\mathfrak{m}^2$을 법으로 한 잉여류들이 $k$-벡터공간
$\mathfrak{m}/\mathfrak{m}^2$의 기저를 이루면, $x_i$들이
$\mathfrak{m}$을 생성함이 알려져 있다 (Bourbaki,
\emph{Alg.~comm.}, chap.~II, \S~3, no~2, prop.~4의 cor.~2).
\end{proof}

\begin{proposition}[16.2.7]
\label{0.16.2.7-fr}
$k$를 체, $B$를 양의 차수의 유한형 등급 $k$-대수라 하자. $B$는
차수 1의 동차 원소들로 생성되고 $B_0=k$라 하자. $E$를 유한형 등급
$B$-가군이라 하자. 그러면 (III, 2.5.4)에 의해 충분히 큰 $n$에
대하여 $\operatorname{rg}_k(E_n)=r(n)$은
$\Phi(n)-\Phi(n-1)$ 꼴이며, 여기서 $\Phi$는 차수가 $d$인 $n$의
다항식이다. 이때 $B$에는 차수가 $\geq1$인 $d$개의 동차 원소로
이루어진 족 $(t_i)_{1\leq i\leq d}$가 존재하여
$E/(\sum_i t_iE)$는 $k$ 위에서 유한 차원이다.
\end{proposition}

\begin{proof}
$\mathfrak{q}=\bigoplus_{n\geq1}B_n$을 $B$의 극대 아이디얼이라 하고,
각 $n$에 대하여 $E'_n=\bigoplus_{h\geq n+1}E_h$를 $E$의
부분-$B$-가군이라 하자. $B-\mathfrak{q}$의 각 원소는 아르틴
$B$-가군 $E/E'_n$에서 단사인 스칼라곱 자기준동형을 정하고, 따라서
이 자기준동형은 전단사이다 (Bourbaki, \emph{Alg.}, chap.~VIII,
\S~2, no~2, lemme~3). 그러므로
$(E/E'_n)_{\mathfrak{q}}$는 $E/E'_n$과 동일시된다.
$B_{\mathfrak{q}}$의 잉여체가 $k$이므로 가정에서
$\operatorname{long}_{B_{\mathfrak{q}}}
((E/E'_n)_{\mathfrak{q}})$이 $n$의 차수 $d$인 다항식임을 얻는다.
한편 $E$는 유한형
$B$-가군이고 $B$는 차수 1의 동차 원소들로 생성되므로 $E$ 위의
여과 $(E'_n)$은 $\mathfrak{q}$-좋은 여과이다 (Bourbaki,
\emph{Alg.~comm.}, chap.~III, \S~1, no~3, lemme~1). 따라서
$E_{\mathfrak{q}}$ 위의 여과 $((E'_n)_{\mathfrak{q}})$는
$\mathfrak{q}B_{\mathfrak{q}}$-좋은 여과이다.
$B_{\mathfrak{q}}$가 뇌터 국소환이므로 (16.2.3)과 (16.2.2.1)에서
$\dim(E_{\mathfrak{q}})=d$를 얻는다. 이제 $d$에 관해 귀납하자.
$d=0$이면 주장은 자명하다. $d\geq1$이라 하자.
$\operatorname{Ass}(E)$의 극소 소 아이디얼 $\mathfrak{p}_i$
$(1\leq i\leq r)$들은 \emph{등급 아이디얼}이고 (Bourbaki,
\emph{Alg.~comm.}, chap.~IV, \S~3, no~1, prop.~1), $d\geq1$이므로
$\mathfrak{q}$는 $\mathfrak{p}_i$들의 합집합과 같지 않다. 따라서
어느 $\mathfrak{p}_i$에도 속하지 않는 \emph{동차} 원소
$t_1\in\mathfrak{q}$가 존재한다 (Bourbaki,
\emph{Alg.~comm.}, chap.~III, \S~1, no~4, prop.~8).
소 아이디얼 $(\mathfrak{p}_i)_{\mathfrak{q}}$들은
$\operatorname{Ass}(E_{\mathfrak{q}})$의 극소 원소들이므로
(Bourbaki, \emph{Alg.~comm.}, chap.~IV, \S~1, no~2, prop.~5),
$\dim(E_{\mathfrak{q}}/t_1E_{\mathfrak{q}})=d-1$이다 (16.3.4).
이제 등급 $B$-가군 $E/t_1E$에 귀납가정을 적용하면
결론이 따른다.
\end{proof}

\subsection{뇌터 국소환의 매개변수계}
\label{subsection:0.16.3-fr}

\begin{proposition}[16.3.1]
\label{0.16.3.1-fr}
$A$를 뇌터 환, $\mathfrak{p}$를 $A$의 소 아이디얼, $n$을 정수라
하자. 다음 조건들은 동치이다.
\begin{enumerate}
  \item[$a)$] $\operatorname{ht}(\mathfrak{p})=
  \dim(A_{\mathfrak{p}})\leq n$이다.
  \item[$b)$] $A$의 원소 $x_i$ $(1\leq i\leq n)$가 존재하여,
  $\mathfrak{p}$는 $x_i$들이 생성하는 아이디얼 $\mathfrak{a}$를
  포함하는 소 아이디얼들의 집합에서 극소이다. 다시 말해
  $V(\mathfrak{p})$는 $V(\mathfrak{a})$의 기약성분이다.
\end{enumerate}
\end{proposition}

\begin{proof}
조건 $b)$이면 $\mathfrak{a}A_{\mathfrak{p}}$는
$A_{\mathfrak{p}}$의 정의 아이디얼이므로 (16.2.5)에 의해 $a)$가
따른다. 역으로 $a)$가 성립하면 $A_{\mathfrak{p}}$에는
$s\in A-\mathfrak{p}$인 $n$개의 원소 $x_i/s$로 생성되는 정의
아이디얼 $\mathfrak{b}$가 존재한다. $A_{\mathfrak{p}}$의 소
아이디얼들과 $\mathfrak{p}$에 포함되는 $A$의 소 아이디얼들 사이의
전단사 대응에 의해, $x_i$들이 생성하는 $A$의 아이디얼
$\mathfrak{a}$는 $b)$를 만족한다.
\end{proof}

\oldpage[0\textsubscript{IV-1}]{29}

$n=1$이면 다음 특수한 경우를 얻는다.

\begin{corollary}[16.3.2]
\label{0.16.3.2-fr}
\textup{(``주 아이디얼 정리'').}
뇌터 환의 소 아이디얼이 높이 $\leq1$일 필요충분조건은 그 소
아이디얼이 어떤 적당한 주 아이디얼을 포함하는 소 아이디얼들의
집합에서 극소인 것이다.
\end{corollary}

\begin{corollary}[16.3.3]
\label{0.16.3.3-fr}
\textup{(아르틴--테이트).}
$A$를 뇌터 정역이라 하자. 다음 조건들은 동치이다.
\begin{enumerate}
  \item[$a)$] $A$는 차원이 $\leq1$인 반국소환이다.
  \item[$b)$] $A_f$가 체가 되는 $f\neq0$인 원소 $f\in A$가
  존재한다.
\end{enumerate}
\end{corollary}

\begin{proof}
조건 $a)$는 $A$가 $0$이 아닌 소 아이디얼
$\mathfrak{m}_i$를 유한 개만 가지며 이 아이디얼들이 모두
극대라는 것과 같다. 여기서 $(0)$은 소 아이디얼이다. $A$가
정역이므로 $\mathfrak{m}_i$들의 곱은 $0$이 아니며, 이 곱의
$0$이 아닌 원소 $f$는 모든 $\mathfrak{m}_i$에 속한다. 따라서
정역 $A_f$의 소 아이디얼은 $(0)$뿐이고, 다시 말해 $A_f$는 체이다.
(증명의 이 부분에서는 $A$가 뇌터라는 사실을 사용하지 않는다.)

이제 $b)$에서 $a)$가 따름을 보이자. 가정 $b)$는 $A$의 모든
$0$이 아닌 소 아이디얼 $\mathfrak{p}$가 $f$를 포함한다는 뜻이다.
$\mathfrak{p}_i$ $(1\leq i\leq n)$를 $f$를 포함하는 소
아이디얼들 가운데 극소인 것들이라 하자. $A/fA$가 뇌터이므로 이들은
유한 개이다. $\mathfrak{p}_i$들이 \emph{극대} 아이디얼임을 보이면
충분하다. 실제로 그러면 모든 $0$이 아닌 소 아이디얼은 어떤
$\mathfrak{p}_i$를 포함하므로 그 $\mathfrak{p}_i$와 같다.
어떤 $\mathfrak{p}_i$를 포함하면서 그와 다른 극대 아이디얼
$\mathfrak{m}$이 존재한다고 가정하자. 그러면 $\mathfrak{m}$은
$\mathfrak{p}_i$들의 합집합에 포함되지 않으므로 (Bourbaki,
\emph{Alg.~comm.}, chap.~II, \S~1, no~1, prop.~2), 어느
$\mathfrak{p}_i$에도 속하지 않는 $g\in\mathfrak{m}$이 존재한다.
$\mathfrak{q}$를 $g$를 포함하는 소 아이디얼들 가운데 극소인 것
하나로 택하자. ``주 아이디얼 정리'' (16.3.2)에 의해
$\mathfrak{q}$의 높이는 $1$이고, 따라서 $\mathfrak{q}\neq0$이다.
그러므로 $\mathfrak{q}$는 $f$를 포함하며, 따라서 어떤
$\mathfrak{p}_i$도 포함한다. 한편 $g\in\mathfrak{q}$이고 $g$는
어느 $\mathfrak{p}_i$에도 속하지 않으므로 $\mathfrak{q}$는 그
$\mathfrak{p}_i$와 다르다. 따라서 $\mathfrak{q}$의 높이는
$\geq2$가 되어 모순이다.
\end{proof}

\begin{proposition}[16.3.4]
\label{0.16.3.4-fr}
$A$를 뇌터 반국소환, $\mathfrak{r}$을 그 근기, $M$을 유한형
$A$-가군이라 하자. $\mathfrak{p}_i$들을
$\dim(A/\mathfrak{p}_i)=\dim(M)$을 만족하는
$\operatorname{Supp}(M)$의 원소들이라 하자. 그러면
$\mathfrak{p}_i$들은 $\operatorname{Ass}(M)$의 극소 원소들이고
$\dim(M/\mathfrak{p}_iM)=\dim(M)$이다. 모든
$x\in\mathfrak{r}$에 대하여
\begin{equation}
\label{0.16.3.4.1-fr}
\tag{16.3.4.1}
  \dim(M/xM)\geq\dim(M)-1.
\end{equation}
또한 (16.3.4.1)의 양변이 같을 필요충분조건은 $x$가 어느
$\mathfrak{p}_i$에도 속하지 않는 것이다.
\end{proposition}

\begin{proof}
첫 번째 주장은 곧바로 따른다. 실제로 $\mathfrak{p}_i$들은 정의에
따라 $\operatorname{Supp}(M)$의 극소 원소들이고 (16.1.6), 이들은
$\operatorname{Ass}(M)$의 극소 원소들이기도 하다 (Bourbaki,
\emph{Alg.~comm.}, chap.~IV, \S~1, no~4, th.~2). 또한
$\operatorname{Supp}(M/\mathfrak{p}_iM)$은
$\mathfrak{p}_i$를 포함하는 소 아이디얼들의 집합이므로
($0_{\mathrm I}$, 1.7.5),
$\dim(M/\mathfrak{p}_iM)=\dim(A/\mathfrak{p}_i)=\dim(M)$이다.

$N=M/xM$이라 두자. $x_1,\ldots,x_n\in\mathfrak{r}$이고
$N/(x_1N+\cdots+x_nN)$이 유한 길이면, 이는
$M/(xM+x_1M+\cdots+x_nM)$이 유한 길이라는 뜻이다. 따라서
(16.2.3)에 의해 부등식 (16.3.4.1)을 얻는다.

$x$가 어느 $\mathfrak{p}_i$에도 속하지 않을 때 (16.3.4.1)의
양변이 같다는 사실은 (16.3.4.1)과 (16.2.3.1)에서 따른다. 역으로
\oldpage[0\textsubscript{IV-1}]{30}
$\dim(M/xM)=\dim(M)-1$이면 어느 $\mathfrak{p}_i$도
$\operatorname{Supp}(M/xM)$에 속할 수 없다. 이 지지집합의 소
아이디얼들은 $\operatorname{Ann}(M)+Ax$를 포함하는 것들이다.
$\mathfrak{p}_i$들은 $\operatorname{Ann}(M)$을 포함하므로 $x$를
포함할 수 없다.
\end{proof}

\begin{corollary}[16.3.5]
\label{0.16.3.5-fr}
(16.3.4)의 표기 아래 모든 아이디얼
$\mathfrak{a}\subset\mathfrak{p}_i$에 대하여
$\dim(M/\mathfrak{a}M)=\dim(M)$이다. 또한
$N=M/\mathfrak{a}M$이라 두면, 관계
$\dim(M/xM)=\dim(M)-1$은
$\dim(N/xN)=\dim(N)-1$을 함의한다.
\end{corollary}

\begin{proof}
실제로 $\operatorname{Supp}(N)$은
$\mathfrak{a}+\operatorname{Ann}(M)$을 포함하는 소 아이디얼들의
집합이다. 따라서 $\mathfrak{p}_i\in\operatorname{Supp}(N)$이고,
$\dim(N)\leq\dim(M)=\dim(A/\mathfrak{p}_i)$이므로
$\dim(A/\mathfrak{p}_i)=\dim(N)$이다. 또한
$\dim(A/\mathfrak{q})=\dim(N)$을 만족하는
$\mathfrak{q}\in\operatorname{Supp}(N)$은 $\mathfrak{p}_j$들 중
일부이다. $x$가 어느 $\mathfrak{p}_j$에도 속하지 않으므로
(16.3.4)에 의해 $\dim(N/xN)=\dim(N)-1$이다.
\end{proof}

\begin{definition}[16.3.6]
\label{0.16.3.6-fr}
$A$를 뇌터 반국소환, $\mathfrak{r}$을 그 근기, $M$을 유한형
$A$-가군이라 하고 $n=\dim(M)$이라 두자. $\mathfrak{r}$의 $n$개
원소로 이루어진 족 $(x_i)_{1\leq i\leq n}$ 중에서
$M/(x_1M+\cdots+x_nM)$이 유한 길이인 것을 \emph{$M$의
매개변수계}라 한다.
\end{definition}

그러한 매개변수계가 항상 존재함은 (16.2.3)에서 보았다.

(16.2.3.3)의 증명에서 다음을 보았다.
$\mathfrak{J}=\sum_{i=1}^n x_iA$이고 $M/\mathfrak{J}M$이 유한
길이면 $\mathfrak{q}=\mathfrak{J}+\operatorname{Ann}(M)$은
$A$의 \emph{정의 아이디얼}이다. 역으로 이것이 성립하면
$M/\mathfrak{J}M$은 아르틴 환 $A/\mathfrak{q}$ 위의 유한형
가군이므로 유한 길이이다. 따라서 $(x_i)_{1\leq i\leq n}$이
\emph{$M$의} 매개변수계라고 하는 것과
\emph{$A/\operatorname{Ann}(M)$의} 매개변수계라고 하는 것은 같다.
달리 말해, $A/\operatorname{Ann}(M)$에서의 상들이 이 환의
매개변수계를 이룬다. $(x_i)_{1\leq i\leq n}$이 $M$의
매개변수계이면 모든 정수 $k\geq1$에 대하여 $(x_i^k)$도 그러하다.
실제로 앞의 설명에 따라 $M=A$인 경우만 고려하면 되고,
$\mathfrak{J}$가 $A$의 정의 아이디얼이면
$\sum_{i=1}^n x_i^kA$는 $\mathfrak{J}^{kn}$을 포함하므로 역시 정의
아이디얼이다.

\begin{proposition}[16.3.7]
\label{0.16.3.7-fr}
$A$를 뇌터 반국소환, $\mathfrak{r}$을 그 근기,
$x_1,\ldots,x_k$를 $\mathfrak{r}$의 원소들, $M$을 유한형
$A$-가군이라 하자. 그러면
\begin{equation}
\label{0.16.3.7.1-fr}
\tag{16.3.7.1}
  \dim(M/(x_1M+\cdots+x_kM))\geq\dim(M)-k.
\end{equation}
또한 (16.3.7.1)의 양변이 같을 필요충분조건은
$\mathfrak{r}$의 원소 $x_i$ $(k+1\leq i\leq n=\dim(M))$가 존재하여
$(x_i)_{1\leq i\leq n}$이 $M$의 매개변수계를 이루는 것이다.
\end{proposition}

\begin{proof}
부등식 (16.3.7.1)은 $k$에 관한 귀납과 (16.3.4.1)에서 따른다.
(16.3.7.1)의 양변이 같고 $x_{k+1},\ldots,x_n$이
$M/(x_1M+\cdots+x_kM)$의 매개변수계이면
$M/(x_1M+\cdots+x_kM+x_{k+1}M+\cdots+x_nM)$이 유한 길이임은
명백하다. 따라서 $(x_i)_{1\leq i\leq n}$은 $M$의 매개변수계이다.
역으로 이 성질을 갖는 $x_i$ $(k+1\leq i\leq n)$가 존재한다고 하고
$N=M/(x_1M+\cdots+x_kM)$이라 두자. 그러면
$N/(x_{k+1}N+\cdots+x_nN)$이 유한 길이이므로
$\dim(N)\leq n-k$이고, (16.3.7.1)에 의해 양변이 같다.
\end{proof}

\oldpage[0\textsubscript{IV-1}]{31}

\begin{proposition}[16.3.8]
\label{0.16.3.8-fr}
$A$를 뇌터 반국소환, $X=\operatorname{Spec}(A)$를 그 스펙트럼,
$\mathfrak{a}\neq A$를
$\operatorname{codim}(V(\mathfrak{a}),X)=r>0$인 $A$의
아이디얼이라 하자. 그러면 $\mathfrak{a}$에는 $A$의 매개변수계의
일부를 이루는 $r$개의 원소 $x_1,\ldots,x_r$가 존재하여
\[
  \operatorname{codim}(V(\mathfrak{a}),
  V(Ax_1+\cdots+Ax_r))=0.
\]
\end{proposition}

\begin{proof}
$\mathfrak{p}_i$ $(1\leq i\leq m)$를 환 $A$의 극소 소
아이디얼들, $\mathfrak{q}_j$ $(1\leq j\leq n)$를
$\mathfrak{a}$를 포함하는 소 아이디얼들 가운데 극소인 것들이라
하자. (14.2.1)에 의해
\[
  \operatorname{codim}(V(\mathfrak{a}),X)
  =\inf_j(\operatorname{codim}(V(\mathfrak{q}_j),X))
  =\inf_j(\dim(A_{\mathfrak{q}_j}))
\]
이다 (16.1.3.2). 가정 $r>0$으로부터 $\mathfrak{a}$는 어느
$\mathfrak{p}_i$에도 포함되지 않는다. 따라서 어느
$\mathfrak{p}_i$에도 속하지 않는
$x\in\mathfrak{a}$가 존재한다 (Bourbaki, \emph{Alg.~comm.},
chap.~II, \S~1, no~1, prop.~2). 앞에서와 같이
\[
  \operatorname{codim}(V(\mathfrak{a}),V(Ax))
  =\inf_j(\dim((A/Ax)_{\mathfrak{q}_j/Ax}))
  =\inf_j(\dim(A_{\mathfrak{q}_j}/A_{\mathfrak{q}_j}x)).
\]
그런데 $x/1$은 $A_{\mathfrak{q}_j}$의 어느 극소 소 아이디얼에도
속하지 않으므로
$\dim(A_{\mathfrak{q}_j}/A_{\mathfrak{q}_j}x)
=\dim(A_{\mathfrak{q}_j})-1$이다 (16.3.4). 따라서
$\operatorname{codim}(V(\mathfrak{a}),V(Ax))=r-1$이다. 같은 이유로
$\dim(A/Ax)=\dim(A)-1$이다.

이제 $r$에 관해 귀납하자. $r>1$이면 $A'=A/Ax$에는
$\mathfrak{a}/Ax$에 속하고 이 환의 매개변수계의 일부를 이루는
$r-1$개의 원소 $x'_i$ $(2\leq i\leq r)$가 존재하여
\[
  \operatorname{codim}(V(\mathfrak{a}/Ax),
  V(A'x'_2+\cdots+A'x'_r))=0.
\]
$(x'_i)_{2\leq i\leq s}$가 $A'$의 매개변수계를 이루도록
$A'$의 원소들로 이루어진 족 $(x'_i)_{r+1\leq i\leq s}$를 택하자.
$2\leq i\leq s$인 각 $i$에 대하여 $x_i$를 $A$에서 $x'_i$의 한
대표원으로 택하되, $2\leq i\leq r$이면
$x_i\in\mathfrak{a}$가 되게 한다. 그러면
$A/(Ax+Ax_2+\cdots+Ax_s)$가 유한 길이임은 명백하다.
$\dim(A)=s$이므로 $x_1=x$와 $i\geq2$인 $x_i$들은 명제의 조건을
만족한다.
\end{proof}

\begin{proposition}[16.3.9]
\label{0.16.3.9-fr}
$A$, $B$를 두 뇌터 국소환, $\mathfrak{m}$을 $A$의 극대 아이디얼,
$k$를 그 잉여체, $\varphi:A\to B$를 국소 준동형이라 하자.
\begin{enumerate}
  \item[\rm{(i)}] 다음 부등식이 성립한다.
  \begin{equation}
  \label{0.16.3.9.1-fr}
  \tag{16.3.9.1}
    \dim(B)\leq\dim(A)+\dim(B\otimes_A k).
  \end{equation}
  \item[\rm{(ii)}] $M\neq0$을 유한형 $A$-가군,
  $N\neq0$을 유한형 $B$-가군이라 하자. 그러면
  \begin{equation}
  \label{0.16.3.9.2-fr}
  \tag{16.3.9.2}
    \dim_B(M\otimes_A N)
    \leq\dim_A(M)+\dim_{B\otimes_A k}(N\otimes_A k).
  \end{equation}
\end{enumerate}
\end{proposition}

\begin{proof}
$m=\dim(A)$라 하고 $(s_i)_{1\leq i\leq m}$을 $A$의 매개변수계라
하자 (16.3.6). $\mathfrak{a}=\sum_i s_iA$라 두면
$A/\mathfrak{a}$는 유한 길이 $A$-가군이다. 따라서 환
$A/\mathfrak{a}$는 아르틴이고, 그 극대 아이디얼
$\mathfrak{m}/\mathfrak{a}$는 멱영이다. $B/\mathfrak{a}B$의
아이디얼 $\mathfrak{m}B/\mathfrak{a}B$는
$(\mathfrak{m}/\mathfrak{a})\otimes_A B$의 상이므로 역시 멱영이다.
따라서 두 환의 스펙트럼이 같은 기저공간을 가지므로
$\dim(B/\mathfrak{a}B)=\dim(B/\mathfrak{m}B)$이다. 한편
$B/\mathfrak{a}B=B/(s_1B+\cdots+s_mB)$이고, $s_i$들의 $B$에서의
상들은 $B$의 극대 아이디얼에 속한다. 그러므로 (16.3.7)에 의해
\[
  \dim(B)\leq m+\dim(B\otimes_A k).
\]
이는 바로 (16.3.9.1)이다.

(16.3.9.2)를 증명하자. $\mathfrak{r}$과 $\mathfrak{s}$를 각각
$A$에서 $M$의 소멸자와 $B$에서 $N$의 소멸자라 하자.
\oldpage[0\textsubscript{IV-1}]{32}
정의 (16.1.7)에 의해 $\dim_A(M)=\dim(A/\mathfrak{r})$이고
$\dim_B(N)=\dim(B/\mathfrak{s})$이다. 한편
$\operatorname{Supp}(M\otimes_A N)$은 $\operatorname{Spec}(B)$의
닫힌 부분집합이며 (I, 9.1.13) 다음과 동일시된다.
\[
  \operatorname{Supp}(M)\times_{\operatorname{Spec}(A)}
  \operatorname{Supp}(N)
  =\operatorname{Spec}((B/\mathfrak{s})\otimes_A(A/\mathfrak{r}))
  =\operatorname{Spec}(B/(\mathfrak{s}+\mathfrak{r}B)).
\]
$N\otimes_A k$는 유한형 $B$-가군이므로 (16.1.9)에 의해
$\dim_{B\otimes_A k}(N\otimes_A k)=\dim_B(N\otimes_A k)$이다.
또 위와 같은 논증으로
$\operatorname{Supp}(N\otimes_A k)
=\operatorname{Spec}(B/(\mathfrak{s}+\mathfrak{m}B))$이다.
(16.3.9.1)을 $A/\mathfrak{r}$과
$B/(\mathfrak{s}+\mathfrak{r}B)$에 적용하면
\[
  \dim(B/(\mathfrak{s}+\mathfrak{r}B))
  \leq\dim(A/\mathfrak{r})+
  \dim(B/(\mathfrak{s}+\mathfrak{m}B))
\]
을 얻으며, 이는 바로 (16.3.9.2)이다.
\end{proof}

\begin{corollary}[16.3.10]
\label{0.16.3.10-fr}
(16.3.9)의 가정 아래 $\mathfrak{m}B$가 $B$의 정의 아이디얼이면
$\dim(B)\leq\dim(A)$이다. 더 나아가 $A$가 정역이고
$\dim(A)=\dim(B)$이면 $\varphi$는 단사이다.
\end{corollary}

\begin{proof}
첫 번째 주장은 이때 $\dim(B/\mathfrak{m}B)=0$이므로 (16.3.9)에서
따른다. 또한 $A$를 $\varphi(A)$로 바꿀 수 있으므로
$\dim(B)\leq\dim(\varphi(A))$이다. 만일
$\mathfrak{a}=\operatorname{Ker}(\varphi)\neq0$이고 $A$가 정역이면
$\dim(\varphi(A))=\dim(A/\mathfrak{a})<\dim(A)$이다 (16.1.2.2).
따라서 $\dim(A)=\dim(B)$일 수 있으려면
$\mathfrak{a}=0$이어야 한다.
\end{proof}

\subsection[깊이와 여깊이]{깊이와
여깊이\protect\footnote{제III장의 제3부에서 깊이의 개념을 코호몰로지
관점과 더 일반적인 틀에서 전개한다.}}
\label{subsection:0.16.4-fr}

\begin{proposition}[16.4.1]
\label{0.16.4.1-fr}
$A$를 뇌터 국소환, $\mathfrak{m}$을 그 극대 아이디얼, $M$을 유한형
$A$-가군이라 하자. $\mathfrak{m}$의 원소들로 이루어진 모든
$M$-정칙열 $(x_i)_{1\leq i\leq r}$은 (15.1.7) $M$의 매개변수계의
일부를 이룬다.
\end{proposition}

\begin{proof}
$r$에 관해 귀납하자. $A$-가군
\[
  N=M/(x_1M+\cdots+x_{r-1}M)
\]
을 생각하자. 가정에 의해 $x_1,\ldots,x_{r-1}$은 $M$의
매개변수계의 일부를 이루므로 $\dim(N)=\dim(M)-(r-1)$이다
(16.3.7). 또 가정에 의해 $N$에서 $x_r$에 의한 스칼라곱
자기준동형은 단사이므로 $x_r$은 $N$의 어느 연관 소 아이디얼에도
속하지 않는다. 따라서 (16.3.4)에 의해
$\dim(N/x_rN)=\dim(N)-1$이고, 이는
\[
  \dim(M/(x_1M+\cdots+x_rM))=\dim(M)-r
\]
과 같다. 그러므로 (16.3.7)에서 결론이 따른다.
\end{proof}

$M\neq0$인 $A$-가군에 대하여, $\mathfrak{m}$의 원소들로 이루어진
$M$-정칙열은 따라서 원소를 \emph{많아야} $\dim(M)$개 갖는다.

\begin{lemma}[16.4.2]
\label{0.16.4.2-fr}
$A$를 뇌터 국소환, $\mathfrak{m}$을 그 극대 아이디얼,
$k=A/\mathfrak{m}$을 그 잉여체, $M$을 $A$-가군이라 하자.
$\mathfrak{m}$의 원소들로 이루어진 $M$-정칙열
$(x_i)_{1\leq i\leq n}$이 극대이기 위한 필요충분조건은
\[
  \operatorname{Hom}_A(k,M/(x_1M+\cdots+x_nM))\neq0
\]
이다.
\end{lemma}

\begin{proof}
$(x_i)$가 극대라는 것은 어떤 $x\in\mathfrak{m}$에 대해서도
$N=M/(x_1M+\cdots+x_nM)$에서 $x$에 의한 스칼라곱 자기준동형이
단사가 아니라는 뜻이다. 동치로, $\mathfrak{m}$은 $N$의 연관 소
아이디얼들의 합집합에 포함된다 (Bourbaki, \emph{Alg.~comm.},
chap.~IV, \S~1, no~1, prop.~2의 cor.~2). 따라서 $\mathfrak{m}$은
그 소 아이디얼들 가운데 하나와 같다. 실제로 $\mathfrak{m}$은
\oldpage[0\textsubscript{IV-1}]{33}
극대 아이디얼이기 때문이다 (Bourbaki, \emph{Alg.~comm.}, chap.~II,
\S~1, no~1, prop.~2). 다시 말해, 소멸자가 $\mathfrak{m}$인
$z\in N$가 존재하며, 따라서 $N$의 부분가군 $Az$는
$A/\mathfrak{m}=k$와 동형이다. 이로부터 보조정리가 따른다.
\end{proof}

\begin{lemma}[16.4.3]
\label{0.16.4.3-fr}
$A$를 뇌터 국소환, $\mathfrak{m}$을 그 극대 아이디얼,
$k=A/\mathfrak{m}$을 그 잉여체, $M$을 $A$-가군,
$(x_i)_{1\leq i\leq n}$을 $\mathfrak{m}$의 원소들로 이루어진
$M$-정칙열이라 하자. 그러면 $A$-가군
\[
  \operatorname{Hom}_A(k,M/(x_1M+\cdots+x_nM))
  \quad\text{과}\quad \operatorname{Ext}^n_A(k,M)
\]
은 서로 동형이고, $n\geq1$이면 이들은
$\operatorname{Ext}^{n-1}_A(k,M/x_1M)$과도 동형이다.
\end{lemma}

\begin{proof}
$n$에 관해 귀납하자. $n=0$이면 명제는 명백하다.
$N=M/x_1M$이라 두자. 열 $(x_i)_{2\leq i\leq n}$은
$N$-정칙열이므로
\[
  \operatorname{Hom}_A(k,M/(x_1M+\cdots+x_nM))
  =\operatorname{Hom}_A(k,N/(x_2N+\cdots+x_nN))
\]
은 귀납가정에 의해 $\operatorname{Ext}^{n-1}_A(k,N)$과 동형이다.
이제 완전열
$0\to M\xrightarrow{x_1}M\to M/x_1M=N\to0$을 생각하자.
이로부터 Ext의 완전열
\begin{equation}
\label{0.16.4.3.1-fr}
\tag{16.4.3.1}
  \operatorname{Ext}^{n-1}_A(k,M)
  \to\operatorname{Ext}^{n-1}_A(k,N)
  \to\operatorname{Ext}^n_A(k,M)
  \xrightarrow{x_1}\operatorname{Ext}^n_A(k,M)
\end{equation}
을 얻는다. 귀납가정에 의해 $\operatorname{Ext}^{n-1}_A(k,M)$은
\[
  \operatorname{Hom}_A(k,M/(x_1M+\cdots+x_{n-1}M))
\]
과 동형이다. 그런데 $(x_i)_{1\leq i\leq n-1}$은 극대
$M$-정칙열이 아니므로 이는 0이다 (16.4.2). 한편
$x_1\in\mathfrak{m}$이므로 임의의 $A$-가군 $T$에 대하여
$A$-가군 $\operatorname{Hom}_A(k,T)$에서 $x_1$에 의한 스칼라곱은
0이고, 따라서 모든 $A$-가군 $\operatorname{Ext}^n_A(k,T)$에서도
$x_1$에 의한 스칼라곱은 0이다. 그러므로 보조정리의 주장들은
완전열 (16.4.3.1)에서 곧바로 따른다.
\end{proof}

\begin{corollary}[16.4.4]
\label{0.16.4.4-fr}
$A$를 뇌터 국소환, $\mathfrak{m}$을 그 극대 아이디얼,
$k=A/\mathfrak{m}$을 그 잉여체, $M\neq0$을 유한형 $A$-가군이라
하자. $\mathfrak{m}$의 원소들로 이루어진 모든 극대 $M$-정칙열은
원소의 개수가 같고, 그 개수는
$\operatorname{Ext}^n_A(k,M)\neq0$이 되는 가장 작은 정수 $n$이다.
\end{corollary}

\begin{definition}[16.4.5]
\label{0.16.4.5-fr}
$A$를 뇌터 국소환, $\mathfrak{m}$을 그 극대 아이디얼, $M$을 유한형
$A$-가군이라 하자. $\mathfrak{m}$의 원소들로 이루어진 $M$-정칙열의
원소 개수의 상한을 $M$의 \emph{깊이}라 하고,
$\operatorname{prof}_A M$ 또는 $\operatorname{prof}(M)$으로 쓴다.
따라서 $\operatorname{prof}(M)=+\infty$일 필요충분조건은 $M=0$이고,
(16.4.1)에 의해
\begin{equation}
\label{0.16.4.5.1-fr}
\tag{16.4.5.1}
  \operatorname{prof}(M)\leq\dim(M)\quad(M\neq0\text{일 때}).
\end{equation}
모든 $k>0$에 대하여 명백히
$\operatorname{prof}(M^k)=\operatorname{prof}(M)$이다.
\end{definition}

\begin{proposition}[16.4.6]
\label{0.16.4.6-fr}
$A$를 뇌터 국소환, $\mathfrak{m}$을 그 극대 아이디얼, $M$을 유한형
$A$-가군이라 하자.
\begin{enumerate}
  \item[\rm{(i)}] $\operatorname{prof}(M)=0$일 필요충분조건은
  $\mathfrak{m}\in\operatorname{Ass}(M)$인 것이다. 동치로, $M$에는
  $A$의 잉여체 $k=A/\mathfrak{m}$와 동형인 부분가군이 존재한다.
  \item[\rm{(ii)}] $\mathfrak{m}$에 속하는 모든 $M$-정칙 원소 $x$에
  대하여
  \begin{equation}
  \label{0.16.4.6.1-fr}
  \tag{16.4.6.1}
    \operatorname{prof}(M/xM)=\operatorname{prof}(M)-1.
  \end{equation}
  \item[\rm{(iii)}] $M\neq0$이면
  \begin{equation}
  \label{0.16.4.6.2-fr}
  \tag{16.4.6.2}
    \operatorname{prof}(M)
    \leq\inf_{\mathfrak{p}\in\operatorname{Ass}(M)}\dim(A/\mathfrak{p})
    \leq\dim(M).
  \end{equation}
  \oldpage[0\textsubscript{IV-1}]{34}
  \item[\rm{(iv)}] 다음 등식이 성립한다.
  $\operatorname{prof}_A(M)=
  \operatorname{prof}_{\widehat{A}}(\widehat{M})$.
\end{enumerate}
\end{proposition}

\begin{proof}
(i) $\operatorname{prof}(M)=0$이라는 것은 $M$에서 $x$에 의한
스칼라곱 자기준동형이 단사인 $x\in\mathfrak{m}$이 하나도 없다는
뜻이다. 다시 말해, 모든 $x\in\mathfrak{m}$은 $M$의 어느 연관 소
아이디얼에 속한다 (Bourbaki, \emph{Alg.~comm.}, chap.~IV, \S~1,
no~1, prop.~2의 cor.~2). 그런데 $\mathfrak{m}$이 $M$의 연관 소
아이디얼들의 합집합이 되려면 그 가운데 하나와 같아야 한다
(Bourbaki, \emph{Alg.~comm.}, chap.~II, \S~1, no~1, prop.~2).

(ii) $(x_i)_{1\leq i\leq m}$이 $\mathfrak{m}$의 원소들로 이루어진
$(M/xM)$-정칙열이면, $x$와 $x_i$들로 이루어진 열은 $M$-정칙열이다.
정의에 의해 $(x_i)_{1\leq i\leq m}$이 극대 $(M/xM)$-정칙열일 때
그 열도 극대이므로 결론을 얻는다.

(iii) $n=\operatorname{prof}(M)$에 관해 귀납하자. $n=0$이면 주장은
명백하다. 다음 보조정리를 사용한다.
\end{proof}

\begin{lemma}[16.4.6.3]
\label{0.16.4.6.3-fr}
$A$를 뇌터 환, $M$을 유한형 $A$-가군, $t\in A$를 $M$-정칙 원소라
하자. 모든 $\mathfrak{p}\in\operatorname{Ass}(M)$에 대하여
$\mathfrak{p}+At$를 포함하는 소 아이디얼들 가운데 극소인 것은
모두 $M/tM$에 연관된다.
\end{lemma}

\begin{proof}
완전열
\[
  0\to M'\to M\to M''\to0
\]
이 존재하여 $\operatorname{Ass}(M')=\{\mathfrak{p}\}$이고
$\operatorname{Ass}(M'')=
\operatorname{Ass}(M)-\{\mathfrak{p}\}$임이 알려져 있다
(Bourbaki, \emph{Alg.~comm.}, chap.~IV, \S~1, no~1, prop.~4).
$t$는 $M$-정칙이므로 $M'$-정칙이고 $M''$-정칙이기도 하다
(\emph{loc.~cit.}, \S~1, no~1, prop.~2의 cor.~2). 따라서 (15.1.18)에
의해 완전열
\[
  0\to M'/tM'\to M/tM\to M''/tM''\to0
\]
을 얻는다. 그러므로
$\operatorname{Ass}(M'/tM')\subset\operatorname{Ass}(M/tM)$이다.
한편 $\mathfrak{p}$를 포함하는 $A$의 소 아이디얼들은
$\operatorname{Supp}(M')$의 점들이다
(\emph{loc.~cit.}, \S~1, no~3, prop.~7). 따라서
$\mathfrak{p}+At$를 포함하는 소 아이디얼들은
$M'/tM'=M'\otimes_A(A/tA)$의 지지의 점들이다
($0_{\mathrm I}$, 1.7.5). 이 소 아이디얼 집합의 극소 원소는
$\operatorname{Ass}(M'/tM')$에 속하고
(Bourbaki, \emph{loc.~cit.}, \S~1, no~3, prop.~7의 cor.~1), 따라서
더욱이 $\operatorname{Ass}(M/tM)$에 속한다.
\end{proof}

\begin{proof}[명제 (16.4.6)의 증명 계속]
보조정리가 증명되었으므로 (iii)의 증명으로 돌아가자.
$x\in\mathfrak{m}$을 $M$-정칙 원소라 하자. (ii)에 의해
$\operatorname{prof}(M/xM)=n-1$이다. 모든
$\mathfrak{p}\in\operatorname{Ass}(M)$에 대하여 보조정리에 의해
$M/xM$에 연관되고 $\mathfrak{p}+xA$를 포함하는 소 아이디얼
$\mathfrak{p}'$이 존재한다. 그런데 $x$는 $M$-정칙이므로
$x\notin\mathfrak{p}$이다. 따라서 $\mathfrak{p}'\neq\mathfrak{p}$이고
$\dim(A/\mathfrak{p}')\leq\dim(A/\mathfrak{p})-1$이다.
귀납가정에 의해 $n-1\leq\dim(A/\mathfrak{p}')$이므로, 모든
$\mathfrak{p}\in\operatorname{Ass}(M)$에 대하여
$n\leq\dim(A/\mathfrak{p})$이다.

(iv) $M\neq0$인 경우로 한정해도 된다. 표준적인 동형 아래
$\widehat{M}=M\otimes_A\widehat{A}$임을 상기하자
($0_{\mathrm I}$, 7.3.3). $\widehat{A}$는 평탄 $A$-가군이므로 표준적인
동형
\[
  \operatorname{Ext}^i_A(k,M)\otimes_A\widehat{A}
  \xrightarrow{\sim}
  \operatorname{Ext}^i_{\widehat{A}}
  (k,M\otimes_A\widehat{A})
\]
들이 존재한다. 실제로 표준적인 동형 아래
$k=k\otimes_A\widehat{A}$이다 ($0_{\mathrm{III}}$, 12.3.5).
그러므로 (iv)는 (16.4.4)와 $\widehat{A}$가 충실 평탄 $A$-가군이라는
사실에서 따른다 ($0_{\mathrm I}$, 7.3.5).
\end{proof}

\begin{env}[16.4.6.4]
\label{0.16.4.6.4-fr}
$\mathfrak{p}$가 $A$의 소 아이디얼일 때
$\operatorname{prof}_{A_{\mathfrak{p}}}(M_{\mathfrak{p}})
<\operatorname{prof}_A(M)$일 수도 있고
$\operatorname{prof}_{A_{\mathfrak{p}}}(M_{\mathfrak{p}})
>\operatorname{prof}_A(M)$일 수도 있음에 주의하자. 첫 번째 경우의
예로는
\oldpage[0\textsubscript{IV-1}]{35}
차원이 $\geq1$인 국소 정역 $A$를 취하면 충분하다. 이때
$\operatorname{prof}(A)\geq1$이지만 $A$의 분수체 $K$에 대해서는
$\operatorname{prof}(K)=0$이다. 두 번째 경우의 예로, 차원이
$\geq2$이고 극대 아이디얼이 $\mathfrak{m}_0$이며 잉여체가
$k=A_0/\mathfrak{m}_0$인 뇌터 국소 정역 $A_0$를 생각하자.
$A=A_0\oplus\mathfrak{J}$를 $A_0$-가군 $k$에 의한 $A_0$의 자명한
확대라 하자 (18.2.3). 여기서 아이디얼 $\mathfrak{J}$는 $A_0$-가군으로
$k$와 동형이어서 $A$의 곱셈은
$(x,z)(x',z')=(xx',xz'+x'z)$로 주어진다. $\mathfrak{J}$가 $A$의
영근기이고 $A/\mathfrak{J}$가 $A_0$와 동형임은 명백하다. 따라서
$A$의 모든 소 아이디얼 $\mathfrak{p}$는
$\mathfrak{p}_0\oplus\mathfrak{J}$ 꼴이며, 여기서
$\mathfrak{p}_0$는 $A_0$의 소 아이디얼이다. 특히
$\mathfrak{m}=\mathfrak{m}_0\oplus\mathfrak{J}$는 $A$의 유일한 극대
아이디얼이다. $\mathfrak{m}_0$의 모든 원소가 $\mathfrak{J}$를
소멸시키므로
$\operatorname{prof}(A)=\operatorname{prof}(A_{\mathfrak{m}})=0$이다.
반대로 $\mathfrak{p}_0\neq\mathfrak{m}_0$이면 $\mathfrak{J}$의
원소들은 $\mathfrak{m}_0-\mathfrak{p}_0$의 원소들에 의해 소멸된다.
따라서 $A_{\mathfrak{p}}$는 표준적으로
$(A_0)_{\mathfrak{p}_0}$와 동형이고, $A_0$가 정역이므로
$\mathfrak{p}_0\neq0$일 때
$\operatorname{prof}(A_{\mathfrak{p}})
=\operatorname{prof}((A_0)_{\mathfrak{p}_0})\geq1$이다.
\end{env}

\begin{corollary}[16.4.7]
\label{0.16.4.7-fr}
축소환인 뇌터 국소환 $A$의 깊이가 0일 필요충분조건은 $A$가 체인 것이다.
\end{corollary}

\begin{proof}
실제로 $A$의 극소 소 아이디얼들의 교집합은 0이므로 $A$에는 극소가
아닌 연관 소 아이디얼이 없다. (16.4.6, (i))에 의해
$\operatorname{prof}(A)=0$이면 $A$의 극대 아이디얼
$\mathfrak{m}$도 극소 소 아이디얼이다. 따라서 이것은 $A$의 유일한
소 아이디얼이고, $A$가 축소환이므로 $\mathfrak{m}=0$이다.
\end{proof}

\begin{proposition}[16.4.8]
\label{0.16.4.8-fr}
$A$, $B$를 두 뇌터 국소환, $\rho:A\to B$를 국소 준동형이라 하고,
$M_{[\rho]}$가 유한형 $A$-가군이 되는 $B$-가군 $M$을 취하자. 그러면
\begin{equation}
\label{0.16.4.8.1-fr}
\tag{16.4.8.1}
  \operatorname{prof}_A(M_{[\rho]})=\operatorname{prof}_B(M).
\end{equation}
\end{proposition}

\begin{proof}
$M\neq0$인 경우로 한정해도 된다.
$\operatorname{prof}_A(M_{[\rho]})=n$이라 하자.
$(x_i)_{1\leq i\leq n}$이 $A$의 극대 아이디얼의 원소들로 이루어진
극대 $(M_{[\rho]})$-정칙열이면, $\rho(x_i)$들은 $B$의 극대 아이디얼의
원소들로 이루어진 $M$-정칙열이다. 따라서
$N=M/(\rho(x_1)M+\cdots+\rho(x_n)M)$이라 두면 (16.4.6, (ii))에 의해
$\operatorname{prof}_B(N)=\operatorname{prof}_B(M)-n$이다. 마찬가지로
$N_{[\rho]}=M_{[\rho]}/(x_1M_{[\rho]}+\cdots+x_nM_{[\rho]})$이고
$\operatorname{prof}_A(N_{[\rho]})=0$이다. 그러므로 $n=0$인 경우에
명제를 증명하면 충분하다. $k$를 $A$의 잉여체라 하자. $k$는 $A$의
몫이므로 $A$-가군
$P=\operatorname{Hom}_A(k,M_{[\rho]})$는
$\operatorname{Hom}_A(A,M_{[\rho]})=M_{[\rho]}$의 부분가군이고,
가정에 의해 $P\neq0$이다 (16.4.6, (i)). 또한 $M_{[\rho]}$가 유한형
$A$-가군이므로 $P$도 유한형 $A$-가군이다. $P$는 $A$의 극대
아이디얼에 의해 소멸되므로 유한차원 $k$-벡터공간이고, 따라서 유한
길이 $A$-가군이다. 한편 $P$는 $B$-가군 $M$의 부분가군이기도 하다.
$A$-가군으로 유한 길이이므로 더욱이 $B$-가군으로도 유한 길이이다.
$P\neq0$이므로 $P$는 단순 $B$-부분가군, 즉 $B$의 잉여체와 동형인
부분가군을 포함한다. 그러므로 (16.4.6, (i))에 의해 실제로
$\operatorname{prof}_B(M)=0$이다.
\end{proof}

\begin{definition}[16.4.9]
\label{0.16.4.9-fr}
$A$를 뇌터 국소환, $M$을 유한형 $A$-가군이라 하자. $M\neq0$이면
유한한 비음 정수
$\dim_A(M)-\operatorname{prof}_A(M)$을 (16.4.5.1) $M$의
\emph{여깊이}라 하고,
$\operatorname{coprof}_A(M)$ 또는 $\operatorname{coprof}(M)$으로
쓴다. $M=0$이면 $\operatorname{coprof}(M)=0$으로 둔다.

모든 $k>0$에 대하여
$\operatorname{coprof}(M^k)=\operatorname{coprof}(M)$이다.
\end{definition}

\begin{proposition}[16.4.10]
\label{0.16.4.10-fr}
$A$를 뇌터 국소환, $M$을 유한형 $A$-가군이라 하자.
\oldpage[0\textsubscript{IV-1}]{36}
\begin{enumerate}
  \item[(i)] $A$의 극대 아이디얼에 속하는 모든 $M$-정칙 원소 $x$에
  대하여 $\operatorname{coprof}(M/xM)=\operatorname{coprof}(M)$이다.
  \item[(ii)] 다음 등식이 성립한다.
  $\operatorname{coprof}_A(M)=
  \operatorname{coprof}_{\widehat A}(\widehat M)$.
\end{enumerate}
\end{proposition}

\begin{proof}
(i)은 (16.3.4)와 (16.4.6, (ii))에서, (ii)는 (16.2.4)와
(16.4.6, (iv))에서 따른다.
\end{proof}

나중에 (IV, 6.11.5)에서 모든 $A$의 소 아이디얼 $\mathfrak{p}$에 대하여
$\operatorname{coprof}_{A_{\mathfrak{p}}}(M_{\mathfrak{p}})
\leq\operatorname{coprof}_A(M)$임을 보일 것이다.

\begin{proposition}[16.4.11]
\label{0.16.4.11-fr}
(16.4.8)의 가정 아래
\begin{equation}
\label{0.16.4.11.1-fr}
\tag{16.4.11.1}
  \operatorname{coprof}_A(M_{[\rho]})=\operatorname{coprof}_B(M)
\end{equation}
이다.
\end{proposition}

\begin{proof}
이는 (16.1.9)와 (16.4.8)에서 따른다.
\end{proof}

\subsection{Cohen--Macaulay 가군}
\label{subsection:0.16.5-fr}

\begin{definition}[16.5.1]
\label{0.16.5.1-fr}
$A$를 뇌터 국소환이라 하자. 유한형 $A$-가군 $M$이
$\operatorname{coprof}(M)=0$을 만족하면, 다시 말해 $M=0$이거나
$M\neq0$이고 $\dim(M)=\operatorname{prof}(M)$이면 $M$을
\emph{Cohen--Macaulay 가군}이라 한다. $A$-가군 $A$가
Cohen--Macaulay 가군이면 $A$를 \emph{Cohen--Macaulay 환}이라 한다.
\end{definition}

유한 길이 $A$-가군은 차원이 0이므로 (16.1.10) Cohen--Macaulay
가군이다. 차원이 1인 국소 축소환 $A$는 Cohen--Macaulay 환이다.
실제로 이때 극대 아이디얼 $\mathfrak{m}$은 $A$에 연관될 수 없으므로
(Bourbaki, \emph{Alg.~comm.}, chap.~IV, \S~2, no~5, prop.~10),
(16.4.6, (i))에 의해 $\operatorname{prof}(A)=1$이다.
$A$가 \emph{차원이 $\geq2$인 정규 국소 정역}이면
$\operatorname{prof}(A)\geq2$이다. 실제로
$x\neq0$이 $\mathfrak{m}$의 원소이면 $A/xA$의 연관 소 아이디얼들은
높이가 1임이 알려져 있다
(Bourbaki, \emph{Alg.~comm.}, chap.~VII, \S~1, no~4, prop.~8).
따라서 이들은 $\mathfrak{m}$과 다르고, 그 합집합도
$\mathfrak{m}$과 다르다
(Bourbaki, \emph{Alg.~comm.}, chap.~II, \S~1, no~1, prop.~2).
그러므로 $\mathfrak{m}$의 원소들로 이루어진 $A$-정칙열 $(x,y)$가
존재하며, 이로부터 주장이 따른다. 따라서 차원이 2인 정규 국소환은
Cohen--Macaulay 환이다.

\begin{proposition}[16.5.2]
\label{0.16.5.2-fr}
$A$를 뇌터 국소환, $M$을 유한형 $A$-가군이라 하자. $M$이
Cohen--Macaulay 가군일 필요충분조건은 $\widehat M$이
Cohen--Macaulay $\widehat A$-가군인 것이다.
\end{proposition}

\begin{proof}
이는 (16.4.10, (ii))에서 곧바로 따른다.
\end{proof}

\begin{proposition}[16.5.3]
\label{0.16.5.3-fr}
(16.4.8)의 가정 아래 $M$이 Cohen--Macaulay $B$-가군일
필요충분조건은 $M_{[\rho]}$가 Cohen--Macaulay $A$-가군인 것이다.
\end{proposition}

\begin{proof}
이는 (16.4.11)에서 따른다.
\end{proof}

\begin{proposition}[16.5.4]
\label{0.16.5.4-fr}
$A$를 뇌터 국소환, $M\neq0$을 유한형 $A$-가군이라 하자. $M$이
Cohen--Macaulay $A$-가군이면, 모든 소 아이디얼
$\mathfrak{p}\in\operatorname{Ass}(M)$에 대하여
$\dim(A/\mathfrak{p})=\dim(M)$이다. 특히 $M$의 연관 소 아이디얼
가운데 극소가 아닌 것은 없다.
\end{proposition}

\begin{proof}
이는 부등식 (16.4.6.2)에서 곧바로 따른다.
\end{proof}

$M$ 또는 $\operatorname{Supp}(M)$이 \emph{등차원}이라고 말할 수도
있다 (16.1.7).

\oldpage[0\textsubscript{IV-1}]{37}

\begin{proposition}[16.5.5]
\label{0.16.5.5-fr}
$A$를 뇌터 국소환, $\mathfrak{m}$을 그 극대 아이디얼,
$M\neq0$을 유한형 $A$-가군, $x$를 $\mathfrak{m}$의 원소라 하자.
$M$이 Cohen--Macaulay 가군이라고 가정하자. 그러면 다음 조건들은
동치이다.
\begin{enumerate}
  \item[a)] $x$는 $M$-정칙이다.
  \item[b)] $\dim(M/xM)=\dim(M)-1$이다.
  \item[c)] $x$는 $\operatorname{Supp}(M)$의 어느 극소 원소에도
  속하지 않는다.
\end{enumerate}
또한 이때 $M/xM$도 Cohen--Macaulay 가군이다.
\end{proposition}

\begin{proof}
(16.5.4)에서 $\operatorname{Ass}(M)$의 모든 원소
$\mathfrak{p}$가 $\operatorname{Supp}(M)$에서 극소이고
$\dim(A/\mathfrak{p})=\dim(M)$임을 보았다. 따라서 a)와 c)의 동치는
Bourbaki, \emph{Alg.~comm.}, chap.~IV, \S~1, no~1, prop.~2의
cor.~2에서 따르고, b)와 c)의 동치는 (16.3.4)에서 따른다. 마지막으로
$M/xM$이 Cohen--Macaulay 가군이라는 사실은 (16.4.10, (i))에서
따른다.
\end{proof}

\begin{corollary}[16.5.6]
\label{0.16.5.6-fr}
(16.5.5)의 가정 아래 $(x_i)_{1\leq i\leq r}$을
$\mathfrak{m}$의 원소들로 이루어진 열이라 하자. 다음 조건들은
동치이다.
\begin{enumerate}
  \item[a)] 열 $(x_i)$는 $M$-정칙이다.
  \item[b)] 다음 등식이 성립한다.
  $\dim(M/(x_1M+\cdots+x_rM))=\dim(M)-r$.
  \item[c)] $x_i$들은 $M$의 매개변수계의 일부를 이룬다.
\end{enumerate}
또한 이 조건들이 성립할 때
$N=M/(x_1M+\cdots+x_rM)$은 Cohen--Macaulay 가군이다. 따라서 모든
아이디얼 $\mathfrak{p}\in\operatorname{Ass}(N)$에 대하여
$\dim(A/\mathfrak{p})=\dim(M)-r$이다.
\end{corollary}

\begin{proof}
b)와 c)가 동치임은 이미 알려져 있고 (16.3.6), $M$에 관한 아무
가정 없이 a)가 c)를 함의한다 (16.4.1). c)가 a)를 함의함을 보이고
명제의 마지막 주장도 증명하기 위해 $r$에 관해 귀납하자. $r=0$이면
주장은 자명하다. $x_1$이 매개변수계의 일부를 이루므로
$\dim(M/x_1M)=\dim(M)-1$이고 (16.3.6), 따라서 (16.5.5)에 의해
$x_1$은 $M$-정칙이며 $M/x_1M$은 Cohen--Macaulay 가군이다.
열 $(x_i)_{2\leq i\leq r}$은 $P=M/x_1M$의 매개변수계의 일부를
이루므로, 귀납가정에 의해 $(x_i)_{2\leq i\leq r}$은 $P$-정칙열이고
$P/(x_2P+\cdots+x_rP)=N$은 Cohen--Macaulay 가군이다. 또한
$(x_i)_{1\leq i\leq r}$은 $M$-정칙열이다.
\end{proof}

\begin{remark}[16.5.7]
\label{0.16.5.7-fr}
(16.5.6)에 의해 $M$이 Cohen--Macaulay 가군이면 \emph{$M$의
매개변수계}와 \emph{$\mathfrak{m}$의 원소들로 이루어진 극대
$M$-정칙열}은 같은 것이다. 역으로 (16.4.1)에 의해, 비영 유한형
$A$-가군 $M$에 $M$의 매개변수계이기도 한 $M$-정칙열이 존재하면
$M$은 Cohen--Macaulay 가군이다. (16.5.6)의 마지막 성질도
Cohen--Macaulay 가군을 \emph{특징짓는다}.
\end{remark}

\begin{proposition}[16.5.8]
\label{0.16.5.8-fr}
$A$를 뇌터 국소환, $M$을 유한형 $A$-가군이라 하자. $M$의
매개변수계의 일부를 이루는 $A$의 원소들로 된 모든 열
$(x_i)_{1\leq i\leq r}$과
$N=M/(x_1M+\cdots+x_rM)$에 연관된 모든 소 아이디얼
$\mathfrak{p}$에 대하여 $\dim(A/\mathfrak{p})=\dim(N)$이라고
가정하자. 그러면 $M$은 Cohen--Macaulay 가군이다.
\end{proposition}

\begin{proof}
$(y_i)_{1\leq i\leq n}$을 $M$의 매개변수계라 하자. 이 열이
$M$-정칙임을 보이면 충분하다. $n=\dim(M)$에 관해 귀납하자.
$n=0$이면 명제는 자명하다. 가정을 $r=0$에 적용하면 $M$에 연관된
소 아이디얼 $\mathfrak{p}_j$들은 모두
$\dim(A/\mathfrak{p}_j)=\dim(M)$을 만족한다. $y_1$은 $M$의
매개변수계의 일부를
\oldpage[0\textsubscript{IV-1}]{38}
이루므로, (16.3.7)과 (16.3.4)에 의해 $y_1$은 어느
$\mathfrak{p}_j$에도 속하지 않는다. 따라서 $y_1$은 $M$-정칙이다.
$P=M/y_1M$이라 두면 $(y_i)_{2\leq i\leq n}$은 $P$의 매개변수계이고,
$P$가 명제의 가정을 만족함은 자명하다. 그러므로 귀납가정을 적용하면
$(y_i)_{2\leq i\leq n}$은 $P$-정칙열이며, 이로써 명제가 증명된다.
\end{proof}

\begin{proposition}[16.5.9]
\label{0.16.5.9-fr}
$A$를 뇌터 국소환, $M$을 유한형 $A$-가군이라 하고, $M$이
Cohen--Macaulay 가군이라고 가정하자.
$\mathfrak{p}\in\operatorname{Supp}(M)$을 취하고
$r=\dim(M)-\dim(M/\mathfrak{p}M)$이라 두자. 그러면
$\mathfrak{p}$의 원소들로 이루어진 $M$-정칙열
$(x_i)_{1\leq i\leq r}$이 존재한다. 이러한 성질을 갖는 모든 열에
대하여
\[
  \dim(M/(x_1M+\cdots+x_rM))
  =\dim(M/\mathfrak{p}M)=\dim(A/\mathfrak{p})
\]
이고, $\mathfrak{p}$는
$\operatorname{Ass}(M/(x_1M+\cdots+x_rM))$의 극소 원소이다.
\end{proposition}

\begin{proof}
먼저 $x_i$들의 존재를 증명하자. $r=0$이면 증명할 것이 없다.
$r>0$일 때 $r$에 관해 귀납하자. $M$에 연관된 모든 소 아이디얼
$\mathfrak{q}_i$에 대하여
$\dim(A/\mathfrak{q}_i)=\dim(M)$이다 (16.5.4). 따라서
$\dim(M/\mathfrak{p}M)<\dim(M)$이라는 가정에 의해
$\mathfrak{p}$는 어느 $\mathfrak{q}_i$에도 포함될 수 없다. 실제로
$\mathfrak{q}_i$들은 $\operatorname{Ann}(M)$을 포함하는 소
아이디얼들 가운데 극소인 것들이다 (16.1.2.2와 16.1.7). 그러므로
$\mathfrak{p}$는 $\mathfrak{q}_i$들의 합집합에도 포함되지 않고
(Bourbaki, \emph{Alg.~comm.}, chap.~II, \S~1, no~1, prop.~2), 따라서
$M$-정칙인 $x_1\in\mathfrak{p}$가 존재한다
(Bourbaki, \emph{Alg.~comm.}, chap.~IV, \S~1, no~1, prop.~2의
cor.~2). (16.5.5)에 의해 $N=M/x_1M$은 차원이
$\dim(M)-1$인 Cohen--Macaulay 가군이다. 또한
$x_1\in\mathfrak{p}$이므로
$N/\mathfrak{p}N=M/\mathfrak{p}M$이고, 따라서
$\dim(N)-\dim(N/\mathfrak{p}N)=r-1$이다. 이제 $N$에 귀납가정을
적용할 수 있다. $(x_i)_{2\leq i\leq r}$이 $\mathfrak{p}$의
원소들로 이루어진 $N$-정칙열이면
$(x_i)_{1\leq i\leq r}$은 $\mathfrak{p}$의 원소들로 이루어진
$M$-정칙열임이 명백하다.

$P=M/(x_1M+\cdots+x_rM)$이라 두자. $x_i$들이
$\mathfrak{p}$에 속하므로
$P/\mathfrak{p}P=M/\mathfrak{p}M$이고, (16.5.6)에 의해
$\dim(P)=\dim(P/\mathfrak{p}P)$이다. 한편 $P$는 Cohen--Macaulay
가군이므로 모든 $\mathfrak{p}'\in\operatorname{Ass}(P)$에 대하여
$\dim(P)=\dim(A/\mathfrak{p}')$이다 (16.5.4).
$\mathfrak{p}\in\operatorname{Supp}(P)$이므로
($0_{\mathrm I}$, 1.7.5), $\mathfrak{p}$는 어느
$\mathfrak{p}'\in\operatorname{Ass}(P)$를 포함한다
(Bourbaki, \emph{Alg.~comm.}, chap.~IV, \S~1, no~4, th.~2). 그런데
\[
  \dim(P)=\dim(P/\mathfrak{p}P)
  \leq\dim(A/\mathfrak{p})
  \leq\dim(A/\mathfrak{p}')=\dim(P)
\]
이므로 반드시 $\mathfrak{p}=\mathfrak{p}'$이다 (16.1.2.2). 이로써
증명이 끝난다.
\end{proof}

\begin{corollary}[16.5.10]
\label{0.16.5.10-fr}
(16.5.9)와 같은 가정 아래 다음이 성립한다.
\begin{enumerate}
  \item[(i)] $M_{\mathfrak{p}}$는 Cohen--Macaulay
  $A_{\mathfrak{p}}$-가군이다.
  \item[(ii)] 다음 등식이 성립한다.
  \begin{equation}
  \label{0.16.5.10.1-fr}
  \tag{16.5.10.1}
    \dim(M)=\dim_A(M/\mathfrak{p}M)
    +\dim_{A_{\mathfrak{p}}}(M_{\mathfrak{p}}).
  \end{equation}
\end{enumerate}
\end{corollary}

\begin{proof}
(16.5.9)의 표기 아래 $x_i'$를 $A_{\mathfrak{p}}$에서 $x_i$의 표준
상이라 하자 $(1\leq i\leq r)$. $x_i'$들은
$A_{\mathfrak{p}}$의 극대 아이디얼에 속하고, 평탄성에 의해
$M_{\mathfrak{p}}$-정칙열을 이룬다 (15.1.14). 따라서
\[
  \operatorname{prof}_{A_{\mathfrak{p}}}(M_{\mathfrak{p}})
  \geq r=\dim(M)-\dim(M/\mathfrak{p}M)
  \geq\dim_{A_{\mathfrak{p}}}(M_{\mathfrak{p}})
\]
이다 (16.1.8.1). 그러나 (16.4.5.1)을 함께 고려하면 이 부등식의 세
항은 반드시 모두 같으므로 따름정리가 증명된다.
\end{proof}

\oldpage[0\textsubscript{IV-1}]{39}

\begin{corollary}[16.5.11]
\label{0.16.5.11-fr}
$A$를 뇌터 국소환이라 하자. 지지가 $\operatorname{Spec}(A)$인 유한형
$A$-가군 $M$ 가운데 Cohen--Macaulay 가군인 것이 존재한다고
가정하자. 그러면 $A$의 모든 소 아이디얼 $\mathfrak{p}$에 대하여
\begin{equation}
\label{0.16.5.11.1-fr}
\tag{16.5.11.1}
  \dim(A)=\dim(A/\mathfrak{p})+\dim(A_{\mathfrak{p}}).
\end{equation}
\end{corollary}

\begin{proof}
실제로 $\operatorname{Supp}(M/\mathfrak{p}M)
=\operatorname{Spec}(A/\mathfrak{p})$이고 ($0_{\mathrm I}$, 1.7.5),
$\operatorname{Supp}(M_{\mathfrak{p}})
=\operatorname{Spec}(A_{\mathfrak{p}})$이다. 따라서 관계 (16.5.11.1)은
(16.5.10.1)의 특수한 경우이다.
\end{proof}

\begin{corollary}[16.5.12]
\label{0.16.5.12-fr}
(16.5.11)의 조건을 만족하는 뇌터 국소환 $A$의 모든 몫환은
\emph{카테너리}이다. 특히 Cohen--Macaulay 국소환의 모든 몫환은
카테너리이다.
\end{corollary}

\begin{proof}
$A$ 자체가 카테너리임을 보이면 충분함은 명백하다. 즉 $A$의 두 소
아이디얼 $\mathfrak{p}\subset\mathfrak{q}$에 대하여 (16.1.4.2)
\[
  \dim(A_{\mathfrak{q}})
  =\dim(A_{\mathfrak{p}})
  +\dim(A_{\mathfrak{q}}/\mathfrak{p}A_{\mathfrak{q}})
\]
임을 보이면 된다. 그런데 (16.5.10, (i))에 의해
$A_{\mathfrak{q}}$는 $A$와 같은 가정을 만족한다. 따라서 앞의 관계는
$A_{\mathfrak{q}}$와 그 소 아이디얼
$\mathfrak{p}A_{\mathfrak{q}}$에 (16.5.11.1)을 적용한 것일 뿐이다.
\end{proof}

\begin{env}[16.5.13]
\label{0.16.5.13-fr}
$A$를 뇌터 환, $M$을 유한형 $A$-가군이라 하자. $A$의 모든 소
아이디얼 $\mathfrak{p}$에 대하여 $M_{\mathfrak{p}}$가
Cohen--Macaulay $A_{\mathfrak{p}}$-가군이면 $M$을
\emph{Cohen--Macaulay $A$-가군}이라 한다. (16.5.10)에 의해 $A$가
\emph{국소환}이면 이 정의는 (16.5.1)과 일치한다. 모든
$A_{\mathfrak{p}}$가 Cohen--Macaulay 환이면 $A$를
\emph{Cohen--Macaulay 환}이라 한다. $M$이 Cohen--Macaulay
$A$-가군이면 임의의 $A$의 곱셈적 부분집합 $S$에 대하여
$S^{-1}M$은 Cohen--Macaulay $S^{-1}A$-가군이다. 마찬가지로 $A$가
Cohen--Macaulay 환이면 $S^{-1}A$도 Cohen--Macaulay 환임은 명백하다.
\end{env}

\section{정칙환}
\label{section:0.17-fr}

\subsection{정칙환의 정의}
\label{subsection:0.17.1-fr}

\begin{proposition}[17.1.1]
\label{0.17.1.1-fr}
$A$를 차원 $n$인 뇌터 국소환, $\mathfrak{m}$을 그 극대 아이디얼,
$k=A/\mathfrak{m}$을 그 잉여체라 하자. 다음 조건들은 서로 동치이다.
\begin{enumerate}
  \item[\rm a)] $k$-등급가군의 정준 전사 준동형 (15.1.1.1)
  \[
    \varphi:\mathbf{S}^*_k(\mathfrak{m}/\mathfrak{m}^2)
    \longrightarrow\operatorname{gr}^*_{\mathfrak{m}}(A)
  \]
  는 전단사이다. 여기서 오른쪽 항은 $\mathfrak{m}$-준아딕 여과를 준
  $A$에 연관된 $k$-등급가군이다.
  \item[\rm b)]
  $\operatorname{rg}_k(\mathfrak{m}/\mathfrak{m}^2)=n$이다.
  \item[\rm c)] 아이디얼 $\mathfrak{m}$에는 $n$개 원소로 이루어진
  생성계가 존재한다.
  \item[\rm d)] 아이디얼 $\mathfrak{m}$에는 $A$-정칙열인 생성계가
  존재한다.
\end{enumerate}
\end{proposition}

\begin{proof}
나카야마 보조정리에 의하면
$\operatorname{rg}_k(\mathfrak{m}/\mathfrak{m}^2)$는
$\mathfrak{m}$의 생성계를 이루는 원소들의 최소 개수이므로 b)와 c)는
동치이다. 한편 $(x_i)_{1\leq i\leq r}$가 $\mathfrak{m}$의 생성계이고
그 원소들의 mod. $\mathfrak{m}^2$에서의 류 $\overline{x}_i$들이
$k$-벡터공간 $\mathfrak{m}/\mathfrak{m}^2$의 기저를 이루면,
$\mathbf{S}^*(\mathfrak{m}/\mathfrak{m}^2)$는 다항식환
\oldpage[0\textsubscript{IV-1}]{40}
$k[\mathrm{T}_1,\ldots,\mathrm{T}_r]$과 동형이다. 모든 유한형
$A$-가군이 $\mathfrak{m}$-준아딕 여과에 대하여 분리라는 사실
($0_{\mathrm I}$, 7.3.5)을 고려하면, (15.1.9)에 의해 a)와 d)는
동치이다. 또한 $\mathfrak{m}$의 원소들로 된 모든 $A$-정칙열은
원소를 많아야 $n$개 가지므로 (16.4.1), d)는 c)를 함의한다.
이제 c)가 a)를 함의함을 보이면 된다.

간단히
$\mathrm{S}=\mathbf{S}^*(\mathfrak{m}/\mathfrak{m}^2)
=k[\mathrm{T}_1,\ldots,\mathrm{T}_n]$,
$\mathrm{G}=\operatorname{gr}^*_{\mathfrak{m}}(A)$라 두고, 완전열
$0\longrightarrow\mathfrak{J}\longrightarrow\mathrm{S}
\xrightarrow{\varphi}\mathrm{G}\longrightarrow0$을 생각하자. 여기서
$\varphi$의 핵 $\mathfrak{J}$은 $\mathrm{S}$의 등급 아이디얼이다.
따라서 모든 정수 $s>0$에 대하여
\begin{equation}
\label{0.17.1.1.1-fr}
\tag{17.1.1.1}
  \binom{s+n-1}{n-1}
  =\operatorname{long}(\mathrm{S}_s)
  =\operatorname{long}(\mathrm{G}_s)
  +\operatorname{long}(\mathfrak{J}_s).
\end{equation}

$\mathfrak{J}\neq0$이라고 가정하자. 그러면 차수 $h>0$인 동차 원소
$u\in\mathfrak{J}$가 존재한다. $s\geq h$이면 $\mathfrak{J}_s$는
$u\mathrm{S}_{s-h}$를 포함하고, $u\neq0$이며 $\mathrm{S}$가
정역이므로 $u\mathrm{S}_{s-h}$는 $k$-벡터공간으로서
$\mathrm{S}_{s-h}$와 동형이다. 따라서
$\operatorname{long}(\mathfrak{J}_s)\geq
\operatorname{long}(\mathrm{S}_{s-h})=\binom{s-h+n-1}{n-1}$이고,
$s\geq h$에 대하여
\begin{equation}
\label{0.17.1.1.2-fr}
\tag{17.1.1.2}
  \operatorname{long}(\mathrm{G}_s)
  \leq\binom{s+n-1}{n-1}-\binom{s-h+n-1}{n-1}.
\end{equation}
그런데 (17.1.1.2)의 오른쪽 항은 $s$에 관한 차수 $\leq n-2$인
다항식이다. 반면 $\mathrm{G}_s=\mathfrak{m}^s/\mathfrak{m}^{s+1}$이고,
충분히 큰 $s$에 대하여 $\operatorname{long}(\mathrm{G}_s)$는 $s$에
관한 차수가 정확히 $n-1$이고 최고차항 계수가 $>0$인 다항식이다
(16.2.1). 이는 부등식 (17.1.1.2)와 모순이다. C.Q.F.D.
\end{proof}

\begin{definition}[17.1.2]
\label{0.17.1.2-fr}
(17.1.1)의 동치인 조건들을 만족하는 뇌터 국소환을 \emph{정칙}이라
한다.
\end{definition}

\begin{corollary}[17.1.3]
\label{0.17.1.3-fr}
정칙 국소환은 정역이고 정수적으로 닫혀 있으며 Cohen--Macaulay
환이다.
\end{corollary}

\begin{proof}
마지막 주장은 (17.1.1, d))에서 따른다. 한편 $A$가 정칙이면
$\operatorname{gr}^*_{\mathfrak{m}}(A)$는
$k[\mathrm{T}_1,\ldots,\mathrm{T}_n]$과 동형이므로 정역이고 완전
정수적으로 닫혀 있다. 따라서 $A$도 이 두 성질을 갖는다 (Bourbaki,
\emph{Alg.~comm.}, chap.~III, \S~2, no~3, cor. de la prop.~1 및 chap.~V,
\S~1, no~5, prop.~15).
\end{proof}

\begin{examples}[17.1.4]
\label{0.17.1.4-fr}
\begin{enumerate}
  \item[\rm (i)] \emph{차원} $0$인 정칙 국소환은 (17.1.3)에 의해
  정역이므로 반드시 \emph{체}이고, 그 역도 성립한다.
  \item[\rm (ii)] \emph{차원} $1$인 뇌터 국소환이 정칙이기 위한
  필요충분조건은 \emph{이산 값매김환}인 것이다. 실제로 차원 $1$인
  뇌터 국소환이 정칙이라는 것은 그 극대 아이디얼이 \emph{주
  아이디얼}이라는 뜻이고, 결론은 Bourbaki, \emph{Alg.~comm.},
  chap.~VI, \S~3, no~6, prop.~9에서 따른다.
  \item[\rm (iii)] $k$를 체라 하자. 형식적 멱급수환
  $A=k[[\mathrm{T}_1,\ldots,\mathrm{T}_n]]$은 \emph{차원 $n$인
  정칙환}이다. 실제로 $\mathrm{T}_i$들 $(1\leq i\leq n)$이 $A$의
  극대 아이디얼 $\mathfrak{m}$을 생성함은 명백하고, 이들은
  $A$-정칙열을 이룬다. 왜냐하면
  $A/(\mathrm{T}_1A+\cdots+\mathrm{T}_iA)$는
  $k[[\mathrm{T}_{i+1},\ldots,\mathrm{T}_n]]$과 동형이기 때문이다.
\end{enumerate}
\end{examples}

\begin{proposition}[17.1.5]
\label{0.17.1.5-fr}
뇌터 국소환 $A$가 정칙이기 위한 필요충분조건은 그 완비화
$\widehat{A}$가 정칙인 것이다.
\end{proposition}

\begin{proof}
\oldpage[0\textsubscript{IV-1}]{41}
실제로 $\widehat{A}$의 극대 아이디얼은
$\mathfrak{m}\widehat{A}$이고,
$\mathfrak{m}^h/\mathfrak{m}^{h+1}$와
$(\mathfrak{m}\widehat{A})^h/(\mathfrak{m}\widehat{A})^{h+1}$가 서로
동형인 $k$-벡터공간임은 알려져 있다. 따라서 (17.1.1)의 조건 a)는
$A$와 $\widehat{A}$에 대하여 같다.
\end{proof}

\begin{definition}[17.1.6]
\label{0.17.1.6-fr}
$A$를 극대 아이디얼이 $\mathfrak{m}$인 뇌터 국소환이라 하자.
$A$의 매개변수계 (16.3.5) 가운데 $\mathfrak{m}$을 생성하는 것을
\emph{$A$의 정칙 매개변수계}라 한다.
\end{definition}

그러한 매개변수계가 존재하는 것은 $A$가 정칙인 것과 동치이다
(17.1.1). $(x_i)_{1\leq i\leq n}$이 정칙 매개변수계이면 이는
$\mathfrak{m}$의 최소 생성계이고, 그 원소들의 mod.
$\mathfrak{m}^2$에서의 류들은 $\mathfrak{m}/\mathfrak{m}^2$의
$k$-기저를 이룬다. (17.1.1, a))와 (15.1.9)에 의해 $(x_i)$는 극대
$A$-정칙열이다 (이로부터 이 용어를 쓴다).

그러나 정칙환에서 $A$의 매개변수계이면서 $A$-정칙열인 것이 반드시
정칙 매개변수계인 것은 아님에 주의해야 한다. 정칙 매개변수계의 각
원소를 $k$제곱하여 얻는 예가 이를 보여 준다 (15.1.20).

\begin{proposition}[17.1.7]
\label{0.17.1.7-fr}
$A$를 뇌터 국소환, $\mathfrak{m}$을 그 극대 아이디얼,
$k=A/\mathfrak{m}$을 그 잉여체, $(x_i)_{1\leq i\leq r}$를
$\mathfrak{m}$의 원소들로 된 열이라 하고,
$\mathfrak{J}=x_1A+\cdots+x_rA$라 하자. 다음 조건들은 서로 동치이다.
\begin{enumerate}
  \item[\rm a)] $A$는 정칙이고 $x_i$들은 $A$의 한 정칙
  매개변수계의 일부를 이룬다.
  \item[\rm a')] $A$는 정칙이고 $x_i$들의 mod.
  $\mathfrak{m}^2$에서의 류 $\overline{x}_i$들은 $k$ 위에서 선형
  독립이다.
  \item[\rm b)] $x_i$들은 $A$의 한 매개변수계의 일부를 이루고
  $A/\mathfrak{J}$은 정칙환이다.
\end{enumerate}
또한 이 조건들이 성립하면 $\mathfrak{J}$은 소 아이디얼이다.
\end{proposition}

\begin{proof}
마지막 주장은 $A/\mathfrak{J}$이 정칙이어서 정역이라는 사실
(17.1.3)에서 바로 따른다. $n=\dim(A)$라 두자. a)와 a')의 동치는
명백하다. 실제로 정칙 매개변수계의 원소들의 mod.
$\mathfrak{m}^2$에서의 류들은 $\mathfrak{m}/\mathfrak{m}^2$의
$k$-기저를 이룬다. 역으로 $\overline{x}_i$ $(1\leq i\leq r)$가
$k$ 위에서 선형 독립이면, $\overline{x}_i$ $(1\leq i\leq n)$이
$\mathfrak{m}/\mathfrak{m}^2$의 기저를 이루도록 하는 $n-r$개의
원소 $x_i\in\mathfrak{m}$ $(r+1\leq i\leq n)$을 찾을 수 있다.
따라서 $(x_i)_{1\leq i\leq n}$은 정칙 매개변수계이다.

a)가 b)를 함의함을 보이자. $x_i$들이 $A$의 한 매개변수계의 일부를
이루므로 $\dim(A/\mathfrak{J})=n-r$이다 (16.3.6).
$\mathfrak{n}=\mathfrak{m}/\mathfrak{J}$을 $A/\mathfrak{J}$의 극대
아이디얼이라 하자. $\mathfrak{n}^2=
(\mathfrak{m}^2+\mathfrak{J})/\mathfrak{J}$이므로 완전열
\begin{equation}
\label{0.17.1.7.1-fr}
\tag{17.1.7.1}
  0\longrightarrow(\mathfrak{m}^2+\mathfrak{J})/\mathfrak{m}^2
  \longrightarrow\mathfrak{m}/\mathfrak{m}^2
  \longrightarrow\mathfrak{n}/\mathfrak{n}^2\longrightarrow0
\end{equation}
을 얻는다. 가정 a)에 의해
$\operatorname{rg}_k((\mathfrak{m}^2+\mathfrak{J})/
\mathfrak{m}^2)=r$이다. 따라서
$\operatorname{rg}_k(\mathfrak{n}/\mathfrak{n}^2)=n-r
=\dim(A/\mathfrak{J})$이고, (17.1.1, b))에 의해 b)가 성립한다.

역으로 b)가 a)를 함의함을 보이자. $x_i$들이 매개변수계의 일부를
이루므로 $\dim(A/\mathfrak{J})=n-r$이고 (16.3.6),
$A/\mathfrak{J}$이 정칙이라는 가정에서
$\operatorname{rg}_k(\mathfrak{n}/\mathfrak{n}^2)=n-r$를 얻는다
(17.1.1). 한편 명백히
$\operatorname{rg}_k((\mathfrak{m}^2+\mathfrak{J})/
\mathfrak{m}^2)\leq r$이다. 따라서 (17.1.7.1)에서
$\operatorname{rg}_k(\mathfrak{m}/\mathfrak{m}^2)
\leq r+(n-r)=n$을 얻고, (16.2.6)과 (17.1.1)에 의해 $A$는
정칙이다. 이제
$\operatorname{rg}_k(\mathfrak{m}/\mathfrak{m}^2)=n$이므로
$\operatorname{rg}_k((\mathfrak{m}^2+\mathfrak{J})/
\mathfrak{m}^2)=r$이고, 따라서 $\overline{x}_i$들은 $k$ 위에서 선형
독립이다. C.Q.F.D.
\end{proof}

\oldpage[0\textsubscript{IV-1}]{42}

\begin{corollary}[17.1.8]
\label{0.17.1.8-fr}
$A$를 뇌터 국소환, $\mathfrak{m}$을 그 극대 아이디얼,
$t\in\mathfrak{m}$이라 하자. 다음 조건들은 서로 동치이다.
\begin{enumerate}
  \item[\rm a)] $A/tA$는 정칙이고 $t$는 $A$에서 영인자가 아니다.
  \item[\rm b)] $A$는 정칙이고 $t\notin\mathfrak{m}^2$이다.
\end{enumerate}
\end{corollary}

\begin{proof}
이는 (17.1.7)에서 $r=1$인 특수한 경우이다. 여기서는
$\mathfrak{m}$에 속하고 영인자가 아닌 원소가 매개변수계의 일부를
이룬다는 사실 (16.4.1)을 함께 사용한다.
\end{proof}

\begin{corollary}[17.1.9]
\label{0.17.1.9-fr}
$A$를 정칙 국소환, $\mathfrak{m}$을 그 극대 아이디얼,
$\mathfrak{J}\subset\mathfrak{m}$을 $A$의 아이디얼이라 하자. 다음
조건들은 서로 동치이다.
\begin{enumerate}
  \item[\rm a)] 환 $A/\mathfrak{J}$은 정칙이다.
  \item[\rm b)] $\mathfrak{J}$은 $A$의 한 정칙 매개변수계의 일부를
  이루는 열 $(x_i)_{1\leq i\leq r}$에 의해 생성된다.
\end{enumerate}
\end{corollary}

\begin{proof}
b)가 a)를 함의함은 (17.1.7)에서 이미 보았다. 역으로
$A/\mathfrak{J}$이 정칙이라고 가정하자. 그러면 특히
$\mathfrak{J}$은 소 아이디얼이다. $n=\dim(A)$,
$n-r=\dim(A/\mathfrak{J})$라 하자. (17.1.7)의 표기를 쓰면 완전열
(17.1.7.1)에서
$\operatorname{rg}_k((\mathfrak{m}^2+\mathfrak{J})/
\mathfrak{m}^2)=r$를 얻는다. 따라서 $\mathfrak{J}$의 원소들로 된 열
$(x_i)_{1\leq i\leq r}$로서 $A$의 한 정칙 매개변수계의 일부를 이루는
것이 존재한다. $\mathfrak{J}'=x_1A+\cdots+x_rA$라 두자. (17.1.7)에
의해 $\mathfrak{J}'$은 $A$의 소 아이디얼이고
$A/\mathfrak{J}'$은 차원 $n-r$인 정칙환이다. 그런데
$A/\mathfrak{J}=(A/\mathfrak{J}')/(\mathfrak{J}/\mathfrak{J}')$이고
$\mathfrak{J}/\mathfrak{J}'$은 정역 $A/\mathfrak{J}'$의 소
아이디얼이므로, $A/\mathfrak{J}$과 $A/\mathfrak{J}'$의 차원이 같으려면
$\mathfrak{J}=\mathfrak{J}'$이어야 한다 (16.1.2.2).
\end{proof}

\subsection{가군의 사영 차원과 단사 차원에 관한 복습}
\label{subsection:0.17.2-fr}

\begin{env}[17.2.1]
\label{0.17.2.1-fr}
$A$를 환, $M$을 $A$-가군이라 하자. $M$의 \emph{사영 차원} (각각
\emph{단사 차원})은 길이 $n$인 $M$의 왼쪽 사영 분해 (각각 오른쪽
단사 분해)가 존재하게 하는 가장 작은 $n$ (정수이거나 $+\infty$)이며
(M, VI, 2), 이를 $\operatorname{dim.proj}(M)$ 또는
$\operatorname{dim.proj}_A(M)$ (각각
$\operatorname{dim.inj}(M)$ 또는 $\operatorname{dim.inj}_A(M)$)로
나타냄을 상기하자. 이는 모든 $A$-가군 $N$에 대하여 모든 $i>n$에서,
또는 $i=n+1$에서만,
$\operatorname{Ext}^i_A(M,N)=0$ (각각
$\operatorname{Ext}^i_A(N,M)=0$)이 되게 하는 가장 작은 수가 $n$이라는
것과 같다 (\emph{loc. cit.}).

$M$의 모든 \emph{직접인자} $M'$에 대하여
$\operatorname{Ext}^i_A(M',N)$이 $\operatorname{Ext}^i_A(M,N)$의
직접인자이므로
$\operatorname{dim.proj}(M')\leq\operatorname{dim.proj}(M)$이다.
마찬가지로
$\operatorname{dim.inj}(M')\leq\operatorname{dim.inj}(M)$이다.
$\operatorname{dim.proj}(M)=n$ (각각
$\operatorname{dim.inj}(M)=n$)과 동치인 조건은, 모든 완전열
\[
  0\longrightarrow R\longrightarrow P_{n-1}\longrightarrow\cdots
  \longrightarrow P_0\longrightarrow M\longrightarrow0
\]
(각각
$0\longrightarrow M\longrightarrow Q_0\longrightarrow\cdots
\longrightarrow Q_{n-1}\longrightarrow C\longrightarrow0$)에서
$0\leq i<n$인 모든 $P_i$가 사영이면 (각각 모든 $Q_i$가 단사이면)
$R$이 사영인 (각각 $C$가 단사인) 가장 작은 수가 $n$이라는 것이다.
\end{env}

\begin{remarks}[17.2.2]
\label{0.17.2.2-fr}
\begin{enumerate}
  \item[\rm (i)] $M$이 \emph{사영} (각각 \emph{단사}) $A$-가군이라는
  것은 $\operatorname{dim.proj}(M)=0$ (각각
  $\operatorname{dim.inj}(M)=0$)이라는 것과 같다.
  \item[\rm (ii)] $T$를 $A$-가군의 범주에서 한 아벨 범주로 가는
  가법적 \emph{공변 함자}라 하자. (17.2.1)의 정의와 유도함자의 정의에서
  곧바로, $\operatorname{dim.proj}(M)\leq n$ (각각
  $\operatorname{dim.inj}(M)\leq n$)이면
  \oldpage[0\textsubscript{IV-1}]{43}
  \emph{모든 $i>n$}에 대하여 $L^iT(M)=0$ (각각 $R^iT(M)=0$)임을
  얻는다.
  \item[\rm (iii)] $A$가 \emph{뇌터}이고 $M$이 \emph{유한형}이라고
  가정하자. (17.2.1)에서 준 사영 차원의 마지막 해석에 의해,
  $\operatorname{dim.proj}(M)=n$이면 $M$에는 유한형 사영 가군들로
  이루어진 길이 $n$의 왼쪽 분해가 존재한다. 더 나아가 $A$가 뇌터
  \emph{국소환}이면 $M$에는 유한형 \emph{자유} 가군들로 이루어진 길이
  $n$의 분해가 존재한다. 실제로 이때 유한형 사영 $A$-가군은 자유이다
  (Bourbaki, \emph{Alg.~comm.}, chap.~II, \S~3, no~2, cor.~2 de la
  prop.~5).
\end{enumerate}
\end{remarks}

\begin{lemma}[17.2.3]
\label{0.17.2.3-fr}
$A$를 환, $M$을 $A$-가군이라 하자.
$\operatorname{dim.inj}(M)\leq n$이기 위한 필요충분조건은 한 원소로
생성되는 모든 $A$-가군 $N$에 대하여
$\operatorname{Ext}^{n+1}_A(N,M)=0$인 것이다.
\end{lemma}

\begin{proof}
(17.2.1)의 표기를 쓰면 $C$가 단사임을 보이는 문제이다. 모든
$A$-가군 $N$에 대하여 $\operatorname{Ext}^{n+1}_A(N,M)$은
$\operatorname{Ext}^1_A(N,C)$와 동형이다 (M, V, 7). 따라서 $A$의
임의의 아이디얼 $\mathfrak{J}$에 대하여 $N=A/\mathfrak{J}$ 꼴인 모든
$A$-가군에서 $\operatorname{Ext}^1_A(N,C)=0$이다. 그러면 Ext 완전열에
의해 모든 아이디얼 $\mathfrak{J}$에 대하여 정준 준동형
$\operatorname{Hom}(A,C)\longrightarrow
\operatorname{Hom}(\mathfrak{J},C)$가 전사이다. 이는 $C$가 단사
$A$-가군임을 함의한다 (M, I, 3.2).
\end{proof}

\begin{lemma}[17.2.4]
\label{0.17.2.4-fr}
$A$를 뇌터 환, $M$을 유한형 $A$-가군이라 하자.
$\operatorname{dim.proj}(M)\leq n$이기 위한 필요충분조건은 한 원소로
생성되는 모든 $A$-가군 $N$에 대하여
$\operatorname{Ext}^{n+1}_A(M,N)=0$인 것이다.
\end{lemma}

\begin{proof}
$\operatorname{dim.proj}(M)\leq n$이라는 조건은 모든 \emph{유한형}
$A$-가군 $N$에 대하여 $\operatorname{Ext}^{n+1}_A(M,N)=0$인 것과
동치임이 알려져 있다 (M, VI, 2.5). 명제의 조건도 충분함을 보이기 위해
$N$의 생성원 수 $m$에 관하여 귀납한다. $m-1$개 원소로 생성되는
$N$의 부분가군 $N_1$로서 $N_2=N/N_1$이 한 원소로 생성되는 것을 택할
수 있다. 완전열
$0\longrightarrow N_1\longrightarrow N\longrightarrow N_2
\longrightarrow0$에서 완전열
$\operatorname{Ext}^{n+1}_A(M,N_1)\longrightarrow
\operatorname{Ext}^{n+1}_A(M,N)\longrightarrow
\operatorname{Ext}^{n+1}_A(M,N_2)$를 얻는다. 귀납가정에 의해 명제의
조건은 실제로 $\operatorname{Ext}^{n+1}_A(M,N)=0$을 함의한다.
\end{proof}

\begin{corollary}[17.2.5]
\label{0.17.2.5-fr}
$A$를 뇌터 환, $M$을 $A$-가군이라 하자. 그러면
\begin{equation}
\label{0.17.2.5.1-fr}
\tag{17.2.5.1}
  \operatorname{dim.inj}_A(M)
  =\sup_{\mathfrak{m}}
  \bigl(\operatorname{dim.inj}_{A_{\mathfrak{m}}}(M_{\mathfrak{m}})\bigr)
\end{equation}
이고, $M$이 \emph{유한형}이면
\begin{equation}
\label{0.17.2.5.2-fr}
\tag{17.2.5.2}
  \operatorname{dim.proj}_A(M)
  =\sup_{\mathfrak{m}}
  \bigl(\operatorname{dim.proj}_{A_{\mathfrak{m}}}(M_{\mathfrak{m}})\bigr)
\end{equation}
이다. 여기서 $\mathfrak{m}$은 $A$의 소 아이디얼들의 집합 (또는 극대
아이디얼들의 집합)을 돈다.
\end{corollary}

\begin{proof}
실제로 $N$이 유한형 $A$-가군이면 (따라서 유한 표시이면), $A$의 모든
곱셈적 부분집합 $S$와 모든 $i\geq0$에 대하여
\[
  S^{-1}\operatorname{Ext}^i_A(N,M)
  =\operatorname{Ext}^i_{S^{-1}A}(S^{-1}N,S^{-1}M)
\]
이다. 이는 평탄성과 $N$의 유한형 자유 분해를 생각하고, 앞의 관계가
$i=0$일 때 성립한다는 사실을 사용하면 얻어진다 (Bourbaki,
\emph{Alg.~comm.}, chap.~II, \S~2, no~7, prop.~19). 특히 모든 소
아이디얼 $\mathfrak{m}$과 모든 아이디얼 $\mathfrak{J}$에 대하여
$\operatorname{Ext}^i_{A_{\mathfrak{m}}}
(A_{\mathfrak{m}}/\mathfrak{J}A_{\mathfrak{m}},M_{\mathfrak{m}})
=(\operatorname{Ext}^i_A(A/\mathfrak{J},M))_{\mathfrak{m}}$이다.
(17.2.3)과 $A_{\mathfrak{m}}$의 모든 아이디얼이 적당한 $A$의
아이디얼 $\mathfrak{J}$에 대하여
$\mathfrak{J}A_{\mathfrak{m}}$ 꼴이라는 사실을 고려하면,
\oldpage[0\textsubscript{IV-1}]{44}
앞의 결과와 Bourbaki, \emph{Alg.~comm.}, chap.~II, \S~3, no~3,
th.~1에서 공식 (17.2.5.1)을 얻는다. (17.2.5.2)는 이번에는 (17.2.4)를
사용하고 $M$과 $N$의 역할을 바꾸어 같은 방식으로 증명한다.
\end{proof}

따라서 뇌터 환과 유한형 가군의 사영 차원 또는 단사 차원 연구는
(17.2.5)에 의해 \emph{국소환}의 경우로 귀착된다. 이때 다음이 성립한다.

\begin{lemma}[17.2.6]
\label{0.17.2.6-fr}
$A$를 뇌터 국소환, $k$를 그 잉여체, $M$을 유한형 $A$-가군이라 하자.
$\operatorname{dim.proj}(M)\leq n$이기 위해서는 모든 $i>n$에 대하여
$\operatorname{Tor}^A_i(M,k)=0$이어야 하고,
$\operatorname{Tor}^A_{n+1}(M,k)=0$이면 충분하다.
\end{lemma}

\begin{proof}
필요성은 공변 함자 $M\mapsto k\otimes_A M$에 주석 (17.2.2 (ii))를
적용한 특수한 경우이다. 조건이 충분함을 보이기 위해, (17.2.1)의 표기를
쓰고 $P_i$들이 유한형이라고 가정했을 때 $R$이 사영임을 보이면 된다.
그런데 $\operatorname{Tor}^A_{n+1}(M,k)$는
$\operatorname{Tor}^A_1(R,k)$와 동형이고 (M, V, 7), $R$이 유한형이므로
$\operatorname{Tor}^A_1(R,k)=0$이면 $R$은 자유이다 (Bourbaki,
\emph{Alg.~comm.}, chap.~II, \S~3, no~2, cor.~2 de la prop.~5).
\end{proof}

\begin{corollary}[17.2.7]
\label{0.17.2.7-fr}
(17.2.6)의 가정 아래 $x$를 $A$의 극대 아이디얼에 속하는
$M$-정칙 원소라 하자. 그러면
\begin{equation}
\label{0.17.2.7.1-fr}
\tag{17.2.7.1}
  \operatorname{dim.proj}(M/xM)=\operatorname{dim.proj}(M)+1.
\end{equation}
\end{corollary}

\begin{proof}
실제로 $M$에서 $x$를 곱하여 정의되는 완전열
$0\longrightarrow M\xrightarrow{x}M\longrightarrow M/xM
\longrightarrow0$에서 모든 $i\geq1$에 대한 Tor 완전열
\[
  \operatorname{Tor}^A_i(M,k)\xrightarrow{x}
  \operatorname{Tor}^A_i(M,k)\longrightarrow
  \operatorname{Tor}^A_i(M/xM,k)\longrightarrow
  \operatorname{Tor}^A_{i-1}(M,k)\xrightarrow{x}
  \operatorname{Tor}^A_{i-1}(M,k)
\]
을 얻는다. 그런데 모든 $A$-가군 $N$에 대하여 $M\otimes_A N$에서
$x$를 곱하는 사상은 $M$에서 $x$를 곱하는 사상으로부터도, $N$에서
$x$를 곱하는 사상으로부터도 유도된다. $k$에서 $x$를 곱하는 사상은
정의상 영이므로, 모든 $i$에 대하여 준동형
$\operatorname{Tor}^A_i(M,k)\xrightarrow{x}
\operatorname{Tor}^A_i(M,k)$는 \emph{영}이다. 다시 말해 완전열
\[
  0\longrightarrow\operatorname{Tor}^A_i(M,k)\longrightarrow
  \operatorname{Tor}^A_i(M/xM,k)\longrightarrow
  \operatorname{Tor}^A_{i-1}(M,k)\longrightarrow0
\]
을 얻는다. $\operatorname{dim.proj}(M)=n$이면 (17.2.6)에 의해
$\operatorname{Tor}^A_n(M,k)\neq0$이고
$\operatorname{Tor}^A_{n+1}(M,k)=
\operatorname{Tor}^A_{n+2}(M,k)=0$이다. 따라서 앞의 결과에서
$\operatorname{Tor}^A_{n+1}(M/xM,k)\neq0$이고
$\operatorname{Tor}^A_{n+2}(M/xM,k)=0$을 얻으므로, 다시 (17.2.6)에
의해 $\operatorname{dim.proj}(M/xM)=n+1$이다.
\end{proof}

\begin{proposition}[17.2.8]
\label{0.17.2.8-fr}
\textnormal{(M.~Auslander)} $A$를 환이라 하자. 다음 조건들은 서로
동치이다.
\begin{enumerate}
  \item[\rm a)] 모든 $A$-가군의 사영 차원이 $\leq n$이다.
  \item[\rm a')] 모든 유한형 $A$-가군의 사영 차원이 $\leq n$이다.
  \item[\rm b)] 모든 $A$-가군의 단사 차원이 $\leq n$이다.
  \item[\rm c)] 모든 $A$-가군의 쌍 $M,N$에 대하여
  $\operatorname{Ext}^{n+1}_A(M,N)=0$이다.
  \item[\rm c')] $M$이 유한형인 (또는 한 원소로 생성되기만 하는)
  모든 $A$-가군의 쌍 $M,N$에 대하여
  $\operatorname{Ext}^{n+1}_A(M,N)=0$이다.
\end{enumerate}
\end{proposition}

\begin{proof}
\oldpage[0\textsubscript{IV-1}]{45}
이는 (17.2.1)과 (17.2.3)에서 곧바로 따른다.
\end{proof}

(17.2.8)의 동치인 조건들을 만족하게 하는 가장 작은 수 $n$ (정수 또는
$+\infty$)을 $A$의 \emph{전역 코호몰로지 차원} (또는 단순히
\emph{코호몰로지 차원})이라 하고 $\operatorname{dim.coh}(A)$로
나타낸다.

\begin{proposition}[17.2.9]
\label{0.17.2.9-fr}
$A$를 뇌터 환이라 하자. 다음 조건들은 서로 동치이다.
\begin{enumerate}
  \item[\rm a)] $\operatorname{dim.coh}(A)\leq n$이다.
  \item[\rm b)] 모든 유한형 $A$-가군의 단사 차원이 $\leq n$이다.
  \item[\rm c)] 모든 유한형 $A$-가군의 쌍 $M,N$에 대하여
  $\operatorname{Ext}^{n+1}_A(M,N)=0$이다.
\end{enumerate}
\end{proposition}

\begin{proof}
이는 정의와 (17.2.4)에서 곧바로 따른다.
\end{proof}

\begin{corollary}[17.2.10]
\label{0.17.2.10-fr}
$A$가 뇌터 환이면
\begin{equation}
\label{0.17.2.10.1-fr}
\tag{17.2.10.1}
  \operatorname{dim.coh}(A)
  =\sup_{\mathfrak{m}}
  \bigl(\operatorname{dim.coh}(A_{\mathfrak{m}})\bigr)
\end{equation}
이다. 여기서 $\mathfrak{m}$은 $A$의 스펙트럼 (또는 극대 아이디얼들의
집합)을 돈다.
\end{corollary}

\begin{proof}
이는 (17.2.9)와 (17.2.5)에서 따른다.
\end{proof}

\begin{proposition}[17.2.11]
\label{0.17.2.11-fr}
$A$를 뇌터 국소환, $k$를 그 잉여체라 하자.
$\operatorname{dim.coh}(A)\leq n$이기 위해서는 모든 $i>n$에 대하여
$\operatorname{Tor}^A_i(k,k)=0$이어야 하고,
$\operatorname{Tor}^A_{n+1}(k,k)=0$이면 충분하다.
\end{proposition}

\begin{proof}
(17.2.6)을 고려하면 $\operatorname{dim.coh}(A)\leq n$과
$\operatorname{dim.proj}_A(k)\leq n$이 동치임을 보이면 충분하다.
정의에 의해 첫째 관계가 둘째 관계를 함의함은 명백하다. 역으로
$\operatorname{dim.proj}_A(k)\leq n$이면 (17.2.2, (ii))에 의해 모든
$A$-가군 $M$과 모든 $i>n$에 대하여
$\operatorname{Tor}^A_i(M,k)=0$이다. 따라서
$\operatorname{dim.proj}(M)\leq n$이고, (17.2.8)에 의해 명제가
따른다.
\end{proof}

\begin{corollary}[17.2.12]
\label{0.17.2.12-fr}
$A$를 뇌터 환이라 하자. $\operatorname{dim.coh}(A)\leq n$이기
위해서는 $A$의 모든 극대 아이디얼 $\mathfrak{m}$과 모든 $i>n$에
대하여
$\operatorname{Tor}^{A_{\mathfrak{m}}}_i
(A/\mathfrak{m},A/\mathfrak{m})=0$이어야 하고, 이 관계들이
$i=n+1$일 때 성립하면 충분하다.
\end{corollary}

\begin{proof}
이는 (17.2.11)과 (17.2.10)에서 곧바로 따른다.
\end{proof}

\begin{proposition}[17.2.13]
\label{0.17.2.13-fr}
$A,B$를 두 뇌터 국소환이라 하고,
$\varphi:A\longrightarrow B$를 $B$가 평탄 $A$-가군이 되게 하는 국소
준동형이라 하자. 그러면
\begin{equation}
\label{0.17.2.13.1-fr}
\tag{17.2.13.1}
  \operatorname{dim.coh}(A)\leq\operatorname{dim.coh}(B).
\end{equation}
\end{proposition}

\begin{proof}
$\operatorname{dim.coh}(B)=n$이 유한이라고 가정하자. 모든 유한형
$A$-가군의 쌍 $(M,N)$과 모든 $i>n$에 대하여
$\operatorname{Tor}^A_i(M,N)=0$임을 보이면 충분하다 (17.2.11).
$B$는 충실 평탄 $A$-가군이므로 ($0_{\rm I}$, 6.6.2), 이는
($0_{\rm I}$, 6.4.1)
\begin{equation}
\label{0.17.2.13.2-fr}
\tag{17.2.13.2}
  \operatorname{Tor}^A_i(M,N)\otimes_A B=0
\end{equation}
임을 보이는 것과 같다. $L_\bullet=(L_j)$을 자유 $A$-가군들로 이루어진
$M$의 오른쪽 분해라 하자. $B$가 평탄 $A$-가군이므로
$L_\bullet\otimes_A B=(L_j\otimes_A B)$는 자유 $B$-가군들로 이루어진
$B$-가군 $M\otimes_A B$의 오른쪽 분해이다. 또한
$(L_j\otimes_A B)\otimes_B(N\otimes_A B)
=(L_j\otimes_A N)\otimes_A B$이다. 따라서 (17.2.13.2)의 왼쪽 항은
\oldpage[0\textsubscript{IV-1}]{46}
$\operatorname{Tor}^B_i(M\otimes_A B,N\otimes_A B)$와 같다. $B$에
대한 가정에 의해 이 $B$-가군은 $i>n$이면 영이므로 (17.2.2, (ii))
결론이 따른다.
\end{proof}

\begin{env}[17.2.14]
\label{0.17.2.14-fr}
$(X,\mathscr{O}_X)$를 환 달린 공간, $\mathscr{F}$를
$\mathscr{O}_X$-가군이라 하자. 수
$\sup_{x\in X}(\operatorname{dim.proj}(\mathscr{F}_x))$ (각각
$\sup_{x\in X}(\operatorname{dim.inj}(\mathscr{F}_x))$)을
$\mathscr{F}$의 \emph{점별 사영 차원} (각각 \emph{점별 단사 차원})이라
하고 $\operatorname{dim.proj}(\mathscr{F})$ (각각
$\operatorname{dim.inj}(\mathscr{F})$)로 나타낸다. 수
$\sup_{x\in X}(\operatorname{dim.coh}(\mathscr{O}_x))$를
\emph{$X$의 점별 코호몰로지 차원}이라 하고
$\operatorname{dim.coh}(X)$로 나타낸다. $A$가 뇌터 환, $M$이 유한형
$A$-가군이고 $X=\operatorname{Spec}(A)$이면 (17.2.5)와 (17.2.10)에
의해
$\operatorname{dim.proj}(\widetilde{M})=\operatorname{dim.proj}(M)$,
$\operatorname{dim.inj}(\widetilde{M})=\operatorname{dim.inj}(M)$,
$\operatorname{dim.coh}(X)=\operatorname{dim.coh}(A)$이다.
$\mathscr{F}$의 \emph{점 $x\in X$에서의 사영 차원} (각각 \emph{단사
차원})이란 $\mathscr{F}_x$의 사영 차원 (각각 단사 차원)을 말하고,
\emph{점 $x$에서의 $X$의 코호몰로지 차원}이란 $\mathscr{O}_x$의
코호몰로지 차원을 말한다.
\end{env}

\begin{proposition}[17.2.15]
\label{0.17.2.15-fr}
$X,Y$를 뇌터 국소환들로 환 달린 두 공간,
$f:X\longrightarrow Y$를 평탄 사상이라 하자.
$\operatorname{dim.coh}(X)\leq n$이면 $Y$는 $f(X)$의 모든 점에서
코호몰로지 차원이 $\leq n$이다.
\end{proposition}

\begin{proof}
이는 (17.2.13)에서 곧바로 따른다.
\end{proof}

\subsection{정칙환의 코호몰로지 이론}
\label{subsection:0.17.3-fr}

\begin{theorem}[17.3.1]
\label{0.17.3.1-fr}
\textnormal{(Hilbert--Serre)} $A$를 뇌터 국소환이라 하자. $A$의
코호몰로지 차원이 유한이기 위한 필요충분조건은 $A$가 정칙인 것이다.
이때
\begin{equation}
\label{0.17.3.1.1-fr}
\tag{17.3.1.1}
  \operatorname{dim.coh}(A)=\dim(A)
\end{equation}
이다.
\end{theorem}

\begin{proof}
$A$가 정칙이라고 가정하자. $\mathfrak{m}$을 그 극대 아이디얼,
$\mathbf{x}=(x_i)_{1\leq i\leq n}$을 $A$의 한 정칙 매개변수계라 하자
(17.1.6). 자유 $A$-가군들로 이루어지고 $i>n$이면
$K_i(\mathbf{x})=0$인 외대수 복합체 $K_\bullet(\mathbf{x})$를 생각하자
(\textbf{III}, 1.1.1). $(x_i)$가 정칙열이므로 $i>0$이면
$H_i(K_\bullet(\mathbf{x}))=0$이고 (\textbf{III}, 1.1.4 및 1.1.3.3),
$H_0(K_\bullet(\mathbf{x}))=A/(x_1A+\cdots+x_nA)=A/\mathfrak{m}=k$이다.
따라서 $K_i(\mathbf{x})$들은 길이 $n$인 $k$의 \emph{자유 분해}를
이룬다. 한편 $x_i\in\mathfrak{m}$이므로 복합체
$K_\bullet(\mathbf{x},k)=K_\bullet(\mathbf{x})\otimes_A k$에서는 모든
차수의 경계 연산자가 영이다. 따라서 정의에 의해
$\operatorname{Tor}^A_i(k,k)=\bigwedge^i(k^n)$이다. 이제 (17.2.11)에서
등식 (17.3.1.1)이 곧바로 따른다 (이 결과가 본질적으로 Hilbert의
``syzygy 정리''이다).

이제 $A$가 극대 아이디얼이 $\mathfrak{m}$인 뇌터 국소환이고
$\operatorname{dim.proj}_A(\mathfrak{m})$이 \emph{유한}이면 $A$가
정칙임을 보이자. 그러면 (17.3.1)의 증명이 끝난다.
$n=\operatorname{rg}_k(\mathfrak{m}/\mathfrak{m}^2)$에 관하여
귀납한다. $n=0$이면 $\mathfrak{m}=0$이므로 주장은 자명하다.
\end{proof}

\begin{lemma}[17.3.1.2]
\label{0.17.3.1.2-fr}
\textnormal{(Nagata)} $\mathfrak{m}-\mathfrak{m}^2$의 모든 원소가
$A$에서 영인자이면, $c\mathfrak{m}=0$인 $c\neq0$가 $A$에 존재한다.
다시 말해 $\mathfrak{m}\in\operatorname{Ass}(A)$이다.
\end{lemma}

\begin{proof}
$\mathfrak{m}\neq0$인 경우로 한정해도 되며, 이때
$\mathfrak{m}\neq\mathfrak{m}^2$이다. 가정에 의해
$\mathfrak{m}-\mathfrak{m}^2$는 $\operatorname{Ass}(A)$에 속하는
아이디얼들 $\mathfrak{p}_i$의 합집합에 포함된다 (Bourbaki,
\emph{Alg.~comm.}, chap.~IV, \S~1, no~1, cor.~3 de la prop.~2).
따라서 $\mathfrak{m}$은 $\mathfrak{m}^2$와 $\mathfrak{p}_i$들의
합집합에 포함되고,
\oldpage[0\textsubscript{IV-1}]{47}
$\mathfrak{m}\not\subset\mathfrak{m}^2$이므로 그중 한
$\mathfrak{p}_i$에 포함된다 (Bourbaki, \emph{Alg.~comm.}, chap.~II,
\S~1, no~1, prop.~2). $\mathfrak{m}$은 극대이므로 보조정리가 따른다.
\end{proof}

\begin{lemma}[17.3.1.3]
\label{0.17.3.1.3-fr}
$a\in\mathfrak{m}-\mathfrak{m}^2$이면 $\mathfrak{m}/Aa$는
$\mathfrak{m}/a\mathfrak{m}$의 한 직접인자와 동형이다.
\end{lemma}

\begin{proof}
$a$를 포함하는 $\mathfrak{m}$의 최소 생성계가 존재한다. 그 원소들의
$\mathfrak{m}/\mathfrak{m}^2$에서의 상들은 이 $k$-벡터공간의 기저를
이룬다. 이 생성계에서 $a$ 이외의 원소들이 생성하는 아이디얼을
$\mathfrak{b}$라 하자. 관계 $xa\in\mathfrak{b}$에서
$x\in\mathfrak{m}$이 따른다. 실제로 $xa$의
$\mathfrak{m}/\mathfrak{m}^2$에서의 상을 생각하면 된다. 따라서
$\mathfrak{b}\cap Aa\subset a\mathfrak{m}$이고, 포함
$\mathfrak{b}\longrightarrow\mathfrak{m}$에서 몫으로 내려가면 단사
준동형
$\mathfrak{b}/(\mathfrak{b}\cap Aa)\longrightarrow
\mathfrak{m}/a\mathfrak{m}$을 얻는다. 또한
$\mathfrak{b}+Aa=\mathfrak{m}$이고, 합성 준동형
\[
  \mathfrak{m}/Aa=(\mathfrak{b}+Aa)/Aa
  \xrightarrow{\sim}\mathfrak{b}/(\mathfrak{b}\cap Aa)
  \longrightarrow\mathfrak{m}/a\mathfrak{m}
  \longrightarrow\mathfrak{m}/Aa
\]
은 항등이다. 이로부터 보조정리가 따른다.
\end{proof}

\begin{lemma}[17.3.1.4]
\label{0.17.3.1.4-fr}
$A$를 극대 아이디얼이 $\mathfrak{m}$인 뇌터 국소환, $E$를 사영
차원이 유한인 유한형 $A$-가군이라 하자. $a\in\mathfrak{m}$이
$A$-정칙이고 $E$-정칙이면, $E/aE$는 사영 차원이 유한이고
$\operatorname{dim.proj}_A(E)$ 이하인 $(A/aA)$-가군이다.
\end{lemma}

\begin{proof}
$h=\operatorname{dim.proj}_A(E)$에 관하여 귀납하자. $h=0$인 경우에는
$E$가 사영 $A$-가군이므로 $E/aE$가 사영 $(A/aA)$-가군이어서 자명하다.
$L$이 자유이고 $N$이 유한형이며
$\operatorname{dim.proj}_A(N)=h-1$인 완전열
\[
  0\longrightarrow N\longrightarrow L\longrightarrow E\longrightarrow0
\]
이 존재한다 (17.2.2, (iii)). 또한 완전열
\[
  0\longrightarrow N/aN\longrightarrow L/aL\longrightarrow E/aE
  \longrightarrow0
\]
을 얻는다 (15.1.18). $a$가 $A$-정칙이고 $L$이 자유이므로 $a$는
$N$-정칙이기도 하다. 귀납가정에 의해 $N/aN$은 사영 차원이
$\leq h-1$인 $(A/aA)$-가군이고, $L/aL$은 자유 $(A/aA)$-가군이다.
따라서 $E/aE$는 사영 차원이 $\leq h$인 $(A/aA)$-가군이다.
\end{proof}

이제 두 경우를 살펴보자.

\medskip\noindent\textbf{I.} --- 먼저
$\mathfrak{m}-\mathfrak{m}^2$의 모든 원소가 $A$에서 영인자라고
가정하자. 그러면 (17.3.1.2)에 의해 $c\mathfrak{m}=0$인
$c\neq0$가 $A$에 존재한다. 이 경우 $\mathfrak{m}=0$임을 보이자.
그렇지 않다고 하자. 우선 $\mathfrak{m}$은 사영 $A$-가군일 수 없다.
그렇다면 자유이기 때문에 (Bourbaki, \emph{Alg.~comm.}, chap.~II,
\S~3, no~2, cor.~2 de la prop.~5) 관계 $c\mathfrak{m}=0$과 모순된다.
따라서 $n=\operatorname{dim.coh}(A)\geq1$이다.
$\mathfrak{m}\in\operatorname{Ass}(A)$이므로 $A$-준동형의 완전열
\[
  0\longrightarrow k\longrightarrow A\longrightarrow E\longrightarrow0
\]
이 존재한다. 그러나 이는 불가능하다. 실제로 $n\geq1$이므로 Tor
완전열에서 완전열
$0\longrightarrow\operatorname{Tor}^A_{n+1}(E,k)
\longrightarrow\operatorname{Tor}^A_n(k,k)\longrightarrow0$을 얻는다.
그런데 $\operatorname{Tor}^A_{n+1}(E,k)=0$이고 (17.2.2, (ii)),
$\operatorname{Tor}^A_n(k,k)\neq0$이므로 (17.2.11) 모순이다.

\medskip\noindent\textbf{II.} --- 따라서
$A$-정칙인 $a\in\mathfrak{m}-\mathfrak{m}^2$가 존재하는 경우로
한정할 수 있고, 이 원소는 $\mathfrak{m}$-정칙이기도 하다.
$A'=A/aA$와 그 극대 아이디얼
$\mathfrak{m}'=\mathfrak{m}/aA$를 생각하자.
\oldpage[0\textsubscript{IV-1}]{48}
$\operatorname{rg}_k(\mathfrak{m}'/\mathfrak{m}'^2)=n-1$임은 명백하다.
(17.3.1.4)에 의해 $\mathfrak{m}/a\mathfrak{m}$은 사영 차원이 유한인
$A'$-가군이고, 따라서 그 직접인자인
$\mathfrak{m}'=\mathfrak{m}/aA$도 그러하다 (17.3.1.3 및 17.2.1).
귀납가정에 의해 $A'=A/aA$는 정칙이고, 따라서 (17.1.8)에 의해 $A$도
정칙이다.

\begin{corollary}[17.3.2]
\label{0.17.3.2-fr}
$A$가 정칙 국소환이면 $A$의 모든 소 아이디얼 $\mathfrak{p}$에 대하여
$A_{\mathfrak{p}}$는 정칙이다.
\end{corollary}

\begin{proof}
실제로 (17.2.10)에서
$\operatorname{dim.coh}(A_{\mathfrak{p}})\leq
\operatorname{dim.coh}(A)$임을 보았으므로, 결론은 (17.3.1)에서 곧바로
따른다.
\end{proof}

\begin{proposition}[17.3.3]
\label{0.17.3.3-fr}
$A$, $B$를 두 뇌터 국소환, $\varphi:A\longrightarrow B$를 국소
준동형, $\mathfrak{m}$ (각각 $\mathfrak{n}$)을 $A$ (각각 $B$)의 극대
아이디얼, $k=A/\mathfrak{m}$ (각각 $k'=B/\mathfrak{n}$)를 $A$
(각각 $B$)의 잉여체라 하자. 그러면 $k'$-벡터공간의 표준 준동형
\begin{equation}
\label{0.17.3.3.1-fr}
\tag{17.3.3.1}
  \psi:(\mathfrak{m}/\mathfrak{m}^2)\otimes_k k'
  \longrightarrow\mathfrak{n}/\mathfrak{n}^2
\end{equation}
을 얻는다.
\begin{enumerate}
  \item[\rm (i)] $B$가 정칙이고 $B$가 평탄 $A$-가군이면 $A$는
  정칙이다.
  \item[\rm (ii)] 다음 조건들은 서로 동치이다.
  \begin{enumerate}
    \item[\rm a)] $B$는 정칙이고 준동형 $\psi$ (17.3.3.1)는 단사이다.
    \item[\rm b)] $A$와 $B$는 정칙이고, $A$의 한 정칙 매개변수계
    $(x_i)_{1\leq i\leq m}$에 대하여 $\varphi(x_i)$들은 $B$의 한 정칙
    매개변수계의 일부를 이룬다 (이 경우 이 성질은 $A$의
    \emph{모든} 정칙 매개변수계에 대하여 성립한다).
    \item[\rm c)] $B$와 $B\otimes_A k$는 정칙이고 $B$는 평탄
    $A$-가군이다.
    \item[\rm d)] $A$와 $B\otimes_A k$는 정칙이고 $B$는 평탄
    $A$-가군이다.
    \item[\rm e)] $A$와 $B\otimes_A k$는 정칙이고
    \begin{equation}
    \label{0.17.3.3.2-fr}
    \tag{17.3.3.2}
      \dim(B)=\dim(A)+\dim(B\otimes_A k)
    \end{equation}
    이다.
  \end{enumerate}
\end{enumerate}
\end{proposition}

\begin{proof}
\textnormal{(i)} (17.2.13)에 의해
$\operatorname{dim.coh}(A)\leq\operatorname{dim.coh}(B)$이므로
(17.3.1)을 적용하면 충분하다.

\textnormal{(ii)} $A$가 정칙이라고 가정하고 $(x_i)$를 $A$의 한 정칙
매개변수계라 하자. $B$가 평탄 $A$-가군이라는 것은 (15.1.21)에 의해
$y_i=\varphi(x_i)$의 열이 $B$-정칙이라는 것과 동치이다
($B\otimes_A k$는 평탄 $k$-가군이기 때문이다). 한편
$\dim(A)=m$이고 열 $(y_i)$가 아이디얼 $\mathfrak{m}B$를 생성하므로,
\[
  \dim(B/\mathfrak{m}B)=\dim(B)-m
\]
으로도 쓸 수 있는 관계 (17.3.3.2)는 열
$(y_i)_{1\leq i\leq m}$이 $B$의 한 매개변수계의 일부를 이룬다는 것과
동치이다 (16.3.7). 따라서 조건 b), d), e)는 각각 다음 조건들과
동치이다.
\begin{enumerate}
  \item[\rm b')] $A$와 $B$는 정칙이고 $(y_i)_{1\leq i\leq m}$은
  $B$의 한 정칙 매개변수계의 일부를 이룬다.
  \item[\rm d')] $A$는 정칙이고
  $B/(\sum_{i=1}^m y_iB)$는 정칙이며 열 $(y_i)$는 $B$-정칙이다.
  \item[\rm e')] $A$는 정칙이고
  $B/(\sum_{i=1}^m y_iB)$는 정칙이며 열 $(y_i)$는 $B$의 한
  매개변수계의 일부를 이룬다.
\end{enumerate}
그런데 (17.1.7)에 의해 b')와 e')는 동치이고, d')는 e')를 함의하며
(16.4.1), b')는 d')를 함의하므로 (17.1.7), 이 세 조건은 서로
동치이다. 이제
\oldpage[0\textsubscript{IV-1}]{49}
b)와 d)의 결합은 자명하게 c)를 함의하고,
(i)에 의해 c)는 d)를 함의한다. 따라서 b), c), d), e)의 동치를
증명하였다. 또한 b)는 a)를 함의함이 명백하다. 실제로 $x_i$들의
$\mathfrak{m}/\mathfrak{m}^2$에서의 류들은 이 $k$-벡터공간의
기저를 이루고, $y_i$들의 $\mathfrak{n}/\mathfrak{n}^2$에서의
류들은 이 $k'$-벡터공간의 선형독립계를 이룬다. 따라서 a)가
e)를 함의함을 보이는 일만 남는다. $V=\mathfrak{m}/\mathfrak{m}^2$,
$W=\mathfrak{n}/\mathfrak{n}^2$라 하자. 그러면
\[
  \psi(V\otimes_k k')=(\mathfrak{n}^2+\mathfrak{m}B)/\mathfrak{n}^2
\]
임은 명백하다. 먼저 (16.3.9)에 의해
\begin{equation}
\label{0.17.3.3.3-fr}
\tag{17.3.3.3}
  \dim(B)\leq\dim(A)+\dim(B\otimes_A k)
\end{equation}
이고, 다음으로 (16.2.6)에 의해
\begin{equation}
\label{0.17.3.3.4-fr}
\tag{17.3.3.4}
  \dim(A)\leq\operatorname{rg}_{k'}(V\otimes_k k')
\end{equation}
이다. 마지막으로 국소환 $B\otimes_A k=B/\mathfrak{m}B$에서
$\mathfrak{n}'=\mathfrak{n}/\mathfrak{m}B$는 극대 아이디얼이고
$k'$은 잉여체이며, $\mathfrak{n}'/\mathfrak{n}'^2$은
\[
  \mathfrak{n}/(\mathfrak{n}^2+\mathfrak{m}B)
  =W/\psi(V\otimes_k k')
\]
와 동형이다. 따라서 (16.2.6)에 의해
\begin{equation}
\label{0.17.3.3.5-fr}
\tag{17.3.3.5}
  \dim(B\otimes_A k)
  \leq\operatorname{rg}_{k'}W
  -\operatorname{rg}_{k'}\psi(V\otimes_k k')
\end{equation}
이다. 끝으로 $B$가 정칙이라고 가정했으므로
$\dim(B)=\operatorname{rg}_{k'}W$이다 (17.1.1). 따라서
(17.3.3.3), (17.3.3.4), (17.3.3.5)에서
\[
  \operatorname{rg}_{k'}W
  \leq\dim(A)+\dim(B\otimes_A k)
  \leq\operatorname{rg}_{k'}W
  +\operatorname{rg}_{k'}(V\otimes_k k')
  -\operatorname{rg}_{k'}\psi(V\otimes_k k')
\]
를 얻는다. 이 부등식의 양 끝 항이 같다는 것은 $\psi$가
\emph{단사}라는 뜻이다. 그러므로 조건 a)는 (17.3.3.3),
(17.3.3.4), (17.3.3.5) 각각에서 양변이 \emph{같음}을 필연적으로
함의한다. 그런데 (17.3.3.4)에서의 등식 (각각 (17.3.3.5)에서의
등식)은 $A$ (각각 $B\otimes_A k$)가 \emph{정칙}이라는 뜻이다
(17.1.1). 이로써 a)가 e)를 함의함을 증명하였다.
\end{proof}

\begin{proposition}[17.3.4]
\label{0.17.3.4-fr}
$A$를 차원이 $n$인 정칙 국소환, $\mathfrak{m}$을 그 극대 아이디얼이라
하자. 영이 아닌 모든 유한형 $A$-가군 $M$에 대하여
\begin{equation}
\label{0.17.3.4.1-fr}
\tag{17.3.4.1}
  \operatorname{prof}(M)+\operatorname{dim.proj}(M)=n
\end{equation}
이다.
\end{proposition}

\begin{proof}
$m=\operatorname{prof}(M)$에 관하여 귀납하자. $m=0$이면 (16.4.6,
(i))에 의해 $M$의 부분 $A$-가군으로서 $k=A/\mathfrak{m}$와 동형인
$N$이 존재한다. 완전열
$0\longrightarrow N\longrightarrow M\longrightarrow M/N\longrightarrow0$
에 Tor 완전열을 적용하면 완전열
\[
  \operatorname{Tor}^A_{n+1}(M/N,k)\longrightarrow
  \operatorname{Tor}^A_n(N,k)\longrightarrow
  \operatorname{Tor}^A_n(M,k)
\]
을 얻는다. $A$가 정칙이라는 가정에서
$\operatorname{Tor}^A_{n+1}(M/N,k)=0$이 따르므로 ((17.3.1.1) 및
(17.2.6)), $\operatorname{Tor}^A_n(k,k)$는
$\operatorname{Tor}^A_n(M,k)$의 한 부분 $A$-가군과 동형이다. 다시
(17.3.1.1)과 (17.2.6)을 적용하면
$\operatorname{Tor}^A_n(M,k)\neq0$이고, 따라서
$\operatorname{dim.proj}(M)\geq n$이다 (17.2.6). 한편 (17.3.1.1)에
의해 $\operatorname{dim.proj}(M)\leq n$이므로
$\operatorname{dim.proj}(M)=n$이다.

이제 $m>0$이라고 가정하고 $\mathfrak{m}$에 속하는 $M$-정칙 원소
$x$를 택하자. 그러면 (16.4.6, (ii))에 의해
$\operatorname{prof}(M/xM)=m-1$이고, 한편 (17.2.7)에 의해
$\operatorname{dim.proj}(M/xM)=\operatorname{dim.proj}(M)+1$이다.
귀납가정에서 관계 (17.3.4.1)이 곧바로 따른다.
\end{proof}

\oldpage[0\textsubscript{IV-1}]{50}

\begin{corollary}[17.3.5]
\label{0.17.3.5-fr}
\begin{enumerate}
  \item[\rm(i)] $A$를 차원이 $n$인 정칙 국소환이라 하자. 영이 아닌
  유한형 $A$-가군 $M$이 자유이기 위한 필요충분조건은 $M$이 차원이
  $n$인 Cohen--Macaulay $A$-가군인 것이다.
  \item[\rm(ii)] $A$를 정칙 국소환, $B$를 국소환,
  $\rho:A\to B$를 $B$가 유한형 $A$-가군이 되게 하는 국소 준동형이라
  하자. $B$가 자유 $A$-가군이기 위한 필요충분조건은 $B$가
  Cohen--Macaulay 환이고 $\rho$가 단사인 것이다 (또는 이와 동치로
  (16.1.5), $\dim(B)=\dim(A)$인 것이다).
\end{enumerate}
\end{corollary}

\begin{proof}
(i) 이는 (17.3.4.1)과, 유한형 $A$-가군에 대하여 사영이라는 것과
자유라는 것이 동치라는 사실에서 따른다 (Bourbaki,
\emph{Alg. comm.}, chap.~II, \S~3, n\textsuperscript{o}~2, cor.~2 de la
prop.~5). 따라서 유한형 자유 $A$-가군 $M$은 관계
$\operatorname{dim.proj}(M)=0$으로 특징지어진다 (17.2.2, (i)).

(ii) $B$가 Cohen--Macaulay 환이라는 것은 $B$가 Cohen--Macaulay
$A$-가군이라는 것과 동치이다 (16.5.3). 또한
$\dim(B)=\dim(A)$이므로 (16.1.5), (i)를 적용하면 충분하다.
\end{proof}

\begin{env}[17.3.6]
\label{0.17.3.6-fr}
뇌터환 $A$의 모든 소 아이디얼 $\mathfrak{p}$에 대하여 국소환
$A_{\mathfrak{p}}$가 정칙이면 $A$를 \emph{정칙}이라 한다. $A$ 자체가
국소환이면 (17.3.2)에 의해 이 정의는 (17.1.2)의 정의와 동치이다.
(17.3.2)에 의해, 모든 극대 아이디얼 $\mathfrak{m}$에 대하여
$A_{\mathfrak{m}}$이 정칙이면 $A$가 정칙이기에 충분하다. 또한 이
정의에서 곧바로, $A$의 모든 곱셈적 부분집합 $S$에 대하여
$S^{-1}A$가 정칙임이 따른다.
\end{env}

\begin{proposition}[17.3.7]
\label{0.17.3.7-fr}
$A$가 정칙 뇌터환이면 모든 다항식환
$A[T_1,\ldots,T_n]$은 정칙이다.
\end{proposition}

\begin{proof}
다항식환 $B=A[T]$가 정칙임을 보이면 충분함은 명백하다. $B$는 자유
$A$-가군이므로, $B$의 모든 소 아이디얼 $\mathfrak{q}$에 대하여
$B_{\mathfrak{q}}$는 $A_{\mathfrak{p}}$-평탄이다. 여기서
$\mathfrak{p}=\mathfrak{q}\cap A$이다 ($0_{\mathrm I}$, 6.3.1).
따라서 (17.3.3, (ii))에 의해
$B_{\mathfrak{q}}/\mathfrak{p}B_{\mathfrak{q}}$가 정칙임을 보이면
충분하다. 이 환은
\[
  A_{\mathfrak{p}}[T]/\mathfrak{p}A_{\mathfrak{p}}[T]=k[T]
\]
의 한 소 아이디얼에서의 국소환이다. 여기서
$k=A_{\mathfrak{p}}/\mathfrak{p}A_{\mathfrak{p}}$는 잉여체이다.
그러므로 $C=k[T]$가 정칙임을 보이면 충분하다. 그런데 $C$는 주
아이디얼환이므로 그 소 아이디얼들에서의 국소환들은 이산 값매김환이거나
체이다 (아이디얼 $(0)$의 경우). 따라서 이들은 정칙이고 (17.1.4),
증명이 끝난다.
\end{proof}

\begin{corollary}[17.3.8]
\label{0.17.3.8-fr}
$A$가 정칙환이면 모든 형식적 멱급수환
$A[[T_1,\ldots,T_n]]$은 정칙이다.
\end{corollary}

\begin{proof}
$\mathfrak{J}$을 다항식환 $A[T_1,\ldots,T_n]$에서 $T_i$들이 생성하는
아이디얼이라 하자. 이 다항식환은 (17.3.7)에 의해 정칙이고,
$A[[T_1,\ldots,T_n]]$은 $A[T_1,\ldots,T_n]$의 $\mathfrak{J}$-준아딕
위상에 대한 완비화이므로 (Bourbaki, \emph{Alg. comm.}, chap.~III,
\S~2, n\textsuperscript{o}~12, \emph{Exemple 1}), 결론은 다음
보조정리에서 따른다.
\end{proof}

\begin{lemma}[17.3.8.1]
\label{0.17.3.8.1-fr}
$A$를 정칙환, $\mathfrak{J}$을 $A$의 한 아이디얼,
$\widehat{A}$를 $A$의 $\mathfrak{J}$-준아딕 위상에 대한 분리
완비화라 하자. 그러면 $\widehat{A}$는 정칙이다.
\end{lemma}

\begin{proof}
실제로 (17.3.6)에 의해, $\widehat{A}$의 모든 극대 아이디얼
$\mathfrak{n}$에 대하여 국소환 $\widehat{A}_{\mathfrak{n}}$이 정칙임을
보이면 충분하다.
\oldpage[0\textsubscript{IV-1}]{51}
$\mathfrak{n}$은 $\mathfrak{J}$을 포함하는 $A$의 한
극대 아이디얼 $\mathfrak{m}$에 대하여
$\mathfrak{n}=\mathfrak{m}\widehat{A}$ 꼴이고 (Bourbaki,
\emph{Alg. comm.}, chap.~III, \S~3, n\textsuperscript{o}~4, prop.~8),
표준 준동형 $A\to\widehat{A}$는 \emph{단사} 준동형
\[
  A_{\mathfrak{m}}\longrightarrow\widehat{A}_{\mathfrak{n}}
\]
을 준다는 것을 알고 있다. 이때
$\widehat{A}_{\mathfrak{n}}$의
$\mathfrak{n}\widehat{A}_{\mathfrak{n}}$-준아딕 위상이
$A_{\mathfrak{m}}$에 유도하는 위상은
$\mathfrak{m}A_{\mathfrak{m}}$-준아딕 위상이고,
$A_{\mathfrak{m}}$은 $\widehat{A}_{\mathfrak{n}}$에서
$\mathfrak{n}\widehat{A}_{\mathfrak{n}}$-준아딕 위상에 관하여
조밀하다. 따라서 뇌터 국소환 $A_{\mathfrak{m}}$과
$\widehat{A}_{\mathfrak{n}}$의 \emph{완비화들}은 표준적으로
\emph{동형}이다. 그런데 가정에 의해 $A_{\mathfrak{m}}$은 정칙이고,
따라서 그 완비화도 정칙이다 (17.1.5). 그러므로
$\widehat{A}_{\mathfrak{n}}$의 완비화가 정칙이고, 다시 (17.1.5)에
의해 $\widehat{A}_{\mathfrak{n}}$도 정칙이다.
\end{proof}

\begin{corollary}[17.3.9]
\label{0.17.3.9-fr}
$A$를 정칙 뇌터환 $B$의 몫인 뇌터환이라 하자. $C$가 유한형
$A$-대수이면 $C$의 모든 분수환 $S^{-1}C$는 한 정칙환의 몫이다.
\end{corollary}

\begin{proof}
실제로 $C$는 한 다항식환 $A[T_1,\ldots,T_n]$의 몫이다.
$A=B/\mathfrak{J}$이고 $\mathfrak{J}$이 $B$의 한 아이디얼이므로
\[
  A[T_1,\ldots,T_n]
  =B[T_1,\ldots,T_n]/\mathfrak{J}B[T_1,\ldots,T_n].
\]
따라서 (17.3.7)에 의해 $C$가 한 정칙환 $B'$의 몫인 경우로 한정할 수
있다. $S'$을 $B'$에서 $S$의 역상이라 하면 $S'$은 $B'$의 곱셈적
부분집합이고, $S^{-1}C$는 $S'^{-1}B'$의 한 몫환이다. 그러므로 결국
$C$가 \emph{정칙}일 때 $S^{-1}C$도 정칙임을 보이면 충분하며, 이는
(17.3.6)에서 보았다.
\end{proof}

정칙환의 몫들이 중요한 이유는 특히 앞의 성질과 이들이
\emph{카테너리 환}이라는 사실에 있다 (16.5.12).

대수기하학에 응용되는 중요한 모든 환은 정칙환의 몫이다.

\section{대수의 확장에 관한 보충}
\label{section:0.18-fr}

이 절은 환 위에서의 몇 가지 함자적 구성을 모은다. 이 구성들은
\S\S~19와 20에서 반복하여 사용될 것이며, 여기에는 비자명한 결과가
하나도 들어 있지 않다.

\subsection{증대환의 역상}
\label{subsection:0.18.1-fr}

\begin{env}[18.1.1]
\label{0.18.1.1-fr}
환 $A$ (반드시 가환일 필요는 없다)가 주어졌을 때, \emph{$A$-환}의
범주의 대상은 환 $B$와 환 준동형 $\rho:A\to B$로 이루어진 쌍
$(B,\rho)$이고, 사상 (\emph{$A$-준동형}이라고도 한다)은 다음 조건을
만족하는 환 준동형 $u:B\to C$이다. 즉 $\rho:A\to B$와
$\sigma:A\to C$가 각각 $B$와 $C$의 $A$-환 구조를 정의하는
준동형 (\emph{구조 준동형}이라고 한다)이면, 도식
\begin{equation}
\label{0.18.1.1.1-fr}
\tag{18.1.1.1}
  \xymatrix{
    B\ar[r]^u & C\\
    & A\ar[lu]^\rho\ar[u]_\sigma
  }
\end{equation}
은 가환이다. $u$의 \emph{핵} $\mathfrak{J}$은 양쪽 아이디얼이고
($\rho$를 통해) $A$-\emph{쌍가군}이다.

$f:A'\to A$가 환 준동형이면, 모든 $A$-환 $(B,\rho)$에 대하여
$(B,\rho\circ f)$는 $A'$-환이다. 또한 $u:B\to C$가 $A$-환
$(B,\rho)$에서 $A$-환 $(C,\sigma)$로 가는 $A$-준동형이면
\oldpage[0\textsubscript{IV-1}]{52}
이는 $A'$-환 $(B,\rho\circ f)$에서 $A'$-환
$(C,\sigma\circ f)$로 가는 $A'$-준동형이기도 하다. 이로써
$A$-환의 범주에서 $A'$-환의 범주로 가는 표준 함자가 정의된다.
\end{env}

\begin{env}[18.1.2]
\label{0.18.1.2-fr}
$B$, $E$, $F$를 세 $A$-환이라 하고, $f:E\to B$, $g:F\to B$를
두 $A$-준동형이라 하자. $E\times_BF$로 나타내는 $E$와 $F$의
$B$ 위 \emph{섬유곱} (즉 $A$-준동형 $f$와 $g$에 관한 섬유곱)은
$f(x)=g(y)$를 만족하는 쌍 $(x,y)$들로 이루어진 곱환 $E\times F$의
부분환 $G$임을 상기하자. 사영 $\operatorname{pr}_1$과
$\operatorname{pr}_2$를 $G$에 제한한 $p_1:G\to E$, $p_2:G\to F$도
\emph{표준 사영}이라 한다. $G$의 $A$-환 구조는 준동형
$\alpha\mapsto(\rho(\alpha),\sigma(\alpha))$로 정의된다. 여기서
$\rho:A\to E$와 $\sigma:A\to F$는 구조 준동형이며, 위 준동형은
(18.1.1)에 의해 실제로 $A$를 $G$ 안으로 보낸다. $E\times_BF$의
특징짓는 성질은 다음과 같다. $f\circ u=g\circ v$를 만족하는 임의의
두 $A$-준동형 $u:C\to E$, $v:C\to F$에 대하여
$u=p_1\circ w$, $v=p_2\circ w$를 만족하는 $A$-준동형 $w:C\to G$가
유일하게 존재한다. 다시 말해 $E\times_BF$는 $A$-환의 범주에서
$B$, $E$, $F$와 $A$-준동형 $f$, $g$로 이루어진 사영계의
\emph{사영극한}이다 ($0_{\mathrm{III}}$, 8.1.9).
\end{env}

\begin{env}[18.1.3]
\label{0.18.1.3-fr}
$\mathfrak{J}$을 $f$의 핵인 $E$의 양쪽 아이디얼이라 하자.
$p_2:G\to F$의 핵 $\mathfrak{J}'$은 $x\in\mathfrak{J}$인 원소
$(x,0)$들로 이루어진 아이디얼임이 자명하다. 따라서 $p_1$의 제한
$i_1:\mathfrak{J}'\to\mathfrak{J}$은 \emph{단위원을 갖지 않는
환으로서의 동형사상}이고, 또한 $A$-쌍가군의 동형사상이다. 마찬가지로
$\mathfrak{R}$이 $g$의 핵이면, $p_1:G\to E$의 핵
$\mathfrak{R}'$은 $y\in\mathfrak{R}$인 원소 $(0,y)$들로 이루어진
아이디얼이고, $p_2$의 제한 $i_2:\mathfrak{R}'\to\mathfrak{R}$은
(같은 의미에서) 동형이다. 끝으로 $E\times_BF$에서 $B$로 가는
$A$-준동형 $q=f\circ p_1=g\circ p_2$의 핵은
$\mathfrak{J}\times\mathfrak{R}=\mathfrak{J}'\oplus\mathfrak{R}'$
임이 명백하므로, 다음 가환 도식을 얻는다.
\begin{equation}
\label{0.18.1.3.1-fr}
\tag{18.1.3.1}
  \xymatrix@C=2.2pc@R=1.8pc{
    0\ar[dr] & & 0\ar[d] & 0\ar[d]\\
    & \mathfrak{J}'\oplus\mathfrak{R}'\ar[r]\ar[d]\ar[dr]
      & \mathfrak{R}'\ar[r]^{i_2}_{\sim}\ar[d]^j
      & \mathfrak{R}\ar[d]^{p_2}\\
    0\ar[r] & \mathfrak{J}'\ar[r]\ar[d]^{i_1}_{\sim}
      & G\ar[r]\ar[d]^{p_1}\ar[dr]^q & F\ar[d]^g\\
    0\ar[r] & \mathfrak{J}\ar[r] & E\ar[r]_f & B
  }
\end{equation}

(18.1.2)와 (18.1.3)의 정의와 결과는 $A$-준동형들의 족
$f_\lambda:E_\lambda\to B$로 정의되는 임의의 $A$-환의 족
$(E_\lambda)_{\lambda\in I}$의 섬유곱에도 곧바로 확장된다. 이
결과들의 정식화는 독자에게 맡긴다.
\end{env}

\begin{env}[18.1.4]
\label{0.18.1.4-fr}
(M, VIII)의 용어법에 따라, 전사 $A$-준동형 $f:E\to B$가 주어진
$A$-환 $E$를 \emph{$B$ 위에 증대된 $A$-환}이라 한다. 이때 $f$를
$E$의 \emph{증대}라 하고, $f$의 핵 $\mathfrak{J}$을
\emph{증대 아이디얼}이라 한다. 증대 $f:E\to B$의 \emph{오른쪽 역}인
$A$-환의 $A$-준동형 $s:B\to E$ (즉 $f\circ s=1_B$)가 존재하면
증대 $A$-환 $E$를 \emph{자명하다}고 한다. 그러면 $A$-쌍가군의 완전열
\[
  0\longrightarrow\mathfrak{J}\longrightarrow E
  \xrightarrow{f}B\longrightarrow0
\]
\oldpage[0\textsubscript{IV-1}]{53}
은 \emph{분할된다}. 다시 말해 $A$-쌍가군 $E$를
$B\times\mathfrak{J}$과 동일시할 수 있고, 이때 $E$에서의 곱셈은
\[
  (b,z)(b',z')=(bb',bz'+zb'+zz')
\]
로 주어진다. 여기서 $\mathfrak{J}$은 $s:B\to E$를 통해
$B$-쌍가군으로 본다.
\end{env}

\begin{env}[18.1.5]
\label{0.18.1.5-fr}
(18.1.2)의 표기에서 $A$-준동형 $f$가 \emph{전사}, 다시 말해
$E$를 $B$ 위에 \emph{증대된} $A$-환으로 만든다고 하자. 그러면
$p_2:G\to F$도 \emph{전사}임이 명백하다. 즉 이 사상은 $G$에
$F$ 위에 \emph{증대된} $A$-환의 구조를 정의하며, 이를
$g:F\to B$에 의한 증대환 $E$의 \emph{역상}이라 한다.
\end{env}

\begin{proposition}[18.1.6]
\label{0.18.1.6-fr}
$g:F\to B$에 의한 증대 $A$-환 $E$의 역상
$G=E\times_BF$가 자명한 증대 $A$-환이기 위한 필요충분조건은 도식
\[
  \xymatrix{
    & F\ar[ld]_u\ar[d]^g\\
    E\ar[r]_f & B
  }
\]
을 가환으로 만드는 $A$-준동형 $u:F\to E$가 존재하는 것이다.
\end{proposition}

\begin{proof}
$p_2$의 오른쪽 역인 $A$-준동형 $s:F\to G$에 대하여
$u=p_1\circ s$로 잡으면 되므로 이 조건이 필요함은 명백하다.
거꾸로 명제의 조건을 만족하는 $A$-준동형 $u$가 존재하면,
$A$-준동형 $u:F\to E$와 $1_F:F\to F$에 섬유곱의 보편 성질
(18.1.2)를 적용하여 $p_2$의 오른쪽 역 $s$가 존재함을 얻는다.

특히 이 결과로부터 $E$가 \emph{자명한} 증대환이면 그 모든 역상도
자명함을 얻는다.
\end{proof}

\begin{env}[18.1.7]
\label{0.18.1.7-fr}
(18.1.2)와 (18.1.3)의 상황으로 돌아가자. $\mathfrak{R}$에는 그
$F$-쌍가군 구조와 환 준동형 $p_2:G\to F$에서 오는 $G$-쌍가군
구조가 명백히 존재한다. 또한 $i_2$가 전단사이므로 단사 사상
$\theta=j\circ i_2^{-1}:\mathfrak{R}\to G$를 얻으며, 이는
$G$-쌍가군의 준동형이다. 거꾸로 $f$와 $g$가 \emph{전사}, 다시 말해
$E$와 $F$가 $B$ 위에 \emph{증대된} $A$-환이라고 더 가정할 때에는,
이러한 준동형 $\theta$가 주어지면 $F$ 위의 증대환 $G$로부터
$B$ 위의 증대환 $E$를 복원할 수 있음을 보이겠다.

더 정확히는 $B$ 위에 증대된 $A$-환 $F$와 $F$ 위에 증대된
$A$-환 $G$가 주어졌다고 하자. 각각의 증대 아이디얼을
$\mathfrak{R}$과 $\mathfrak{J}'$으로 나타내면 다음 도식을 얻는다.
\begin{equation}
\label{0.18.1.7.1-fr}
\tag{18.1.7.1}
  \xymatrix@C=2pc@R=1.5pc{
    & & 0\ar[d] & & \\
    & & \mathfrak{R}\ar[ld]^\theta\ar[d] & & \\
    0\ar[r] & \mathfrak{J}'\ar[r] & G\ar[r]_h & F\ar[r]\ar[d]^g & 0\\
    & & & B\ar[d] & \\
    & & & 0 &
  }
\end{equation}

\oldpage[0\textsubscript{IV-1}]{54}
$h$는 $F$의 모든 아이디얼 (특히 $\mathfrak{R}$) 위에 $G$-쌍가군
구조를 정의한다. 도식 (18.1.7.1)을 가환으로 만드는 $G$-쌍가군
준동형 $\theta:\mathfrak{R}\to G$가 주어졌다고 하자. 그러면
$\theta$는 \emph{단사}이고
$\mathfrak{J}'\cap\theta(\mathfrak{R})=0$이다. 또한
$\theta(\mathfrak{R})$은 $G$의 양쪽 아이디얼이고
$h(\theta(\mathfrak{R}))=\mathfrak{R}$이므로, 합성 준동형
$g\circ h:G\to B$는
$G\xrightarrow{q}G/\theta(\mathfrak{R})\xrightarrow{f}B$로 분해된다.
여기서 $f$는 전사이다. 더 나아가
$E=G/\theta(\mathfrak{R})$에서 $\mathfrak{J}'$의 상
$\mathfrak{J}$은 $f$의 핵이고 ($g\circ h$의 핵은
$\mathfrak{J}'\oplus\theta(\mathfrak{R})$이다), 표준 준동형
$q:G\to E$를 $\mathfrak{J}'$에 제한한 것은 단사이다. 따라서
\emph{섬유곱} $G'=E\times_BF$를 만들 수 있다. 두 $A$-준동형
$q:G\to E$와 $h:G\to F$가 $f\circ q=g\circ h$를 만족하므로,
$G'$의 보편 성질 (18.1.2)에 따라 유일한 $A$-준동형
$u:G\to G'$을 정의한다. 이제 $u$가 \emph{전단사}임을 보이겠다.
$u$를 $A$-쌍가군의 준동형으로 보아 이를 증명하면 충분하다. 이때
$u$는 각각 $G$, $\mathfrak{J}'$과 $G'$,
$\mathfrak{J}''=\operatorname{Ker}(G'\to F)$으로 이루어진 유한
여과들과 양립한다. 또한 (18.1.3)에서 보았듯이
$\operatorname{gr}_0u:F\to F$는 항등사상이고
$\operatorname{gr}_1u:\mathfrak{J}'\to\mathfrak{J}''$은
전단사이다. 따라서 $u$ 자체가 전단사이다 (Bourbaki,
\emph{Alg. comm.}, chap.~III, \S~2, n\textsuperscript{o}~8,
cor.~3 du th.~1).
\end{env}

\subsection{쌍가군에 의한 환의 확장}
\label{subsection:0.18.2-fr}

\begin{env}[18.2.1]
\label{0.18.2.1-fr}
$E$를 $B$ 위에 증대된 $A$-환, $f:E\to B$를 그 증대,
$\mathfrak{J}=\operatorname{Ker}(f)$를 증대 아이디얼이라 하자.
$\mathfrak{J}^2=0$이면 $\mathfrak{J}$은 $E$-쌍가군일 뿐 아니라
$B$가 $E/\mathfrak{J}$과 동형이므로 $B$-\emph{쌍가군}이기도 하다.
더 정확히 말하면, 모든 $b\in B$는 어떤 $x\in E$에 대하여
$f(x)$ 꼴이고, $z,z'\in\mathfrak{J}$이면
$(x+z)z'=xz'$ (각각 $z'(x+z)=z'x$)이다. 따라서 $xz'$ (각각
$z'x$)의 값은 $f^{-1}(b)$의 원소 $x$에 의존하지 않으므로 이를
$bz'$ (각각 $z'b$)로 쓸 수 있으며, 이것이 위의 $B$-쌍가군 구조를
정의한다. 거꾸로 $x\in E$, $z'\in\mathfrak{J}$에 대하여
$xz'=f(x)z'$와 $z'x=z'f(x)$가 되도록 $\mathfrak{J}$에
$B$-쌍가군 구조가 주어지면 $\mathfrak{J}^2=0$임이 명백하다.
\end{env}

\begin{env}[18.2.2]
\label{0.18.2.2-fr}
$B$-쌍가군 $L$에 의한 $A$-환 $B$의 \emph{$A$-확장}이라 함은
$A$-쌍가군 준동형들의 \emph{완전열}
\[
  0\longrightarrow L\xrightarrow{j}E\xrightarrow{f}B\longrightarrow0
\]
로서, $E$는 $A$-환이고 $f$는 환의 $A$-준동형이며, 모든
$x\in E$, $z\in L$에 대하여
\[
  j(f(x)z)=xj(z),\qquad j(zf(x))=j(z)x
\]
가 성립하는 것을 말한다. 이로부터 (18.2.1)에 의해 $j(L)$은
$E$의 제곱영인 양쪽 아이디얼이다. 용어를 남용하여 $E$ 자체도
$L$에 의한 $B$의 확장이라고 한다. $L$에 의한 $B$의 두
$A$-확장 $E$, $E'$에 대하여 도식
\begin{equation}
\label{0.18.2.2.1-fr}
\tag{18.2.2.1}
  \xymatrix@C=2.1pc@R=1.2pc{
    & & E\ar[dr]\ar[dd]_u\ar@{}[dd]|{\wr} & & \\
    0\ar[r] & L\ar[ur]\ar[dr] & & B\ar[r] & 0\\
    & & E'\ar[ur] & &
  }
\end{equation}
을 가환으로 만드는 $A$-환의 동형
$u:E\xrightarrow{\sim}E'$ (\emph{$A$-확장의 $A$-동치}라고도 한다)가
존재하면 $E$와 $E'$을 \emph{$A$-동치}라고 한다.
\end{env}

\oldpage[0\textsubscript{IV-1}]{55}

\begin{env}[18.2.3]
\label{0.18.2.3-fr}
$B$-쌍가군 $L$에 의한 $B$의 $A$-확장 $E$가 자명한 (즉 $B$ 위에
증대된) $A$-환이면 이 확장을 \emph{$A$-자명}하다고 한다 (18.1.4).
곱 $A$-쌍가군 $B\times L$ 위에
$(x,s)(y,t)=(xy,xt+sy)$로 놓아 $A$-환 구조를 정의할 수 있음은
곧바로 확인된다. 표준 사상 $j:L\to B\times L$과
$f:B\times L\to B$는 $L$에 의한 $B$의 $A$-자명한 $A$-확장을
정의하고, 표준 사상 $g:B\to B\times L$은 $f$의 오른쪽 역인
$A$-준동형이다. 이 확장을 $L$에 의한 $B$의 \emph{표준 자명 확장}이라
하고 $D_B(L)$로 나타낸다. $L$에 의한 $B$의 모든 $A$-자명한
$A$-확장은 $D_B(L)$과 $A$-동치임이 명백하다.

$A$-환 $A$ 자체의 $A$-쌍가군에 의한 모든 확장은 반드시
$A$-자명하다는 것에 유의하자.
\end{env}

\begin{env}[18.2.4]
\label{0.18.2.4-fr}
두 $A$-확장
$L\xrightarrow{j}E\xrightarrow{f}B$와
$L'\xrightarrow{j'}E'\xrightarrow{f'}B'$이 주어졌다고 하자. 첫째
확장에서 둘째 확장으로 가는 \emph{사상}은 정의상 도식
\begin{equation}
\label{0.18.2.4.1-fr}
\tag{18.2.4.1}
  \xymatrix{
    0\ar[r] & L\ar[r]^j\ar[d]_w & E\ar[r]^f\ar[d]_u
      & B\ar[r]\ar[d]_v & 0\\
    0\ar[r] & L'\ar[r]_{j'} & E'\ar[r]_{f'} & B'\ar[r] & 0
  }
\end{equation}
을 가환으로 만드는 $A$-쌍가군 준동형들의 삼중항 $(u,v,w)$을 말한다.
여기서 $u$와 $v$는 환의 $A$-준동형이고, $z\in L$, $b\in B$에
대하여 $w(bz)=v(b)w(z)$, $w(zb)=w(z)v(b)$이다. 다시 말해
쌍 $(v,w)$는 $B$-쌍가군 $L$에서 $B'$-쌍가군 $L'$로 가는
\emph{디-준동형}을 이룬다. $(u',v',w')$이
$L'\to E'\to B'$에서 $A$-확장 $L''\to E''\to B''$으로 가는
사상이면 $(u'\circ u,v'\circ v,w'\circ w)$도 사상임이 명백하며,
이것이 위 용어법을 정당화한다.

도식 (18.2.4.1)의 두 가환 정사각형을 고찰하면 $A$-환의 확장들에
대한 두 연산에 이르게 된다.
\end{env}

\begin{env}[18.2.5]
\label{0.18.2.5-fr}
먼저 $L'$에 의한 $B'$의 $A$-확장 $E'$
\[
  0\longrightarrow L'\xrightarrow{j'}E'\xrightarrow{f'}B'\longrightarrow0
\]
과 환의 $A$-준동형 $v:B\to B'$을 생각하자. 증대 $A$-환 $E'$의
$v$에 의한 \emph{역상}을 $F=E'\times_{B'}B$라 하면 (18.1.5),
다음 가환 도식을 얻는다.
\[
  \xymatrix{
    0\ar[r] & L_0\ar[r]\ar[d]^i\ar@{}[d]|{\wr} & F\ar[r]^{p_2}\ar[d]^{p_1}
      & B\ar[r]\ar[d]^v & 0\\
    0\ar[r] & L'\ar[r]_{j'} & E'\ar[r]_{f'} & B'\ar[r] & 0
  }
\]
두 행은 완전열이고 $p_1$, $p_2$는 표준 준동형이다. (18.1.3)에서
$i$가 전단사임을 보았고, (18.1.2)의 정의로부터 $L_0^2=0$도
따른다. 따라서 $F$를 $L'$에 의한 $B$의 $A$-확장으로 볼 수 있으며,
이를 $L'$에 의한 $B'$의 확장 $E'$의 \emph{$v$에 의한 역상}이라
한다. 여기서 $L'$은 환 준동형 $v$를 통해 자연스럽게 $B$-쌍가군으로
\oldpage[0\textsubscript{IV-1}]{56}
본다. 또한 각 인자에 관한 섬유곱의 함자성으로부터, $B'$의 두 확장
사이의 사상
\[
  \xymatrix{
    0\ar[r] & L'_1\ar[r]\ar[d]^h & E'_1\ar[r]\ar[d]^g
      & B'\ar[r]\ar[d]^{1_{B'}} & 0\\
    0\ar[r] & L'_2\ar[r] & E'_2\ar[r] & B'\ar[r] & 0
  }
\]
이 주어지면 그 \emph{$v$에 의한 역상} 사상
\[
  \xymatrix{
    0\ar[r] & L'_1\ar[r]\ar[d]^h & E'_1\times_{B'}B\ar[r]
      \ar[d]^{g\times_{B'}1_B} & B\ar[r]\ar[d]^{1_B} & 0\\
    0\ar[r] & L'_2\ar[r] & E'_2\times_{B'}B\ar[r] & B\ar[r] & 0.
  }
\]
을 얻는다. 특히 $E'_1$과 $E'_2$가 $L'$에 의한 $B'$의
$A$-동치인 $A$-확장들이면, 이들의 $v$에 의한 역상들도 $L'$에 의한
$B$의 $A$-동치인 $A$-확장들이다.

섬유곱의 정의에 따르면 $A$-확장들의 사상 (18.2.4.1)은 $E'$의
$v$에 의한 역상을 거쳐 \emph{인수분해된다}. 더 정확히 말하면 도식
\[
  \xymatrix{
    0\ar[r] & L\ar[r]^j\ar[d]^{w_0} & E\ar[r]^f\ar[d]^{u_0}
      & B\ar[r]\ar[d]^{1_B} & 0\\
    0\ar[r] & L_0\ar[r]\ar[d]^i\ar@{}[d]|{\wr} & F\ar[r]^{p_2}\ar[d]^{p_1}
      & B\ar[r]\ar[d]^v & 0\\
    0\ar[r] & L'\ar[r]_{j'} & E'\ar[r]_{f'} & B'\ar[r] & 0
  }
\]
을 가환으로 만드는 유일한 $A$-준동형
$u_0:E\to F=E'\times_{B'}B$가 존재한다. 여기서 $w_0$은
$u_0$를 $L$에 제한한 것이고 $p_1\circ u_0=u$, $i\circ w_0=w$이다.
\end{env}

\begin{env}[18.2.6]
\label{0.18.2.6-fr}
특히 전사 $A$-준동형에 의한 확장의 역상들을 살펴보자. 전사
$A$-준동형 $v:B\to B'$, $B$-쌍가군 $L$에 의한 $B$의 $A$-확장
$F$, 그리고 $v$의 핵인 $B$의 양쪽 아이디얼 $\mathfrak{R}$을
생각하자. $\mathfrak{R}$은 증대 $F\to B$를 통해 $F$-쌍가군으로
볼 수 있다. 도식
\[
  \xymatrix{
    & \mathfrak{R}\ar[ld]^\theta\ar[d]\\
    F\ar[r] & B
  }
\]
을 가환으로 만드는 모든 $F$-쌍가군 준동형
$\theta:\mathfrak{R}\to F$는, 그 $v:B\to B'$에 의한 역상이
$F$와 동치인 $L$에 의한 $B'=B/\mathfrak{R}$의 확장 $E'$을
정한다는 것을 (18.1.7)에서 보았다. 또한 이 마지막 성질을 가지는
$L$에 의한 $B'$의 모든 $A$-확장 $E'$은 (하나의 $A$-동치를
제외하고) 이와 같이 얻어진다. 더 나아가 두 $F$-쌍가군 준동형
$\theta_1,\theta_2:\mathfrak{R}\to F$가 $L$에 의한 $B'$의
$A$-동치인 두 $A$-확장을 주기 위한 필요충분조건은
$\theta_2=u\circ\theta_1$을 만족하는 $A$-확장 $F$의 자기
$A$-동치 $u$가 존재하는 것이다.
\oldpage[0\textsubscript{IV-1}]{57}
이는 (18.2.5)에서 본 것과 $F$에서 섬유곱
$E'\times_{B'}B$로 가는 표준 전단사의 정의 (18.1.7)로부터 곧바로
따른다.
\end{env}

\begin{env}[18.2.7]
\label{0.18.2.7-fr}
이제 (18.2.4.1)의 왼쪽 사각형을 생각하자. 먼저 $A$-쌍가군의
범주에서 \emph{밂}의 개념을 상기하자. 세 $A$-쌍가군
$X$, $Y$, $Z$와 두 $A$-준동형 $f:X\to Y$, $g:X\to Z$가
주어졌을 때, 밂 $Y\oplus_XZ$는 $A$-쌍가군의 범주에서
$X$, $Y$, $Z$와 $A$-준동형 $f$, $g$로 이루어진 귀납계의
\emph{귀납극한}이다 ($0_{\mathrm{III}}$, 8.1.11). 이 $A$-쌍가군은
곱 $Y\times Z$를 준동형
$x\mapsto(f(x),-g(x))$에 의한 $X$의 상인 부분 $A$-쌍가군
$M$으로 나눈 몫으로 정의된다. 이를 특징짓는 성질은 다음과 같다.
$u\circ f=v\circ g$를 만족하는 임의의 두 $A$-쌍가군 준동형
$u:Y\to T$, $v:Z\to T$에 대하여, $u=w\circ j_1$ 및
$v=w\circ j_2$를 만족하는 준동형 $w:Y\oplus_XZ\to T$가
유일하게 존재한다. 여기서
$j_1:Y\to Y\oplus_XZ$와 $j_2:Z\to Y\oplus_XZ$는 표준 사상이다.
\end{env}

\begin{env}[18.2.8]
\label{0.18.2.8-fr}
이제 $B$-쌍가군 $L$에 의한 $B$의 $A$-확장 $E$
\[
  0\longrightarrow L\xrightarrow{j}E\xrightarrow{f}B\longrightarrow0
\]
과 $B$-쌍가군 준동형 $w:L\to L'$을 생각하자. $H$를 밂인
$A$-쌍가군 $E\oplus_LL'$이라 하자. 이 $A$-쌍가군에 $A$-환 구조를
주고 $A$-확장
\[
  0\longrightarrow L'\longrightarrow H\longrightarrow B\longrightarrow0
\]
을 정의할 수 있음을 보이자.

이를 위해 $L'$에는 증대 준동형 $E\to B$를 통해 $E$-\emph{쌍가군}
구조가 주어져 있음을 주목하자. 따라서 표준 자명 $A$-확장
$G=D_E(L')$을 만들 수 있다 (18.2.3). 이제 $L$에서 $G$로 가는 사상
$\theta:z\mapsto(j(z),-w(z))$를 생각하자. 이는 $G$-쌍가군의
준동형이다. 여기서 $L$은 표준 준동형 $\rho:G\to E$를 통해
$G$-쌍가군으로 본다. 실제로 $(x,z')\in G$와 $z\in L$에 대하여
\[
  (j(z),-w(z))(x,z')=(j(z)x,j(z)z'-w(z)f(x)).
\]
$L'$ 위의 $E$-쌍가군 구조의 정의에 의해
$j(z)z'=f(j(z))z'=0$이고, $L'$ 위의 $B$-쌍가군 구조의 정의에
의해 $w(z)f(x)=w(zf(x))$이다. 같은 방식으로 $\theta$가 왼쪽
$G$-가군의 준동형임을 확인한다. 이제 가환도식
\[
  \xymatrix@C=2pc@R=1.4pc{
    & & 0\ar[d] & & \\
    & & L\ar[ld]^\theta\ar[d]^j & & \\
    0\ar[r] & L'\ar[r] & G\ar[r]_p & E\ar[r]\ar[d]^f & 0\\
    & & & B\ar[d] & \\
    & & & 0 &
  }
\]
에 (18.1.7)의 결과를 적용할 수 있다. 정의에 의해
$H=G/\theta(L)$이므로, 우리의 주장은 (18.1.7)에서 곧바로 따른다.

\oldpage[0\textsubscript{IV-1}]{58}

$L'$에 의한 $B$의 $A$-확장 $H$를 준동형 $w:L\to L'$을 통해
$E$에서 \emph{유도된} 확장이라 한다. 밂의 각 항에 대한
함자성으로부터, 확장들의 사상
\[
  \xymatrix{
    0\ar[r] & L\ar[r]\ar[d]^{1_L}\ar@{}[d]|{\wr}
      & E_1\ar[r]\ar[d]^g & B_1\ar[r]\ar[d]^h & 0\\
    0\ar[r] & L\ar[r] & E_2\ar[r] & B_2\ar[r] & 0
  }
\]
이 주어지면 표준적으로 확장들의 사상
\[
  \xymatrix{
    0\ar[r] & L'\ar[r]\ar[d]^{1_{L'}}\ar@{}[d]|{\wr}
      & E_1\oplus_LL'\ar[r]\ar[d]^{g\oplus_L1_{L'}}
      & B_1\ar[r]\ar[d]^h & 0\\
    0\ar[r] & L'\ar[r] & E_2\oplus_LL'\ar[r] & B_2\ar[r] & 0.
  }
\]
을 얻는다. 특히 $E_1$과 $E_2$가 $L$에 의한 $B$의 $A$-동치인
$A$-확장들이면, $w$를 통해 이들로부터 유도된 $L'$에 의한 $B$의
확장들도 $A$-동치이다.

$A$-확장들의 사상 (18.2.4.1)이 주어졌을 때, 이 사상은 준동형
$w:L\to L'$을 통해 $E$에서 유도된 $L'$에 의한 $B$의 $A$-확장
$H$를 거쳐 \emph{인수분해된다}. 여기서 $L'$은 준동형
$v:B\to B'$를 통해 $B$-쌍가군으로 본다. 실제로 밂의 정의에
따르면 도식
\[
  \xymatrix{
    0\ar[r] & L\ar[r]^j\ar[d]^{w_0} & E\ar[r]^f\ar[d]^{j_1}
      & B\ar[r]\ar[d]^{1_B}\ar@{}[d]|{\wr} & 0\\
    0\ar[r] & L'_0\ar[r]^{j_2}\ar[d]^{j_0}\ar@{}[d]|{\wr}
      & H\ar[r]^{f_0}\ar[d]^{u_0} & B\ar[r]\ar[d]^v & 0\\
    0\ar[r] & L'\ar[r]_{j'} & E'\ar[r]_{f'} & B'\ar[r] & 0
  }
\]
을 가환으로 만드는 유일한 $A$-쌍가군 준동형 $u_0:H\to E'$이
존재하며, $u=u_0\circ j_1$, $w=j_0\circ w_0$이다. 여기서
$j_1$과 $j_2$는 표준 준동형이다. $u_0$가 환 준동형이기도 함은
곧바로 확인된다.

끝으로 자명한 확장에 관한 함자적 성질들을 기록하자.
\end{env}

\begin{proposition}[18.2.9]
\label{0.18.2.9-fr}
$B$, $B'$을 두 $A$-환, $L$을 $B$-쌍가군, $L'$을 $B'$-쌍가군,
$v:B\to B'$을 환의 $A$-준동형, $w:L\to L'$을 $(v,w)$가
쌍가군의 디-준동형을 이루는 $A$-쌍가군 준동형이라 하자. 그러면
도식들
\[
  \xymatrix@C=2.2pc@R=1.6pc{
    D_B(L)\ar[r]^u & D_{B'}(L') & & D_B(L)\ar[r]^u & D_{B'}(L')\\
    B\ar[u]\ar[r]_v & B'\ar[u] & & L\ar[u]\ar[r]_w & L'\ar[u]
  }
\]
을 가환으로 만드는 유일한 환의 $A$-준동형
$u:D_B(L)\to D_{B'}(L')$이 존재한다. 여기서 세로 화살표들은 표준
단사이다.
\end{proposition}

\oldpage[0\textsubscript{IV-1}]{59}

\begin{proof}
실제로 $u$는 사상 $(x,s)\mapsto(v(x),w(s))$일 수밖에 없으며, 이것이
환의 $A$-준동형임을 확인하는 일만 남는다. 이는 정의 (18.2.3)에서
곧바로 따른다. 또한 $u$는 도식
\[
  \xymatrix@C=2.2pc@R=1.5pc{
    D_B(L)\ar[r]^u\ar[d] & D_{B'}(L')\ar[d]\\
    B\ar[r]_v & B'
  }
\]
도 가환으로 만듦을 주목하자. 여기서는 세로 화살표들이 증대이다.
\end{proof}

\begin{proposition}[18.2.10]
\label{0.18.2.10-fr}
$B$를 $A$-환, $L$을 $B$-쌍가군, $E$를 $L$에 의한 $B$의
$A$-확장이라 하자. $B$에서 $E$로 가는 $B$-환 준동형들의 집합
$G$ (달리 말해, 증대 $E\to B$의 오른쪽 역인 $A$-준동형들의 집합)와
$D_B(L)$에서 $E$로 가는 $A$-동치사상들의 집합 $G'$ 사이의 전단사
사상을 다음과 같이 정의한다. 각 $g\in G$에는 $A$-동치사상
$g':(x,s)\mapsto g(x)+s$를 대응시키고, 역함수는 각 $g'\in G'$에
$B$-준동형 $x\mapsto g'(x,0)$을 대응시킨다.
\end{proposition}

\begin{proof}
이는 정의들에서 곧바로 따른다.
\end{proof}

\subsection{$A$-확장류의 군}
\label{subsection:0.18.3-fr}

\begin{env}[18.3.1]
\label{0.18.3.1-fr}
$A$-환 $B$와 $B$-쌍가군 $L$을 고정하자. 그러면 $L$에 의한 $B$의
두 $A$-확장 $E$, $E'$ 사이의 ``$E$와 $E'$은 $A$-동치이다''라는
관계는 동치관계이며, 이 관계에 관하여 $L$에 의한 $B$의
$A$-확장들의 $A$-동치류의 \emph{집합}을 말할 수 있다.
이를 보이려면 다음을 주목하면 충분하다. 각 $x\in B$에 대하여
$c_x$를 $B$에서의 상이 $x$인 $E$의 원소라 하면, 모든 $z\in E$는
$t\in L$인 꼴 $c_x+t$로 유일하게 쓰이며,
\[
  c_{x+y}=c_x+c_y+\varphi(x,y),\qquad
  c_{xy}=c_xc_y+\psi(x,y)
\]
로 쓸 수 있다. 여기서 $\varphi(x,y)$와 $\psi(x,y)$는 $L$의
원소들이고, $B\times B$에서 $L$로 가는 사상 $\varphi$와 $\psi$는
$E$가 $A$-환임을 나타내는 조건들을 만족해야 하지만, 그 조건들을
여기에 쓸 필요는 없다. 따라서 $L$에 의한 $B$의 모든 $A$-확장은
바탕집합이 $B\times L$인 어떤 $A$-확장과 $A$-동치이며, 이로부터
주장이 곧바로 따른다.
\end{env}

\begin{env}[18.3.2]
\label{0.18.3.2-fr}
$A$-환 $B$를 고정하고, $L$에 의한 $B$의 $A$-확장류의 집합을
잠정적으로 $T(L)$로 나타내자. 모든 $B$-쌍가군 준동형
$w:L\to L'$에 대하여, $L$에 의한 $B$의 $A$-확장 $E$의 류에
(18.2.8)에 따라 $w$를 통해 $E$에서 유도된 $A$-확장
$E\oplus_LL'$의 류를 대응시킴으로써 사상
$T(w):T(L)\to T(L')$을 표준적으로 정의한다. $w':L'\to L''$가
두 번째 $B$-쌍가군 준동형이면 또한
\begin{equation}
\label{0.18.3.2.1-fr}
\tag{18.3.2.1}
  T(w'\circ w)=T(w')\circ T(w)
\end{equation}
이다. 실제로 귀납극한의 일반 성질에 의해 ($\mathrm{I}$, 3.3.9 참조)
$A$-쌍가군의 표준 동형사상
\[
  E\oplus_LL''\xrightarrow{\sim}(E\oplus_LL')\oplus_{L'}L''
\]
이 존재함을 알고 있으며, 이것이 실제로 $A$-확장들의 $A$-동치사상임은
곧바로 확인된다. 따라서

\oldpage[0\textsubscript{IV-1}]{60}

관계 (18.3.2.1)을 얻는다. $\mathbf C_B$가 $B$-쌍가군들의 범주를
나타내면, 이로써 \emph{공변 함자}
$T:\mathbf C_B\to\mathbf{Ens}$가 정의되었음을 알 수 있다.
\end{env}

\begin{env}[18.3.3]
\label{0.18.3.3-fr}
이제 $B$-쌍가군들의 족 $(L_\alpha)_{\alpha\in I}$와 그 곱
$L=\prod_{\alpha\in I}L_\alpha$를 생각하자. 사영
$\operatorname{pr}_\alpha:L\to L_\alpha$는 사상
\[
  T(\operatorname{pr}_\alpha):T(L)\to T(L_\alpha)
\]
을 정의하고, 따라서 표준 사상
\begin{equation}
\label{0.18.3.3.1-fr}
\tag{18.3.3.1}
  \prod_\alpha T(\operatorname{pr}_\alpha):T(L)
  \longrightarrow\prod_{\alpha\in I}T(L_\alpha)
\end{equation}
을 정의한다. 이 사상이 \emph{전단사}임을 보이자. 실제로 각
$\alpha\in I$에 대하여 $E_\alpha$를 $L_\alpha$에 의한 $B$의
$A$-확장이라 하고, $F$를 $B$ 위에서 $E_\alpha$들의
\emph{섬유곱}이라 하자 (18.1.3). 그러면 $F$가
$L=\prod_{\alpha\in I}L_\alpha$에 의한 $B$의 $A$-확장임은
곧바로 알 수 있다. 또한 각 $E_\alpha$를 이와 $A$-동치인
$A$-확장 $E'_\alpha$로 바꾸면, $B$ 위에서 $E'_\alpha$들의
섬유곱 $F'$은 $F$와 $A$-동치이다. 이로써 사상
$\prod_{\alpha\in I}T(L_\alpha)\to T(L)$이 정의되며, 이 사상이
(18.3.3.1)의 역함수임은 명백하다. 따라서 주장이 따른다. 이제
(18.3.3.1)을 통해 $T(L)$을 $\prod_{\alpha\in I}T(L_\alpha)$와
표준적으로 동일시한다. 나아가 $(L'_\alpha)_{\alpha\in I}$가
$B$-쌍가군들의 또 다른 족이고 각 $\alpha$에 대하여
$w_\alpha:L_\alpha\to L'_\alpha$가 $B$-쌍가군 준동형이면,
$w=\prod_\alpha w_\alpha:L\to
L'=\prod_{\alpha\in I}L'_\alpha$로 둘 때 앞의 동일시 아래에서 $T(w)$은
$\prod_{\alpha\in I}T(w_\alpha)$와 동일시됨도 곧바로 확인된다.
\end{env}

\begin{env}[18.3.4]
\label{0.18.3.4-fr}
이제 $B$-쌍가군 $L$에 대하여 덧셈은 $B$-쌍가군 준동형
$s:L\times L\to L$이고, $L$의 덧셈법칙의 역원 사상 $t:L\to L$도
마찬가지이다. (18.3.3)에 따라 이로부터 $T(L)$ 위의 합성법칙
\[
  T(s):T(L)\times T(L)\longrightarrow T(L)
\]
을 얻는다. 가환도식들을 이용한 군 대상의 정의로부터 알 수 있듯이
($\mathbf 0_{\mathrm{III}}$, 8.2.5와 8.2.6), 이 합성법칙은
$T(L)$에 \emph{가환군} 구조를 주고 그 역원 사상은 $T(t)$이다.
이렇게 정의된
가환군을 $\operatorname{Exan}_A(B,L)$로 나타내고, 이를
\emph{$L$에 의한 $B$의 $A$-확장류의 군}이라 한다.
\end{env}

\begin{env}[18.3.5]
\label{0.18.3.5-fr}
대상이 다음과 같은 삼중항 $(A,B,L)$인 범주를 $\mathbf K$로
나타내자. 여기서 $A$는 환, $B$는 $A$-환, $L$은 $B$-쌍가군이다.
$(A,B,L)$에서 $(A',B',L')$로 가는 이 범주의 \emph{사상}은
삼중항 $(u,v,w)$이다. 여기서
$u:A'\to A$와 $v:B'\to B$는 도식
\[
  \xymatrix@C=2pc@R=1.5pc{
    A\ar[r] & B & L\ar[d]^w\\
    A'\ar[u]^u\ar[r] & B'\ar[u]^v & L'
  }
\]
의 왼쪽 사각형을 가환으로 만드는 두 환 준동형이고, 가로 화살표들은
구조 준동형이다. 끝으로 $w:L\to L'$은 모든 $z\in L$과
$b'\in B'$에 대하여 $w(v(b')z)=b'w(z)$ 및

\oldpage[0\textsubscript{IV-1}]{61}

$w(zv(b'))=w(z)b'$을 만족하는 가환군 준동형이다. 달리 말하면,
$L$에 $v$가 정의하는 $B'$-쌍가군 구조를 줄 때 $w$는
$B'$-쌍가군 준동형이다.

사상들의 합성은
$(u',v',w')\circ(u,v,w)
=(u\circ u',v\circ v',w'\circ w)$로 정의되며, 이는 곧바로
정당화된다. 우리는
\begin{equation}
\label{0.18.3.5.1-fr}
\tag{18.3.5.1}
  (A,B,L)\longmapsto\operatorname{Exan}_A(B,L)
\end{equation}
이 범주 $\mathbf K$에서 가환군들의 범주 $\mathbf{Ab}$로 가는
\emph{공변 함자}임을 보이고자 한다. 따라서 위와 같은 각 삼중항
$(u,v,w)$에 대하여 가환군 준동형
\[
  (u,v,w)_*:\operatorname{Exan}_A(B,L)
  \longrightarrow\operatorname{Exan}_{A'}(B',L')
\]
을 정의해야 한다. $\mathbf K$의 사상의 정의에 따라
\[
  (u,v,w)=(1_{A'},1_{B'},w)\circ(1_{A'},v,1_L)
  \circ(u,1_B,1_L)
\]
로 쓸 수 있다. 여기서 왼쪽 첫째 사상 $(1_{A'},1_{B'},w)$에서는
$L$에 $v$가 정의하는 $B'$-쌍가군 구조를 준다. 따라서 먼저
$u$, $v$, $w$ 가운데 둘이
항등사상인 경우에 $(u,v,w)_*$를 정의하겠다.
\end{env}

\begin{env}[18.3.6]
\label{0.18.3.6-fr}
먼저 $(1_A,1_B,w)_*$로는 (18.3.2)에서 $T(w)$로 나타낸 사상
\begin{equation}
\label{0.18.3.6.1-fr}
\tag{18.3.6.1}
  w_*:\operatorname{Exan}_A(B,L)\longrightarrow
  \operatorname{Exan}_A(B,L')
\end{equation}
을 취한다. 이것이 군 준동형임은 곧바로 확인된다. 이 성질은
$L$과 $L'$에 대한 같은 형태의 도식에 $T$를 적용하여 얻은 도식들의
가환성으로 표현되기 때문이다.

사상 $(1_A,v,1_L)_*$는 다음과 같이 정의되는 사상
\begin{equation}
\label{0.18.3.6.2-fr}
\tag{18.3.6.2}
  v^*:\operatorname{Exan}_A(B,L)\longrightarrow
  \operatorname{Exan}_A(B',L)
\end{equation}
이다. $E$가 $L$에 의한 $B$의 $A$-확장이면
$E\times_BB'$은 $L$에 의한 $B'$의 $A$-확장임을 (18.2.5)에서
보았다. 또한 $E$를 이와 $A$-동치인 $A$-확장 $E'$로 바꾸면
$E'\times_BB'$은 $E\times_BB'$과 $A$-동치이다. $v^*$는 $E$의
류를 $E\times_BB'$의 류로 보낸다. 이제 $w:L\to L'$이
$B$-쌍가군 준동형이면 도식
\begin{equation}
\label{0.18.3.6.3-fr}
\tag{18.3.6.3}
  \xymatrix@C=3pc@R=2.4pc{
    \operatorname{Exan}_A(B,L)\ar[r]^{v^*}\ar[d]_{w_*}
      & \operatorname{Exan}_A(B',L)\ar[d]^{w_*}\\
    \operatorname{Exan}_A(B,L')\ar[r]_{v^*}
      & \operatorname{Exan}_A(B',L')
  }
\end{equation}
이 가환임은 곧바로 확인된다. 여기서 오른쪽 열에서는 $L$과 $L'$을
$v$를 통하여 $B'$-쌍가군으로 본다. $L$과 $L'$을 각각
$L\times L$과 $L$로 바꾸고 $w$를 $L$의 덧셈 $s$로 바꾸면,
$v^*$가 실제로 \emph{군 준동형}임을 얻는다.
\end{env}

\oldpage[0\textsubscript{IV-1}]{62}

\begin{env}[18.3.6.4]
\label{0.18.3.6.4-fr}
끝으로 사상 $(u,1_B,1_L)_*$는 다음 사상이다.
\[
  u^*:\operatorname{Exan}_A(B,L)\longrightarrow
  \operatorname{Exan}_{A'}(B,L)
\]
이는 $L$에 의한 $B$의 $A$-확장 $E$에, $u$를 통하여
$A'$-\emph{환}으로 본 환 $E$를 대응시켜 얻는다 (18.1.1). 이 환은
명백히 $L$에 의한 $B$의 $A'$-확장이다. 여기서 $B$도 $u$를 통하여
$A'$-환으로 본다. $A$-동치사상은 또한 $A'$-동치사상이므로
사상 (18.3.6.4)을 얻는다. 모든 $B$-쌍가군 준동형
$w:L\to L'$에 대하여 이 사상은 다시 도식
\begin{equation}
\label{0.18.3.6.5-fr}
\tag{18.3.6.5}
  \xymatrix@C=3pc@R=2.4pc{
    \operatorname{Exan}_A(B,L)\ar[r]^{u^*}\ar[d]_{w_*}
      & \operatorname{Exan}_{A'}(B,L)\ar[d]^{w_*}\\
    \operatorname{Exan}_A(B,L')\ar[r]_{u^*}
      & \operatorname{Exan}_{A'}(B,L')
  }
\end{equation}
을 가환으로 만든다. 따라서 위와 마찬가지로 $u^*$가 군 준동형임을
얻는다.

이제
$(u,v,w)_*=(1_{A'},1_{B'},w)_*\circ(1_{A'},v,1_L)_*
\circ(u,1_B,1_L)_*$로 정의한다. 도식 (18.3.6.3)과
(18.3.6.5)의 가환성에 따라
\[
  (u\circ u',v\circ v',w'\circ w)_*
  =(u',v',w')_*\circ(u,v,w)_*
\]
임을 쉽게 확인한다. 이로써 (18.3.5.1)이 함자임을 모두 증명하였다.

군 준동형 (18.3.6.1)의 존재로부터 특히 다음이 따른다. $E$가
$L$에 의한 $B$의 \emph{자명한} $A$-확장이면, (18.2.8)에서 정의한
$L'$에 의한 $B$의 확장 $E\oplus_LL'$도 자명하다. 이는 직접
확인하기도 어렵지 않다.

한편 $B$의 \emph{중심} $Z$의 모든 원소 $z$에 대하여
$h_z:y\mapsto yz$라는 스칼라곱 자기준동형은 $B$-쌍가군 $L$의
자기준동형이다. 따라서 $(h_z)_*$는 가환군
$\operatorname{Exan}_A(B,L)$의 자기준동형이고, 함자성에 의해 이
자기준동형들은 $\operatorname{Exan}_A(B,L)$에 표준적인
$Z$-\emph{가군} 구조를 정의한다.
\end{env}

\begin{env}[18.3.7]
\label{0.18.3.7-fr}
$A$, $A'$을 두 환, $u:A'\to A$를 준동형, $B$를 $A$-환,
$L$을 $B$-쌍가군이라 하자. 군 준동형
\[
  u^*:\operatorname{Exan}_A(B,L)\longrightarrow
  \operatorname{Exan}_{A'}(B,L)
\]
의 \emph{핵}은 정의상 다음과 같은 $L$에 의한 $B$의
$A$-확장류들로 이루어진다. 즉 이 확장들을 $u$를 통해
$A'$-확장으로 볼 때 $A'$-\emph{자명}한 것들이다. 혼동의 여지가
없을 때 이 핵을 $\operatorname{Exan}_{A/A'}(B,L)$로 나타낸다.

$\Lambda$가 환이고 $\Lambda$-환들의 $\Lambda$-\emph{준동형}으로
이루어진 가환도식
\begin{equation}
\label{0.18.3.7.1-fr}
\tag{18.3.7.1}
  \xymatrix@C=2pc@R=1.5pc{
    B'\ar[r] & B\\
    A'\ar[u]\ar[r] & A\ar[u]
  }
\end{equation}

\oldpage[0\textsubscript{IV-1}]{63}

이 주어지면, 이로부터 표준적으로 군 준동형들
\begin{equation}
\label{0.18.3.7.2-fr}
\tag{18.3.7.2}
  \operatorname{Exan}_{A/\Lambda}(B,L)\longrightarrow
  \operatorname{Exan}_{A'/\Lambda}(B,L)\longrightarrow
  \operatorname{Exan}_{A'/\Lambda}(B',L)
\end{equation}
을 얻는다. 이들은 도식
\[
  \xymatrix@C=2.8pc@R=2.3pc{
    \operatorname{Exan}_A(B,L)\ar[r]\ar[d]
      & \operatorname{Exan}_{A'}(B,L)\ar[r]\ar[d]
      & \operatorname{Exan}_{A'}(B',L)\ar[d]\\
    \operatorname{Exan}_{\Lambda}(B,L)\ar[r]
      & \operatorname{Exan}_{\Lambda}(B,L)\ar[r]
      & \operatorname{Exan}_{\Lambda}(B',L)
  }
\]
의 가환성에서 나온다. 여기서 각 화살표는 (18.3.7.1)의 화살표들로부터
함자성에 따라 유도된다.
\end{env}

\begin{proposition}[18.3.8]
\label{0.18.3.8-fr}
$B$를 환, $\mathfrak J$를 $B$의 양쪽 아이디얼,
$C=B/\mathfrak J$를 몫환이라 하자. 그러면
$\mathfrak J/\mathfrak J^2$에는 표준적으로 $C$-쌍가군 구조가
주어진다. 모든 $C$-쌍가군 $L$에 대하여
$\operatorname{Hom}_C(\mathfrak J/\mathfrak J^2,L)$을
$\mathfrak J/\mathfrak J^2$에서 $L$로 가는 $C$-쌍가군 준동형들의
덧셈군이라 하자. 다음과 같이 가환군의 표준 동형사상
\begin{equation}
\label{0.18.3.8.1-fr}
\tag{18.3.8.1}
  \eta_L:\operatorname{Hom}_C(\mathfrak J/\mathfrak J^2,L)
  \xrightarrow{\sim}\operatorname{Exan}_B(C,L)
\end{equation}
을 정의한다. 각 $C$-쌍가군 준동형
$w:\mathfrak J/\mathfrak J^2\to L$은 따라서 특히
$B$-쌍가군 준동형이기도 하며, 이에 (18.2.8)에 따라
$\mathfrak J/\mathfrak J^2$에 의한 $C$의 확장
$B/\mathfrak J^2$에서 $w$를 통해 유도된 확장
$(B/\mathfrak J^2)\oplus_{(\mathfrak J/\mathfrak J^2)}L$의 류를
대응시킨다. 역동형사상은 $L$에 의한 $C$의 $B$-확장
$E$의 류에 다음 성질을 만족하는 준동형 $w$를 대응시킨다.
$j$를 통해 $L$을 $E$의 핵과 동일시할 때 합성
$\mathfrak J\to\mathfrak J/\mathfrak J^2\xrightarrow{w}L$은 구조
준동형 $B\to E$를 $\mathfrak J$에 제한한 사상과 같다.
\end{proposition}

\begin{proof}
$F=B/\mathfrak J^2$라 하자. 이는
$\mathfrak J/\mathfrak J^2$에 의한 $C$의 $B$-확장이다.
$L$에 의한 $C$의 임의의 $B$-확장
\[
  0\longrightarrow L\xrightarrow{j}E\xrightarrow{p}C\longrightarrow0
\]
에 대하여, 구조 준동형 $f:B\to E$는 합성
$B\xrightarrow{f}E\xrightarrow{p}C$가 표준 준동형
$B\to B/\mathfrak J=C$이 되게 한다.
$p\circ f$에 의한 $\mathfrak J$의 상은 영이므로
$f(\mathfrak J)$은 $p$의 핵, 즉 $j(L)$에 포함된다.
또한 $j(L)$은 제곱영이므로 $f(\mathfrak J^2)=0$이다.
따라서 $f$는 $B\to B/\mathfrak J^2=F\xrightarrow{u}E$를 거쳐
인수분해된다. $w$를 $u$의
$\mathfrak J/\mathfrak J^2$에 대한 제한이라 하면 가환도식
\[
  \xymatrix@C=2pc@R=1.5pc{
    0\ar[r] & \mathfrak J/\mathfrak J^2\ar[r]\ar[d]^w
      & F\ar[r]^g\ar[d]^u
      & C\ar[r]\ar[d]^{1_C}\ar@{}[d]|{\wr} & 0\\
    0\ar[r] & L\ar[r]_j & E\ar[r]_p & C\ar[r] & 0
  }
\]
을 얻는다. 달리 말하면 $(u,1_C,w)$는 확장들의 사상이다 (18.2.4).
따라서 확장들의 사상
\[
  \xymatrix@C=2pc@R=1.5pc{
    0\ar[r] & L\ar[r]\ar[d]^{1_L}\ar@{}[d]|{\wr}
      & E'\ar[r]\ar[d]^{u'}
      & C\ar[r]\ar[d]^{1_C}\ar@{}[d]|{\wr} & 0\\
    0\ar[r] & L\ar[r]_j & E\ar[r]_p & C\ar[r] & 0
  }
\]
을 얻는다.

\oldpage[0\textsubscript{IV-1}]{64}

여기서 $E'=F\oplus_{(\mathfrak J/\mathfrak J^2)}L$이다 (18.2.8).
$\eta_L$이 전단사임을 보이려면 $u'$이 $B$-\emph{동치사상}임을
보이면 된다. (18.2.8)의 구성과 (18.1.7)에 따라, 이는 표준 준동형
$g:F\to C$에 의한 $C$의 확장 $E$의 역상 확장 $G$가
$F$-자명함을 보이는 것으로 귀착된다. 그런데 $g$는
$F\xrightarrow{u}E\xrightarrow{p}C$를 거쳐 인수분해되므로 이는
명백하다 (18.1.6).

$\eta_L$이 \emph{군} 준동형임을 보이는 일만 남았다. 모든
$B$-쌍가군 준동형 $h:L\to L'$에 대하여 도식
\[
  \xymatrix@C=3pc@R=2.4pc{
    \operatorname{Hom}_C(\mathfrak J/\mathfrak J^2,L)
      \ar[r]^{\eta_L}\ar[d]_{\operatorname{Hom}(1,h)}
      & \operatorname{Exan}_B(C,L)\ar[d]^{h_*}\\
    \operatorname{Hom}_C(\mathfrak J/\mathfrak J^2,L')
      \ar[r]_{\eta_{L'}}^{\sim}
      & \operatorname{Exan}_B(C,L')
  }
\]
이 가환임은 곧바로 확인된다. 이제 덧셈을 정의하는 준동형
$L\times L\to L$에 이 사실을 적용하면 결론을 얻는다.
\end{proof}

\subsection{대수의 확장}
\label{subsection:0.18.4-fr}

\begin{env}[18.4.1]
\label{0.18.4.1-fr}
$A$를 \emph{가환}환이라 하자. $A$-\emph{대수}들의 범주는
$A$-환들의 범주의 충만한 부분범주이며, 구조 준동형
$\rho:A\to B$의 상이 $B$의 \emph{중심}에 포함된다는 조건으로
특징지어진다.

$B$가 $A$-대수이고 $L$이 $B$를 통해 유도되는 두 $A$-가군 구조가
일치하는 $B$-쌍가군이면, 표준 자명 $A$-확장 $D_B(L)$이
$A$-대수임은 명백하다. 마찬가지로 (18.2.5)의 구성에서
$B$, $B'$, $E'$이 $A$-대수이면 $E'\times_{B'}B$도 $A$-대수이고,
(18.2.8)의 구성에서 $B$와 $E$가 $A$-대수이면
$E\oplus_LL'$도 $A$-대수이다. 끝으로 $B$가 $A$-대수이고,
$B$-쌍가군 $L$에 의한 $B$의 $A$-확장 $E$가 $A$-대수이면,
$E$와 $A$-\emph{동치}인 $L$에 의한 $B$의 모든 $A$-확장도
$A$-대수이다. 따라서 정의 (18.3.4)로부터, $A$-\emph{대수}인
$L$에 의한 $B$의 $A$-확장들의 $A$-동치류들이
$\operatorname{Exan}_A(B,L)$의 \emph{부분군}을 이룸을 곧바로
알 수 있다. 이 부분군을 $\operatorname{Exal}_A(B,L)$로 나타낸다.

(18.3.5)에서 정의한 범주 $\mathbf K$의 충만한 부분범주로서,
대상 $(A,B,L)$이 다음 조건을 만족하는 범주를 $\mathbf K'$이라 하자.
즉 $A$는 가환환이고, $B$는 $A$-대수이며, $B$를 통해 $L$에
유도되는 두 $A$-가군 구조는 일치한다. 그러면 앞의 논의에 따라
\[
  (A,B,L)\longmapsto\operatorname{Exal}_A(B,L)
\]
은 $\mathbf K'$에서 $\mathbf{Ab}$로 가는 \emph{공변 함자}이다.
(18.3.7)의 결과들은 모든 곳에서 $\operatorname{Exan}$을
$\operatorname{Exal}$로 바꾸어도 그대로 성립한다.
\end{env}

\begin{env}[18.4.2]
\label{0.18.4.2-fr}
계속해서 환 $A$가 가환이라고 가정하자. (18.4.1)의 논의에서
``대수''를 ``가환 대수''로, ``쌍가군''을 ``가군''으로 바꾸어도
그대로 성립한다. $B$가 가환 $A$-대수이고 $L$이 $B$-가군이면,

\oldpage[0\textsubscript{IV-1}]{65}

$A$-\emph{가환 대수}인 $L$에 의한 $B$의 $A$-확장들의
$A$-동치류들, 또는 이와 같은 말로 $A$-\emph{가환환}인 확장들의
$A$-동치류들은 $\operatorname{Exal}_A(B,L)$의 \emph{부분군}을
이룬다. 이를 $\operatorname{Exalcom}_A(B,L)$로 나타낸다.
$\mathbf K''$를 $\mathbf K'$의 충만한 부분범주로서, $A$와 $B$가
가환이고 $L$이 $B$-가군인 삼중항 $(A,B,L)$들로 이루어진 범주라
하면,
\[
  (A,B,L)\longmapsto\operatorname{Exalcom}_A(B,L)
\]
은 역시 $\mathbf K''$에서 $\mathbf{Ab}$로 가는 \emph{공변 함자}이다.
또한 (18.3.7)의 모든 곳에서 $\operatorname{Exan}$을
$\operatorname{Exalcom}$으로 바꿀 수 있다. 끝으로 (18.3.8)에서
$B$가 \emph{가환}환이고 $L$이 $C$-\emph{가군}이라고 가정하면,
같은 논증으로 표준 동형사상
\begin{equation}
\label{0.18.4.2.1-fr}
\tag{18.4.2.1}
  \operatorname{Hom}_C(\mathfrak J/\mathfrak J^2,L)
  \xrightarrow{\sim}\operatorname{Exalcom}_B(C,L)
\end{equation}
을 얻는다. 여기서 왼쪽 항은 $C$-\emph{가군} 준동형들의 군이다.
\end{env}

\begin{env}[18.4.3]
\label{0.18.4.3-fr}
$A$를 가환환이라 하자. $A$-대수의 확장에서 중요한 한 경우는
다음과 같은 $A$-대수 $E$들로 이루어진다. 즉 $E$는
$B$-쌍가군 $L$에 의한 $A$-대수 $B$의 확장이고, $A$-\emph{가군}으로
볼 때 $A$-가군 $L$에 의한 $A$-가군 $B$의 \emph{자명한} 확장이다.
달리 말해 $A$-\emph{가군}의 완전열
$0\to L\to E\to B\to0$은 \emph{분할된다}. 이때 $E$를
$L$에 의한 $B$의 \emph{호흐실트 확장}이라 한다. $B$가
\emph{사영} $A$-가군이면 언제나 이 경우에 해당하며, 특히 $A$가
가환 \emph{체}이면 그러하다. $A$-가군으로서 $E$를 $B\times L$과
동일시할 수 있고, $E$에서의 곱셈은
$(x,s)(y,t)=(xy,xt+sy+f(x,y))$로 주어진다. 여기서 $f(x,y)\in L$이다.
이 곱셈이 $B\times L$ 위에 $A$-대수 구조를 정의한다는 조건을 쓰면
(M, XIV, 2), $f$는 $B\times B$에서 $L$로 가는
$A$-\emph{쌍선형} 사상으로서
\begin{equation}
\label{0.18.4.3.1-fr}
\tag{18.4.3.1}
  f(xy,z)+f(x,y)z=xf(y,z)+f(x,yz)
\end{equation}
를 만족해야 함을 얻는다. 달리 말해 $f$는 호흐실트의 의미에서
$B$ 위에서 $L$에 값을 갖는 $2$-\emph{공순환}이다.
확장 $E$가 $A$-자명하기 위한 필요충분조건은
\begin{equation}
\label{0.18.4.3.2-fr}
\tag{18.4.3.2}
  f(x,y)=xg(y)-g(xy)+g(x)y
\end{equation}
가 성립하는 것이다. 여기서 $g$는 $B$에서 $L$로 가는
$A$-\emph{선형} 사상이다. 달리 말해 $f$는 호흐실트의 의미에서
$2$-\emph{공경계}여야 한다. 따라서 $L$에 의한 $B$의 호흐실트
확장류들은 $\operatorname{Exal}_A(B,L)$의 부분군을 이루며,
이 부분군은 호흐실트 \emph{코호몰로지군}
$\mathrm H_A^2(B,L)$과 동형이다.

$B$가 가환 $A$-대수이고 $L$이 $B$-가군이면, $L$에 의한 $B$의
\emph{가환 호흐실트 확장}은 $f(y,x)=f(x,y)$를 만족하는
\emph{대칭} $2$-공순환에 대응한다. 따라서 $L$에 의한 $B$의 가환
호흐실트 확장류들은 $\operatorname{Exalcom}_A(B,L)$의 부분군을
이룬다. 이 부분군은 대칭 $2$-공순환들의 군이
$\mathrm H_A^2(B,L)$ 안에 갖는 상과 동형이며, 그 상을
$\mathrm H_A^2(B,L)^s$로 나타낸다.
\end{env}

\begin{env}[18.4.4]
\label{0.18.4.4-fr}
$A$와 $B$가 \emph{가환}이고 $L$이 $B$-가군인 경우로 한정하여,
이 경우 호흐실트 코호몰로지군 $\mathrm H_A^i(B,L)$의 동치인 정의를
상기하자 (M, IX, 6). $B$-가군의 복합체
$\mathbf P_\bullet=(\mathbf P_n)_{n\geqslant0}$를 생각하자. 여기서
$\mathbf P_n=B^{\otimes(n+1)}=B\otimes_A B\otimes_A\cdots\otimes_A B$
이고 인자는 $n+1$개이며, $B$-가군 구조는

\oldpage[0\textsubscript{IV-1}]{66}

다음과 같이 정의된다.
\[
  x(y_1\otimes y_2\otimes\cdots\otimes y_{n+1})
  =(xy_1)\otimes y_2\otimes\cdots\otimes y_{n+1}\,;
\]
차수 $-1$의 미분 $d_n:\mathbf P_n\to\mathbf P_{n-1}$은
\begin{align*}
  d_n(x_1\otimes x_2\otimes\cdots\otimes x_{n+1})
    &=(x_1x_2)\otimes x_3\otimes\cdots\otimes x_{n+1}
      -x_1\otimes(x_2x_3)\otimes\cdots\otimes x_{n+1}+\cdots\\
    &\quad+(-1)^{n-1}x_1\otimes x_2\otimes\cdots\otimes x_{n-1}
      \otimes(x_nx_{n+1})
      +(-1)^n(x_{n+1}x_1)\otimes x_2\otimes\cdots\otimes x_n,
\end{align*}
으로 정의된다. $B$가 가환이므로 이는 실제로 $B$-선형이다.

이 복합체의 $L$에 값을 갖는 $2$-공순환은
$h(d_3(x\otimes y\otimes z\otimes t))=0$을 만족하는
$B\otimes_A B\otimes_A B$에서 $L$로 가는 $B$-선형 사상 $h$이다.
그런데 $h(x\otimes y\otimes z)=xh(1\otimes y\otimes z)$이므로,
공순환 $h$는 $B\times B$에서 $L$로 가는 $A$-쌍선형 사상
$(y,z)\longmapsto f(y,z)=h(1\otimes y\otimes z)$에 의해 결정된다.
$h$에 대한 앞의 조건을 쓰면 $f$에 대하여 조건 (18.4.3.1)을 다시
얻는다. 마찬가지로 $2$-공경계는
$x\otimes y\otimes z\longmapsto h'(d_2(x\otimes y\otimes z))$ 꼴의
사상이다. 여기서 $h':B\otimes_A B\to L$은 $B$-선형이며, 다시
$h'$은 $B$에서 $L$로 가는 $A$-선형 사상
$g:y\longmapsto h'(1\otimes y)$에 의해 결정된다. 그러면
$h'(d_2(1\otimes x\otimes y))=xg(y)-g(xy)+yg(x)$를 얻으며,
이는 (18.4.3.2)를 다시 준다. 모든 $i$에 대해서도 같은 방식으로
진행하면
\begin{equation}
\label{0.18.4.4.1-fr}
\tag{18.4.4.1}
  \mathrm H_A^i(B,L)=\mathrm H^i(\operatorname{Hom}_B(\mathbf P_\bullet,L)).
\end{equation}
을 얻는다.
\end{env}

\begin{env}[18.4.5]
\label{0.18.4.5-fr}
(18.4.4)의 조건 아래에서 군 $\mathrm H_A^2(B,L)^s$도 같은 방식으로
해석할 수 있다 (18.4.3). 이를 위해 차수 $3$에서 복합체
$\mathbf P_\bullet$을 바꾸어 새로운 복합체
\[
  \mathbf P'_\bullet:\mathbf P'_3\xrightarrow{d'_3}\mathbf P_2
  \xrightarrow{d_2}\mathbf P_1\,;
\]
를 생각하자. $\mathbf P'_3=\mathbf P_3\oplus
(B\otimes_A B\otimes_A B)$로 놓고, $d'_3$은 $\mathbf P_3$ 위에서는
$d_3$과 같으며 $B\otimes_A B\otimes_A B$ 위에서는
\[
  d'_3(x\otimes y\otimes z)=x\otimes y\otimes z-x\otimes z\otimes y
\]
로 주어진다.

관계 $d_2d'_3=0$은 $B$의 가환성에서 따른다. 위에서 도입한 표기
아래에서, 이제 $\mathbf P'_\bullet$의 $2$-공순환은
\emph{대칭}이고 (18.4.3.1)을 만족하는 $A$-쌍선형 사상
$f:B\times B\to L$에 대응한다. 따라서
\[
  \mathrm H_A^2(B,L)^s
  =\mathrm H^2(\operatorname{Hom}_B(\mathbf P'_\bullet,L)).
\]
\end{env}

\begin{env}[18.4.6]
\label{0.18.4.6-fr}
특히 가환 \emph{체} $k$와 $k$의 \emph{확대체} $K$를 생각하자.
임의의 $K$-벡터공간 $L$, $M$에 대하여
$\operatorname{Ext}_K^1(M,L)=0$이다. 따라서 (M, VI, 3.3 $a$)에 의해
표준 동형사상
\begin{equation}
\label{0.18.4.6.1-fr}
\tag{18.4.6.1}
  \mathrm H_k^i(K,L)\xrightarrow{\sim}
  \operatorname{Hom}_K(\mathrm H_i(\mathbf P_\bullet),L)
\end{equation}
을 얻고, 마찬가지로
\begin{equation}
\label{0.18.4.6.2-fr}
\tag{18.4.6.2}
  \mathrm H_k^2(K,L)^s\xrightarrow{\sim}
\operatorname{Hom}_K(\mathrm H_2(\mathbf P'_\bullet),L).
\end{equation}
을 얻는다.
\end{env}

\subsection{위상환의 경우}
\label{subsection:0.18.5-fr}

\begin{env}[18.5.1]
\label{0.18.5.1-fr}
$A$, $B$를 두 선형 위상환, $\rho:A\to B$를 연속 준동형,
$L$을 위상 $B$-쌍가군이라 하자. $B$의 \emph{열린} 양쪽 아이디얼
$\mathfrak R_0$가 존재하여
$\mathfrak R_0.L=L.\mathfrak R_0=0$이라고 가정하자. 그러면 $L$은

\oldpage[0\textsubscript{IV-1}]{67}

모든 열린 아이디얼 $\mathfrak R\subset\mathfrak R_0$에 대하여
$(B/\mathfrak R)$-쌍가군으로 볼 수 있다. 이제
$\mathfrak R\subset\mathfrak R_0$를 $B$의 열린 양쪽 아이디얼,
$\mathfrak J$를 $\rho(\mathfrak J)\subset\mathfrak R$을 만족하는
$A$의 열린 양쪽 아이디얼이라 하자. 그러면 $B/\mathfrak R$을
이산 $(A/\mathfrak J)$-환으로 볼 수 있다. 앞의 논의로부터 군
$\operatorname{Exan}_{A/\mathfrak J}(B/\mathfrak R,L)$이 정의된다.
더 나아가 $\mathfrak R'\subset\mathfrak R$과
$\mathfrak J'\subset\mathfrak J$이 각각 $B$와 $A$의 열린 양쪽
아이디얼이고 $\rho(\mathfrak J')\subset\mathfrak R'$이면,
(18.3.5.1)에 의해 표준 준동형
\begin{equation}
\label{0.18.5.1.1-fr}
\tag{18.5.1.1}
  \operatorname{Exan}_{A/\mathfrak J}(B/\mathfrak R,L)
  \longrightarrow
  \operatorname{Exan}_{A/\mathfrak J'}(B/\mathfrak R',L)
\end{equation}
을 얻는다.

$\rho(\mathfrak J)\subset\mathfrak R\subset\mathfrak R_0$을 만족하는
열린 아이디얼의 쌍 $(\mathfrak J,\mathfrak R)$들의 집합은
``$\mathfrak J\supset\mathfrak J'$이고
$\mathfrak R\supset\mathfrak R'$''이라는 관계에 관하여 자명하게
오른쪽 여과 순서집합을 이룬다. 사상 (18.5.1.1)들은 이 집합을
지표집합으로 하는 가법군의 \emph{귀납계}를 정의한다. 이제 표기를 남용하여
(확장들의 집합과 자연스럽게 전단사 대응하는 확장류의 군이 더 이상
아니므로)
\begin{equation}
\label{0.18.5.1.2-fr}
\tag{18.5.1.2}
  \operatorname{Exantop}_A(B,L)
  =\varinjlim\operatorname{Exan}_{A/\mathfrak J}(B/\mathfrak R,L)
\end{equation}
로 놓는다.

(18.5.1.2)의 우변이 \emph{영}이라는 것은 다음을 뜻한다. 즉
$\rho(\mathfrak J)\subset\mathfrak R\subset\mathfrak R_0$을 만족하는
열린 아이디얼의 모든 쌍 $\mathfrak J\subset A$,
$\mathfrak R\subset B$와 $L$에 의한 $B/\mathfrak R$의 모든
$(A/\mathfrak J)$-확장 $E$에 대하여, 열린 아이디얼
$\mathfrak J'\subset\mathfrak J$과
$\mathfrak R'\subset\mathfrak R$이 존재하여
$\rho(\mathfrak J')\subset\mathfrak R'$이고, 준동형
$B/\mathfrak R'\to B/\mathfrak R$에 의한 $E$의 역상 확장이
\emph{자명}하다.

$A$가 가환이고 $B$가 위상 $A$-대수(각각 위상 가환 $A$-대수)인
경우에는 $\operatorname{Exal}_{A/\mathfrak J}(B/\mathfrak R,L)$
(각각 $\operatorname{Exalcom}_{A/\mathfrak J}(B/\mathfrak R,L)$)로부터
귀납극한 $\operatorname{Exaltop}_A(B,L)$
(각각 $\operatorname{Exalcotop}_A(B,L)$)을 유사하게 정의한다.
그 정의는 독자에게 맡긴다.
\end{env}

\begin{env}[18.5.2]
\label{0.18.5.2-fr}
환들의 연속 준동형으로 이루어진 가환도식
\[
  \xymatrix@C=3pc@R=2pc{
    B'\ar[r] & B\\
    A'\ar[u]\ar[r] & A\ar[u]
  }
\]
이 주어지면, (18.3.5.1)에서 귀납극한으로 넘어가 표준적으로 가법군의
두 준동형
\[
  \operatorname{Exantop}_A(B,L)
  \longrightarrow\operatorname{Exantop}_{A'}(B,L)
  \longrightarrow\operatorname{Exantop}_{A'}(B',L)
\]
을 얻는다.

가환군의 범주에서 함자 $\varinjlim$이 완전하므로, 준동형
\[
  \operatorname{Exantop}_A(B,L)
  \longrightarrow\operatorname{Exantop}_{A'}(B,L)
\]
의 핵은 준동형들
\[
  \operatorname{Exan}_{A/\mathfrak J}(B/\mathfrak R,L)
  \longrightarrow
  \operatorname{Exan}_{A'/\mathfrak J'}(B/\mathfrak R,L)
\]
의 핵들의 귀납극한이다. 여기서 $\mathfrak J'$은
$\mathfrak J$의 역상으로 잡았다. 이 핵을
$\operatorname{Exantop}_{A/A'}(B,L)$로 나타낸다. 마찬가지로
$\operatorname{Exaltop}_{A/A'}(B,L)$과
$\operatorname{Exalcotop}_{A/A'}(B,L)$을 정의한다. 끝으로 연속
$B$-쌍가군 준동형 $L\to L'$이 주어지면,

\oldpage[0\textsubscript{IV-1}]{68}

(18.3.6.1)에서 귀납극한으로 넘어가 표준적으로 가법군 준동형
\[
  \operatorname{Exantop}_A(B,L)\longrightarrow
  \operatorname{Exantop}_A(B,L')
\]
을 얻는다.
\end{env}

\begin{env}[18.5.3]
\label{0.18.5.3-fr}
위상환 $C$와 두 위상 $C$-쌍가군 $M$, $N$이 주어졌을 때,
$M$에서 $N$으로 가는 \emph{연속} $C$-쌍가군 준동형들의 가법군을
$\operatorname{Hom.\ cont}_C(M,N)$으로 나타낸다.
\end{env}

\begin{lemma}[18.5.3.1]
\label{0.18.5.3.1-fr}
$C$를 위상환, $E$, $L$을 두 위상 $C$-쌍가군이라 하자. 이 위상들이
선형이고, $L$은 이산이며 $C$의 어떤 열린 양쪽 아이디얼에 의해
소멸된다고 가정하자. 그러면 표준 동형사상
\begin{equation}
\label{0.18.5.3.2-fr}
\tag{18.5.3.2}
  \varinjlim\operatorname{Hom}_{C/\mathfrak R}(E/V,L)
  \xrightarrow{\sim}\operatorname{Hom.\ cont}_C(E,L)
\end{equation}
이 성립한다. 여기서 왼쪽 항의 귀납극한은 다음 조건을 만족하는
쌍 $(\mathfrak R,V)$들의 여과 순서집합을 따라 취한다.
$\mathfrak R$은 $C$의 열린 양쪽 아이디얼이고, $V$는 $E$의 열린
$C$-부분쌍가군이며,
$\mathfrak R.L=L.\mathfrak R=0$, $\mathfrak R.E\subset V$,
$E.\mathfrak R\subset V$이다.
\end{lemma}

\begin{proof}
$C/\mathfrak R$과 $E/V$가 이산이므로 표준 준동형
\[
  w_{\mathfrak R,V}:\operatorname{Hom}_{C/\mathfrak R}(E/V,L)
  \longrightarrow\operatorname{Hom.\ cont}_C(E,L)
\]
들이 존재하며, 이들은 귀납계를 이룬다. 따라서 귀납극한으로 넘어가면
(18.5.3.2)의 준동형을 얻는다. $V'\subset V$일 때 준동형
$E/V'\to E/V$가 전사이므로, 정의로부터 준동형
\[
  \operatorname{Hom}_{C/\mathfrak R}(E/V,L)
  \longrightarrow\operatorname{Hom}_{C/\mathfrak R}(E/V',L)
\]
은 단사임이 곧 따른다. 여기서는
$\mathfrak R.L=L.\mathfrak R=0$, $\mathfrak R.E\subset V'$,
$E.\mathfrak R\subset V'$이다. 마찬가지로
$\mathfrak R'\subset\mathfrak R$일 때 준동형
\[
  \operatorname{Hom}_{C/\mathfrak R}(E/V,L)
  \longrightarrow\operatorname{Hom}_{C/\mathfrak R'}(E/V,L)
\]
도 단사임이 명백하다. 따라서 준동형 (18.5.3.2)는 단사이다.
한편 $u$가 $E$에서 $L$로 가는 연속 $C$-쌍가군 준동형이면, 그 핵은
$E$의 열린 부분쌍가군 $V_0$이다. $C$의 열린 양쪽 아이디얼
$\mathfrak R_0$을
$\mathfrak R_0.L=L.\mathfrak R_0=0$,
$\mathfrak R_0.E\subset V_0$, $E.\mathfrak R_0\subset V_0$을
만족하도록 잡자. 그러면 $u$는 $E/V_0$에서 $L$로 가는 어떤
$(C/\mathfrak R_0)$-준동형의 표준 상임이 명백하다. 따라서
(18.5.3.2)는 전사이다.
\end{proof}

이제 명제 (18.3.8)은 다음과 같이 위상환에 일반화된다.

\begin{proposition}[18.5.4]
\label{0.18.5.4-fr}
$B$를 선형 위상환, $\mathfrak J$을 $B$의 양쪽 아이디얼,
$C=B/\mathfrak J$을 몫 위상환이라 하자. $\mathfrak J$에는 $B$에서
유도되는 위상을 주고, $\mathfrak J/\mathfrak J^2$에는
$\mathfrak J$의 위상에서 오는 몫위상을 준다. 그러면
$\mathfrak J/\mathfrak J^2$에는 표준적으로 위상 $C$-쌍가군 구조가
주어진다. $C$의 어떤 열린 아이디얼에 의해 소멸되는 모든 이산
$C$-쌍가군 $L$에 대하여 표준 동형사상
\begin{equation}
\label{0.18.5.4.1-fr}
\tag{18.5.4.1}
  \operatorname{Hom.\ cont}_C(\mathfrak J/\mathfrak J^2,L)
  \xrightarrow{\sim}\operatorname{Exantop}_B(C,L).
\end{equation}
이 존재한다.
\end{proposition}

\begin{proof}
실제로 $(\mathfrak J+\mathfrak R)/\mathfrak J$이 $L$을 소멸시키는
$B$의 모든 열린 아이디얼 $\mathfrak R$에 대하여, (18.3.8)에 의해
표준 동형사상
\[
  \operatorname{Hom}_{B/(\mathfrak J+\mathfrak R)}
    ((\mathfrak J+\mathfrak R)/(\mathfrak J^2+\mathfrak R),L)
  \xrightarrow{\sim}
  \operatorname{Exan}_{B/\mathfrak R}(B/(\mathfrak J+\mathfrak R),L)
\]
을 얻는다. 이제 (18.5.3.1)을 고려하여 귀납극한으로 넘어가면 충분하다.
\end{proof}

\section{형식적으로 매끄러운 대수와 Cohen 환}
\label{section:0.19-fr}
\label{0.19-fr}

\setcounter{subsection}{-1}
\subsection{서론}
\label{subsection:0.19.0-fr}
\label{0.19.0-fr}

\oldpage[0\textsubscript{IV-1}]{69}

\begin{env}[19.0.1]
\label{0.19.0.1-fr}
제IV장에서는 특히 준스킴의 중요한 사상 부류인 \emph{매끄러운 사상}을
도입하고 연구할 것이다.\footnote{최근의 몇몇 연구에서는 고전적인
``단순점''이라는 용어에서 착안하여 이를 \emph{단순} 사상이라고도
한다. 그러나 이 용어는 특히 대수군론에서 혼동을 일으킨다.}
그 기본 성질 가운데 하나이며, 유한성 조건과 결합하면 이 사상들을
특징짓는 성질은 \emph{사상의 올림 성질}이다. 즉 $f:X\to Y$가
매끄러운 사상이고 $g:Y'\to Y$가 사상이라 하자. 또
$h:Y''\to Y'$가 $Y''$을 $Y'$과 ``조금밖에 다르지 않은'' 준스킴으로
만드는 사상이면, \emph{모든 $Y$-사상 $Y''\to X$는}
$Y''\xrightarrow{h}Y'\to X$로 \emph{인수분해된다}. 이를 정확히
설명하기 위해 $Y=\operatorname{Spec}(A)$와
$X=\operatorname{Spec}(B)$가 \emph{아핀}인 경우로 한정하자.
모든 $A$-대수 $C$와 $C$의 모든 \emph{멱영} 아이디얼
$\mathfrak J$에 대하여, 모든 $A$-준동형 $B\to C/\mathfrak J$을
$B\to C\to C/\mathfrak J$로 \emph{올릴 수 있으면} $B$를
\emph{형식적으로 매끄러운} $A$-대수라고 한다. 달리 말하면 사상
\[
  \operatorname{Hom}_A(B,C)\longrightarrow
  \operatorname{Hom}_A(B,C/\mathfrak J)
\]
이 \emph{전사}이다. 많은 응용에서 $X$는 $Y$-준스킴의 범주에서
집합의 범주로 가는 \emph{표현 가능한 반변 함자}
$Y'\mapsto F(Y')$의 표현 대상으로 나타나므로
($\mathbf{0}_{\mathrm{III}}$, 8.1.8)
$F(Y')=\operatorname{Hom}_Y(Y',X)$이다. 아핀인 경우
$F^0(C)=F(\operatorname{Spec}(C))$로 놓자. 위의 조건 아래에서
$F^0(C)\to F^0(C/\mathfrak J)$가 \emph{전사}임을 확인하면
(이 확인은 $F$가 표현 가능함을 \emph{알지 못하더라도} 할 수 있다),
추가적인 유한성 조건이 충족되는 한 $f$가 매끄러움이 따른다.
\end{env}

\begin{env}[19.0.2]
\label{0.19.0.2-fr}
위상환 $A$ 위의 \emph{위상 대수}에 대해서도 이와 유사한
\emph{형식적으로 매끄러운} 대수의 개념이 있으나, 여기서는 이를
정확히 정의하지 않겠다(19.3.1 참조). 이 개념들은 먼저 §~19에서
초등적인 관점으로 연구한다. 이어 §§~20과 21에서 전개할 미분의
성질들을 이용하여 §~22에서 더 정교한 정리들을 증명할 것이다.
형식적으로 매끄러운 대수에 관한 주요 결과들을 여기서 요약하자.

\medskip\noindent\textbf{I.} --- $A$, $B$를 두 뇌터 국소환,
$\varphi:A\to B$를 국소 준동형, $k$를 $A$의 잉여체라 하고,
$B_0=B\otimes_A k$라 하자. $A$, $B$, $B_0$에는 각각 준아딕 위상을,
$k$에는 이산 위상을 준다. 그러면 $B$가 형식적으로 매끄러운
$A$-대수일 필요충분조건은 $B$가 \emph{평탄} $A$-가군이고
$B_0$가 \emph{형식적으로 매끄러운} $k$-대수인 것이다(19.7.1).
따라서 이 정리는 뇌터 국소환의 형식적 매끄러움에 관한 문제를
\emph{체} 위의 대수인 뇌터 국소환에 관한 같은 문제로 환원한다.

\medskip\noindent\textbf{II.} --- $k$를 체, $A$를 $k$-대수인 뇌터
국소환이라 하자. $A$가 형식적으로 매끄러울 필요충분조건은
$A$가 $k$ 위에서 \emph{기하학적으로 정칙}인 것이다. 이는 $k$의
모든 \emph{유한} 확대체 $k'$에 대하여 반국소환
$A\otimes_k k'$이

\oldpage[0\textsubscript{IV-1}]{70}

\noindent 정칙이라는 뜻이다(22.5.8). 이때 $A$의 \emph{완비화}
$\widehat A$는 형식적 멱급수환 $K[[T_1,\ldots,T_n]]$과 동형인
환이다(19.6.5). 더욱이 $A$가 \emph{완비}이고 $k$ 위에서 형식적으로
매끄럽다고 가정하면, $A$의 $k$-대수 구조는 $A$의 잉여체 $K$와
$A$의 \emph{차원}에 의해 완전히 결정된다. 이 차원은 부등식
\[
  \dim(A)\geq \operatorname{rg}_k(\Gamma_{K/k})
\]
을 만족하는 한 임의로 주어질 수 있다. 여기서 $\Gamma_{K/k}$는
$K$의 ``불완전성 가군''이다(22.2.6).

특히 $k$의 \emph{확대체} $K$가 형식적으로 매끄러울 필요충분조건은
$K$가 $k$의 \emph{분리가능} 확대체인 것이다(19.6.1).

\medskip\noindent\textbf{III.} --- $A$를 뇌터 국소환,
$\mathfrak J\ne A$를 $A$의 아이디얼, $A_0=A/\mathfrak J$이라 하자.
$B_0$를 \emph{완비} 뇌터 국소환이라 하고, $A_0\to B_0$를
$B_0$가 형식적으로 매끄러운 $A_0$-대수가 되게 하는 국소 준동형이라
하자. 그러면 \emph{완비} 뇌터 국소환 $B$와 $B$를 평탄
$A$-가군으로 만드는 국소 준동형 $A\to B$, 그리고 $A_0$-\emph{동형사상}
\[
  B\otimes_A A_0\xrightarrow{\sim}B_0
\]
가 존재한다. 따라서 I에 의해 $B$는 형식적으로 매끄러운
$A$-대수이다. 더욱이 이 성질들은 $B$를 동형을 제외하면 유일하게
결정한다(19.7.2). 특히 이 정리는 \emph{완비} 뇌터 국소환의 구조에
관한 \emph{Cohen의 정리들}(19.8)을 포함하며, 이 정리들은 제IV장의
§§~6과 7에서 중요한 역할을 할 것이다.

\medskip\noindent\textbf{IV.} --- 다른 뇌터 국소환 위에서 형식적으로
매끄러운 뇌터 국소환들을 연구하는 의의는 매끄러운 사상에 대한
다음의 “점별” 특징짓기에서 비롯된다. $X$와 $Y$가 국소 뇌터
준스킴이고 $f:X\to Y$가 국소 유한형 사상이면, $f$가
매끄러울 필요충분조건은 모든 $x\in X$에 대하여 환
$\mathcal O_x$가 $\mathcal O_{f(x)}$ 위에서 형식적으로 매끄러운
것이다. 특히 $k$가 완전체이고 $Y=\operatorname{Spec}(k)$이면,
$f$가 매끄럽다는 것은 $X$가 \emph{정칙} 준스킴이라는 것과 동치이다.

\medskip\noindent\textbf{V.} --- 끝으로 §§~20, 21, 22에서는
형식적으로 매끄러운 대수의 개념이 \emph{켈러 미분} 이론에 자연스럽게
등장하며, 두 이론이 서로를 밝혀 준다는 것을 볼 것이다.
\end{env}

\begin{env}[19.0.3]
\label{0.19.0.3-fr}
이 절과 뒤의 절들 전체에서 위상환과 위상 가군은 \emph{선형 위상을
갖는다}고 가정한다($\mathbf{0}_{\mathrm I}$, 7.1.1). 또한 명시적으로
반대로 말하지 않는 한, 다루는 위상환은 \emph{가환}이라고 가정한다.
$A$, $B$가 두 위상환이고 $\rho:A\to B$가 $B$에 $A$-대수 구조를
정하는 환 준동형임을 상기하자. 이때 주어진 위상들에 대하여
$\rho$가 \emph{연속}이면 $B$를 \emph{위상 $A$-대수}라고 한다.

간단히 쓰기 위해, 위상환 $A$(각각 위상 $A$-가군 $M$)에서
“\emph{열린 아이디얼들의 기본계}(각각
\emph{열린 부분가군들의 기본계})”
라고 하면 “아이디얼들(각각 부분가군들)로 이루어진 $0$의 근방
기본계”를 뜻하기로 한다.

위상환 $A$와 $A$-가군 $M$이 주어졌다고 하자. $\mathfrak J$가
열린 아이디얼들의 기본계를 달릴 때 부분가군 $\mathfrak J M$들은
$M$을 위상 $A$-가군으로 만드는 한 위상의 열린 부분가군 기본계를
이룬다. 이 위상을 $A$의 위상에서 \emph{유도된} 위상이라고 한다.

위상 $A$-가군 $M$의 위상이 $A$의 위상에서 유도된 위상보다
\emph{덜 섬세하다}고 하자. 그러면 $N$이 $M$의 \emph{열린}
부분가군일 때 이산 $A$-가군 $M/N$은 $A$의 어떤 \emph{열린
아이디얼에 의해 소멸된다}. 실제로 가정에 의해
$\mathfrak R.M\subset N$인 열린 아이디얼 $\mathfrak R$이 존재한다.

\oldpage[0\textsubscript{IV-1}]{71}

$M$, $N$이 두 위상 $A$-가군이고 두 위상 모두 $A$의 위상에서
유도되었다고 하자. 그러면 모든 $A$-준동형 $u:M\to N$은
\emph{연속}이다. 실제로 $N$에서 $0$의 모든 근방 $V$에 대하여
$\mathfrak JN\subset V$인 $A$의 열린 아이디얼 $\mathfrak J$이
가정에 의해 존재하므로,
$u(\mathfrak JM)\subset\mathfrak JN\subset V$이다.

$B$가 위상 $A$-대수이면 $A$의 위상에서 \emph{유도된} $B$의 위상은
주어진 위상보다 \emph{더 섬세하다}. 실제로 $B$의 모든 열린
아이디얼 $\mathfrak R$에 대하여
$\mathfrak JB\subset\mathfrak R$인 $A$의 열린 아이디얼
$\mathfrak J$이 가정에 의해 존재한다.
\end{env}

\subsection{형식적 전사 준동형과 형식적 단사 준동형}
\label{subsection:0.19.1-fr}

\begin{proposition}[19.1.1]
\label{0.19.1.1-fr}
$A$를 위상환, $M$, $N$을 두 위상 $A$-가군,
$(W_\lambda)$를 $N$ 안의 열린 부분가군들의 기본계,
$u:M\to N$을 연속 $A$-준동형이라 하자. $\widehat M$,
$\widehat N$을 각각 $M$, $N$의 분리 완비화라 하고,
$\widehat u:\widehat M\to\widehat N$을 $u$의 분리 완비화들로의
연속 연장이라 하자.
\begin{enumerate}
  \item[(i)] 다음 조건들은 서로 동치이다.
  \begin{enumerate}
    \item[a)] $u(M)$은 $N$에서 \emph{조밀하다}.
    \item[a$'$)] $\widehat u(\widehat M)$은 $\widehat N$에서
    \emph{조밀하다}.
    \item[a$''$)] 모든 $\lambda$에 대하여 합성 준동형
    $M\xrightarrow{u}N\to N/W_\lambda$는 \emph{전사}이다.
  \end{enumerate}
  \item[(ii)] 다음 조건들은 서로 동치이다.
  \begin{enumerate}
    \item[b)] $u$에 의한 $N$의 위상의 역상은 $M$의 위상과 같다.
    \item[b$'$)] $\widehat u$는 위상 $\widehat A$-가군
    $\widehat M$에서 $\widehat N$의 위상 $\widehat A$-부분가군
    $\widehat u(\widehat M)$ 위로 가는 \emph{동형사상}이다
    (이 부분가군은 반드시 \emph{닫혀} 있다).
    \item[b$''$)] $u^{-1}(W_\lambda)$들은 $M$에서 $0$의 근방
    기본계를 이룬다.
  \end{enumerate}
\end{enumerate}
이는 분리 완비화의 정의와 $N/W_\lambda$가 이산이라는 사실에서
곧바로 따른다.
\end{proposition}

\begin{env}[19.1.2]
\label{0.19.1.2-fr}
(19.1.1)의 (i)(각각 (ii))의 동치인 조건들이 성립하면 $u$를
\emph{형식적 전사 준동형}(각각 \emph{형식적 단사 준동형})이라고
한다. $u$가 형식적 단사 준동형이면서 형식적 전사 준동형이면 이를
\emph{형식적 전단사 준동형}이라고 한다. (19.1.1)에 의해 이는
$\widehat u$가 위상 $\widehat A$-가군 $\widehat M$에서 위상
$\widehat A$-가군 $\widehat N$ 위로 가는 \emph{동형사상}이라는
것과 같다.
\end{env}

\begin{proposition}[19.1.3]
\label{0.19.1.3-fr}
$A$를 위상환, $M$, $N$을 두 위상 $A$-가군,
$u:M\to N$을 연속 $A$-준동형이라 하자. $M$, $N$ 안에 각각 열린
부분가군들의 기본계 $(V_\lambda)$, $(W_\lambda)$가 존재하고, 두
기본계의 지표집합이 같으며 모든 $\lambda$에 대하여
$u(V_\lambda)\subset W_\lambda$라고 가정하자. 몫으로 내려가
$u$에서 유도되는 준동형을
$u_\lambda:M/V_\lambda\to N/W_\lambda$라 하자. 그러면 다음이
성립한다.
\begin{enumerate}
  \item[(i)] $u$가 형식적 전사 준동형일 필요충분조건은 모든
  $\lambda$에 대하여 $u_\lambda$가 전사인 것이다.
  \item[(ii)] $u$가 형식적 단사 준동형일 필요충분조건은 모든
  $\lambda$에 대하여 $V_\mu\subset V_\lambda$이고
  $\operatorname{Ker}(u_\mu)$가 표준 사상
  $M/V_\mu\to M/V_\lambda$의 핵 $V_\lambda/V_\mu$에 포함되는
  $\mu$가 존재하는 것이다.
  \item[(iii)] $u$가 형식적 전단사 준동형일 필요충분조건은 모든
  $\lambda$에 대하여 $u_\lambda$가 전사이고,
  $V_\mu\subset V_\lambda$이며 표준 준동형
  $M/V_\mu\to M/V_\lambda$가

  \oldpage[0\textsubscript{IV-1}]{72}

  \noindent
  $M/V_\mu\xrightarrow{u_\mu}N/W_\mu\to M/V_\lambda$로
  인수분해되는 $\mu$가 존재하는 것이다. 여기서 준동형
  $N/W_\mu\to M/V_\lambda$는 반드시 유일하다.
\end{enumerate}
\end{proposition}

\begin{proof}
(i)은 자명하다. 한편
$\operatorname{Ker}(u_\mu)=u^{-1}(W_\mu)/V_\mu$이고, 표준 사상
$M/V_\mu\to M/V_\lambda$의 핵은 $V_\lambda/V_\mu$이다.
따라서 (ii)의 조건은 $u^{-1}(W_\mu)\subset V_\lambda$라는 것과
같고, 이로부터 (ii)가 곧 따른다. 끝으로 $u_\mu$가 전사일 때,
$\operatorname{Ker}(u_\mu)$가 $V_\lambda/V_\mu$에 포함된다는
것은 $M/V_\mu\to M/V_\lambda$가 $v\circ u_\mu$로
인수분해된다는 것과 같다. 여기서 $v$는
$N/W_\mu\to M/V_\lambda$인 준동형이다. 실제로 이때
$N/W_\mu$는
$(M/V_\mu)/\operatorname{Ker}(u_\mu)$와 동일시된다.
\end{proof}

\begin{corollary}[19.1.4]
\label{0.19.1.4-fr}
$A$를 위상환, $M$, $N$을 위상이 $A$의 위상에서 유도된 두 위상
$A$-가군, $u:M\to N$을 형식적 전사 준동형이라 하자.
$(\mathfrak J_\lambda)$를 $A$ 안의 열린 아이디얼들의 기본계라 하자.
$u$가 형식적 전단사 준동형일 필요충분조건은 모든 $\lambda$에
대하여 몫으로 내려가 $u$에서 유도되는 준동형
$u_\lambda:M/\mathfrak J_\lambda M\to
N/\mathfrak J_\lambda N$이 전단사인 것이다.
\end{corollary}

\begin{proof}
실제로
$u(\mathfrak J_\lambda M)\subset\mathfrak J_\lambda N$이고
(19.1.3, (iii))의 판정법을 적용할 수 있다. 그런데 인수분해
\[
  M/\mathfrak J_\mu M\xrightarrow{u_\mu}
  N/\mathfrak J_\mu N\xrightarrow{v}M/\mathfrak J_\lambda M
\]
가 있으면, $u_\mu$에 의한
$\mathfrak J_\lambda M/\mathfrak J_\mu M$의 상은
$\mathfrak J_\lambda N/\mathfrak J_\mu N$이다. 따라서
$v$에 의한 $\mathfrak J_\lambda N/\mathfrak J_\mu N$의 상은
$0$이고, $v$는
\[
  N/\mathfrak J_\mu N\longrightarrow N/\mathfrak J_\lambda N
  \xrightarrow{v'}M/\mathfrak J_\lambda M
\]
로 인수분해된다. 그러면 $M/\mathfrak J_\lambda M$의 항등
자기동형사상은
\[
  M/\mathfrak J_\lambda M\xrightarrow{u_\lambda}
  N/\mathfrak J_\lambda N\xrightarrow{v'}M/\mathfrak J_\lambda M
\]
로 인수분해된다. 따라서 $u_\lambda$는 단사이고, 이미 전사임을
알고 있으므로 따름정리가 증명된다.
\end{proof}

\begin{proposition}[19.1.5]
\label{0.19.1.5-fr}
$A$를 위상환, $M$, $N$을 두 위상 $A$-가군,
$u:M\to N$을 연속 $A$-준동형이라 하자. 다음 조건들은 서로 동치이다.
\begin{enumerate}
  \item[a)] 모든 이산 위상 $A$-가군 $E$와 모든 연속 $A$-준동형
  $v:M\to E$에 대하여, $v=w\circ u$인 연속 $A$-준동형
  $w:N\to E$가 존재한다.
  \item[b)] $M$의 모든 열린 부분가군 $V'$에 대하여, $N$의 열린
  부분가군 $W''$, $V''\subset V'\cap u^{-1}(W'')$인 열린 부분가군
  $V''$ 및 준동형 $h:N/W''\to M/V'$가 존재하여 표준 준동형
  $M/V''\to M/V'$가
  \[
    M/V''\xrightarrow{u''}N/W''\xrightarrow{h}M/V'
  \]
  으로 인수분해된다. 여기서 $u''$은 몫으로 내려가 $u$에서 유도된다.
\end{enumerate}
\end{proposition}

\begin{proof}
a)가 b)를 함의함을 보이려면 $E=M/V'$와 표준 준동형
$v:M\to M/V'$에 a)를 적용하면 충분하다. 그러면
$W''=w^{-1}(0)$은 $N$의 열린 부분가군이고
$V''=u^{-1}(W'')$은 $V'$에 포함되는 $M$의 열린 부분가군이다.
이 부분가군들과 몫으로 내려가 $w$에서 유도되는 준동형
$h:N/W''\to M/V'$가 조건을 만족한다. 역으로 b)가 a)를 함의함을
보이기 위해서는 $E$를 $v(M)$으로 바꾸어 $v$가 전사인 경우로
한정할 수 있다. 이때 $V'=\operatorname{Ker}(v)$는 $M$의 열린
부분가군이므로 $E$는 이산 $A$-가군 $M/V'$과 동형이다. 따라서
b)를 적용하고

\oldpage[0\textsubscript{IV-1}]{73}

합성 $N\to N/W''\xrightarrow{h}M/V'$을 취하면 원하는 연속
$A$-준동형 $w$를 얻는다. 실제로 도식
\[
  \xymatrix{
    M\ar[r]^u\ar[d] & N\ar[d]\\
    M/V''\ar[r]_{u''} & N/W''
  }
\]
은 가환한다.
\end{proof}

(19.1.5)의 동치인 조건들이 성립하면 $u$가
\emph{형식적으로 왼쪽 가역}이라고 한다. (19.1.5)의 조건 b)는
$u^{-1}(W'')\subset V'$을 함의하므로, 이때 $u$는 형식적 단사
준동형이다. 다음 따름정리들은 이 용어를 더욱 정당화한다.

\begin{corollary}[19.1.6]
\label{0.19.1.6-fr}
$s\circ\widehat u=1_{\widehat M}$인 연속 $\widehat A$-준동형
$s:\widehat N\to\widehat M$이 존재하면, $u$는 형식적으로 왼쪽
가역이다.
\end{corollary}

이때 $\widehat u$는 $\widehat M$에서 $\widehat N$의 한
\emph{위상적 직합인자} 위로 가는 위상동형사상임을 주목하자.
(19.1.6)을 증명하려면 (19.1.5, a))의 표기에서 $E$가 이산이므로
$v:M\to E$가 연속 $\widehat A$-준동형
$\widehat v:\widehat M\to E$로 연장된다는 사실을 주목하면
충분하다. $j:M\to\widehat M$과 $j':N\to\widehat N$을 표준
준동형이라 하자. 그러면
$w=\widehat v\circ s\circ j'$로 놓을 때
$w\circ u=\widehat v\circ s\circ\widehat u\circ j
=\widehat v\circ j=v$이다.

\begin{corollary}[19.1.7]
\label{0.19.1.7-fr}
$M$과 $N$의 위상이 $A$의 위상에서 유도되었다고 하고,
$(\mathfrak J_\lambda)$를 $A$ 안의 열린 아이디얼들의 기본계라 하자.
$u$가 형식적으로 왼쪽 가역일 필요충분조건은 모든 $\lambda$에
대하여 몫으로 내려가 $u$에서 유도되는 준동형
$u_\lambda:M/\mathfrak J_\lambda M\to N/\mathfrak J_\lambda N$이
왼쪽 가역인 것이다. 다시 말해, $u_\lambda$가
$M/\mathfrak J_\lambda M$에서 $N/\mathfrak J_\lambda N$의 한
직합인자 위로 가는 동형사상인 것이다.
\end{corollary}

실제로 이 조건은 충분하다. (19.1.5)의 조건 b)에서
$V'=\mathfrak J_\lambda M$일 때
$W''=\mathfrak J_\lambda N$, $V''=V'$로 놓고
$h\circ u_\lambda$가 $M/\mathfrak J_\lambda M$의 항등사상이 되는
$h$를 택하면 되기 때문이다. 역으로 $u$가 형식적으로 왼쪽
가역이면, 모든 $\lambda$에 대하여 (19.1.5, b))에 의해
$\mathfrak J_\mu\subset\mathfrak J_\lambda$인 $\mu$와 준동형
$h:N/\mathfrak J_\mu N\to M/\mathfrak J_\lambda M$이 존재하여 표준
준동형 $M/\mathfrak J_\mu M\to M/\mathfrak J_\lambda M$이
\[
  M/\mathfrak J_\mu M\xrightarrow{u_\mu}N/\mathfrak J_\mu N
  \xrightarrow{h}M/\mathfrak J_\lambda M
\]
으로 인수분해된다. 그런데
$h(\mathfrak J_\lambda(N/\mathfrak J_\mu N))
\subset\mathfrak J_\lambda(M/\mathfrak J_\lambda M)=0$이므로,
$h$는 표준적으로
\[
  N/\mathfrak J_\mu N\longrightarrow N/\mathfrak J_\lambda N
  \xrightarrow{s}M/\mathfrak J_\lambda M
\]
으로 인수분해된다. $s$가 $u_\lambda$의 왼쪽 역임은 즉시 알 수 있다.

\begin{proposition}[19.1.8]
\label{0.19.1.8-fr}
$A$를 감소하는 열린 아이디얼들의 가산 기본계
$(\mathfrak J_n)$을 갖는 위상환이라 하자. $M$, $N$을 위상이
$A$의 위상에서 유도되는 두 위상 $A$-가군, $u:M\to N$을
$A$-준동형이라 하자. $M_n=M/\mathfrak J_nM$,
$N_n=N/\mathfrak J_nN$으로 놓고, 몫으로 내려가 $u$에서 유도되는
$(A/\mathfrak J_n)$-준동형을 $u_n:M_n\to N_n$이라 하자. 모든
$n$에 대하여 $P_n=\operatorname{Coker}(u_n)$이
$(A/\mathfrak J_n)$-사영 가군이고 $M$이 분리 완비라고 하자.
그러면 $u$가 형식적으로 왼쪽 가역일 필요충분조건은 $u$가 왼쪽
가역인 것이다. 이때 $u$는 $M$에서 $N$의 한 위상적 직합인자
위로 가는 위상동형사상이다.
\end{proposition}

\oldpage[0\textsubscript{IV-1}]{74}

\begin{proof}
(19.1.7)에 따라 다음 가환 도식들이 있다.
\[
  \xymatrix{
    0\ar[r] & M_n\ar[r]^{u_n} & N_n\ar[r]^{p_n} & P_n\ar[r] & 0\\
    0\ar[r] & M_{n+1}\ar[u]^f\ar[r]_{u_{n+1}} &
      N_{n+1}\ar[u]^g\ar[r]_{p_{n+1}} & P_{n+1}\ar[u]^h\ar[r] & 0
  }
\]
여기서 각 행은 \emph{분할 완전열}이다. 다시 말해, 각 $n$에
대하여 $p_n\circ s_n=1_{P_n}$인 준동형 $s_n:P_n\to N_n$이
존재한다. $n$에 대한 귀납법으로
$p_n\circ s'_n=1_{P_n}$인 준동형 $s'_n:P_n\to N_n$을 정하여 도식
\[
  \xymatrix{
    P_n\ar[r]^{s'_n} & N_n\\
    P_{n+1}\ar[u]^h\ar[r]_{s'_{n+1}} & N_{n+1}\ar[u]^g
  }
\]
이 가환하도록 할 수 있음을 보이자. 실제로
$g\circ s_{n+1}-s'_n\circ h$는 $P_{n+1}$에서
\[
  u_n(M_n)=u_{n+1}(M_{n+1})/
  \mathfrak J_nu_{n+1}(M_{n+1})
\]
으로 가는 준동형이다. $P_{n+1}$은
$(A/\mathfrak J_{n+1})$-사영 가군이므로 이 준동형은
\[
  P_{n+1}\xrightarrow{t_{n+1}}u_{n+1}(M_{n+1})
  \longrightarrow
  u_{n+1}(M_{n+1})/\mathfrak J_nu_{n+1}(M_{n+1})
\]
으로 인수분해된다. 따라서
$s'_{n+1}=s_{n+1}-t_{n+1}$로 놓으면 원하는 조건을 만족한다.
이제 $N_n$을 $u_n(M_n)$과 $s'_n(P_n)$의 직합으로 분해하면,
$u_n$의 왼쪽 역인 준동형 $w_n:N_n\to M_n$을 얻고 도식
\[
  \xymatrix{
    N_n\ar[r]^{w_n} & M_n\\
    N_{n+1}\ar[u]^g\ar[r]_{w_{n+1}} & M_{n+1}\ar[u]^f
  }
\]
은 가환한다. 따라서 사영계 $(w_n)$에는 사영극한
$\widehat w:\widehat N\to\widehat M=M$이 있다. 이를 표준 준동형
$N\to\widehat N$과 합성하여 준동형 $w:N\to M$을 얻는다. 모든
$n$에 대하여 몫으로 내려가 $w\circ u$에서 유도되는
$M_n=M/\mathfrak J_nM$의 자기준동형 $(w\circ u)_n$은 항등이다.
$M$이 분리 완비이므로 이는 $w\circ u=1_M$을 함의한다.
\end{proof}

\begin{proposition}[19.1.9]
\label{0.19.1.9-fr}
$A$를 준허용 위상환 ($\mathbf{0}_{\mathrm I}$, 7.1.2),
$\mathfrak Q$를 $A$의 정의 아이디얼,
$(\mathfrak J_\lambda)$를 $A$ 안의 열린 아이디얼들의 기본계라 하자.
$M$, $N$을 위상이 $A$의 위상에서 유도되는 두 위상 $A$-가군이라
하고, 모든 $\lambda$에 대하여 $N/\mathfrak J_\lambda N$이
$(A/\mathfrak J_\lambda)$-사영 가군이라고 하자 (19.2.3 참조).
$u:M\to N$을 $A$-준동형이라 하자. 그러면 다음 조건들은 서로
동치이다.
\begin{enumerate}
  \item[a)] $u$는 형식적으로 왼쪽 가역이다.

  \oldpage[0\textsubscript{IV-1}]{75}

  \item[b)] 몫으로 내려가 $u$에서 유도되는 준동형
  $u_0:M/\mathfrak QM\to N/\mathfrak QN$은 왼쪽 가역이다.
\end{enumerate}
\end{proposition}

\begin{proof}
(19.1.7)에서 조건 a)는 모든 $\lambda$에 대하여 $u_\lambda$가 왼쪽
가역이라는 것과 동치임을 보았다. $\mathfrak Q$는 열린 아이디얼이므로
어떤 $\mathfrak J_\lambda$를 포함하고, 따라서 $u_0$이 왼쪽 가역임을
즉시 얻는다. 역으로 b)가 a)를 함의함을 보이기 위해, 모든
$\lambda$에 대하여 $\mathfrak Q/\mathfrak J_\lambda$가 가정상
$A/\mathfrak J_\lambda$의 멱영 아이디얼임을 주목하자. 우리의 주장은
다음 명제에서 따른다.
\end{proof}

\begin{proposition}[19.1.10]
\label{0.19.1.10-fr}
$A$를 환, $M$, $N$을 두 $A$-가군이라 하고 $N$이 사영 가군이라고
하자. $u:M\to N$을 $A$-준동형, $\mathfrak J$를 $A$의 아이디얼이라
하자. 다음 조건들 가운데 하나가 성립한다고 가정한다.
\begin{enumerate}
  \item[(i)] $\mathfrak J$는 멱영이다.
  \item[(ii)] $\mathfrak J$는 $A$의 근기에 포함되고 $M$은 유한형이다.
\end{enumerate}
그러면 $u$가 왼쪽 가역일 필요충분조건은 몫으로 내려가 $u$에서
유도되는 $(A/\mathfrak J)$-가군의 준동형
$u_0:M/\mathfrak JM\to N/\mathfrak JN$이 왼쪽 가역인 것이다.
\end{proposition}

\begin{proof}
조건이 필요함은 명백하므로 충분함을 증명하자. $v_0$을 $u_0$의 왼쪽
역이라 하자. $N$이 사영 가군이므로 합성 준동형
$N\to N/\mathfrak JN\xrightarrow{v_0}M/\mathfrak JM$은
$N\xrightarrow{v}M\to M/\mathfrak JM$으로 인수분해된다. 그러면
$w=v\circ u$는 $M$의 자기준동형이고, 몫으로 내려가 얻는
$M/\mathfrak JM$의 자기준동형 $w_0$은 항등이다. $w$ 자체가
전단사임을 보이면 충분하다. 실제로 이때 $w^{-1}\circ v$는 $u$의
왼쪽 역이다. 이제 두 경우를 나누자.

(i) 모든 $n$에 대하여 다음 도식은 가환한다.
\[
  \xymatrix{
    (\mathfrak J^n/\mathfrak J^{n+1})\otimes_{A/\mathfrak J}
      (M/\mathfrak JM)\ar[r]\ar[d]_{1\otimes\operatorname{gr}^0(w)} &
      \mathfrak J^nM/\mathfrak J^{n+1}M\ar[d]^{\operatorname{gr}^n(w)}\\
    (\mathfrak J^n/\mathfrak J^{n+1})\otimes_{A/\mathfrak J}
      (M/\mathfrak JM)\ar[r] & \mathfrak J^nM/\mathfrak J^{n+1}M
  }
\]
여기서 가로 화살표들은 \emph{전사}이고
$\operatorname{gr}^0(w)=w_0$은 항등이므로
$\operatorname{gr}^n(w)$도 항등이며, 특히 전단사이다.
$\mathfrak J$가 멱영이므로 $M$ 위의 $\mathfrak J$-준아딕 여과는
유한하다. 따라서 $w$는 전단사이다 (Bourbaki, \emph{Alg. comm.},
chap.~III, §~2, n\textsuperscript{o}~8, th.~1의 cor.~3).

(ii) 모든 $A$의 극대 아이디얼 $\mathfrak m$에 대하여
$M_{\mathfrak m}$의 자기준동형 $w_{\mathfrak m}$이 전단사임을
보이면 충분하다 (Bourbaki, \emph{Alg. comm.}, chap.~II, §~3,
n\textsuperscript{o}~3, th.~1). 그런데
$\mathfrak JA_{\mathfrak m}\subset\mathfrak mA_{\mathfrak m}$이고
$A_{\mathfrak m}/\mathfrak JA_{\mathfrak m}
=(A/\mathfrak J)_{\mathfrak m}$이므로, $A$가 국소환인 경우로
귀착된다. 또한 $\mathfrak J$가 $A$의 극대 아이디얼
$\mathfrak J_0$이라고 가정할 수 있다. 실제로 $u_0$이 왼쪽

\oldpage[0\textsubscript{IV-1}]{76}

가역이면, $u_0$을 $1_{A/\mathfrak J_0}$과 텐서하여 얻는
$u_{00}:M/\mathfrak J_0M\to N/\mathfrak J_0N$도 왼쪽 가역이다.
여기서는
\[
  M/\mathfrak J_0M=(M/\mathfrak JM)
  \otimes_{A/\mathfrak J}(A/\mathfrak J_0)
\]
이기 때문이다. 이제 $\mathfrak J$가 극대라고 하자. 그러면
$A/\mathfrak J$는 체이다. 명제의 조건 아래 $M$이 자유 $A$-가군임을
보이면 충분함은 명백하다. 실제로 $\det(w_0)$은 $\det(w)$의
$A/\mathfrak J$ 안의 표준상이고, 따라서 $\det(w)$는
$\mathfrak J$에 속하지 않으므로 가역이다. 이제
$M/\mathfrak JM$은 유한 차원 자유 $(A/\mathfrak J)$-벡터공간이므로,
유한 생성 자유 $A$-가군 $L$과 $A$-준동형 $f:L\to M$을 택하여
몫으로 내려가 $f$에서 유도되는 준동형
$f_0:L/\mathfrak JL\to M/\mathfrak JM$이 전단사가 되게 할 수 있다.
$M$이 유한형이므로 나카야마 보조정리에 따라 먼저 $f$가 전사임을
얻는다 (Bourbaki, \emph{Alg. comm.}, chap.~II, §~3,
n\textsuperscript{o}~2, prop.~4의 cor.~1). 또한 $g=u\circ f$라 놓으면,
몫으로 내려가 얻는 준동형
$g_0:L/\mathfrak JL\to N/\mathfrak JN$은 왼쪽 가역이다. 여기서
$L$은 자유이므로 증명 첫머리의 논의에 따라 $g$ 자체도 왼쪽
가역이다. 이는 $f$가 \emph{단사}임을 명백히 함의하며, 증명이
끝난다.
\end{proof}

(19.1.10)의 다음 따름정리들도 덧붙여 두자.

\begin{corollary}[19.1.11]
\label{0.19.1.11-fr}
$A$를 국소환, $k$를 그 잉여체, $M$을 유한형 $A$-가군,
$N$을 사영 $A$-가군, $u:M\to N$을 준동형이라 하자. $u$가 왼쪽
가역일 필요충분조건은 $M$의 생성계 $(x_i)_{1\leq i\leq m}$가
존재하여 $u\otimes1:M\otimes_Ak\to N\otimes_Ak$에 의한
$x_i\otimes1$들의 상들이 $N\otimes_Ak$ 안에서 선형독립인 것이다.
이때 $x_i$들은 $M$의 기저를 이룬다.
\end{corollary}

조건이 필요함은 명백하다. 실제로 $u$가 왼쪽 가역이면 $M$은
유한형 사영 $A$-가군이고, 따라서 자유이다. 역으로 조건이 성립하면
$x_i\otimes1$들은 $M\otimes_Ak$의 기저를 이루고
$u\otimes1$은 왼쪽 가역임이 명백하다. 따라서 $A$의 극대 아이디얼
$\mathfrak J$에 (19.1.10)을 적용하면 충분하다.

\begin{corollary}[19.1.12]
\label{0.19.1.12-fr}
$A$를 환, $M$을 유한형 $A$-가군, $N$을 사영 $A$-가군,
$u:M\to N$을 준동형이라 하자. 모든 소 아이디얼
$\mathfrak p\in\operatorname{Spec}(A)$에 대하여 다음 조건들은 서로
동치이다.
\begin{enumerate}
  \item[a)] 준동형
  $u_{\mathfrak p}:M_{\mathfrak p}\to N_{\mathfrak p}$는 왼쪽
  가역이다.
  \item[b)] 준동형
  $u\otimes1:M\otimes_Ak(\mathfrak p)\to
  N\otimes_Ak(\mathfrak p)$는 단사이다. 여기서
  $k(\mathfrak p)$는 $\mathfrak p$에서 $A$의 잉여체이다.
  \item[c)] $x_i\in M$인 유한 원소계 $(1\leq i\leq m)$가 존재하여
  그 $M_{\mathfrak p}$ 안의 상들이 $M_{\mathfrak p}$를 생성하고,
  $N$ 위의 $m$개의 선형 형식 $y_i^*$ $(1\leq i\leq m)$가 존재하여
  $\det(\langle y_i^*,u(x_j)\rangle)\notin\mathfrak p$이다.
  \item[d)] 어떤 $f\in A-\mathfrak p$에 대하여 준동형
  $u_f:M_f\to N_f$가 왼쪽 가역이다.
\end{enumerate}
또한 이 조건들을 만족하는 $\mathfrak p\in\operatorname{Spec}(A)$들의
집합은 $\operatorname{Spec}(A)$에서 열린집합이다.
\end{corollary}

마지막 주장은 d)에서 자명하게 따른다. $N$은 사영 가군이므로 자유
$A$-가군 $A^{(I)}$의 직합인자이다. 또한 $M$은 유한형이므로
$u(M)$은 유한한 $n$에 대한 $A^{(I)}$의 부분가군 $A^n$에 포함된다.
a), b), c), d)의 각 명제는 $N$을 $N\oplus P$로 바꾼 대응하는
명제와 동치이므로, 여기서도 $u$는 $M$을 $N$으로 보낸다고 하면서
$N$이 유한 생성 자유 가군인 경우로 귀착된다. d)가 a)를 함의함은
자명하고, a)와 b)는 (19.1.11)에 따라 동치이다. 또한 a)가 성립하면
$M_{\mathfrak p}$는 자유이고 (19.1.11), 그 상들이 이

\oldpage[0\textsubscript{IV-1}]{77}

$A_{\mathfrak p}$-가군의 기저를 이루는 원소족을 $x_i$들로 택할 수
있다. $N$이 유한 생성 자유 가군이므로 $N_{\mathfrak p}$ 위의 모든
선형 형식은 어떤 $f\in A-\mathfrak p$와 $N$ 위의 선형 형식 $v$에
대하여 $v/f$로 쓸 수 있다. 따라서 a)는 c)를 함의한다. c)가 b)를
함의함은 명백하므로, a)가 d)를 함의함만 보이면 된다. $N$은 유한
표시이므로
$(\operatorname{Hom}_A(N,M))_{\mathfrak p}$는 표준적으로
$\operatorname{Hom}_{A_{\mathfrak p}}
(N_{\mathfrak p},M_{\mathfrak p})$와 동일시된다
(Bourbaki, \emph{Alg. comm.}, chap.~II, §~2,
n\textsuperscript{o}~7, prop.~19). $w_{\mathfrak p}$를
$u_{\mathfrak p}$의 왼쪽 역이라 하자. 그러면 준동형 $w:N\to M$과
원소 $f\in A-\mathfrak p$가 존재하여
$w_{\mathfrak p}=w\otimes(1/f)1_{A_{\mathfrak p}}$이다. 따라서 관계
$w_{\mathfrak p}\circ u_{\mathfrak p}=1_{M_{\mathfrak p}}$는
$(w\circ u)\otimes1_{A_{\mathfrak p}}
=f.1_{M_{\mathfrak p}}$로도 쓸 수 있다. $M$은 유한형이므로 어떤
$g\in A-\mathfrak p$가 존재하여 자기준동형
$g((w\circ u)-f.1_M)$이 $M$의 모든 생성원에서 소멸하고, 따라서
영이다. $h=fg$, $u_h=u\otimes1_{A_h}$,
$w_h=w\otimes1_{A_h}$로 놓으면 모든 $z\in M_h$에 대하여
$w_h(u_h(z))=(f/1)z$이다. 그런데 $h/1$이 $A_h$에서 가역이므로
$f/1=f'$도 가역이고, 따라서 $f'^{-1}w_h$는 $u_h$의 왼쪽 역이다.
이로써 증명이 끝난다.

\begin{proposition}[19.1.13]
\label{0.19.1.13-fr}
(19.1.9)의 가정들이 성립하고, $(\mathfrak J_\lambda)$가 감소하는 열
$(\mathfrak J_n)$이며 $M$이 분리 완비라고 더 가정하자. 그러면
(19.1.9)의 조건 a)와 b)는 다음 조건과도 동치이다.
\begin{enumerate}
  \item[c)] $u$는 왼쪽 가역이다.
\end{enumerate}
\end{proposition}

c)가 a)를 함의함은 이미 알고 있다 (19.1.6). 역으로 a)가 성립하면
(19.1.8)의 표기에서 $M_n$은 사영 $(A/\mathfrak J_n)$-가군
$N_n$의 직합인자이다. 따라서 $P_n$도 $N_n$의 한 직합인자와
동형이고, 그러므로 사영 가군이다. 이제 (19.1.8)을 적용하면
충분하다.

끝으로 다음 명제를 주목하자.

\begin{proposition}[19.1.14]
\label{0.19.1.14-fr}
$A$를 환, $M$을 유한형 $A$-가군, $N$을 사영 $A$-가군,
$u:M\to N$을 $A$-준동형이라 하자.
\begin{enumerate}
  \item[(i)] $u$가 왼쪽 가역일 필요충분조건은 $A$의 모든 극대
  아이디얼 $\mathfrak m$에 대하여
  $u_{\mathfrak m}:M_{\mathfrak m}\to N_{\mathfrak m}$이 왼쪽
  가역인 것이다.
  \item[(ii)] $A'$을 충실 평탄 $A$-가군인 $A$-대수라 하자. $u$가
  왼쪽 가역일 필요충분조건은
  $u\otimes1:M\otimes_AA'\to N\otimes_AA'$이 왼쪽 가역인 것이다.
\end{enumerate}
\end{proposition}

(19.1.12)에서와 같이 $N$이 유한 생성 자유 가군인 경우로 한정할 수
있다. 이때 $u$가 왼쪽 가역이라는 것은 $u$가 단사이고 몫가군
$P=N/u(M)$이 \emph{사영}이라는 뜻이다. 실제로 이때 $u(M)$은
$N$의 직합인자이다. 또한 $M$이 유한형이므로 $P$는 \emph{유한
표시}이다.

(i) 조건이 필요함은 명백하다. 역으로 조건이 성립하면 $u$는
단사이고 (Bourbaki, \emph{Alg. comm.}, chap.~II, §~3,
n\textsuperscript{o}~3, th.~1), 모든 $\mathfrak m$에 대하여
$P_{\mathfrak m}=N_{\mathfrak m}/
u_{\mathfrak m}(M_{\mathfrak m})$이 사영이므로 $P$가 사영임을 안다
(\emph{loc. cit.}, §~5, n\textsuperscript{o}~2, th.~1).

(ii) 여기서도 조건이 필요함은 자명하다. 역으로 조건이 성립하면
$u$는 단사이고 ($\mathbf{0}_{\mathrm I}$, 6.4.1),
$P\otimes_AA'=\operatorname{Coker}(u\otimes1)$은 사영이며 따라서
평탄하다. 그러므로 $P$는 평탄 $A$-가군이고
($\mathbf{0}_{\mathrm I}$, 6.6.3), 유한 표시이므로 사영이다
(Bourbaki, \emph{Alg. comm.}, chap.~II, §~5,
n\textsuperscript{o}~2, th.~1의 cor.~2).

\oldpage[0\textsubscript{IV-1}]{78}

\begin{remark}[19.1.15]
\label{0.19.1.15-fr}
형식적으로 왼쪽 가역인 준동형이라는 개념과 ``쌍대''인 개념은
\emph{형식적으로 오른쪽 가역}인 준동형이라는 개념이다. 그러한 연속
$A$-준동형 $u:M\to N$은 정의에 따라 다음 조건을 만족한다. $N$의
임의의 열린 부분가군 $W'$에 대하여 $M$의 열린 부분가군
$V'\subset u^{-1}(W')$, $N$의 열린 부분가군 $W''\subset W'$, 그리고
준동형 $h:N/W''\to M/V'$이 존재하여 표준 준동형 $N/W''\to N/W'$이
다음과 같이 인수분해된다.
\[
  N/W''\xrightarrow{h}M/V'\xrightarrow{u'}N/W'.
\]
여기서 $u'$은 $u$에서 몫으로 내려가 얻은 준동형이다. 이 조건으로부터
$u$가 형식적 전사 준동형임이 따른다. 또한
$\widehat u\circ r=1_{\widehat N}$인 연속 $\widehat A$-준동형
$r:\widehat N\to\widehat M$이 존재하면 $u$가 형식적으로 오른쪽
가역임을 곧바로 확인할 수 있다. (19.1.7)의 조건들이 충족되는 경우,
$u$가 형식적으로 오른쪽 가역일 필요충분조건은 모든 $u_\lambda$가
오른쪽 가역인 것이다. 즉 $\operatorname{Ker}(u_\lambda)$는
$M/\mathfrak J_\lambda M$의 직합인자이고, 어떤 부분가군 $L_\lambda$에
대하여 $M/\mathfrak J_\lambda M=
\operatorname{Ker}(u_\lambda)\oplus L_\lambda$이며 $u_\lambda$의
$L_\lambda$에 대한 제한은 $L_\lambda$에서
$N/\mathfrak J_\lambda N$ 위로 가는 동형이다. (19.1.8)과 (19.1.9)에
대응하는 명제의 서술과 증명은 독자에게 맡긴다. (19.1.8)의 대응
명제에서는 $M$이 분리되고 완비이며 $M_n$이
$(A/\mathfrak J_n)$-사영가군이라고 가정해야 한다. (19.1.9)의 대응
명제에서는 $N/\mathfrak J_\lambda N$에 관한 가정을 제거하는 대신,
모든 $\lambda$에 대하여 $M/\mathfrak J_\lambda M$이
$(A/\mathfrak J_\lambda)$-사영가군이라고 가정해야 한다.
\end{remark}

\subsection{형식적으로 사영인 가군}
\label{subsection:0.19.2-fr}

\begin{definition}[19.2.1]
\label{0.19.2.1-fr}
$A$를 위상환이라 하자. 위상 $A$-가군 $M$이 다음 조건을 만족할 때
$M$을 \emph{형식적으로 사영}이라고 한다. $A$의 임의의 열린 아이디얼
$\mathfrak J$, 임의의 두 이산 $(A/\mathfrak J)$-가군 $P$, $Q$, 임의의
전사 $A$-준동형 $u:P\to Q$, 그리고 임의의 연속 $A$-준동형
$v:M\to Q$에 대하여 $v=u\circ w$인 연속 $A$-준동형 $w:M\to P$가
존재한다.
\end{definition}

\begin{env}[19.2.2]
\label{0.19.2.2-fr}
(19.2.1)의 조건을 확인할 때에는 $Q$를 $v(M)$으로, $P$를
$u^{-1}(v(M))$으로 바꾸어 $v$ 자체가 전사인 경우로 한정해도 됨은
명백하다. 그러면 $Q$는 $M/V$와 동형이다. 여기서 $V$는
$\mathfrak JM\subset V$인 $M$의 열린 부분가군이다. 따라서 (19.2.1)의
조건은 다음 조건과 동치이다. 임의의 이산 $(A/\mathfrak J)$-가군 $P$와
임의의 전사 $A$-준동형 $u:P\to M/V$에 대하여, $M$ 안에서
$V'\subset V$인 열린 부분가군 $V'$과 $A$-준동형 $w:M/V'\to P$가
존재하여 표준 준동형 $M/V'\to M/V$이 다음과 같이 인수분해된다.
\[
  M/V'\xrightarrow{w}P\xrightarrow{u}M/V.
\]
이 조건은 $\mathfrak J$가 아이디얼들로 이루어진 $A$의 $0$의 근방들의
한 기본계를 달리는 경우에만 확인하면 충분함에 유의하자.
\end{env}

\begin{env}[19.2.3]
\label{0.19.2.3-fr}
$A$ 안의 열린 아이디얼들의 한 기본계 $(\mathfrak J_\lambda)$와 $M$의
열린 부분가군들의 한 기본계 $(V_\lambda)$가 존재한다고 가정하자.
두 기본계는 \emph{같은} 지표집합을 가지며, 모든 $\lambda$에 대하여
$M/V_\lambda$는 사영 $(A/\mathfrak J_\lambda)$-가군이라고 하자. 그러면
$M$은 형식적으로 사영이다. 실제로 (19.2.2)의 기호 아래
$\mathfrak J_\lambda\subset\mathfrak J$이고 $V_\lambda\subset V$인
$\lambda$를 택하면 충분하다. $P$도 $(A/\mathfrak J_\lambda)$-가군이므로,
표준 준동형 $M/V_\lambda\to M/V$을
$M/V_\lambda\to P\to M/V$로 인수분해할 수 있다. 이는 그 표준
준동형이 $(A/\mathfrak J_\lambda)$-준동형이고 $M/V_\lambda$가 사영
$(A/\mathfrak J_\lambda)$-가군이라고 가정했기 때문이다.

이 번호의 더 강한 조건이 만족될 때 $M$을 \emph{엄밀하게 형식적으로
사영}이라고 한다.
\end{env}

\begin{proposition}[19.2.4]
\label{0.19.2.4-fr}
$A$를 위상환, $M$을 그 위상이 $A$의 위상에서 유도된 위상
$A$-가군이라 하자. $M$이 형식적으로 사영일 필요충분조건은 $M$이
엄밀하게 형식적으로 사영인 것이다.
\end{proposition}

\oldpage[0\textsubscript{IV-1}]{79}

\begin{proof}
이 조건이 충분함은 방금 보았다. 역으로 $M$이 형식적으로 사영이라고
가정하고, $(\mathfrak J_\lambda)$를 $A$의 열린 아이디얼들의 한
기본계라 하자. 각 $\lambda$에 대하여 $P$를 자유
$(A/\mathfrak J_\lambda)$-가군이라 하고,
$u:P\to M/\mathfrak J_\lambda M$을 전사
$(A/\mathfrak J_\lambda)$-준동형이라 하자. 그러면
$\mathfrak J_\mu\subset\mathfrak J_\lambda$이고 표준 준동형
$M/\mathfrak J_\mu M\to M/\mathfrak J_\lambda M$이 다음과 같이
인수분해되는 $\mathfrak J_\mu$가 존재한다.
\[
  M/\mathfrak J_\mu M\xrightarrow{w}P\xrightarrow{u}
  M/\mathfrak J_\lambda M.
\]
그런데
$w(\mathfrak J_\lambda(M/\mathfrak J_\mu M))
\subset\mathfrak J_\lambda P=0$이므로 $w$는
$M/\mathfrak J_\lambda M\xrightarrow{v}P$를 거쳐 인수분해된다.
또한 $u\circ v$가 $M/\mathfrak J_\lambda M$의 항등사상임은 명백하다.
따라서 이 $(A/\mathfrak J_\lambda)$-가군은 사영이다.
\end{proof}

\begin{proposition}[19.2.5]
\label{0.19.2.5-fr}
$A$를 위상환, $M$을 위상 $A$-가군이라 하자.
\begin{enumerate}
  \item[(i)] $M$이 형식적으로 사영일 (각각 엄밀하게 형식적으로
  사영일) 필요충분조건은 위상 $\widehat A$-가군 $\widehat M$이
  그러한 것이다.
  \item[(ii)] $A'$을 위상 $A$-대수라 하자. $M$이 형식적으로
  사영이면 (각각 엄밀하게 형식적으로 사영이면), 텐서곱 위상을 준
  $M\otimes_A A'$은 형식적으로 사영인 (각각 엄밀하게 형식적으로
  사영인) 위상 $A'$-가군이다.
\end{enumerate}
\end{proposition}

\begin{proof}
\begin{enumerate}
  \item[(i)] $\mathfrak J$가 $A$의 열린 아이디얼 전체를 달릴 때
  그 분리 완비화 $\widehat{\mathfrak J}$는 $\widehat A$의 열린
  아이디얼 전체를 달리고, $V$가 $M$의 열린 부분가군 전체를 달릴 때
  그 분리 완비화 $\widehat V$는 $\widehat M$의 열린 부분가군 전체를
  달린다. 또한 표준 동형 아래
  $\widehat A/\widehat{\mathfrak J}=A/\mathfrak J$이고
  $\widehat M/\widehat V=M/V$이다. 형식적으로 사영인 (각각 엄밀하게
  형식적으로 사영인) 가군의 개념은 $A/\mathfrak J$들과 $M/V$들만을
  사용하므로 (i)가 곧바로 따른다.

  \item[(ii)] 먼저 $M$이 형식적으로 사영이라고 가정하고
  $M'=M\otimes_A A'$이라 놓자. $\mathfrak J'$을 $A'$의 열린
  아이디얼, $P'$, $Q'$을 두 이산 $(A'/\mathfrak J')$-가군,
  $u':P'\to Q'$을 전사 $A'$-준동형, $v':M'\to Q'$을 연속
  $A'$-준동형이라 하자. $\mathfrak JA'\subset\mathfrak J'$인 $A$의
  열린 아이디얼 $\mathfrak J$가 존재하므로, $P'$과 $Q'$은 이산
  $(A/\mathfrak J)$-가군이기도 하다. 표준 준동형 $j:M\to M'$을 써서
  연속 $A$-준동형
  $v:M\xrightarrow{j}M'\xrightarrow{v'}Q'$을 생각하면, 가정에 따라
  $v=u'\circ w$인 연속 $A$-준동형 $w:M\to P'$이 존재한다. 그런데
  $P'$이 위상 $A'$-가군이므로 $w$는
  $M\xrightarrow{j}M'\xrightarrow{w'}P'$을 거쳐 인수분해된다. 여기서
  $w'$은 연속 $A'$-준동형이다. 이제
  $v'\circ j=(u'\circ w')\circ j$이므로 $v'=u'\circ w'$이다.

  다음으로 $M$이 엄밀하게 형식적으로 사영이라고 가정하자.
  $(\mathfrak J_\lambda)$와 $(V_\lambda)$를 각각 $A$의 열린
  아이디얼들과 $M$의 열린 부분가군들의 기본계로서, 각 $\lambda$에
  대하여 $M/V_\lambda$가 사영 $(A/\mathfrak J_\lambda)$-가군이 되게
  하는 기본계라 하자. 또 $(\mathfrak J'_\mu)$를 $A'$의 열린
  아이디얼들의 한 기본계라 하자. $M'$에서 다음 부분가군들을 취하면
  $0$의 근방들의 한 기본계를 얻는다.
  \[
    \operatorname{Im}(M\otimes_A\mathfrak J'_\mu)
    +\operatorname{Im}(V_\lambda\otimes_A A')=W_{\lambda\mu},
  \]
  여기서 $\lambda$와 $\mu$는
  $\mathfrak J_\lambda A'\subset\mathfrak J'_\mu$를 만족한다. 또한
  \[
    (M\otimes_A A')/W_{\lambda\mu}
    =(M/V_\lambda)\otimes_{A/\mathfrak J_\lambda}
      (A'/\mathfrak J'_\mu)
  \]
  이고 $M/V_\lambda$는 사영 $(A/\mathfrak J_\lambda)$-가군이므로,
  $M'/W_{\lambda\mu}$는 사영 $(A'/\mathfrak J'_\mu)$-가군이다.
\end{enumerate}
\end{proof}

\subsection{형식적으로 매끄러운 대수}
\label{subsection:0.19.3-fr}

\begin{definition}[19.3.1]
\label{0.19.3.1-fr}
$A$를 위상환, $B$를 위상 $A$-대수라 하자. 임의의 이산 위상
$A$-대수 $C$, $C$의 임의의 멱영 아이디얼 $\mathfrak J$, 그리고 임의의
연속 $A$-준동형 $u:B\to C/\mathfrak J$에 대하여 $u$가 다음과 같이
인수분해될 때, $B$를 \emph{형식적으로 매끄러운} $A$-대수라고 한다.
\[
  B\xrightarrow{v}C\xrightarrow{\varphi}C/\mathfrak J,
\]
여기서 $v$는 연속 준동형이고 $\varphi$는 표준 준동형이다.
\end{definition}

\oldpage[0\textsubscript{IV-1}]{80}

정의 (19.3.1)는 다음 성질이 성립한다는 것과 같다.
\begin{enumerate}
  \item[(P)] $A$-대수 $B$의 임의의 열린 아이디얼 $\mathfrak R$과
  임의의 $A$-준동형 $u':B/\mathfrak R\to C/\mathfrak J$에 대하여,
  $B$의 열린 아이디얼 $\mathfrak R'\subset\mathfrak R$이 존재하여
  준동형 $B\to B/\mathfrak R\xrightarrow{u'}C/\mathfrak J$이 다음과
  같이 인수분해된다.
  \[
    B\to B/\mathfrak R'\xrightarrow{v'}C\to C/\mathfrak J,
  \]
  여기서 $v'$은 $A$-준동형이다.
\end{enumerate}
실제로 $u:B\to C/\mathfrak J$가 연속 $A$-준동형이면 그 핵
$\mathfrak R$은 $B$의 열린 아이디얼이므로 $u$는
$B\to B/\mathfrak R\xrightarrow{u'}C/\mathfrak J$으로 인수분해된다.
(P)가 성립하면 $u'$에 (P)를 적용하고, (19.3.1)의 조건을 만족시키는
$v$로 합성 $B\to B/\mathfrak R'\xrightarrow{v'}C$를 취하면 충분하다.
역으로 $B$가 형식적으로 매끄러운 $A$-대수라고 가정하자. $B$의 열린
아이디얼 $\mathfrak R$과 $A$-준동형
$u':B/\mathfrak R\to C/\mathfrak J$이 주어졌다고 하고, 합성
$u:B\to B/\mathfrak R\xrightarrow{u'}C/\mathfrak J$에 정의 (19.3.1)을
적용하자. $u$가
$B\xrightarrow{v}C\to C/\mathfrak J$으로 인수분해되는 연속
$A$-준동형 $v:B\to C$를 취하면, $B$의 아이디얼
$\mathfrak R'=\operatorname{Ker}(v)\cap\mathfrak R$은 열리고
$\mathfrak R$에 포함된다. 따라서 $v$는
$B\to B/\mathfrak R'\xrightarrow{v'}C$로 인수분해되고, 이 $v'$은
성질 (P)를 만족한다.

\begin{proposition}[19.3.2]
\label{0.19.3.2-fr}
$A$를 이산환, $V$를 사영 $A$-가군이라 하자. 대칭 대수
$B=\mathbf S_A^\bullet(V)$에 이산 위상을 주면 $B$는 형식적으로
매끄러운 $A$-대수이다.
\end{proposition}

\begin{proof}
$C$와 $\mathfrak J$의 뜻을 (19.3.1)에서와 같이 하고,
$u:B\to C/\mathfrak J$를 $A$-대수 준동형이라 하자. 이를
$V=\mathbf S_A^1(V)$에 제한하면 $A$-가군 준동형
$u_1:V\to C/\mathfrak J$을 얻는다. $V$가 사영이므로 $u_1$은
$V\xrightarrow{v_1}C\xrightarrow{\varphi}C/\mathfrak J$으로
인수분해된다. 또한 $v_1$은 $A$-대수 준동형
$v:\mathbf S_A^\bullet(V)\to C$로 연장된다. 합성
$\varphi\circ v$는 $V$ 위에서 $u$와 일치하므로
$\varphi\circ v=u$이다. 이로써 명제가 증명된다.
\end{proof}

\begin{corollary}[19.3.3]
\label{0.19.3.3-fr}
$A$가 이산환이면, 이산 위상을 준 모든 다항식 대수
$B=A[T_\alpha]_{\alpha\in I}$는 형식적으로 매끄러운 $A$-대수이다.
\end{corollary}

\begin{proposition}[19.3.4]
\label{0.19.3.4-fr}
$A$를 위상환이라 하고,
$B=A[[T_\alpha]]_{\alpha\in I}$를 형식적 멱급수
$\sum_{\lambda\in\mathbb{N}^{(I)}}c_\lambda T^\lambda$의 환이라 하자.
즉 $B$는 모노이드 $\mathbb{N}^{(I)}$의 원소들을 지표로 하는 모든
형식합의 $A$-대수이며, $A$-가군으로서 곱
$A^{\mathbb{N}^{(I)}}$와 동일시된다. $B$에 곱위상을 주면 $B$는
형식적으로 매끄러운 $A$-대수이다.
\end{proposition}

$(\mathfrak L_\mu)_{\mu\in M}$을 $A$의 열린 아이디얼들의 한 기본계라
하자. $I$의 임의의 유한 부분집합 $H$, 임의의 $\mu\in M$, 임의의 정수
$n$에 대하여 $\mathfrak K_{H,\mu,n}$을 다음 $(c_\lambda)$들로 이루어진
$B$ 안의 $0$의 근방이라 하자. 곧
$\lambda=(\lambda_\alpha)_{\alpha\in I}$가 $\alpha\notin H$에 대하여
$\lambda_\alpha=0$이고
$\sum_{\alpha\in H}\lambda_\alpha\leq n$이면
$c_\lambda\in\mathfrak L_\mu$인 것들이다.
$\mathfrak K_{H,\mu,n}$들이 $B$ 안의 $0$의 근방들의 한 기본계를 이루고
$B$의 아이디얼임을 곧바로 확인할 수 있다. 따라서 곱위상은 $B$의
$A$-대수 구조와 양립한다.

먼저 같은 기호 아래 다음을 기록하자.
\begin{lemma}[19.3.4.1]
\label{0.19.3.4.1-fr}
$E$를 이산 $A$-대수라 하자.
\begin{enumerate}
  \item[(i)] $f:B\to E$가 연속 $A$-준동형이면, $I$의 유한 부분집합
  $H$가 존재하여 $\alpha\notin H$일 때 $f(T_\alpha)=0$이고, 모든
  $\alpha\in H$에 대하여 $f(T_\alpha)$는 $E$에서 멱영이다.
  \item[(ii)] 역으로 $H$를 $I$의 유한 부분집합,
  $(z_\alpha)_{\alpha\in H}$를 $E$의 멱영 원소들의 족이라 하자. 그러면
  $\alpha\in H$일 때 $g(T_\alpha)=z_\alpha$이고 $\alpha\notin H$일 때
  $g(T_\alpha)=0$인 연속 $A$-준동형 $g:B\to E$가 유일하게 존재한다.
\end{enumerate}
\end{lemma}

\oldpage[0\textsubscript{IV-1}]{81}

\begin{proof}
(i)는 $f^{-1}(0)$이 곱 $A$-가군 $A^{\mathbb{N}^{(I)}}$ 안의 $0$의
근방이라는 사실에서 곧바로 따른다. 따라서 유한 개의
$\lambda\in\mathbb{N}^{(I)}$를 제외하면 $f(T^\lambda)=0$이다.
(ii)를 증명하려면 다항식환 $B'=A[T_\alpha]_{\alpha\in I}$가 $B$에서
조밀함에 유의하면 충분하다. $g|B'$의 존재와 유일성은 자명하고,
$\alpha\in H$인 $z_\alpha$들이 멱영이라는 가정에서 그 연속성이
따른다. 실제로 모든 $\alpha\in H$에 대하여 $z_\alpha^n=0$이면,
유한 지지를 갖는 $\lambda=(\lambda_\alpha)_{\alpha\in I}$ 가운데
$\alpha\notin H$일 때 $\lambda_\alpha=0$이고 $\alpha\in H$일 때
$\lambda_\alpha\leq n$인 유한 개의 $\lambda$를 제외하면
$g(T^\lambda)=0$이다.
\end{proof}

\begin{proof}[명제~{\rm(19.3.4)}의 증명]
이 보조정리를 적용하자. $C$와 $\mathfrak J$의 뜻을 (19.3.1)에서와
같이 하면, $I$의 어떤 유한 부분집합 $H$ 밖에서는
$u(T_\alpha)=0$이고, $\alpha\in H$에 대한
$z_\alpha=u(T_\alpha)$는 $C/\mathfrak J$에서 멱영이다.
$\mathfrak J$가 멱영이므로, $C/\mathfrak J$에서의 표준 상이 각각
$z_\alpha$인 $C$의 멱영 원소들의 족
$(x_\alpha)_{\alpha\in H}$가 존재한다. $\alpha\notin H$일 때
$v(T_\alpha)=0$이고 $\alpha\in H$일 때
$v(T_\alpha)=x_\alpha$인 연속 $A$-준동형 $v:B\to C$를 취하면,
$u$가 $B\xrightarrow{v}C\to C/\mathfrak J$으로 인수분해됨은 명백하다.
\end{proof}

\begin{proposition}[19.3.5]
\label{0.19.3.5-fr}
$A$를 위상환이라 하자.
\begin{enumerate}
  \item[(i)] $A$는 형식적으로 매끄러운 $A$-대수이다.
  \item[(ii)] $B$가 형식적으로 매끄러운 $A$-대수이고 $C$가 형식적으로
  매끄러운 $B$-대수이면, $C$는 형식적으로 매끄러운 $A$-대수이다.
  \item[(iii)] $B$를 형식적으로 매끄러운 $A$-대수, $A'$을 위상
  $A$-대수라 하자. 그러면 위상 $A'$-대수 $B\otimes_A A'$은
  형식적으로 매끄럽다 ($\mathbf 0_{\rm I}$, 7.7.5와 7.7.6).
  \item[(iv)] $B$를 위상 $A$-대수, $S$와 $T$를 각각 $A$와 $B$의
  곱셈적 부분집합으로서 $S$의 $B$에서의 표준 상이 $T$에 포함되게
  하는 것이라 하자. $B$가 형식적으로 매끄러운 $A$-대수이면
  $T^{-1}B$는 형식적으로 매끄러운 $S^{-1}A$-대수이다.
  \item[(v)] $B_i$ ($1\leq i\leq n$)를 위상 $A$-대수들이라 하자.
  $\prod_{i=1}^n B_i$가 형식적으로 매끄러운 $A$-대수일 필요충분조건은
  각 $B_i$가 그러한 것이다.
\end{enumerate}
\end{proposition}

\oldpage[0\textsubscript{IV-1}]{82}

\begin{proof}
\begin{enumerate}
  \item[(i)] $C$를 이산 $A$-대수,
  $\varphi:C\to C/\mathfrak J$를 $C$에서 그 임의의 $A$-몫대수로 가는
  표준 준동형이라 하자. $A$에서 $C/\mathfrak J$로 가는
  \emph{유일한} $A$-준동형은 합성
  $A\xrightarrow{\psi}C\xrightarrow{\varphi}C/\mathfrak J$이다. 여기서
  $\psi$는 $C$의 $A$-대수 구조를 정의하는 준동형이다. $\psi$가
  연속이므로 (19.3.1)의 조건은 자명하게 충족된다.

  \item[(ii)] $\alpha:A\to B$와 $\beta:B\to C$를 각각 $B$의
  $A$-대수 구조와 $C$의 $B$-대수 구조를 정의하는 연속 준동형이라
  하자. 따라서 $\beta\circ\alpha$는 $C$의 $A$-대수 구조를 정의한다.
  $E$를 이산 $A$-대수, $\mathfrak L$을 $E$의 멱영 아이디얼,
  $u:C\to E/\mathfrak L$을 연속 $A$-준동형이라 하자. 그러면
  $u\circ\beta\circ\alpha$는 $E/\mathfrak L$의 $A$-대수 구조를
  정의하는 준동형이다. $B$가 형식적으로 매끄러운 $A$-대수이므로
  연속 $A$-준동형 $u\circ\beta:B\to E/\mathfrak L$은
  $B\xrightarrow{v}E\to E/\mathfrak L$으로 인수분해된다. 여기서 $v$는
  연속 $A$-준동형이다. 이제 $v$와 $u\circ\beta$는 각각 $E$와
  $E/\mathfrak L$에 위상 $B$-대수 구조를 주며, 이 구조 아래
  $E/\mathfrak L$은 여전히 $B$-대수 $E$의 이산 몫대수이다. 더욱이
  $u$는 연속 $B$-준동형이므로
  $C\xrightarrow{w}E\to E/\mathfrak L$으로 인수분해된다. 여기서 $w$는
  연속 $B$-준동형이다. 합성 $w\circ\beta\circ\alpha$가 $E$의
  $A$-대수 구조를 정의하는 준동형이므로 $w$는 실제로 연속
  $A$-준동형이다.

  \item[(iii)] $C$를 이산 위상 $A'$-대수, $\mathfrak J$를 $C$의 멱영
  아이디얼, $u:B\otimes_A A'\to C/\mathfrak J$를 연속
  $A'$-준동형이라 하자. $u$를 표준 준동형
  $\rho:B\to B\otimes_A A'$과 합성하면 연속 $A$-준동형을 얻는다.
  실제로 $A\to B\xrightarrow{\rho}B\otimes_A A'$은
  $A\to A'\xrightarrow{\sigma}B\otimes_A A'$과 같다. 따라서 가정에
  따라 이 합성은 $B\xrightarrow{v}C\to C/\mathfrak J$으로
  인수분해된다. 여기서 $v$는 연속 $A$-준동형이다. $w:A'\to C$를
  $C$의 $A'$-대수 구조를 정의하는 준동형이라 하자. 두 합성
  $A\to B\xrightarrow{v}C$와 $A\to A'\xrightarrow{w}C$가 같으므로,
  $v=f\circ\rho$이고 $w=f\circ\sigma$인 연속 환 준동형
  $f:B\otimes_A A'\to C$가 존재한다
  ($\mathbf 0_{\rm I}$, 7.7.6). 이제 합성
  $B\xrightarrow{\rho}B\otimes_A A'\xrightarrow{f}C\to C/\mathfrak J$과
  $B\xrightarrow{\rho}B\otimes_A A'\xrightarrow{u}C/\mathfrak J$은 같고,
  마찬가지로 합성
  $A'\xrightarrow{\sigma}B\otimes_A A'\xrightarrow{f}C
  \to C/\mathfrak J$과
  $A'\xrightarrow{\sigma}B\otimes_A A'\xrightarrow{u}C/\mathfrak J$도
  같다. 그러므로
  $u:B\otimes_A A'\xrightarrow{f}C\to C/\mathfrak J$이라는 인수분해를
  얻고, (iii)가 증명된다.
  \item[(iv)] $S^{-1}A$와 $T^{-1}B$에 주는 위상은 자연스럽게 각각
  다음 $0$의 근방 기본계를 갖는 위상이다. 곧 $\mathfrak J$가 $A$의
  열린 아이디얼들의 한 기본계를 달릴 때의 $S^{-1}\mathfrak J$들과,
  $\mathfrak R$이 $B$의 열린 아이디얼들의 한 기본계를 달릴 때의
  $T^{-1}\mathfrak R$들이다 ($\mathbf 0_{\rm I}$, 7.6.1).
  $\alpha:A\to B$가 표준 준동형이면, 여기서 유도되는 표준 준동형
  $\alpha':S^{-1}A\to T^{-1}B$는 연속임이 명백하다. 이제 $C$를 이산
  위상 $S^{-1}A$-대수, $\mathfrak J$를 $C$의 멱영 아이디얼,
  $u:T^{-1}B\to C/\mathfrak J$를 연속 $S^{-1}A$-준동형이라 하자.
  합성 $B\to T^{-1}B\xrightarrow{u}C/\mathfrak J$은 연속
  $A$-준동형이므로 가정에 따라
  $B\xrightarrow{v}C\to C/\mathfrak J$으로 인수분해된다. 여기서 $v$는
  연속 $A$-준동형이다. 모든 $t\in T$에 대하여 $t/1$은 $T^{-1}B$에서
  가역이므로 $u(t/1)$은 $C/\mathfrak J$에서 가역이다.
  $\mathfrak J$가 멱영이므로 $C$에서 $u(t/1)$을 나타내는 모든 원소,
  특히 $v(t)$는 $C$에서 가역이다 ($\mathbf 0_{\rm I}$, 7.1.12).
  따라서 $v$는 $B\to T^{-1}B\xrightarrow{w}C$로 인수분해된다. $v$가
  연속이므로 $w$도 연속이다 ($\mathbf 0_{\rm I}$, 7.6.6). 또한 합성
  $A\to B\to T^{-1}B\xrightarrow{w}C$가
  \[
    A\to S^{-1}A\xrightarrow{\alpha'}T^{-1}B\xrightarrow{w}C
  \]
  와 같으므로 $w$는 $S^{-1}A$-준동형이다. 즉
  $S^{-1}A\xrightarrow{\alpha'}T^{-1}B\xrightarrow{w}C$는 $C$의
  $S^{-1}A$-대수 구조를 정의하는 표준 준동형이다. 끝으로 합성
  $B\to T^{-1}B\xrightarrow{w}C\to C/\mathfrak J$과
  $B\to T^{-1}B\xrightarrow{u}C/\mathfrak J$이 같으므로, 같은 논증에
  따라 $u$는 실제로 합성
  $T^{-1}B\xrightarrow{w}C\to C/\mathfrak J$과 같다. 이로써 (iv)가
  증명된다.

  \item[(v)] 먼저 필요성은 (iv)에서 따른다. 실제로 각 $B_i$는
  $B=\prod_{i=1}^nB_i$를 그 표준 멱등원에서 국소화한 환이다. 역으로
  $C$를 이산 $A$-대수, $\mathfrak J$를 $C$의 멱영 아이디얼,
  $u:B\to C/\mathfrak J$를 연속 $A$-준동형이라 하자. $B$의 표준
  멱등원들의 $u$에 의한 상들은 합이 $1$인 유한 개의 쌍마다 직교하는
  멱등원들을 이룬다. 이 족을 $C$ 안에서 합이 $1$인 쌍마다 직교하는
  멱등원들 $e_i$로 올리자. $C_i=e_iC$와
  $\mathfrak J_i=e_i\mathfrak J$라 놓으면, $u$의 $B_i$에 대한 제한은
  연속 $A$-준동형 $B_i\to C_i/\mathfrak J_i$이다. $B_i$의 형식적
  매끄러움에 의하여 이는 연속 $A$-준동형 $B_i\to C_i$로 올라간다.
  이 올림들의 곱은 $u$를 올리는 연속 $A$-준동형 $B\to C$이다.
\end{enumerate}
\end{proof}

\oldpage[0\textsubscript{IV-1}]{83}

\begin{proposition}[19.3.6]
\label{0.19.3.6-fr}
$A$를 위상환, $B$를 위상 $A$-대수라 하고, $\widehat A$와
$\widehat B$를 각각 $A$와 $B$의 분리 완비화라 하자. 이때
$\widehat B$는 위상 $\widehat A$-대수이다. 다음 조건들은 서로
동치이다.
\begin{enumerate}
  \item[a)] $B$는 형식적으로 매끄러운 $A$-대수이다.
  \item[b)] $\widehat B$는 형식적으로 매끄러운 $A$-대수이다.
  \item[c)] $\widehat B$는 형식적으로 매끄러운 $\widehat A$-대수이다.
\end{enumerate}
\end{proposition}

\begin{proof}
$\varphi:A\to B$가 $B$의 $A$-대수 구조를 정의하는 준동형이면,
$\widehat B$의 $\widehat A$-대수 구조는 물론 준동형
$\widehat\varphi$로 정의된다. 모든 이산 $A$-대수 $C$는 분리이고
완비이므로, $A$에서 $C$로 가는 준동형의 표준 연장에 의하여
$\widehat A$-대수이기도 하다. 또한 연속 $A$-준동형 $B\to C$는 표준
연장에 의하여 $\widehat A$-준동형(따라서 특히 $A$-준동형)
$\widehat B\to C$를 준다. 다시 말해 모든 연속 $A$-준동형 $B\to C$는
유일한 방법으로 $B\to\widehat B\to C$로 인수분해된다. 이 사실들과
정의 (19.3.1)에서 명제가 곧바로 따른다.
\end{proof}

\begin{corollary}[19.3.7]
\label{0.19.3.7-fr}
(19.3.5), (iv)의 가정 아래 위상 $A\{S^{-1}\}$-대수
$B\{T^{-1}\}$은 형식적으로 매끄럽다.
\end{corollary}

이는 정의들($\mathbf 0_{\rm I}$, 7.6.1과 7.6.7), (19.3.5), (iv) 및
(19.3.6)에서 따른다.

\begin{proposition}[19.3.8]
\label{0.19.3.8-fr}
$A$를 위상환, $B$를 위상 $A$-대수라 하고, $B$의 모든 열린 아이디얼
$\mathfrak R$에 대하여 $\mathfrak R^2$도 열려 있다고 가정하자.
$A'$과 $B'$을 각각 $A$와 $B$에 원래 위상보다 더 섬세한 위상을 주어
얻은 위상환이라 하고, 표준 준동형 $\varphi:A\to B$가 여전히
$A'$에서 $B'$으로 가는 연속 준동형이라고 가정하자. 이때 $B'$이
형식적으로 매끄러운 $A'$-대수이면, $B$는 형식적으로 매끄러운
$A$-대수이다.
\end{proposition}

다음 보조정리를 적용하면 충분하다.

\begin{lemma}[19.3.8.1]
\label{0.19.3.8.1-fr}
$C$를 이산 $A$-대수, $\mathfrak J$를 $C$의 멱영 아이디얼이라 하자.
$B$의 모든 열린 아이디얼의 제곱이 열려 있다고 가정하자. 이때 합성
준동형
$v:B\xrightarrow{u}C\to C/\mathfrak J$가 연속이면, 준동형 $u$도
연속이다.
\end{lemma}

\begin{proof}
실제로 $\mathfrak R=u^{-1}(\mathfrak J)$은 $v$의 핵이므로 가정에 따라
열려 있다. 어떤 정수 $n$에 대하여 $\mathfrak J^n=0$이므로
$\mathfrak R^n\subset\operatorname{Ker}(u)$이다. 한편 $h$에 대한
귀납법과 가정에 의하여 $\mathfrak R^{2^h}$은 열려 있고, 따라서 모든
$n$에 대하여 $\mathfrak R^n$도 열려 있다. 그러므로
$\operatorname{Ker}(u)$도 열려 있으며, 이는 주장을 증명한다.
\end{proof}

$B$에 대한 가정은 $B$의 위상이 열린 아이디얼 $\mathfrak R$ 전체를
따라 얻는 $\mathfrak R$-준아딕 위상들의 상한이라는 뜻임에 유의하자.
$B$가 준아딕환이면($\mathbf 0_{\rm I}$, 7.1.9), 이 조건은 항상
충족된다.

\begin{env}[19.3.9]
\label{0.19.3.9-fr}
이제 형식적으로 매끄럽다는 성질이 정의 (19.3.1)의 조건보다 더
일반적인 조건 아래 준동형의 ``올림'' 성질들을 함의함을 보이겠다.
\end{env}

\oldpage[0\textsubscript{IV-1}]{84}

\begin{proposition}[19.3.10]
\label{0.19.3.10-fr}
$A$를 위상환, $B$를 형식적으로 매끄러운 $A$-대수라 하자. $C$를 위상
$A$-대수, $\mathfrak J$를 다음 조건들을 만족하는 $C$의 아이디얼이라
하자.
\begin{enumerate}
  \item[$1^\circ$] $C$는 거리화 가능하고 완비이다.
  \item[$2^\circ$] $\mathfrak J$는 닫혀 있고 수열
  $(\mathfrak J^n)$은 $0$으로 수렴한다.
\end{enumerate}
그러면 모든 연속 $A$-준동형 $u:B\to C/\mathfrak J$는
$B\xrightarrow{v}C\to C/\mathfrak J$로 인수분해된다. 여기서 $v$는
연속 $A$-준동형이다.
\end{proposition}

\begin{proof}
$(\mathfrak Q_n)$을 $C$ 안의 $0$의 근방들의 기본계를 이루는 감소하는
아이디얼들의 수열이라 하자. 각 $n$에 대하여 이산 $A$-대수
$C/\mathfrak Q_n$과 그 아이디얼
$(\mathfrak J+\mathfrak Q_n)/\mathfrak Q_n$을 생각하자. 어떤 $k$에
대하여 $\mathfrak J^k\subset\mathfrak Q_n$이므로
$(\mathfrak J+\mathfrak Q_n)/\mathfrak Q_n$은 멱영이다. 각 $n$에
대하여 $u_n$을 다음 연속 준동형이라 하자.
\[
  B\xrightarrow{u}C/\mathfrak J\to C/(\mathfrak J+\mathfrak Q_n)
  =(C/\mathfrak Q_n)/((\mathfrak J+\mathfrak Q_n)/\mathfrak Q_n).
\]
가정에 따라 $u_n$은
\[
  B\xrightarrow{v_n}C/\mathfrak Q_n
  \xrightarrow{\varphi_n}C/(\mathfrak J+\mathfrak Q_n)
\]
으로 인수분해된다. 여기서 $v_n$은 연속 $A$-준동형이다. 이제 모든
$n$에 대하여 $v_n$이
\[
  B\xrightarrow{v_{n+1}}C/\mathfrak Q_{n+1}\to C/\mathfrak Q_n
\]
으로 인수분해되도록 $v_n$들을 차례로 선택할 수 있음을 보이자. 다시
말해 $v_n$들이 준동형들의 \emph{사영계}를 이루도록 하겠다. 이는 다음
보조정리에서 따른다.
\end{proof}

\begin{lemma}[19.3.10.1]
\label{0.19.3.10.1-fr}
$B$를 형식적으로 매끄러운 $A$-대수, $E$와 $E'$을 두 이산 위상
$A$-대수, $\mathfrak R$과 $\mathfrak R'$을 각각 $E$와 $E'$의 멱영
아이디얼이라 하자. $f:E\to E'$을
$f(\mathfrak R)\subset\mathfrak R'$을 만족하는 전사
$A$-준동형이라 하고,
$f':E/\mathfrak R\to E'/\mathfrak R'$을 몫으로 내려가 $f$에서
유도되는 준동형이라 하자. $u:B\to E/\mathfrak R$을 연속
$A$-준동형, $u'$을 합성
$B\xrightarrow{u}E/\mathfrak R\xrightarrow{f'}E'/\mathfrak R'$이라
하고, $u'$이 $B\xrightarrow{v'}E'\to E'/\mathfrak R'$으로
인수분해되게 하는 연속 $A$-준동형 $v':B\to E'$이 주어졌다고 하자.
그러면 $v'$이 $B\xrightarrow{v}E\xrightarrow{f}E'$으로
인수분해되게 하는 연속 $A$-준동형 $v:B\to E$가 존재한다.
\end{lemma}

가환도식
\[
\xymatrix{
  E \ar[r]^f \ar[d] & E' \ar[d] \\
  E/\mathfrak R \ar[r]_{f'} & E'/\mathfrak R'
}
\]
의 모든 준동형은 전사이다. $\mathfrak L$이 $f$의 핵이면, $E'$은
$E/\mathfrak L$과 동일시되고 $\mathfrak R'$은
$(\mathfrak R+\mathfrak L)/\mathfrak L$과 동일시된다. 이제 다음의
기초적인 보조정리를 사용하자.

\begin{lemma}[19.3.10.2]
\label{0.19.3.10.2-fr}
$F$를 반드시 가환일 필요는 없는 환, $\mathfrak a$와 $\mathfrak b$를
$F$의 두 양쪽 아이디얼이라 하자. 그러면 섬유곱
\[
  (F/\mathfrak a)\times_{F/(\mathfrak a+\mathfrak b)}(F/\mathfrak b)
\]
은 $F/(\mathfrak a\cap\mathfrak b)$와 표준적으로 동일시된다.
\end{lemma}

\begin{proof}
이는 (18.1.7)의 특수한 적용에 지나지 않는다. 즉 도식 (18.1.7.1)에서
$G$를 $F/(\mathfrak a\cap\mathfrak b)$로, $\mathfrak J'$을
$\mathfrak b/(\mathfrak a\cap\mathfrak b)$로, 그 도식의 $F$를
$F/\mathfrak b$로, $\mathfrak R$을
$(\mathfrak a+\mathfrak b)/\mathfrak b$로, $B$를
$F/(\mathfrak a+\mathfrak b)$로, $\theta$를 표준 전단사
$F$-쌍가군 준동형
$(\mathfrak a+\mathfrak b)/\mathfrak b
\to\mathfrak a/(\mathfrak a\cap\mathfrak b)$으로 바꾸어 적용한다.
\end{proof}

\oldpage[0\textsubscript{IV-1}]{85}

\begin{proof}[보조정리~{\rm(19.3.10.1)}의 증명 계속]
이 보조정리를 (19.3.10.1)의 상황에 적용하면, 도식
\[
\xymatrix{
  B \ar[r]^{v'} \ar[d]_u & E'=E/\mathfrak L \ar[d] \\
  E/\mathfrak R \ar[r]_{f'} & E'/\mathfrak R'
}
\]
의 가환성으로부터 $v'$과 $u$가 각각
$B\xrightarrow{w}E/(\mathfrak R\cap\mathfrak L)\to E/\mathfrak L$과
$B\xrightarrow{w}E/(\mathfrak R\cap\mathfrak L)\to E/\mathfrak R$로
인수분해되게 하는 준동형
$w:B\to E/(\mathfrak R\cap\mathfrak L)$이 존재한다. $w$의 핵은
$v'$의 핵과 $u$의 핵의 교집합을 포함하므로 $B$에서 열려 있고,
따라서 $w$는 연속이다. 끝으로 $\mathfrak R\cap\mathfrak L$이
멱영임은 명백하므로, $w$는
$B\xrightarrow{v}E\to E/(\mathfrak R\cap\mathfrak L)$로
인수분해된다. 여기서 $v$는 요구한 연속 $A$-준동형이다.
\end{proof}

\begin{proof}[명제~{\rm(19.3.10)}의 증명 계속]
보조정리 (19.3.10.1)에서 $E$, $E'$, $\mathfrak R$,
$\mathfrak R'$을 각각 $C/\mathfrak Q_{n+1}$, $C/\mathfrak Q_n$,
$(\mathfrak J+\mathfrak Q_{n+1})/\mathfrak Q_{n+1}$,
$(\mathfrak J+\mathfrak Q_n)/\mathfrak Q_n$으로 바꾸고, $u$, $u'$,
$v'$을 각각 $u_{n+1}$, $u_n$, $v_n$으로 바꾸어 적용하면
$v_{n+1}$의 존재를 얻는다. $C$는 분리이고 완비이므로
$C=\varprojlim_n(C/\mathfrak Q_n)$이고,
$v=\varprojlim_n v_n$은 $B$에서 $C$로 가는 연속 $A$-준동형임에
유의하자. $\mathfrak J$가 $C$에서 닫혀 있으므로
$\mathfrak J=\bigcap_n(\mathfrak J+\mathfrak Q_n)$이다. 또한
$C/\mathfrak J$는 거리화 가능하고 완비이며
$(\mathfrak J+\mathfrak Q_n)/\mathfrak J$들이 $C/\mathfrak J$ 안의
$0$의 근방들의 기본계를 이루므로
\[
  C/\mathfrak J=\varprojlim_n C/(\mathfrak J+\mathfrak Q_n).
\]
따라서 $\varprojlim_n\varphi_n$은 표준 사상
$\varphi:C\to C/\mathfrak J$이고,
$\varprojlim_n u_n=u$이므로 $u=\varphi\circ v$이다. 증명 끝.
\end{proof}

\begin{corollary}[19.3.11]
\label{0.19.3.11-fr}
$A$를 위상환, $B$를 형식적으로 매끄러운 $A$-대수, $C$를 완비 뇌터
국소환, $\mathfrak J$를 $C$의 진 아이디얼이라 하자. 또한
$\varphi:A\to C$를 $C$에 위상 $A$-대수 구조를 주는 연속 준동형이라
하자. 그러면 모든 연속 $A$-준동형
$u_0:B\to C_0=C/\mathfrak J$은
$B\xrightarrow{u}C\to C/\mathfrak J$로 인수분해된다. 여기서 $u$는
연속 $A$-준동형이다.
\end{corollary}

실제로 (19.3.10)의 모든 조건이 충족된다
($\mathbf 0_{\rm I}$, 7.3.5).

\begin{proposition}[19.3.12]
\label{0.19.3.12-fr}
$A$를 위상환, $B$를 형식적으로 매끄러운 $A$-대수, $C$를 위상
$A$-대수라 하고, $\mathfrak J$를 다음 조건들을 만족하는 $C$의
아이디얼이라 하자.
\begin{enumerate}
  \item[$1^\circ$] $C$에는 열린 아이디얼들의 기본계
  $(\mathfrak Q_\lambda)$가 존재하여, 각
  $C_\lambda=C/\mathfrak Q_\lambda$는 아르틴환이고 표준 준동형
  $C\to\varprojlim_\lambda C_\lambda$는 위상환의 동형사상이다.
  \item[$2^\circ$] $\mathfrak J$는 $C$에서 닫혀 있고 위상적으로
  멱영이다.
  \item[$3^\circ$] $B$의 모든 열린 아이디얼의 제곱은 열려 있다.
\end{enumerate}
이 조건들 아래 모든 연속 $A$-준동형 $u:B\to C/\mathfrak J$은
$B\xrightarrow{v}C\to C/\mathfrak J$로 인수분해된다. 여기서 $v$는
연속 $A$-준동형이다.
\end{proposition}

\begin{proof}
$\mathfrak J_\lambda=(\mathfrak J+\mathfrak Q_\lambda)/
\mathfrak Q_\lambda$를 $C_\lambda$에서의 $\mathfrak J$의 표준 상이라
하자. 이는 멱영 아이디얼이다. 또한 $u_\lambda$를 합성 준동형
$B\xrightarrow{u}C/\mathfrak J\to C_\lambda/\mathfrak J_\lambda
=C/(\mathfrak J+\mathfrak Q_\lambda)$라 하자. $u$가
$B\xrightarrow{w}C\to C/\mathfrak J$으로 인수분해되게 하는 임의의
$A$-준동형 $w:B\to C$는 반드시 연속임을 먼저 보이자. 실제로 각
$\lambda$에 대하여 $B$의 열린 아이디얼 $\mathfrak R$이 존재하여
$u(\mathfrak R)\subset
(\mathfrak J+\mathfrak Q_\lambda)/\mathfrak J$이고, 따라서
$w(\mathfrak R)\subset\mathfrak J+
\mathfrak Q_\lambda$이다. 그러므로
$w(\mathfrak R^{2^m})\subset\mathfrak Q_\lambda$가 되는
$\mathfrak R$의 거듭제곱 $\mathfrak R^{2^m}$이 존재한다.
$\mathfrak R^{2^m}$이 $B$에서 열려 있으므로 이로써 주장이
증명된다. 따라서 연속성을 따로 염려하지 않고 위 인수분해 성질을
갖는 $A$-준동형 $w$를 찾으면 충분하다.

모든 $A$-대수 준동형들의 집합에 대하여
$\operatorname{Hom}_A(B,C)=
\varprojlim_\lambda\operatorname{Hom}_A(B,C_\lambda)$이다. 각
$\operatorname{Hom}_A(B,C_\lambda)$는 곱위상을 준
$C_\lambda$-가군 $C_\lambda^B$ 안에서 닫혀 있고,
$C_\lambda$가 아르틴환이므로 이 위상에 관하여 선형 콤팩트이다. 각
$\lambda$에 대하여 $W_\lambda$를

\oldpage[0\textsubscript{IV-1}]{86}

$u_\lambda$가
$B\xrightarrow{w_\lambda}C_\lambda\to
C_\lambda/\mathfrak J_\lambda$로 인수분해되게 하는
$w_\lambda\in\operatorname{Hom}_A(B,C_\lambda)$들의 집합이라 하자.
$B$가 형식적으로 매끄러우므로 $W_\lambda$는 비어 있지 않으며,
$C_\lambda$-가군 $\operatorname{Hom}_A(B,C_\lambda)$ 안의 닫힌 아핀
부분공간이다. 곱
$E_\lambda=\prod_{\mu\leq\lambda}
\operatorname{Hom}_A(B,C_\mu)$ 안에서 $\mu\leq\lambda$일 때
$w_\mu=\varphi_{\mu\lambda}\circ w_\lambda$이고
$w_\lambda\in W_\lambda$인 족
$(w_\mu)_{\mu\leq\lambda}$들로 이루어진 닫힌 아핀 부분공간을
$U_\lambda$라 하자. 여기서
$\varphi_{\mu\lambda}:C_\lambda\to C_\mu$는 표준 준동형이다. 끝으로
$\prod_\mu\operatorname{Hom}_A(B,C_\mu)$ 안에서 $U_\lambda$와
$\mu\not\leq\lambda$에 대한
$\operatorname{Hom}_A(B,C_\mu)$들의 곱인 닫힌 아핀 부분공간을
$V_\lambda$라 하자. 이제 $V_\lambda$들의 교집합이 비어 있지 않음을
보이면 된다. 실제로 이 교집합의 원소 $w$는 정의상
$\varprojlim_\lambda\operatorname{Hom}_A(B,C_\lambda)$에 속한다.
\footnote{인쇄된 증명에는 이 지점에 빈틈이 있다. 대수 준동형들의
집합은 점별 덧셈 아래 가군이 아니며, 일반 멱영 아이디얼에 대한
올림의 섬유는 아핀 부분공간일 필요가 없다. 따라서 제시된 선형
콤팩트성 논증만으로 이 교집합이 비어 있지 않음은 증명되지 않는다.
별도의 양립하는 올림 논증이 필요하다.}

또한 $\mathfrak J$가 닫혀 있고
$C=\varprojlim_\lambda C_\lambda$가 선형 콤팩트이므로
$C/\mathfrak J$는
$\varprojlim_\lambda C_\lambda/\mathfrak J_\lambda$와 동일시되고,
표준 사상 $\psi:C\to C/\mathfrak J$는
$\varprojlim_\lambda\psi_\lambda$와 동일시된다. 여기서
$\psi_\lambda:C_\lambda\to C_\lambda/\mathfrak J_\lambda$는 표준
사상이다. 따라서 (19.3.10)에서와 같이 $\psi\circ w=u$를 얻는다.
\end{proof}

\subsection{형식적 매끄러움의 첫 판정법들}
\label{subsection:0.19.4-fr}

\begin{proposition}[19.4.1]
\label{0.19.4.1-fr}
$A$를 위상환, $B$를 위상 $A$-대수라 하자. 각각 $A$와 $B$의
아이디얼들로 이루어진 두 감소하는 필터족
$(\mathfrak J_\alpha)_{\alpha\in I}$와
$(\mathfrak R_\alpha)_{\alpha\in I}$가 존재하여 다음 조건들을
만족한다고 하자.
\begin{enumerate}
  \item[$1^\circ$] $(\mathfrak J_\alpha)$는 $A$에서 $0$으로 수렴하고
  $(\mathfrak R_\alpha)$는 $B$에서 $0$으로 수렴한다.
  \item[$2^\circ$] 모든 $\alpha\in I$에 대하여
  $\mathfrak J_\alpha B\subset\mathfrak R_\alpha$이다. 따라서
  $B/\mathfrak R_\alpha$는 위상
  $(A/\mathfrak J_\alpha)$-대수이다.
  \item[$3^\circ$] 모든 $\alpha\in I$에 대하여
  $B/\mathfrak R_\alpha$는 형식적으로 매끄러운
  $(A/\mathfrak J_\alpha)$-대수이다.
\end{enumerate}
그러면 $B$는 형식적으로 매끄러운 $A$-대수이다.
\end{proposition}

\begin{proof}
$C$를 이산 $A$-대수, $\mathfrak L$을 $C$의 멱영 아이디얼,
$\mathfrak R$을 $B$의 열린 아이디얼이라 하자. 가정에 따라
$\mathfrak R_\alpha\subset\mathfrak R$인 $\alpha\in I$가 존재하므로
$B/\mathfrak R$은 $B/\mathfrak R_\alpha$의 몫이다. 모든
$A$-준동형 $u':B/\mathfrak R\to C/\mathfrak L$은 또한
$(A/\mathfrak J_\alpha)$-준동형이다. 따라서 $B$의 열린 아이디얼
$\mathfrak R'$이 존재하여
$\mathfrak R_\alpha\subset\mathfrak R'\subset\mathfrak R$이고, 합성
준동형
\[
  B/\mathfrak R_\alpha\longrightarrow B/\mathfrak R
  \xrightarrow{u'}C/\mathfrak L
\]
은
\[
  B/\mathfrak R_\alpha\longrightarrow B/\mathfrak R'
  \longrightarrow C\longrightarrow C/\mathfrak L
\]
로 인수분해된다. 여기서 가운데 준동형
$B/\mathfrak R'\to C$는
$(A/\mathfrak J_\alpha)$-준동형이다. 이제 (19.3.1) 뒤의 주의에
따라 결론을 얻는다.
\end{proof}

\begin{corollary}[19.4.2]
\label{0.19.4.2-fr}
$A$를 위상환, $B$를 위상 $A$-대수라 하고,
$(\mathfrak J_\alpha)_{\alpha\in I}$를 $A$에서 $0$으로 수렴하는
아이디얼들의 감소하는 필터족이라 하자. $B$가 형식적으로 매끄러운
$A$-대수이기 위한 필요충분조건은, 모든 $\alpha\in I$에 대하여
\[
  A_\alpha=A/\mathfrak J_\alpha,\qquad
  B_\alpha=B/\mathfrak J_\alpha B
\]
로 놓을 때 $B_\alpha$가 형식적으로 매끄러운
$A_\alpha$-대수인 것이다.
\end{corollary}

이 조건은 (19.4.1)에 따라 충분하고, (19.3.5), (iii)에 따라
필요하기도 하다.

\begin{proposition}[19.4.3]
\label{0.19.4.3-fr}
$A$를 위상환, $B$를 위상 $A$-대수라 하자. 모든 이산 $A$-대수
$C$와 $\mathfrak J^2=0$인 $C$의 모든 아이디얼 $\mathfrak J$에
대하여, 모든 연속 $A$-준동형 $u:B\to C/\mathfrak J$이
\[
  B\xrightarrow{v}C\longrightarrow C/\mathfrak J
\]
로 인수분해된다고 하자. 여기서 $v$는 연속 $A$-준동형이다. 그러면
$B$는 형식적으로 매끄러운 $A$-대수이다.
\end{proposition}

\begin{proof}
같은 표기 아래 $\mathfrak R$을 $C$의 임의의 멱영 아이디얼이라
하고, 연속 $A$-준동형 $u':B\to C/\mathfrak R$을 생각하자.
$\mathfrak R^m=0$이라 하고, $1\leq i\leq m$에 대하여
$C_i=C/\mathfrak R^i$로 놓자. 그러면 $C_m=C$이고,
$\mathfrak R^{i-1}/\mathfrak R^i$는 $C_i$의 제곱이 영인 아이디얼
$\mathfrak J_i$이며,
$C_i=C_{i+1}/\mathfrak J_{i+1}$이다. 이제 가정을 $i$에 관하여
귀납적으로 적용하면, $u_1=u'$이고 각 $u_i$가
\[
  B\xrightarrow{u_{i+1}}C_{i+1}\longrightarrow
  C_{i+1}/\mathfrak J_{i+1}=C_i
\]
로 인수분해되게 하는 연속 $A$-준동형 $u_i:B\to C_i$들이
존재한다. 따라서 $u_m:B\to C$는 $u'$의 연속 올림이며, 결론을
얻는다.
\end{proof}

\begin{proposition}[19.4.4]
\label{0.19.4.4-fr}
$A$를 위상환, $B$를 가환 위상 $A$-대수라 하자. $B$가 형식적으로
매끄러운 $A$-대수이기 위한 필요충분조건은, $B$의 열린 아이디얼에
의하여 소멸되는 모든 이산 위상 $B$-가군 $L$에 대하여
(18.5.1 참조)

\oldpage[0\textsubscript{IV-1}]{87}

\[
  \operatorname{Exalcotop}_A(B,L)=0
\]
인 것이다.
\end{proposition}

\begin{proof}
$(\mathfrak R_\lambda)$를 $B$의 열린 아이디얼들의 감소하는
기본계라 하고 $B_\lambda=B/\mathfrak R_\lambda$로 놓자. 이산
위상 $A$-대수 $C$와 $\mathfrak J^2=0$인 $C$의 아이디얼
$\mathfrak J$를 생각하자. 그러면 $C$는 $\mathfrak J$에 의한
$C/\mathfrak J$의 $A$-확장이다. $A$-준동형
$u:B_\lambda\to C/\mathfrak J$이 주어졌다고 하고, $u$에 의한
$C$의 역상인, $\mathfrak J$에 의한 $B_\lambda$의 $A$-확장
$E_\lambda$를 만들자 (18.2.5). 이는 이산 위상을 준 위상
$A$-대수이다.

가정한 $\operatorname{Exalcotop}$ 소멸 조건을
$B_\lambda$-가군 $\mathfrak J$에 적용하면, 정의 (18.5.1)에 따라
$\mathfrak R_\mu\subset\mathfrak R_\lambda$이고 역상 확장
$E_\mu$가 $A$-자명인 $\mu$가 존재한다. 이는 (18.1.6)에 따라
표준 준동형 $B_\mu\to B_\lambda$가
\[
  B_\mu\xrightarrow{v}E_\lambda\longrightarrow B_\lambda
\]
로 인수분해되게 하는 연속 $A$-준동형
$v:B_\mu\to E_\lambda$가 존재한다는 뜻이다. 따라서
$B_\mu\to B_\lambda\xrightarrow{u}C/\mathfrak J$은 곧
\[
  B_\mu\xrightarrow{v}E_\lambda\longrightarrow C
  \longrightarrow C/\mathfrak J
\]
로 인수분해된다. (19.3.1)과 (19.4.3)에 따라 이로부터 $B$가
형식적으로 매끄러운 $A$-대수임을 얻는다.

역은 즉시 따른다. 실제로 $C$가 $L$에 의한 $B_\lambda$의
$A$-확장인 위상 $A$-대수인 경우와 항등 준동형
$B_\lambda\to B_\lambda=C/L$에 (19.3.1)을 적용하면 된다.
\end{proof}

$A$와 $B$가 이산환이면 형식적 매끄러움의 판정법 (19.4.4)는
다음으로 환원된다.
\begin{equation}
\tag{19.4.4.1}
\label{0.19.4.4.1-fr}
\operatorname{Exalcom}_A(B,L)=0
\qquad\text{모든 $B$-가군 $L$에 대하여}.
\end{equation}
달리 말해, \emph{$B$-가군에 의한 $B$의 모든 가환 $A$-확장은
$A$-자명해야 한다}.

\begin{corollary}[19.4.5]
\label{0.19.4.5-fr}
$A$를 위상환, $B$를 가환 위상 $A$-대수라 하자.
\begin{enumerate}
  \item[(i)] $B$가 형식적으로 매끄러운 $A$-대수라고 하자. 그러면
  $B$의 모든 열린 아이디얼 $\mathfrak R$, 모든
  $(B/\mathfrak R)$-가군 $L$, 그리고
  $(B/\mathfrak R)\times(B/\mathfrak R)$에서 $L$로 가는 모든
  대칭인 $A$-쌍선형 호흐실트 $2$-공순환 $f$에 대하여 (18.4.3),
  열린 아이디얼 $\mathfrak R'\subset\mathfrak R$이 존재하여
  다음이 성립한다. $\varphi:B/\mathfrak R'\to B/\mathfrak R$를
  표준 준동형이라 하면
  $f\circ(\varphi\times\varphi)$는
  $(B/\mathfrak R')\times(B/\mathfrak R')$에서 $L$로 가는
  $A$-쌍선형 호흐실트 $2$-공경계이다.
  \item[(ii)] $B$가 형식적으로 사영인 $A$-가군이고 (19.2.1),
  (i)의 조건이 성립하면 $B$는 형식적으로 매끄러운 $A$-대수이다.
\end{enumerate}
\end{corollary}

\begin{proof}
\begin{enumerate}
  \item[(i)] $2$-공순환 $f$는 $L$에 의한 $B/\mathfrak R$의
  호흐실트 $A$-확장 $C$를 정의한다 (18.4.3). $C$와 $C$의 제곱이
  영인 아이디얼 $L$, 그리고 항등 준동형
  $B/\mathfrak R\to B/\mathfrak R=C/L$에 (19.3.1)을 적용하면,
  (18.4.3)에 따라 (i)의 조건을 얻는다.

  \item[(ii)] 판정법 (19.4.3)을 적용하자. $B$의 열린 아이디얼
  $\mathfrak R$, $\mathfrak JB\subset\mathfrak R$인 $A$의 열린
  아이디얼 $\mathfrak J$, 그리고 $L$에 의한
  $B/\mathfrak R$의 $(A/\mathfrak J)$-확장 $E$를 생각하자.
  $B$가 형식적으로 사영인 $A$-가군이므로, 표준 연속 $A$-선형
  사상 $\rho:B\to B/\mathfrak R$은
  \[
    B\xrightarrow{w}E\longrightarrow B/\mathfrak R
  \]
  로 인수분해된다. 여기서 $w$는 연속 $A$-선형 사상이다 (19.2.1).
  따라서 $w$ 자체도
  \[
    B\longrightarrow B/\mathfrak R'\xrightarrow{w'}E
  \]
  로 인수분해된다. 여기서 $\mathfrak R'$은
  $\mathfrak R'\subset\mathfrak R$인 $B$의 열린 아이디얼이다.
  $\mathfrak J$를 더 작은 아이디얼로 바꾸어
  $\mathfrak JB\subset\mathfrak R'$이라고 가정할 수 있다. 이제
  표준 준동형 $B/\mathfrak R'\to B/\mathfrak R$에 의한,
  $L$에 의한 $B/\mathfrak R$의 확장 $E$의 역상은 $L$에 의한
  $B/\mathfrak R'$의 $(A/\mathfrak J)$-\emph{호흐실트 확장}
  $E'$이다.

  \oldpage[0\textsubscript{IV-1}]{88}

  이 확장을 정의하는 공순환에 (i)를 적용하면 (18.4.3), $B$의
  열린 아이디얼 $\mathfrak R''\subset\mathfrak R'$이 존재하여
  $B/\mathfrak R''\to B/\mathfrak R$에 의한 $E$의 역상
  $E''$이 $A$-자명임을 얻는다. 증명 끝.
\end{enumerate}
\end{proof}

\begin{corollary}[19.4.6]
\label{0.19.4.6-fr}
$A$를 위상환, $B$를 형식적으로 사영인 $A$-가군인 위상
$A$-대수라 하자. $A'$을 $A$의 위상에서 유도된 위상을 준
$A$-대수라 하자. 또한 $A'$이 충실 평탄 $A$-가군이고 다음 조건들
가운데 하나가 성립한다고 하자.
\begin{enumerate}
  \item[$1^\circ$] $A$의 열린 아이디얼들의 기본계
  $(\mathfrak J_\lambda)$와 $B$의 열린 아이디얼들의 기본계
  $(\mathfrak M_\lambda)$가 존재하고, 두 기본계는 같은 지표집합을
  가지며, 모든
  $\lambda$에 대하여
  $\mathfrak J_\lambda B\subset\mathfrak M_\lambda$이고
  $B/\mathfrak M_\lambda$가 유한형 사영
  $(A/\mathfrak J_\lambda)$-가군이다.
  \item[$2^\circ$] $A'$은 유한형 사영 $A$-가군이다.
\end{enumerate}
그러면 텐서곱 위상을 준
$B'=B\otimes_A A'$이 형식적으로 매끄러운 $A'$-대수이기 위한
필요충분조건은 $B$가 형식적으로 매끄러운 $A$-대수인 것이다.
\end{corollary}

\begin{proof}
이 조건의 충분성은 $B$나 $A'$에 아무 추가 가정도 두지 않은
(19.3.5), (iii)에 들어 있다. 역을 증명하기 위하여 판정법
(19.4.5)를 적용하자. 조건 $2^\circ$ 아래에서는
$(\mathfrak M_\lambda)$를 다시 $B$의 열린 아이디얼들의 기본계라
하고, 각 $\lambda$에 대하여
$\mathfrak J_\lambda B\subset\mathfrak M_\lambda$인 $A$의
열린 아이디얼 $\mathfrak J_\lambda$를 택한다. 두 경우 모두
\[
  A_\lambda=A/\mathfrak J_\lambda,\qquad
  B_\lambda=B/\mathfrak M_\lambda,\qquad
  A'_\lambda=A_\lambda\otimes_A A',\qquad
  B'_\lambda=B_\lambda\otimes_A A'
\]
로 놓자. $f$를 $B_\lambda$-가군 $L$에 값을 갖는,
$B_\lambda\times B_\lambda$의 대칭인
$A_\lambda$-쌍선형 호흐실트 $2$-공순환이라 하자. 스칼라를
확장하면 $L'=L\otimes_A A'$에 값을 갖는,
$B'_\lambda\times B'_\lambda$의 대칭인
$A'_\lambda$-쌍선형 호흐실트 $2$-공순환
$f'=f\otimes1$을 얻는다.

가정에 따라 $B'$은 형식적으로 매끄러운 $A'$-대수이므로,
(19.4.5), (i)에 따라
$\mathfrak M_\mu\subset\mathfrak M_\lambda$이고 다음 조건을
만족하는 지표 $\mu$가 존재한다. 표준 사상을
$\varphi':B'_\mu\to B'_\lambda$라 하면
$f'\circ(\varphi'\times\varphi')$는
$B'_\mu\times B'_\mu$에서 $L'$에 값을 갖는 호흐실트
$2$-공경계이다. 달리 말해, 호흐실트 군
\[
  \mathrm H^2_{A'_\mu}(B'_\mu,L')^s
\]
에서 그 상 $c'$은 영이다. 이제
$\varphi:B_\mu\to B_\lambda$를 표준 준동형이라 하고,
$c$를 호흐실트 군
$\mathrm H^2_{A_\mu}(B_\mu,L)^s$에서
$f\circ(\varphi\times\varphi)$의 류라 하면, $c'$은 $c$의 표준
상임이 명백하다.

$P_\bullet$을 (18.4.5)에서 정의한 $A_\mu$와 $B_\mu$에 대한
복합체, 곧 $\mathrm H^2_{A_\mu}(B_\mu,L)^s$의 계산에 쓰이는
복합체라 하자. $A'_\mu$와 $B'_\mu$에 대한 같은 복합체는
명백히 $P_\bullet\otimes_A A'$이다. 조건 $1^\circ$ 아래에서는
$P_\bullet$의 구성으로부터 그 각 항이 유한형 사영
$A_\mu$-가군임을 알 수 있다. 따라서 두 조건 어느 것 아래에서도
Bourbaki, \emph{Alg.}, chap.~II, 3\textsuperscript{e} éd., §~5,
n\textsuperscript{o}~3, prop.~7에 따라 표준 동형사상 아래
\[
  \operatorname{Hom}_{A'_\mu}(P_\bullet\otimes_A A',
  L\otimes_A A')
  =
  \operatorname{Hom}_{A_\mu}(P_\bullet,L)\otimes_A A'
\]
을 얻는다. $A'$은 평탄 $A$-가군이므로 (18.4.5)에 따라
\[
  \mathrm H^2_{A'_\mu}(B'_\mu,L')^s
  =
  \mathrm H^2_{A_\mu}(B_\mu,L)^s\otimes_A A'
\]
이고, 따라서 $c'=c\otimes1$이라 쓸 수 있다. 그런데 $A'$은 충실
평탄 $A$-가군이므로 $c'=0$에서 $c=0$이 따른다. 이로써 증명이
끝난다.
\end{proof}

\begin{proposition}[19.4.7]
\label{0.19.4.7-fr}
$A$를 준허용 위상환, $\mathfrak J$를 $A$의 정의 아이디얼
($\mathbf 0_{\rm I}$, 7.1.2), $B$를 형식적으로 사영인
$A$-가군인 위상 $A$-대수라 하자. 몫위상환
\[
  A_0=A/\mathfrak J,\qquad B_0=B/\mathfrak JB
\]
를 생각하자. $B$가 형식적으로 매끄러운 $A$-대수이기 위한
필요충분조건은 $B_0$가 형식적으로 매끄러운 $A_0$-대수인 것이다.
\end{proposition}

\oldpage[0\textsubscript{IV-1}]{89}

\begin{proof}
조건의 필요성은 (19.4.2)에서 따른다. 충분성을 보이기 위하여 먼저
$\mathfrak J$에 포함되는 아이디얼들
$\mathfrak J_\alpha$로 이루어진 $A$의 $0$의 열린 근방 기본계를
생각하자. $B/\mathfrak J_\alpha B$가 형식적으로 사영인
$(A/\mathfrak J_\alpha)$-가군이므로 (19.2.5), (19.4.2)에 따라
$A$가 이산인 경우로 환원할 수 있다. 그러면 준허용환의 정의에
따라 $\mathfrak J$는 \emph{멱영}이다.

또한 $\mathfrak J^2=0$인 경우에 명제를 증명하는 것으로 충분하다.
실제로 $\mathfrak J^m=0$이면
\[
  A_k=A/\mathfrak J^k,\qquad B_k=B/\mathfrak J^kB
\]
와 $A_k$의 아이디얼
$\mathfrak J^{k-1}/\mathfrak J^k$에 대하여 $k=2,3,\ldots,m$의
순서로 명제를 적용한다. 이때
$B_k=B\otimes_A A_k$가 형식적으로 사영인 $A_k$-가군임에
유의한다 (19.2.5).

이제 $\mathfrak J^2=0$이라 하고 판정법 (19.4.5), (ii)를
적용하자. $B$의 열린 아이디얼 $\mathfrak R$과
$(B/\mathfrak R)$-가군 $L$에 값을 갖는,
$(B/\mathfrak R)\times(B/\mathfrak R)$의 대칭인
$A$-쌍선형 호흐실트 $2$-공순환 $f$를 생각한다.

먼저 $\mathfrak JL=0$인 특별한 경우를 보자. 이때 $L$은
$(B_0/\mathfrak R B_0)$-가군으로도 볼 수 있고, $f$는
\[
  (B/\mathfrak R)\times(B/\mathfrak R)
  \longrightarrow
  (B_0/\mathfrak R B_0)\times(B_0/\mathfrak R B_0)
  \xrightarrow{f_0}L
\]
로 인수분해된다. 여기서 $f_0$는 대칭인 쌍선형 호흐실트
$2$-공순환이다. 가정에 따라 $B$의 열린 아이디얼
$\mathfrak R'\subset\mathfrak R$과 $A_0$-선형 사상
\[
  g_0:B_0/\mathfrak R'B_0\longrightarrow L
\]
이 존재하여, $B_0/\mathfrak R'B_0$의 $x,y$에 대하여
\[
  f_0(\varphi_0(x),\varphi_0(y))
  =xg_0(y)-g_0(xy)+g_0(x)y
\]
이다. 여기서 $\varphi_0$는 표준 준동형
$B_0/\mathfrak R'B_0\to B_0/\mathfrak R B_0$이다. 따라서 합성
$A$-선형 사상
\[
  g:B/\mathfrak R'\longrightarrow
  B_0/\mathfrak R'B_0\xrightarrow{g_0}L
\]
은 표준 준동형
$\varphi:B/\mathfrak R'\to B/\mathfrak R$에 대하여
$dg=f\circ(\varphi\times\varphi)$를 만족한다.

이제 일반적인 경우로 가자. 먼저 $(B/\mathfrak R)$-가군
$L'=L/\mathfrak JL$을 생각하면 $\mathfrak JL'=0$이다. 합성
$A$-쌍선형 사상
\[
  (B/\mathfrak R)\times(B/\mathfrak R)
  \xrightarrow{f}L\longrightarrow L'
\]
을 $f'$이라 하면, $f'$도 대칭인 호흐실트 $2$-공순환이다. 앞에서
보인 바에 따라 $B$의 열린 아이디얼
$\mathfrak R'\subset\mathfrak R$과 $A$-선형 사상
$g':B/\mathfrak R'\to L'$이 존재하여, 표준 사상
$\varphi':B/\mathfrak R'\to B/\mathfrak R$에 대하여
\[
  dg'=f'\circ(\varphi'\times\varphi')
\]
이다. $B$가 형식적으로 사영인 $A$-가군이므로, 열린 아이디얼
$\mathfrak R'_1\subset\mathfrak R'$과 $A$-선형 사상
$g_1:B/\mathfrak R'_1\to L$이 존재하여, 합성
\[
  B/\mathfrak R'_1\longrightarrow
  B/\mathfrak R'\xrightarrow{g'}L'
\]
이
$B/\mathfrak R'_1\xrightarrow{g_1}L\to L'$로 인수분해된다.
표준 준동형을
$\varphi_1:B/\mathfrak R'_1\to B/\mathfrak R$라 하고,
호흐실트 $2$-공순환
\[
  f_1(x,y)=f(\varphi_1(x),\varphi_1(y))
  -xg_1(y)+g_1(xy)-g_1(x)y
\]
를 생각하자. 이는
$(B/\mathfrak R'_1)\times(B/\mathfrak R'_1)$에서 $L$로 가는
대칭인 $A$-쌍선형 사상이다. $g_1$의 선택에 따라 $f_1$의 값은
$\mathfrak JL$에 속한다. $\mathfrak J(\mathfrak JL)=0$이므로
첫 번째 경우를 다시 적용할 수 있다. 따라서 열린 아이디얼
$\mathfrak R''\subset\mathfrak R'_1$과 $A$-선형 사상
$g_2:B/\mathfrak R''\to\mathfrak JL$이 존재하여
\[
  f_1(\varphi_2(x),\varphi_2(y))
  =xg_2(y)-g_2(xy)+g_2(x)y
\]
이다. 여기서
$\varphi_2:B/\mathfrak R''\to B/\mathfrak R'_1$은 표준
사상이다. 그러므로 $A$-선형 사상
\[
  g=g_2+g_1\circ\varphi_2:B/\mathfrak R''\longrightarrow L
\]
은 표준 사상
$\varphi:B/\mathfrak R''\to B/\mathfrak R$에 대하여
$dg=f\circ(\varphi\times\varphi)$를 만족한다. 증명 끝.
\end{proof}

\subsection{형식적 매끄러움과 연관 등급환}
\label{subsection:0.19.5-fr}

\begin{env}[19.5.1]
\label{0.19.5.1-fr}
\oldpage[0\textsubscript{IV-1}]{90}
$C$를 가환 위상환, $V$를 위상 $C$-가군이라 하고 대칭 대수
\[
  \mathbf S_C^*(V)=\bigoplus_{n\geq0}\mathbf S_C^n(V)
\]
를 생각하자. 이제 이 대수에 $C$-대수 구조와 양립하는
\emph{선형 위상}을 표준적으로 주겠다. 이를 위하여 $V$의 열린
부분가군 $U$와, $U$가 $\mathbf S_C^*(V)$ 안에서 생성하는
\emph{등급 아이디얼} $U.\mathbf S_C^*(V)$를 생각한다.
$\mathbf S_C^*(V)$에서 $0$의 근방들의 기본계로
\[
  \mathfrak a.\mathbf S_C^*(V)+U.\mathbf S_C^*(V)
\]
꼴의 합들을 택한다. 여기서 $\mathfrak a$는 $C$의 열린
아이디얼들의 기본계를, $U$는 $V$의 열린 부분가군들의 기본계를
각각 움직인다.

$V$의 위상이 $C$의 위상에서 유도된 위상보다 섬세하지 않다면,
$\mathfrak aV\subset U$인 쌍 $(\mathfrak a,U)$만을 택해도 된다.
이 경우
\[
  \mathfrak a.\mathbf S_C^*(V)+U.\mathbf S_C^*(V)
  =\mathfrak a.1_C+U.\mathbf S_C^*(V).
\]
따라서 $n\geq1$일 때 각 $\mathbf S_C^n(V)$에 유도되는 위상은
$V$의 열린 부분가군 $U$에 대한
$U.\mathbf S_C^{n-1}(V)$들을 $0$의 근방 기본계로 갖는다. 특히
$V=\mathbf S_C^1(V)$ 위에서는 주어진 위상과 일치한다. 일반적으로는
여기서 유도되는 위상이 주어진 위상보다 섬세하지 않다.

또한 어느 경우에나 이렇게 정의한 \emph{위상 대수}
$\mathbf S_C^*(V)$는 위상 $C$-가군과 위상 $C$-대수의 범주에서,
$C$-가군과 $C$-대수의 범주에서와 같은 보편 문제의 해이다. 실제로
$E$를 위상 $C$-대수, $u:V\to E$를 $C$-가군 준동형이라 하고,
$\mathbf S(u)=\widetilde u$를 $\mathbf S_C^*(V)$로부터의 표준
확장이라 하자. $u$가 연속이라고 하자. $\mathfrak b$가 $E$의 열린
아이디얼이면 그 역상 $U=u^{-1}(\mathfrak b)$는 $V$의 열린
$C$-부분가군이고,
\[
  \widetilde u\bigl(U.\mathbf S_C^*(V)\bigr)
  \subset\mathfrak b.
\]
한편 $\mathfrak aE\subset\mathfrak b$인 $C$의 열린 아이디얼
$\mathfrak a$가 존재하므로
\[
  \widetilde u\bigl(\mathfrak a.\mathbf S_C^*(V)\bigr)
  \subset\mathfrak b.
\]
따라서 $\widetilde u$는 연속이다. 역으로 $\widetilde u$가
연속이고 $\mathfrak b$가 $E$의 열린 아이디얼이면, $V$의 열린
부분가군 $U$가 존재하여
$\widetilde u(U.\mathbf S_C^*(V))\subset\mathfrak b$이다. 특히
$\widetilde u(U.1_C)\subset\mathfrak b$, 곧
$u(U)\subset\mathfrak b$이므로 $u$는 연속이다.

끝으로 이산 위상을 준 위상 $C$-가군들의 표준 동형사상
\begin{equation}
\label{0.19.5.1.1-fr}
  \mathbf S_C^n(V)/U.\mathbf S_C^{n-1}(V)
  \xrightarrow{\sim}\mathbf S_C^n(V/U)
  \tag{19.5.1.1}
\end{equation}
이 있음을 상기하자.
\end{env}

\begin{env}[19.5.2]
\label{0.19.5.2-fr}
$A$를 위상환, $B$를 가환 위상 $A$-대수, $\mathfrak J$를
$B$의 아이디얼이라 하자. 여기서 $\mathfrak J$는 열려 있거나
닫혀 있다고 가정하지 않는다. 이하 $\mathfrak J$와
$\mathfrak J^n$ ($n\geq2$)에는 $B$의 위상에서 유도된 위상을,
몫
\[
  C=B/\mathfrak J,\qquad
  \mathfrak J^n/\mathfrak J^{n+1}
  =\gr_{\mathfrak J}^n(B)
\]
에는 몫위상을 준다. 그러면
$\mathfrak J^n/\mathfrak J^{n+1}$들은 위상 $C$-가군이다. 표준
단사 준동형
\[
  \mathfrak J/\mathfrak J^2\longrightarrow
  \gr_{\mathfrak J}^*(B)
\]
은 $C$-대수 준동형
\[
  \varphi:\mathbf S_C^*(\mathfrak J/\mathfrak J^2)
  \longrightarrow\gr_{\mathfrak J}^*(B)
\]
으로 확장된다. 우선 이 준동형에는 위상을 고려하지 않는다. 각
$n$에서 이는 $C$-가군의 \emph{전사} 준동형
\begin{equation}
\label{0.19.5.2.1-fr}
  \varphi_n:\mathbf S_C^n(\mathfrak J/\mathfrak J^2)
  \longrightarrow\mathfrak J^n/\mathfrak J^{n+1}
  =\gr_{\mathfrak J}^n(B)
  \tag{19.5.2.1}
\end{equation}
을 준다.

\oldpage[0\textsubscript{IV-1}]{91}

$\mathbf S_C^n(\mathfrak J/\mathfrak J^2)$에 (19.5.1)에서 정의한
위상을 주면 $\varphi_n$들은 연속이다. 실제로 (19.5.1)의 보편
성질을, $\mathfrak J^k/\mathfrak J^{k+1}$들의 위상들의 곱위상을
준 위상 $C$-대수
\[
  E=\gr_{\mathfrak J/\mathfrak J^n}(B/\mathfrak J^n)
\]
와 표준 단사 준동형
$\mathfrak J/\mathfrak J^2\to E$에 적용하면 된다. 여기서
$\mathfrak J/\mathfrak J^2$의 위상은 $C$의 위상에서 유도된
위상보다 섬세하지 않다. 이 $C$에서 유도된 위상은
$\mathfrak J/\mathfrak J^2$ 위에서 $B$의 위상으로부터 유도된
위상과도 같다.
\end{env}

\begin{theorem}[19.5.3]
\label{0.19.5.3-fr}
$A$를 위상환, $B$를 위상 $A$-대수, $\mathfrak J$를 $B$의
아이디얼이라 하고, $C=B/\mathfrak J$를 몫위상 $A$-대수라 하자.
이산 $A$-대수 $C$가 형식적으로 매끄럽다고 가정한다. 그러면 다음이
성립한다.
\begin{enumerate}
  \item[\rm(i)] $B$가 형식적으로 매끄러운 $A$-대수이면,
  $\mathfrak J/\mathfrak J^2$는 형식적으로 사영인 위상
  $C$-가군이다 (19.2.1).
  \item[\rm(ii)] $B$가 형식적으로 매끄러운 $A$-대수이면서
  준허용환이면 ($\mathbf 0_{\rm I}$, 7.1.2), (19.5.2)의
  준동형 $\varphi_n$들은 형식적 전단사 준동형이다 (19.1.2).
  \item[\rm(iii)] 역으로 $B$가 준허용환이고, 열
  $(\mathfrak J^n)$이 $B$에서 $0$으로 수렴하며,
  $\mathfrak J/\mathfrak J^2$가 형식적으로 사영인 $C$-가군이고,
  $\varphi_n$들이 형식적 전단사 준동형이라고 하자. 그러면 $B$는
  형식적으로 매끄러운 $A$-대수이다.
\end{enumerate}
\end{theorem}

이 정리의 증명은 길고 기술적인 세부사항이 많다. 따라서 먼저
주된 생각이 더 분명히 드러나는 다음의 간단한 따름정리를 증명한다.
이 따름정리는 뒤에서 가장 자주 쓰일 정리 (19.5.3)의 특별한
경우이기도 하다.

\begin{corollary}[19.5.4]
\label{0.19.5.4-fr}
$A$를 위상환, $B$를 위상 $A$-대수, $\mathfrak J$를 $B$의
아이디얼이라 하고, $B$의 위상이 $\mathfrak J$-준아딕 위상이라고
하자. 이산 $A$-대수 $C=B/\mathfrak J$가 형식적으로 매끄럽다고
가정한다. 그러면 다음 세 조건은 서로 동치이다.
\begin{enumerate}
  \item[a)] $B$는 형식적으로 매끄러운 $A$-대수이다.
  \item[b)] $\mathfrak J/\mathfrak J^2$는 사영 $C$-가군이고,
  표준 준동형
  \[
    \varphi:\mathbf S_C^*(\mathfrak J/\mathfrak J^2)
    \longrightarrow\gr_{\mathfrak J}^*(B)
  \]
  은 전단사이다.
  \item[c)] $B$의 분리 완비화 $\widehat B$는
  $\widehat A$-대수
  \[
    \widehat{B'},\qquad B'=\mathbf S_C^*(V)
  \]
  꼴의 위상 $\widehat A$-대수와 동형이다. 여기서 $V$는 사영
  $C$-가군이고, $B'$에는 그 증대 아이디얼 ${B'}^+$에 대한
  ${B'}^+$-준아딕 위상을 준다.
\end{enumerate}
\end{corollary}

\begin{env}[19.5.4.1]
\label{0.19.5.4.1-fr}
먼저 $a)$가 성립하면 $\mathfrak J/\mathfrak J^2$가 사영
$C$-가군임을 증명하자. $P,Q$를 두 $C$-가군,
$u:P\to Q$를 전사 $C$-준동형,
$v:\mathfrak J/\mathfrak J^2\to Q$를 $C$-준동형이라 하자.
$B/\mathfrak J^2$을 $\mathfrak J/\mathfrak J^2$에 의한 $C$의
$B$-확장으로 보면, $v$를 통해 그로부터 $Q$에 의한 $C$의
$B$-확장
\[
  E=(B/\mathfrak J^2)\oplus_{(\mathfrak J/\mathfrak J^2)}Q
\]
를 얻는다 (18.2.8). $C$는 이산이며 형식적으로 매끄러운
$A$-대수이므로, 확장 $E$는 $A$-자명하고 (19.4.4), 따라서
$\mathrm D_C(Q)$와 동일시할 수 있다 (18.2.3). 그러면 전사 준동형
$u:P\to Q$는 표준적으로 전사 $A$-준동형
$\psi:\mathrm D_C(P)\to\mathrm D_C(Q)$를 정한다 (18.2.9). 그 핵은
확장 $E'=\mathrm D_C(P)$의 아이디얼 $\mathfrak L$이고 $P$에
포함되므로, 더욱이 제곱영이다. $E$의 $B$-대수 구조를 정하는
준동형을 $f:B\to E=E'/\mathfrak L$라 하자. $\mathfrak L$이
제곱영이고 $B$가 형식적으로 매끄러운 $A$-대수이므로 $f$는
\[
  B\xrightarrow{g}E'\longrightarrow E'/\mathfrak L
\]
로 인수분해된다. 여기서 $g$는 연속 $A$-준동형이다.

\oldpage[0\textsubscript{IV-1}]{92}

$v'$를 $v$로부터 유도되는 사상이라 하면 (18.2.8), 도식
\[
  \xymatrix{
    E' \ar[r]^{\psi} & E \\
    B \ar[u]^{g} \ar[r]_{\theta} & B/\mathfrak J^2 \ar[u]_{v'}
  }
\]
은 가환이다. 또한 $\psi\circ g$에 의한 $\mathfrak J$의 상은
$Q$에 포함되므로, $g$에 의한 $\mathfrak J$의 상은
$P+\mathfrak L=P$에 포함되고, $g$에 의한 $\mathfrak J^2$의 상은
0이다. $g$와 $\theta$를 $\mathfrak J$에 제한하면 가환도식
\[
  \xymatrix{
    P \ar[r]^{u} & Q \\
    \mathfrak J/\mathfrak J^2 \ar[u]^{w} \ar[ur]_{v}
  }
\]
을 얻는다. 이로써 $C$-가군 $\mathfrak J/\mathfrak J^2$가
사영임이 증명된다.
\end{env}

\begin{env}[19.5.4.2]
\label{0.19.5.4.2-fr}
둘째로 $a)$가 성립하면 $\varphi$가 전단사임을 증명하자. 이로써
$a)$가 $b)$를 함의한다는 증명이 끝난다. 다음과 같이 놓는다.
\[
  E_n=B/\mathfrak J^{n+1},
\]
그리고 $F_n$을 $\mathbf S_C^*(\mathfrak J/\mathfrak J^2)$를 그
증대 아이디얼의 $(n+1)$제곱으로 나눈 몫이라 하자.
$\mathfrak J/\mathfrak J^{n+1}$은 $E_n$에서 멱영이고 $C$는 이산이며
형식적으로 매끄러운 $A$-대수이므로, $C$의 항등 자기동형사상은
\[
  C\xrightarrow{f}E_n\longrightarrow
  C=E_n/(\mathfrak J/\mathfrak J^{n+1})
\]
로 인수분해된다. 여기서 $f$는 $A$-준동형이다. 한편 (19.5.4.1)에
따라 $\mathfrak J/\mathfrak J^2$가 사영 $C$-가군이므로 그 항등
자기동형사상은
\[
  \mathfrak J/\mathfrak J^2\xrightarrow{g}
  \mathfrak J/\mathfrak J^{n+1}
  \longrightarrow\mathfrak J/\mathfrak J^2
\]
로 인수분해된다. 여기서 $g$는 $C$-선형 사상이다. $f$와 $g$로부터
표준적으로 $C$-대수 준동형
\[
  \mathbf S_C^*(\mathfrak J/\mathfrak J^2)\longrightarrow E_n
\]
을 얻는다. 이 준동형은 $g$의 정의에 따라
$\mathbf S_C^*(\mathfrak J/\mathfrak J^2)$의 증대 아이디얼의
$(n+1)$제곱 위에서 0이 된다. 따라서 몫을 취하면 전사 $A$-대수
준동형 $v:F_n\to E_n$을 얻고,
$\gr^0(v)$와 $\gr^1(v)$는 각각 $C$와
$\mathfrak J/\mathfrak J^2$의 항등 자기동형사상이다. 표준 준동형
$\varphi$의 정의에 따라 모든 $j\leq n$에 대하여
$\gr^j(v)=\varphi_j$이다.

이제 $v$의 핵 $\mathfrak R$은 $F_n$의 멱영 아이디얼이므로
$E_n$을 $F_n/\mathfrak R$과 동일시할 수 있다. $B$가 형식적으로
매끄러운 $A$-대수이므로 연속인 표준 $A$-준동형
$p_n:B\to E_n=B/\mathfrak J^{n+1}$은
\[
  B\xrightarrow{w}F_n\xrightarrow{v}E_n
\]
으로 인수분해된다. 여기서 $w$는 연속 $A$-준동형이다.
$\gr^0(v)$가 항등사상이므로 $w(\mathfrak J)$는 $F_n$의 증대
아이디얼에 포함되고, 따라서 $w(\mathfrak J^{n+1})=0$이다. 그러므로
$w$는
\[
  B\xrightarrow{p_n}E_n\xrightarrow{w'}F_n
\]
으로 인수분해되며, 합성
$E_n\xrightarrow{w'}F_n\xrightarrow{v}E_n$은 항등사상이다. 또한
$\gr^0(v)$와 $\gr^1(v)$가 각각 $C$와
$\mathfrak J/\mathfrak J^2$의 항등 자기동형사상이므로,
$\gr^0(w')$와 $\gr^1(w')$도 그러하다. $\gr^*(E_n)$은
$\gr^1(E_n)$으로 생성되므로, 합성
\[
  \gr^j(F_n)\xrightarrow{\varphi_j}\gr^j(E_n)
  \xrightarrow{\gr^j(w')}\gr^j(F_n)
\]
은 $j=0,1$일 때 항등사상이라는 사실로부터 모든 $j\leq n$에
대하여 항등사상이다. $j=n$으로 놓으면 $\varphi_n$이 단사임을
얻는다. 이로써 $a)$가 $b)$를 함의함이 증명된다.
\end{env}

\begin{env}[19.5.4.3]
\label{0.19.5.4.3-fr}
\oldpage[0\textsubscript{IV-1}]{93}
이제 $b)$가 $a)$를 함의함을 증명하자. (19.5.4.2)의 처음과 같은
논법으로, 모든 $j$에 대하여 $\gr^j(v)=\varphi_j$를 만족하는 전사
$A$-대수 준동형 $v:F_n\to E_n$이 존재한다. 모든 $\varphi_j$가
전단사이고 $F_n$과 $E_n$의 여과가 유한하므로, $v$가
\emph{전단사}라고 결론지을 수 있다 (Bourbaki,
\emph{Alg. comm.}, chap.~III, §~2, n\textsuperscript{o}~8,
cor.~3 du th.~1).

이제 $G$를 이산 위상 $A$-대수, $\mathfrak R$을 $G$의 제곱영
아이디얼, $f:B\to G/\mathfrak R$을 연속 $A$-대수 준동형이라 하자.
$G$가 이산이므로 $f$가 $\mathfrak J^m$ 위에서 0이 되는 정수 $m$이
존재한다. 따라서 $n=2m$으로 놓으면 $f$는
\[
  B\longrightarrow E_n\xrightarrow{f_n}G/\mathfrak R
\]
으로 인수분해된다. 그러므로 합성
\[
  r:C\longrightarrow F_n\xrightarrow{v}E_n
  \xrightarrow{f_n}G/\mathfrak R
\]
은 연속 $A$-대수 준동형이다. $C$는 이산이며 형식적으로 매끄러운
$A$-대수이므로 $r$은
\[
  C\xrightarrow{r'}G\longrightarrow G/\mathfrak R
\]
로 인수분해된다. 이에 따라 $r'$를 통해 $G$에 $C$-대수 구조를
준다. 한편 준동형
\[
  f'_n:F_n\xrightarrow{v}E_n\xrightarrow{f_n}G/\mathfrak R
\]
을 $\mathfrak J/\mathfrak J^2$에 제한하면 $C$-선형 사상
$g:\mathfrak J/\mathfrak J^2\to G/\mathfrak R$을 얻는다.
$\mathfrak J/\mathfrak J^2$는 사영 $C$-가군이므로 $g$는
\[
  \mathfrak J/\mathfrak J^2\xrightarrow{h}G
  \longrightarrow G/\mathfrak R
\]
로 인수분해된다. 여기서 $h$는 $C$-선형 사상이다. 이를 확장하면
$C$-대수 준동형
\[
  w:\mathbf S_C^*(\mathfrak J/\mathfrak J^2)\longrightarrow G
\]
를 얻는다. 구성상 차수가 $m$인 모든 원소의 $w$에 의한 상은
$\mathfrak R$에 속한다. 따라서 $\mathfrak R^2=0$이므로 차수가
$n=2m$인 모든 원소의 상은 0이다. 다시 말해 $w$는
\[
  \mathbf S_C^*(\mathfrak J/\mathfrak J^2)
  \longrightarrow F_n\xrightarrow{w'}G
\]
로 인수분해된다. 구성상 합성
$F_n\xrightarrow{w'}G\xrightarrow{\theta}G/\mathfrak R$은 $C$와
$\mathfrak J/\mathfrak J^2$ 위에서 $f'_n$과 일치하므로 $f'_n$과
같다. 끝으로 합성
\[
  B\longrightarrow E_n\xrightarrow{v^{-1}}F_n
  \xrightarrow{w'}G\longrightarrow G/\mathfrak R
\]
이 $f$와 같으므로, $B$는 형식적으로 매끄러운 $A$-대수이다.
\end{env}

\begin{env}[19.5.4.4]
\label{0.19.5.4.4-fr}
$a)$와 $c)$의 동치를 증명하는 일만 남았다. 먼저 $c)$는 $a)$를
함의한다. 실제로 $B'$은 이산 위상에 대하여 형식적으로 매끄러운
$C$-대수이고 (19.3.2), 따라서 그 증대 아이디얼 ${B'}^+$에 대한
${B'}^+$-준아딕 위상에 대해서도 그러하다 (19.3.8). $C$가
형식적으로 매끄러운 $A$-대수이므로 $B'$도 형식적으로 매끄러운
$A$-대수이고 (19.3.5, (ii)), 끝으로 $\widehat{B'}$은 형식적으로
매끄러운 $\widehat A$-대수이다 (19.3.6). 따라서 $B$는 형식적으로
매끄러운 $A$-대수이다 (19.3.6).

이제 $a)$가 $c)$를 함의함을 보이면 된다. 먼저 $C$가 형식적으로
매끄러운 $A$-대수이므로 항등 준동형 $C\to B/\mathfrak J$는
\[
  C\longrightarrow\widehat B\longrightarrow B/\mathfrak J
\]
로 인수분해된다. 여기서 $C\to\widehat B$는 $A$-준동형이다
(19.3.10). 따라서 $\widehat B$, 나아가 모든
$B/\mathfrak J^{n+1}$에 $C$-대수 구조가 주어진다. 한편 $b)$에 따라
$\mathfrak J/\mathfrak J^2$는 사영 $C$-가군이므로, 표준 단사
준동형 $\mathfrak J/\mathfrak J^2\to B/\mathfrak J^2$으로부터
$n\geq1$에 대한 $C$-준동형들의 사영계
\[
  u_n:\mathfrak J/\mathfrak J^2\longrightarrow
  B/\mathfrak J^{n+1}
\]
를 만들 수 있다. 따라서 대칭 대수의 보편 성질에 의해 $C$-대수
준동형들의 사영계
\[
  v_n:\mathbf S_C^*(\mathfrak J/\mathfrak J^2)=B'
  \longrightarrow B/\mathfrak J^{n+1}
\]
도 얻는다. 또한 $v_n$은 $(${B'}^+)^{n+1}$ 위에서 0이 되고,
$\gr^0(v_n)=\varphi_0$, $\gr^1(v_n)=\varphi_1$이므로
$0\leq j\leq n$인 모든 $j$에 대하여
$\gr^j(v_n)=\varphi_j$이다. $b)$에 따라 $\varphi_j$들은
동형사상이고
\[
  B'/(${B'}^+)^{n+1},\qquad B/\mathfrak J^{n+1}
\]
의 여과들은 유한하므로 $v_n$은 모든 $n$에 대하여
\emph{전단사}이다.

\oldpage[0\textsubscript{IV-1}]{94}

이는 Bourbaki, \emph{Alg. comm.}, chap.~III, §~2,
n\textsuperscript{o}~8, cor.~3 du th.~1에서 따른다. 사영극한을
취하면 $c)$를 얻는다.
\end{env}

\begin{remark}[19.5.5]
\label{0.19.5.5-fr}
(19.5.4)의 일반 가정 아래에서, 더 나아가 $\mathfrak J$가 $B$의
\emph{극대} 아이디얼이어서 $C=B/\mathfrak J$가 체이고
$\mathfrak J/\mathfrak J^2$가 \emph{유한} 차원 $C$-벡터공간이라고
가정하자. 그러면 (19.5.4)의 조건 $a)$, $b)$, $c)$는 다음 조건과도
동치이다.
\begin{enumerate}
  \item[d)] \emph{체 $C$ 위의 다항식 대수
  $E=C[T_1,\ldots,T_n]$이 주어지고, $\mathfrak n$이
  $T_i$ ($1\leq i\leq n$)들로 생성되는 $E$의 극대 아이디얼이라
  하자. 그러면 $\mathfrak n^k$를 포함하는 $E$의 모든 아이디얼
  $\mathfrak b$에 대하여, $C$-증대된 $A$-대수의 모든 준동형
  $v:B\to E/\mathfrak b$는}
  \[
    B\xrightarrow{w}E/\mathfrak n^k\longrightarrow E/\mathfrak b
  \]
  \emph{로 인수분해된다. 여기서 $w$는 $A$-준동형이다.}
\end{enumerate}

실제로 여기서는 $\mathfrak J/\mathfrak J^2$가 자유 $C$-가군이므로,
조건 $d)$가 표준 준동형
\[
  \varphi:\mathbf S_C^*(\mathfrak J/\mathfrak J^2)
  \longrightarrow\gr_{\mathfrak J}^*(B)
\]
의 전단사성을 함의함을 확인하면 충분하다. 이 경우
$\mathbf S_C^*(\mathfrak J/\mathfrak J^2)$는 다항식 대수
$E=C[T_1,\ldots,T_n]$과 동일시된다. 여기서
\[
  n=\operatorname{rg}_C(\mathfrak J/\mathfrak J^2),
\]
이고 $E$의 증대 아이디얼은 $T_i$들로 생성되는 아이디얼
$\mathfrak n$이다. 모든 정수 $k\geq0$에 대하여, 이산 $A$-대수
$C$의 형식적 매끄러움으로부터 (19.5.4.2)에서와 같이 전사
$A$-준동형
\[
  f:E/\mathfrak n^{k+1}\longrightarrow B/\mathfrak J^{k+1}
\]
이 존재하고, 모든 $j\leq k$에 대하여
$\gr^j(f)=\varphi_j$이다. 만일
$\mathfrak b/\mathfrak n^{k+1}=\operatorname{Ker}(f)$이면
$B/\mathfrak J^{k+1}$을 $E/\mathfrak b$와 동일시할 수 있다. 조건
$d)$에 따라 표준 준동형 $B\to B/\mathfrak J^{k+1}$은
\[
  B\xrightarrow{w}E/\mathfrak n^{k+1}
  \xrightarrow{f}B/\mathfrak J^{k+1}
\]
로 인수분해된다. $\gr^0(f)$가 항등사상이므로
\[
  w(\mathfrak J)\subset\mathfrak n/\mathfrak n^{k+1},
\]
따라서 $w(\mathfrak J^{k+1})=0$이다. (19.5.4.2)에서와 같이
모든 $k$에 대하여 $\varphi_k$가 단사라고 결론지을 수 있다. 이로써
주장이 증명된다.
\end{remark}

\begin{env}[19.5.6]
\label{0.19.5.6-fr}
\emph{정리 (19.5.3)의 증명}. --- $(\mathfrak b_\alpha)$를 $B$의
열린 아이디얼들의 감소하는 기본계라 하자. 다음과 같이 놓는다.
\[
  B_\alpha=B/\mathfrak b_\alpha,\qquad
  C_\alpha=B/(\mathfrak b_\alpha+\mathfrak J),\qquad
  \mathfrak J_\alpha=
  (\mathfrak b_\alpha+\mathfrak J)/\mathfrak b_\alpha.
\]
그러면
\[
  C_\alpha=B_\alpha/\mathfrak J_\alpha,\qquad
  \mathfrak J_\alpha^n/\mathfrak J_\alpha^{n+1}
  =(\mathfrak b_\alpha+\mathfrak J^n)/
   (\mathfrak b_\alpha+\mathfrak J^{n+1}).
\]
$\mathfrak J^n/\mathfrak J^{n+1}$의 $C$-부분가군
\[
  ((\mathfrak b_\alpha\cap\mathfrak J^n)+\mathfrak J^{n+1})/
  \mathfrak J^{n+1}
\]
들은 위상 $C$-가군 $\mathfrak J^n/\mathfrak J^{n+1}$의 열린
부분가군들의 기본계를 이룬다. 또한
\[
  (\mathfrak b_\alpha\cap\mathfrak J^n)+\mathfrak J^{n+1}
  =\mathfrak J^n\cap
   (\mathfrak b_\alpha+\mathfrak J^{n+1})
\]
이므로
\[
\begin{aligned}
  \mathfrak J^n/
  ((\mathfrak b_\alpha\cap\mathfrak J^n)+\mathfrak J^{n+1})
  &=
  \mathfrak J^n/
  (\mathfrak J^n\cap(\mathfrak b_\alpha+\mathfrak J^{n+1}))\\
  &=
  (\mathfrak b_\alpha+\mathfrak J^n)/
  (\mathfrak b_\alpha+\mathfrak J^{n+1})\\
  &=\mathfrak J_\alpha^n/\mathfrak J_\alpha^{n+1}.
\end{aligned}
\]
\end{env}

\begin{env}[19.5.6.1]
\label{0.19.5.6.1-fr}
$P,Q$를 두 이산 $C_\alpha$-가군이라 하고,
$u:P\to Q$를 전사 $C_\alpha$-준동형이라 하자. 이산환
$B/(\mathfrak b_\alpha+\mathfrak J^2)$은 제곱영 아이디얼
$\mathfrak J_\alpha/\mathfrak J_\alpha^2$에 의한 $C_\alpha$의
$B$-확장이다. $v:\mathfrak J/\mathfrak J^2\to Q$를 연속
$C$-준동형이라 하자. 필요하면 $\mathfrak b_\alpha$를 $B$의 더 작은
열린 아이디얼로 바꾸어, $v$의 핵이 열린 $C$-부분가군
\[
  ((\mathfrak b_\alpha\cap\mathfrak J)+\mathfrak J^2)/
  \mathfrak J^2
\]
을 포함한다고 가정할 수 있다. 몫을 취하면 이산 가군들의
$C_\alpha$-준동형
\[
  v':\mathfrak J_\alpha/\mathfrak J_\alpha^2\longrightarrow Q
\]
을 얻는다. $E_\alpha$를 $v'$를 사용하여
$B/(\mathfrak b_\alpha+\mathfrak J^2)$에서 유도되는, $Q$에 의한
$C_\alpha$의 $B_\alpha$-확장이라 하자. 또
\[
  v'_*:B/(\mathfrak b_\alpha+\mathfrak J^2)\longrightarrow E_\alpha
\]
를 이에 대응하는 $B_\alpha$-준동형이라 하자 (18.2.8).
$E_\alpha$는 이산 위상 $A$-대수이고, 표준 동형사상
$C_\alpha\xrightarrow{\sim}E_\alpha/Q$로부터 합성에 의해 연속
$A$-준동형
\[
  f:C\longrightarrow C_\alpha\longrightarrow E_\alpha/Q
\]
를 얻는다. 그런데 $Q$는 $E_\alpha$에서 제곱영이고 $C$는 형식적으로
매끄러운 $A$-대수이므로 $f$는
\[
  C\xrightarrow{g}E_\alpha\longrightarrow E_\alpha/Q
\]
로 인수분해된다. 여기서 $g$는 연속 $A$-준동형이다. $g$가 연속이고
$E_\alpha$가 이산이므로, 어떤 $\beta\geq\alpha$가 존재하여 $g$는
\[
  C\longrightarrow C_\beta\xrightarrow{g'}E_\alpha
\]
로 인수분해된다.

한편 $E_\beta$를 표준 준동형 $C_\beta\to C_\alpha$에 의한
$E_\alpha$의 역상, 곧 $Q$에 의한 $C_\beta$의 $A$-확장이라 하자.
합성 $C_\beta\to E_\alpha\to C_\alpha$가 표준 준동형이 되게 하는
$g'$의 존재는 $E_\beta$가 $A$-\emph{자명한} 확장이라는 것과
동치이다. 따라서 이를 $\mathrm D_{C_\beta}(Q)$와 동일시할 수 있다.
이제 전사 준동형 $u:P\to Q$는 표준적으로 전사 준동형
\[
  \mathrm D_{C_\beta}(P)\longrightarrow
  \mathrm D_{C_\beta}(Q)
\]
을 정한다 (18.2.9). 그 핵은 확장
$E'_\beta=\mathrm D_{C_\beta}(P)$의 아이디얼 $\mathfrak L$이고
$P$에 포함되므로, 더욱이 제곱영이다. $E_\beta$의 위상
$B$-대수 구조를 정하는 연속 $A$-준동형을
\[
  f':B\longrightarrow E_\beta=E'_\beta/\mathfrak L
\]
라 하자. $\mathfrak L$이 제곱영이고 $B$가 형식적으로 매끄러운
$A$-대수이므로 $f'$는
\[
  B\xrightarrow{h'}E'_\beta\longrightarrow E'_\beta/\mathfrak L
\]
로 인수분해된다. 여기서 $h'$은 연속 $A$-준동형이다.
$E'_\beta$의 구성으로부터 합성에 의해 $A$-준동형
\[
  \psi:E'_\beta\longrightarrow E_\beta\longrightarrow E_\alpha
\]
를 얻으며, 도식
\[
  \xymatrix{
    E'_\beta \ar[r]^{\psi} & E_\alpha \\
    B \ar[u]^{h'} \ar[r]_{\theta} &
    B/(\mathfrak b_\alpha+\mathfrak J^2) \ar[u]_{v'_*}
  }
\]
은 가환이다. 또한 $\psi\circ h'$에 의한 $\mathfrak J$의 상은
$Q$에 포함되므로, $h'$에 의한 $\mathfrak J$의 상은 $P$에
포함되고, $h'$에 의한 $\mathfrak J^2$의 상은 0이다. $h'$과
$\theta$를 $\mathfrak J$에 제한하면 가환도식
\[
  \xymatrix{
    P \ar[r]^{u} & Q \\
    \mathfrak J/\mathfrak J^2 \ar[u]^{w} \ar[ur]_{v}
  }
\]
을 얻는다. 여기서 $w$는 연속이다. 이로써
$\mathfrak J/\mathfrak J^2$가 형식적으로 사영인 $C$-가군임이
증명된다.
\end{env}

\oldpage[0\textsubscript{IV-1}]{95}

\begin{env}[19.5.6.2]
\label{0.19.5.6.2-fr}
모든 정수 $n\geq0$에 대하여
\[
  E_n=B/\mathfrak J^{n+1},
  \qquad\text{따라서}\qquad E_0=C
\]
로 놓는다. 아이디얼들
\[
  (\mathfrak b_\alpha+\mathfrak J^{n+1})/\mathfrak J^{n+1}
\]
은 $E_n$의 열린 아이디얼들의 기본계를 이룬다. 또
\[
  E_{\alpha,n}
  =B/(\mathfrak b_\alpha+\mathfrak J^{n+1})
  =B_\alpha/\mathfrak J_\alpha^{n+1}
\]
로 놓으면, 이는 이산 몫환이다. 마찬가지로
\[
  \gr_{\mathfrak J}^n(B)
  =\mathfrak J^n/\mathfrak J^{n+1}
\]
에서
\[
  ((\mathfrak b_\alpha\cap\mathfrak J^n)+\mathfrak J^{n+1})/
  \mathfrak J^{n+1}
\]
들이 0의 근방들의 기본계를 이룸을 이미 보았다. 이 부분가군으로
$\mathfrak J^n/\mathfrak J^{n+1}$을 나눈 몫은 표준적으로
$\mathfrak J_\alpha^n/\mathfrak J_\alpha^{n+1}$과 동일시된다.
대칭 대수
\[
  \mathbf S_{C_\alpha}^*
  (\mathfrak J_\alpha/\mathfrak J_\alpha^2)
\]
를 생각하자. $F_{\alpha,n}$을 이 대칭 대수를 그 증대 아이디얼의
$(n+1)$제곱으로 나눈 몫이라 하자. 고정된 $n\geq1$에 대하여
(19.5.1.1)로부터
\[
  \mathbf S_{C_\alpha}^n
  (\mathfrak J_\alpha/\mathfrak J_\alpha^2)
\]
들은 위상 $C$-가군
$\mathbf S_C^n(\mathfrak J/\mathfrak J^2)$을 그 열린 부분가군들의
기본계로 나눈 몫들임을 얻는다.

용어를 줄이기 위하여, $\alpha\leq\beta$일 때의 표준 준동형
\[
\begin{gathered}
  B_\beta\longrightarrow B_\alpha,\qquad
  C_\beta\longrightarrow C_\alpha,\qquad
  \mathfrak J_\beta/\mathfrak J_\beta^2
  \longrightarrow\mathfrak J_\alpha/\mathfrak J_\alpha^2,\\
  E_{\beta,n}\longrightarrow E_{\alpha,n},\qquad
  F_{\beta,n}\longrightarrow F_{\alpha,n},\quad\ldots
\end{gathered}
\]
들을 \emph{전이 준동형}이라 부르겠다.
\end{env}

\begin{lemma}[19.5.6.3]
\label{0.19.5.6.3-fr}
$A$-대수 $C$가 형식적으로 매끄럽고, 환 $B$가 준허용환이며,
$C$-가군 $\mathfrak J/\mathfrak J^2$가 형식적으로 사영이라고
가정하자. 그러면 다음이 성립한다.
\begin{enumerate}
  \item[\rm(i)] 모든 $\alpha$에 대하여 $\beta\geq\alpha$와 전사
  $A$-대수 준동형
  \[
    v_{\alpha\beta}:F_{\beta,n}\longrightarrow E_{\alpha,n}
  \]
  이 존재하여,
  \[
    \gr^0(v_{\alpha\beta}):C_\beta\longrightarrow C_\alpha
  \]
  및
  \[
    \gr^1(v_{\alpha\beta}):
    \mathfrak J_\beta/\mathfrak J_\beta^2
    \longrightarrow\mathfrak J_\alpha/\mathfrak J_\alpha^2
  \]
  는 전이 준동형이다.
  \item[\rm(ii)] $\beta$와 $v_{\alpha\beta}$가 \emph{(i)}의 조건을
  만족한다고 하자. 그러면 모든 $\gamma\geq\beta$에 대하여
  $\delta\geq\gamma$와 전사 $A$-대수 준동형
  \[
    v_{\gamma\delta}:F_{\delta,n}\longrightarrow E_{\gamma,n}
  \]
  이 존재한다. 이들은 $\gamma,\delta$에 대하여 \emph{(i)}의
  조건을 만족하고, 다음 도식을 가환으로 만든다.
  \begin{equation}
  \label{0.19.5.6.4-fr}
  \xymatrix@C=3em{
    F_{\beta,n} \ar[r]^{v_{\alpha\beta}} & E_{\alpha,n} \\
    F_{\delta,n} \ar[u] \ar[r]_{v_{\gamma\delta}} &
    E_{\gamma,n} \ar[u]
  }
  \tag{19.5.6.4}
  \end{equation}
  여기서 수직 화살표들은 전이 준동형이다.
\end{enumerate}
\end{lemma}

\begin{env}[19.5.6.4]
\label{0.19.5.6.4-proof-fr}
\begin{enumerate}
  \item[\rm(i)] 이산 위상 $A$-대수 $E_{\alpha,n}$에서 아이디얼
  $\mathfrak J_\alpha/\mathfrak J_\alpha^{n+1}$은 멱영이고, 자연스러운
  동일시 동형
  $C_\alpha\xrightarrow{\sim}
  E_{\alpha,n}/(\mathfrak J_\alpha/\mathfrak J_\alpha^{n+1})$을 합성하면
  연속 $A$-준동형
  $C\to C_\alpha\to
  E_{\alpha,n}/(\mathfrak J_\alpha/\mathfrak J_\alpha^{n+1})$을 얻는다.

\oldpage[0\textsubscript{IV-1}]{96}

  $C$는 형식적으로 매끄러운 $A$-대수이므로 이 준동형은

  \[
    C\xrightarrow{f_\alpha}E_{\alpha,n}\to
    E_{\alpha,n}/(\mathfrak J_\alpha/\mathfrak J_\alpha^{n+1})
  \]

  으로 인수분해된다. 여기서 $f_\alpha$는 연속 $A$-준동형이다. 따라서
  $\mathfrak J_\alpha/\mathfrak J_\alpha^{n+1}$은 $f_\alpha$를 통하여
  $C$의 한 열린 아이디얼에 의해 소멸되는 이산 위상 $C$-가군이 된다.
  $\mathfrak J/\mathfrak J^2$가 형식적으로 사영인 $C$-가군이라는
  가정으로부터, 다음 도식을 가환으로 만드는 연속 $C$-선형 사상

  \[
    g_\alpha:\mathfrak J/\mathfrak J^2\longrightarrow
    \mathfrak J_\alpha/\mathfrak J_\alpha^{n+1}
  \]

  이 존재한다.
  \[
    \xymatrix{
      \mathfrak J_\alpha/\mathfrak J_\alpha^{n+1} \ar[r] &
      \mathfrak J_\alpha/\mathfrak J_\alpha^2 \\
      & \mathfrak J/\mathfrak J^2 \ar[ul]^{g_\alpha} \ar[u]
    }
  \]
  $f_\alpha$와 $g_\alpha$가 연속이므로, 이 준동형들이 각각

  \[
    C\to C_\beta\xrightarrow{f'_{\alpha\beta}}E_{\alpha,n},
  \]

  \[
    \mathfrak J/\mathfrak J^2\to
    \mathfrak J_\beta/\mathfrak J_\beta^2
    \xrightarrow{g'_{\alpha\beta}}
    \mathfrak J_\alpha/\mathfrak J_\alpha^{n+1}
  \]

  으로 인수분해되게 하는 $\beta\geq\alpha$가 존재한다.
  $f'_{\alpha\beta}$와 $g'_{\alpha\beta}$로부터 표준적으로
  $C_\beta$-대수 준동형
  \[
    \mathbf S_{C_\beta}^*(\mathfrak J_\beta/\mathfrak J_\beta^2)
    \longrightarrow E_{\alpha,n}
  \]
  을 얻는다. 이 준동형은 $g'_{\alpha\beta}$의 정의에 따라
  $\mathbf S_{C_\beta}^*(\mathfrak J_\beta/\mathfrak J_\beta^2)$의 증대
  아이디얼의 $(n+1)$제곱 위에서 $0$이 된다. 이 $(n+1)$제곱을 법으로
  한 몫으로 내려가면, $f_\alpha$와 $g_\alpha$의 정의에 비추어 원하는
  준동형 $v_{\alpha\beta}$를 얻는다. 실제로
  $\gr^0(v_{\alpha\beta})$와 $\gr^1(v_{\alpha\beta})$가 모두
  전사이므로 $\gr(v_{\alpha\beta})$가 전사이다. 이는 등급 대수
  $\gr^*(E_{\alpha,n})$가 $\gr^0(E_{\alpha,n})$와
  $\gr^1(E_{\alpha,n})$에 의해 생성되기 때문이다. 이제 여기서 다루는
  여과들은 유한하므로 Bourbaki, \emph{Alg. comm.}, chap.~III, §~2,
  n\textsuperscript{o}~8, cor.~3 du th.~1을 적용할 수 있고,
  $v_{\alpha\beta}$의 전사성이 따른다.

  \item[\rm(ii)] $B$가 \emph{준허용환}이라는 가정은 모든
  $\mathfrak b_\alpha$가, 그 거듭제곱들이 $0$으로 \emph{수렴하는} 하나의
  $\mathfrak b_{\alpha_0}$ 안에 들도록 할 수 있음을 뜻한다. 특히 전이
  준동형 $E_{\gamma,n}\to E_{\alpha,n}$의 핵
  \[
    (\mathfrak b_\alpha+\mathfrak J^{n+1})/
    (\mathfrak b_\gamma+\mathfrak J^{n+1})
  \]
  은 \emph{멱영}이다. 보조정리 (19.3.10.1)을 적용하면 다음 도식이
  가환이 되도록 $f_\gamma$를 선택할 수 있다.
  \[
    \xymatrix{
      & E_{\alpha,n} \\
      C \ar[ur]^{f_\alpha} \ar[r]_{f_\gamma} & E_{\gamma,n} \ar[u]
    }
  \]
  여기서 수직 화살표는 전이 준동형이다. 한편
  $\mathfrak J/\mathfrak J^2$가 형식적으로 사영인 $C$-가군이라는
  가정에 따라, 다음 도식이 가환이 되도록

  \[
    g_\gamma:\mathfrak J/\mathfrak J^2\longrightarrow
    \mathfrak J_\gamma/\mathfrak J_\gamma^{n+1}
  \]

  을 선택할 수 있다.
  \[
    \xymatrix@C=5em@R=3em{
      \mathfrak J_\alpha/\mathfrak J_\alpha^{n+1} \ar[r] &
      \mathfrak J_\alpha/\mathfrak J_\alpha^2 \\
      \mathfrak J_\gamma/\mathfrak J_\gamma^{n+1} \ar[u] \ar[r] &
      \mathfrak J_\gamma/\mathfrak J_\gamma^2 \ar[u] \\
      & \mathfrak J/\mathfrak J^2
        \ar[ul]_{g_\gamma} \ar[u]_{p_\gamma} \ar[uul]^{g_\alpha}
    }
  \]
  실제로 $\mathfrak J/\mathfrak J^2$의 $(g_\alpha,p_\gamma)$에 의한
  상은 곱가군
  \[
    (\mathfrak J_\alpha/\mathfrak J_\alpha^{n+1})
    \times(\mathfrak J_\gamma/\mathfrak J_\gamma^2)
  \]
  안에서, $g_\alpha$의 정의와 관계 $\beta\leq\gamma$에 따라
  $P=\mathfrak J_\gamma/\mathfrak J_\gamma^{n+1}$의 이 곱가군 안에서의
  표준 상 $Q$에 포함된다. 이제 전사 준동형 $P\to Q$와
  준동형
  \[
    (g_\alpha,p_\gamma):\mathfrak J/\mathfrak J^2\longrightarrow Q
  \]
  에 정의 (19.2.1)을 적용하면 된다. 이렇게 $f_\gamma$와
  $g_\gamma$를 선택하면 \emph{(i)}에서와 같이 $v_{\gamma\delta}$를
  구성하면서 동시에 도식 (19.5.6.4)를 가환으로 만들 수 있다.
\end{enumerate}
\end{env}

\oldpage[0\textsubscript{IV-1}]{97}

\begin{lemma}[19.5.6.5]
\label{0.19.5.6.5-fr}
$A$-대수 $B$와 $C$가 형식적으로 매끄럽고 $B$가 준허용환이라고
가정하자. 따라서 (19.5.6.1)에 의하여 (19.5.6.3)의 조건들이 성립한다.
두 지표 $\alpha\leq\beta$와 (19.5.6.3, \emph{(i)})의 조건을 만족하는
준동형 $v_{\alpha\beta}$의 모든 조에 대하여, 지표
$\lambda\geq\beta$와 $A$-대수 준동형
\[
  w_{\beta\lambda}:E_{\lambda,n}\longrightarrow F_{\beta,n}
\]
이 존재하여 다음을 만족한다.
$1^\circ$ $\gr^0(w_{\beta\lambda}):C_\lambda\to C_\beta$와
\[
  \gr^1(w_{\beta\lambda}):
  \mathfrak J_\lambda/\mathfrak J_\lambda^2\longrightarrow
  \mathfrak J_\beta/\mathfrak J_\beta^2
\]
는 전이 준동형이다.
$2^\circ$ 합성
\[
  E_{\lambda,n}\xrightarrow{w_{\beta\lambda}}F_{\beta,n}
  \xrightarrow{v_{\alpha\beta}}E_{\alpha,n}
\]
은 전이 준동형이다.
\end{lemma}

\begin{proof}
보조정리 (19.5.6.3, \emph{(ii)})를 $\gamma=\beta$로 적용하면,
$\delta\geq\beta$와
\[
  v_{\beta\delta}:F_{\delta,n}\longrightarrow E_{\beta,n}
\]
을 얻는다. $v_{\beta\delta}$는 전사임을 상기하자. 한편 그 핵
$\mathfrak R$은 \emph{멱영}이다. 실제로 $F_{\delta,n}$의 증대 아이디얼
$F_{\delta,n}^+$는 정의에 따라 멱영이고,
\[
  \gr^0(v_{\beta\delta}):C_\delta\longrightarrow C_\beta
\]
의 핵도 $B$가 준허용환이라는 가정에 따라 멱영이다. 따라서
$E_{\beta,n}$은 $F_{\delta,n}/\mathfrak R$과 동일시된다. $B$가
형식적으로 매끄러운 $A$-대수이므로, 연속인 표준 $A$-준동형
$p_{\beta,n}:B\to E_{\beta,n}$은
\[
  B\xrightarrow{w}F_{\delta,n}
  \xrightarrow{v_{\beta\delta}}E_{\beta,n}
\]
으로 인수분해된다. 여기서 $w$는 연속 $A$-준동형이다.
$F_{\delta,n}$은 이산이므로, $w$가 $\mathfrak b_\lambda$ 위에서
$0$이 되게 하는 $\lambda\geq\delta$가 존재한다. 따라서 $w$는
\[
  B\longrightarrow B/\mathfrak b_\lambda
  \xrightarrow{w'_\lambda}F_{\delta,n}
\]
으로 인수분해된다. 합성
\[
  B/\mathfrak b_\lambda\xrightarrow{w'_\lambda}F_{\delta,n}
  \xrightarrow{q_{\beta\delta}}F_{\beta,n}
\]
을 생각하자. 여기서 $q_{\beta\delta}$는 전이 준동형이다. 또한
\[
  \gr^0(q_{\beta\delta})=\gr^0(v_{\beta\delta}):
  C_\delta\longrightarrow C_\beta
\]
는 전이 준동형이다. 합성
$v_{\alpha\beta}\circ q_{\beta\delta}\circ w'_\lambda$는 표준
준동형
\[
  B/\mathfrak b_\lambda\longrightarrow
  B/(\mathfrak b_\alpha+\mathfrak J^{n+1})
\]
이므로, $q_{\beta\delta}\circ w'_\lambda$에 의한
\[
  \mathfrak J_\lambda=
  (\mathfrak b_\lambda+\mathfrak J)/\mathfrak b_\lambda
\]
의 상은 증대 아이디얼 $F_{\beta,n}^+$에 포함된다. 따라서
\[
  \mathfrak J_\lambda^{n+1}=
  (\mathfrak b_\lambda+\mathfrak J^{n+1})/\mathfrak b_\lambda
\]
의 상은 $0$이다. 다시 말해 $q_{\beta\delta}\circ w'_\lambda$는
\[
  B/\mathfrak b_\lambda\longrightarrow
  B/(\mathfrak b_\lambda+\mathfrak J^{n+1})=E_{\lambda,n}
  \xrightarrow{w_{\beta\lambda}}F_{\beta,n}
\]
으로 인수분해되며, 합성
\[
  E_{\lambda,n}\xrightarrow{w_{\beta\lambda}}F_{\beta,n}
  \xrightarrow{v_{\alpha\beta}}E_{\alpha,n}
\]
은 전이 준동형이다. 앞의 논증은
$\gr^0(w_{\beta\lambda})$도 전이 준동형 $C_\lambda\to C_\beta$임을
보인다. 실제로 이 사상은 합성
\[
  B/(\mathfrak b_\lambda+\mathfrak J)longrightarrow
  \gr^0(F_{\delta,n})
  \xrightarrow{\gr^0(q_{\beta\delta})}\gr^0(F_{\beta,n})
\]
이고, $v_{\beta\delta}\circ w$는 표준 준동형이다. 또한
(19.5.6.3)에 따라
\[
  \gr^1(q_{\beta\delta})=\gr^1(v_{\beta\delta})
\]
이므로 같은 논증으로 $\gr^1(w_{\beta\lambda})$가 전이 준동형임이
따른다.
\end{proof}

\begin{env}[19.5.6.6]
\label{0.19.5.6.6-fr}
\emph{(19.5.3, \emph{(ii)})의 증명}. --- $\varphi_n$이 형식적 전단사
준동형임을 보이기 위하여 (19.1.3, \emph{(iii)})의 판정법을 적용하자.
가정에 따라 보조정리 (19.5.6.5)의 조건들이 성립하므로, 모든 지표
$\alpha$에 대하여 그 보조정리의 결론을 만족하는
$v_{\alpha\beta}$와 $w_{\beta\lambda}$를 정한다. 준동형
\[
  \gr^n(v_{\alpha\beta}):\gr^n(F_{\beta,n})
  \longrightarrow\gr^n(E_{\alpha,n})
\]
은 표준 준동형 $\varphi_n$ (19.5.2.1)에서 몫으로 내려가 얻는 준동형
\[
  \varphi_{\alpha\beta,n}:
  \mathbf S_{C_\beta}^n(\mathfrak J_\beta/\mathfrak J_\beta^2)
  \longrightarrow\mathfrak J_\alpha^n/\mathfrak J_\alpha^{n+1}
\]
과 다르지 않다. 실제로 (19.5.6.3)에 따라
$\gr^0(v_{\alpha\beta})$와 $\gr^1(v_{\alpha\beta})$는 각각
$\varphi_{\alpha\beta,0}$ 및 $\varphi_{\alpha\beta,1}$과 일치하고,
표준 준동형 $\varphi$의 정의로부터 $j\leq n$에 관한 귀납법에 의해
모든 $j$에 대하여
\[
  \gr^j(v_{\alpha\beta})=\varphi_{\alpha\beta,j}
\]
임이 따른다. 이제 $\varphi_n$은 전사이므로, 당연히 형식적 전사
준동형이다. 또한 $j\leq n$에 대하여 합성
\[
  \gr^j(F_{\lambda,n})
  \xrightarrow{\varphi_{\lambda\lambda,j}}
  \gr^j(E_{\lambda,n})
  \xrightarrow{\gr^j(w_{\beta\lambda})}
  \gr^j(F_{\beta,n})
\]
은 전이 준동형이다. 보조정리 (19.5.6.5)에 의해 $j=0,1$일 때 그러하고,
$\gr^*(E_{\lambda,n})$가 $\gr^1(E_{\lambda,n})$에 의해 생성되므로
$j$에 관한 귀납법으로 주장이 따른다. 여기에 전이 준동형
$\gr^j(F_{\beta,n})\to\gr^j(F_{\alpha,n})$을 합성하면, $j=n$에
대하여 전이 준동형
$\gr^n(F_{\lambda,n})\to\gr^n(F_{\alpha,n})$을
\[
  \gr^n(F_{\lambda,n})
  \xrightarrow{\varphi_{\lambda\lambda,n}}
  \gr^n(E_{\lambda,n})
  \longrightarrow\gr^n(F_{\alpha,n})
\]
으로 인수분해한 것이다. 이는 $\varphi_n$이 형식적 전단사 준동형이
되기 위한 (19.1.3)의 판정 조건이다.
\end{env}

\oldpage[0\textsubscript{IV-1}]{98}

\begin{lemma}[19.5.6.7]
\label{0.19.5.6.7-fr}
$B$가 준허용환이고, $\mathfrak J/\mathfrak J^2$가 형식적으로 사영인
$C$-가군이며, $C$가 형식적으로 매끄러운 $A$-대수이고, 모든
$\varphi_n$이 형식적 전단사 준동형이라고 가정하자. 그러면
$\alpha\leq\beta$인 모든 지표쌍과 (19.5.6.3, \emph{(i)})의 조건을
만족하는 모든 준동형
$v_{\alpha\beta}:F_{\beta,n}\to E_{\alpha,n}$에 대하여, 어떤 지표
$\gamma\geq\beta$가 존재하여 모든 $\mu\geq\gamma$에 대해
$A$-준동형들의 가환도식
\[
\xymatrix@C=4em@R=3em{
  F_{\beta,n}\ar[r]^{v_{\alpha\beta}} & E_{\alpha,n} \\
  & E_{\mu,n}\ar[ul]_{u_{\beta\mu}}\ar[u]_{p_{\alpha\mu}}
}
\]
이 존재한다. 여기서 $p_{\alpha\mu}$는 전이 준동형이다.
\end{lemma}

\begin{proof}
각 $\varphi_j$에 $0\leq j\leq n$에 대하여 판정법 (19.1.3)을
적용하면, 지표 $\gamma\geq\beta$와 유일하게 정해지는 전사
$C_\gamma$-가군 준동형
\[
  w_{\beta\gamma}^j:\gr^j(E_{\gamma,n})longrightarrow
  \gr^j(F_{\beta,n})
\]
들이 존재하여 합성
\[
  \gr^j(F_{\gamma,n})
  \xrightarrow{\varphi_{\gamma\gamma,j}}\gr^j(E_{\gamma,n})
  \xrightarrow{w_{\beta\gamma}^j}\gr^j(F_{\beta,n})
\]
이 전이 준동형이 된다. $\varphi_j$들은 유한 개이므로 모든
$\varphi_j$에 대하여 같은 지표 $\gamma$를 선택할 수 있다. 또한
$w_{\beta\gamma}^j$들의 유일성과 $\varphi$가 등급 대수 준동형이라는
사실로부터
\[
  w_{\beta\gamma}=(w_{\beta\gamma}^j)_{0\leq j\leq n}
\]
은 등급 $C_\gamma$-대수 준동형
\[
  \gr^*(E_{\gamma,n})\longrightarrow
  \gr^*(F_{\beta,n})=F_{\beta,n}
\]
이다. 나아가 $\varphi_{\gamma\gamma,0}$과
$\varphi_{\gamma\gamma,1}$은 항등 준동형이므로,
\[
  w_{\beta\gamma}^0:C_\gamma\longrightarrow C_\beta,
  \qquad
  w_{\beta\gamma}^1:
  \mathfrak J_\gamma/\mathfrak J_\gamma^2\longrightarrow
  \mathfrak J_\beta/\mathfrak J_\beta^2
\]
는 전이 준동형이다. 따라서
$\gr^0(v_{\alpha\beta}\circ w_{\beta\gamma})$와
$\gr^1(v_{\alpha\beta}\circ w_{\beta\gamma})$도 전이 준동형이다.
$\gr^*(E_{\gamma,n})$는 $\gr^1(E_{\gamma,n})$에 의해 생성되므로,
$0\leq j\leq n$인 모든 $j$에 대하여
$\gr^j(v_{\alpha\beta}\circ w_{\beta\gamma})$도 전이 준동형이다.

이제 $\alpha$, $\beta$, $\gamma$에 보조정리 (19.5.6.3,
\emph{(ii)})를 적용하면 $\delta\geq\gamma$인 도식 (19.5.6.4)를 얻는다.
그런 다음 처음의 논증에서 $\alpha$, $\beta$를 각각 $\gamma$,
$\delta$로 바꾸어 되풀이하자. 그러면 지표 $\lambda\geq\delta$와
다음 준동형들의 가환도식을 얻는다.
\[
\xymatrix@C=4em@R=4em{
  \gr^*(E_{\gamma,n})\ar[r]^{w_{\beta\gamma}} &
  F_{\beta,n}\ar[r]^{v_{\alpha\beta}} & E_{\alpha,n} \\
  \gr^*(E_{\lambda,n})\ar[u]_{\gr^*(p_{\gamma\lambda})}
    \ar[r]_{w_{\delta\lambda}} &
  F_{\delta,n}\ar[u]_{q_{\beta\delta}}\ar[r]_{v_{\gamma\delta}} &
  E_{\gamma,n}\ar[u]_{p_{\alpha\gamma}}
    \ar@{-->}[ul]_{u_{\beta\gamma}}
}
\]
여기서 수직 화살표들은 전이 준동형이다. 이제 도식을 가환으로 만드는
준동형 $u_{\beta\gamma}$가 존재함을 보이면 충분하다. 이를 위해서는
$v_{\gamma\delta}$와 $q_{\beta\delta}$가 모두 전사이므로
\[
  \operatorname{Ker}(v_{\gamma\delta})
  \subset\operatorname{Ker}(q_{\beta\delta})
\]
임을 보이면 된다. $w_{\delta\lambda}$가 전사이므로 이 관계는
\[
\begin{split}
  \operatorname{Ker}(v_{\gamma\delta}\circ w_{\delta\lambda})
  &\subset\operatorname{Ker}(q_{\beta\delta}\circ w_{\delta\lambda})\\
  &=\operatorname{Ker}\bigl(w_{\beta\gamma}\circ
  \gr(p_{\gamma\lambda})\bigr)
\end{split}
\]
과 동치이다. 그런데 위에서
\[
  \gr(p_{\gamma\lambda})=
  \gr(v_{\gamma\delta}\circ w_{\delta\lambda})
\]
임을 보았고, 다음 포함관계는 명백하다.
\[
\begin{split}
  \operatorname{Ker}(v_{\gamma\delta}\circ w_{\delta\lambda})
  &\subset\operatorname{Ker}\bigl(
  \gr(v_{\gamma\delta}\circ w_{\delta\lambda})\bigr)\\
  &=\operatorname{Ker}(\gr(p_{\gamma\lambda}))\\
  &\subset\operatorname{Ker}\bigl(w_{\beta\gamma}\circ
  \gr(p_{\gamma\lambda})\bigr).
\end{split}
\]
이로써 보조정리의 증명이 끝난다. 실제로 모든 $\mu\geq\gamma$에
대하여 $u_{\beta\mu}$를 합성
\[
  E_{\mu,n}\xrightarrow{p_{\gamma\mu}}E_{\gamma,n}
  \xrightarrow{u_{\beta\gamma}}F_{\beta,n}
\]
으로 정의하면 된다.
\end{proof}

\begin{env}[19.5.6.8]
\label{0.19.5.6.8-fr}
\emph{(19.5.3, \emph{(iii)})의 증명}. --- $G$를 이산 위상
$A$-대수, $\mathfrak R$을 $G$의 제곱영 아이디얼,
$f:B\to G/\mathfrak R$을 연속 $A$-대수 준동형이라 하자.
$G/\mathfrak R$이 이산이므로 $f$가 $\mathfrak b_\alpha$ 위에서
$0$이 되는 지표 $\alpha$가 존재한다. 가정에 따라 어떤 정수 $m$에
대하여 $\mathfrak J^m\subset\mathfrak b_\alpha$이므로, $n=2m$으로
놓으면 $f$는
\[
  B\xrightarrow{p_\alpha}E_{\alpha,n}
  \xrightarrow{f_{\alpha,n}}G/\mathfrak R
\]
으로 인수분해된다. 보조정리 (19.5.6.3)의 가정들이 성립하므로 먼저
$\beta\geq\alpha$와 다음 합성 $A$-준동형을 얻는다.
\begin{equation}
\label{0.19.5.6.9-fr}
  f_{\alpha,n}\circ v_{\alpha\beta}:F_{\beta,n}
  \xrightarrow{v_{\alpha\beta}}E_{\alpha,n}
  \xrightarrow{f_{\alpha,n}}G/\mathfrak R.
  \tag{19.5.6.9}
\end{equation}
$F_{\beta,n}$은 $C_\beta$-대수이고, 따라서 특히 이산 위상
$C$-대수이다. 그러므로 합성하여 연속 $A$-준동형
\[
  r:C\longrightarrow F_{\beta,n}\longrightarrow G/\mathfrak R
\]
을 얻는다. 한편 $C$는 형식적으로 매끄러운 $A$-대수이므로,

\oldpage[0\textsubscript{IV-1}]{99}

이 준동형은
\[
  r:C\xrightarrow{r'}G\longrightarrow G/\mathfrak R
\]
으로 인수분해된다. 여기서 $r'$은 연속 $A$-준동형이다. 이로써
$G$에는 이산 위상 $C$-대수 구조가 주어진다. 다른 한편 표준 준동형
\[
  \mathfrak J/\mathfrak J^2\longrightarrow
  \mathfrak J_\beta/\mathfrak J_\beta^2\longrightarrow F_{\beta,n}
\]
과 합성하여 (19.5.6.9)에서 연속 $C$-준동형
\[
  g:\mathfrak J/\mathfrak J^2\longrightarrow G/\mathfrak R
\]
을 얻는다. $\mathfrak J/\mathfrak J^2$는 형식적으로 사영인
$C$-가군이므로 $g$는
\[
  \mathfrak J/\mathfrak J^2\xrightarrow{h}G
  \longrightarrow G/\mathfrak R
\]
으로 인수분해된다. 여기서 $h$는 연속 $C$-준동형이다. $G$가
이산이므로, 다음 두 상이 모두 $0$이 되게 하는 지표
$\gamma\geq\beta$가 존재한다.
\[
  h\left(((\mathfrak b_\gamma\cap\mathfrak J)+\mathfrak J^2)/
  \mathfrak J^2\right),
  \qquad
  r'\left((\mathfrak b_\gamma+\mathfrak J)/\mathfrak J\right).
\]
따라서 $h$는
\[
  \mathfrak J/\mathfrak J^2\longrightarrow
  \mathfrak J_\gamma/\mathfrak J_\gamma^2
  \xrightarrow{h'}G
\]
으로 인수분해되며, $h'$은 $C_\gamma$-준동형이다. 이를 연장하면
$h'$에서 $C_\gamma$-대수 준동형
\[
  w:\mathbf S_{C_\gamma}^*(\mathfrak J_\gamma/\mathfrak J_\gamma^2)
  \longrightarrow G
\]
을 얻는다.
구성상 차수 $m$인 모든 원소의 $w$에 의한 상은 $\mathfrak R$에
속한다. 따라서 차수 $n=2m$인 모든 원소의 상은
$\mathfrak R^2=0$이므로 $0$이다. 다시 말해 $w$는
\[
  \mathbf S_{C_\gamma}^*(\mathfrak J_\gamma/\mathfrak J_\gamma^2)
  \longrightarrow F_{\gamma,n}\xrightarrow{w'_\gamma}G
\]
로 인수분해된다. 이제 가정이 성립하는 준동형
\[
  v_{\alpha\gamma}:F_{\gamma,n}
  \xrightarrow{q_{\beta\gamma}}F_{\beta,n}
  \xrightarrow{v_{\alpha\beta}}E_{\alpha,n}
\]
에 보조정리 (19.5.6.7)을 적용한다. 그러면
$\delta\geq\gamma$와 준동형
$u_{\gamma\delta}:E_{\delta,n}\to F_{\gamma,n}$가 존재하여, 합성
$v_{\alpha\gamma}\circ u_{\gamma\delta}$는 전이 준동형
$p_{\alpha\delta}:E_{\delta,n}\to E_{\alpha,n}$이 된다. 마침내 연속
$A$-준동형들의 다음 가환 도식을 얻는다.
\[
\xymatrix@C=3em@R=2.5em{
  && E_{\alpha,n}\ar[r]^{f_{\alpha,n}} & G/\mathfrak R \\
  && F_{\beta,n}\ar[u]_{v_{\alpha\beta}} & \\
  B\ar[r]_{p_\delta} & E_{\delta,n}\ar[uur]^{p_{\alpha\delta}}
    \ar[r]_{u_{\gamma\delta}} &
  F_{\gamma,n}\ar[u]_{q_{\beta\gamma}}\ar[r]^{w'_\gamma} & G\ar[uu]
}
\]
그러므로 아래 행의 준동형들의 합성 $f':B\to G$에 대하여 $f$는
$B\xrightarrow{f'}G\to G/\mathfrak R$로 인수분해된다. 따라서 $B$는
형식적으로 매끄러운 $A$-대수이다.

이로써 정리 (19.5.3)의 증명이 끝났다.
\end{env}

\begin{corollary}[19.5.7]
\label{0.19.5.7-ko}
$A$를 위상환, $B$를 위상 $A$-대수,
$(\mathfrak b_\lambda)$를 $B$ 안의 열린 아이디얼들의 기본계,
$\mathfrak J$를 $B$의 아이디얼이라 하고, $C=B/\mathfrak J$를 몫 위상을
갖춘 $A$-대수라 하자. $C_\lambda=B/(\mathfrak b_\lambda+\mathfrak J)$로
놓는다. 다음을 가정한다. $1^\circ$ 모든 $n$에 대하여 $B$의 위상에서
$\mathfrak J^n$에 유도되는 위상은 $B$의 위상으로부터 얻어지는
$B$-가군 $\mathfrak J^n$의 위상이기도 하다 (19.0.2). 특히 $B$가
뇌터환이고 그 위상이 준아딕이면 이 조건이 성립한다
($\mathbf 0_{\mathrm I}$, 7.3.2). $2^\circ$ $C$는 형식적으로 매끄러운
$A$-대수이다. 이 조건 아래에서 다음이 성립한다.
\begin{enumerate}
  \item[\rm(i)] $B$가 형식적으로 매끄러운 $A$-대수이면, 모든
  $\lambda$에 대하여
  $(\mathfrak J/\mathfrak J^2)\otimes_C C_\lambda$는 사영
  $C_\lambda$-가군이다.
  \item[\rm(ii)] $B$가 형식적으로 매끄러운 $A$-대수이자 준허용환이면,
  모든 $\lambda$에 대하여 표준 준동형
  \begin{equation}
  \label{0.19.5.7.1-ko}
    \varphi_\lambda=\varphi\otimes 1_{C_\lambda}:
    \mathbf S_{C_\lambda}^*((\mathfrak J/\mathfrak J^2)\otimes_C C_\lambda)
    \longrightarrow\gr_{\mathfrak J}^*(B)\otimes_C C_\lambda
    \tag{19.5.7.1}
  \end{equation}
  은 전단사이다.

  \oldpage[0\textsubscript{IV-1}]{100}

  \item[\rm(iii)] 역으로 $B$가 준허용환이고 열
  $(\mathfrak J^n)$이 $0$으로 수렴하며, 모든 $\lambda$에 대하여
  $(\mathfrak J/\mathfrak J^2)\otimes_C C_\lambda$가 사영
  $C_\lambda$-가군이고 준동형 (19.5.7.1)이 전단사라고 하자. 그러면
  $B$는 형식적으로 매끄러운 $A$-대수이다.
\end{enumerate}
\end{corollary}

\begin{proof}
실제로 $\mathfrak J/\mathfrak J^2$의 위상과
$\mathbf S_C^n(\mathfrak J/\mathfrak J^2)$의 위상도 $B$의 위상에서
얻어진다 (19.5.1). 따라서 이 결론들은 이러한 위상에 특수화한
(19.5.3)과 (19.1.4), (19.2.4)의 판정법에서 곧바로 따른다.
\end{proof}

\begin{remark}[19.5.8]
\label{0.19.5.8-ko}
$\mathfrak J/\mathfrak J^2$가 유한형 $C$-가군이고, $B$의 위상에서 얻은
몫 위상에 대하여 $C$가 자리스키 환이라고 하자. $\mathfrak r$를
$C$의 정의 아이디얼이라 하여, $\mathfrak r^n$들이 $C$에서 $0$의
근방 기본계를 이루게 하자. 그러면 (19.5.7)의 (i), (ii)의 결론은
각각 다음으로 바꾸어 쓸 수 있다.
\begin{enumerate}
  \item[\rm(i')] $\mathfrak J/\mathfrak J^2$는 사영 $C$-가군이다.
  \item[\rm(ii')] 표준 준동형
  $\varphi:\mathbf S_C^*(\mathfrak J/\mathfrak J^2)
  \to\gr_{\mathfrak J}^*(B)$는 전단사이다.
\end{enumerate}
실제로 (i')이 (19.5.7)의 (i)의 결론을 함의하는 것은 명백하다.
역으로 모든 $\lambda$에 대하여
$(\mathfrak J/\mathfrak J^2)\otimes_C C_\lambda$가 사영
$C_\lambda$-가군이면, 모든 $n$에 대하여
$(\mathfrak J/\mathfrak J^2)\otimes_C(C/\mathfrak r^n)$은 사영이고
따라서 평탄한 $(C/\mathfrak r^n)$-가군이다. 그러므로
$\mathfrak J/\mathfrak J^2$는 평탄한 $C$-가군이다
($\mathbf 0_{\mathrm{III}}$, 10.2.2). 이것은 유한 표시 가군이므로
사영 가군이다 (Bourbaki, \emph{Alg. comm.}, chap.~II, \textsection5,
n\textsuperscript{o}~2, cor.~2 du th.~1). 한편 $C$-가군
$\mathbf S_C^n(\mathfrak J/\mathfrak J^2)$와
$\gr_{\mathfrak J}^n(B)$는 유한형이다. $C$가 자리스키 환일 때에는
$\varphi_n$이 전단사라는 것과 $\widehat\varphi_n$이 전단사라는 것이
동치임이 알려져 있다 (Bourbaki, \emph{Alg. comm.}, chap.~III,
\textsection3, n\textsuperscript{o}~5, prop.~9). 따라서 (ii)와 (ii')은
동치이다.
\end{remark}

\subsection{체 위의 대수의 경우}
\label{subsection:0.19.6-ko}

\begin{theorem}[19.6.1 (Cohen)]
\label{0.19.6.1-ko}
$k$를 체, $K$를 $k$의 확대체라 하고, $k$와 $K$에 이산 위상을 주자.
$K$가 형식적으로 매끄러운 $k$-대수일 필요충분조건은 $K$가 $k$의
분리 확대체인 것이다.
\end{theorem}

\begin{proof}
이 조건의 필요성은 (19.6.5.1)에서 증명할 것이며, 물론 그때까지는
사용하지 않는다. 여기서는 조건이 \emph{충분함}만을 증명한다. 두
경우로 나눈다.

\emph{I. --- $K$가 $k$의 유한형 분리 확대체인 경우.} 이때
(Bourbaki, \emph{Alg.}, chap.~V, \textsection9,
n\textsuperscript{o}~3, th.~2)에 따라 $K$의 순수 초월 부분확대체
$K'=k(T_1,\ldots,T_n)$가 존재하여, $K$는 $K'$의 유한 분리 대수
확대체가 된다. 따라서 (19.3.5, (iii))을 고려하면 $K=K'$인 경우와
$K$가 $k$의 유한 대수 확대체인 경우만 다루면 된다. 첫 번째 경우에는
$A=k[T_1,\ldots,T_n]$이 형식적으로 매끄러운 $k$-대수이므로
(19.3.3), $S=A-\{0\}$에 대하여
$k(T_1,\ldots,T_n)=S^{-1}A$도 그러하다 (19.3.5, (iv)). 두 번째
경우에는 임의의 $(K,K)$-쌍가군 $L$과 모든 $n$에 대하여 호흐실트
코호몰로지군 $\mathrm H_k^n(K,L)$이 $0$이다. 실제로 텐서곱
$C=K\otimes_k K$를 $k$-대수로 보면 $L$은 왼쪽 $C$-가군이고,
$K$도 $(K,K)$-쌍가군으로 보았을 때 호흐실트 코호몰로지군
$\mathrm H_k^n(K,L)$은 $\operatorname{Ext}_C^n(K,L)$과 같다.

\oldpage[0\textsubscript{IV-1}]{101}

따라서 $K$는 왼쪽 $C$-가군이기도 하다 (M, IX, 4). 그런데 $K$는
$k$의 유한 분리 확대체이므로 $K\otimes_k K$는 $k$의 확대체들의
직접곱이고, 그 가운데 하나는 $K$ 자체이다 (Bourbaki, \emph{Alg.},
chap.~VIII, \textsection8, prop.~3). 따라서 $K$는 사영 $C$-가군이며,
위 주장이 따른다. 그러므로 핵이 $L$인 $K$의 모든 $k$-확장은
$k$-자명하고, 더욱이 모든 가환 $k$-확장도 자명하다. 따라서 이
경우의 정리가 성립한다 (19.4.4).

\emph{II. --- 일반적인 경우.} --- (18.4.5)의 표기를 쓰면, 모든
$K$-벡터공간 $L$에 대하여
$\operatorname{Hom}_K(\mathrm H_2(P'_\bullet),L)=0$임을 보이면 된다.
이는 명백히 $\mathrm H_2(P'_\bullet)=0$이라는 뜻이다. $K$가 $k$의
부분확대체 $K^{(\alpha)}$들의 증가 여과족의 합집합이면,
$k$-가군의 범주에서 함자 $\varinjlim$이 텐서곱과 가환하므로
$P'_\bullet$은 대응하는 복합체 $P'^{(\alpha)}_\bullet$들의
\emph{귀납극한}이다. 이 범주에서 $\varinjlim$ 함자가 완전하므로
$\mathrm H_2(P'_\bullet)
=\varinjlim\mathrm H_2(P'^{(\alpha)}_\bullet)$이다. $K$가 분리
확대체이면 그 모든 부분확대체도 그러하고, $K$는
유한형 부분확대체들의 합집합이다. 따라서 증명의 첫 부분에 의해
$K$는 형식적으로 매끄러운 $k$-대수이다. C.Q.F.D.
\end{proof}

\begin{corollary}[19.6.2]
\label{0.19.6.2-ko}
$A$를 분리이고 완비인 국소환, $K$를 그 잉여체, $k$를 $K$의
부분체라 하자. $K$가 $k$의 분리 확대체라고 가정한다. 그러면 $k$를
포함하는 $A$의 부분체 $K'$가 존재하여, 표준 준동형 $A\to K$를
$K'$에 제한한 것은 $K'$에서 $K$로의 동형사상이다.
\end{corollary}

\begin{proof}
$\mathfrak m$을 $A$의 극대 아이디얼이라 하자. 가정과 (19.6.1)에 의해
$K$는 형식적으로 매끄러운 $k$-대수이다. (19.3.10)에서 $C$를 $A$로,
$\mathfrak J$를 $\mathfrak m$으로 바꾸어 적용하면
$K=A/\mathfrak m$의 항등 자기동형사상이
$K\xrightarrow{u}A\to A/\mathfrak m$으로 인수분해된다. 여기서 $u$는
$k$-준동형이므로, $K$에서 $k$를 포함하는 $A$의 부분체 $K'$로의
$k$-동형사상이다.
\end{proof}

\begin{corollary}[19.6.3]
\label{0.19.6.3-ko}
$A$를 뇌터 완비 국소환, $\mathfrak m$을 그 극대 아이디얼, $k$를 그
잉여체라 하자. 다음 조건들은 동치이다.
\begin{enumerate}
  \item[a)] $A$는 부분체를 포함한다.
  \item[b)] $p$가 $k$의 표수이면 $pA=0$이다.
  \item[c)] 어떤 체 $K$가 존재하여 $A$는 형식적 멱급수환
  $B=K[[T_1,\ldots,T_n]]$의 몫환과 동형이다.
\end{enumerate}
이 조건들이 성립할 때에는 $A$가 $B/\mathfrak b$와 동형이고
$\mathfrak b$가 $B$의 극대 아이디얼의 제곱에 포함되도록 할 수 있다.
\end{corollary}

\begin{proof}
$A$는 국소환이므로 $A\neq0$이다. 따라서 $c)$가 $a)$를 함의하는 것은
곧바로 알 수 있다. 실제로 $K$는 체이고 합성 준동형
$K\xrightarrow{\varphi}B\to A$는 영 준동형이 아니므로 단사이다. $a)$가
$b)$를 함의함을 보려면, $K$가 $A$의 부분체일 때 $K$가 $k$의
부분체와 동형이고 따라서 표수 $p$가 같다는 것만 관찰하면 된다.
$K$에서 $p\mathbin{.}1=0$이고 따라서 $A$에서도 그러하므로 $pA=0$이다.
역으로 $b)$는 $a)$를 함의한다. 표준 준동형들의 합성
$\mathbf Z\xrightarrow{j}A\to k$의 핵은 $p\mathbf Z$이고
$j(p\mathbf Z)=0$이므로 $\operatorname{Ker}(j)=p\mathbf Z$이다. 따라서
$A$는 $\mathbf Z/p\mathbf Z$와 동형이고 $\mathfrak m$과의 교집합이
$0$인 부분환 $A'$을 포함한다. $p>0$이면 $A'$은 체이다. 그렇지 않으면

\oldpage[0\textsubscript{IV-1}]{102}

$A'$의 모든 비영 원소가 $A$에서 가역이므로 $A'$의 분수체
($\mathbf Q$와 동형)도 $A$에 포함된다.

마지막으로 $a)$가 $c)$를 함의함을 증명하자. 조건 $a)$에 따라 $A$는
$k$의 소체와 동형인 소부분체 $k_0$를 포함하고, 따라서 $k$는
$k_0$의 분리 확대체이다. (19.6.2)를 적용하면 $k$에서 $A$의 부분체
$K$로의 동형사상 $f$가 존재한다. 한편 $(x_i)_{1\leq i\leq n}$을
$\mathfrak m$의 원소들의 계라 하여, $x_i$들의 $\mathfrak m^2$ 법의
류들이 $\mathfrak m/\mathfrak m^2$의 $k$-기저를 이루게 하자. $A$는
완비이므로 연속 환 준동형
$u:k[[T_1,\ldots,T_n]]\to A$가 존재하여 $k$ 위에서는 $f$와 같고 모든
$i$에 대하여 $u(T_i)=x_i$이다. $x_i$들을 이렇게 택했으므로 이
준동형은 전사이다 (Bourbaki, \emph{Alg. comm.}, chap.~III,
\textsection2, n\textsuperscript{o}~9, prop.~11).
\end{proof}

\begin{theorem}[19.6.4]
\label{0.19.6.4-ko}
$k$를 체, $A$를 국소 뇌터 $k$-대수, $\mathfrak m$을 그 극대
아이디얼, $K=A/\mathfrak m$을 그 잉여체라 하고, $A$에
$\mathfrak m$-준아딕 위상을 주자. $K$가 $k$의 분리 확대체라고
가정한다. 그러면 다음 조건들은 동치이다.
\begin{enumerate}
  \item[a)] $A$는 형식적으로 매끄러운 $k$-대수이다.
  \item[b)] $A$의 완비화 $\widehat A$는 형식적 멱급수환
  $K[[T_1,\ldots,T_n]]$과 $k$-동형이다. 여기서 $k$-대수 구조는 합성
  준동형
  \[
    k\xrightarrow{\varphi}K\xrightarrow{\psi}K[[T_1,\ldots,T_n]]
  \]
  으로 정의하며, $\psi$는 표준 포함사상이고 $\varphi$는 $K$의
  $k$-확대체 구조를 정의하는 준동형이다.
  \item[b')] $\widehat A$는 형식적 멱급수환
  $K[[T_1,\ldots,T_n]]$과 환으로서 동형이다.
  \item[c)] $A$는 정칙 국소환이다.
\end{enumerate}
\end{theorem}

\begin{proof}
$a)$가 $b)$를 함의한다는 것은 (19.6.1)을 고려하여 (19.5.4)의
$a)$와 $c)$의 동치를 $A=k$, $B=A$, $\mathfrak J=\mathfrak m$,
$C=K=A/\mathfrak m$으로 치환해 적용한 특수한 경우이다. $b)$가
$b')$를 함의하는 것은 명백하고, $b')$는 (17.1.1)에 의해 $c)$를
함의한다. 마지막으로 $c)$가 성립하면 (17.1.1, $a)$)와 (19.5.4)의
$b)$와 $a)$의 동치를 같은 치환으로 적용하여 $A$가 형식적으로
매끄러운 $k$-대수임을 얻는다.
\end{proof}

\begin{corollary}[19.6.5]
\label{0.19.6.5-ko}
$k$를 체, $A$를 국소 뇌터 $k$-대수, $\mathfrak m$을 그 극대
아이디얼이라 하고, $A$에 $\mathfrak m$-준아딕 위상을 주자. $A$가
형식적으로 매끄러운 $k$-대수라고 가정한다. 그러면 $A$는 $k$ 위에서
기하학적으로 정칙이다. 다시 말해 (cf.~$\mathbf{IV}$, 6.7.6), $k$의
모든 유한 확대체 $k'$에 대하여 반국소환 $A'=A\otimes_k k'$은
정칙이다 (17.3.6). 더 나아가 $K$가 $A$의 잉여체이면 $\widehat A$는
형식적 멱급수환 $K[[T_1,\ldots,T_n]]$과 동형이다.
\end{corollary}

\begin{proof}
$\widehat A'=\widehat A\otimes_k k'$이므로, (19.3.6)과 $A'$의 극대
아이디얼들에 대응하는 국소환들에 적용한 (17.1.5)에 따라 $A$가
완비인 경우만 다루면 된다. 그러면 $A'$은 형식적으로 매끄러운
$k'$-대수이고 (19.3.5, (iii)), 형식적으로 매끄러운 국소
$k'$-대수들의 직접곱이다 (19.3.5, (v)). 따라서 $A$가
\emph{정칙}임만 증명하면 된다. $k_0$를 $k$의 소체라 하자. $k$는
$k_0$의 분리 확대체이므로 형식적으로 매끄러운 $k_0$-대수이다
(19.6.1). 가정과 (19.3.5, (ii))에 의해 $A$도 그러하다. $A$의
잉여체 $K$도 $k_0$의 분리 확대체이므로 $A$와 $k_0$에 (19.6.4)를
적용할 수 있고, 결론이 따른다.
\end{proof}

\oldpage[0\textsubscript{IV-1}]{103}

\begin{corollary}[19.6.5.1]
\label{0.19.6.5.1-ko}
$A$를 국소 뇌터환, $k$를 $A$의 부분체라 하고, $A$가 준아딕 위상에
대하여 형식적으로 매끄러운 $k$-대수라고 하자. 그러면
$k\subset K\subset A$를 만족하는 모든 체 $K$는 $k$의 분리
확대체이다.
\end{corollary}

\begin{proof}
$k$의 모든 유한 확대체 $k'$에 대하여 $K\otimes_k k'$은
$A'=A\otimes_k k'$의 부분환과 동일시된다. $A'$은 정칙환이므로
축소환이고, 따라서 $K\otimes_k k'$도 축소환이다. 그러므로 $K$는
$k$의 분리 확대체이다 (Bourbaki, \emph{Alg.}, chap.~VIII,
\textsection7, n\textsuperscript{o}~3, th.~1).
\end{proof}

이것으로 (19.6.1)의 명제에 있는 조건이 \emph{필요함}도 증명되었다.

\begin{remark}[19.6.5.2]
\label{0.19.6.5.2-ko}
$K$가 $k$의 \emph{비분리} 확대체이면 형식적 멱급수환
$K[[T_1,\ldots,T_n]]$은 (19.6.4)에서 상기한 통상적인 $k$-대수
구조에 대하여 형식적으로 매끄러운 $k$-대수가 \emph{아니다}. 한편
형식적으로 매끄러운 $k$-대수 $A$ 가운데 잉여체 $K$가 $k$의
\emph{비분리} 확대체인 완비 국소 뇌터환들도 있다. 예를 들어
$B=k[T_1,\ldots,T_n]$의 극대 아이디얼 $\mathfrak n$ 가운데
$B/\mathfrak n$이 $k$의 유한 비분리 대수 확대체가 되는 것에 대한
국소화들의 완비화가 그러하다. 따라서 이러한 환은 (19.6.5)에 따라
$K[[T_1,\ldots,T_n]]$과 환으로서는 동형이지만, 그것과 $k$-동형일
수는 없다.

어떤 경우에는 따름정리 (19.6.5)를 다음과 같이 더 정밀하게 할 수
있다.
\end{remark}

\begin{proposition}[19.6.6]
\label{0.19.6.6-ko}
$k$를 체, $A$를 국소 뇌터 $k$-대수, $\mathfrak m$을 그 극대
아이디얼이라 하고, $A$에 $\mathfrak m$-준아딕 위상을 주자. $A$의
잉여체 $K=A/\mathfrak m$에 대하여 다음 성질을 만족하는 $k$의 유한
순수 비분리 확대체 $k_0$가 존재한다고 가정한다. 즉 국소환
$K\otimes_k k_0$의 잉여체가 $k_0$의 분리 확대체라고 하자. 나중에는
이 성질을 $K$가 $k$ 위에서 \emph{유한 순수 비분리 중복도}를 갖는다고
표현할 것이다 ($\mathbf{IV}$, 4.7.4). 또한 $K$가 $k$의 \emph{유한형
확대체}이면 이 조건이 언제나 성립함을 볼 것이다. 그러면 다음
조건들은 동치이다.
\begin{enumerate}
  \item[a)] $A$는 형식적으로 매끄러운 $k$-대수이다.
  \item[b)] $k$의 어떤 유한 순수 비분리 확대체 $k'$에 대하여
  $\widehat A\otimes_k k'$은 형식적 멱급수 대수
  $K'[[T_1,\ldots,T_n]]$과 $k'$-동형이다. 여기서 $K'$은 $k'$의 분리
  확대체이다.
  \item[b')] $k$의 어떤 유한 확대체 $k'$에 대하여 반국소환
  $\widehat A\otimes_k k'$은 형식적 멱급수 대수
  $K'[[T_1,\ldots,T_n]]$과 $k'$-동형인 국소 성분을 적어도 하나
  갖는다. 여기서 $K'$은 $k'$의 분리 확대체이다.
  \item[c)] $A$는 $k$ 위에서 기하학적으로 정칙이다 (19.6.5).
\end{enumerate}
\end{proposition}

\begin{proof}
먼저 $k'$이 $k$의 유한 순수 비분리 확대체이면
$A'=A\otimes_k k'$에는 $\mathfrak m$ 위에 놓이는 극대 아이디얼이
하나뿐임을 주목하자. 이 아이디얼은 어떤 적당한 $h$에 대하여
$p^h$제곱이 $\mathfrak m A'$에 속하는 원소들로 이루어진다. 여기서
$p$는 $k$의 표수 지수이다 (Bourbaki, \emph{Alg. comm.}, chap.~V,
\textsection2, n\textsuperscript{o}~3, lemme~4). 따라서 $A'$은
국소환이고
$K\otimes_k k'=(A\otimes_k k')/(\mathfrak m\otimes_k k')$도
국소환이며, 이 두 국소환의 잉여체는 같다. 한편 $K$가 $k$의 분리
확대체이면 $k$의 모든 유한 확대체 $k''$에 대하여
$K\otimes_k k''$은 체들의 직접곱이다 (Bourbaki, \emph{Alg.},
chap.~VIII, \textsection7, n\textsuperscript{o}~3, cor.~1 du th.~1).
따라서 $\mathfrak m\otimes_k k''$은 $A\otimes_k k''$의 근기이고,
$K\otimes_k k''$의 체 성분들은 $A\otimes_k k''$의 극대 아이디얼들에
대한 잉여체들이다.

\oldpage[0\textsubscript{IV-1}]{104}

더 나아가 이 체들은 $k''$의 분리 확대체이다. 실제로 $k''$의 모든
확대체 $k_1$에 대하여
$K\otimes_k k_1=(K\otimes_k k'')\otimes_{k''}k_1$은 축소환이다
(\emph{loc. cit.}, th.~1). 이 사실들을 $k$ 위에서 유한 순수 비분리
중복도를 갖는 체 $K$에 적용하면, $k_0$의 \emph{모든} 유한 확대체
$k''$에 대하여 $(K\otimes_k k'')_{\mathrm{red}}$가 $k''$의 분리
대수임을 알 수 있다.

이제 (19.6.6)을 증명하자. $b)$가 $b')$를 함의하는 것은 자명하다.
$b')$가 $a)$를 함의함을 보이자. $B'=K'[[T_1,\ldots,T_n]]$로 놓는다.
$K'$이 $k'$의 분리 확대체라는 가정과 앞의 사실들에 의해, $k'$의
모든 유한 확대체 $k''$에 대하여 완비 반국소환
$B'\otimes_{k'}k''$, 즉 형식적 멱급수환
$(K'\otimes_{k'}k'')[[T_1,\ldots,T_n]]$의 성분들은
$K_i''[[T_1,\ldots,T_n]]$들이다. 여기서 $K_i''$들은
$K'\otimes_{k'}k''$의 체 성분들이다. $B'\otimes_{k'}k''$의 국소
성분들은
\[
  \widehat A\otimes_k k''
  =(\widehat A\otimes_k k')\otimes_{k'}k''
\]
의 국소 성분들이기도 하다. 따라서 $k'$을 그 임의의 유한 확대체로
바꾸어도 가정 $b')$는 여전히 성립한다.

특히 $k'$을 \emph{준갈루아} 확대체로 택할 수 있다. 즉 $k'$은 $k$의
어떤 순수 비분리 확대체 $k'_0$의 갈루아 확대체이다. 이때
$A'=\widehat A\otimes_k k'$은
$A'_0\otimes_{k'_0}k'$으로 쓸 수 있으며,
$A'_0=\widehat A\otimes_k k'_0$은 완비 국소환이다. $k'/k'_0$의
갈루아 군은 $A'$의 국소 성분들에 추이적으로 작용하므로, 이 성분들은
서로 갈루아 켤레이다 (Bourbaki, \emph{Alg.}, chap.~VIII,
\textsection8, prop.~4). 가정에 의해 그 가운데 하나는 $k'$의 분리
확대체 $K'$에 대한 형식적 멱급수환 $K'[[T_1,\ldots,T_n]]$과
$k'$-동형이고, 따라서 (19.6.4)에 의해 형식적으로 매끄러운
$k'$-대수이다. 그 갈루아 켤레들도 $k'$ 위에서 형식적으로 매끄럽고,
따라서 그 직접곱 $A'$도 형식적으로 매끄러운 $k'$-대수이다
(19.3.5, (v)). $k'$은 $k'_0$ 위에서 충실하게 평탄하고 사영인 유한형
가군이므로, (19.4.6)에 의해 $A'_0$는 형식적으로 매끄러운
$k'_0$-대수이다. 끝으로 $k'_0$도 $k$ 위에서 같은 성질들을 가지므로
(19.4.6)을 다시 적용하여 $\widehat A$가 형식적으로 매끄러운
$k$-대수임을 얻는다. 따라서 (19.3.6)에 의해 $A$도 형식적으로
매끄러운 $k$-대수이다.

$a)$가 $c)$를 함의함은 이미 (19.6.5)에서 보았다. 이제 $c)$가 $b)$를
함의함만 확인하면 된다. $c)$가 성립하면 $k_0$를 포함하는 $k$의 모든
유한 순수 비분리 확대체 $k'$에 대하여 $A\otimes_k k'$은 정칙
국소환이고, 그 잉여체 $K'$은 $k'$의 분리 확대체이다. 따라서
(19.6.4)에 의해 그 완비화 $\widehat A\otimes_k k'$은 형식적
멱급수환 $K'[[T_1,\ldots,T_n]]$과 $k'$-동형이다.
\end{proof}

\begin{remarks}[19.6.7]
\label{0.19.6.7-ko}
\begin{enumerate}
  \item[\rm(i)] $K$가 $k$ 위에서 유한 순수 비분리 중복도를 갖는다는
  가정은 $b')\Rightarrow a)$의 증명에서만 사용되었다.
  \item[\rm(ii)] 나중에 (22.5.8)에서 $K/k$에 관한 \emph{아무 가정
  없이도} $a)$와 $c)$가 동치임을 증명할 것이다.
\end{enumerate}
\end{remarks}

\subsection{국소 준동형의 경우; 존재성과 유일성 정리}
\label{subsection:0.19.7-ko}

이 절에서 반국소환을 위상환으로 볼 때에는 언제나 그 근기
$\mathfrak r$에 대한 $\mathfrak r$-준아딕 위상을 주는 것으로 이해한다.
따라서 국소환 사이의 모든 국소 준동형은 자동으로 \emph{연속}이다.

\begin{theorem}[19.7.1]
\label{0.19.7.1-ko}
$A$, $B$를 국소 뇌터환, $\mathfrak m$, $\mathfrak n$을 각각의 극대
아이디얼, $k=A/\mathfrak m$을 $A$의 잉여체라 하자. $A$와 $B$에 각각
$\mathfrak m$-준아딕 위상과 $\mathfrak n$-준아딕 위상을 준다고
가정한다. $\varphi:A\to B$를 국소 준동형이라 하고
$B_0=B\otimes_A k$로 놓자. 다음 성질들은 동치이다.
\begin{enumerate}
  \item[a)] $B$는 형식적으로 매끄러운 $A$-대수이다.

  \oldpage[0\textsubscript{IV-1}]{105}

  \item[b)] $B$는 평탄한 $A$-가군이고, 몫 위상을 준
  $B_0=B/\mathfrak mB$는 형식적으로 매끄러운 $k$-대수이다.
\end{enumerate}
\end{theorem}

\begin{proof}
증명을 여러 단계로 나눈다.

\medskip
\noindent\phantomsection\label{0.19.7.1.1-ko}\textbf{(19.7.1.1)}\enspace
먼저 $b)$가 $a)$를 함의함을 증명하자. $\mathfrak J=\mathfrak m$으로
놓고 판정법 (19.4.7)을 적용한다. $b)$의 두 번째 가정에 의해 $B$가
형식적으로 사영인 $A$-가군임을 보이기만 하면 된다. 가정에 따라 모든
$h>0$에 대하여 $B/\mathfrak m^hB$는 평탄한
$(A/\mathfrak m^h)$-가군이다 ($\mathbf 0_{\mathrm{III}}$, 10.2.1).
$\mathfrak m^h$들이 $A$에서 $0$의 근방 기본계를 이루고
\[
  (B/\mathfrak m^hB)\otimes_{A/\mathfrak m^h}k=B_0
\]
이므로, $A$와 $B$를 각각 $A/\mathfrak m^h$와
$B/\mathfrak m^hB$로 바꿀 수 있다. 따라서 $A$가 \emph{아르틴환}이고
그 위상이 이산인 경우만 생각하면 된다. $B_0$는 형식적으로 매끄러운
$k$-대수이므로 정칙환이다 (19.6.5). $(x_i^0)_{1\leq i\leq n}$을
$B_0$의 정칙 매개변수계라 하고 (17.1.6), 각 $i$에 대하여 $B$의 원소
$x_i$를 택하여 $B_0=B/\mathfrak mB$에서 그 상이 $x_i^0$이 되게 하자.
$x_i^0$들이 $B_0$의 극대 아이디얼
$\mathfrak n_0=\mathfrak n/\mathfrak mB$를 생성하므로, 아이디얼
\[
  \mathfrak Q_m=\sum_{i=1}^n(x_i^0)^mB_0 \qquad (m>0)
\]
들은 $B_0$에서 $0$의 근방 기본계를 이룬다. 실제로
$\mathfrak Q_m\subset\mathfrak n_0^m$이고
$\mathfrak Q_m\supset\mathfrak n_0^{mn}$이다. 모든 $m>0$에 대하여
$\mathfrak J_m=\sum_{i=1}^n x_i^mB$로 놓자. 명백히
$\mathfrak n=\mathfrak J_1+\mathfrak mB$이다. 어떤 $h>0$에 대하여
$\mathfrak m^h=0$이므로 $m>h$일 때
$\mathfrak n^m\subset\mathfrak J_1^{m-h}
\subset\mathfrak n^{m-h}$이다. 위에서
$\mathfrak J_m\supset\mathfrak J_1^{mn}$임을 보았으므로
$\mathfrak J_m$들은 $B$에서 $0$의 근방 기본계를 이룬다. 따라서
$B/\mathfrak J_m$들이 자유 $A$-가군임을 증명하면 충분하다. 이는
그들이 평탄한 $A$-가군임을 증명하는 것과 같다
($\mathbf 0_{\mathrm{III}}$, 10.1.3). 그런데 $(x_i^0)$가 $B_0$의 극대
아이디얼에 속하는 $B_0$-정칙열이라는 가정은 모든 $m>0$에 대하여
$((x_i^0)^m)_{1\leq i\leq n}$에도 같은 성질을 함의한다 (15.1.20).
따라서 결론은 (15.1.16, $b)$와 $c)$)에서 따른다.
\end{proof}

\begin{lemma}[19.7.1.2]
\label{0.19.7.1.2-ko}
$A$를 위상환, $B$, $C$를 국소 뇌터환인 두 위상 $A$-대수라 하고,
$\mathfrak m$, $\mathfrak n'$을 각각의 극대 아이디얼이라 하자. 더
나아가 $C$는 완비이고 $B$의 잉여체 $B/\mathfrak m$은 유한 생성
$A$-가군이라고 가정한다. $E$를 완비 텐서곱
$B\widehat\otimes_A C$라 하자. 그러면 다음이 성립한다.
\begin{enumerate}
  \item[(i)] $E$는 완비 반국소 뇌터환이다.
  \item[(ii)] 아이디얼 $\mathfrak mE$는 $E$의 근기에 포함되고, 모든
  $h>0$에 대하여 $E/\mathfrak m^hE$는
  $(B/\mathfrak m^h)\otimes_A C$와 동형이다.
  \item[(iii)] $C$가 평탄한 $A$-가군이면 $E$는 평탄한 $B$-가군이다.
\end{enumerate}
\end{lemma}

\begin{proof}
정의상 $E$는 아이디얼
$\operatorname{Im}(\mathfrak m^i\otimes_A C)+
\operatorname{Im}(B\otimes_A\mathfrak n'{}^j)$들이 정의하는 위상에
대한 $B\otimes_A C$의 분리 완비화이다
($\mathbf 0_{\mathrm I}$, 7.7.5). 다음과 같이 놓자.
\[
  \mathfrak r=\operatorname{Im}(\mathfrak m\otimes_A C)+
  \operatorname{Im}(B\otimes_A\mathfrak n').
\]
그러면
\[
  \mathfrak r^{2i}\subset
  \operatorname{Im}(\mathfrak m^i\otimes_A C)+
  \operatorname{Im}(B\otimes_A\mathfrak n'{}^i)
  \subset\mathfrak r^i
\]
이므로 $E$는 $\mathfrak r$-준아딕 위상에 대한 $B\otimes_A C$의 분리
완비화이기도 하다. 가정에 의해
$(B\otimes_A C)/\mathfrak r=(B/\mathfrak m)\otimes_A
(C/\mathfrak n')$은 유한 생성 $(C/\mathfrak n')$-가군이고, 따라서
아르틴환이다. 또한 $\mathfrak r/\mathfrak r^2$은
\[
  ((\mathfrak m/\mathfrak m^2)\otimes_A C)
  \oplus(B\otimes_A(\mathfrak n'/\mathfrak n'{}^2))
\]
의 몫이므로 유한 생성 $(B\otimes_A C)$-가군이다. 따라서
($\mathbf 0_{\mathrm I}$, 7.2.11)을 적용하면 $E$가 \emph{뇌터환}임을
알 수 있다. 더 나아가 $E/\mathfrak rE$는
$(B\otimes_A C)/\mathfrak r$와 동형이고 아르틴환이므로 $E$는
\emph{반국소환}이다.

이제
\[
  E\simeq
  \varprojlim_{i,j}((B/\mathfrak m^i)\otimes_A(C/\mathfrak n'{}^j))
\]
이고, 이중 사영극한 정리에 의해
\[
  E\simeq\varprojlim_i\left(\varprojlim_j
  ((B/\mathfrak m^i)\otimes_A(C/\mathfrak n'{}^j))\right).
\]

\oldpage[0\textsubscript{IV-1}]{106}

그런데
$\varprojlim_j((B/\mathfrak m^i)\otimes_A(C/\mathfrak n'{}^j))$는
$(B/\mathfrak m^i)\otimes_A C$의 분리 완비화이다. $C$가 완비이고
$B/\mathfrak m^i$가 유한 생성 $A$-가군이므로 이는
$(B/\mathfrak m^i)\otimes_A C$ 자체이다. 여기서
$\mathfrak m/\mathfrak m^i$는 유한 생성 $(B/\mathfrak m)$-가군이므로
$B/\mathfrak m^i$가 유한 생성 $A$-가군이라는 사실을 사용했다
($\mathbf 0_{\mathrm I}$, 7.3.6). 따라서
\[
  E=\varprojlim_i((B/\mathfrak m^i)\otimes_A C).
\]
모든 정수 $h>0$에 대하여 $\mathfrak m^hE$는 $E$의 아이디얼이므로
$E$에서 닫혀 있고 ($\mathbf 0_{\mathrm I}$, 7.3.5), 따라서 완비이다.
한편 이는
\[
  \varprojlim_i
  (\operatorname{Im}(\mathfrak m^h/\mathfrak m^{h+i})\otimes_A C)
\]
에서 명백히 조밀하므로 이 사영극한과 같다. 더 나아가 여기서 생각한
모든 사영계의 전이 준동형은 \emph{전사}이다. 따라서
($\mathbf 0_{\mathrm{III}}$, 13.2.2)에 의해
\[
  E/\mathfrak m^hE\simeq(B/\mathfrak m^h)\otimes_A C
  =(B\otimes_A C)/\operatorname{Im}(\mathfrak m^h\otimes_A C).
\]
특히 $\operatorname{Im}(\mathfrak m\otimes_A C)\subset\mathfrak r$이므로
$\mathfrak mE\subset\mathfrak rE$이고, 따라서 $\mathfrak mE$는 $E$의
근기에 포함된다. 마지막으로 $C$가 평탄한 $A$-가군이면 모든 $h>0$에
대하여
$(B/\mathfrak m^h)\otimes_A C=E/\mathfrak m^hE$는 평탄한
$(B/\mathfrak m^h)$-가군이다. $B$와 $E$가 뇌터환이고
$\mathfrak mE$가 $E$의 근기에 포함되므로
($\mathbf 0_{\mathrm{III}}$, 10.2.2)에 의해 $E$는 평탄한
$B$-가군이다.
\end{proof}

\begin{lemma}[19.7.1.3]
\label{0.19.7.1.3-ko}
$A$를 국소 뇌터환, $\mathfrak m$을 그 극대 아이디얼, $k$를 그
잉여체라 하고, $B_0$를 $k$-대수라 하자. $B_0$가 완비 정칙 국소
뇌터환이라고 가정한다. 그러면 완비 국소 뇌터환인 위상 $A$-대수
$B$가 존재하여, $B$는 평탄한 $A$-가군이고 $B_0$는
$B\otimes_A k=B/\mathfrak mB$와 $k$-동형이다.
\end{lemma}

\begin{proof}
$\widehat A$는 $A$ 위에서 평탄하고 잉여체도 같으므로 $A$가
\emph{완비}인 경우만 다루면 된다.

$K$를 $B_0$의 잉여체라 하고 두 경우로 나눈다.

\medskip
\noindent I) $K$가 $k$의 \emph{분리} 확대체인 경우. (19.6.4)에 의해
$B_0$는 형식적 멱급수환 $K[[T_1,\ldots,T_n]]$과 $k$-동형이다.
$B_0=K$인 경우에는 이 보조정리가 이미 증명되었다
($\mathbf 0_{\mathrm{III}}$, 10.3.1). 완비 국소 뇌터환 $C$를 택하여
평탄한 $A$-가군이고 $C\otimes_A k$가 $K$와 동형이 되게 하자.
$n\geq1$이면 앞의 표기로 $B=C[[T_1,\ldots,T_n]]$를 택하면 충분하다.
실제로 $B$는 평탄한 $C$-가군이고 (Bourbaki, \emph{Alg. comm.},
chap.~III, \textsection3, n\textsuperscript{o}~4, cor.~3 du th.~1), 따라서
평탄한 $A$-가군이기도 하다. 또한 곧바로
\[
  C[[T_1,\ldots,T_n]]\otimes_A k
  \simeq(C/\mathfrak mC)[[T_1,\ldots,T_n]]=B_0
\]
를 얻는다.

\medskip
\noindent II) $K$의 \emph{표수}가 $p>0$인 경우. 그러면 $k$의 표수도
$p$이다. $P$를 소체 $\mathbf F_p$라 하고, $W(P)$를 잉여체가 $P$인
$p$-아딕 정수환 $\mathbf Z_p$라 하자. 이는 완비 국소환인 이산
값매김환이고, 따라서 정칙환이다. 먼저 $A$를 위상 $W(P)$-대수로
만드는 연속 환 준동형 $W(P)\to A$가 존재함을 보이자. 표준 준동형
$j:\mathbf Z\to A$에 대하여 가정에서
$j(p\mathbf Z)\subset\mathfrak m$이므로
$j^{-1}(\mathfrak m)=p\mathbf Z$이다. 따라서 $j$는
$\mathbf Z\to\mathbf Z_{p\mathbf Z}\xrightarrow{f}A$로 인수분해된다.
$f$는 국소 준동형이므로 연속이고, $A$가 완비이므로 연속성에 의해
원하는 준동형 $W(P)\to A$로 연장된다.

$k$는 $P$의 분리 확대체이므로 경우 I)에 의해 국소 준동형
$W(P)\to W(k)$가 존재한다. 여기서 $W(k)$는 완비 국소 뇌터환이고
평탄한 $W(P)$-가군이며,
$W(k)\otimes_{W(P)}P$는 $k$와 동형이다. 더 나아가 $W(P)$의
균등화원 $p$는 평탄성에 의해 $W(k)$-정칙 원소이고
($\mathbf 0_{\mathrm I}$, 6.3.4), $W(k)/pW(k)=k$이므로 $pW(k)$는
$W(k)$의 극대 아이디얼이다.

\oldpage[0\textsubscript{IV-1}]{107}

따라서 $W(k)$는 \emph{완비 이산 값매김환}이다 (Bourbaki,
\emph{Alg. comm.}, chap.~VI, \textsection3, n\textsuperscript{o}~5,
prop.~9). 또한 $k$는 $P$의 분리 확대체, 곧 형식적으로 매끄러운
$P$-대수이므로 (19.6.1), (19.7.1.1)에 의해 $W(k)$는 형식적으로
매끄러운 $W(P)$-대수이다. 연속 $W(P)$-준동형 $W(k)\to k$는
$W(k)\to A\to k$로 인수분해된다 (19.3.11). 이에 따라 $A$를 위상
$W(k)$-대수로 볼 수 있다.

이제 $B_0$를 $P$-대수로 보고 $W(P)$에 경우 I)을 적용하면, 완비 국소
뇌터환인 $W(P)$-대수 $B_P$가 존재하여 평탄한 $W(P)$-가군이고
$B_P\otimes_{W(P)}P$가 $B_0$와 $P$-동형이 된다. $W(k)$가 형식적으로
매끄러운 $W(P)$-대수라는 사실을 다시 사용하면 합성 준동형
$W(k)\to k\to B_0$는 $W(k)\to B_P\to B_0$로 인수분해된다
(19.3.11). 또한 $k=W(k)/pW(k)$이므로
\[
  B_P\otimes_{W(k)}k=B_P/pB_P
  =B_P\otimes_{W(P)}P=B_0.
\]
$B_P$가 평탄한 $W(k)$-가군임을 보이자. $W(k)$는 균등화원이 $p$인
이산 값매김환이므로 $B_P$가 꼬임 없는 $W(k)$-가군임을
확인하면 충분하다 ($\mathbf 0_{\mathrm I}$, 6.3.4). 이는 $p$가
$B_P$-정칙 원소라는 것과 같고, $B_P$가 평탄한 $W(P)$-가군이라는
사실에서 따른다 ($\mathbf 0_{\mathrm I}$, 6.3.4).

이제
\[
  B=B_P\widehat\otimes_{W(k)}A
\]
로 놓자. $A$의 잉여체는 $W(k)$의 잉여체와 같으므로 특히 유한 생성
$W(k)$-가군이다. (19.7.1.2)에 의해 먼저 $B$는 완비 반국소
뇌터환이고 $\mathfrak mB$는 $B$의 근기에 포함된다. 또한
$B/\mathfrak mB$는
$B_P\otimes_{W(k)}k=B_0$와 $k$-동형이므로 $B$는 실제로 국소환이다.
마지막으로 $B_P$가 평탄한 $W(k)$-가군이므로 (19.7.1.2)에 의해 $B$는
평탄한 $A$-가군이다. C.Q.F.D.
\end{proof}

\begin{lemma}[19.7.1.4]
\label{0.19.7.1.4-ko}
$A$를 환, $\mathfrak J$를 $A$의 아이디얼, $M$, $N$을
$\mathfrak J$-준아딕 위상에 대하여 분리인 두 $A$-가군이라 하자. 더
나아가 $N$은 $\mathfrak J$-준아딕 위상에 대하여 완비이고 $M$은
평탄한 $A$-가군이라고 가정한다. $u:N\to M$을 $A$-준동형이라 하자.
준동형
\[
  u\otimes1:N\otimes_A(A/\mathfrak J)
  \longrightarrow M\otimes_A(A/\mathfrak J)
\]
이 전단사이면 $u$도 전단사이다.
\end{lemma}

\begin{proof}
연관 등급 가군들은 $\mathfrak J$-준아딕 여과에 관하여 취한다. $M$과
$N$의 $\mathfrak J$-준아딕 위상에 관한 가정에 따라
\[
  \gr(u):\gr_\bullet(N)\longrightarrow\gr_\bullet(M)
\]
이 전단사임을 증명하면 충분하다 (Bourbaki, \emph{Alg. comm.},
chap.~III, \textsection2, n\textsuperscript{o}~8, cor.~3 du th.~1). 그런데
다음 가환 도식이 있다.
\[
\xymatrix{
  \gr_0(N)\otimes_{A_0}\gr_\bullet(A)
    \ar[r]^-{\gr_0(u)\otimes1}\ar[d]_{\varphi_N} &
  \gr_0(M)\otimes_{A_0}\gr_\bullet(A)\ar[d]^{\varphi_M} \\
  \gr_\bullet(N)\ar[r]_{\gr(u)} & \gr_\bullet(M)
}
\]

\oldpage[0\textsubscript{IV-1}]{108}

여기서 $A_0=A/\mathfrak J=\gr_0(A)$이고, $\varphi_M$과
$\varphi_N$은 표준 사상들이다
($\mathbf 0_{\mathrm{III}}$, 10.1.1.2). 가정에 따라
$\gr_0(u)$는 전단사이고 $\varphi_M$도 전단사이며
($\mathbf 0_{\mathrm{III}}$, 10.2.1), $\varphi_N$은 전사이다. 따라서
먼저 $\gr_0(u)\otimes1$이 전단사이고, 이어서 $\varphi_N$이 단사이므로
전단사이며, 마침내 $\gr(u)$가 전단사임을 얻는다.
\end{proof}

\begin{lemma}[19.7.1.5]
\label{0.19.7.1.5-ko}
$A$를 뇌터환, $\mathfrak J$를 $A$의 아이디얼, $B$, $B'$을 국소
뇌터환인 두 $A$-대수라 하자. 준동형 $A\to B$, $A\to B'$은 $A$의
$\mathfrak J$-준아딕 위상에 대하여 연속이라고 가정한다. 다음을
가정하자.
\begin{enumerate}
  \item[$1^{\mathrm o}$] $B$와 $B'$은 $\mathfrak J$-준아딕 위상에
  대하여 완비이다.
  \item[$2^{\mathrm o}$] $B$는 형식적으로 매끄러운 $A$-대수이다.
  \item[$3^{\mathrm o}$] $B'$은 평탄한 $A$-가군이다.
\end{enumerate}
$A_0=A/\mathfrak J$로 놓고
$u_0:B\otimes_A A_0\to B'\otimes_A A_0$를 $A_0$-동형사상이라 하자.
그러면 $u_0=u\otimes1$을 만족하는 $A$-동형사상 $u:B\to B'$이
존재한다. 특히 $B'$은 형식적으로 매끄러운 $A$-대수이고 $B$는
평탄한 $A$-가군이다.
\end{lemma}

\begin{proof}
$B_0=B\otimes_A A_0$, $B'_0=B'\otimes_A A_0$로 놓자.
$\mathfrak m$, $\mathfrak m'$을 각각 $B$, $B'$의 극대 아이디얼이라
하면 $\mathfrak JB\subset\mathfrak m$이고
$\mathfrak JB'\subset\mathfrak m'$이므로 $B$와 $B'$의
$\mathfrak J$-준아딕 위상은 분리이다. 또한 $\mathfrak JB'$은
$B'$의 $\mathfrak J$-준아딕 위상에서 닫혀 있다. 따라서
$\mathfrak J$-준아딕 위상에 대하여 연속인 합성 준동형
$B\to B_0\xrightarrow{u_0}B'_0$는
\[
  B\xrightarrow{u}B'\longrightarrow B'_0
\]
로 인수분해된다. 여기서 $u$는 연속 $A$-준동형이다 (19.3.10).
명백히 $u_0=u\otimes1$이고, $u_0$가 전단사라는 가정과 (19.7.1.4)에
의해 $u$도 전단사이다.
\end{proof}

\medskip
\noindent\phantomsection\label{0.19.7.1.6-ko}\textbf{(19.7.1.6)}\enspace
\emph{증명의 끝.} --- (19.7.1)의 증명을 끝내려면 $a)$가 $b)$를
함의함을 보이면 된다. $a)$가 $B_0$를 형식적으로 매끄러운
$k$-대수로 만든다는 것은 이미 알고 있다 (19.3.5, (iii)). 따라서
$B$가 평탄한 $A$-가군임만 증명하면 된다. 이는 $\widehat B$가 평탄한
$\widehat A$-가군임을 증명하는 것과 같다 (Bourbaki,
\emph{Alg. comm.}, chap.~III, \textsection5, n\textsuperscript{o}~4,
prop.~4). 또한 $\widehat B$는 형식적으로 매끄러운 $\widehat A$-대수이다
(19.3.6). 그러므로 $A$와 $B$가 완비인 경우만 생각하면 된다.
$B_0$는 형식적으로 매끄러운 $k$-대수이므로 정칙환이고 (19.6.5),
완비이다 ($\mathbf 0_{\mathrm I}$, 6.3.5). (19.7.1.3)에 따라 완비
국소 뇌터환인 $A$-대수 $B'$이 존재하여, $B'$은 평탄한 $A$-가군이고
$A\to B'$은 국소 준동형이며
$B_0\simeq B'_0=B'\otimes_A k$인 $k$-동형사상이 존재한다. 이제
$\mathfrak J$를 $A$의 극대 아이디얼로 하여 (19.7.1.5)를 적용하면
$B$는 $B'$과 $A$-동형이고, 따라서 평탄한 $A$-가군이다. C.Q.F.D.

\begin{theorem}[19.7.2]
\label{0.19.7.2-ko}
$A$를 뇌터 국소환, $\mathfrak J$를 $A$의 극대 아이디얼에 포함되는
아이디얼, $A_0=A/\mathfrak J$라 하자. 또 $B_0$를 완비 뇌터 국소환이라
하고, 국소 준동형 $A_0\to B_0$가 $B_0$를 형식적으로 매끄러운
$A_0$-대수로 만든다고 하자. 그러면 완비 뇌터 국소환 $B$, $B$를
평탄한 $A$-가군으로 만드는 국소 준동형 $A\to B$, 그리고
$A_0$-동형사상
\[
  u:B\otimes_A A_0\xrightarrow{\sim}B_0
\]
가 존재한다. $(B',u')$이 $(B,u)$와 같은 조건을 만족하는 쌍이면,
다음 도표를 가환하게 하는 $A$-동형사상
$v:B\xrightarrow{\sim}B'$이 존재한다:
\[
\xymatrix{
  B\otimes_A A_0\ar[rr]^{v\otimes1}_{\sim}\ar[dr]^{\sim}_u &&
  B'\otimes_A A_0\ar[dl]^{u'}_{\sim} \\
  & B_0
}
\]
\end{theorem}

\oldpage[0\textsubscript{IV-1}]{109}

\begin{proof}
$\mathfrak m$을 $A$의 극대 아이디얼이라 하자. 그러면
$\mathfrak m_0=\mathfrak m/\mathfrak J$는 $A_0$의 극대 아이디얼이고,
$A$와 $A_0$의 잉여체는 같은 체 $k$이다.
$B_{00}=B_0\otimes_{A_0}k$로 놓자. $B_{00}$은 형식적으로 매끄러운
$k$-대수이므로 (19.3.5, (iii)) 정칙 국소환이다 (19.6.5). 따라서
(19.7.1.3)에 의해, 완비 뇌터 국소환이면서 평탄한 $A$-가군인 위상
$A$-대수 $B$와 $k$-동형사상
\[
  u_{00}:B\otimes_A k\xrightarrow{\sim}B_{00}
\]
가 존재한다. (19.7.1)에 의해 $B$는 형식적으로 매끄러운
$A$-대수이다. 따라서
$B\otimes_A A_0=B/\mathfrak JB$는 형식적으로 매끄러운
$A_0$-대수이자 (19.3.5, (iii)) 완비 뇌터 국소환이다. 한편 가정과
(19.7.1)에 의해 $B_0$는 평탄한 $A_0$-가군이다. 이제 $k$-동형사상
\[
  u_{00}:B\otimes_A k
  =(B\otimes_A A_0)\otimes_{A_0}k
  \xrightarrow{\sim}B_0\otimes_{A_0}k=B_{00}
\]
가 있으므로, (19.7.1.5)에서 $A$를 $A_0$로, $\mathfrak J$를
$\mathfrak m_0$로 바꾸어 적용하면
$u_{00}=u\otimes1$을 만족하는 $A_0$-동형사상
\[
  u:B\otimes_A A_0\xrightarrow{\sim}B_0
\]
를 얻는다.

유일성에 관한 주장을 보자. 아이디얼 $\mathfrak J^hB$와
$\mathfrak J^hB'$은 각각 $B$와 $B'$에서 닫혀 있다
($\mathbf 0_{\mathrm I}$, 7.3.5). 따라서 $B$와 $B'$은
$\mathfrak J$-준아딕 위상에 대하여 분리이고 완비이다 (Bourbaki,
\emph{Top. gén.}, chap.~III, 3\textsuperscript{e} éd., \textsection3,
n\textsuperscript{o}~5, cor.~2 de la prop.~9). 가정에 의해
\[
  v_0:B\otimes_A A_0\xrightarrow{\sim}B'\otimes_A A_0,
  \qquad u'\circ v_0=u
\]
인 $A_0$-동형사상이 있다. $B$는 형식적으로 매끄러운 $A$-대수이고
$B'$은 평탄한 $A$-가군이므로 (19.7.1.5)를 적용하여 요구한
$A$-동형사상 $v$를 얻는다.
\end{proof}

\begin{remarks}[19.7.3]
\label{0.19.7.3-ko}
\begin{enumerate}
  \item[(i)] (19.7.2)의 유일성 주장은 $B$와 $B'$이
  \emph{$\mathfrak J$-준아딕 위상에 대하여 완비}라고만 가정해도
  여전히 성립한다. 존재성 주장도 같은 방식으로 강화할 수 있는지는
  알려져 있지 않다. 다시 말해 국소환 $B_0$가 그 극대 아이디얼을
  $\mathfrak n_0$라 할 때 $\mathfrak n_0$-준아딕 위상에 대하여
  \emph{완비}라는 가정을 버리고, $B$가
  \emph{$\mathfrak J$-준아딕 위상에 대하여 완비}라고만 요구할 수
  있는지는 알려져 있지 않다. $A_0$가 $\mathfrak m$-준아딕 위상에
  대하여 완비인 경우, 이 문제는 다음 문제와 동치임을 보일 수 있다:
  소체 $\mathbf F_p=\mathbf Z/p\mathbf Z$를 포함하는 정칙 뇌터
  국소환 $B_0$가 주어졌을 때(완비일 필요는 없다), 모든 $n>1$에
  대하여 $B/pB$가 $B_0$와 동형인 평탄한
  $(\mathbf Z/p^n\mathbf Z)$-대수 $B$가 존재하는가?
  \item[(ii)] 일반적으로 (19.7.2)에서 존재가 주장된 동형사상 $v$는
  \emph{유일하지 않다} (19.8.7 참조).
\end{enumerate}
\end{remarks}

\subsection{Cohen 대수와 Cohen의 $p$-환; 완비 국소환의 구조에 대한 응용}
\label{subsection:0.19.8-ko}

이 절의 결과들은 (19.7)의 정리들을 곧바로 적용한 것이지만, 실제로
중요하므로 명시해 둘 가치가 있다.

\begin{definition}[19.8.1]
\label{0.19.8.1-ko}
$A$, $B$를 뇌터 국소환, $\mathfrak m$을 $A$의 극대 아이디얼,
$k=A/\mathfrak m$을 그 잉여체라 하고, $\varphi:A\to B$를 $B$에
$A$-대수 구조를 주는 국소 준동형이라 하자. 다음 조건들이 성립하면
$B$를 \emph{Cohen $A$-대수}라 한다.
\begin{enumerate}
  \item[(i)] $B$는 \emph{완비환}이다.
  \item[(ii)] $B$는 평탄한 $A$-가군이다.
  \item[(iii)] $B\otimes_A k$는 $k$의 \emph{분리 확대체}인
  \emph{체}이다. 다시 말해 $\mathfrak mB$는 $B$의 극대
  아이디얼이다.
\end{enumerate}
\end{definition}

\oldpage[0\textsubscript{IV-1}]{110}

\begin{theorem}[19.8.2]
\label{0.19.8.2-ko}
$A$를 뇌터 국소환, $k$를 그 잉여체라 하자.
\begin{enumerate}
  \item[(i)] $B$가 Cohen $A$-대수이면, $B$는 형식적으로 매끄러운
  $A$-대수이다. 따라서 임의의 완비 뇌터 국소환 $C$, 임의의 국소
  준동형 $A\to C$, 그리고 $C$의 임의의 아이디얼
  $\mathfrak J\neq C$에 대하여, 모든 $A$-준동형
  $B\to C/\mathfrak J$는
  \[
    B\xrightarrow{v}C\longrightarrow C/\mathfrak J
  \]
  로 인수분해된다. 여기서 $v$는 $A$-준동형이며, 반드시 국소
  준동형이다.
  \item[(ii)] $k$의 임의의 분리 확대체 $K$에 대하여,
  $B\otimes_A k$가 $K$와 $k$-동형인 Cohen $A$-대수 $B$가
  존재하고, 그러한 $A$-대수는 동형을 제외하면 유일하다.
\end{enumerate}
\end{theorem}

\begin{proof}
$B\otimes_A k$는 $k$의 분리 확대체이므로 형식적으로 매끄러운
$k$-대수이다 (19.6.1). 따라서 (19.7.1)에 의해 (i)가 성립한다.
(ii)를 증명할 때에는 $A$가 \emph{완비}인 경우만 생각해도 된다.
실제로 $B$가 평탄한 $A$-가군인 것과 평탄한 $\widehat A$-가군인
것은 동치이고 ($\mathbf 0_{\mathrm{III}}$, 10.2.3),
$\mathfrak mB=\mathfrak m\widehat AB$이며 $k$는 $\widehat A$의
잉여체이다. 이제 (19.7.2)에서 $\mathfrak J=\mathfrak m$,
$B_0=K$로 놓고 (19.6.1)을 사용하면 된다.
\end{proof}

\begin{definition}[19.8.3]
\label{0.19.8.3-ko}
$p\mathbf Z$가 $\mathbf Z$의 소 아이디얼일 때
$\mathbf Z_{p\mathbf Z}$ 꼴의 국소환을 \emph{소 국소환}이라 한다.
소 국소환의 완비화를 \emph{완비 소 국소환}이라 한다.
\end{definition}

소 국소환에는 다음 두 종류가 있다.
\begin{enumerate}
  \item[$1^{\mathrm o}$] $0$이 아닌 소수 $p$에 대응하는 극대
  아이디얼 $p\mathbf Z$의 경우, $\mathbf Z_{p\mathbf Z}$는 이산 값매김환이다.
  그 완비화는 보통 $\mathbf Z_p$로 나타내는
  \emph{$p$-아딕 정수환}이다\footnote{현재 보편적으로 쓰이는 이
  표기는 이 경우 ($\mathbf 0_{\mathrm I}$, 1.2.3)에서 채택한
  $A_f$ 표기와 충돌한다. 실제로 $A=\mathbf Z$, $f=p$이면
  $A_f$는 $k\in\mathbf Z$, 정수 $n\geq0$에 대하여 $k/p^n$ 꼴인
  유리수들의 환을 뜻한다. 우리는 언제나 이 후자의 환을 나타내는
  데 $\mathbf Z_p$ 표기를 쓰지 않을 것이다.}.
  \item[$2^{\mathrm o}$] 소 아이디얼 $p\mathbf Z=(0)$의 경우,
  $\mathbf Z_{p\mathbf Z}$는 유리수체 $\mathbf Q$이고 자기 완비화와
  같다. 여기서 위상은 자연스럽게 뇌터 국소환의 위상이므로 이
  경우에는 이산 위상이다.
\end{enumerate}

(19.8.3)의 용어는 ‘소체’라는 용어와 마찬가지 방식으로
정당화된다. 임의의 국소환 $A$에 대하여 표준 준동형
$\varphi:\mathbf Z\to A$를 생각하고, $A$의 극대 아이디얼
$\mathfrak m$의 역상을 $p\mathbf Z$라 하자. 이는 $\mathbf Z$의
소 아이디얼이므로 위 준동형은
\[
  \mathbf Z\longrightarrow\mathbf Z_{p\mathbf Z}
  \xrightarrow{\psi}A
\]
로 인수분해된다. 더욱이 $\varphi$는 $\mathbf Z$에서 $A$로 가는
\emph{유일한} 준동형이므로 $p$와 $\psi$는 \emph{유일하게}
정해진다. 다시 말해 임의의 국소환 $A$에 대하여, 어떤 소 국소환
$P$에서 $A$로 가는 \emph{유일한 준동형} $\psi:P\to A$가 있다.
더 나아가 $A$가 분리이고 완비이면 이 준동형을 완비화에 의해
연장할 수 있으므로, 완비 소 국소환 $\widehat P$에서 $A$로 가는
\emph{유일한 준동형}
$\widehat\psi:\widehat P\to A$가 있다. 또한 몫을 취하면 $\psi$는
$p>0$일 때 잉여체 $\mathbf F_p$에서, $p=0$일 때 $\mathbf Q$에서
$A$의 잉여체 $k$로 가는 준동형을 주므로, $p$는 $k$의
\emph{표수}이다.

특히 $A$를 소 국소환(각각 완비 소 국소환)으로 택하면, 그러한 환의
자기 준동형은 항등 준동형 하나뿐임을 알 수 있다.

\oldpage[0\textsubscript{IV-1}]{111}

\begin{definition}[19.8.4]
\label{0.19.8.4-ko}
$A$를 국소환, $P\to A$를 소 국소환 $P$에서 $A$로 가는 유일한
준동형, $p$를 $P$와 $A$의 잉여체들의 표수라 하자. $A$가 Cohen
$P$-대수일 때, 즉 (19.8.1)에 따라 다음 조건들을 만족할 때 $A$를
\emph{Cohen 환}이라 한다.
\begin{enumerate}
  \item[$1^{\mathrm o}$] $A$는 뇌터이고 완비이다.
  \item[$2^{\mathrm o}$] $A$는 평탄한 $P$-가군이다. 이는 $A$가
  평탄한 $\widehat P$-가군이라는 것과도 동치이다 (Bourbaki,
  \emph{Alg. comm.}, chap.~III, \textsection5,
  n\textsuperscript{o}~4, prop.~4).
  \item[$3^{\mathrm o}$] $A/pA$는 체이다. 이 체는 $P$의 잉여체
  위에서 반드시 분리인데, 그 잉여체가 소체이기 때문이다.
\end{enumerate}
\end{definition}

$p=0$이면 이 조건들은 $A$가 \emph{표수 $0$인 체}라는 것과
동치이다. $p>0$이면 반드시 $pA\neq0$이다. 조건
$3^{\mathrm o}$는 $pA$가 $A$의 극대 아이디얼 $\mathfrak m$이라는
뜻이고, 조건 $2^{\mathrm o}$는 $P$가 이산 값매김환이므로
$p$가 $A$-정칙이라는 뜻이다 ($\mathbf 0_{\mathrm I}$, 6.3.4).
따라서 $A$는 차원 $1$인 정칙환이고 (17.1.1, $d)$), 그러므로
이산 값매김환이며 $1^{\mathrm o}$에 의해 완비이다. 요약하면 다음과
같다.

\begin{proposition}[19.8.5]
\label{0.19.8.5-ko}
Cohen 환은 표수 $0$인 체이거나, 잉여체의 표수가 $p>0$이고
극대 아이디얼이 $p\cdot1$로 생성되는 완비 이산 값매김환이다.
여기서 $1$은 환의 단위원이다.
\end{proposition}

두 번째 경우에는 $p$가 $A$-정칙이므로 $p\cdot1\neq0$이다. 따라서
$p\cdot1$을 정수 $p$와 동일시할 수 있다. 표준 준동형
$\mathbf Z_p\to A$는 단사이고, $p\cdot1$을 $\mathbf Z_p$의 원소
$p$와 동일시한다. 이 경우 $A$를 \emph{Cohen의 $p$-환}이라 한다.

\begin{theorem}[19.8.6]
\label{0.19.8.6-ko}
\textup{(Cohen).} ---
\begin{enumerate}
  \item[(i)] $W$를 Cohen 환, $C$를 완비 뇌터 국소환,
  $\mathfrak J$를 $C$와 다른 $C$의 아이디얼이라 하자. 그러면 모든
  국소 준동형 $u:W\to C/\mathfrak J$는
  \[
    W\xrightarrow{v}C\longrightarrow C/\mathfrak J
  \]
  로 인수분해된다. 여기서 $v$는 국소 준동형이다.
  \item[(ii)] 임의의 체 $K$에 대하여 잉여체가 $K$와 동형인 Cohen 환
  $W$가 존재한다. $W'$이 또 다른 Cohen 환이고 $K'$이 그 잉여체이면,
  모든 동형사상 $u:K\xrightarrow{\sim}K'$은 몫을 취하여 얻어지는
  어떤 동형사상 $v:W\xrightarrow{\sim}W'$에서 유도된다.
\end{enumerate}
\end{theorem}

\begin{proof}
이는 (19.8.2)에서 $A$가 소 국소환인 경우를 적용한 것에 지나지 않는다.
\end{proof}

\begin{remarks}[19.8.7]
\label{0.19.8.7-ko}
\begin{enumerate}
  \item[(i)] $K$의 표수가 $0$이면 (19.8.6)의 (ii)는 자명해진다.
  \item[(ii)] (19.8.6, (i))의 준동형 $v$는 $u$에 의해
  \emph{반드시 유일하게 정해지는 것은 아니다}. 이미 $W$가 표수
  $0$인 체이고 $\mathfrak J^2=0$이며 $u$가 동형사상인 경우가 이를
  보여 준다 (21.5.5 참조). 마찬가지로 (19.8.6, (ii))의 동형사상
  $v$도 $u$에 의해 \emph{반드시 유일하게 정해지는 것은 아니다}
  (21.5.5 참조).

  그러나 $K$가 표수 $p>0$인 \emph{완전체}이면, (19.8.6, (ii))의
  동형사상 $v$가 유일함을 (21.5.5)에서 볼 것이다. 또한 이 경우
  $W$가 $K$ 위의 \emph{무한 길이 비트 벡터}의 환
  $W_\infty(K)$와 동일시됨을 뒤에서 볼 것이다.
  \item[(iii)] (19.8.6, (i))에서는 (19.3.10)과 (19.3.12)를
  이용하여 $C$에 대한 가정을 약화할 수 있다.
\end{enumerate}
\end{remarks}

\begin{theorem}[19.8.8]
\label{0.19.8.8-ko}
\textup{(Cohen).} --- $A$를 완비 뇌터 국소환, $k$를 그 잉여체라 하자.

\oldpage[0\textsubscript{IV-1}]{112}

\begin{enumerate}
  \item[(i)] 어떤 Cohen 환 $W$가 존재하여 $A$는 형식적 멱급수환
  $W[[T_1,\ldots,T_n]]$의 한 몫환과 동형이다. 특히 $A$는 한 완비
  정칙 국소환의 몫과 동형이다 (17.3.8). $A$가 체를 포함하면
  $A$는 $k[[T_1,\ldots,T_n]]$의 한 몫환과 동형이다.
  \item[(ii)] 더 나아가 $A$가 정역이라고 하자. 그러면 다음 조건을
  만족하는 $A$의 부분환 $B$가 존재한다. $1^{\mathrm o}$ $B$는 어떤
  체 또는 Cohen 환인 환 $C$ 위의 형식적 멱급수환과 동형이다. 따라서
  $B$는 완비 정칙 국소환이다 (17.3.8). $2^{\mathrm o}$ $B$와 $A$의
  잉여체는 같고, 포함사상 $B\to A$는 국소 준동형이다.
  $3^{\mathrm o}$ $A$는 유한 $B$-대수이다.
\end{enumerate}
\end{theorem}

\begin{proof}
$\mathfrak m$을 $A$의 극대 아이디얼이라 하자. 잉여체가 $k$와 동형인
Cohen 환 $W$가 존재한다 (19.8.6, (ii)). 따라서 국소 준동형
$W\to A/\mathfrak m$이 있고, 이는
$W\xrightarrow{u}A\to A/\mathfrak m$으로 인수분해된다. 여기서 $u$는
국소 준동형이다 (19.8.6, (i)). $\mathfrak m$의 원소들의 임의의 유한족
$(x_i)_{1\leq i\leq n}$에 대하여, $u$를 연장하고 모든 $i$에 대해
$v(T_i)=x_i$를 만족하는 국소 준동형
\[
  v:W[[T_1,\ldots,T_n]]\longrightarrow A
\]
가 존재한다 (Bourbaki, \emph{Alg. comm.}, chap.~III, \textsection4,
n\textsuperscript{o}~5, prop.~6). $A$가 체를 포함하면, $A$는 소체
$P$를 포함하고 $k$는 $P$의 확대체이다. 이 확대체는 반드시
분리가능하므로 $A$는 $k$와 동형인 체를 포함한다 (19.6.2). 따라서
위의 $v$의 정의에서 $W$를 $k$로 바꿀 수 있다.

\medskip
\noindent(i)\enspace
먼저 $x_i$들을 $\mathfrak m$의 한 생성계로 택하자. $W$와 $A$의
잉여체가 같고, $x_i$들의 $\gr_\bullet(A)$에서의 류들이
$\gr_\bullet(A)$를 $k$-대수로서 생성하므로
\[
  \gr(v):\gr_\bullet(W[[T_1,\ldots,T_n]])
  \longrightarrow\gr_\bullet(A)
\]
는 전사이다. 따라서 $v$ 자체도 전사이다 (Bourbaki,
\emph{Alg. comm.}, chap.~III, \textsection2,
n\textsuperscript{o}~8, th.~1의 cor.~2). $A$가 체를 포함하는 경우는
이미 증명했으며 여기서는 상기하기 위해서만 적었음을 덧붙인다
(19.6.3).

\medskip
\noindent(ii)\enspace
$A$가 체를 포함하면, 앞에서 본 바와 같이 $A$는 $k$와 동형인 체
$k'$을 포함한다. $A$의 한 매개변수계
$(y_j)_{1\leq j\leq m}$을 택하고 (16.3.6),
$B=k[[T_1,\ldots,T_m]]$로 놓자. 그리고 $k$ 위에서는 동형사상
$k\to k'$과 일치하고 $1\leq j\leq m$에 대하여 $w(T_j)=y_j$를
만족하는 국소 준동형 $w:B\to A$를 생각하자.

$A$가 체를 포함하지 않으면, 유일한 준동형
$\mathbf Z_{p\mathbf Z}\to A$는 반드시 단사이다 (19.8.3). 실제로
그렇지 않으면 $A$가 정역이므로 그 핵은 $\mathbf Z_{p\mathbf Z}$의
극대 아이디얼이고 그 상은 체와 동형이 된다. 또한 가정에 의해
$p>0$이다. $p=0$이면 $\mathbf Z_{p\mathbf Z}$가 체이기 때문이다.
$A$의 원소 $p\cdot1$을 $p$와 동일시하자. 이 원소는 $A$에서 영인자가
아니고 $\mathfrak m$에 속한다. 따라서 (16.3.4)와 (16.3.7)에 의해,
$p$와 함께 $A$의 한 매개변수계를 이루는 원소족
$(z_i)_{1\leq i\leq m-1}$이 존재한다. 증명 첫머리에서 택한 Cohen 환
$W$는 잉여체가 $k$인 이산 값매김환이고, 그 극대 아이디얼은 $p$로
생성된다 (19.8.5). 또한 앞서 정의한 유일한 준동형 $u:W\to A$는
$p$를 자기 자신으로 보낸다. 이제 $B=W[[T_1,\ldots,T_{m-1}]]$로 놓고,
$W$ 위에서는 $u$와 일치하며 $1\leq i\leq m-1$에 대하여
$w(T_i)=z_i$를 만족하는 국소 준동형 $w:B\to A$를 생각하자.

두 경우 모두 $\mathfrak n$을 $B$의 극대 아이디얼이라 하면
$\mathfrak nA$는 분명 $A$의 정의 아이디얼이다. 더욱이
$B/\mathfrak n$과 $A/\mathfrak m$이 동형이므로 $A$는 준유한
$B$-가군이다 ($\mathbf 0_{\mathrm I}$, 7.4.4). 따라서 $B$가 완비이고
$A$가 $\mathfrak n$-준아딕 위상에 대하여 분리이므로 $A$는 유한 생성
$B$-가군이다 ($\mathbf 0_{\mathrm I}$, 7.4.1).

\oldpage[0\textsubscript{IV-1}]{113}

한편 두 경우 모두 $\dim(B)=\dim(A)=m$이다. 첫 번째 경우에는
(17.1.4, (iii))에서 따르고, 두 번째 경우에는 $p$와 $T_i$
($1\leq i\leq m-1$)가 $\mathfrak n$을 생성하는 $B$-정칙열을 이룸을
직접 보면 된다. 또는 이 원소들이 $\mathfrak n$을 생성한다는 것과
(16.3.10)에 따른 부등식 $\dim(B)\geq\dim(A)$를 써도 된다. $A$와
$B$가 정역이므로 (16.3.10)에서 마침내 $w$가 단사임을 얻는다. 이로써
증명이 끝난다.
\end{proof}

\begin{corollary}[19.8.9]
\label{0.19.8.9-ko}
$A$를 체 $k_0$를 포함하는 완비 뇌터 국소 정역이라 하고, $k$를 $A$의
잉여체라 하자. $k$가 $k_0$ 위에서 유한이라고 가정하자. 그러면
(19.8.8, (ii))의 결론에서 조건 $1^{\mathrm o}$와 $2^{\mathrm o}$를
다음 조건으로 바꿀 수 있다. 즉 $B$는 $k_0[[T_1,\ldots,T_n]]$ 꼴이고,
표준 포함사상 $B\to A$는 $B$의 통상적인 $k_0$-대수 구조에 대한 국소
$k_0$-준동형이다.
\end{corollary}

\begin{proof}
실제로 (19.8.8, (ii))의 증명을 되풀이하되, 이번에는
\[
  w:k_0[[T_1,\ldots,T_n]]\longrightarrow A
\]
가 $k_0$ 위에서는 항등사상과 일치하고 $1\leq j\leq n$에 대하여
$T_j$를 $y_j$로 보내도록 정의한다. $k$가 $k_0$ 위에서 유한 차수라는
가정에 의해 $A$는 여전히 준유한 $B$-가군이고, 따라서
($\mathbf 0_{\mathrm I}$, 7.4.1)에 의해 유한 생성 $B$-가군이다.
이제 (19.8.8)에서와 같이 결론을 얻는다.
\end{proof}

\begin{corollary}[19.8.10]
\label{0.19.8.10-ko}
$A$를 극대 아이디얼 $\mathfrak m$의 제곱이 $0$인 아르틴 국소환이라
하자. 그러면 극대 아이디얼이 $\mathfrak n$인 뇌터 정칙 국소환 $B$가
존재하여 $A$는 $B/\mathfrak n^2$과 동형이다.
\end{corollary}

\begin{proof}
$K=A/\mathfrak m$을 $A$의 잉여체라 하고,
$n=\dim_K(\mathfrak m/\mathfrak m^2)=\dim_K(\mathfrak m)$이라 하자.
$A$가 체를 포함하면 (19.6.3)에 의해 $A$는 $B/\mathfrak b$와 동형이다.
여기서 $B=K[[T_1,\ldots,T_n]]$이고 $\mathfrak b$는 $B$의 극대
아이디얼 $\mathfrak n$의 제곱에 포함된다. 그런데
\[
  \operatorname{long}(B/\mathfrak n^2)=n+1
  =\operatorname{long}(A)
\]
이므로 반드시 $\mathfrak b=\mathfrak n^2$이다.

이제 $A$가 체를 포함하지 않는다고 하자. 그러면 $K$의 표수는
$p>0$이고 $A$에서 $p\cdot1\neq0$이다 (19.6.3). 따라서
$p\cdot1\in\mathfrak m$이고, $p\cdot1$과 함께 $\mathfrak m$의
$K$-기저를 이루는 $n-1$개의 다른 원소
$x_i$ ($2\leq i\leq n$)가 존재한다. 잉여체가 $K$와 동형인 Cohen 환
$W$를 택하자. $W$는 $p$가 극대 아이디얼을 생성하는 이산 값매김환이다.
(19.8.8)의 증명에서, $p$를 자기 자신으로 보내고 몫을 취하면 $K$ 위의
항등사상을 주는 준동형 $u:W\to A$가 있음을 보았다. 이제
$B=W[[T_2,\ldots,T_n]]$로 놓고, $W$ 위에서는 $u$와 일치하며
$2\leq i\leq n$에 대하여 $w(T_i)=x_i$를 만족하는 국소 준동형
$w:B\to A$를 생각하자. $w$는 분명 전사이고, 그 핵 $\mathfrak b$는
$B$의 극대 아이디얼
\[
  \mathfrak n=pB+BT_2+\cdots+BT_n
\]
의 제곱에 포함된다. 다시
$\operatorname{long}(B/\mathfrak n^2)=n+1=\operatorname{long}(A)$이므로
$\mathfrak b=\mathfrak n^2$이다.
\end{proof}

\begin{proposition}[19.8.11]
\label{0.19.8.11-ko}
$A$를 아르틴 국소환, $\mathfrak m$을 그 극대 아이디얼, $k$를 그
잉여체라 하자. $A$가 어떤 Cohen 환의 몫환과 동형이기 위한 필요충분조건은
$\mathfrak m$이 $p\cdot1$로 생성되는 것이다. 여기서 $p$는 $k$의
표수이다.
\end{proposition}

\begin{proof}
이 조건이 필요함은 분명하다 (19.8.5). 충분함을 보이기 위하여
(19.8.8)의 증명 첫머리에서와 같이, 잉여체가 $k$와 동형인 Cohen 환
$W$와 국소 준동형 $u:W\to A$가 존재함을 관찰하자. 합성 준동형
\[
  \mathbf Z_{p\mathbf Z}\longrightarrow W\xrightarrow{u}A
\]
을 생각하면, 이는 반드시 $\mathbf Z_{p\mathbf Z}$에서 $A$로 가는
유일한 준동형이다. 따라서 $W$의 원소 $p\cdot1$의 $u$에 의한 상은

\oldpage[0\textsubscript{IV-1}]{114}

$A$의 원소 $p\cdot1$이다. $W$의 원소 $p\cdot1$이 $W$의 극대
아이디얼을 생성하므로, 가정에서 곧바로
\[
  \gr(u):\gr_\bullet(W)\longrightarrow\gr_\bullet(A)
\]
가 전사임을 얻는다. 따라서 $u$도 전사이다 (Bourbaki,
\emph{Alg. comm.}, chap.~III, \textsection2,
n\textsuperscript{o}~8, th.~1의 cor.~2).
\end{proof}

\subsection{상대적으로 형식적으로 매끄러운 대수}
\label{subsection:0.19.9-ko}

\begin{definition}[19.9.1]
\label{0.19.9.1-ko}
$\Lambda$를 위상환, $A$를 위상 $\Lambda$-대수, $B$를 위상
$A$-대수라 하자. 임의의 이산 위상 $A$-대수 $C$와 $C$의 임의의
멱영 아이디얼 $\mathfrak J$에 대하여 다음 조건이 성립할 때, $B$를
\emph{$\Lambda$에 대하여 상대적으로 형식적으로 매끄러운 $A$-대수}라
한다. 즉 연속 $A$-준동형 $u_0:B\to C/\mathfrak J$가
\[
  B\xrightarrow{u}C\longrightarrow C/\mathfrak J
\]
로 인수분해되고 $u$가 연속 $\Lambda$-준동형이면, $u_0$는 또한
\[
  B\xrightarrow{v}C\longrightarrow C/\mathfrak J
\]
로 인수분해된다. 여기서 $v$는 연속 $A$-준동형이다.
\end{definition}

이 정의에서 곧 다음이 따른다. $B$가 \emph{형식적으로 매끄러운}
$A$-대수이면, $A$ 위에 정의된 임의의 위상 $\Lambda$-대수 구조,
다시 말해 임의의 연속 환 준동형 $\Lambda\to A$에 대하여 $B$는
$\Lambda$에 대하여 상대적으로 형식적으로 매끄럽다.

\begin{proposition}[19.9.2]
\label{0.19.9.2-ko}
$\Lambda$를 위상환, $A$를 위상 $\Lambda$-대수라 하자.
\begin{enumerate}
  \item[(i)] $A$는 $\Lambda$에 대하여 상대적으로 형식적으로
  매끄러운 $\Lambda$-대수이다.
  \item[(ii)] $B$가 $\Lambda$에 대하여 상대적으로 형식적으로
  매끄러운 $A$-대수이고 $C$가 $\Lambda$에 대하여 상대적으로
  형식적으로 매끄러운 $B$-대수이면, $C$는 $\Lambda$에 대하여
  상대적으로 형식적으로 매끄러운 $A$-대수이다.
  \item[(iii)] $B$를 $\Lambda$에 대하여 상대적으로 형식적으로
  매끄러운 $A$-대수, $A'$을 위상 $A$-대수라 하자. 그러면 위상
  $A'$-대수 $B\otimes_A A'$은 $\Lambda$에 대하여 상대적으로
  형식적으로 매끄럽다.
  \item[(iv)] $B$를 위상 $A$-대수, $S$와 $T$를 각각 $A$와 $B$의
  곱셈적 부분집합이라 하고, $S$의 $B$에서의 표준 상이 $T$에
  포함된다고 하자. $B$가 $\Lambda$에 대하여 상대적으로 형식적으로
  매끄러운 $A$-대수이면, $T^{-1}B$는 $\Lambda$에 대하여 상대적으로
  형식적으로 매끄러운 $S^{-1}A$-대수이다.
  \item[(v)] $B_i$ ($1\leq i\leq n$)를 위상 $A$-대수들이라 하자.
  $\prod_{i=1}^n B_i$가 $\Lambda$에 대하여 상대적으로 형식적으로
  매끄러운 $A$-대수이기 위한 필요충분조건은 각 $B_i$가 그러한
  $A$-대수인 것이다.
\end{enumerate}
\end{proposition}

\begin{proof}
(i)는 자명하고, 나머지 명제들의 증명은 (19.3.5)의 증명들을 그대로
따라가므로 독자에게 맡긴다.
\end{proof}

\begin{corollary}[19.9.3]
\label{0.19.9.3-ko}
$A$를 위상환, $A'$과 $B$를 두 위상 $A$-대수라 하자. 그러면 위상
$A$-대수 $B'=B\otimes_A A'$은 $A$에 대하여 상대적으로 형식적으로
매끄러운 $B$-대수이다.
\end{corollary}

\begin{proof}
이는 (19.9.2, (i), (iii))에서 따른다.
\end{proof}

\begin{proposition}[19.9.4]
\label{0.19.9.4-ko}
$\Lambda$를 위상환, $A$를 위상 $\Lambda$-대수, $B$를 위상
$A$-대수라 하자. 다음 조건들은 서로 동치이다.
\begin{enumerate}
  \item[a)] $B$는 $\Lambda$에 대하여 상대적으로 형식적으로
  매끄러운 $A$-대수이다.
  \item[b)] $\widehat B$는 $\Lambda$에 대하여 상대적으로 형식적으로
  매끄러운 $A$-대수이다.
  \item[c)] $\widehat B$는 $\Lambda$에 대하여 상대적으로 형식적으로
  매끄러운 $\widehat A$-대수이다.
  \item[d)] $\widehat B$는 $\widehat\Lambda$에 대하여 상대적으로
  형식적으로 매끄러운 $\widehat A$-대수이다.
\end{enumerate}
\end{proposition}

\begin{proof}
이번에도 (19.3.6)의 증명을 그대로 따른다. 증명은 독자에게 맡긴다.
\end{proof}

\oldpage[0\textsubscript{IV-1}]{115}

\begin{env}[19.9.5]
\label{0.19.9.5-ko}
마찬가지로 (19.3.8)의 명제에서 ‘형식적으로 매끄러운’을
‘$\Lambda$에 대하여 상대적으로 형식적으로 매끄러운’으로 바꾸어도
같은 명제와 같은 증명이 성립한다. (19.3.10)의 명제에서 같은 방식으로
바꾸면 결론은 다음으로 대체되고, 증명은 본질적으로 변하지 않는다.
\emph{연속 $A$-준동형 $u_0:B\to C/\mathfrak J$가
$B\xrightarrow{u}C\to C/\mathfrak J$로 인수분해되고 $u$가 연속
$\Lambda$-준동형이면, $u_0$는 또한
$B\xrightarrow{v}C\to C/\mathfrak J$로 인수분해된다. 여기서 $v$는
연속 $A$-준동형이다.}
\end{env}

\begin{env}[19.9.6]
\label{0.19.9.6-ko}
형식적 매끄러움의 판정법 (19.4.1)과 (19.4.2)에서 ‘형식적으로
매끄러운’을 ‘$\Lambda$에 대하여 상대적으로 형식적으로 매끄러운’으로
바꾸어도 이 판정법들은 성립하며, 증명들도 거의 변하지 않는다.
\end{env}

\begin{proposition}[19.9.7]
\label{0.19.9.7-ko}
$\Lambda$를 위상환, $A$를 위상 $\Lambda$-대수, $B$를 위상
$A$-대수라 하자. 임의의 이산 $A$-대수 $C$와 $\mathfrak J^2=0$인
$C$의 임의의 아이디얼 $\mathfrak J$에 대하여 다음이 성립한다고 하자.
즉 연속 $A$-준동형 $u_0:B\to C/\mathfrak J$가
\[
  B\xrightarrow{u}C\longrightarrow C/\mathfrak J
\]
로 인수분해되고 $u$가 연속 $\Lambda$-준동형이면, $u_0$는 또한
\[
  B\xrightarrow{v}C\longrightarrow C/\mathfrak J
\]
로 인수분해되며 $v$는 연속 $A$-준동형이다. 그러면 $B$는
$\Lambda$에 대하여 상대적으로 형식적으로 매끄러운 $A$-대수이다.
\end{proposition}

\begin{proof}
(19.4.3)의 증명을 그대로 옮기면 된다.
\end{proof}

\begin{proposition}[19.9.8]
\label{0.19.9.8-ko}
$\Lambda$를 위상환, $A$를 위상 $\Lambda$-대수, $B$를 위상
$A$-대수라 하자. $B$가 $\Lambda$에 대하여 상대적으로 형식적으로
매끄러운 $A$-대수이기 위한 필요충분조건은, $B$의 열린 아이디얼에
의하여 소멸되는 모든 이산 위상 $B$-가군 $L$에 대하여
(18.4.2 참조)
\[
  \operatorname{Exalcotop}_{A/\Lambda}(B,L)=0
\]
인 것이다.
\end{proposition}

\begin{proof}
(19.4.4)의 증명에서와 같은 표기를 쓰자. 여기서는 $B_\Lambda$의 확대
$E_\Lambda$가 $\Lambda$-자명이라고 가정할 수 있다는 점만 관찰하면
된다. 나머지 증명은 변하지 않는다.
\end{proof}

$\Lambda$, $A$, $B$가 이산환이면 판정법 (19.9.8)은 다음으로
환원된다.
\begin{equation}
\tag{19.9.8.1}
\label{0.19.9.8.1-ko}
  \operatorname{Exalcom}_{A/\Lambda}(B,L)=0
  \quad\text{모든 $B$-가군 $L$에 대하여}.
\end{equation}
달리 말해, \emph{$B$-가군에 의한 $B$의 모든 가환 $A$-확대가
$\Lambda$-자명이면 그 확대는 $A$-자명하다.}

\subsection{형식적으로 비분기인 대수와 형식적으로 에탈인 대수}
\label{subsection:0.19.10-ko}
\label{0.19.10-ko}

\begin{env}[19.10.1]
\label{0.19.10.1-ko}
형식적으로 매끄러운 대수라는 개념과 밀접하게 연관된 두 개념도
사용하게 될 것이다.
\end{env}

\begin{definition}[19.10.2]
\label{0.19.10.2-ko}
$A$를 위상환, $B$를 위상 $A$-대수라 하자. 핵이 멱영인 이산
$A$-대수들의 임의의 $A$-준동형 $p:E\to C$와 임의의 연속
$A$-준동형 $u:B\to C$에 대하여, $u$가
\[
  B\xrightarrow{v}E\xrightarrow{p}C
\]
로 인수분해되게 하는 연속 $A$-준동형 $v:B\to E$가
\emph{많아야 하나}일 때, $B$를 \emph{형식적으로 비분기인
$A$-대수}라 한다.

$B$가 형식적으로 매끄럽고 형식적으로 비분기일 때, $B$를
\emph{형식적으로 에탈인 $A$-대수}라 한다. 달리 말해, 핵이 멱영인
이산 $A$-대수들의 임의의 전사 $A$-준동형 $p:E\to C$와 임의의 연속
$A$-준동형 $u:B\to C$에 대하여, $u$가
\[
  B\xrightarrow{v}E\xrightarrow{p}C
\]
로 인수분해되게 하는 연속 $A$-준동형 $v:B\to E$가
\emph{정확히 하나}일 때와 같다.
\end{definition}

\oldpage[0\textsubscript{IV-1}]{116}

\begin{examples}[19.10.3]
\label{0.19.10.3-ko}
여기서는 \emph{이산환}에 한정한다.
\begin{enumerate}
  \item[(i)] 구조 준동형 $\rho:A\to B$가 \emph{전사}이면 $B$는
  \emph{형식적으로 비분기인} $A$-대수이다. 실제로 (19.10.2)의
  표기 아래 합성 준동형
  \[
    A\xrightarrow{\rho}B\xrightarrow{v}E
  \]
  는 구조 준동형이어야 하므로, $v$가 존재한다면 유일하다. 한편
  $B$가 형식적으로 매끄럽기 위한 필요충분조건은
  $\ker(\rho)=\ker(\rho)^2$인 것이다.
  \item[(ii)] $A$를 환, $S$를 $A$의 곱셈적 부분집합이라 하자. 그러면
  $B=S^{-1}A$는 \emph{형식적으로 에탈인} $A$-대수이다. 실제로
  $j:A\to B$를 표준 준동형, $\rho:A\to E$를 환 준동형,
  $\mathfrak J$를 $E$의 멱영 아이디얼,
  $p:E\to C=E/\mathfrak J$를 표준 준동형이라 하자. $u:B\to C$가
  $A$-준동형이면, $s\in S$에 대하여 $s/1$이 $B$에서 가역이므로
  $u(s/1)$은 $C$에서 가역이다. 그런데 $u(s/1)=p(\rho(s))$이고
  $\mathfrak J$가 멱영이므로, $\rho(s)$의 $E/\mathfrak J$에서의 상이
  가역이면 $\rho(s)$ 자체가 $E$에서 가역이다
  ($\mathbf 0_{\mathrm I}$, 7.1.12). 따라서 $\rho$는 유일한 방식으로
  \[
    A\xrightarrow{j}B\xrightarrow{v}E
  \]
  로 인수분해된다. 또한 분명
  $u\circ j=(p\circ v)\circ j$이므로 $u=p\circ v$이다. 이로써 명제가
  증명된다.
  \item[(iii)] 체 $k$의 유한 분리가능 확대체 $k'$은 형식적으로
  매끄러운 $k$-대수이다 (19.6.1). 체 $k$의 \emph{유한 생성} 확대체
  $K$가 형식적으로 에탈인 $k$-대수이기 위한 필요충분조건은 $K$가
  형식적으로 비분기인 것이며, 이는 다시 $K$가 $k$의 유한 분리가능
  확대체인 것과 동치임을 뒤에서 볼 것이다 (21.7.4).
  \item[(iv)] 환 $A$ 위의 다항식 대수 $A[X_\alpha]_{\alpha\in I}$는
  $A$ 위에서 \emph{형식적으로 매끄럽지만} (19.3.3), $I$가 공집합이
  아니면 정의에서 곧 알 수 있듯이 \emph{형식적으로 에탈이지 않다}.
\end{enumerate}
\end{examples}

\begin{remark}[19.10.4]
\label{0.19.10.4-ko}
(19.4.3)에서와 같은 논법에 의해, 위상 $A$-대수 $B$가 형식적으로
비분기임을 확인할 때 정의 (19.10.2)에서 $p:E\to C$의 핵의
\emph{제곱이 $0$}인 경우만 생각해도 된다. 또한 $u(B)\subset p(E)$인
경우만 생각해도 된다. 그렇지 않으면 $u$를 인수분해하는 $v$가 전혀
없기 때문이다. 끝으로 $C$를 $u(B)$로, $E$를 $p^{-1}(u(B))$로 바꾸어
$u$와 $p$가 전사라고 가정할 수 있다.
\end{remark}

\section{미분과 미분 가군}
\label{section:0.20-ko}
\label{0.20-ko}

이 절에서 도입하는 개념들은 제IV장 \textsection16에서 기하학적 형태로
다시 다루며, 준스킴의 연구에서 중요한 역할을 할 것이다. 이 장에서
이 개념들이 중요한 첫 번째 이유는 형식적으로 매끄러운 대수라는 개념과
맺는 관계에 있다. 특히 정리 (20.4.9), (20.5.7), (20.5.12)는 제IV장에서
매끄러운 사상을 다루는 절에서 기하학적 언어로 옮겨지며, 응용에서
끊임없이 사용된다. 한편 미분에 관한 개념들은 \textsection22에서 중요한
정칙성 판정법들을 증명하는 데 쓰인다. 이 판정법들은 제IV장
\textsection7에서 이루어지는 뇌터 국소환의 심층 연구에서 핵심적인
역할을 할 것이다.

\subsection{미분과 대수의 확대}
\label{subsection:0.20.1-ko}

\oldpage[0\textsubscript{IV-1}]{117}
\begin{proposition}[20.1.1]
\label{0.20.1.1-ko}
$A$를 반드시 가환일 필요는 없는 환, $B$, $E$, $C$를 $A$-환들이라
하고, $p:E\to C$를 핵 $\mathfrak Q$의 제곱이 $0$인 $A$-준동형,
$u:B\to C$를 $A$-준동형이라 하자. $u$가
\[
  B\xrightarrow{v_0}E\xrightarrow{p}C
\]
로 인수분해되게 하는 $A$-준동형 $v_0:B\to E$가 존재한다고 하자.
그러면 $u=p\circ v$를 만족하는 $A$-준동형 $v:B\to E$들의 집합은
다음 두 조건을 만족하는 사상 $D:B\to\mathfrak Q$에 대한 사상
$v_0+D$들의 집합과 같다.
\begin{enumerate}
  \item[(i)] $D$는 $A$-쌍가군의 준동형이다.
  \item[(ii)] 임의의 $f,g\in B$에 대하여 다음이 성립한다.
\end{enumerate}
\begin{equation}
\tag{20.1.1.1}
\label{0.20.1.1.1-ko}
  D(fg)=f.D(g)+D(f).g.
\end{equation}
\end{proposition}

\begin{proof}
$f\in B$에 대하여 $p(v(f))=p(v_0(f))$라는 것은
$D(f)=v(f)-v_0(f)$가 $\mathfrak Q$에 속한다는 뜻이다.
$v(fg)=v(f)v(g)$를 쓰고 $\mathfrak Q^2=0$을 이용하면 관계식
(20.1.1.1)을 얻는다. 또한 $v$와 $v_0$가 $A$-준동형이므로 조건 (i)가
따른다.

$\rho:A\to B$를 구조 준동형이라 하자. (20.1.1.1)에서 모든
$a\in A$와 $f\in B$에 대하여
\[
  D(\rho(a)f)=\rho(a)D(f)+D(\rho(a))f
\]
를 얻는다. 그런데 (i)에 의해 $D(\rho(a)f)=\rho(a)D(f)$여야 하므로
$f=1$로 놓으면
\[
  D(\rho(a))=0\qquad\text{모든 $a\in A$에 대하여}
\]
를 얻는다. 역으로 $D$가 $\rho(A)$ 위에서 $0$이고 (20.1.1.1)을
만족하면 조건 (i)도 만족한다.
\end{proof}

\begin{definition}[20.1.2]
\label{0.20.1.2-ko}
$A$-환 $B$와 $B$-쌍가군 $L$이 주어졌다고 하자. (20.1.1)의 조건
(i), (ii)를 만족하는 사상 $D:B\to L$을 $B$에서 $L$로 가는
\emph{$A$-미분}이라 한다.
\end{definition}

(20.1.1)에서 $A$-미분 $D:B\to L$의 \emph{핵}은 $B$의 부분
$A$-환임을 알 수 있다.

때로는 $B$에서 $L$로 가는 \emph{미분}을 $\mathbf Z$-미분, 즉
(20.1.1.1)을 만족하는 $B$에서 $L$로 가는 \emph{가법} 사상으로
정의한다. 그러한 사상에 대해서는 $D(1)=0$이다. $B$가 \emph{소체}
$P$ 위의 대수이면 $B$의 모든 $\mathbf Z$-미분은 $P$-미분이다.
$P$의 표수가 양수이면 앞의 결과에서 바로 따르고, 그렇지 않으면
각 $0\neq n\in\mathbf Z$에 대한 관계식 $n.n^{-1}=1$에서
\[
  nD(n^{-1})=0
\]
을 얻는다. $P$의 표수가 $0$이므로 $D(n^{-1})=0$이다.

정의 (20.1.2)에서 곧 다음이 따른다. $D,D'$이 $B$에서 $L$로 가는
두 $A$-미분이면 $D-D'$도 $A$-미분이다. 다시 말해 $B$에서 $L$로
가는 $A$-미분들의 집합에는 \emph{가법군} 구조가 있으며, 이 군을
$\operatorname{D\acute{e}r}_A(B,L)$로 나타낸다. $A$가 \emph{가환},
$B$가 가환 $A$-\emph{대수}, $L$이 $B$-가군이면, 모든 $A$-미분
$D$와 모든 $a\in A$에 대하여 $aD$도 $A$-미분이다. 따라서 이 경우
$\operatorname{D\acute{e}r}_A(B,L)$에는 $A$-\emph{가군} 구조가 있다.

\oldpage[0\textsubscript{IV-1}]{118}
명제 \sref{0.20.1.1-ko}은 다음과 같이 해석할 수 있다.

\begin{corollary}[20.1.3]
\label{0.20.1.3-ko}
$A$-환의 두 $A$-준동형 $p:E\to C$, $u:B\to C$가 주어지고, 첫째
준동형의 핵 $\mathfrak Q$의 제곱이 $0$이라고 하자. $u=p\circ v$를
만족하는 $A$-준동형 $v:B\to E$들의 집합은 공집합이거나 군
$\operatorname{D\acute{e}r}_A(B,\mathfrak Q)$에 대한 주동차공간이다.
\end{corollary}

특히 다음이 성립한다.

\begin{corollary}[20.1.4]
\label{0.20.1.4-ko}
$A$를 환, $C$를 $A$-환, $L$을 $C$-쌍가군, $E$를 $L$에 의한
$C$의 $A$-확장이라 하고, $p:E\to C$를 증대라 하자. 각 미분
$D\in\operatorname{D\acute{e}r}_A(C,L)$에 사상
$x\mapsto x+D(p(x))$를 대응시키는 사상은 군
$\operatorname{D\acute{e}r}_A(C,L)$에서 $E$의 자기 $A$-동치사상들의
군으로 가는 동형사상이다.
\end{corollary}

\begin{proof}
\sref{0.20.1.1-ko}을 $B=E$, $v_0=\mathrm I_E$에 적용하자.
$p\circ v=p$를 만족하는 $A$-준동형 $v:E\to E$들의 집합은
$D'\in\operatorname{D\acute{e}r}_A(E,L)$인 사상
$v=\mathrm I_E+D'$들의 집합과 같다. 그러한 $A$-준동형 $v$가
$A$-동치라는 것은 $v$를 $L$에 제한하면 $L$에서 $E$로 가는 표준
포함사상이 된다는 것, 달리 말해 모든 $x\in L$에 대하여 $L$에서
$D'(x)=0$이라는 것과 동치이다. 이는 다시 $D'$이 어떤 $A$-미분
$D$에 대하여
\[
  E\xrightarrow{p}C\xrightarrow{D}L
\]
로 인수분해된다는 뜻이다. 따라서 따름정리가 성립한다.
\end{proof}

자명한 확장에 대해서는 \sref{0.20.1.3-ko}으로부터 다음을 얻는다.

\begin{corollary}[20.1.5]
\label{0.20.1.5-ko}
$A$를 환, $B$, $C$를 두 $A$-환, $L$을 $C$-쌍가군,
$u:B\to C$를 $A$-준동형이라 하자. 각 미분
$D\in\operatorname{D\acute{e}r}_A(B,L)$에 사상
$\alpha:x\mapsto(u(x),D(x))$를 대응시키는 사상은 $u$가
\[
  B\xrightarrow{v}D_C(L)\longrightarrow C
\]
로 인수분해되게 하는 $A$-준동형 $v:B\to D_C(L)$들의 집합
(18.2.3)으로 가는 전단사이다.
\end{corollary}

더 특별히 다음이 성립한다.

\begin{corollary}[20.1.6]
\label{0.20.1.6-ko}
$A$를 환, $B$를 $A$-환, $L$을 $B$-쌍가군이라 하자. 각 미분
$D\in\operatorname{D\acute{e}r}_A(B,L)$에 다음 두 사상을 대응시키자.
$1^\circ$ 확장 $D_B(L)$의 자기 $A$-동치사상
$(x,y)\mapsto(x,y+D(x))$; $2^\circ$ 증대 준동형
$D_B(L)\to B$의 오른쪽 역인 $A$-준동형
$x\mapsto(x,D(x))$ ($B\to D_B(L)$). 그러면 다음 세 집합 사이에
표준 전단사들이 정의된다.
\begin{enumerate}
  \item[(i)] 집합 $\operatorname{D\acute{e}r}_A(B,L)$;
  \item[(ii)] $D_B(L)$의 자기 $A$-동치사상들의 집합;
  \item[(iii)] 증대 준동형 $D_B(L)\to B$의 오른쪽 역인
  $A$-준동형들의 집합.
\end{enumerate}
\end{corollary}

이렇게 얻은 (i)와 (ii) 사이의 전단사는 군 구조들을 보존하며,
여기서 유도되는 (ii)와 (iii) 사이의 전단사는 (18.3.8)에서 이미
정의한 전단사와 다름없다.

\begin{corollary}[20.1.7]
\label{0.20.1.7-ko}
$A$를 환, $C$를 $A$-환, $L$을 $C$-쌍가군, $E$를 $L$에 의한
$C$의 자명한 $A$-확장이라 하고, $p:E\to C$를 증대라 하자.
모든 $x\in L$에 대하여 $L$에서 $D(x)=x$를 만족하는 $E$에서
$L$로 가는 $A$-미분 $D$들의 집합은 사상
$x\mapsto x-w(p(x))$들의 집합과 같다. 여기서 $w$는 $p$의 오른쪽
역인 $A$-준동형 전체를 달린다.
\end{corollary}

\begin{proof}
다시 \sref{0.20.1.1-ko}을 $B=E$, $v_0=\mathrm I_E$에 적용하면
충분하다. $v=\mathrm I_E-D$라 하자. 모든 $x\in L$에 대하여
$D(x)=x$라는 조건은 모든 $x\in L$에 대하여 $v(x)=0$이라는 것과
동치이며, 이는 어떤 $A$-준동형 $w:C\to E$에 대하여
\[
  v=w\circ p
\]
로 인수분해된다는 뜻이다. 또한 조건 $p\circ v=p$는
$p\circ w=\mathrm I_C$와 동치이다. 다시 말해 $w$는 $p$의 오른쪽
역이다.
\end{proof}

\subsection{미분의 함자적 성질}
\label{subsection:0.20.2-ko}

\oldpage[0\textsubscript{IV-1}]{119}
\begin{env}[20.2.1]
\label{0.20.2.1-ko}
$A$를 환, $B$를 $A$-환, $L$을 $B$-쌍가군이라 하자. $L'$이 또
다른 $B$-쌍가군이고 $w:L\to L'$이 $B$-쌍가군의 준동형이면, 사상
$D\mapsto w\circ D$는 가법군의 준동형
\begin{equation}
\tag{20.2.1.1}
\label{0.20.2.1.1-ko}
  w_0:\operatorname{D\acute{e}r}_A(B,L)
  \longrightarrow\operatorname{D\acute{e}r}_A(B,L')
\end{equation}
임이 명백하다. $w':L'\to L''$이 또 다른 $B$-쌍가군의 준동형이면
$(w'\circ w)_0=w'_0\circ w_0$이다. $A$가 가환이고 $B$가 가환
$A$-대수이며 $L$이 $B$-가군일 때, (20.2.1.1)은 $B$-가군의
준동형이다.

둘째로 $B'$을 $A$-환이라 하고, $v:B'\to B$를 $L$에 $B'$-쌍가군
구조를 주는 $A$-준동형이라 하자. 그러면 사상 $D\mapsto D\circ v$는
가법군의 준동형
\begin{equation}
\tag{20.2.1.2}
\label{0.20.2.1.2-ko}
  v^0:\operatorname{D\acute{e}r}_A(B,L)
  \longrightarrow\operatorname{D\acute{e}r}_A(B',L)
\end{equation}
이다. 이는 \sref{0.20.1.1.1-ko}에서 따른다. $v':B''\to B'$이 또
다른 $A$-준동형이면 $(v\circ v')^0=v'^0\circ v^0$이다. $A$, $B$,
$B'$이 가환이고 $L$이 $B$-가군일 때, (20.2.1.2)는 $B$-가군의
준동형이다.

끝으로 $u:A'\to A$를 $B$에 $A'$-환 구조를 주는 환 준동형이라
하자. 모든 $A$-미분 $D\in\operatorname{D\acute{e}r}_A(B,L)$은
$A'$-미분이기도 하므로 가환군의 표준 단사 준동형
\begin{equation}
\tag{20.2.1.3}
\label{0.20.2.1.3-ko}
  u^0:\operatorname{D\acute{e}r}_A(B,L)
  \longrightarrow\operatorname{D\acute{e}r}_{A'}(B,L)
\end{equation}
을 얻는다. $u':A''\to A'$이 또 다른 환 준동형이면
$(u\circ u')^0=u'^0\circ u^0$이다. $A$, $A'$, $B$가 가환이고
$L$이 $B$-가군일 때, (20.2.1.3)은 $B$-가군의 준동형이다.

이를 달리 말하면
\[
  (A,B,L)\leadsto\operatorname{D\acute{e}r}_A(B,L)
\]
은 (18.3.5)에서 정의한 범주 $\boldsymbol K$에서 가환군의 범주
$\boldsymbol{Ab}$로 가는 \emph{공변 함자}이며, $\boldsymbol K$의
사상을 이루는 각 삼중항 $(u,v,w)$에는 준동형
$w_0\circ v^0\circ u^0$를 대응시킨다. 함자성은 임의의
$B$-쌍가군 준동형 $w:L\to L'$에 대하여 다음 두 도식이 가환한다는
데서 따른다.
\[
\xymatrix{
  \operatorname{D\acute{e}r}_A(B,L) \ar[r]^{v^0} \ar[d]_{w_0} &
    \operatorname{D\acute{e}r}_A(B',L) \ar[d]^{w_0} \\
  \operatorname{D\acute{e}r}_A(B,L') \ar[r]^{v^0} &
    \operatorname{D\acute{e}r}_A(B',L')
}
\qquad
\xymatrix{
  \operatorname{D\acute{e}r}_A(B,L) \ar[r]^{u^0} \ar[d]_{w_0} &
    \operatorname{D\acute{e}r}_{A'}(B,L) \ar[d]^{w_0} \\
  \operatorname{D\acute{e}r}_A(B,L') \ar[r]^{u^0} &
    \operatorname{D\acute{e}r}_{A'}(B,L')
}
\]
\end{env}

\oldpage[0\textsubscript{IV-1}]{120}
\begin{theorem}[20.2.2]
\label{0.20.2.2-ko}
$u:A\to B$, $v:B\to C$를 두 환 준동형이라 하고, $L$을
$C$-쌍가군이라 하자. 다음과 같은 가환군의 표준 완전열이 있다.
\begin{align}
\tag{20.2.2.1}
\label{0.20.2.2.1-ko}
  0\to\operatorname{D\acute{e}r}_B(C,L)\xrightarrow{u^0}
  \operatorname{D\acute{e}r}_A(C,L)&\xrightarrow{v^0}
  \operatorname{D\acute{e}r}_A(B,L)\xrightarrow{\partial} \\
  &\to\operatorname{Exan}_B(C,L)\xrightarrow{u^1}
  \operatorname{Exan}_A(C,L)\xrightarrow{v^1}
  \operatorname{Exan}_A(B,L) \notag
\end{align}
여기서 $u^0$, $v^0$는 각각 (20.2.1.3), (20.2.1.2)의 준동형이고,
$u^1$, $v^1$은 각각 (18.3.4.1), (18.3.3.1)에서 정의한 준동형이다.
$\partial$는 다음과 같이 정의한다. $B$에서 $L$로 가는 각 $A$-미분
$D$에 대하여 $\partial(D)$는 $A$-준동형
$\alpha:x\mapsto(v(x),D(x))$가 환 $D_C(L)$ 위에 정하는, $L$에 의한
$C$의 $B$-확장의 류이다 (\sref{0.20.1.5-ko} 참조). 또한 완전열
\sref{0.20.2.2.1-ko}은 $L$에 관하여 함자적이다. 여기서는 각각
(20.2.1.1), (18.3.1.1)에서 정의한 준동형들을 사용한다.
\end{theorem}

\begin{proof}
$D$는 $B$에서 $L$로 가는 $A$-미분이고, 따라서 특히
$\mathbf Z$-미분이므로, $A$-준동형
$\alpha:x\mapsto(v(x),D(x))$는 $D_C(L)$ 위에 실제로 $B$-확장 구조를
정의한다. 따라서 $\partial(D)$는 잘 정의된다
(\sref{0.20.1.5-ko}). 이제 다섯 곳에서 완전성을 확인해야 한다.

\medskip
\noindent 1)\enspace
$\operatorname{D\acute{e}r}_B(C,L)$에서의 완전성은 자명하다
(\sref{0.20.2.1-ko} 참조).

\medskip
\noindent 2)\enspace
정의 \sref{0.20.2.1-ko}에 따라 $v^0$의 핵은 $v(B)$ 위에서 $0$인,
$C$에서 $L$로 가는 $A$-미분들의 집합이다. 이는 곧 $B$-미분이기도
한 $A$-미분들의 집합이다 (\sref{0.20.1.1-ko}). 따라서
$\operatorname{D\acute{e}r}_A(C,L)$에서 완전하다.

\medskip
\noindent 3)\enspace
$\partial$의 핵은 $\alpha:x\mapsto(v(x),D(x))$가 정하는 $B$-확장이
$B$-자명하게 되는 미분
$D\in\operatorname{D\acute{e}r}_A(B,L)$들로 이루어진다. 이는
(18.2.3)에 따라, $\alpha$가 $D_C(L)$에 정하는 $B$-환 구조에 관하여
$C$에서 $D_C(L)$로 가는 $B$-준동형 $z\mapsto(z,w(z))$가 존재한다는
뜻이다. 그런데 그러한 준동형은 특히 $A$-준동형이므로
\sref{0.20.1.6-ko}에 따라
$D'\in\operatorname{D\acute{e}r}_A(C,L)$인
$z\mapsto(z,D'(z))$의 꼴이다. 이것이 $B$-준동형이라는 조건을 쓰면
모든 $x\in B$, $z\in C$에 대하여
\[
  D(x)z+v(x)D'(z)=D'(v(x)z)=v(x)D'(z)+D'(v(x))z
\]
이고, 따라서 $D'(v(x))=D(x)$이다. 그러므로 $\partial$의 핵은
$v^0$의 상이다.

\medskip
\noindent 4)\enspace
$u^1$의 핵은 $A$-자명한, $L$에 의한 $C$의 $B$-확장들의 류들로
이루어진다 (18.3.7). 따라서 동치를 제외하면 이 확장들은
$D_C(L)$의 꼴이며, 그 $A$-환 구조는 준동형
$t\mapsto(v(u(t)),0)$으로 정해진다. 한편 $D_C(L)$ 위의 모든
$B$-확장 구조는 $B$에서 $L$로 가는 어떤 $\mathbf Z$-미분 $D$에
대한 준동형 $\alpha:x\mapsto(v(x),D(x))$로 정해진다
(\sref{0.20.1.5-ko}). 이 $B$-확장의 $A$-환 구조가 $u:A\to B$를
통해 그 $B$-환 구조에서 유도된다는 것은 모든 $t\in A$에 대하여
$D(u(t))=0$이라는 뜻이다. 이는 $D$가 $A$-미분이라는 것, 달리 말해
고려한 $B$-확장의 류가 $\partial(D)$의 꼴이라는 것과 같다. 따라서
$\operatorname{Exan}_B(C,L)$에서 완전하다.

\medskip
\noindent 5)\enspace
$v^1$의 핵은 $v(B)$ 위에서 자명하게 되는, $L$에 의한 $C$의
$A$-확장 $E$들의 류들로 이루어진다. 즉 $v$가
\[
  B\xrightarrow{w}E\longrightarrow C
\]
로 인수분해되게 하는 $A$-준동형 $w:B\to E$가 존재하는 확장들의
류이다. 그런데 그러한 $A$-준동형은 $E$ 위에 $B$-확장 구조를
정의하고, 그 류를 $u^1$로 보낸 상은 주어진 $A$-확장의 류이다.
역은 자명하므로 $\operatorname{Exan}_A(C,L)$에서 완전함이
증명되었다.

끝으로 $L$에 관한 함자성은 정의들에서 곧바로 따른다.
\end{proof}

\oldpage[0\textsubscript{IV-1}]{121}
\begin{corollary}[20.2.3]
\label{0.20.2.3-ko}
$A$, $B$, $C$를 세 가환환, $u:A\to B$, $v:B\to C$를 두 환
준동형, $L$을 $C$-가군이라 하자. 다음과 같은 $A$-가군의 표준
완전열이 있다.
\begin{align}
\tag{20.2.3.1}
\label{0.20.2.3.1-ko}
  0\to\operatorname{D\acute{e}r}_B(C,L)\xrightarrow{u^0}
  \operatorname{D\acute{e}r}_A(C,L)&\xrightarrow{v^0}
  \operatorname{D\acute{e}r}_A(B,L)\xrightarrow{\partial} \\
  &\to\operatorname{Exalcom}_B(C,L)\xrightarrow{u^1}
  \operatorname{Exalcom}_A(C,L)\xrightarrow{v^1}
  \operatorname{Exalcom}_A(B,L) \notag
\end{align}
이 완전열은 $L$에 관하여 함자적이다.
\end{corollary}

\begin{proof}
각 미분 $D\in\operatorname{D\acute{e}r}_A(B,L)$에 대하여
$\partial(D)$가 실제로 $L$에 의한 $C$의 $B$-확장으로서 가환
$B$-대수인 것의 류임을 확인하고 나면, 논법은
\sref{0.20.2.2-ko}과 같다. 그런데 이는 $C$의 가환성과 $L$이
$C$-가군이라는 사실에서 곧바로 따른다.
\end{proof}

\begin{corollary}[20.2.4]
\label{0.20.2.4-ko}
\sref{0.20.2.2-ko}의 가정 아래에서 (각각
\sref{0.20.2.3-ko}의 가정 아래에서), 다음과 같은 $L$에 관하여
함자적인 표준 완전열이 있다.
\begin{equation}
\tag{20.2.4.1}
\label{0.20.2.4.1-ko}
  0\to\operatorname{D\acute{e}r}_B(C,L)\xrightarrow{u^0}
  \operatorname{D\acute{e}r}_A(C,L)\xrightarrow{v^0}
  \operatorname{D\acute{e}r}_A(B,L)\xrightarrow{\partial}
  \operatorname{Exan}_{B/A}(C,L)\to0
\end{equation}
(각각
\begin{equation}
\tag{20.2.4.2}
\label{0.20.2.4.2-ko}
  0\to\operatorname{D\acute{e}r}_B(C,L)\xrightarrow{u^0}
  \operatorname{D\acute{e}r}_A(C,L)\xrightarrow{v^0}
  \operatorname{D\acute{e}r}_A(B,L)\xrightarrow{\partial}
  \operatorname{Exalcom}_{B/A}(C,L)\to0\text{).}
\end{equation}
이는 $\operatorname{Exan}_{B/A}(C,L)$ (각각
$\operatorname{Exalcom}_{B/A}(C,L)$)의 정의에서 따른다
((18.3.7), (18.4.2)).
\end{corollary}

\begin{remark}[20.2.5]
\label{0.20.2.5-ko}
환 준동형들의 다음 가환 도식이 주어졌다고 하자.
\[
\xymatrix{
  A \ar[r] & B \ar[r] & C \\
  A' \ar[u] \ar[r] & B' \ar[u] \ar[r] & C' \ar[u]
}
\]
그러면 다음 도식도 가환한다.
\[
\xymatrix@C=0.5em{
  0 \ar[r] & \operatorname{D\acute{e}r}_B(C,L) \ar[r] \ar[d] &
  \operatorname{D\acute{e}r}_A(C,L) \ar[r] \ar[d] &
  \operatorname{D\acute{e}r}_A(B,L) \ar[r]^-\partial \ar[d] &
  \operatorname{Exan}_B(C,L) \ar[r] \ar[d] &
  \operatorname{Exan}_A(C,L) \ar[r] \ar[d] &
  \operatorname{Exan}_A(B,L) \ar[d] \\
  0 \ar[r] & \operatorname{D\acute{e}r}_{B'}(C',L) \ar[r] &
  \operatorname{D\acute{e}r}_{A'}(C',L) \ar[r] &
  \operatorname{D\acute{e}r}_{A'}(B',L) \ar[r]^-\partial &
  \operatorname{Exan}_{B'}(C',L) \ar[r] &
  \operatorname{Exan}_{A'}(C',L) \ar[r] &
  \operatorname{Exan}_{A'}(B',L)
}
\]
완전열 \sref{0.20.2.3.1-ko}, \sref{0.20.2.4.1-ko},
\sref{0.20.2.4.2-ko}에 대해서도 마찬가지이다.
\end{remark}

\subsection{위상환의 연속 미분}
\label{subsection:0.20.3-ko}

\begin{env}[20.3.1]
\label{0.20.3.1-ko}
두 위상환 $A$, $B$가 주어졌다고 하자. 늘 그렇듯 이들은 선형
위상환이다. $A$에서 $B$로 가는 연속 준동형들의 집합을
$\operatorname{Hom.\ cont}(A,B)$로 나타내겠다. 위상환 $A$가 주어졌을
때, 위상 $A$-환들의 범주는 $A$-환들의 범주 (18.1.4)의 정의에서
어디서나 ``환''을 ``위상환''으로, ``준동형''을 ``연속 준동형''으로
바꾸어 정의한다. $B$와

\oldpage[0\textsubscript{IV-1}]{122}
$C$가 두 위상 $A$-환이면, $B$에서 $C$로 가는 연속
$A$-준동형들의 집합을 $\operatorname{Hom.\ cont}_A(B,C)$로 나타낸다.

$A$를 위상환, $B$를 위상 $A$-환, $L$을 위상 $B$-쌍가군이라 하자.
$B$에서 $L$로 가는 \emph{연속} $A$-미분들의 집합을
$\operatorname{D\acute{e}r.\ cont}_A(B,L)$로 나타낸다. 이는 분명
$\operatorname{D\acute{e}r}_A(B,L)$의 \emph{부분군}이다. 또한 $A$와
$B$가 가환이고 $L$이 $B$-가군일 때에는 부분 $A$-가군이다.
\end{env}

\begin{env}[20.3.2]
\label{0.20.3.2-ko}
명제 \sref{0.20.1.1-ko}에서 ``환''을 ``위상환''으로,
``준동형''을 ``연속 준동형''으로 바꾸어도 그 명제는 그대로 성립한다.
마찬가지로 \sref{0.20.1.3-ko}, \sref{0.20.1.4-ko}에서
$\operatorname{D\acute{e}r}$를 $\operatorname{D\acute{e}r.\ cont}$로,
``환'' (각각 ``쌍가군'')을 ``위상환'' (각각 ``위상 쌍가군'')으로,
``준동형''을 ``연속 준동형''으로 바꾸어도 결과들이 성립한다. 물론
\sref{0.20.1.4-ko}에서는 $p$가 연속이라고 가정하고
``$A$-동치사상''을 ``연속 $A$-동치사상''으로 바꾸어야 한다.
\sref{0.20.1.5-ko}, \sref{0.20.1.6-ko}에 대해서도 같은 결과가
성립한다. 다만 $D_C(L)$ (각각 $D_B(L)$)에는 $C\times L$ (각각
$B\times L$)의 \emph{곱위상}을 주어야 한다.
\end{env}

\begin{proposition}[20.3.3]
\label{0.20.3.3-ko}
$A$를 위상환, $B$를 위상 $A$-환, $L$을 $B$의 어떤 열린 양쪽
아이디얼에 의해 소멸되는 이산 위상 $B$-쌍가군이라 하자. $B$의 모든
열린 양쪽 아이디얼의 제곱이 열려 있으면
\[
  \operatorname{D\acute{e}r.\ cont}_A(B,L)
  =\operatorname{D\acute{e}r}_A(B,L).
\]
\end{proposition}

\begin{proof}
실제로 $\mathfrak R$을 $L$을 소멸시키는 $B$의 열린 양쪽 아이디얼이라
하고, $D$를 $B$에서 $L$로 가는 $A$-미분이라 하자. 그러면
\sref{0.20.1.1.1-ko}에 의해
\[
  D(\mathfrak R^2)\subset
  \mathfrak R.D(\mathfrak R)+D(\mathfrak R).\mathfrak R=0.
\]
따라서 $D$는 연속이다.
\end{proof}

\begin{env}[20.3.4]
\label{0.20.3.4-ko}
\sref{0.20.2.1-ko}의 모든 결과에서 ``환''을 ``위상환''으로,
``쌍가군''을 ``위상 쌍가군''으로, ``준동형''을 ``연속 준동형''으로,
$\operatorname{D\acute{e}r}$를 $\operatorname{D\acute{e}r.\ cont}$로
바꾸어도 그대로 성립한다.
\end{env}

\begin{proposition}[20.3.5]
\label{0.20.3.5-ko}
$A$를 위상환, $B$를 위상 $A$-환, $L$을 $B$의 어떤 열린 양쪽
아이디얼에 의해 소멸되는 이산 위상 $B$-쌍가군이라 하자. 그러면
표준 동형사상
\begin{equation}
\tag{20.3.5.1}
\label{0.20.3.5.1-ko}
  \varinjlim\operatorname{D\acute{e}r}_{A/\mathfrak J}
  (B/\mathfrak R,L)\xrightarrow{\sim}
  \operatorname{D\acute{e}r.\ cont}_A(B,L)
\end{equation}
이 있다. 왼쪽의 귀납극한은 다음 조건들을 만족하는 양쪽 아이디얼의
쌍 $(\mathfrak J,\mathfrak R)$들의 여과 순서집합을 따라 취한다.
\[
  \mathfrak R.L=L.\mathfrak R=0,\qquad
  \mathfrak J.B\subset\mathfrak R,\qquad
  B.\mathfrak J\subset\mathfrak R.
\]
\end{proposition}

\begin{proof}
$A/\mathfrak J$와 $B/\mathfrak R$은 이산이므로, 표준 준동형
\[
  w_{\mathfrak J,\mathfrak R}:
  \operatorname{D\acute{e}r}_{A/\mathfrak J}(B/\mathfrak R,L)
  \longrightarrow\operatorname{D\acute{e}r.\ cont}_A(B,L)
\]
들이 있고, 이들은 귀납계를 이룬다 (\sref{0.20.3.4-ko}). 따라서
귀납극한을 취하면 준동형 \sref{0.20.3.5.1-ko}을 얻는다.

$\mathfrak R\supset\mathfrak R'$일 때 준동형
$B/\mathfrak R'\to B/\mathfrak R$은 전사이므로, 정의에서 곧바로
준동형
\[
  \operatorname{D\acute{e}r}_{A/\mathfrak J}(B/\mathfrak R,L)
  \longrightarrow
  \operatorname{D\acute{e}r}_{A/\mathfrak J}(B/\mathfrak R',L)
\]
이 단사임을 알 수 있다. 여기서는
$\mathfrak R.L=L.\mathfrak R=0$이고,
$\mathfrak J.B\subset\mathfrak R'$,
$B.\mathfrak J\subset\mathfrak R'$이다. 마찬가지로
$\mathfrak J'\subset\mathfrak J$일 때 준동형
\[
  \operatorname{D\acute{e}r}_{A/\mathfrak J}(B/\mathfrak R,L)
  \longrightarrow
  \operatorname{D\acute{e}r}_{A/\mathfrak J'}(B/\mathfrak R,L)
\]
도 분명 단사이다 (\sref{0.20.2.1-ko}). 그러므로 준동형
\sref{0.20.3.5.1-ko}은 단사이다.

다른 한편 $D$가 $B$에서 $L$로 가는 연속 $A$-미분이면, 그 핵은
$B$의 어떤 열린 양쪽 아이디얼 $\mathfrak R_0$을 포함한다.
$\mathfrak J_0B\subset\mathfrak R_0$이고
$B\mathfrak J_0\subset\mathfrak R_0$인 $A$의 열린 양쪽 아이디얼
$\mathfrak J_0$을 택하면, $D$는 $B/\mathfrak R_0$에서 $L$로 가는
어떤 $(A/\mathfrak J_0)$-미분의 표준 상임이 명백하다. 따라서
\sref{0.20.3.5.1-ko}은 전사이다.
\end{proof}

\oldpage[0\textsubscript{IV-1}]{123}
\begin{proposition}[20.3.6]
\label{0.20.3.6-ko}
$u:A\to B$, $v:B\to C$를 위상환의 두 연속 준동형이라 하고,
$L$을 $C$의 어떤 열린 양쪽 아이디얼에 의해 소멸되는 이산
$C$-쌍가군이라 하자. 다음과 같은 표준 완전열이 있다.
\begin{align}
\tag{20.3.6.1}
\label{0.20.3.6.1-ko}
  0\to\operatorname{D\acute{e}r.\ cont}_B(C,L)\xrightarrow{u^0}
  \operatorname{D\acute{e}r.\ cont}_A(C,L)&\xrightarrow{v^0}
  \operatorname{D\acute{e}r.\ cont}_A(B,L)\xrightarrow{\partial} \\
  &\to\operatorname{Exantop}_B(C,L)\xrightarrow{u^1}
  \operatorname{Exantop}_A(C,L)\xrightarrow{v^1}
  \operatorname{Exantop}_A(B,L) \notag
\end{align}
여기서 $\partial$는 \sref{0.20.2.2.1-ko}의 준동형 $\partial$에서
귀납극한을 취하여 정의한다. 이 완전열은 열린 양쪽 아이디얼들에 의해
소멸되는 이산 $C$-쌍가군들의 범주에서 $L$에 관하여 함자적이다.
\end{proposition}

\begin{proof}
이는 \sref{0.20.2.2-ko}과 함자 $\varinjlim$의 완전성에서 따른다.
\end{proof}

\begin{corollary}[20.3.7]
\label{0.20.3.7-ko}
$A$, $B$, $C$를 세 가환 위상환, $u:A\to B$, $v:B\to C$를 두
연속 준동형, $L$을 $C$의 어떤 열린 아이디얼에 의해 소멸되는 이산
$C$-가군이라 하자. 다음과 같은 $L$에 관하여 함자적인 $A$-가군의
표준 완전열이 있다.
\begin{align}
\tag{20.3.7.1}
\label{0.20.3.7.1-ko}
  0\to\operatorname{D\acute{e}r.\ cont}_B(C,L)\xrightarrow{u^0}
  \operatorname{D\acute{e}r.\ cont}_A(C,L)&\xrightarrow{v^0}
  \operatorname{D\acute{e}r.\ cont}_A(B,L)\xrightarrow{\partial} \\
  &\to\operatorname{Exalcotop}_B(C,L)\xrightarrow{u^1}
  \operatorname{Exalcotop}_A(C,L)\xrightarrow{v^1}
  \operatorname{Exalcotop}_A(B,L). \notag
\end{align}
\end{corollary}

\begin{corollary}[20.3.8]
\label{0.20.3.8-ko}
\sref{0.20.3.6-ko}의 가정 아래에서 (각각
\sref{0.20.3.7-ko}의 가정 아래에서), 다음과 같은 $L$에 관하여
함자적인 표준 완전열이 있다.
\begin{equation}
\tag{20.3.8.1}
\label{0.20.3.8.1-ko}
  0\to\operatorname{D\acute{e}r.\ cont}_B(C,L)\xrightarrow{u^0}
  \operatorname{D\acute{e}r.\ cont}_A(C,L)\xrightarrow{v^0}
  \operatorname{D\acute{e}r.\ cont}_A(B,L)\to
  \operatorname{Exantop}_{B/A}(C,L)\to0
\end{equation}
(각각
\begin{equation}
\tag{20.3.8.2}
\label{0.20.3.8.2-ko}
  0\to\operatorname{D\acute{e}r.\ cont}_B(C,L)\xrightarrow{u^0}
  \operatorname{D\acute{e}r.\ cont}_A(C,L)\xrightarrow{v^0}
  \operatorname{D\acute{e}r.\ cont}_A(B,L)\to
  \operatorname{Exalcotop}_{B/A}(C,L)\to0\text{).}
\end{equation}
\end{corollary}

\sref{0.20.2.5-ko}의 도식들과 유사한 도식들을 적는 일은 독자에게
맡긴다.

\subsection{주요 부분과 미분}
\label{subsection:0.20.4-ko}

이 항의 나머지와 이어지는 세 항에서는 모든 환이 \emph{가환}이라고
가정한다.

\begin{env}[20.4.1]
\label{0.20.4.1-ko}
$A$를 위상환, $B$를 위상 $A$-대수라 하자. $A$-대수
$B\otimes_A B$에는 \emph{텐서곱 위상}을 주며, 이로써 이는 위상
$A$-대수가 된다. $p$ (또는 $p_{B/A}$)로 다음 표준 전사
$A$-준동형을 나타내겠다.
\begin{equation}
\tag{20.4.1.1}
\label{0.20.4.1.1-ko}
  p:B\otimes_A B\to B,
  \qquad p(b\otimes b')=bb'.
\end{equation}
$p$가 연속임은 명백하다. $p$의 \emph{핵}을
$\mathfrak J_{B/A}$로 나타내며, 혼동의 우려가 없으면 단순히
$\mathfrak J$로 나타낸다. 또한 두 표준 $A$-준동형
$j_1,j_2:B\to B\otimes_A B$를
\[
  j_1(b)=b\otimes1,\qquad j_2(b)=1\otimes b
\]
로 정의한다. 이들은 연속이다.
\end{env}

\oldpage[0\textsubscript{IV-1}]{124}
\begin{definition}[20.4.2]
\label{0.20.4.2-ko}
위상 $A$-대수 몫
\begin{equation}
\tag{20.4.2.1}
\label{0.20.4.2.1-ko}
  \mathrm P^1_{B/A}=(B\otimes_A B)/\mathfrak J^2
\end{equation}
에 다음 구조를 준 것을 $B$의 $A$에 대한 \emph{1차 주요 부분의
증대 $B$-대수}라 하고 $\mathrm P^1_{B/A}$로 나타낸다. 그 위상
$B$-대수 구조는 $j_1$과 표준 준동형
$B\otimes_A B\to\mathrm P^1_{B/A}$의 합성에서 얻는 준동형
$\bar j_1:B\to\mathrm P^1_{B/A}$가 정의하며, $B$-대수의 증대
$\varepsilon:\mathrm P^1_{B/A}\to B$ (또는 $\varepsilon_{B/A}$)는
$p$를 몫으로 내려 얻는다.

정의상 $p(b\otimes1)=b$이므로 $\varepsilon$가 실제로 $B$-대수의
증대임은 명백하다.
\end{definition}

\begin{definition}[20.4.3]
\label{0.20.4.3-ko}
증대 $\varepsilon:\mathrm P^1_{B/A}\to B$의 핵
\begin{equation}
\tag{20.4.3.1}
\label{0.20.4.3.1-ko}
  \Omega^1_{B/A}=\mathfrak J_{B/A}/(\mathfrak J_{B/A})^2
\end{equation}
에 $\mathrm P^1_{B/A}$의 위상이 유도하는 위상을 주면 이는 위상
$B$-가군이 된다. 이를 $B$의 $A$에 대한 \emph{1차 미분 가군}
(또는 단순히 \emph{미분 가군})이라 한다.
\end{definition}

$\Omega^1_{B/A}$의 위상은 $B\otimes_A B$의 위상이
$\mathfrak J_{B/A}$에 유도하는 위상의 몫위상이기도 하다 (Bourbaki,
\emph{Top. gén.}, chap.~III, 3\textsuperscript{e} éd., §~2,
n\textsuperscript{o}~7, prop.~20). $B$가 이산이면
$\Omega^1_{B/A}$도 이산이다. 위상 $B$-가군 $\Omega^1_{B/A}$의
\emph{분리 완비화}를 $\widehat{\Omega}^1_{B/A}$로 나타낸다.

모든 위상환 $B$는 이산 위상을 준 $\mathbf Z$ 위의 위상 대수로
볼 수 있으므로 위상 $B$-가군 $\Omega^1_{B/\mathbf Z}$를 정의할 수
있다. 이를 때로 $B$ 위의 \emph{절대 미분 가군}이라 하며
$\Omega^1_B$로도 나타낸다. $B$가 이산 소체 $P$ 위의 위상 대수이면
\[
  B\otimes_PB=B\otimes_{\mathbf Z}B,
  \qquad \Omega^1_{B/P}=\Omega^1_B.
\]

\begin{lemma}[20.4.4]
\label{0.20.4.4-ko}
$A$를 환, $B$를 $A$-대수라 하자. $B\otimes_A B$의 아이디얼
$\mathfrak J_{B/A}$는 $A$-대수 $B$의 한 생성원 집합을 $s$가 달릴
때의 원소
\[
  1\otimes s-s\otimes1
\]
들로 생성된다.
\end{lemma}

\begin{proof}
모든 $x\in B$에 대하여
$x\otimes1-1\otimes x\in\mathfrak J$임은 명백하다. 한편
$x,y\in B$이면
\[
  x\otimes y=xy\otimes1+(x\otimes1)(1\otimes y-y\otimes1).
\]
$\sum_i(x_i\otimes y_i)\in\mathfrak J$이면 정의상
$\sum_i x_iy_i=0$이므로
\begin{equation}
\tag{20.4.4.1}
\label{0.20.4.4.1-ko}
  \sum_i(x_i\otimes y_i)
  =\sum_i(x_i\otimes1)(1\otimes y_i-y_i\otimes1).
\end{equation}
따라서 $\mathfrak J$는
$1\otimes x-x\otimes1$ 꼴의 원소들로 생성된다. 더욱이 $x=st$이면
\begin{equation}
\tag{20.4.4.2}
\label{0.20.4.4.2-ko}
  x\otimes1-1\otimes x
  =(s\otimes1)(t\otimes1-1\otimes t)
   +(s\otimes1-1\otimes s)(1\otimes t),
\end{equation}
이므로 귀납법으로 증명이 곧 끝난다.
\end{proof}

\begin{proposition}[20.4.5]
\label{0.20.4.5-ko}
$A$를 위상환, $B$를 위상 $A$-대수라 하자.
$\Omega^1_{B/A}$의 위상은 $B$의 위상으로부터 얻어지는 위상
\sref{0.19.0.2-ko}보다 섬세하지 않다. $B$의 모든 열린 아이디얼의
제곱이 열려 있으면 이 두 위상은 일치한다.
\end{proposition}

\begin{proof}
첫째 주장은 명백하다. $B\otimes_A B$의 텐서곱 위상은 $B$의
위상으로부터 얻어지는 위상보다 섬세하지 않기 때문이다. 따라서
$B\otimes_A B$의 위상이
$\mathfrak J=\mathfrak J_{B/A}$에 유도하는 위상도 $B$의
위상으로부터 $\mathfrak J$에 얻어지는 위상보다 섬세하지 않다.

\oldpage[0\textsubscript{IV-1}]{125}
둘째 주장을 보이자. $B$의 두 $A$-부분가군 $M,N$에 대하여 표기의
남용으로 $\operatorname{Im}(M\otimes_A N)$을 $M\otimes N$이라
쓰겠다. $j_1$으로 $B$-가군 구조를 준 $B\otimes_A B$ 안에서의 관계
\[
\begin{aligned}
  (xy)\otimes z-x\otimes(yz)
  &=(x\otimes1)(1\otimes z)(y\otimes1-1\otimes y)\\
  &=xz\cdot(y\otimes1-1\otimes y)
    -x\cdot(z\otimes1-1\otimes z)(y\otimes1-1\otimes y)
\end{aligned}
\]
와 \sref{0.20.4.4-ko}를 고려하면, $B$의 아이디얼 $\mathfrak R$에
대하여 곧
\[
  ((\mathfrak R^2\otimes B)+(B\otimes\mathfrak R^2))
  \cap\mathfrak J
  \subset
  (\mathfrak R\otimes\mathfrak R)\cap\mathfrak J
  +\mathfrak R\mathfrak J+\mathfrak J^2
\]
를 얻는다. 또 $x_i,y_i\in\mathfrak R$이고
$\sum_i(x_i\otimes y_i)\in\mathfrak J$이면
\sref{0.20.4.4.1-ko}에서
$\sum_i(x_i\otimes y_i)\in\mathfrak R\mathfrak J$를 얻는다. 따라서
마침내
\begin{equation}
\tag{20.4.5.1}
\label{0.20.4.5.1-ko}
  (\mathfrak R^2\otimes B+B\otimes\mathfrak R^2)\cap\mathfrak J
  \subset\mathfrak R\mathfrak J+\mathfrak J^2.
\end{equation}
한편 $B\otimes_A B$의 위상이 $\mathfrak J$에 유도하는 위상에서,
$\mathfrak R$가 열린 아이디얼들을 달릴 때
\[
  (\mathfrak R\otimes B+B\otimes\mathfrak R)\cap\mathfrak J
\]
들은 $\mathfrak J$ 안의 $0$의 근방 기본계를 이룬다.
$B$의 위상으로부터 얻어지는 $\Omega^1_{B/A}$의 위상 역시
$B$의 위상으로부터 $\mathfrak J$에 얻어지는 위상을
$\mathfrak J^2$로 나눈 몫위상이다. 그러므로 $B$의 열린
아이디얼들에 대한 가정과 \sref{0.20.4.5.1-ko}가 둘째 주장의 증명을
완성한다.
\end{proof}

\begin{definition}[20.4.6]
\label{0.20.4.6-ko}
합성사상
\[
  B\xrightarrow{j_1}B\otimes_A B\longrightarrow\mathrm P^1_{B/A},
  \qquad
  B\xrightarrow{j_2}B\otimes_A B\longrightarrow\mathrm P^1_{B/A}
\]
을 각각 $p_1,p_2$라 하자. 이들은
$\varepsilon\circ p_1=\varepsilon\circ p_2=\mathrm I_B$를 만족하는
연속 $A$-준동형이다. $A$-가군의 연속 $A$-준동형
\begin{equation}
\tag{20.4.6.1}
\label{0.20.4.6.1-ko}
  d_{B/A}=p_1-p_2:B\to\Omega^1_{B/A}
\end{equation}
를 $B$의 $A$에 대한 \emph{외미분}이라 한다. 각 $x\in B$에 대하여
$d_{B/A}(x)$ ($d(x)$ 또는 $dx$로도 쓴다)를 $x$의 $A$에 대한
\emph{미분}이라 한다.
\end{definition}

$A=\mathbf Z$이면 $d_{B/\mathbf Z}$ 대신 $d_B$라 쓴다. $B$가 소체
$P$ 위의 대수이면 $d_B=d_{B/P}$이다.

\begin{proposition}[20.4.7]
\label{0.20.4.7-ko}
$A$-대수 $B$의 한 생성원 집합을 $x$가 달릴 때의 원소
$d_{B/A}(x)$들이 $B$-가군 $\Omega^1_{B/A}$를 생성한다.
\end{proposition}

\begin{proof}
$d_{B/A}(x)$는 $\Omega^1_{B/A}$에서
$x\otimes1-1\otimes x$의 표준상이다. 따라서 주장은
\sref{0.20.4.4-ko}의 즉각적인 귀결이다.
\end{proof}

\begin{theorem}[20.4.8]
\label{0.20.4.8-ko}
$A$를 위상환, $B$를 위상 $A$-대수라 하자.
\begin{enumerate}
  \item[\rm(i)] $\Omega^1_{B/A}$ 위에서 항등사상이 되는 증대 위상
  $B$-대수의 동형사상
\begin{equation}
\tag{20.4.8.1}
\label{0.20.4.8.1-ko}
  \varphi:\mathrm P^1_{B/A}\xrightarrow{\sim}
  D_B(\Omega^1_{B/A})
\end{equation}
  이 유일하게 존재한다.
  \item[\rm(ii)] 준동형 $d_{B/A}$는 $B$에서
  $\Omega^1_{B/A}$로 가는 $A$-미분이며 다음 보편 성질을 갖는다.
  모든 위상 $B$-가군 $L$에 대하여 $u\mapsto u\circ d_{B/A}$는
  $A$-가군의 동형사상
\begin{equation}
\tag{20.4.8.2}
\label{0.20.4.8.2-ko}
  \operatorname{Hom.\ cont}_B(\Omega^1_{B/A},L)
  \xrightarrow{\sim}
  \operatorname{D\acute{e}r.\ cont}_A(B,L)
\end{equation}
  이다.
\end{enumerate}
\end{theorem}

\begin{proof}
\oldpage[0\textsubscript{IV-1}]{126}
(i) \sref{0.20.4.6-ko}의 표기를 쓰면 $\varphi$는 반드시
\[
  z\longmapsto
  \bigl(\varepsilon(z),z-p_1(\varepsilon(z))\bigr)
\]
이고 그 역동형사상은 $(b,x)\mapsto p_1(b)+x$이다. 두 사상이
연속이므로 첫째 주장이 성립한다. 이로부터 $p_1$을 통해 $B$의
위상이 $B\otimes_A B$의 위상을 $\mathfrak J_{B/A}$로 나눈
\emph{몫위상}과 동일시됨도 알 수 있다.

(ii) $d_{B/A}$가 $B$에서 $\Omega^1_{B/A}$로 가는 $A$-미분이라는
사실은 정의 \sref{0.20.4.6-ko}와 \sref{0.20.1.1-ko}에서 나온다.

보편 성질을 증명하자.
$\operatorname{D\acute{e}r.\ cont}_A(B,L)$는 다음 연속
$A$-대수 준동형들의 집합 $G$와 표준적으로 동일시된다.
\[
  u:B\to D_B(L),\qquad q\circ u=\mathrm I_B,
\]
여기서 $q(b,x)=b$는 사영이다
(\sref{0.20.1.6-ko}, \sref{0.20.3.2-ko} 참조). 한편 동형사상
$\varphi$에 의하여
$\operatorname{Hom.\ cont}_B(\Omega^1_{B/A},L)$는 합성사상
\[
  \mathrm P^1_{B/A}\xrightarrow{v}D_B(L)\xrightarrow{q}B
\]
이 증대 $\varepsilon$인 연속 $B$-대수 준동형
$v:\mathrm P^1_{B/A}\to D_B(L)$들의 집합과 표준적으로 동일시된다.
정의상 $p_2=p_1-d_{B/A}$이므로 모든 $u\in G$가 다음과 같이
분해됨을 보이면 충분하다.
\[
\xymatrix{
  B \ar[r]^u \ar[d]_{p_2} & D_B(L) \\
  \mathrm P^1_{B/A} \ar[ur]_v &
}
\]
여기서 $v$는 연속 $B$-대수 준동형이어야 한다.

이미 $G$에 속하는 연속 $A$-대수 준동형
$j:B\to D_B(L)$, $b\mapsto(b,0)$가 있다. 위상 대수의 위상
텐서곱의 정의 ($\mathbf 0_{\mathrm I}$, 7.7.6)에 따라 다음 도식을
가환으로 만드는 연속 $A$-대수 준동형
$w:B\otimes_A B\to D_B(L)$가 존재한다.
\[
\xymatrix{
  B\otimes_A B \ar[dr]^w & B \ar[l]_{j_1} \ar[d]^j \\
  B \ar[u]^{j_2} \ar[r]_u & D_B(L)
}
\]
그러면 정의상
\[
  w(b\otimes1-1\otimes b)=j(b)-u(b)\in L.
\]
\sref{0.20.4.4-ko}에 따라 $w(\mathfrak J)\subset L$이고, $L$은
$D_B(L)$에서 제곱이 영인 아이디얼이므로 $w(\mathfrak J^2)=0$이다.
따라서 $w$는
\[
  B\otimes_A B\longrightarrow\mathrm P^1_{B/A}
  \xrightarrow{v}D_B(L)
\]
로 분해된다. 여기서 $v$는 연속 $A$-대수 준동형이다. 또한
$v\circ p_1=j$가 $B$-대수 준동형이므로
$\mathrm P^1_{B/A}$의 $B$-대수 구조의 정의에 따라 $v$도
$B$-대수 준동형이다. 정의상 $u=v\circ p_2$이므로 증명이 끝난다.
\end{proof}

\begin{theorem}[20.4.9]
\label{0.20.4.9-ko}
$B$가 형식적으로 매끄러운 위상 $A$-대수라고 가정하자. 그러면 위상
$B$-가군 $\Omega^1_{B/A}$는 \emph{형식적으로 사영}이다.
\end{theorem}

\begin{proof}
$j_1$이 정의하는 위상 $B$-대수 구조를 준 $B\otimes_A B$는 가정에
따라 형식적으로 매끄러운 위상 $B$-대수이다
\sref{0.19.3.5-ko}, (iii). 따라서 이는 형식적으로 매끄러운 위상
$A$-대수이기도 하다 \sref{0.19.3.5-ko}, (ii). 한편 $B$는
$A$-대수 몫
\[
  (B\otimes_A B)/\mathfrak J
\]
과 위상적으로 동형이고, 형식적으로 매끄러운 $A$-대수이다. 그러므로
결론은 \sref{0.19.5.3-ko}에서 나온다.
\end{proof}

\oldpage[0\textsubscript{IV-1}]{127}
\begin{corollary}[20.4.10]
\label{0.20.4.10-ko}
더 나아가 $B$의 모든 열린 아이디얼의 제곱이 열려 있다고 가정하자.
그러면 $B$의 모든 열린 아이디얼 $\mathfrak R$에 대하여
\[
  \Omega^1_{B/A}\otimes_B(B/\mathfrak R)
  =\Omega^1_{B/A}/\mathfrak R\Omega^1_{B/A}
\]
는 사영 $(B/\mathfrak R)$-가군이다.
\end{corollary}

\begin{proof}
이때 $\Omega^1_{B/A}$의 위상은 $B$의 위상으로부터 얻어진다
\sref{0.20.4.5-ko}. 따라서 \sref{0.19.2.4-ko}를 적용하면 된다.
\end{proof}

\begin{corollary}[20.4.11]
\label{0.20.4.11-ko}
$A,B$를 두 뇌터 국소환, $\rho:A\to B$를 $B$를 형식적으로
매끄러운 위상 $A$-대수로 만드는 국소 준동형이라 하자 (선아딕
위상에 대하여). 그러면 $B$의 모든 정의 아이디얼 $\mathfrak b$에
대하여
\[
  \Omega^1_{B/A}\otimes_B(B/\mathfrak b)
\]
는 자유 $(B/\mathfrak b)$-가군이다.
\end{corollary}

\begin{proof}
\sref{0.20.4.10-ko}에 따라 이 가군은 사영이다. $B/\mathfrak b$는
아르틴환이므로 모든 사영 $(B/\mathfrak b)$-가군은 자유이다
($\mathbf 0_{\mathrm{III}}$, 10.1.3).
\end{proof}

\begin{proposition}[20.4.12]
\label{0.20.4.12-ko}
$A$를 위상환, $B$를 위상 $A$-대수라 하자. 구조 준동형
$A\to B$가 전사이면
\[
  \Omega^1_{B/A}=0.
\]
\end{proposition}

\begin{proof}
\sref{0.20.1.1-ko}에 따라 모든 $B$-가군 $L$에 대하여
$\operatorname{D\acute{e}r}_A(B,L)=0$이다. 따라서 주장은
\sref{0.20.4.8-ko}에서 곧 나온다.
\end{proof}

\begin{examples}[20.4.13]
\label{0.20.4.13-ko}
\begin{enumerate}
  \item[\rm(i)] $A$를 환,
  $B=A[X_\alpha]_{\alpha\in I}$를 $A$ 위의 다항식 대수라 하자.
  그러면 $\Omega^1_{B/A}$는 $dX_\alpha$들을 기저로 하는 자유
  $B$-가군이다. 실제로 $dX_\alpha$들은 이 $B$-가군을 생성한다
  \sref{0.20.4.7-ko}. 한편 같은 지표집합 $I$를 갖는 기저
  $(z_\alpha)_{\alpha\in I}$를 가진 자유 $B$-가군 $L$을 택하면,
  모든 $\alpha\in I$에 대하여
  \[
    u(X_\alpha)=(X_\alpha,z_\alpha)
  \]
  인 $A$-대수 준동형 $u:B\to D_B(L)$가 존재한다. 따라서
  \sref{0.20.1.5-ko}에 의해 모든 $\alpha$에 대하여
  $D(X_\alpha)=z_\alpha$인 $A$-미분 $D:B\to L$이 존재한다.
  \sref{0.20.4.8.1-ko}에 따라 이는 $dX_\alpha$들이 선형 독립임을
  보인다.

  \item[\rm(ii)] $A$를 환, $L$을 $A$-가군으로 놓고
  $B=D_A(L)$을 제곱영 $A$-대수라 하자. 모든 $B$-가군 $M$에 대하여
  $L$로의 제한은 $\operatorname{D\acute{e}r}_A(B,M)$을 다음 조건을
  만족하는 $A$-선형사상 $u:L\to M$들과 동일시한다.
  \[
    (0,x)u(y)+(0,y)u(x)=0
    \qquad(x,y\in L).
  \]
  따라서 $\Omega^1_{B/A}$는 다음 표시를 갖는다.
  \[
    B\otimes_A(L\otimes_A L)\longrightarrow
    B\otimes_A L\longrightarrow
    \Omega^1_{B/A}\longrightarrow0.
  \]
  여기서 첫째 화살표는
  \[
    b\otimes(x\otimes y)\longmapsto
    b\bigl((0,x)\otimes y+(0,y)\otimes x\bigr)
  \]
  이고, 둘째 화살표는
  $b\otimes x\mapsto b\,d_{B/A}(0,x)$이다.

  \item[\rm(iii)] $B$가 두 위상 $A$-대수 $B_1,B_2$의 곱이라고
  가정하고 이들을 $B$의 아이디얼과 동일시하자. 준동형 $p$
  \sref{0.20.4.1.1-ko}에 의한
  $B_1\otimes_A B_2$와 $B_2\otimes_A B_1$의 상은 영이다. 따라서
  곧 $\mathrm P^1_{B/A}$는 위상적 \emph{곱}
  \[
    \mathrm P^1_{B_1/A}\times\mathrm P^1_{B_2/A}
  \]
  과 동일시되고, $B$-가군 $\Omega^1_{B/A}$는
  $\Omega^1_{B_1/A}$와 $\Omega^1_{B_2/A}$의 위상적
  \emph{직합}이다. 이 두 가군은 각각 $B_2$와 $B_1$에 의해
  소멸한다.

  \item[\rm(iv)] $A\subset B$인 두 정역에서 $B$가 $A$ 위에서
  정수적이고, $B$의 분수체가 $A$의 분수체 $K$의 분리 확대라고
  하자. 그러면 $B$-가군 $\Omega^1_{B/A}$는 \emph{꼬임 가군}이다.
  실제로 $x\in B$를 택하고 $f(T)$를 $x$의 $K$ 위 최소다항식이라
  하자. $af(T)\in A[T]$가 되도록 $0\neq a\in A$를 택할 수 있다.
  $x$가 $K$ 위에서 분리적이므로 $f'(x)\neq0$이고, 관계
  $af(x)=0$에서
  \[
    a f'(x)d_{B/A}x=0
  \]
  을 얻는다. 여기서 $a f'(x)\neq0$이다. 따라서
  \sref{0.20.4.7-ko}에 의해 주장이 성립한다.
\end{enumerate}
\end{examples}

\begin{remarks}[20.4.14]
\label{0.20.4.14-ko}
\begin{enumerate}
  \item[\rm(i)] 위상을 잊은 $B$-가군 $\Omega^1_{B/A}$는 $A$와
  $B$의 위상에 의존하지 않는다.
  \item[\rm(ii)] 뒤에서 $B$의 $A$에 대한 “$n$차 주요 부분의
  대수”
  \[
    \mathrm P^n_{B/A}
    =(B\otimes_A B)/(\mathfrak J_{B/A})^{n+1}
  \]
  를 도입하겠다. 이는 “$n$차 미분 계산”의 기초가 된다.
\end{enumerate}
\end{remarks}

\subsection{\texorpdfstring{$\Omega^1_{B/A}$}{Omega1(B/A)}의 기본
함자적 성질}
\label{subsection:0.20.5-ko}

\oldpage[0\textsubscript{IV-1}]{128}
\begin{env}[20.5.1]
\label{0.20.5.1-ko}
이 항과 다음 항에서는 명시적으로 달리 말하지 않는 한, 고려하는
모든 환과 가군에 \emph{이산} 위상을 주었다고 가정한다.
\end{env}

\begin{env}[20.5.2]
\label{0.20.5.2-ko}
$A$를 환, $B,C$를 두 $A$-대수,
$u:B\to C$를 $A$-준동형이라 하자. 가환도식
\[
\xymatrix{
  B\otimes_A B \ar[r]^{u\otimes u} \ar[d]_{p_{B/A}} &
  C\otimes_A C \ar[d]^{p_{C/A}} \\
  B \ar[r]^u & C
}
\tag{20.5.2.1}
\label{0.20.5.2.1-ko}
\]
에서 몫으로 내려가면 대수의 $A$-준동형
\begin{equation}
\tag{20.5.2.2}
\label{0.20.5.2.2-ko}
  u':\mathrm P^1_{B/A}\to\mathrm P^1_{C/A}
\end{equation}
을 얻으며, 다음 도식은 가환이다.
\[
\xymatrix{
  \mathrm P^1_{B/A} \ar[r]^{u'} & \mathrm P^1_{C/A} \\
  B \ar[u]^{p_1} \ar[r]^u & C \ar[u]_{p_1}
}
\]
$u\otimes u$가 $\mathfrak J_{B/A}$를 $\mathfrak J_{C/A}$로
보내므로 $u'$를 $\Omega^1_{B/A}$에 제한하여 사상
\begin{equation}
\tag{20.5.2.3}
\label{0.20.5.2.3-ko}
  u'':\Omega^1_{B/A}\to\Omega^1_{C/A}
\end{equation}
을 얻는다. $B$-가군 $\Omega^1_{B/A}$와 $C$-가군
$\Omega^1_{C/A}$에 대하여 쌍 $(u'',u)$는 \emph{디-준동형}이다.
따라서 표준적으로 $C$-가군 준동형
\begin{equation}
\tag{20.5.2.4}
\label{0.20.5.2.4-ko}
  u_{C/B/A}:\Omega^1_{B/A}\otimes_B C\to\Omega^1_{C/A}
\end{equation}
을 얻는다.

\oldpage[0\textsubscript{IV-1}]{129}
또한 도식
\[
\xymatrix{
  \mathrm P^1_{B/A} \ar[r]^{u'} & \mathrm P^1_{C/A} \\
  B \ar[u]^{p_2} \ar[r]^u & C \ar[u]_{p_2}
}
\tag{20.5.2.5}
\label{0.20.5.2.5-ko}
\]
도 가환이므로 다음 도식이 가환임을 얻는다.
\[
\xymatrix{
  \Omega^1_{B/A} \ar[r]^{u''} & \Omega^1_{C/A} \\
  B \ar[u]^{d_{B/A}} \ar[r]^u & C \ar[u]_{d_{C/A}}
}
\tag{20.5.2.6}
\label{0.20.5.2.6-ko}
\]
끝으로 $w:C\to D$가 둘째 $A$-대수 준동형이면 정의에서 곧 다음
추이성을 얻는다.
\begin{equation}
\tag{20.5.2.7}
\label{0.20.5.2.7-ko}
  (w\circ u)_{D/B/A}
  =w_{D/C/A}\circ
  (u_{C/B/A}\otimes\mathrm I_D).
\end{equation}
\end{env}

\begin{env}[20.5.3]
\label{0.20.5.3-ko}
이제 $A,B$를 두 환, $v:A\to B$를 환 준동형, $C$를 $v$를 통해
$A$-대수가 되는 $B$-대수라 하자. 표준 사상
\[
  v_0:C\otimes_A C\to C\otimes_B C
\]
은 $v:A\to B$에 관한 대수의 전사 \emph{디-준동형}이며 다음 도식은
가환이다.
\[
\xymatrix{
  C\otimes_A C \ar[r]^{v_0} \ar[d]_{p_{C/A}} &
  C\otimes_B C \ar[d]^{p_{C/B}} \\
  C \ar[r]^{\mathrm I_C} & C
}
\tag{20.5.3.1}
\label{0.20.5.3.1-ko}
\]
몫으로 내려가면 대수의 디-준동형
\begin{equation}
\tag{20.5.3.2}
\label{0.20.5.3.2-ko}
  v':\mathrm P^1_{C/A}\to\mathrm P^1_{C/B}
\end{equation}
을 얻으며, 다음 도식은 가환이다.
\[
\xymatrix{
  \mathrm P^1_{C/A} \ar[r]^{v'} & \mathrm P^1_{C/B} \\
  C \ar[u]^{p_1} \ar[r]^{\mathrm I_C} & C \ar[u]_{p_1}
}
\]
$v_0$가 $\mathfrak J_{C/A}$를 $\mathfrak J_{C/B}$로 보내므로
$v'$를 $\Omega^1_{C/A}$에 제한하여 $C$-가군 준동형
\begin{equation}
\tag{20.5.3.3}
\label{0.20.5.3.3-ko}
  v_{C/B/A}:\Omega^1_{C/A}\to\Omega^1_{C/B}
\end{equation}
을 얻는다.

\oldpage[0\textsubscript{IV-1}]{130}
또한 도식
\[
\xymatrix{
  \mathrm P^1_{C/A} \ar[r]^{v'} & \mathrm P^1_{C/B} \\
  C \ar[u]^{p_2} \ar[r]^{\mathrm I_C} & C \ar[u]_{p_2}
}
\tag{20.5.3.4}
\label{0.20.5.3.4-ko}
\]
도 가환이므로 다음 도식이 가환임을 얻는다.
\[
\xymatrix{
  \Omega^1_{C/A} \ar[r]^{v_{C/B/A}} & \Omega^1_{C/B} \\
  C \ar[u]^{d_{C/A}} \ar[r]^{\mathrm I_C} & C \ar[u]_{d_{C/B}}
}
\tag{20.5.3.5}
\label{0.20.5.3.5-ko}
\]
끝으로 $s:A'\to A$가 둘째 환 준동형이면 다음 추이성이 성립한다.
\begin{equation}
\tag{20.5.3.6}
\label{0.20.5.3.6-ko}
  (v\circ s)_{C/B/A'}
  =v_{C/B/A}\circ s_{C/A/A'}.
\end{equation}
\end{env}

\begin{env}[20.5.4]
\label{0.20.5.4-ko}
이제 환 준동형의 가환도식
\[
\xymatrix{
  B \ar[r]^v & B' \\
  A \ar[u]^f \ar[r]^u & A' \ar[u]_g
}
\]
이 주어졌다고 하자. \sref{0.20.5.2.4-ko}와
\sref{0.20.5.3.3-ko}를 합성하여 $B'$-가군 준동형
\begin{equation}
\tag{20.5.4.1}
\label{0.20.5.4.1-ko}
  \Omega^1_{B/A}\otimes_B B'
  \xrightarrow{v_{B'/B/A}}\Omega^1_{B'/A}
  \xrightarrow{u_{B'/A'/A}}\Omega^1_{B'/A'}
\end{equation}
을 얻는다. 이때 다음 $A$-준동형들의 도식은 가환이다.
\[
\xymatrix{
  \Omega^1_{B/A}\otimes_B B' \ar[r] & \Omega^1_{B'/A'} \\
  B \ar[u]^{x\mapsto d_{B/A}x\otimes1} \ar[r]^v &
  B' \ar[u]_{d_{B'/A'}}
}
\tag{20.5.4.2}
\label{0.20.5.4.2-ko}
\]
또한 준동형 \sref{0.20.5.4.1-ko}는 $B$-가군의 디-준동형
\begin{equation}
\tag{20.5.4.3}
\label{0.20.5.4.3-ko}
  \Omega^1_{B/A}\to\Omega^1_{B'/A'}
\end{equation}
에 대응한다.
\end{env}

\begin{proposition}[20.5.5]
\label{0.20.5.5-ko}
$A',B$가 두 $A$-대수이고 $B'=B\otimes_A A'$이면 표준 준동형
\sref{0.20.5.4.1-ko}
\begin{equation}
\tag{20.5.5.1}
\label{0.20.5.5.1-ko}
  \Omega^1_{B/A}\otimes_B B'\to\Omega^1_{B'/A'}
\end{equation}
은 전단사이다.
\end{proposition}

\begin{proof}
\oldpage[0\textsubscript{IV-1}]{131}
\sref{0.20.5.5.1-ko}의 왼쪽 항은
$\Omega^1_{B/A}\otimes_A A'$에 다름 아니다. 표준 동형을 통해
\[
  B'\otimes_{A'}B'=(B\otimes_A B)\otimes_A A'
\]
로 쓸 수 있고, 이때 $p_{B'/A'}$는
$p_{B/A}\otimes\mathrm I_{A'}$와 동일시된다. $p_{B/A}$가
전사이므로
\[
  \mathfrak J_{B'/A'}
  =\operatorname{Im}(\mathfrak J_{B/A}\otimes_A A'),
  \qquad
  \mathfrak J^2_{B'/A'}
  =\operatorname{Im}(\mathfrak J^2_{B/A}\otimes_A A')
\]
이다. 따라서 증대 아이디얼들을 서로 보내는 표준 동형을 통해
\[
  \mathrm P^1_{B'/A'}
  =\mathrm P^1_{B/A}\otimes_A A'
\]
이다. $A$-가군 $\mathrm P^1_{B/A}$는 표준적으로
$B\oplus\Omega^1_{B/A}$와 동일시되므로 같은 동형을 통해
\begin{equation}
\tag{20.5.5.2}
\label{0.20.5.5.2-ko}
  \Omega^1_{B/A}\otimes_A A'=\Omega^1_{B'/A'}
\end{equation}
을 얻는다. 이 동형과 표준 동형
\[
  \Omega^1_{B/A}\otimes_B B'
  \xrightarrow{\sim}\Omega^1_{B/A}\otimes_A A'
\]
의 합성이 바로 \sref{0.20.5.5.1-ko}임을 곧 확인할 수 있다.
\end{proof}

\begin{env}[20.5.6]
\label{0.20.5.6-ko}
표준 준동형 \sref{0.20.5.2.4-ko}와
\sref{0.20.5.3.3-ko}는 함자성에 따라 모든 $C$-가군 $L$에 대하여
표준 준동형
\begin{align}
\tag{20.5.6.1}
\label{0.20.5.6.1-ko}
  &\operatorname{Hom}_C(\Omega^1_{C/A},L)
  \to\operatorname{Hom}_C(\Omega^1_{B/A}\otimes_B C,L)
  =\operatorname{Hom}_B(\Omega^1_{B/A},L),\\
\tag{20.5.6.2}
\label{0.20.5.6.2-ko}
  &\operatorname{Hom}_C(\Omega^1_{C/B},L)
  \to\operatorname{Hom}_C(\Omega^1_{C/A},L)
\end{align}
을 준다. \sref{0.20.4.8.2-ko}와 가환도식
\sref{0.20.5.2.6-ko}, \sref{0.20.5.3.5-ko}를 고려하면 이
준동형들은 표준적 동일시 아래 각각
\sref{0.20.2.1.2-ko}와 \sref{0.20.2.1.3-ko}의 준동형에 다름
아니다.
\end{env}

\begin{theorem}[20.5.7]
\label{0.20.5.7-ko}
$u:A\to B$, $v:B\to C$를 두 환 준동형이라 하자.
\begin{enumerate}
  \item[\rm(i)] $C$-가군의 열
\begin{equation}
\tag{20.5.7.1}
\label{0.20.5.7.1-ko}
  \Omega^1_{B/A}\otimes_B C
  \xrightarrow{v_{C/B/A}}\Omega^1_{C/A}
  \xrightarrow{u_{C/B/A}}\Omega^1_{C/B}\to0
\end{equation}
은 완전하다.
  \item[\rm(ii)] $v_{C/B/A}$가 왼쪽 가역이기 위한 필요충분조건은
  $C$가 $A$에 대하여 상대적으로 형식적으로 매끄러운
  $B$-대수인 것이다(이산 위상에 대하여
  (cf. \sref{0.19.9.1-ko})). 특히 $C$가 형식적으로 매끄러운
  $B$-대수이면 충분하다(이산 위상에 대하여).
\end{enumerate}
\end{theorem}

\begin{proof}
(i) \sref{0.20.2.4.2-ko}의 열이 완전하다는 사실과
\sref{0.20.5.6-ko}에 의하여 모든 $C$-가군 $L$에 대하여 열
\[
  0\to\operatorname{Hom}_C(\Omega^1_{C/B},L)
  \to\operatorname{Hom}_C(\Omega^1_{C/A},L)
  \to\operatorname{Hom}_C(\Omega^1_{B/A}\otimes_B C,L)
\]
이 완전하다. 이로부터 열 \sref{0.20.5.7.1-ko}이 완전함이 따른다는
것은 알려져 있다 (Bourbaki, \emph{Alg.}, chap.~II,
3\textsuperscript{e} éd., §~2, n\textsuperscript{o}~1, th.~1).

(ii) \sref{0.20.5.7.1-ko}이 완전하므로 $v_{C/B/A}$가 왼쪽
가역이라는 것은 열
\begin{equation}
\tag{20.5.7.2}
\label{0.20.5.7.2-ko}
  0\to\Omega^1_{B/A}\otimes_B C
  \xrightarrow{v_{C/B/A}}\Omega^1_{C/A}
  \xrightarrow{u_{C/B/A}}\Omega^1_{C/B}\to0
\end{equation}
이 \emph{완전하고 분할된다}는 뜻이다. 이것은 모든 $C$-가군
$L$에 대하여 열
\[
  0\to\operatorname{Hom}_C(\Omega^1_{C/B},L)
  \to\operatorname{Hom}_C(\Omega^1_{C/A},L)
  \to\operatorname{Hom}_C(\Omega^1_{B/A}\otimes_B C,L)\to0
\]
이 완전하다는 것과 동치임이 알려져 있다
(Bourbaki, \emph{loc. cit.}, n\textsuperscript{o}~1, prop.~1).
\sref{0.20.5.6-ko}과 \sref{0.20.2.4.2-ko}에 의하면 이 조건은
모든 $C$-가군 $L$에 대하여
\[
  \operatorname{Exalcom}_{B/A}(C,L)=0
\]
이라는 것과 동치이다. 따라서 결론은
\sref{0.19.9.8.1-ko}에서 따른다.
\end{proof}

\oldpage[0\textsubscript{IV-1}]{132}
또한 환 준동형의 가환도식
\[
\xymatrix{
  A' \ar[r] & B' \ar[r] & C' \\
  A \ar[u] \ar[r] & B \ar[u] \ar[r] & C \ar[u]
}
\]
이 주어지면 가환도식
\[
\xymatrix{
  \Omega^1_{B/A}\otimes_B C \ar[r] \ar[d] &
  \Omega^1_{C/A} \ar[r] \ar[d] &
  \Omega^1_{C/B} \ar[r] \ar[d] & 0 \\
  \Omega^1_{B'/A'}\otimes_{B'} C' \ar[r] &
  \Omega^1_{C'/A'} \ar[r] &
  \Omega^1_{C'/B'} \ar[r] & 0
}
\tag{20.5.7.3}
\label{0.20.5.7.3-ko}
\]
을 얻는다. 여기서 세로 화살표들은
\sref{0.20.5.4.3-ko}의 디-준동형들로부터 유도된다.

\begin{corollary}[20.5.8]
\label{0.20.5.8-ko}
준동형 $v:B\to C$에 의해 $C$가 형식적으로 에탈인 $B$-대수라고
하자(이산 위상에 대하여 \sref{0.19.10.2-ko}). 그러면 준동형
\sref{0.20.5.2.4-ko}
\[
  v_{C/B/A}:\Omega^1_{B/A}\otimes_B C\to\Omega^1_{C/A}
\]
은 전단사이다.
\end{corollary}

\begin{proof}
실제로 $C$가 이산 위상에 대하여 형식적으로 비분기인
$B$-대수이면 \sref{0.19.10.4-ko}, \sref{0.20.4.8-ko},
\sref{0.20.1.1-ko}에 의하여 모든 $C$-가군 $L$에 대하여
\[
  \operatorname{Hom}_C(\Omega^1_{C/B},L)=0
\]
이므로 $\Omega^1_{C/B}=0$이다
(cf. \sref{0.20.7.4-ko}). 한편 $C$가 이산 위상에 대하여
형식적으로 매끄러운 $B$-대수이면 열
\sref{0.20.5.7.2-ko}은 완전하다. 따라서 따름정리를 얻는다.
\end{proof}

\begin{corollary}[20.5.9]
\label{0.20.5.9-ko}
$A$를 환, $B$를 $A$-대수, $S$를 $B$의 곱셈적 부분집합이라
하자. 그러면 표준 준동형
\begin{equation}
\tag{20.5.9.1}
\label{0.20.5.9.1-ko}
  S^{-1}\Omega^1_{B/A}\to\Omega^1_{S^{-1}B/A}
\end{equation}
은 전단사이다.
\end{corollary}

\begin{proof}
$C=S^{-1}B$에 \sref{0.20.5.8-ko}을 적용하면 충분하다. 실제로
$C$는 이산 위상에 대하여 형식적으로 에탈인 $B$-대수이다
\sref{0.19.10.3-ko}, (ii).
\end{proof}

\sref{0.20.5.5-ko}을 고려하면 표준 동형들을 제외하고
\begin{equation}
\tag{20.5.9.2}
\label{0.20.5.9.2-ko}
  \Omega^1_{S^{-1}B/S^{-1}A}
  =S^{-1}\Omega^1_{B/A}
  =\Omega^1_{S^{-1}B/A}
\end{equation}
로 쓸 수 있다.

\begin{corollary}[20.5.10]
\label{0.20.5.10-ko}
$k$가 체이고 $K=k(X_\alpha)_{\alpha\in I}$가 $k$의 순수 초월
확대이면, $dX_\alpha$들은 $K$-벡터공간
$\Omega^1_{K/k}$의 기저를 이룬다.
\end{corollary}

\begin{proof}
$K$는 다항식환 $k[X_\alpha]_{\alpha\in I}$의 분수체이므로
\sref{0.20.4.13-ko}, (i)과 \sref{0.20.5.9-ko}에서 결론이 따른다.
\end{proof}

\oldpage[0\textsubscript{IV-1}]{133}
\begin{env}[20.5.11]
\label{0.20.5.11-ko}
$A$를 환, $B$를 $A$-대수, $\mathfrak R$을 $B$의 아이디얼,
$C=B/\mathfrak R$을 몫 $A$-대수라 하고, $A$-가군의 합성 준동형
\begin{equation}
\tag{20.5.11.1}
\label{0.20.5.11.1-ko}
  \delta:\mathfrak R\longrightarrow
  B\xrightarrow{d_{B/A}}\Omega^1_{B/A}
\end{equation}
을 생각하자. 첫째 화살표는 표준 포함사상이다.
$d(xy)=xdy+ydx$이므로
\[
  \delta(\mathfrak R^2)\subset
  \mathfrak R.\Omega^1_{B/A}
\]
이다. 따라서 몫을 취하면 $A$-가군 준동형
\begin{equation}
\tag{20.5.11.2}
\label{0.20.5.11.2-ko}
  \delta_{C/B/A}:\mathfrak R/\mathfrak R^2\longrightarrow
  \Omega^1_{B/A}\otimes_B C
  =\Omega^1_{B/A}/(\mathfrak R.\Omega^1_{B/A})
\end{equation}
을 얻는다. 실제로 $\delta_{C/B/A}$는 $C$-\emph{가군} 준동형이다.
왜냐하면 $x\in B$, $y\in\mathfrak R$에 대하여
$ydx\in\mathfrak R.\Omega^1_{B/A}$이므로
\[
  d(xy)\equiv xdy
  \pmod{\mathfrak R.\Omega^1_{B/A}}
\]
이기 때문이다. 이는 먼저 \sref{0.20.5.11.2-ko}가 $B$-가군
준동형임을 보이고, $\mathfrak R$이 양쪽 가군을 소멸시키므로
주장이 따른다.

$B'$을 또 하나의 $A$-대수, $u:B\to B'$을 $A$-준동형,
$\mathfrak R'$을 $u(\mathfrak R)\subset\mathfrak R'$을 만족하는
$B'$의 아이디얼, $C'=B'/\mathfrak R'$을 몫대수라 하자. 그러면
가환도식
\begin{equation}
\tag{20.5.11.3}
\label{0.20.5.11.3-ko}
\xymatrix{
  \mathfrak R/\mathfrak R^2
  \ar[r]^-{\delta_{C/B/A}} \ar[d] &
  \Omega^1_{B/A}\otimes_B C \ar[d] \\
  \mathfrak R'/\mathfrak R'^2
  \ar[r]_-{\delta_{C'/B'/A}} &
  \Omega^1_{B'/A}\otimes_{B'} C'
}
\end{equation}
을 얻는다. 여기서 세로 화살표들은 $u$로부터 유도된다
\sref{0.20.5.2.4-ko}.
\end{env}

\begin{theorem}[20.5.12]
\label{0.20.5.12-ko}
$A$를 환, $B$를 $A$-대수, $C=B/\mathfrak R$을 몫
$A$-대수, $u:B\to C$를 표준 준동형이라 하자.
\begin{enumerate}
  \item[\rm(i)] $C$-가군의 완전열
\begin{equation}
\tag{20.5.12.1}
\label{0.20.5.12.1-ko}
  \mathfrak R/\mathfrak R^2
  \xrightarrow{\delta_{C/B/A}}\Omega^1_{B/A}\otimes_B C
  \xrightarrow{u_{C/B/A}}\Omega^1_{C/A}\longrightarrow0
\end{equation}
이 존재한다. 여기서 $u_{C/B/A}$와 $\delta_{C/B/A}$는 각각
\sref{0.20.5.2.4-ko}와 \sref{0.20.5.11.2-ko}로 정의된다.
  \item[\rm(ii)] $E=B/\mathfrak R^2$로 놓으면 표준 준동형
  \sref{0.20.5.2.4-ko}
\[
  \Omega^1_{B/A}\otimes_B C\longrightarrow
  \Omega^1_{E/A}\otimes_E C
\]
은 전단사이다.
  \item[\rm(iii)] 다음 세 조건은 동치이다.
  \begin{enumerate}
    \item[\rm a)] $\delta_{C/B/A}$는 왼쪽 가역이다.
    \item[\rm b)] $C$-가군 $L$에 의한 $C$의 모든 $A$-확장 중
    $u:B\to C$에 의한 역상이 $A$-자명한 것은 그 자체도
    $A$-자명하다.
    \item[\rm c)] $A$-대수 $E=B/\mathfrak R^2$은
    $\mathfrak R/\mathfrak R^2$에 의한 $C$의 $A$-자명한 확장이다.
  \end{enumerate}
  \item[\rm(iv)] $\delta_{C/B/A}$의 왼쪽 역들과 표준 준동형
  $E\to C$의 오른쪽 역들 사이에는 표준 전단사 대응이 존재한다.
\end{enumerate}
\end{theorem}

\begin{proof}
\oldpage[0\textsubscript{IV-1}]{134}
(i) $u$가 전사이므로 \sref{0.20.1.1-ko}에 의하여 모든
$C$-가군 $L$에 대하여
$\operatorname{D\acute{e}r}_B(C,L)=0$이다. 따라서 완전열
\sref{0.20.2.3.1-ko}은
\begin{equation}
\tag{20.5.12.2}
\label{0.20.5.12.2-ko}
  0\longrightarrow\operatorname{D\acute{e}r}_A(C,L)
  \xrightarrow{u^0}\operatorname{D\acute{e}r}_A(B,L)
  \xrightarrow{\partial}\operatorname{Exalcom}_B(C,L)
  \xrightarrow{v^1}\operatorname{Exalcom}_A(C,L)
  \xrightarrow{u^1}\operatorname{Exalcom}_A(B,L)
\end{equation}
이 된다. 여기서 $v$는 준동형 $A\to B$이다. 한편
\sref{0.18.3.8-ko}에 따라
$\operatorname{Exalcom}_B(C,L)$은
$\operatorname{Hom}_C(\mathfrak R/\mathfrak R^2,L)$과 표준적으로
동일시된다. 따라서 \sref{0.20.5.12.2-ko}와
\sref{0.20.4.8-ko}로부터 완전열
\begin{equation}
\tag{20.5.12.3}
\label{0.20.5.12.3-ko}
\begin{split}
  0\longrightarrow\operatorname{Hom}_C(\Omega^1_{C/A},L)
  \longrightarrow\operatorname{Hom}_C
    (\Omega^1_{B/A}\otimes_B C,L)
  \xrightarrow{\varphi}\operatorname{Hom}_C
    (\mathfrak R/\mathfrak R^2,L)\\
  \xrightarrow{\psi}\operatorname{Ker}(u^1)\longrightarrow0
\end{split}
\end{equation}
을 얻는다. 여기서
$\varphi=\partial\circ\eta^{-1}$,
$\psi=v^1\circ\eta$이다. $\partial$의 정의
\sref{0.20.2.2-ko}와 $\eta$의 정의
\sref{0.18.3.8-ko}로 돌아가면 $\varphi$는 바로 준동형
\[
  \operatorname{Hom}(\delta_{C/B/A},1_L)
\]
임을 곧 확인한다. 따라서 \sref{0.20.5.12.3-ko}의 처음 네 항으로
이루어진 완전열이 존재하므로 열 \sref{0.20.5.12.1-ko}은 완전하다
(Bourbaki, \emph{Alg.}, chap.~II, 3\textsuperscript{e} éd.,
\S~2, n\textsuperscript{o}~1, th.~1).

\medskip
(ii) $B$와 아이디얼 $\mathfrak R^2$에 완전열
\sref{0.20.5.12.1-ko}을 적용하면
\begin{equation}
\tag{20.5.12.4}
\label{0.20.5.12.4-ko}
  \mathfrak R^2/\mathfrak R^4\longrightarrow
  \Omega^1_{B/A}\otimes_B E\longrightarrow
  \Omega^1_{E/A}\longrightarrow0
\end{equation}
을 얻는다. 여기에 $E$-대수로 보는 $C$를 텐서하여 완전열
\[
  \mathfrak R^2/\mathfrak R^4\otimes_E C\longrightarrow
  \Omega^1_{B/A}\otimes_B C\longrightarrow
  \Omega^1_{E/A}\otimes_E C\longrightarrow0
\]
을 얻는다. 이제 $x,y\in\mathfrak R$이고 $\zeta$가
$xy$의 $\mathfrak R^4$을 법으로 한 류이면, 정의에 따라
\[
  \delta'(\zeta\otimes1)
  =d_{B/A}(xy)\otimes1
  =(x\,d_{B/A}(y)+y\,d_{B/A}(x))\otimes1.
\]
$C$에서 $x,y$의 상은 영이므로
$\Omega^1_{B/A}\otimes_B C$에서
$d_{B/A}(xy)\otimes1=0$이다. 이로써 주장이 증명된다.

\medskip
(iii) \sref{0.20.5.12.1-ko}의 완전성에 의하여
$\delta_{C/B/A}$가 왼쪽 가역이라는 것은 열
\begin{equation}
\tag{20.5.12.5}
\label{0.20.5.12.5-ko}
  0\longrightarrow\mathfrak R/\mathfrak R^2
  \xrightarrow{\delta_{C/B/A}}\Omega^1_{B/A}\otimes_B C
  \xrightarrow{u_{C/B/A}}\Omega^1_{C/A}\longrightarrow0
\end{equation}
이 \emph{완전하고 분할된다}는 것과 동치이다. 이는 다시 모든
$L$에 대하여 완전열 \sref{0.20.5.12.3-ko}에서
$\operatorname{Ker}(u^1)=0$이라고 하는 것과 동치이다
(Bourbaki, \emph{loc. cit.}). 따라서 조건 a)와 b)는 동치이다
(cf. \sref{0.18.3.6.2-ko}).

b)가 c)를 함의하는 것은 다음에서 따른다. 즉
$\mathfrak R/\mathfrak R^2$에 의한 $C$의 $A$-확장 $E$를
$u:B\to C$에 따라 역상하면 $A$-자명하다. 실제로 $u$는 표준
준동형들 $B\to E\to C$의 합성이다
\sref{0.18.1.6-ko}. 역으로 c)는 b)를 함의한다. 모든
$L$에 의한 $C$의 $B$-확장은
\[
  E\oplus_{\mathfrak R/\mathfrak R^2}L
\]
과 $B$-동치이기 때문이다 \sref{0.18.3.8-ko}. 달리 말하면 그
확장의 류는 $B$-확장으로 본 $E$의 류를 $B$-준동형
$w:\mathfrak R/\mathfrak R^2\to L$에 대응하는 준동형
\[
\oldpage[0\textsubscript{IV-1}]{135}
  w_*:\operatorname{Exan}_B
    (C,\mathfrak R/\mathfrak R^2)
  \longrightarrow\operatorname{Exan}_B(C,L)
\]
로 보낸 상이다. 이제 도식 \sref{0.18.3.6.5-ko}
\[
\xymatrix{
  \operatorname{Exan}_B(C,\mathfrak R/\mathfrak R^2)
    \ar[r]^{w_*} \ar[d]_{u^*} &
  \operatorname{Exan}_B(C,L) \ar[d]^{u^*} \\
  \operatorname{Exan}_A(C,\mathfrak R/\mathfrak R^2)
    \ar[r]_{w_*} &
  \operatorname{Exan}_A(C,L)
}
\]
이 가환하므로 c)가 b)를 함의한다.

\medskip
(iv) \sref{0.20.1.7-ko}에서 $E\to C$의 오른쪽 역들은
\[
  D\in\operatorname{D\acute{e}r}_A
  (E,\mathfrak R/\mathfrak R^2)
\]
중 $\mathfrak R/\mathfrak R^2$의 모든 $x$에 대하여
$D(x)=x$를 만족하는 것들의 집합과 표준적으로 전단사 대응함을
보았다. 따라서 \sref{0.20.4.8-ko}에 의하여 이는 합성
\[
  \mathfrak R/\mathfrak R^2\xrightarrow{d_{E/A}}
  \Omega^1_{E/A}\xrightarrow{h}
  \mathfrak R/\mathfrak R^2
\]
이 항등사상이 되게 하는 $E$-준동형
\[
  h:\Omega^1_{E/A}\to\mathfrak R/\mathfrak R^2
\]
들의 집합과도 전단사 대응한다. $\mathfrak R/\mathfrak R^2$이
$C$-가군이므로 $C$를 텐서하면 합성
\[
  \mathfrak R/\mathfrak R^2
  \xrightarrow{\delta_{C/E/A}}
  \Omega^1_{E/A}\otimes_E C
  \xrightarrow{h\otimes1}
  \mathfrak R/\mathfrak R^2
\]
이 항등사상임을 얻는다. 또한
\[
  h\longmapsto h\otimes1
\]
은 집합
\[
  \operatorname{Hom}_E
  (\Omega^1_{E/A},\mathfrak R/\mathfrak R^2)
\]
에서
\[
  \operatorname{Hom}_C
  (\Omega^1_{E/A}\otimes_E C,\mathfrak R/\mathfrak R^2)
\]
로 가는 동형이다. 한편 (ii)에 의하여
$\Omega^1_{E/A}\otimes_E C$와
$\Omega^1_{B/A}\otimes_B C$를 표준적으로 동일시할 수 있고, 이때
$\delta_{C/E/A}$는 $\delta_{C/B/A}$와 동일시된다.
\end{proof}

\begin{example}[20.5.13]
\label{0.20.5.13-ko}
$B=A[X_\alpha]_{\alpha\in I}$를 $A$ 위의 다항식 대수,
$\mathfrak R$을 $B$의 아이디얼, $(P_\lambda)$를
$\mathfrak R$의 생성계, $C=B/\mathfrak R$이라 하자.
$\Omega^1_{B/A}$는 $dX_\alpha$들을 기저로 갖는 자유
$B$-가군이므로 \sref{0.20.4.13-ko}, (i), $dX_\alpha$들은 자유
$C$-가군 $\Omega^1_{B/A}\otimes_B C$의 기저이기도 하다. 한편
정의에서 곧바로, $\delta_{C/B/A}$에 의한
$\mathfrak R/\mathfrak R^2$의 상은
\[
  dP_\lambda=\sum_\alpha
  \frac{\partial P_\lambda}{\partial X_\alpha}\,dX_\alpha
\]
들로 생성되는 $C$-부분가군임을 알 수 있다. 따라서
$\Omega^1_{C/A}$는 $dX_\alpha$들을 기저로 갖는 자유
$C$-가군을 $dP_\lambda$들로 생성되는 $C$-부분가군으로 나눈
몫과 동형이다. 모든 $A$-대수 $C$는 위와 같은 방법으로 얻을 수
있으므로, 이는 \emph{임의의} 대수의 미분 가군을 기술한다.
\end{example}

\begin{corollary}[20.5.14]
\label{0.20.5.14-ko}
$C$가 형식적으로 매끄러운 $A$-대수이면(이산 위상에 대하여) 열
\begin{equation}
\tag{20.5.14.1}
\label{0.20.5.14.1-ko}
  0\longrightarrow\mathfrak R/\mathfrak R^2
  \xrightarrow{\delta_{C/B/A}}\Omega^1_{B/A}\otimes_B C
  \xrightarrow{u_{C/B/A}}\Omega^1_{C/A}\longrightarrow0
\end{equation}
은 완전하고 분할된다.
\end{corollary}

\begin{proof}
실제로 이때 $C$-가군에 의한 $C$의 모든 $A$-확장은 자명하다
\sref{0.19.4.4.1-ko}.
\end{proof}

\oldpage[0\textsubscript{IV-1}]{136}
\begin{remark}[20.5.15]
\label{0.20.5.15-ko}
$u:A\to B$를 환의 \emph{전사} 준동형이라 하자. 그러면 모든 환
준동형 $v:B\to C$에 대하여 표준 준동형
\begin{equation}
\tag{20.5.15.1}
\label{0.20.5.15.1-ko}
  u_{C/B/A}:\Omega^1_{C/A}\longrightarrow\Omega^1_{C/B}
\end{equation}
은 \emph{전단사}이다. 실제로
$\Omega^1_{B/A}=0$이므로 \sref{0.20.4.12-ko}, 이는 완전열
\sref{0.20.5.7.1-ko}에서 따른다.
\end{remark}

\subsection{불완전성 가군과 특성 준동형}
\label{subsection:0.20.6-ko}
\label{0.20.6-ko}

\begin{definition}[20.6.1]
\label{0.20.6.1-ko}
두 환 준동형 $u:A\to B$, $v:B\to C$가 주어졌다고 하자.
준동형
\[
  v_{C/B/A}:\Omega^1_{B/A}\otimes_B C\to\Omega^1_{C/A}
\]
의 핵인 $C$-가군을 $A$에 대한 $B$-대수 $C$의
\emph{불완전성 가군}이라 하고 $\Upsilon_{C/B/A}$로 나타낸다.
따라서 정의에 의해 (cf. \sref{0.20.5.7-ko}) 완전열
\begin{equation}
\tag{20.6.1.1}
\label{0.20.6.1.1-ko}
  0\longrightarrow\Upsilon_{C/B/A}\longrightarrow
  \Omega^1_{B/A}\otimes_B C\xrightarrow{v_{C/B/A}}
  \Omega^1_{C/A}\xrightarrow{u_{C/B/A}}\Omega^1_{C/B}
  \longrightarrow0
\end{equation}
을 얻는다.
\end{definition}

$A=\mathbb Z$이면 $\Omega^1_{B/A}$와 $\Omega^1_{C/A}$는 각각
“절대” 미분 가군 $\Omega^1_B$, $\Omega^1_C$이므로
$\Upsilon_{C/B/\mathbb Z}$ 대신 $\Upsilon_{C/B}$로 쓴다.
$B$와 $C$가 소체 $P$ 위의 대수이면
\[
  \Upsilon_{C/B/P}=\Upsilon_{C/B}
\]
이다.

$R,S$를 각각 $B,C$의 곱셈적 부분집합이라 하고, $R$의 상이
$S$에 포함된다고 하자. 그러면 완전열
\sref{0.20.6.1.1-ko}과 \sref{0.20.5.9-ko}로부터
\begin{equation}
\tag{20.6.1.2}
\label{0.20.6.1.2-ko}
  \Upsilon_{S^{-1}C/R^{-1}B/A}
  =S^{-1}\Upsilon_{C/B/A}
\end{equation}
를 얻는다.

\begin{proposition}[20.6.2]
\label{0.20.6.2-ko}
$C$가 $A$에 대하여 상대적으로 형식적으로 매끄러운
$B$-대수이면(특히 $C$가 형식적으로 매끄러운 $B$-대수이면)
\[
  \Upsilon_{C/B/A}=0
\]
이다.
\end{proposition}

\begin{proof}
이는 \sref{0.20.5.7-ko}, (ii)에서 따른다.
\end{proof}

\begin{proposition}[20.6.3]
\label{0.20.6.3-ko}
$k$를 체, $K$를 $k$의 확대체라 하자. $K$가 $k$의 분리
확대체이기 위한 필요충분조건은
\[
  \Upsilon_{K/k}=0
\]
인 것이다. 달리 말하면 표준 준동형
\[
  \Omega^1_k\otimes_k K\to\Omega^1_K
\]
이 단사이거나, 또 달리 말하면 \sref{0.20.4.8-ko}에 의하여
$k$에서 $K$로 가는 모든 미분이 $K$에서 그 자신으로 가는
미분으로 연장된다.
\end{proposition}

\begin{proof}
$K$가 $k$ 위에서 분리인 것과 형식적으로 매끄러운
$k$-대수인 것은 동치이다 \sref{0.19.6.1-ko}. 한편 $P$를
$k$의 소체라 하면 $K$는 $P$ 위에서 분리이므로 형식적으로
매끄러운 $P$-대수이다. 따라서 $K$가 형식적으로 매끄러운
$k$-대수인 것과 $P$에 대하여 상대적으로 형식적으로 매끄러운
$k$-대수인 것은 동치이다. 끝으로 \sref{0.20.5.7-ko}, (ii)에
의하여 이 마지막 조건은 준동형
\[
  v_{K/k/P}:\Omega^1_k\otimes_k K\to\Omega^1_K
\]
가 왼쪽 가역이라는 것과 동치이다. $K$가 체이므로 이는
$v_{K/k/P}$의 핵, 곧 $\Upsilon_{K/k}$가 영이라는 것과 동치이다.
\end{proof}

\begin{env}[20.6.4]
\label{0.20.6.4-ko}
가환환 준동형의 가환도식
\begin{equation}
\tag{20.6.4.1}
\label{0.20.6.4.1-ko}
\xymatrix{
  A' \ar[r]^{u'} & B' \ar[r]^{v'} & C' \\
  A \ar[u]^f \ar[r]_u & B \ar[u]^g \ar[r]_v & C \ar[u]^h
}
\end{equation}
을 생각하자.

\oldpage[0\textsubscript{IV-1}]{137}
이에 대응하는 도식 \sref{0.20.5.7.3-ko}의 가환성으로부터
\sref{0.20.6.4.1-ko}에 의해 표준적으로 유도되는 유일한
$C$-준동형
\begin{equation}
\tag{20.6.4.2}
\label{0.20.6.4.2-ko}
  \Upsilon_{C/B/A}\longrightarrow\Upsilon_{C'/B'/A'}
\end{equation}
을 얻으며, 이는 도식
\begin{equation}
\tag{20.6.4.3}
\label{0.20.6.4.3-ko}
\xymatrix{
  0 \ar[r] & \Upsilon_{C/B/A} \ar[r] \ar[d] &
  \Omega^1_{B/A}\otimes_B C
    \ar[r]^{v_{C/B/A}} \ar[d] &
  \Omega^1_{C/A} \ar[r]^{u_{C/B/A}} \ar[d] &
  \Omega^1_{C/B} \ar[r] \ar[d] & 0 \\
  0 \ar[r] & \Upsilon_{C'/B'/A'} \ar[r] &
  \Omega^1_{B'/A'}\otimes_{B'} C'
    \ar[r]_{v_{C'/B'/A'}} &
  \Omega^1_{C'/A'} \ar[r]_{u_{C'/B'/A'}} &
  \Omega^1_{C'/B'} \ar[r] & 0
}
\end{equation}
을 가환하게 한다.

준동형 \sref{0.20.6.4.2-ko}을 주는 것은, 표준 준동형
\[
  \Upsilon_{C/B/A}\to\Upsilon_{C/B/A}\otimes_C C'
\]
과 합성하면 \sref{0.20.6.4.2-ko}이 되는 $C'$-준동형
\begin{equation}
\tag{20.6.4.4}
\label{0.20.6.4.4-ko}
  \Upsilon_{C/B/A}\otimes_C C'
  \longrightarrow\Upsilon_{C'/B'/A'}
\end{equation}
을 주는 것과 동치이다. \sref{0.20.6.4.2-ko}은 명백한 추이성을
가지므로, $\Upsilon_{C/B/A}$는 삼중항 $(A,B,C)$에 관한
\emph{함자}라고 말할 수 있다.
\end{env}

\begin{env}[20.6.5]
\label{0.20.6.5-ko}
이후에는 \sref{0.20.6.1-ko}의 조건 아래에서, 차수 $0$과 $1$을
제외하면 항들이 영이고 그 두 차수에서는
\begin{equation}
\tag{20.6.5.1}
\label{0.20.6.5.1-ko}
\left\{
\begin{aligned}
  \mathrm K_0(C/B/A)&=\Omega^1_{C/A},\\
  \mathrm K_1(C/B/A)&=\Omega^1_{B/A}\otimes_B C
\end{aligned}
\right.
\end{equation}
인 $C$-가군의 \emph{사슬 복합체}
$\mathrm K_\bullet(C/B/A)$를 도입하는 것이 편리하다. 미분 연산자
$\mathrm K_1\to\mathrm K_0$는 $v_{C/B/A}$이다. 이로써 표준
동형들을 제외하고 $\Omega^1_{C/B}$와 $\Upsilon_{C/B/A}$를 이
복합체의 \emph{호몰로지 가군}으로
\begin{equation}
\tag{20.6.5.2}
\label{0.20.6.5.2-ko}
  \mathrm H_0(\mathrm K_\bullet(C/B/A))=\Omega^1_{C/B},
  \qquad
  \mathrm H_1(\mathrm K_\bullet(C/B/A))=\Upsilon_{C/B/A}
\end{equation}
와 같이 쓸 수 있다.
\end{env}

마찬가지로 다음이 성립한다.
\begin{proposition}[20.6.6]
\label{0.20.6.6-ko}
\sref{0.20.6.1-ko}의 가정 아래 모든 $C$-가군 $L$에 대하여
표준 $C$-동형
\begin{align}
\tag{20.6.6.1}
\label{0.20.6.6.1-ko}
  \mathrm H^0(\mathrm K_\bullet(C/B/A),L)
  &\simeq\operatorname{D\acute{e}r}_B(C,L),\\
\tag{20.6.6.2}
\label{0.20.6.6.2-ko}
  \mathrm H^1(\mathrm K_\bullet(C/B/A),L)
  &\simeq\operatorname{Exalcom}_{B/A}(C,L)
\end{align}
이 존재한다.
\end{proposition}

\begin{proof}
\sref{0.20.4.8-ko}과 \sref{0.20.5.6-ko}에 의하여 공사슬 복합체
\[
  \operatorname{Hom}_C(\mathrm K_\bullet(C/B/A),L)
\]
는 복합체
\[
  \cdots\longrightarrow0\longrightarrow
  \operatorname{D\acute{e}r}_A(C,L)
  \longrightarrow\operatorname{D\acute{e}r}_A(B,L)
  \longrightarrow0\longrightarrow\cdots
\]
에 다름 아니다.
\oldpage[0\textsubscript{IV-1}]{138}
여기서 미분 연산자는 \sref{0.20.2.1-ko}의 표기에 따른
$v^0$이다. 따라서 명제는 완전열
\sref{0.20.2.4.2-ko}과 코호몰로지 가군의 정의
\begin{equation}
\tag{20.6.6.3}
\label{0.20.6.6.3-ko}
  \mathrm H^i(\mathrm K_\bullet,L)
  =\mathrm H^i
  (\operatorname{Hom}_C(\mathrm K_\bullet,L))
\end{equation}
에서 따른다.
\end{proof}

환 준동형의 가환도식
\[
\xymatrix{
  A' \ar[r] & B' \ar[r] & C' \\
  A \ar[u] \ar[r] & B \ar[u] \ar[r] & C \ar[u]
}
\]
이 주어지면 디-준동형들 \sref{0.20.5.4.3-ko}은 가군
복합체의 디-준동형
\begin{equation}
\tag{20.6.6.4}
\label{0.20.6.6.4-ko}
  \mathrm K_\bullet(C/B/A)\longrightarrow
  \mathrm K_\bullet(C'/B'/A')
\end{equation}
을 정의한다. 이로부터 호몰로지 또는 코호몰로지에 유도되는
디-준동형들은 공식 \sref{0.20.6.5.2-ko},
\sref{0.20.6.6.1-ko}, \sref{0.20.6.6.2-ko} 아래에서 이미
\sref{0.20.5.4.3-ko}, \sref{0.20.6.4.2-ko},
\sref{0.20.2.1-ko}, \sref{0.18.3.7.2-ko}에 정의한
디-준동형들과 동일시된다. 마지막 동일시에서는 완전열
\sref{0.20.2.4.2-ko}의 연산자 $\partial$의 정의를 고려한다.

\begin{env}[20.6.7]
\label{0.20.6.7-ko}
$C$-가군의 복합체 $\mathrm K_\bullet$와 $C$-가군 $L$에 대하여
표준 준동형
\[
  \alpha_i:\mathrm H^i(\mathrm K_\bullet,L)
  \longrightarrow
  \operatorname{Hom}_C(\mathrm H_i(\mathrm K_\bullet),L)
\]
이 존재함은 알려져 있다 (M, IV, 6). 여기서 표준 준동형
\[
  \alpha_1:\mathrm H^1(\mathrm K_\bullet(C/B/A),L)
  \longrightarrow
  \operatorname{Hom}_C
  (\mathrm H_1(\mathrm K_\bullet(C/B/A)),L)
\]
은 제한 준동형
\[
  \operatorname{Hom}_C(\mathrm K_1(C/B/A),L)
  \longrightarrow
  \operatorname{Hom}_C
  (\mathrm H_1(\mathrm K_\bullet(C/B/A)),L)
\]
에서 $\operatorname{Hom}_C(\mathrm K_0,L)$의 상을 법으로 몫을
취하여 곧바로 정의된다. 실제로
$\mathrm H_1(\mathrm K_\bullet(C/B/A))$는
$\mathrm K_1\to\mathrm K_0$의 핵이다.
\sref{0.20.6.6.2-ko}와 \sref{0.20.6.5.2-ko}을 고려하면 표준
$C$-준동형
\begin{equation}
\tag{20.6.7.1}
\label{0.20.6.7.1-ko}
  \operatorname{Exalcom}_{B/A}(C,L)
  \longrightarrow
  \operatorname{Hom}_C(\Upsilon_{C/B/A},L)
\end{equation}
을 얻는다. 이를 구체적으로 쓰면 다음과 같다.
\sref{0.20.2.4.2-ko}에 의하여 $A$-자명한, $L$에 의한 $C$의
모든 $B$-확장은 $B$에서 $L$로 가는 어떤 $A$-미분 $D$, 따라서
\sref{0.20.4.8-ko}에 의하여 어떤 $C$-준동형
\[
  f:\Omega^1_{B/A}\otimes_B C\to L
\]
로부터 얻어진다. 이 확장의 류에는 $f$를
$\Upsilon_{C/B/A}$에 \emph{제한}한 것을 대응시킨다. 이는 실제로
$D$의 선택이 아니라 확장의 류에만 의존한다.
\end{env}

\begin{definition}[20.6.8]
\label{0.20.6.8-ko}
$u:A\to B$, $v:B\to C$를 두 환 준동형이라 하자. $C$-가군
$L$에 의한 $C$의 모든 $B$-확장 $E$ 중 $A$-자명한 것
\sref{0.18.3.7-ko}에 대하여, $E$의 류를 표준 준동형
\sref{0.20.6.7.1-ko}으로 보낸 상인 $C$-준동형
\[
  \Upsilon_{C/B/A}\to L
\]
을 $E$의 \emph{특성 준동형}이라 하고 $\chi_E$로 나타낸다.
\end{definition}

\oldpage[0\textsubscript{IV-1}]{139}
준동형 $\chi_E$는 다음과 같은 다른 방법으로 정의할 수 있다.
\begin{proposition}[20.6.9]
\label{0.20.6.9-ko}
$E$를 $A$-자명한, $C$-가군 $L$에 의한 $C$의 $B$-확장이라 하자
\sref{0.18.3.7-ko}. 그러면 도식
\begin{equation}
\tag{20.6.9.1}
\label{0.20.6.9.1-ko}
\xymatrix{
  0 \ar[r] & \Upsilon_{C/B/A} \ar[r] \ar[d]^{\chi_E} &
  \Omega^1_{B/A}\otimes_B C
    \ar[r]^{v_{C/B/A}}
    \ar[d]^{q_{E/B/A}\otimes1_C} &
  \Omega^1_{C/A} \ar[d]^1 \\
  0 \ar[r] & L \ar[r]_{\delta_{C/E/A}} &
  \Omega^1_{E/A}\otimes_E C \ar[r]_{p_{C/E/A}} &
  \Omega^1_{C/A} \ar[r] & 0
}
\end{equation}
은 가환하고 두 행은 완전하다. 여기서 $q:B\to E$는 $E$의
$B$-확장 구조를 정의하고, $p:E\to C$는 증대 준동형이다.
\end{proposition}

\begin{proof}
도식의 아래 행은 두 준동형
\[
  A\xrightarrow{q\circ u}E,\qquad p:E\to C=E/L
\]
에 대한 열 \sref{0.20.5.12.5-ko}이다. $p$가 전사이고 $E$가
$A$-자명한 확장이므로, 이 열은 \sref{0.20.5.12-ko}, (iii)에
의하여 완전하고 분할된다. \sref{0.20.6.9.1-ko}의 오른쪽 사각형의
가환성은 관계 $v=p\circ q$에서 따른다
\sref{0.20.5.2.7-ko}. 따라서 $v_{C/B/A}$의 핵
$\Upsilon_{C/B/A}$를 $q_{E/B/A}\otimes1_C$로 보낸 상은
$p_{C/E/A}$의 핵 $L$에 포함된다.

한편 $h:C\to E$를 $p$의 오른쪽 역인 $A$-준동형,
$j:L\to E$를 표준 포함사상이라 하자. 그러면 $b\in B$에 대하여
\[
  q(b)=h(v(b))+j(D(b))
\]
이다. 여기서 $D$는 $B$-확장 $E$를 정의하는, $B$에서 $L$로
가는 $A$-미분이다. 어떤 $B$-준동형
$f:\Omega^1_{B/A}\to L$에 대하여
\[
  D=f\circ d_{B/A}
\]
로 쓸 수 있다. \sref{0.20.5.2.6-ko}에 의하여 $b\in B$에 대해
\[
\begin{split}
  q_{E/B/A}(d_{B/A}(b))
  &=d_{E/A}(q(b))\\
  &=d_{E/A}(h(v(b)))
    +d_{E/A}(j(f(d_{B/A}(b)))).
\end{split}
\]
이제 $b_i\in B$, $c_i\in C$이고
\[
  z=\sum_i d_{B/A}(b_i)\otimes c_i
  \in\Omega^1_{B/A}\otimes_B C
\]
라 하자. 그러면
\[
\begin{split}
  (q_{E/B/A}\otimes1_C)(z)
  &=
  \sum_i d_{E/A}(h(v(b_i)))\otimes c_i\\
  &\quad+
  \sum_i d_{E/A}(j(f(d_{B/A}(b_i))))\otimes c_i.
\end{split}
\]
첫째 합에서
\[
  d_{E/A}(h(v(b_i)))
  =h_{E/C/A}(d_{C/A}(v(b_i)))
\]
이므로 \sref{0.20.5.2.6-ko}, 이 합은
\[
  (h_{E/C/A}\otimes1_C)(v_{C/B/A}(z))
\]
와 같다. 따라서 $z\in\Upsilon_{C/B/A}$이면 이 합은 영이고,
$\delta_{C/E/A}$의 정의에 의하여
\[
\begin{split}
  (q_{E/B/A}\otimes1_C)(z)
  &=\sum_i\delta_{C/E/A}
    (c_i f(d_{B/A}(b_i)))\\
  &=\delta_{C/E/A}(f(z)).
\end{split}
\]
이는 \sref{0.20.6.9.1-ko}의 왼쪽 사각형이 가환함을 보인다.
\end{proof}

$\delta_{C/E/A}$가 \emph{단사}이기만 하다고 가정하는 경우에도
(반드시 왼쪽 가역일 필요는 없다), 이 해석을 이용하여
$q_{E/B/A}\otimes1_C$를 $\Upsilon_{C/B/A}$에 제한한 것을
$\chi_E$로 정의할 수 있음에 유의하자.

\begin{corollary}[20.6.10]
\label{0.20.6.10-ko}
$\mathfrak R$이 $B$의 아이디얼이고
\[
  C=B/\mathfrak R,\qquad E=B/\mathfrak R^2
\]
이라 하자. $E$가 $\mathfrak R/\mathfrak R^2$에 의한 $C$의
$A$-자명한 $B$-확장이면 \sref{0.18.3.7-ko}, 특성 준동형
\[
  \chi_E:\Upsilon_{C/B/A}\to\mathfrak R/\mathfrak R^2
\]
은 전단사이다.
\end{corollary}

\begin{proof}
\oldpage[0\textsubscript{IV-1}]{140}
실제로 도식 \sref{0.20.6.9.1-ko}에서 오른쪽의 두 세로 화살표는
전단사 준동형이다 \sref{0.20.5.12-ko}, (ii).
\end{proof}

\begin{theorem}[20.6.11]
\label{0.20.6.11-ko}
$u:A\to B$, $v:B\to C$를 두 환 준동형, $L$을 $C$-가군이라 하자.
다음 조건 중 하나가 성립한다고 가정하자.
\begin{enumerate}
  \item[\rm(i)] $L$은 단사 $C$-가군이다.
  \item[\rm(ii)] $\Upsilon_{C/B/A}$는 $C$-가군
$\Omega^1_{B/A}\otimes_B C$의 직합인자이고,
$u_{C/B/A}:\Omega^1_{C/A}\to\Omega^1_{C/B}$는 오른쪽 가역이다.
\end{enumerate}
그러면 표준 준동형 \sref{0.20.6.7.1-ko}
\[
  \operatorname{Exalcom}_{B/A}(C,L)\longrightarrow
  \operatorname{Hom}_C(\Upsilon_{C/B/A},L)
\]
은 전단사이다.

특히 $\Omega^1_{C/B}$와 $\Omega^1_{C/A}$가 사영 $C$-가군이면,
표준 준동형 \sref{0.20.6.7.1-ko}은 전단사이다.
\end{theorem}

\begin{proof}
조건 (i), (ii)가 각각 \sref{0.20.6.7.1-ko}의 전단사성을 함의하는
사실은 두 경우 모두 $\alpha_1$의 정의에서 따른다. 더 나아가
조건 (ii)는 준동형 \sref{0.20.6.7.1-ko}이 \emph{모든} $C$-가군
$L$에 대하여 전단사이기 위한 \emph{필요충분조건}임에 유의하자
(Bourbaki, \emph{Alg.}, chap. II, 3\textsuperscript{e} éd.,
\S~2, n\textsuperscript{o}~1, prop.~1).
$\Omega^1_{C/B}$와 $\Omega^1_{C/A}$가 사영 $C$-가군이라고 가정하면,
완전열 \sref{0.20.6.1.1-ko}에서
$\operatorname{Ker}(u_{C/B/A})$는 사영 $C$-가군이다. 실제로
\[
  0\longrightarrow\operatorname{Ker}(u_{C/B/A})
  \longrightarrow\Omega^1_{C/A}
  \longrightarrow\Omega^1_{C/B}\longrightarrow0
\]
은 $\Omega^1_{C/B}$가 사영이므로 분할되는 완전열이다. 또한
$\operatorname{Ker}(u_{C/B/A})=\operatorname{Im}(v_{C/B/A})$이므로,
완전열
\[
  0\longrightarrow\Upsilon_{C/B/A}\longrightarrow
  \Omega^1_{B/A}\otimes_B C\longrightarrow
  \operatorname{Im}(v_{C/B/A})\longrightarrow0
\]
도 분할된다.
\end{proof}

\begin{corollary}[20.6.12]
\label{0.20.6.12-ko}
$C$가 형식적으로 매끄러운 $A$-대수라 하자. 그러면 표준 준동형
\begin{equation}
\tag{20.6.12.1}
\label{0.20.6.12.1-ko}
  \operatorname{Exalcom}_B(C,L)\longrightarrow
  \operatorname{Hom}_C(\Upsilon_{C/B/A},L)
\end{equation}
이 존재한다. 더욱이 \sref{0.20.6.11-ko}의 조건 (i), (ii) 중
하나가 성립하면 이 준동형은 전단사이다.
\end{corollary}

\begin{proof}
실제로 $C$에 대한 가정으로부터
$\operatorname{Exalcom}_B(C,L)=\operatorname{Exalcom}_{B/A}(C,L)$이고,
준동형 \sref{0.20.6.12.1-ko}은 바로 \sref{0.20.6.7.1-ko}이다.
\end{proof}

\begin{corollary}[20.6.13]
\label{0.20.6.13-ko}
$C$가 형식적으로 매끄러운 $A$-대수이고 $\Omega^1_{C/B}$가
사영 $C$-가군이면, 준동형 \sref{0.20.6.12.1-ko}은 전단사이다.
\end{corollary}

\begin{proof}
실제로 이때 이산 위상을 생각하면 $\Omega^1_{C/A}$가 사영
$C$-가군임은 알려져 있다 \sref{0.20.4.9-ko}.
\end{proof}

\begin{env}[20.6.14]
\label{0.20.6.14-ko}
기호는 그대로 두고, 이제 더 나아가 $A$, $B$, $C$가
$\Lambda$-대수이고 $u$, $v$가 $\Lambda$-준동형이라고 가정하자.
이는 세 환 준동형
\[
  \Lambda\xrightarrow{s}A\xrightarrow{u}B\xrightarrow{v}C
\]
을 주는 것과 같다.

\oldpage[0\textsubscript{IV-1}]{141}
그러면 불완전성 가군 $\Upsilon_{C/B/A}$ 이외에도 불완전성 가군
$\Upsilon_{B/A/\Lambda}$, $\Upsilon_{C/A/\Lambda}$,
$\Upsilon_{C/B/\Lambda}$가 있고, 표준 $C$-가군 준동형
\sref{0.20.6.4.2-ko}
\begin{align}
\tag{20.6.14.1}
\label{0.20.6.14.1-ko}
  u':\Upsilon_{C/A/\Lambda}&\longrightarrow\Upsilon_{C/B/\Lambda},\\
\tag{20.6.14.2}
\label{0.20.6.14.2-ko}
  s':\Upsilon_{C/B/\Lambda}&\longrightarrow\Upsilon_{C/B/A}
\end{align}
은 이미 정의되어 있다.

가환 도식 \sref{0.20.5.7.3-ko}
\begin{equation}
\tag{20.6.14.3}
\label{0.20.6.14.3-ko}
\xymatrix{
  \Omega^1_{A/\Lambda}\otimes_A B \ar[r] \ar[d] &
  \Omega^1_{B/\Lambda} \ar[r] \ar[d] &
  \Omega^1_{B/A} \ar[r] \ar[d] & 0 \\
  \Omega^1_{A/\Lambda}\otimes_A C \ar[r] &
  \Omega^1_{C/\Lambda} \ar[r] & \Omega^1_{C/A} \ar[r] & 0
}
\end{equation}
에서와 같이 아래 행은 $C$-가군들로 이루어져 있으므로, 텐서화하여
가환 도식
\begin{equation}
\tag{20.6.14.4}
\label{0.20.6.14.4-ko}
\xymatrix{
  \Omega^1_{A/\Lambda}\otimes_A C \ar[r] \ar@{=}[d] &
  \Omega^1_{B/\Lambda}\otimes_B C \ar[r] \ar[d] &
  \Omega^1_{B/A}\otimes_B C \ar[r] \ar[d] & 0 \\
  \Omega^1_{A/\Lambda}\otimes_A C \ar[r] &
  \Omega^1_{C/\Lambda} \ar[r] & \Omega^1_{C/A} \ar[r] & 0
}
\end{equation}
을 얻는다. 여기서 첫째 행도 여전히 완전하고 왼쪽 세로 화살표는
항등사상이다. 이제
\begin{equation}
\tag{20.6.14.5}
\label{0.20.6.14.5-ko}
  \Upsilon^C_{B/A/\Lambda}=
  \operatorname{Ker}(u_{B/A/\Lambda}\otimes1_C)=
  \operatorname{Ker}(\Omega^1_{A/\Lambda}\otimes_A C\longrightarrow
       \Omega^1_{B/\Lambda}\otimes_B C)
\end{equation}
라 두면, $\Upsilon_{C/A/\Lambda}$의 정의를 고려하여 유일한
$C$-준동형
\begin{equation}
\tag{20.6.14.6}
\label{0.20.6.14.6-ko}
  v':\Upsilon^C_{B/A/\Lambda}\longrightarrow\Upsilon_{C/A/\Lambda}
\end{equation}
이 있어서 도식
\begin{equation}
\tag{20.6.14.7}
\label{0.20.6.14.7-ko}
\xymatrix{
  0 \ar[r] & \Upsilon^C_{B/A/\Lambda} \ar[r] \ar[d]^{v'} &
  \Omega^1_{A/\Lambda}\otimes_A C \ar[r] \ar@{=}[d] &
  \Omega^1_{B/\Lambda}\otimes_B C \ar[r] \ar[d] &
  \Omega^1_{B/A}\otimes_B C \ar[r] \ar[d] & 0 \\
  0 \ar[r] & \Upsilon_{C/A/\Lambda} \ar[r] &
  \Omega^1_{A/\Lambda}\otimes_A C \ar[r] &
  \Omega^1_{C/\Lambda} \ar[r] & \Omega^1_{C/A} \ar[r] & 0
}
\end{equation}
이 가환함을 알 수 있다. 이 도식의 두 행은 완전하다.

$\Lambda=\mathbb{Z}$이면 $\Upsilon^C_{B/A/\mathbb{Z}}$ 대신
$\Upsilon^C_{B/A}$라 쓴다. $\Lambda$가 소체이면
$\Upsilon^C_{B/A/\Lambda}=\Upsilon^C_{B/A}$이다.
\end{env}

\begin{env}[20.6.15]
\label{0.20.6.15-ko}
앞의 가군들 사이의 관계를 연구하기 위하여, 한편으로 $B$-가군의
복합체 $\mathrm K_\bullet(B/A/\Lambda)$를, 다른 한편으로 $C$-가군의
복합체
\oldpage[0\textsubscript{IV-1}]{142}
$\mathrm K_\bullet(C/A/\Lambda)$와
$\mathrm K_\bullet(C/B/\Lambda)$를 도입하고
\sref{0.20.6.5-ko}, 더 나아가 다음 $C$-가군의 복합체들을
도입한다. 먼저
\begin{equation}
\tag{20.6.15.1}
\label{0.20.6.15.1-ko}
  \mathrm K^C_\bullet(B/A/\Lambda)=
  \mathrm K_\bullet(B/A/\Lambda)\otimes_B C
\end{equation}
라 둔다. 또한 차수 $0$과 $1$을 제외한 항은 영이고
\begin{equation}
\tag{20.6.15.2}
\label{0.20.6.15.2-ko}
  \mathrm T_0(C/B/\Lambda)=\mathrm T_1(C/B/\Lambda)=
  \Omega^1_{B/\Lambda}\otimes_B C
\end{equation}
이며 미분 연산자가 항등사상인 $C$-가군 복합체를
$\mathrm T_\bullet(C/B/\Lambda)$로 나타낸다. 따라서 이 복합체는
$0$과 호모토픽이다. 마지막으로
\begin{equation}
\tag{20.6.15.3}
\label{0.20.6.15.3-ko}
  \mathrm K'_\bullet(C/A/\Lambda)=
  \mathrm K_\bullet(C/A/\Lambda)\oplus
  \mathrm T_\bullet(C/B/\Lambda)
\end{equation}
라 두자. $\mathrm T_\bullet$가 자명하므로
\begin{equation}
\tag{20.6.15.4}
\label{0.20.6.15.4-ko}
  \mathrm H^i(\mathrm K'_\bullet,L)=
  \mathrm H^i(\mathrm K_\bullet,L)
  \quad\text{그리고}\quad
  \mathrm H_i(\mathrm K'_\bullet)=
  \mathrm H_i(\mathrm K_\bullet)
\end{equation}
임은 명백하다. 첫째 등식은 모든 $C$-가군 $L$과 모든 $i$에
대하여, 둘째 등식은 모든 $i$에 대하여 성립한다.
\end{env}

\begin{env}[20.6.16]
\label{0.20.6.16-ko}
이제 각 차수에서 \emph{분할되는} 복합체의 완전열
\begin{equation}
\tag{20.6.16.1}
\label{0.20.6.16.1-ko}
  0\longrightarrow\mathrm K^C_\bullet(B/A/\Lambda)
  \xrightarrow{j}\mathrm K'_\bullet(C/A/\Lambda)
  \xrightarrow{p}\mathrm K_\bullet(C/B/\Lambda)\longrightarrow0
\end{equation}
을 다음과 같이 정의하자. 잠시 표준 준동형
$u_{B/A/\Lambda}\otimes1_C$와 $v_{C/B/\Lambda}$를 각각
\[
  f:\Omega^1_{A/\Lambda}\otimes_A C\longrightarrow
    \Omega^1_{B/\Lambda}\otimes_B C,
  \qquad
  g:\Omega^1_{B/\Lambda}\otimes_B C\longrightarrow
    \Omega^1_{C/\Lambda}
\]
라 나타내자. 그 합성은
$g\circ f=(v\circ u)_{C/A/\Lambda}$이다
\sref{0.20.6.14.4-ko}. 이제
\[
\begin{gathered}
  j_1(x)=(x,f(x)),\qquad p_1(y,z)=z-f(y),\\
  j_0(x)=(g(x),x),\qquad p_0(y,z)=g(z)-y
\end{gathered}
\]
로 잡는다. 그러면
$\operatorname{Im}(j_1)=\operatorname{Ker}(p_1)$은 $f$의 그래프이며
$\{0\}\oplus\mathrm T_1$의 직합보충이고,
$\operatorname{Im}(j_0)=\operatorname{Ker}(p_0)$은 $g$의 그래프이며
$\mathrm K_0(C/A/\Lambda)\oplus\{0\}$의 직합보충이다. 도식
\[
\xymatrix{
  0 \ar[r] & \mathrm K^C_1(B/A/\Lambda) \ar[r]^{j_1} \ar[d] &
  \mathrm K'_1(C/A/\Lambda) \ar[r]^{p_1} \ar[d] &
  \mathrm K_1(C/B/\Lambda) \ar[r] \ar[d] & 0 \\
  0 \ar[r] & \mathrm K^C_0(B/A/\Lambda) \ar[r]_{j_0} &
  \mathrm K'_0(C/A/\Lambda) \ar[r]_{p_0} &
  \mathrm K_0(C/B/\Lambda) \ar[r] & 0
}
\]
에서 세로 화살표는 미분 연산자이다. 이 도식의 가환성은 곧바로
확인된다.
\end{env}

\begin{theorem}[20.6.17]
\label{0.20.6.17-ko}
$C$-가군의 완전열
\begin{equation}
\tag{20.6.17.1}
\label{0.20.6.17.1-ko}
  0\longrightarrow\Upsilon^C_{B/A/\Lambda}\xrightarrow{v'}
  \Upsilon_{C/A/\Lambda}\xrightarrow{u'}
  \Upsilon_{C/B/\Lambda}\xrightarrow{\partial}
  \Omega^1_{B/A}\otimes_B C\xrightarrow{v_{C/B/A}}
  \Omega^1_{C/A}\xrightarrow{u_{C/B/A}}
  \Omega^1_{C/B}\longrightarrow0
\end{equation}
이 존재한다.
\oldpage[0\textsubscript{IV-1}]{143}
여기서 경계 연산자 $\partial$는 합성
\begin{equation}
\tag{20.6.17.2}
\label{0.20.6.17.2-ko}
  \Upsilon_{C/B/\Lambda}\xrightarrow{s'}\Upsilon_{C/B/A}
  \longrightarrow\Omega^1_{B/A}\otimes_B C
\end{equation}
이고, 둘째 화살표는 표준 단사 준동형이다.
\end{theorem}

\begin{proof}
복합체의 완전열 \sref{0.20.6.16.1-ko}에 대한 긴 호몰로지 완전열을
쓰면 \sref{0.20.6.17.1-ko}을 얻는다. 호몰로지는 차수 $0$과 $1$을
제외하면 영이고, 이 완전열의 준동형들 가운데 $j$와 $p$의 함자성에서
나오는 것들이 바로 명제에 적힌 준동형들임은 명백하다.
이제 $\partial$가 \sref{0.20.6.17.2-ko}과 같음을 확인하면 된다.
$z\in\Upsilon_{C/B/\Lambda}$는 $(0,z)$를 $p_1$로 보낸 상이므로,
$\partial(z)$는 $z$를 표준 준동형
\[
  s_{B/A/\Lambda}\otimes1_C:
  \Omega^1_{B/\Lambda}\otimes_B C\longrightarrow
  \Omega^1_{B/A}\otimes_B C
\]
로 보낸 상이다. 따라서 우리의 주장은 가환 도식
\begin{equation}
\tag{20.6.17.3}
\label{0.20.6.17.3-ko}
\xymatrix{
  \Upsilon_{C/B/\Lambda} \ar[r]^{s'} \ar[dd] &
  \Upsilon_{C/B/A} \ar[dd] \\
  & \\
  \Omega^1_{B/\Lambda}\otimes_B C \ar[r] &
  \Omega^1_{B/A}\otimes_B C
}
\end{equation}
에서 따른다 \sref{0.20.6.4.3-ko}.
\end{proof}

\begin{corollary}[20.6.18]
\label{0.20.6.18-ko}
\begin{enumerate}
  \item[\rm(i)] $C$-가군의 열
\begin{equation}
\tag{20.6.18.1}
\label{0.20.6.18.1-ko}
  0\longrightarrow\Upsilon^C_{B/A/\Lambda}\xrightarrow{v'}
  \Upsilon_{C/A/\Lambda}\xrightarrow{u'}
  \Upsilon_{C/B/\Lambda}\xrightarrow{s'}
  \Upsilon_{C/B/A}\longrightarrow0
\end{equation}
은 완전하다.
  \item[\rm(ii)] $B$가 $\Lambda$에 대하여 형식적으로 매끄러운
  $A$-대수이면 $\Upsilon^C_{B/A/\Lambda}=0$이다.
\end{enumerate}
\end{corollary}

\begin{proof}
(i) \sref{0.20.6.17.1-ko}에서 $\Upsilon_{C/B/\Lambda}$를
$\partial$로 보낸 상은 $v_{C/B/A}$의 핵, 곧 정의에 의하여
$\Upsilon_{C/B/A}$이다.

(ii) 가정에 의하여 완전열
\[
  0\longrightarrow\Omega^1_{A/\Lambda}\otimes_A B
  \longrightarrow\Omega^1_{B/\Lambda}
  \longrightarrow\Omega^1_{B/A}\longrightarrow0
\]
은 분할된다 \sref{0.20.5.7-ko}. 따라서 $C$와 텐서화한 열
\[
  0\longrightarrow\Omega^1_{A/\Lambda}\otimes_A C
  \longrightarrow\Omega^1_{B/\Lambda}\otimes_B C
  \longrightarrow\Omega^1_{B/A}\otimes_B C\longrightarrow0
\]
도 여전히 완전하고, 이로부터 주장이 따른다.
\end{proof}

\begin{corollary}[20.6.19]
\label{0.20.6.19-ko}
$K$를 체, $E$, $F$를 $K\subset E\subset F$인 두 $K$의
확대체라 하자.
\begin{enumerate}
  \item[\rm(i)] $F$가 $E$의 분리 확대체이면
  $\Upsilon_{F/E/K}=0$이다. 다시 말해 표준 준동형
  $\Omega^1_{E/K}\otimes_E F\to\Omega^1_{F/K}$는 단사이다.
  \item[\rm(ii)] 역으로 $F$가 $K$의 분리 확대체이고
  $\Upsilon_{F/E/K}=0$이면, $F$는 $E$의 분리 확대체이다.
\end{enumerate}
\end{corollary}

\begin{proof}
$P$를 $K$의 소체라 하면 완전열 \sref{0.20.6.18-ko}
\[
  \Upsilon_{F/K/P}\longrightarrow\Upsilon_{F/E/P}
  \longrightarrow\Upsilon_{F/E/K}\longrightarrow0
\]
을 얻는다. $\Upsilon_{F/E/P}=\Upsilon_{F/E}=0$이면
$\Upsilon_{F/E/K}=0$이므로 \sref{0.20.6.3-ko}에 의하여 (i)가
따른다. 역으로 $\Upsilon_{F/E/K}=0$이고
$\Upsilon_{F/K/P}=\Upsilon_{F/K}=0$이면
$\Upsilon_{F/E/P}=0$이므로 \sref{0.20.6.3-ko}에 의하여 (ii)가
따른다.
\end{proof}

\oldpage[0\textsubscript{IV-1}]{144}
\begin{corollary}[20.6.20]
\label{0.20.6.20-ko}
\begin{enumerate}
  \item[\rm(i)] $K$가 $k$의 대수적 분리 확대체이면
  $\Omega^1_{K/k}=0$이다.
  \item[\rm(ii)] $k$를 표수 $0$인 체, $K$를 $k$의 확대체라 하자.
  $\Omega^1_{K/k}=0$이기 위한 필요충분조건은 $K$가 $k$의 대수적
  확대체인 것이다. 특히 표수 $0$인 체 $K$에 대하여
  $\Omega^1_K=0$이기 위한 필요충분조건은 $K$가 $\mathbb{Q}$의
  대수적 확대체인 것이다 (21.4.4와 21.7.5 참조).
\end{enumerate}
\end{corollary}

\begin{proof}
(i) 각 $x\in K$에 대하여 $f$를 $k$ 위에서 $x$의 최소다항식이라
하자. $f'(x)\ne0$이고
$d_{K/k}(f(x))=f'(x)d_{K/k}(x)=0$이므로 $d_{K/k}(x)=0$이다.
따라서 주장은 \sref{0.20.4.7-ko}에서 따른다.

(ii) $k\subset L\subset K$이고 $K$가 $L$의 대수적 확대체인
$k$의 순수 초월 확대체 $L$이 존재한다. $K$는 $L$ 위에서
분리이므로 \sref{0.20.6.19-ko}, (i)에 의하여 열
\[
  0\longrightarrow\Omega^1_{L/k}\otimes_L K
  \longrightarrow\Omega^1_{K/k}
  \longrightarrow\Omega^1_{K/L}\longrightarrow0
\]
은 완전하고, (i)에 의하여 $\Omega^1_{K/L}=0$이다. 따라서
$\Omega^1_{K/k}=0$은 $\Omega^1_{L/k}=0$과 동치이다. $L$은
$k$의 순수 초월 확대체이므로, \sref{0.20.5.10-ko}에 의하여
$\Omega^1_{L/k}=0$은 $L=k$와 동치이다.
\end{proof}

\begin{remarks}[20.6.21]
\label{0.20.6.21-ko}
\begin{enumerate}
  \item[\rm(i)] $\Upsilon_{B/A/\Lambda}$는
  \[
    u_{B/A/\Lambda}:\Omega^1_{A/\Lambda}\otimes_A B
    \longrightarrow\Omega^1_{B/\Lambda}
  \]
  의 핵이고 $\Upsilon^C_{B/A/\Lambda}$는
  $u_{B/A/\Lambda}\otimes1_C$의 핵이므로, 표준 준동형
\begin{equation}
\tag{20.6.21.1}
\label{0.20.6.21.1-ko}
  \Upsilon_{B/A/\Lambda}\otimes_B C
  \longrightarrow\Upsilon^C_{B/A/\Lambda}
\end{equation}
  이 존재한다. 완전열
\begin{equation}
\tag{20.6.21.2}
\label{0.20.6.21.2-ko}
  0\longrightarrow\Upsilon_{B/A/\Lambda}\otimes_B C
  \longrightarrow\Omega^1_{A/\Lambda}\otimes_A C
  \longrightarrow\Omega^1_{B/\Lambda}\otimes_B C
  \longrightarrow\Omega^1_{B/A}\otimes_B C\longrightarrow0
\end{equation}
  이 성립하면 이 준동형은 전단사이다. 이는 다음 경우에 성립한다.
\begin{enumerate}
  \item[$1^\circ$] $C$가 평탄 $B$-가군인 경우.
  \item[$2^\circ$] $B$-가군 $\Omega^1_{B/\Lambda}$와
  $\Omega^1_{B/A}$가 \emph{평탄}인 경우. 실제로 이때
  $\operatorname{Ker}(\Omega^1_{B/\Lambda}\to\Omega^1_{B/A})$도
  평탄하고 ($\mathbf{0_I}$, 6.1.2), 따라서
  \sref{0.20.6.21.2-ko}은 ($\mathbf{0_I}$, 6.1.2)에 의하여
  완전하다.
\end{enumerate}
  \item[\rm(ii)] 환 준동형의 가환도식
  \[
  \xymatrix{
    \Lambda' \ar[r] & A' \ar[r] & B' \ar[r] & C' \\
    \Lambda \ar[u] \ar[r] & A \ar[u] \ar[r] &
    B \ar[u] \ar[r] & C \ar[u]
  }
  \]
  을 생각하자. \sref{0.20.6.16-ko}의 정의에 의하여 복합체의
  가환도식
  \[
  \xymatrix{
    0 \ar[r] & \mathrm K^C_\bullet(B/A/\Lambda) \ar[r] \ar[d] &
    \mathrm K'_\bullet(C/A/\Lambda) \ar[r] \ar[d] &
    \mathrm K_\bullet(C/B/\Lambda) \ar[r] \ar[d] & 0 \\
    0 \ar[r] & \mathrm K^{C'}_\bullet(B'/A'/\Lambda') \ar[r] &
    \mathrm K'_\bullet(C'/A'/\Lambda') \ar[r] &
    \mathrm K_\bullet(C'/B'/\Lambda') \ar[r] & 0
  }
  \]
  을 얻는다. 여기서 세로 화살표는 \sref{0.20.6.6.4-ko}에서
  나온다.
\oldpage[0\textsubscript{IV-1}]{145}
  따라서 호몰로지를 취하면 \emph{가환도식}
\begin{equation}
\tag{20.6.21.3}
\label{0.20.6.21.3-ko}
\xymatrix@C=0.7em{
  0 \ar[r] & \Upsilon^C_{B/A/\Lambda} \ar[r] \ar[d] &
  \Upsilon_{C/A/\Lambda} \ar[r] \ar[d] &
  \Upsilon_{C/B/\Lambda} \ar[r] \ar[d] &
  \Omega^1_{B/A}\otimes_B C \ar[r] \ar[d] &
  \Omega^1_{C/A} \ar[r] \ar[d] &
  \Omega^1_{C/B} \ar[r] \ar[d] & 0 \\
  0 \ar[r] & \Upsilon^{C'}_{B'/A'/\Lambda'} \ar[r] &
  \Upsilon_{C'/A'/\Lambda'} \ar[r] &
  \Upsilon_{C'/B'/\Lambda'} \ar[r] &
  \Omega^1_{B'/A'}\otimes_{B'} C' \ar[r] &
  \Omega^1_{C'/A'} \ar[r] & \Omega^1_{C'/B'} \ar[r] & 0
}
\end{equation}
  을 얻는다. \sref{0.20.6.18.1-ko}에 대해서도 이와 비슷한
  가환도식을 얻는다.
\end{enumerate}
\end{remarks}

\begin{proposition}[20.6.22]
\label{0.20.6.22-ko}
$s:\Lambda\to A$, $u:A\to B$를 두 환 준동형,
$\mathfrak R$을 $B$의 아이디얼, $C$를 몫환 $B/\mathfrak R$이라
하자. $E=B/\mathfrak R^2$가 $\mathfrak R/\mathfrak R^2$에 의한
$C$의 $\Lambda$-자명한 $B$-확장이라고 가정하자. 그러면 완전열
\begin{equation}
\tag{20.6.22.1}
\label{0.20.6.22.1-ko}
  0\longrightarrow\Upsilon^C_{B/A/\Lambda}\xrightarrow{v'}
  \Upsilon_{C/A/\Lambda}\xrightarrow{\chi_E}
  \mathfrak R/\mathfrak R^2\xrightarrow{\delta_{C/B/A}}
  \Omega^1_{B/A}\otimes_B C\xrightarrow{v_{C/B/A}}
  \Omega^1_{C/A}\longrightarrow0
\end{equation}
이 존재한다.
\end{proposition}

\begin{proof}
실제로 $v:B\to C$는 전사이므로 $\Omega^1_{C/B}=0$이다
\sref{0.20.4.12-ko}. 또한 \sref{0.20.6.10-ko}에 의하여
$\Upsilon_{C/B/\Lambda}$는 $\mathfrak R/\mathfrak R^2$와 표준적으로
동일시된다. 따라서 완전열 \sref{0.20.6.17.1-ko}과
\sref{0.20.6.18.1-ko}을 적용하고, 가환도식
\[
\xymatrix{
  \Upsilon_{C/A/\Lambda} \ar[r] \ar[d]_{\chi_E} &
  \Omega^1_{A/\Lambda}\otimes_A C \ar[d] \\
  \Upsilon_{C/E/\Lambda}=\Upsilon_{C/B/\Lambda}=\mathfrak R/\mathfrak R^2
    \ar[r]^-{\delta_{C/B/A}} \ar[d]_{\ell} &
  \Omega^1_{B/A}\otimes_B C=\Omega^1_{E/A}\otimes_E C \ar[d] \\
  \mathfrak R/\mathfrak R^2 \ar[r]^{\delta_{C/B/A}} &
  \Omega^1_{B/A}\otimes_B C
}
\]
과 도식 \sref{0.20.6.17.3-ko}의 가환성을 이용하면 된다.
\end{proof}

이로써 $A$가 $\Lambda$-대수이고 $B/\mathfrak R^2$가
$\Lambda$-자명하게 되는 환 $\Lambda$가 존재하는 경우에
$\delta_{C/B/A}$의 \emph{핵}을 구체화하였다. 특히 $C$가 이산
위상에 대하여 \emph{형식적으로 매끄러운 $\Lambda$-대수}이면
이 조건이 성립함에 유의하자.

더 나아가 환 준동형의 가환도식
\[
\xymatrix{
  \Lambda' \ar[r] & A' \ar[r] & B' \\
  \Lambda \ar[u] \ar[r] & A \ar[u] \ar[r] & B \ar[u]^f
}
\]
이 주어졌다고 가정하자.
\oldpage[0\textsubscript{IV-1}]{146}
$\mathfrak R'$을 $f(\mathfrak R)\subset\mathfrak R'$인 $B'$의
아이디얼이라 하고, $E'=B'/\mathfrak R'^2$이
$\mathfrak R'/\mathfrak R'^2$에 의한
$C'=B'/\mathfrak R'$의 $\Lambda'$-자명한 $B'$-확장이라고 하자.
그러면 가환도식
\begin{equation}
\tag{20.6.22.2}
\label{0.20.6.22.2-ko}
\xymatrix@C=0.7em{
  0 \ar[r] & \Upsilon^{C'}_{B'/A'/\Lambda'} \ar[r] &
  \Upsilon_{C'/A'/\Lambda'} \ar[r] &
  \mathfrak R'/\mathfrak R'^2 \ar[r] &
  \Omega^1_{B'/A'}\otimes_{B'} C' \ar[r] &
  \Omega^1_{C'/A'} \ar[r] & 0 \\
  0 \ar[r] & \Upsilon^C_{B/A/\Lambda} \ar[r] \ar[u] &
  \Upsilon_{C/A/\Lambda} \ar[r] \ar[u] &
  \mathfrak R/\mathfrak R^2 \ar[r] \ar[u] &
  \Omega^1_{B/A}\otimes_B C \ar[r] \ar[u] &
  \Omega^1_{C/A} \ar[r] \ar[u] & 0
}
\end{equation}
을 얻는다. 이는 \sref{0.20.6.21.3-ko}과
\sref{0.20.5.11.3-ko}에서 따른다.

\begin{corollary}[20.6.23]
\label{0.20.6.23-ko}
\sref{0.20.6.22-ko}의 가정 아래, 더 나아가 $B$가 형식적으로
매끄러운 $A$-대수라고 가정하자. 그러면 완전열
\begin{equation}
\tag{20.6.23.1}
\label{0.20.6.23.1-ko}
  0\longrightarrow\Upsilon_{C/A/\Lambda}\xrightarrow{\chi_E}
  \mathfrak R/\mathfrak R^2\xrightarrow{\delta_{C/B/A}}
  \Omega^1_{B/A}\otimes_B C\longrightarrow
  \Omega^1_{C/A}\longrightarrow0
\end{equation}
이 존재한다.
\end{corollary}

\begin{proof}
이는 \sref{0.20.6.18-ko}, (ii)에서 따른다.
\end{proof}

\begin{env}[20.6.24]
\label{0.20.6.24-ko}
\sref{0.20.6.22-ko}의 가정이 성립할 때 특성 준동형 $\chi_E$를
\emph{$\Lambda$와 아이디얼 $\mathfrak R$에 대한 $A$-대수 $B$의
특성 준동형}이라고도 한다. 혼동의 염려가 없으면 뒤의 조건들을
생략하며, 때로는 이를 $\chi_B$ 또는 $\chi_{B/A}$로 나타낸다.
\end{env}

\begin{proposition}[20.6.25]
\label{0.20.6.25-ko}
$s:\Lambda\to A$, $u:A\to B$, $v:B\to C$를 세 환 준동형,
$L$을 $C$-가군이라 하자. 그러면 완전열
\begin{equation}
\tag{20.6.25.1}
\label{0.20.6.25.1-ko}
\begin{split}
  0\to\operatorname{D\acute{e}r}_B(C,L)\xrightarrow{u^1}
  \operatorname{D\acute{e}r}_A(C,L)\xrightarrow{v^1}
  \operatorname{D\acute{e}r}_A(B,L)\xrightarrow{\partial} \\
  \to\operatorname{Exalcom}_{B/\Lambda}(C,L)\xrightarrow{u^1}
  \operatorname{Exalcom}_{A/\Lambda}(C,L)\xrightarrow{v^1}
  \operatorname{Exalcom}_{A/\Lambda}(B,L)\to0
\end{split}
\end{equation}
이 존재한다. 여기서 $u^1$, $v^1$은 각각
\sref{0.18.3.6.4-fr}와 \,\sref{0.18.3.6.2-fr}에서 정의한
준동형이고, $\partial$는 \sref{0.20.2.2-ko}에서와 같이 정의한다.
\end{proposition}

\begin{proof}
실제로 완전열 \sref{0.20.6.16.1-ko}이 분할되므로 완전열
\[
  0\to\operatorname{Hom}_C(\mathrm{K}_\bullet(C/B/\Lambda),L)\to
  \operatorname{Hom}_C(\mathrm{K}'_\bullet(C/A/\Lambda),L)\to
  \operatorname{Hom}_C(\mathrm{K}^C_\bullet(B/A/\Lambda),L)\to0
\]
을 얻는다. 이 복합체에 코호몰로지 완전열을 적용하고
\sref{0.20.6.15.4-ko}와 \sref{0.20.6.6-ko}을 고려하면
\sref{0.20.6.25.1-ko}을 얻는다. 실제로 정의에 의하여
\[
  \operatorname{Hom}_C(\mathrm{K}^C_\bullet(B/A/\Lambda),L)
  =\operatorname{Hom}_B(\mathrm{K}_\bullet(B/A/\Lambda),L)
\]
이고, $u^1$, $v^1$이 각각 \sref{0.18.3.6.4-fr}와
\sref{0.18.3.6.2-fr}의 준동형과 동일함은
\sref{0.20.6.6.4-ko}에서 따른다.
\end{proof}

\begin{remark}[20.6.26]
\label{0.20.6.26-ko}
이 번호에서 복합체 $\mathrm{K}_\bullet(C/B/A)$는 몇몇 함자적 거동의
서술을 간단히 하기 위한 기술적 장치로 등장하였다. 실제로 이 복합체들을
``호모토피를 제외하고'' 본 $C$-가군 복합체들의 범주의
\oldpage[0\textsubscript{IV-1}]{147}
대상으로 생각하면(즉 사상은 서로 호모토픽인 준동형사상들을 동일시하여
얻는 류이다), 이들은 $\Omega^1_{C/B}$와 $\Upsilon_{C/B/A}$로 이루어진
쌍보다 더 세밀한 주목할 만한 불변량이다. $k$가 소체이고 $C$가
형식적으로 매끄러운 $k$-대수일 때(예컨대 $k$의 어떤 확대체 위 유한형인
정칙환일 때; ($\mathbf{IV}$, 6.8.6) 참조), 복합체
$\mathrm{K}_\bullet(C/A/k)$를 $k$를 제외하고 $C$와 $A$만을 사용하여
기술할 수 있음을 보일 수 있다. 실제로 $C$를 $A$ 위의 다항식 대수
$B$의 한 아이디얼 $\mathfrak Q$에 의한 몫으로 나타내고, 0이 아닌 항이
두 개인 복합체 $\mathrm{F}_\bullet(C/A)$
\[
  \cdots\to0\to\mathfrak Q/\mathfrak Q^2\to
  \Omega^1_{B/A}\otimes_B C\to0\to\cdots
\]
를 생각한다. 그 호몰로지는 \sref{0.20.6.23.1-ko}에 의하여 실제로
$\mathrm{K}_\bullet(C/A/k)$의 호몰로지와 일치한다. 호몰로지 대수의
관점에서 $\operatorname{Spec}(C)$의 $\operatorname{Spec}(A)$ 위
여접다발의 역할을 하는 이 복합체 $\mathrm{F}_\bullet(C/A)$는 이 책에서
연접층의 쌍대성과 리만--로흐 정리를 다루는 장들에서 중요한 자리를
차지할 것이다.
\end{remark}

\subsection{위상환으로의 일반화}
\label{subsection:0.20.7-ko}

\begin{env}[20.7.1]
\label{0.20.7.1-ko}
\sref{0.20.5.2-ko}와 \sref{0.20.5.3-ko}에서 환 $A$, $B$, $C$를
위상환이라 하고 환 준동형 $u$, $v$를 연속이라 하면, 정의에서 곧
준동형 $u_{C/B/A}$와 $v_{C/B/A}$가 \emph{연속}임이 따른다. 여기서
$\Omega^1_{B/A}\otimes_B C$에는 당연히 텐서곱 위상을 주어야 한다.

더욱이
$u_{C/B/A}:\Omega^1_{C/A}\to\Omega^1_{C/B}$는 위상 $C$-가군의
\emph{엄밀 전사 준동형}이다. 실제로 텐서곱 위상의 정의를 고려하면
표준 준동형 $C\otimes_A C\to C\otimes_B C$가 그러하다. 따라서
몫으로 이행하여 대응하는 준동형
$\mathrm{P}^1_{C/A}\to\mathrm{P}^1_{C/B}$도 엄밀 전사이고,
$u_{C/B/A}$는 이를 $\Omega^1_{C/A}$에 제한한 것이다.

\sref{0.20.5.5-ko}에서 $A'$과 $B$가 위상 $A$-대수이고 $B'$과
$\Omega^1_{B/A}\otimes_B B'$에 텐서곱 위상을 주면, 표준 준동형
\sref{0.20.5.5.1-ko}은 위상동형사상이다. 이는
$\mathrm{P}^1_{B/A}$가 $B$와 $\Omega^1_{B/A}$의 위상적 직합이라는
사실을 고려하면 따른다 \sref{0.20.4.8-ko}.
\end{env}

\begin{proposition}[20.7.2]
\label{0.20.7.2-ko}
$u:A\to B$, $v:B\to C$를 위상환의 두 연속 준동형이라 하자. 연속
준동형
\[
  v_{C/B/A}:\Omega^1_{B/A}\otimes_B C\longrightarrow\Omega^1_{C/A}
\]
가 형식적으로 왼쪽 가역일 \sref{0.19.1.5-fr} 필요충분조건은 $C$가
$A$에 대하여 상대적으로 형식적으로 매끄러운 $B$-대수인 것이다
\sref{0.19.9.1-ko}. 특히 $C$가 형식적으로 매끄러운 $B$-대수이면
충분하다.
\end{proposition}

\begin{proof}
$\Omega^1_{C/A}$와 $\Omega^1_{B/A}\otimes_B C$의 위상은 $C$의
위상으로부터 얻어지는 위상보다 섬세하지 않으므로
\sref{0.20.4.5-ko}, $v_{C/B/A}$가 형식적으로 왼쪽 가역이라는 것은
$C$의 열린 아이디얼에 의하여 소멸되는 모든 이산 $C$-가군 $L$에
대하여 표준 준동형
\[
  \operatorname{Hom.\ cont}_C(\Omega^1_{C/A},L)\longrightarrow
  \operatorname{Hom.\ cont}_C(\Omega^1_{B/A}\otimes_B C,L)
\]
\oldpage[0\textsubscript{IV-1}]{148}
이 전사라는 뜻이다 \sref{0.19.1.5-fr}. 텐서곱 위상의 정의에 의하여
\[
  \operatorname{Hom.\ cont}_C(\Omega^1_{B/A}\otimes_B C,L)
  =\operatorname{Hom.\ cont}_B(\Omega^1_{B/A},L)
\]
이므로, \sref{0.20.4.8-ko}에 의하여 이는 표준 준동형
\[
  \operatorname{D\acute{e}r.\ cont}_A(C,L)\longrightarrow
  \operatorname{D\acute{e}r.\ cont}_A(B,L)
\]
이 전사라는 것과 같다. 그런데 완전열 \sref{0.20.3.8.2-ko}은 이 조건이
$\operatorname{Exalcotop}_{B/A}(C,L)=0$과 동치임을 보인다. 이는 정확히
$C$가 $A$에 대하여 상대적으로 형식적으로 매끄럽다는 뜻이다
\sref{0.19.9.8-ko}.
\end{proof}

\begin{corollary}[20.7.3]
\label{0.20.7.3-ko}
$B$와 $C$에서 모든 열린 아이디얼의 제곱이 열려 있다고 하자. $C$가
$A$에 대하여 상대적으로 형식적으로 매끄러운 $B$-대수일 필요충분조건은
다음과 같다. $C$의 아이디얼들로 이루어진 $0$의 근방 기본계를
$(\mathfrak R_\lambda)$라 나타낼 때, 모든 $\lambda$에 대하여 준동형
\begin{equation}
\tag{20.7.3.1}
\label{0.20.7.3.1-ko}
  v_{C/B/A}\otimes 1_{C/\mathfrak R_\lambda}:
  \Omega^1_{B/A}\otimes_B(C/\mathfrak R_\lambda)\longrightarrow
  \Omega^1_{C/A}\otimes_C(C/\mathfrak R_\lambda)
\end{equation}
이 왼쪽 가역이다.
\end{corollary}

\begin{proof}
실제로 이때 $\Omega^1_{B/A}$의 위상(각각 $\Omega^1_{C/A}$의 위상)은
$B$의 위상(각각 $C$의 위상)으로부터 유도된다
\sref{0.20.4.5-ko}. 따라서 $\Omega^1_{B/A}\otimes_B C$의 위상도
$C$의 위상으로부터 유도된다. 이제 따름정리는 \sref{0.20.7.2-ko}와
\sref{0.19.1.7-fr}에서 따른다.
\end{proof}

\begin{proposition}[20.7.4]
\label{0.20.7.4-ko}
$A$를 위상환, $B$를 위상 $A$-대수라 하자. $B$가 형식적으로
비분기인 $A$-대수일 필요충분조건은 분리 완비화
$\widehat{\Omega}^1_{B/A}$가 $0$인 것이다.
\end{proposition}

\begin{proof}
실제로 \sref{0.19.10.4-ko}와 \sref{0.20.1.1-ko}에서 곧 다음을
얻는다. $B$가 형식적으로 비분기일 필요충분조건은 $B$의 모든 열린
아이디얼 $\mathfrak R$, $\mathfrak HB\subset\mathfrak R$인 $A$의 모든
열린 아이디얼 $\mathfrak H$, 그리고 모든 $(B/\mathfrak R)$-가군 $L$에
대하여
\[
  \operatorname{D\acute{e}r}_{A/\mathfrak H}
  (B/\mathfrak R,L)=0
\]
인 것이다. 이는 열린 아이디얼에 의하여 소멸되는 모든 이산
$B$-가군 $L$에 대하여
$\operatorname{D\acute{e}r.\ cont}_A(B,L)=0$인 것과 같다.
\sref{0.20.4.8-ko}에 의하여 이는 그러한 모든 $B$-가군에 대하여
$\operatorname{Hom.\ cont}_B(\Omega^1_{B/A},L)=0$인 것과 동치이고,
여기서 명제가 곧 따른다.

$B$가 \emph{이산}일 때 \sref{0.20.7.4-ko}의 조건은
$\Omega^1_{B/A}=0$과 동치이다.
\end{proof}

\begin{corollary}[20.7.5]
\label{0.20.7.5-ko}
$A$를 환, $\mathfrak m$을 $A$의 아이디얼, $B$를 $A$-대수라 하자.
$A$에는 $\mathfrak m$-준아딕 위상을, $B$에는 $A$의 위상에서 유도된
위상을 주고
\[
  A_0=A/\mathfrak m,\qquad
  B_0=B/\mathfrak mB=B\otimes_A A_0
\]
라 두자. 그러면 $B$가 형식적으로 비분기인 $A$-대수일 필요충분조건은
$\Omega^1_{B_0/A_0}=0$인 것이다. 이는 $B_0$가 형식적으로 비분기인
$A_0$-대수라는 것과도 같다.
\end{corollary}

\begin{proof}
\sref{0.20.5.5-ko}에 의하여
\[
  \Omega^1_{B/A}\otimes_B B_0
  =\Omega^1_{B/A}/\mathfrak m\Omega^1_{B/A}
  =\Omega^1_{B_0/A_0}.
\]
$\Omega^1_{B/A}$의 모든 열린 부분가군이 $\Omega^1_{B/A}$와 같다는
것은 \sref{0.20.4.5-ko}에 의하여
$\mathfrak m\Omega^1_{B/A}=\Omega^1_{B/A}$라고 하는 것과 동치이므로
결론이 따른다.
\end{proof}

\begin{proposition}[20.7.6]
\label{0.20.7.6-ko}
$u:A\to B$, $v:B\to C$를 위상환의 두 연속 준동형이라 하고, $v$에
의하여 $C$가 형식적으로 에탈인 $B$-대수라고 가정하자. 그러면
\[
  v_{C/B/A}:\Omega^1_{B/A}\otimes_B C\longrightarrow\Omega^1_{C/A}
\]
는 형식적 전단사 준동형이다 \sref{0.19.1.2-fr}.
\end{proposition}

\begin{proof}
$C$의 열린 아이디얼에 의하여 소멸되는 모든 이산 $C$-가군 $L$에
대하여 \sref{0.20.7.4-ko}에 의해
$\operatorname{D\acute{e}r.\ cont}_B(C,L)=0$이다. 따라서
\sref{0.20.3.6.1-ko}에 의해 표준 준동형
\oldpage[0\textsubscript{IV-1}]{149}
\[
  \operatorname{D\acute{e}r.\ cont}_A(C,L)\longrightarrow
  \operatorname{D\acute{e}r.\ cont}_A(B,L)
\]
은 단사이다. 한편 \sref{0.20.7.2-ko}에 의하여 전사이기도 하므로
전단사이다. 따라서 $v_{C/B/A}$의 상은 $\Omega^1_{C/A}$에서 조밀하다.
그렇지 않으면 $\Omega^1_{C/A}$를 $\operatorname{Im}(v_{C/B/A})$의
폐포로 나눈 몫은 분리이고 $0$이 아니므로 $0$이 아닌 이산 몫
$L$을 가질 것이며, 이는 방금 얻은 결과에 모순된다
\sref{0.20.4.8-ko}. 그러므로 \sref{0.20.7.2-ko}에 의하여 형식적 단사
준동형인 $v_{C/B/A}$는 형식적 전사 준동형이기도 하다
\sref{0.19.1.2-fr}. 따라서 이는 형식적 전단사 준동형이다.
\end{proof}

\begin{corollary}[20.7.7]
\label{0.20.7.7-ko}
$B$와 $C$에서 모든 열린 아이디얼의 제곱이 열려 있다고 하자. $C$가
형식적으로 에탈인 $B$-대수이면, $C$의 모든 열린 아이디얼
$\mathfrak R_\lambda$에 대하여 준동형 \sref{0.20.7.3.1-ko}은
전단사이다.
\end{corollary}

\begin{proof}
실제로 \sref{0.19.1.1-fr}에 의하여 이 준동형은 전사이고,
\sref{0.20.7.3-ko}에 의하여 단사이다.
\end{proof}

\begin{proposition}[20.7.8]
\label{0.20.7.8-ko}
$u:A\to B$를 위상환의 연속 준동형, $\mathfrak R$을 $B$의 아이디얼,
$C$를 위상 몫환 $B/\mathfrak R$, $v:B\to C$를 표준 준동형이라 하자.
그러면 다음이 성립한다.
\begin{enumerate}
  \item[\rm(i)] 완전열 \sref{0.20.5.12.1-ko}
\[
  \mathfrak R/\mathfrak R^2\xrightarrow{\delta_{C/B/A}}
  \Omega^1_{B/A}\otimes_B C\xrightarrow{v_{C/B/A}}
  \Omega^1_{C/A}\to0
\]
  에서 준동형 $\delta_{C/B/A}$는 연속이고 $v_{C/B/A}$는 위상
  $C$-가군의 엄밀 사상이다.
  \item[\rm(ii)] $\delta_{C/B/A}$가 형식적으로 왼쪽 가역일
  \sref{0.19.1.5-fr} 필요충분조건은 $C$의 열린 아이디얼에 의하여
  소멸되는 모든 이산 $C$-가군 $L$에 대하여 표준 준동형
\begin{equation}
\tag{20.7.8.1}
\label{0.20.7.8.1-ko}
  \operatorname{Exalcotop}_A(C,L)\longrightarrow
  \operatorname{Exalcotop}_A(B,L)
\end{equation}
  이 단사인 것이다.
\end{enumerate}
\end{proposition}

\begin{proof}
(i) 첫째 주장은 명백하다. 한편 텐서곱 위상의 정의에 의하여 표준
준동형
\[
  v\otimes v:B\otimes_A B\longrightarrow C\otimes_A C
\]
는 엄밀 사상이고, 따라서 $v_{C/B/A}$도 그러함이 곧 따른다
(cf. \sref{0.20.5.2-ko}).

(ii) $\delta_{C/B/A}$가 형식적으로 왼쪽 가역이라는 것은 $C$의 열린
아이디얼에 의하여 소멸되는 모든 이산 $C$-가군 $L$에 대하여 표준
준동형
\[
  \operatorname{Hom.\ cont}_C(\Omega^1_{B/A}\otimes_B C,L)
  \longrightarrow
  \operatorname{Hom.\ cont}_C(\mathfrak R/\mathfrak R^2,L)
\]
이 전사라는 뜻이다. \sref{0.18.4.3-fr}과
\sref{0.20.4.8-ko}을 고려하면 이는 표준 준동형
\[
  \operatorname{D\acute{e}r.\ cont}_A(B,L)\longrightarrow
  \operatorname{Exalcotop}_B(C,L)
\]
이 전사라는 것과 같다. 따라서 완전열 \sref{0.20.3.6.1-ko}에서 결론이
따른다.

\sref{0.18.4.1-fr}의 양변의 정의를 고려하면
\sref{0.20.7.8.1-ko}의 단사성은 다음과 같이도 표현된다. $A$의 모든
열린 아이디얼 $\mathfrak M$, $\mathfrak MB\subset\mathfrak N$인
$B$의 모든 열린 아이디얼 $\mathfrak N$, 그리고
$B/(\mathfrak R+\mathfrak N)$의
$(B/(\mathfrak R+\mathfrak N))$-가군 $L$에 의한 모든
$(A/\mathfrak M)$-확장 $E$를 생각하자. 표준 준동형
\[
  B/\mathfrak N\longrightarrow B/(\mathfrak R+\mathfrak N)
\]
에 의한 $E$의 역상이 $(A/\mathfrak M)$-자명하다고 가정하자. 그러면
$A$의 열린 아이디얼 $\mathfrak M'\subset\mathfrak M$과
\oldpage[0\textsubscript{IV-1}]{150}
$B$의 열린 아이디얼 $\mathfrak N'\subset\mathfrak N$이 존재하여
$\mathfrak M'B\subset\mathfrak N'$이고, 표준 준동형
\[
  B/(\mathfrak R+\mathfrak N')\longrightarrow
  B/(\mathfrak R+\mathfrak N)
\]
에 의한 $E$의 역상은 $(A/\mathfrak M')$-자명하다.
\end{proof}

\begin{corollary}[20.7.9]
\label{0.20.7.9-ko}
위상 $A$-대수 $C=B/\mathfrak R$이 형식적으로 매끄러우면 표준 준동형
$\delta_{C/B/A}$는 형식적으로 왼쪽 가역이다.
\end{corollary}

\begin{env}[20.7.10]
\label{0.20.7.10-ko}
\sref{0.20.6.1-ko}에서 $A$, $B$, $C$가 위상환이고 $u$, $v$가 연속
준동형일 때 $\Upsilon_{C/B/A}$에는
$\Omega^1_{B/A}\otimes_B C$의 위상에서 유도된 위상을 준다. 이때 도식
\sref{0.20.6.4.1-ko}의 준동형들이 연속이면 준동형
\sref{0.20.6.4.2-ko}와 \sref{0.20.6.4.4-ko}도 연속이다.

더욱이 \sref{0.20.6.7-ko}에서 $L$이 $C$의 열린 아이디얼에 의하여
소멸되는 이산 $C$-가군이라고 가정하면, 귀납극한으로 넘어가서 표준
$C$-준동형
\begin{equation}
\tag{20.7.10.1}
\label{0.20.7.10.1-ko}
  \operatorname{Exalcotop}_{B/A}(C,L)\longrightarrow
  \operatorname{Hom.\ cont}_C(\Upsilon_{C/B/A},L)
\end{equation}
을 얻는다. 실제로 \sref{0.18.5.3.1-fr}을 고려하자. $A$의 모든 열린
아이디얼 $\mathfrak M$, $\mathfrak MB\subset\mathfrak N$인 $B$의 모든
열린 아이디얼 $\mathfrak N$, $\mathfrak NC\subset\mathfrak P$이고
$\mathfrak P L=0$인 $C$의 모든 열린 아이디얼 $\mathfrak P$에 대하여,
$L$에 의한 $C/\mathfrak P$의 $(B/\mathfrak N)$-확장 $E$가
$(A/\mathfrak M)$-자명하면 이는 연속 $C$-준동형
\[
  f:\Omega^1_{B/A}\otimes_B C\longrightarrow L
\]
의 데이터에서 나온다. 준동형 \sref{0.20.7.10.1-ko}은 $E$의 류의
$\operatorname{Exalcotop}_{B/A}(C,L)$에서의 상에 $f$를
$\Upsilon_{C/B/A}$에 제한한 준동형을 대응시킨다. 이를 $E$의
\emph{특성 준동형}이라 하며 여전히 $\chi_E$로 나타낸다.
\end{env}

\begin{proposition}[20.7.11]
\label{0.20.7.11-ko}
$C$의 모든 열린 아이디얼의 제곱이 열려 있다고 하자. $C$가
형식적으로 매끄러운 위상 $A$-대수이고 $\Omega^1_{C/B}$가 형식적으로
사영인 $C$-가군이면, $C$의 열린 아이디얼에 의하여 소멸되는 모든 이산
$C$-가군 $L$에 대하여 표준 동형사상
\begin{equation}
\tag{20.7.11.1}
\label{0.20.7.11.1-ko}
  \operatorname{Exalcotop}_B(C,L)\xrightarrow{\sim}
  \operatorname{Hom.\ cont}_C(\Upsilon_{C/B/A},L)
\end{equation}
이 존재한다.
\end{proposition}

\begin{proof}
실제로 완전열 \sref{0.20.3.7.1-ko}에서
$\operatorname{Exalcotop}_A(C,L)=0$이다 \sref{0.19.4.4-fr}. 따라서
\[
  \operatorname{Exalcotop}_B(C,L)
  =\operatorname{Exalcotop}_{B/A}(C,L),
\]
이고 준동형 \sref{0.20.7.8.1-ko}은 바로
\sref{0.20.7.10.1-ko}이다. 이것이 전단사임은
\sref{0.20.6.13-ko}에서 귀납극한으로 넘어가면 얻는다. 여기서는
$\Omega^1_{C/B}$의 위상이 $C$의 위상에서 유도된다는 것
\sref{0.20.4.5-ko}과 \sref{0.19.2.4-fr}을 고려한다.
\end{proof}

\begin{env}[20.7.12]
\label{0.20.7.12-ko}
\sref{0.20.6.14-ko}에서도 $\Upsilon^C_{B/A/\Lambda}$에
$\Omega^1_{A/\Lambda}\otimes_A C$의 위상에서 유도된 위상을 준다.
해당 환들이 위상환이고 환 준동형들이 연속이면 준동형
\sref{0.20.6.14.6-ko}은 연속이다.
\end{env}

\begin{env}[20.7.13]
\label{0.20.7.13-ko}
\sref{0.20.6.23-ko}에서 환 $\Lambda$, $A$, $B$, $C$가 위상환이고
준동형 $s$, $u$, $v$가 연속이며 $L$이 $C$의 열린 아이디얼에 의하여
소멸되는 이산 $C$-가군이라고 하자. 그러면 \sref{0.20.3.6-ko}에서와
같이 귀납극한으로 넘어가서 완전열
\begin{equation}
\tag{20.7.13.1}
\label{0.20.7.13.1-ko}
\begin{split}
  0\to\operatorname{D\acute{e}r.\ cont}_B(C,L)\to
  \operatorname{D\acute{e}r.\ cont}_A(C,L)\to
  \operatorname{D\acute{e}r.\ cont}_A(B,L)\to \\
  \to\operatorname{Exalcotop}_{B/\Lambda}(C,L)\to
  \operatorname{Exalcotop}_{A/\Lambda}(C,L)\to
  \operatorname{Exalcotop}_{A/\Lambda}(B,L)\to0
\end{split}
\end{equation}
을 얻는다.
\end{env}

\begin{env}[20.7.14]
\label{0.20.7.14-ko}
$A$를 위상환, $B$를 위상 $A$-대수라 하고, $\mathfrak M'$을 $A$의
열린 아이디얼, $\mathfrak N'$을 $B$의 열린 아이디얼이라 하되
$\mathfrak M'B\subset\mathfrak N'$이라 하자. 다음과 같이 놓는다.
\[
  A'=A/\mathfrak M',
\]
\oldpage[0\textsubscript{IV-1}]{151}
\[
  B'=B/\mathfrak N'.
\]
준동형 $B\otimes_A B\to B'\otimes_{A'}B'$의 핵은
\[
  \mathfrak U'
  =\operatorname{Im}(\mathfrak N'\otimes B+B\otimes\mathfrak N')
\]
이므로, 준동형
\begin{equation}
\tag{20.7.14.1}
\label{0.20.7.14.1-ko}
  \varphi_{(\mathfrak M',\mathfrak N')}:
  \Omega^1_{B/A}\longrightarrow\Omega^1_{B'/A'}
\end{equation}
의 핵이
\[
  ((\mathfrak J\cap\mathfrak U')+\mathfrak J^2)/\mathfrak J^2
\]
임을 곧바로 얻는다. 한편 준동형 \sref{0.20.7.14.1-ko}은
\sref{0.20.4.7-ko}에 의하여 전사이다.

$\mathfrak M''$을 $A$의 또 다른 열린 아이디얼, $\mathfrak N''$을
$B$의 또 다른 열린 아이디얼이라 하되
\[
  \mathfrak M''\subset\mathfrak M',\qquad
  \mathfrak N''\subset\mathfrak N',\qquad
  \mathfrak M''B\subset\mathfrak N''
\]
이라 하고 $A''=A/\mathfrak M''$, $B''=B/\mathfrak N''$이라 놓으면
마찬가지로 전사 준동형
\[
  \varphi_{(\mathfrak M'',\mathfrak N''),(\mathfrak M',\mathfrak N')}:
  \Omega^1_{B''/A''}\longrightarrow\Omega^1_{B'/A'}
\]
을 얻으며, 이 준동형들은 명백히 \emph{사영계}를 이룬다.
또
\[
  ((\mathfrak J\cap\mathfrak U')+\mathfrak J^2)/\mathfrak J^2
\]
들이 $\Omega^1_{B/A}$에서 $0$의 근방 기본계를 이룸을 주목하면, 위상
$B$-가군 $\Omega^1_{B/A}$의 \emph{분리 완비화}
$\widehat{\Omega}^1_{B/A}$는 표준 동형사상의 의미에서
\begin{equation}
\tag{20.7.14.2}
\label{0.20.7.14.2-ko}
  \widehat{\Omega}^1_{B/A}
  =\varprojlim(\Omega^1_{B'/A'})
\end{equation}
로 주어진다.

더욱이 표준 준동형
$j:\Omega^1_{B/A}\to\widehat{\Omega}^1_{B/A}$는
$\varphi_{(\mathfrak M',\mathfrak N')}$들의 사영계의 사영극한이다.
따라서 $\varphi_{(\mathfrak M',\mathfrak N')}$는
\[
  \Omega^1_{B/A}\xrightarrow{j}\widehat{\Omega}^1_{B/A}
  \longrightarrow\Omega^1_{B'/A'}
\]
로 인수분해된다. $\varphi_{(\mathfrak M',\mathfrak N')}$가 전사이므로 모든 쌍
$(\mathfrak M',\mathfrak N')$에 대하여 표준 준동형
\begin{equation}
\tag{20.7.14.3}
\label{0.20.7.14.3-ko}
  \widehat{\Omega}^1_{B/A}\longrightarrow\Omega^1_{B'/A'}
\end{equation}
도 전사이다.

위의 모든 곳에서 $A'$을 $A$로 바꾸어도 된다.

마지막으로 $L$이 \emph{분리되고 완비인} 위상 $B$-가군이면,
$\Omega^1_{B/A}$에서 $L$로 가는 모든 연속 $B$-준동형은
$\widehat{\Omega}^1_{B/A}$에서 $L$로 가는 연속 $\widehat B$-준동형으로
유일하게 연장된다. 역으로 후자의 준동형을
$\Omega^1_{B/A}\to\widehat{\Omega}^1_{B/A}$와 합성하면 연속
$B$-준동형을 얻는다. 따라서 표준 동형사상
\[
  \operatorname{Hom.\ cont}_B(\Omega^1_{B/A},L)
  \xrightarrow{\sim}
  \operatorname{Hom.\ cont}_{\widehat B}
    (\widehat{\Omega}^1_{B/A},L)
\]
이 존재한다. 특히 $L$이 $B$의 열린 아이디얼에 의하여 소멸되는 이산
$B$-가군이면 표준 동형사상 \sref{0.20.4.8.2-ko}은
\begin{equation}
\tag{20.7.14.4}
\label{0.20.7.14.4-ko}
  \operatorname{Hom.\ cont}_{\widehat B}
    (\widehat{\Omega}^1_{B/A},L)
  \xrightarrow{\sim}\operatorname{D\acute{e}r.\ cont}_A(B,L)
\end{equation}
로도 쓸 수 있다.
\end{env}

\begin{proposition}[20.7.15]
\label{0.20.7.15-ko}
$A$를 이산환, $B$를 아딕 위상 $A$-대수
($\mathbf{0}_{\mathrm I}$, 7.1.9), $\mathfrak n$을 $B$의 정의
아이디얼이라 하고 $B_0=B/\mathfrak n$이라 하자.
$\Omega^1_{B_0/A}$와 $\mathfrak n/\mathfrak n^2$이 유한형
$B_0$-가군이라고 가정하자. 그러면 $\widehat{\Omega}^1_{B/A}$는
유한형 $\widehat B$-가군이다.
\end{proposition}

\begin{proof}
$B$의 모든 열린 아이디얼의 제곱이 열려 있으므로
$\Omega^1_{B/A}$의 위상은 $\mathfrak n$-준아딕 위상이다
\sref{0.20.4.5-ko}. $B$가 아딕환이라는 가정을 고려하면
($\mathbf{0}_{\mathrm I}$, 7.2.7과 7.2.9)에 의하여 다음을 보이면
충분하다.
\oldpage[0\textsubscript{IV-1}]{152}
\[
  \Omega^1_{B/A}/\mathfrak n\Omega^1_{B/A}
  =\Omega^1_{B/A}\otimes_B B_0
\]
이 유한형 $B_0$-가군이다. 이는 가정과 완전열
\sref{0.20.5.12.1-ko}
\[
  \mathfrak n/\mathfrak n^2\longrightarrow
  \Omega^1_{B/A}\otimes_B B_0\longrightarrow
  \Omega^1_{B_0/A}\longrightarrow0
\]
에서 따른다.
\end{proof}

\begin{env}[20.7.16]
\label{0.20.7.16-ko}
예를 들어 $A$가 체 $k$이고, $B=k'[[T_1,\ldots,T_n]]$이 통상적인
위상을 갖춘 형식적 멱급수 대수이며, $k'$이 $k$의 \emph{유한}
확대체이면 명제 \sref{0.20.7.15-ko}를 적용할 수 있다
\sref{0.21.9.2-fr}. $B$의 분수체 $K$의 $k$ 위 초월차수는
무한이다. 따라서 $k$의 표수가 $0$이면, 특히
\sref{0.20.4.13-ko}(i)와 \sref{0.20.5.9-ko}을 사용하고 표수 $0$인
체의 모든 미분은 임의의 확대체로 연장된다는 사실도 사용하여
$\Omega^1_{B/k}$가 \emph{유한형 $B$-가군이 아님}을 곧바로 얻는다.
\end{env}

\begin{env}[20.7.17]
\label{0.20.7.17-ko}
$A$, $B$, $C$를 위상환, $u:A\to B$, $v:B\to C$를 연속
준동형이라 하자. $A$, $B$, $C$를 각각 열린 아이디얼에 의한 몫
\[
  A'=A/\mathfrak M',\qquad
  B'=B/\mathfrak N',\qquad
  C'=C/\mathfrak P'
\]
으로 바꾸되
\[
  u(\mathfrak M')\subset\mathfrak N',\qquad
  v(\mathfrak N')\subset\mathfrak P'
\]
이라 하자. 그러면 준동형 $u':A'\to B'$, $v':B'\to C'$을 얻고,
따라서 표준 준동형 $u'_{C'/B'/A'}$, $v'_{C'/B'/A'}$을 얻는다.
\sref{0.20.5.4-ko}에 의하면 이들은 사영계들을 이루므로, 극한으로
넘어가서 완전열 \sref{0.20.5.7.1-ko}의 준동형들
\[
  \Omega^1_{B/A}\otimes_B C
  \xrightarrow{v_{C/B/A}}\Omega^1_{C/A}
  \xrightarrow{u_{C/B/A}}\Omega^1_{C/B}\longrightarrow0
\]
을 분리 완비화들로 연장하는 표준 준동형
\begin{align}
  \widehat v_{C/B/A}:
  \widehat{\Omega}^1_{B/A}
  \widehat\otimes_{\widehat B}\widehat C
  &\longrightarrow\widehat{\Omega}^1_{C/A},
  \tag{20.7.17.1}\label{0.20.7.17.1-ko}\\
  \widehat u_{C/B/A}:
  \widehat{\Omega}^1_{C/A}
  &\longrightarrow\widehat{\Omega}^1_{C/B}.
  \tag{20.7.17.2}\label{0.20.7.17.2-ko}
\end{align}
을 얻는다. 열
\begin{equation}
\tag{20.7.17.3}
\label{0.20.7.17.3-ko}
  \widehat{\Omega}^1_{B/A}
  \widehat\otimes_{\widehat B}\widehat C
  \xrightarrow{\widehat v_{C/B/A}}
  \widehat{\Omega}^1_{C/A}
  \xrightarrow{\widehat u_{C/B/A}}
  \widehat{\Omega}^1_{C/B}\longrightarrow0
\end{equation}
에서 연속하는 두 준동형의 합성은 $0$이지만, 이 열은
\emph{반드시 완전인 것은 아니다}. 다만 $B$와 $C$가
\emph{거리화 가능}이면 $\widehat u_{C/B/A}$는 \emph{전사}이고
$\operatorname{Im}(\widehat v_{C/B/A})$는
$\operatorname{Ker}(\widehat u_{C/B/A})$에서 \emph{조밀}하다.
실제로 ($\mathbf{0}_{\mathrm I}$, 7.3.1), $B$와 $C$의 아이디얼로
이루어진 감소열 $(\mathfrak M_k)$와 $(\mathfrak N_k)$가 각각
$0$의 근방 기본계를 이루고
\[
  B_k=B/\mathfrak M_k,\qquad C_k=C/\mathfrak N_k
\]
라 놓으면 전이 준동형
\[
  \Omega^1_{C_{k+1}/B_{k+1}}\longrightarrow\Omega^1_{C_k/B_k},
  \qquad
  \Omega^1_{C_{k+1}/A}\longrightarrow\Omega^1_{C_k/A},
  \qquad
  \Omega^1_{B_{k+1}/A}\longrightarrow\Omega^1_{B_k/A}
\]
이 모두 전사라는 사실에서 곧바로 따른다 \sref{0.20.7.14-ko}.
\end{env}

\begin{proposition}[20.7.18]
\label{0.20.7.18-ko}
$A$, $B$, $C$를 위상환, $u:A\to B$, $v:B\to C$를 연속
준동형이라 하자. $B$와 $C$가 허용가능하고
($\mathbf{0}_{\mathrm I}$, 7.1.2) 거리화 가능하다고 가정하자.
표준 준동형
\[
  \widehat v_{C/B/A}:
  \widehat{\Omega}^1_{B/A}
  \widehat\otimes_{\widehat B}\widehat C
  \longrightarrow\widehat{\Omega}^1_{C/A}
\]
가 연속 $C$-준동형인 왼쪽 역을 가질 필요충분조건은 $C$가 $A$에
대하여 상대적으로 형식적으로 매끄러운 $B$-대수인 것이다.
\end{proposition}

\begin{proof}
필요성은 \sref{0.20.7.2-ko}와 \sref{0.19.1.6-fr}에서 따른다.
충분성을 보이기 위하여 위상 $C$-가군
\[
  L=\widehat{\Omega}^1_{B/A}
  \widehat\otimes_{\widehat B}\widehat C
\]
이 가정에 의하여 거리화 가능하고
\oldpage[0\textsubscript{IV-1}]{153}
완비임을 주목하자. 따라서 곱위상을 갖춘 $E=D_C(L)$도 거리화
가능하고 완비이다. 또한 $E$는 허용가능하다. 실제로
$\mathfrak R$이 $C$의 정의 아이디얼이면
\[
  (\mathfrak R\oplus L)^n
  =\mathfrak R^n\oplus\mathfrak R^{n-1}L
\]
은 $0$으로 수렴한다. 합성사상
\[
  D:B\xrightarrow{d_{B/A}}\Omega^1_{B/A}\longrightarrow
  \widehat{\Omega}^1_{B/A}
  \widehat\otimes_{\widehat B}\widehat C=L
\]
은 $B$에서 $L$로 가는 연속 $A$-미분이므로, 연속 $A$-준동형
\[
  f:x\longmapsto(v(x),D(x))
\]
는 $E$에 $B$-확장 구조를 정한다. $L$은 $E$의 닫힌 아이디얼이므로
\sref{0.19.9.5-fr}와 가정에 의하여 항등사상 $C\to E/L$은
\[
  C\xrightarrow{g}E\longrightarrow E/L
\]
로 인수분해된다. 여기서 $g$는 $g\circ v=f$를 만족하는 연속
준동형이다. 따라서 \sref{0.20.1.3-ko}에 의하여 $g$는
$y\mapsto(y,D'(y))$의 꼴이며, $D'$은 $C$에서 $L$로 가는 연속
$B$-미분이다. 다시 말해 $D'\circ v=D$이다.
\sref{0.20.7.14.4-ko}을 고려하면
\[
  D'=h\circ\widehat d_{C/A}
\]
이며, 여기서
$h:\widehat{\Omega}^1_{C/A}\to L$은 연속 $C$-준동형이다. 한편 정의에
의하여
\[
  \widehat d_{C/A}\circ v=\widehat v_{C/B/A}\circ D
\]
이고, $D$에 의한 $B$의 상은 $C$-가군 $L$을 위상적으로 생성하므로
\sref{0.20.4.7-ko}, 등식 $D'\circ v=D$에서
\[
  h\circ\widehat v_{C/B/A}=1_L
\]
을 얻는다.
\end{proof}

\begin{corollary}[20.7.19]
\label{0.20.7.19-ko}
\sref{0.20.7.18-ko}의 가정 아래 $B$와 $C$에서 모든 열린 아이디얼의
제곱도 열린다고 더 가정하자. 이때 $C$가 $A$에 대하여 상대적으로
형식적으로 매끄러운 $B$-대수일 필요충분조건은
$\widehat v_{C/B/A}$가 왼쪽 가역인 것이다.
\end{corollary}

\begin{proof}
실제로
$\widehat{\Omega}^1_{B/A}\widehat\otimes_{\widehat B}\widehat C$와
$\widehat{\Omega}^1_{C/A}$의 위상은 이 경우 $C$의 위상에서 유도되며
\sref{0.20.4.5-ko}, 한쪽에서 다른 쪽으로 가는 모든 $C$-준동형은
반드시 연속이다.
\end{proof}

\begin{env}[20.7.20]
\label{0.20.7.20-ko}
$A$를 위상환, $B$를 \emph{거리화 가능하고 완비인} 위상
$A$-대수, $\mathfrak R$을 $B$의 \emph{닫힌} 아이디얼이라 하고,
$C=B/\mathfrak R$을 거리화 가능하고 완비인 몫위상환이라 하자.
$(\mathfrak M_k)$를 아이디얼들로 이루어진 $B$에서 $0$의 감소하는
근방 기본계라 하고
\[
  B_k=B/\mathfrak M_k,\qquad
  \mathfrak R_k=(\mathfrak R+\mathfrak M_k)/\mathfrak M_k,\qquad
  C_k=B_k/\mathfrak R_k
\]
라 놓자. 준동형들의 사영계
\[
  \delta_{C_k/B_k/A}:
  \mathfrak R_k/\mathfrak R_k^2\longrightarrow
  \Omega^1_{B_k/A}\otimes_{B_k}C_k
\]
가 존재하므로 \sref{0.20.5.11.3-ko}, 극한으로 넘어가서 표준 준동형
\[
  \widehat\delta_{C/B/A}:
  \mathfrak R/\mathfrak R^2\longrightarrow
  \widehat{\Omega}^1_{B/A}
  \widehat\otimes_{\widehat B}\widehat C
\]
을 얻는다. 또한 \sref{0.20.7.17-ko}에서와 같이 추론하면 표준 준동형
\[
  \widehat v_{C/B/A}:
  \widehat{\Omega}^1_{B/A}
  \widehat\otimes_{\widehat B}\widehat C
  \longrightarrow\widehat{\Omega}^1_{C/A}
\]
는 \emph{전사}이고
$\operatorname{Im}(\widehat\delta_{C/B/A})$는
$\operatorname{Ker}(\widehat v_{C/B/A})$에서 \emph{조밀}하다.
\end{env}
\section{표수 $p$인 환에서의 미분}
\label{section:0.21-ko}
\label{0.21-ko}

이 절의 결과들은 \textsection19, \textsection20, \textsection22의 결과들보다
더 특수하고 기술적이며, 제IV장의 전개에서는 예외적으로만 사용된다.
여기서 주된 역할은 \textsection22의 세 정리 (22.3.3), (22.5.8),
(22.7.3)의 증명에 있다. 그 가운데 첫째와 셋째 정리는 제IV장
\textsection7의 뇌터 국소환의 ``정밀한'' 이론에서 본질적으로 쓰인다.
\oldpage[0\textsubscript{IV-1}]{154}

\subsection{$p$-생성원계와 $p$-기저}
\label{subsection:0.21.1-ko}
\label{0.21.1-ko}

\begin{env}[21.1.1]
\label{0.21.1.1-ko}
$p$가 $0$ 또는 소수라고 하자. 표수 $p$인 소체 $P$에서 환 $A$로 가는
환 준동형이 존재하면 $A$를 \emph{표수 $p$인 환}이라 한다. 이 준동형은
유일하다. 실제로 합성
\[
  \mathbf Z\longrightarrow P\longrightarrow A
\]
은 $\mathbf Z$에서 $A$로 가는 유일한 준동형이다. $A\neq0$이면 이는
$A$가 표수 $p$인 체를 포함한다는 것과 동치이다. 이 경우 $P$의 상은
$P$와 동형인 체이며, 실제로 $A$ 안에서 $P$와 동형인 유일한 부분체이다.
\end{env}

\begin{env}[21.1.2]
\label{0.21.1.2-ko}
$p>0$이면 $A$가 표수 $p$라는 것은 $A$에서 $p\cdot1=0$이라는 것,
또는 $pA=0$이라는 것과 동치이다. $p=0$이면 $A$가 표수 $p$라는 것은
모든 정수 $n\neq0$에 대하여 $n\cdot1$이 $A$에서 가역이라는 것과
동치이다. $A\neq0$이면 $A$가 표수 $p$가 되게 하는 $p$는 소수 또는
$0$ 가운데 많아야 하나이다. 이는 서로 다른 두 소수 $p,q$에 대한
베주 항등식 $ap+bq=1$에서 따른다. 반면 $0$환은 모든 $p$에 대하여
표수 $p$이다.
\end{env}

\begin{env}[21.1.3]
\label{0.21.1.3-ko}
$A$가 표수 $p$이면 모든 $A$-대수도 표수 $p$이다. 특히 $A$의 모든
소 아이디얼 $\mathfrak p$에 대하여 $\mathfrak p$에서의 $A$의 잉여체는
표수 $p$이다. 역으로 $p=0$이고 $A$의 모든 극대 아이디얼
$\mathfrak m$에서 잉여체의 표수가 $0$이면 $A$도 표수 $0$이다.
실제로 모든 정수 $n\neq0$에 대하여 $n\cdot1$은 모든
$A_{\mathfrak m}$에서 가역이고, 따라서 $A$에서도 가역이다. 반면
$p>0$일 때 국소환의 잉여체는 표수 $p$이지만 그 국소환 자체는 표수
$p$가 아닐 수 있다. 예를 들어 소 국소환
$\mathbf Z_{p\mathbf Z}$ \sref{0.19.8.3-ko}이나 $n\geq2$인 아르틴
국소환 $\mathbf Z/p^n\mathbf Z$는 체를 포함하지 않는다.

마지막으로 환이 축소환인 경우에도 서로 다른 소 아이디얼에서의
잉여체들은 서로 다른 표수를 가질 수 있다. $\mathbf Z$가 그 예이다.
\end{env}

\begin{env}[21.1.4]
\label{0.21.1.4-ko}
이 절의 나머지에서는 소수 $p$를 하나 고정하며, 달리 명시하지 않는 한
모든 환은 표수 $p$라고 가정한다. 그러한 환 $A$에 대하여 사상
$x\mapsto x^p$은 $A$의 자기준동형이며, 이를 $F_A$로 나타낸다.
$A$가 축소환이면 $F_A$는 단사이다. 다음과 같이 놓는다.
\[
  A^p=F_A(A).
\]
이는 $x\in A$인 $x^p$들로 이루어진 $A$의 부분환이며, $A$를 자연스럽게
$A^p$-대수로 볼 수 있다.

또한 $A$에서 $A$로 가는 준동형 $F_A:x\mapsto x^p$을 통하여 $A$를
$A$-대수로 볼 수 있다. 다시 말해, 이는 $x\in A$와 스칼라
$\lambda\in A$의 곱 $\lambda\cdot x$이 환 $A$에서의 곱
$\lambda^p x$인 대수 $A^{(p)}$이다. 모든 환 준동형 $u:A\to B$에
대하여 쌍 $(u,u)$는 대수들의 디-준동형
$A^{(p)}\to B^{(p)}$이다.

모든 $A$-가군 $E$에 대하여
\[
  E^{(p)}=E\otimes_A A^{(p)}
\]
라 놓는다. 여기서 $A^{(p)}$는 $(A,A)$-쌍가군으로 본다. 왼쪽
$A$-가군 구조는 방금 정의한 것이고, 오른쪽 $A$-가군 구조는 $A$의
곱셈으로 정의한다. 따라서 $x\in E$, $\alpha,\beta\in A$에 대하여
\[
  \alpha(x\otimes\beta)=x\otimes(\alpha\beta),\qquad
  (\alpha x)\otimes\beta=x\otimes\alpha^p\beta
\]
이다.
\end{env}
\oldpage[0\textsubscript{IV-1}]{155}
$x^{(p)}=x\otimes1$이라 놓으면
\[
  (\alpha x)^{(p)}=\alpha^p x^{(p)}
\]
이다. $F_A$가 단사이면 $A^{(p)}$를 부분환 $A^p$ 위의 대수로 본 환
$A$와 동일시할 수 있다.

\begin{proposition}[21.1.5]
\label{0.21.1.5-ko}
$A$, $B$를 환, $u:A\to B$를 준동형이라 하자.
\begin{enumerate}
  \item[\rm(i)] $B$에서 $B$-가군 $L$로 가는 모든 $A$-미분은
  $A[B^p]$에서 $0$이다. 따라서 이는 $A[B^p]$-미분이다.
  \item[\rm(ii)] $A[B^p]$의 모든 부분 $A$-대수 $A'$에 대하여,
  $j:A'\to B$가 표준 단사 준동형이면 표준 준동형
  \[
    j_{B/A'/A}:\Omega^1_{B/A}\longrightarrow\Omega^1_{B/A'}
  \]
  은 전단사이다.
  \item[\rm(iii)] 어떤 정수 $s\geq0$에 대하여
  $u(A)\supset B^{p^s}$라고 가정하자. 그러면 환 $B\otimes_A B$에서
  $\mathfrak J_{B/A}$는 멱영원 아이디얼이다.
\end{enumerate}
\end{proposition}

\begin{proof}
(i) \sref{0.20.1.1.1-ko}에서 귀납적으로, 모든 $A$-미분
$D:B\to L$에 대하여
\[
  D(x^n)=nx^{n-1}D(x)
\]
임을 얻으며, 특히 $D(x^p)=0$이다.

(ii) \sref{0.20.4.8-ko}을 고려하면 (ii)는 (i)를 달리 쓴 것에
불과하다. 실제로 (i)는 모든 $B$-가군 $L$에 대하여
\[
  \operatorname{Der}_{A[B^p]}(B,L)
  =\operatorname{Der}_A(B,L)
\]
이라고 쓸 수 있다.

(iii) 모든 $x\in B$에 대하여
\[
  (x\otimes1-1\otimes x)^{p^s}
  =x^{p^s}\otimes1-1\otimes x^{p^s}=0
\]
이다. 실제로 $x^{p^s}\in u(A)$이다. 원소들
$1\otimes x-x\otimes1$이 $\mathfrak J_{B/A}$를 생성하므로 결론을
얻는다 \sref{0.20.4.4-ko}.
\end{proof}

\sref{0.21.1.5-ko}에서 곧바로 다음을 얻는다. 환 준동형들의 모든 쌍
$A\to B\to C$와 $B$ 안의 모든 부분 $A$-대수
$A'\subset A[B^p]$에 대하여
\begin{equation}
\tag{21.1.5.1}
\label{0.21.1.5.1-ko}
  \Upsilon_{C/B/A}=\Upsilon_{C/B/A'}
\end{equation}
이다. 또한 \sref{0.21.1.5-ko}에 의하여 ``절대'' 미분 가군에 대해서는
특히
\begin{equation}
\tag{21.1.5.2}
\label{0.21.1.5.2-ko}
  \Omega^1_A=\Omega^1_{A/A^p}
\end{equation}
이다.

\begin{corollary}[21.1.6]
\label{0.21.1.6-ko}
$B$가 유한형 $A$-대수이고 어떤 정수 $s\geq0$에 대하여
$u(A)\supset B^{p^s}$라고 가정하자. $\Omega^1_{B/A}=0$이면 $u$는
전사이다.
\end{corollary}

\begin{proof}
\oldpage[0\textsubscript{IV-1}]{156}
\sref{0.21.1.5-ko}(iii)에 의하여 $\mathfrak J_{B/A}$는 멱영원
아이디얼이다. 또한 \sref{0.20.4.4-ko}과 가정에 의하여
$\mathfrak J_{B/A}$는 유한형 아이디얼이므로 멱영 아이디얼이다.
관계식 $\Omega^1_{B/A}=0$은
$\mathfrak J_{B/A}=\mathfrak J_{B/A}^2$이라는 뜻이므로
$\mathfrak J_{B/A}=0$이다. 다시 말해
\[
  \pi:B\otimes_A B\longrightarrow B
\]
는 전단사이다. 한편 $B$의 모든 원소의 $p^s$제곱은 $u(A)$에
속하므로 $B$는 $A$ 위에서 정수적이다. 또한 유한형이므로 $B$는 유한
$A$-대수이다. 따라서 다음 보조정리를 증명하면 충분하다. 이
보조정리에서는 환들이 표수 $p$라고 가정하지 않는다.
\end{proof}

\begin{lemma}[21.1.6.1]
\label{0.21.1.6.1-ko}
$R$을 환, $S$를 유한 $R$-대수라 하자. 표준 준동형
\[
  \pi:S\otimes_R S\longrightarrow S
\]
가 전단사이면 구조 준동형 $u:R\to S$는 전사이다.
\end{lemma}

\begin{proof}
$R$의 모든 극대 아이디얼 $\mathfrak m$에 대하여 $T=R-\mathfrak m$이라
놓을 때 준동형
\[
  u_{\mathfrak m}:R_{\mathfrak m}\longrightarrow T^{-1}S
\]
가 전사임을 보이면 충분하다(Bourbaki, \emph{Alg. comm.}, chap.~II,
§~3, n$^\circ$~3, th.~1). 그런데 가정에 의하여 준동형
\[
  (T^{-1}S)\otimes_{R_{\mathfrak m}}(S/\mathfrak mS)
  \longrightarrow T^{-1}S
\]
는 전단사이고 \sref{0.1.3.4-ko}, $T^{-1}S$는 유한
$R_{\mathfrak m}$-대수이므로 $R$이 국소환인 경우로 환원할 수 있다.

그 극대 아이디얼을 다시 $\mathfrak m$이라 하자. $S$가 유한 생성
$R$-가군이므로 나카야마 보조정리에 의하여
\[
  u\otimes1:R/\mathfrak m\longrightarrow S/\mathfrak mS
\]
이 전사임을 증명하면 충분하다. 표준 준동형
\[
  (S/\mathfrak mS)\otimes_{R/\mathfrak m}(S/\mathfrak mS)
  \longrightarrow S/\mathfrak mS
\]
가 전단사이고 $S/\mathfrak mS$는 유한 $(R/\mathfrak m)$-대수이므로,
마침내 $R$이 체인 경우로 환원된다. 이 경우 $R$ 위에서
$S\otimes_R S$와 $S$의 랭크가 같으려면 $S$의 랭크가 $0$ 또는
$1$이어야 한다. 따라서 $u:R\to S$는 전사이다.
\end{proof}

\begin{proposition}[21.1.7]
\label{0.21.1.7-ko}
$A$를 환, $B$를 $A$-대수, $(x_\alpha)_{\alpha\in I}$를 $B$의
원소족이라 하자. 다음 성질들을 생각하자.
\begin{enumerate}
  \item[\rm a)] $B=A[B^p,(x_\alpha)]$, 다시 말해 $A[B^p]$-대수
  $B$는 원소족 $(x_\alpha)_{\alpha\in I}$로 생성된다.
  \item[\rm b)] $A[B^p]$-가군 $B$는 단항식
  $\prod_\alpha x_\alpha^{n(\alpha)}$들로 생성된다. 여기서
  $(n(\alpha))_{\alpha\in I}$는 유한 지지를 갖고 모든 $\alpha\in I$에
  대하여 $0\leq n(\alpha)<p$인 정수족이다.
  \item[\rm c)] $B$-가군 $\Omega^1_{B/A}$는
  $d_{B/A}(x_\alpha)$들 $(\alpha\in I)$로 생성된다.
\end{enumerate}
그러면 a)와 b)는 동치이고 c)를 함의한다. 더 나아가 $B$가 유한형
$A[B^p]$-대수이면 c)는 a), b)와 동치이다.
\end{proposition}

\begin{proof}
b)가 a)를 함의함은 명백하다. 역으로 a)는 b)를 함의한다. 실제로
$m(\alpha)\geq0$인 유한 지지 정수족에 대한 모든 단항식
$\prod_\alpha x_\alpha^{m(\alpha)}$은, 각 $m(\alpha)$를 $p$로 나누어
\[
  m(\alpha)=p q(\alpha)+r(\alpha),\qquad 0\leq r(\alpha)<p
\]
라 쓰면
\[
  \left(\prod_\alpha x_\alpha^{q(\alpha)}\right)^p
  \prod_\alpha x_\alpha^{r(\alpha)}
\]
의 꼴로 쓸 수 있다. a)가 c)를 함의한다는 것은
\sref{0.21.1.5-ko}(ii)와 \sref{0.20.4.7-ko}에서 따른다.

역으로 c)가 성립하고 $B$가 유한형 $A[B^p]$-대수라고 하자. $B'$를
$x_\alpha$들로 생성되는 $B$의 부분 $A[B^p]$-대수라 하자. 완전열
\sref{0.20.5.7.1-ko}
\[
  \Omega^1_{B'/A[B^p]}\otimes_{B'}B\longrightarrow
  \Omega^1_{B/A[B^p]}\longrightarrow\Omega^1_{B/B'}\longrightarrow0
\]
에서 가정과 \sref{0.21.1.5-ko}(ii)에 의하여 왼쪽 화살표는
전사이다. 따라서 $\Omega^1_{B/B'}=0$이다. 한편
$B^p\subset B'\subset B$이고 $B$는 유한형 $B'$-대수이므로
\sref{0.21.1.6-ko}에 의하여 반드시 $B'=B$이다.
\end{proof}

\begin{remark}[21.1.8]
\label{0.21.1.8-ko}
$B$가 체이면 유한성 가정 없이도 a), b), c)의 동치를 증명할 것이다
\sref{0.21.4.5-ko}.
\end{remark}

\begin{definition}[21.1.9]
\label{0.21.1.9-ko}
$A$를 표수 $p$인 환, $B$를 $A$-대수라 하자. $B$의 원소족
$(x_\alpha)_{\alpha\in I}$에 대하여 단항식족
\[
  \prod_\alpha x_\alpha^{n(\alpha)}
  \qquad
  \bigl(0\leq n(\alpha)<p,
  \ (n(\alpha))_{\alpha\in I}\text{는 유한 지지}\bigr)
\]
이 $A[B^p]$-가군 $B$ 안에서 자유족이면 $(x_\alpha)$를 $A$ 위
$p$-자유족이라 한다. 이 단항식족이 생성원계이면 $(x_\alpha)$를
$A$ 위 $B$의 $p$-생성원계라 하고, 기저이면 $A$ 위 $B$의
$p$-기저라 한다.
\end{definition}

$B$의 원소들의 집합 $M$에 대해서도 표준 단사 $M\to B$가 정하는
원소족을 사용하여 같은 정의를 한다. $A$로 소체 $\mathbf F_p$를
취하면 $A[B^p]=B^p$이며, 이 경우 앞의 정의에서 $A$에 대한 언급을
생략한다. 또는 원소족이 ``절대적으로'' $p$-자유라고, 절대
$p$-생성원계라고, 절대 $p$-기저라고도 한다.

\oldpage[0\textsubscript{IV-1}]{157}
\begin{lemma}[21.1.10]
\label{0.21.1.10-ko}
$C$를 $B^p\subset C$인 $B$의 부분 $A$-대수, $M$을 $C$의 부분집합,
$N$을 $B$의 부분집합이라 하자.
\begin{enumerate}
  \item[(i)] $M$이 $A$ 위 $C$의 $p$-생성원계이고 $N$이 $C$ 위
  $B$의 $p$-생성원계이면, $M\cup N$은 $A$ 위 $B$의
  $p$-생성원계이다.
  \item[(ii)] $M$이 $A$ 위 $C$의 $p$-기저라고 하자. 그러면 $N$이
  $C$ 위 $p$-자유이기 위한 필요충분조건은 $M\cup N$이 $A$ 위
  $p$-자유인 것이다.
\end{enumerate}
\end{lemma}

\begin{proof}
$B^p\subset A\subset C\subset B$인 경우로 환원할 수 있다. 이때 (i)는
다음 일반적 사실의 특수한 경우이다. $P$가 $A$-가군 $C$의 생성원계이고
$Q$가 $C$-가군 $B$의 생성원계이면, $x\in P$, $y\in Q$인 곱
$xy$들의 집합은 $A$-가군 $B$의 생성원계이다.

같은 기호를 유지하자. $P$가 $A$-가군 $C$의 기저이면, $Q$가 $C$
위 자유족이라는 것은
\[
  \sum_{\lambda,\mu}a_{\lambda\mu}x_\lambda y_\mu=0
\]
이라는 관계에서 $a_{\lambda\mu}\in A$, $x_\lambda\in P$,
$y_\mu\in Q$이고 $x_\lambda$들 및 $y_\mu$들이 각각 서로 다를 때,
모든 쌍 $(\lambda,\mu)$에 대하여 $a_{\lambda\mu}=0$인 것과 동치이다.
이는 다시 모든 $\mu$에 대하여
\[
  \sum_\lambda a_{\lambda\mu}x_\lambda=0
\]
인 것과 동치이다. 따라서 (ii)를 얻는다.
\end{proof}

\subsection{$p$-기저와 형식적 매끄러움}
\label{subsection:0.21.2-ko}

\begin{theorem}[21.2.1]
\label{0.21.2.1-ko}
$B$를 환, $A$를 $B^p\subset A\subset B$인 $B$의 부분환,
$(x_\alpha)_{\alpha\in I}$를 $A$ 위 $B$의 $p$-기저라 하자.
$E$를 $A$-대수, $u:A\to E$를 구조 준동형,
$q:E\to B$를 전사 $A$-준동형이라 하고, 그 핵을 $\mathfrak R$이라
하자. $\mathfrak R^p=0$이라고 가정하자. 그러면 다음이 성립한다.
\begin{enumerate}
  \item[(i)] 준동형 $q:E\to B$의 오른쪽 역인 $A$-준동형
  $v:B\to E$가 존재하기 위한 필요충분조건은 모든 $z\in E$에 대하여
  \[
    u(q(z)^p)=z^p
  \]
  인 것이다.
  \item[(ii)] (i)의 조건이 성립한다고 하자. 모든 $\alpha\in I$에
  대하여 $q(z_\alpha)=x_\alpha$인 $E$의 임의의 원소족
  $(z_\alpha)_{\alpha\in I}$에 대하여, 모든 $\alpha$에 대해
  $v(x_\alpha)=z_\alpha$이고 $q$의 오른쪽 역인 $A$-준동형
  $v:B\to E$가 유일하게 존재한다.
\end{enumerate}
\end{theorem}

\begin{proof}
$q$의 오른쪽 역인 $A$-준동형 $v:B\to E$가 존재한다고 하자. 모든
$a\in A\subset B$에 대하여 $v(a)=u(a)$이고 $B^p\subset A$이므로,
모든 $z\in E$에 대하여
\[
  v(q(z)^p)=u(q(z)^p)
\]
이다. 한편 $v(q(z)^p)=(v(q(z)))^p$이고, $q\circ v=1_B$이므로
$v(q(z))\equiv z\pmod{\mathfrak R}$이다. 따라서
$\mathfrak R^p=0$에 의하여 $(v(q(z)))^p=z^p$이다. 이로써 (i)의
필요성을 얻는다. (i) 조건의 충분성은 (ii)에서 따른다.

(ii)의 가정 아래에서 단항식 $\prod_\alpha x_\alpha^{n(\alpha)}$들이
$A$-가군 $B$를 생성하므로 $v$의 유일성은 명백하다. 더욱이 유한
지지를 갖고 $0\leq n(\alpha)<p$인 원소족에 대한 이 단항식들은
$A$-가군 $B$의 기저를 이루므로, 다음을 만족하는 $A$-선형 사상
$v:B\to E$가 존재한다.
\[
  v\left(\prod_\alpha x_\alpha^{n(\alpha)}\right)
  =\prod_\alpha z_\alpha^{n(\alpha)}.
\]
$v$가 환 준동형임을 보이는 일만 남는다. 다음과 같이 쓸 수 있다.
\[
  \left(\prod_\alpha x_\alpha^{m(\alpha)}\right)
  \left(\prod_\alpha x_\alpha^{n(\alpha)}\right)
  =a\prod_\alpha x_\alpha^{r(\alpha)},
\]
여기서 $m(\alpha)+n(\alpha)<p$이면
$r(\alpha)=m(\alpha)+n(\alpha)$이고, 그렇지 않으면
$r(\alpha)=m(\alpha)+n(\alpha)-p$이다. 또한 $a\in A$는
\oldpage[0\textsubscript{IV-1}]{158}
$m(\alpha)+n(\alpha)\geq p$인 $\alpha\in I$에 대한
$x_\alpha^p$들의 곱이다. 이제
\[
  u(a)=\prod_\alpha z_\alpha^p
\]
임을 보이면 충분한데, 이는 $x_\alpha=q(z_\alpha)$와 (i)의 가정에서
따른다.
\end{proof}

\begin{corollary}[21.2.2]
\label{0.21.2.2-ko}
$B$를 환, $A$를 $B^p\subset A\subset B$인 $B$의 부분환이라 하자.
$A$ 위 $B$의 $p$-기저가 존재하면, $B$는 이산 위상에 대하여
$B^p$에 관해 상대적으로 형식적으로 매끄러운 $A$-대수이다.
\end{corollary}

\begin{proof}
$E$를 $B$-가군 $L$에 의한 $B$의 $A$-확장, $q:E\to B$를 증대,
$u:A\to E$를 구조 준동형이라 하자. $E$가 $B^p$-자명하다는 것은
모든 $b\in B$에 대하여
\[
  q(v(b))=b,\qquad v(b^p)=u(b^p)
\]
를 만족하는 환 준동형 $v:B\to E$가 존재한다는 뜻이다. 따라서
$z\in E$에 대하여
\[
  u(q(z)^p)=v(q(z)^p)=(v(q(z)))^p
\]
이다. 그런데 $L^2=0$이므로 $L^p=0$이고,
$v(q(z))-z\in L$이므로 $(v(q(z)))^p=z^p$이다. 따라서
\sref{0.21.2.1-ko}(i)의 조건이 성립하고 $E$는 $B^p$-자명하다.
\end{proof}

\begin{corollary}[21.2.3]
\label{0.21.2.3-ko}
$B$를 환, $A$를 $B^p\subset A\subset B$인 $B$의 부분환,
$(x_\alpha)_{\alpha\in I}$를 $A$ 위 $B$의 $p$-기저라 하고,
$L$을 $B$-가군이라 하자. 그러면 다음이 성립한다.
\begin{enumerate}
  \item[(i)] $A$에서 $L$로 가는 미분 $D$가 $B$에서 $L$로 가는
  미분으로 연장되기 위한 필요충분조건은 $D$가 $B^p$에서 $0$인
  것이다.
  \item[(ii)] $D$가 $B^p$에서 $0$이라고 하자. $L$의 임의의 원소족
  $(y_\alpha)_{\alpha\in I}$에 대하여, $D$를 연장하고 모든
  $\alpha\in I$에 대해 $D'(x_\alpha)=y_\alpha$인 미분
  $D':B\to L$이 유일하게 존재한다.
\end{enumerate}
\end{corollary}

\begin{proof}
$A$에서 $L$로 가는 미분 $D$가 주어졌다고 하자. 환 $E=D_B(L)$와
\[
  u:A\longrightarrow E,\qquad u(a)=(a,D(a))
\]
를 생각하자. 표준 준동형 $q:E\to B$의 오른쪽 역인 $A$-준동형
$v:B\to E$는
\[
  x\longmapsto(x,D'(x))
\]
의 꼴이며, 여기서 $D'$은 $D$를 연장하는 $B$에서 $L$로 가는
미분이다 \sref{0.20.1.5-ko}. 또한 $z\in E$에 대하여
\[
  u(q(z)^p)=(q(z)^p,D(q(z)^p))
\]
이므로, $E$에 \sref{0.21.2.1-ko}을 적용하면 따름정리를 곧바로
얻는다.
\end{proof}

\sref{0.21.2.3-ko}에 의하여 완전열
\begin{equation}
\tag{21.2.3.1}
\label{0.21.2.3.1-ko}
0\longrightarrow\operatorname{Der}_A(B,L)
\longrightarrow\operatorname{Der}_{B^p}(B,L)
\longrightarrow\operatorname{Der}_{B^p}(A,L)\longrightarrow0
\end{equation}
을 얻는다 \sref{0.20.2.2.1-ko}. 또한 \sref{0.21.2.3-ko}에서
$D=0$으로 놓으면 사상
\begin{equation}
\tag{21.2.3.2}
\label{0.21.2.3.2-ko}
D'\longmapsto(D'(x_\alpha))_{\alpha\in I}
\end{equation}
은 $B$-가군 $\operatorname{Der}_A(B,L)$에서 곱 $B$-가군 $L^I$로
가는 동형사상이다.

\begin{corollary}[21.2.4]
\label{0.21.2.4-ko}
\sref{0.21.2.3-ko}의 가정 아래에서 완전열
\begin{equation}
\tag{21.2.4.1}
\label{0.21.2.4.1-ko}
0\longrightarrow\Omega^1_{A/B^p}\otimes_A B
\longrightarrow\Omega^1_{B/B^p}
\longrightarrow\Omega^1_{B/A}\longrightarrow0
\end{equation}
은 분할되고, 원소족
$(d_{B/A}(x_\alpha))_{\alpha\in I}$는 $B$-가군
$\Omega^1_{B/A}$의 기저이다.
\end{corollary}

\begin{proof}
이는 \sref{0.21.2.3-ko}과 공식
\[
  \operatorname{Der}_A(B,L)
  =\operatorname{Hom}_B(\Omega^1_{B/A},L)
\]
에서 곧바로 따른다 \sref{0.20.4.8-ko}.
\end{proof}

\begin{corollary}[21.2.5]
\label{0.21.2.5-ko}
$A$를 환, $B$를 $A$-대수, $(x_\alpha)_{\alpha\in I}$를 $A$ 위
$B$의 $p$-기저라 하자. 그러면
$(d_{B/A}(x_\alpha))_{\alpha\in I}$는 $B$-가군
$\Omega^1_{B/A}$의 기저이다.
\end{corollary}

\begin{proof}
\sref{0.21.1.5-ko}(ii)를 사용하면 $A=A[B^p]$인 경우로 환원되며,
이제 \sref{0.21.2.4-ko}을 적용하면 된다.
\end{proof}

\oldpage[0\textsubscript{IV-1}]{159}
\begin{env}[21.2.6]
\label{0.21.2.6-ko}
$A,B$를 두 환, $u:A\to B$를 환 준동형이라 하자.
\sref{0.21.1.4-ko}의 기호 아래에서 가환 도식
\begin{equation}
\tag{21.2.6.1}
\label{0.21.2.6.1-ko}
\xymatrix{
 A^{(p)} \ar[r]^{u} & B^{(p)}\\
 A \ar[u]^{F_A} \ar[r]^{u} & B \ar[u]_{F_B}
}
\end{equation}
을 얻는다. 이로부터 표준적으로 $B$-대수 준동형
\begin{equation}
\tag{21.2.6.2}
\label{0.21.2.6.2-ko}
u\otimes1:A^{(p)}\otimes_A B\longrightarrow B^{(p)}
\end{equation}
을 얻으며, 그 상은 환 $A[B^p]$이다. 여기서 $A[B^p]$는 준동형
$F_B:B\to B^p$에 의한 $B$-대수이다.
\end{env}

\begin{theorem}[21.2.7]
\label{0.21.2.7-ko}
$A$를 환, $B$를 $A$-대수, $u:A\to B$를 구조 준동형이라 하자.
다음 조건들을 가정하자.
\begin{enumerate}
  \item[(i)] $u$에서 유도되는 준동형 \sref{0.21.2.6.2-ko}가
  단사이다.
  \item[(ii)] $B$는 $A$에 대한 $p$-기저
  $(x_\alpha)_{\alpha\in I}$를 갖는다.
\end{enumerate}
그러면 $B$는 이산 위상에 대하여 형식적으로 매끄러운 $A$-대수이다.
더 정확히 말해 $A,B$에 이산 위상을 주자. $E$를 허용가능한 위상
$A$-대수 \sref{0.7.1.2-ko}, $\mathfrak R$을 $E$의 정의 아이디얼,
$C=E/\mathfrak R$, $v:B\to C$를 $A$-준동형,
$q:E\to C$를 증대라 하자. 모든 $\alpha\in I$에 대하여
$q(z_\alpha)=v(x_\alpha)$인 $E$의 임의의 원소족
$(z_\alpha)_{\alpha\in I}$에 대해,
\[
  v=q\circ w,\qquad w(x_\alpha)=z_\alpha\quad(\alpha\in I)
\]
를 만족하는 $A$-준동형 $w:B\to E$가 유일하게 존재한다.
\end{theorem}

\begin{proof}
먼저 $E$가 이산적인 경우를 생각하자. 이때 $\mathfrak R$은 멱영
아이디얼이다. $v$가 전사인 경우로 환원할 수 있고, 더 나아가
\sref{0.19.4.3-ko}의 논증에 의하여 $\mathfrak R^2=0$, 따라서
$\mathfrak R^p=0$이라고 가정할 수 있다. 끝으로 $\mathfrak R$에 의한
$C$의 확장 $E$의 $v$에 의한 역상을 생각하면, $C=B$이고 $v$가
항등인 경우로 환원할 수 있다 \sref{0.19.4.4-ko}.
환 준동형 $F_E:z\rightsquigarrow z^p$는 $\mathfrak R$에서 $0$이므로
$E\xrightarrow q B\xrightarrow{F'}E$로 인수분해된다. $F'$를 $B$에서
$E^{(p)}$로 가는 준동형으로 생각하면, 합성 준동형
$A\xrightarrow r E\xrightarrow{F_E}E^{(p)}$를 사용하여 정한 $E^{(p)}$의
$A$-대수 구조의 정의에 의해 $F'$는 \emph{$A$-준동형}이다. 여기서
$r:A\to E$는 $A$-대수 $E$의 구조 준동형이다. 또한
$r:A^{(p)}\to E^{(p)}$도 $A$-준동형이다. 따라서 $f$와 표준 준동형
$A^{(p)}\to A^{(p)}\otimes_A B$, $B\to A^{(p)}\otimes_A B$의 합성이
각각 $r$, $F'$가 되는 $A$-준동형 $f:A^{(p)}\otimes_A B\to E^{(p)}$가
유일하게 존재한다. 그런데 가정에 의해 $A^{(p)}\otimes_A B$를
$A[B^p]$와 동일시할 수 있으며, 이때 표준 준동형
$A^{(p)}\to A^{(p)}\otimes_A B$, $B\to A^{(p)}\otimes_A B$는 각각
$u$, $F_B$와 동일시된다. 따라서 이제 준동형 $f$를 통해 $E$를
$A[B^p]$-대수로 볼 수 있으며, 구성상 모든 $z\in E$에 대하여
$f(q(z)^p)=z^p$이다. 그러므로 \sref{0.21.2.1-ko}의 적용 조건이
성립하여 이 경우 정리를 얻는다.

일반적인 경우로 넘어가기 위해 $E$의 열린 아이디얼들의 기본계
$(\mathfrak J_\lambda)$를 생각하고
$E_\lambda=E/\mathfrak J_\lambda$,
$\mathfrak R_\lambda=(\mathfrak R+\mathfrak J_\lambda)/\mathfrak J_\lambda$,
$C_\lambda=E_\lambda/\mathfrak R_\lambda=E/(\mathfrak R+\mathfrak J_\lambda)$라
놓자. $E$가 허용가능하므로 $E=\varprojlim E_\lambda$이다
\sref{0.7.2.2-ko}. 각 쌍 $(\alpha,\lambda)$에 대하여
$z_{\alpha\lambda}$를 $z_\alpha$의 $E_\lambda$ 안의 표준 상이라 하고,
$p_\lambda:C\to C_\lambda$를 표준 준동형이라 하며,
\oldpage[0\textsubscript{IV-1}]{160}
$q_\lambda$를 표준 준동형 $E_\lambda\to C_\lambda$라 하자.
가정에 의하여 $\mathfrak R_\lambda$는 $E_\lambda$에서 멱영 아이디얼이므로,
증명의 첫 부분은 $p_\lambda\circ v=q_\lambda\circ w_\lambda$이고 모든
$\alpha$에 대하여 $w_\lambda(x_\alpha)=z_{\alpha\lambda}$인
$A$-준동형 $w_\lambda:B\to E_\lambda$의 존재성과 유일성을 보인다.
$w_\lambda$들의 유일성에서 $(w_\lambda)$가 준동형들의 사영계를 이룸이
곧바로 따르며, $w=\varprojlim w_\lambda$가 원하는 준동형이다.
\end{proof}

\begin{remark}[21.2.8]
\label{0.21.2.8-ko}
모든 $\alpha$에 대하여 $w(x_\alpha)=z_\alpha$인 준동형
$w:B\to E$의 존재성과 유일성은, $B$가 형식적으로 매끄러운
$A$-대수이고 $d_{B/A}(x_\alpha)$들이 $B$-가군
$\Omega^1_{B/A}$의 기저를 이룬다는 사실 \sref{0.21.2.4-ko}에서
$p$-기저라는 사실을 사용하지 않고도 직접 얻을 수 있다. 따라서 이
결과는 환들이 표수 $p$라고 가정하지 않아도 성립한다.

실제로 $\mathfrak R^2=0$인 경우로 환원할 수 있다. 등식
\[
  \operatorname{Der}_A(B,\mathfrak R)
  =\operatorname{Hom}_B(\Omega^1_{B/A},\mathfrak R)
\]
에 의하여 \sref{0.20.4.8-ko}, $\mathfrak R$의 임의의 원소족
$(z'_\alpha)$에 대해 모든 $\alpha$에 대하여
$D(x_\alpha)=z'_\alpha$인 $B$에서 $\mathfrak R$로 가는
$A$-미분 $D$가 유일하게 존재한다. 결론은 \sref{0.20.1.1-ko}에서
따른다.
\end{remark}

\subsection{$p$-기저와 불완전성 가군}
\label{subsection:0.21.3-ko}

\begin{env}[21.3.1]
\label{0.21.3.1-ko}
$A,B$를 표수 $p$인 두 환이라 하고, 두 환 준동형
$i:A\to B$, $j:B\to A$가 다음을 만족한다고 하자.
\begin{align}
\tag{21.3.1.1}\label{0.21.3.1.1-ko}
j(i(a))&=a^p\quad\text{모든 }a\in A\text{에 대하여},\\
\tag{21.3.1.2}\label{0.21.3.1.2-ko}
i(j(b))&=b^p\quad\text{모든 }b\in B\text{에 대하여}.
\end{align}
대부분의 경우 $i$는 단사이므로 $i:A\to B$를 통해 $A$를 $B$의
부분환과 동일시한다. 이 동일시 아래에서 $j$의 존재는
$B^p\subset A$를 함의하고, $j$는 $F_B$와 동일시된다.
\end{env}

\begin{env}[21.3.2]
\label{0.21.3.2-ko}
$h:A^p\to A$를 표준 단사 준동형이라 하자.
\sref{0.21.3.1.1-ko}과 \sref{0.21.3.1.2-ko}에 의하여 가환 도식
\begin{equation}
\tag{21.3.2.1}\label{0.21.3.2.1-ko}
\xymatrix{
 B \ar[r]^{j} & A\\
 A \ar[u]^{i} \ar[r]^{F_A} & A^p \ar[u]_{h}
}
\end{equation}
을 얻는다. 따라서 쌍 $(j,F_A)$를 $i$에 의한 $A$-대수 $B$에서
$h$에 의한 $A^p$-대수 $A$로 가는 디-준동형으로 볼 수 있다. 이로부터
표준 $A$-가군 준동형
\begin{equation}
\tag{21.3.2.2}\label{0.21.3.2.2-ko}
\pi_{B/A}:\Omega^1_{B/A}\otimes_B A_{[j]}
\longrightarrow\Omega^1_{A/A^p}=\Omega^1_A
\end{equation}
을 얻는다. 여기서 \sref{0.20.5.4-ko}을 사용하였으며,
$\Omega^1_{A/A^p}$와 절대 미분 가군 $\Omega^1_A$의 동일시는
\sref{0.21.1.5-ko}(ii)에서 따른다. 다음과 같이 놓는다.
\begin{align}
\tag{21.3.2.3}\label{0.21.3.2.3-ko}
\Theta_{B/A}&=\Omega^1_{B/A}\otimes_B A_{[j]},\\
\tag{21.3.2.4}\label{0.21.3.2.4-ko}
\Xi_{B/A}&=\operatorname{Ker}(\pi_{B/A}).
\end{align}
\oldpage[0\textsubscript{IV-1}]{161}
따라서 완전열
\begin{equation}
\tag{21.3.2.5}\label{0.21.3.2.5-ko}
0\longrightarrow\Xi_{B/A}\longrightarrow\Theta_{B/A}
\xrightarrow{\pi_{B/A}}\Omega^1_A
\end{equation}
을 얻는다. 기호 $\Theta_{B/A}$와 $\Xi_{B/A}$는 $A,B$뿐 아니라
$i,j$에도 의존하므로 혼동을 일으킬 수 있음에 유의하자.
\end{env}

\begin{env}[21.3.3]
\label{0.21.3.3-ko}
\sref{0.21.3.1.2-ko}에 의하여 $F_B=i\circ j$이므로, 모든
$B$-가군 $M$에 대하여 \sref{0.21.1.4-ko}
\begin{equation}
\tag{21.3.3.1}\label{0.21.3.3.1-ko}
M^{(p)}=M\otimes_B B^{(p)}
=(M\otimes_B A_{[j]})\otimes_A B_{[i]}
\end{equation}
라고 쓸 수 있다. 특히
\begin{equation}
\tag{21.3.3.2}\label{0.21.3.3.2-ko}
(\Omega^1_{B/A})^{(p)}=\Theta_{B/A}\otimes_A B_{[i]}
\end{equation}
이다. 따라서 \sref{0.21.3.2.2-ko}에서 표준 $B$-가군 준동형
\begin{equation}
\tag{21.3.3.3}\label{0.21.3.3.3-ko}
\pi_{B/A}\otimes1_B:(\Omega^1_{B/A})^{(p)}
\longrightarrow\Omega^1_A\otimes_A B_{[i]}
\end{equation}
을 얻는다. \sref{0.20.5.4.2-ko}을 고려하면 이 준동형의 상은
$b\in B$에 대한 원소 $d_A(j(b))\otimes1$들로 생성되는
$B$-가군이다. 이 원소들의 $i$에서 유도되는 표준 준동형
\[
  i_{B/A}:\Omega^1_A\otimes_A B\longrightarrow\Omega^1_B
\]
에 의한 상은 \sref{0.20.5.2-ko}, $b\in B$에 대한
$d_B(i(j(b)))$들로 생성되는 $\Omega^1_B$의 부분 $B$-가군 안에 있다.
\sref{0.21.3.1.2-ko}에 의하여 이 상은 $0$이다. 다시 말해 연속하는
두 준동형의 합성이 $0$인 준동형열
\begin{equation}
\tag{21.3.3.4}\label{0.21.3.3.4-ko}
0\longrightarrow\Xi_{B/A}\otimes_A B_{[i]}
\longrightarrow(\Omega^1_{B/A})^{(p)}
\longrightarrow\Upsilon_{B/A}
\end{equation}
을 얻는다. 이 열은 반드시 완전한 것은 아니다.
\end{env}

\begin{proposition}[21.3.4]
\label{0.21.3.4-ko}
$B$가 $i$에 관하여 평탄한 $A$-가군이면 열
\sref{0.21.3.3.4-ko}는 완전하다.
\end{proposition}

\begin{proof}
이는 평탄성의 정의에서 따른다. 특히 $B^p\subset A\subset B$이고
$i,j$가 각각 표준 단사 준동형과 $F_B$이며, $B$가 $A$ 위
$p$-기저를 갖는 경우에 명제를 적용할 수 있다. 실제로 이때 $B$는
자유 $A$-가군이다. 그러나 이 경우에도 핵 $\Xi_{B/A}$는 반드시
$0$인 것은 아니다. 다만 다음 명제가 성립한다.
\end{proof}

\begin{proposition}[21.3.5]
\label{0.21.3.5-ko}
$B$를 축소환, $A$를 $B^p\subset A\subset B$인 $B$의 부분환이라 하자.
$A$ 위 $B$의 $p$-기저 $(x_\lambda)_{\lambda\in L}$와 $B^p$ 위
$A$의 $p$-기저 $(y_\mu)_{\mu\in M}$가 존재한다고 하자. 그러면
\sref{0.21.3.3.4-ko}의 표준 준동형
\begin{equation}
\tag{21.3.5.1}\label{0.21.3.5.1-ko}
(\Omega^1_{B/A})^{(p)}\longrightarrow\Upsilon_{B/A}
\end{equation}
은 전단사이고, 원소 $d_A(x_\lambda^p)\otimes1$들은 $B$-가군
$\Upsilon_{B/A}$의 기저를 이룬다.
\end{proposition}

\begin{proof}
\oldpage[0\textsubscript{IV-1}]{162}
$B$가 축소환이므로 $x\mapsto x^p$는 $B$에서 $B^p$로 가는 환
동형이다. 이 동형을 따라 구조를 옮기면
$(x_\lambda^p)_{\lambda\in L}$와 $y_\mu$들이 $A^p$ 위 $B^p$의
$p$-기저를 이룸을 알 수 있다 \sref{0.21.1.10-ko}. 따라서
\sref{0.21.1.5-ko}(ii)와 \sref{0.21.2.5-ko}에 의하여
$d_A(x_\lambda^p)$들과 $d_A(y_\mu)$들은 $A$-가군
$\Omega^1_A$의 기저를 이룬다. 한편 $d_A(x_\lambda^p)\otimes1$의
$i_{B/A}$에 의한 상은 $d_B(x_\lambda^p)=0$이다. 다른 한편
$d_A(y_\mu)$의 $i_{B/A}$에 의한 상은 $d_B(i(y_\mu))$이고,
$y_\mu$들과 $x_\lambda$들이 $B^p$ 위 $B$의 $p$-기저를 이루므로
\sref{0.21.1.10-ko}, $d_B(y_\mu)$들은 $\Omega^1_B$의 한 기저의
일부를 이룬다 \sref{0.21.2.5-ko}. 따라서 이들은 $B$ 위에서
선형독립이다. 그러므로 $i_{B/A}$의 핵은
$d_A(x_\lambda^p)\otimes1$들을 기저로 갖는다. 이 원소들은
\sref{0.21.3.5.1-ko}에 의한 $(\Omega^1_{B/A})^{(p)}$의 대응하는
원소들의 상이고, 그 대응하는 원소들은 $B$-가군
$(\Omega^1_{B/A})^{(p)}$의 기저를 이루므로
\sref{0.21.2.5-ko}, $B$-가군 준동형 \sref{0.21.3.5.1-ko}는
전단사이다.
\end{proof}

\begin{remarks}[21.3.6]
\label{0.21.3.6-ko}
\begin{enumerate}
  \item[(i)] $A',B'$을 표수 $p$인 두 환이라 하자.
  \sref{0.21.3.1.1-ko}과 \sref{0.21.3.1.2-ko}를 만족하는 두 준동형
  $i':A'\to B'$, $j':B'\to A'$ 및 다음 도식을 가환이 되게 하는 환
  준동형 $f:A\to A'$, $g:B\to B'$이 주어졌다고 하자.
  \[
  \xymatrix{
  A'\ar[r]^{i'}&B'\ar[r]^{j'}&A'\\
  A\ar[u]^f\ar[r]^i&B\ar[u]^g\ar[r]^j&A\ar[u]^f
  }
  \]
  그러면 표준 준동형
  $\Omega^1_{B/A}\to\Omega^1_{B'/A'}$
  \sref{0.20.5.4.3-ko}은 다음 도식을 가환이 되게 하는 표준
  디-준동형 $\Theta_{B/A}\to\Theta_{B'/A'}$을 유도한다.
  \[
  \xymatrix{
  0\ar[r]&\Xi_{B/A}\ar[r]\ar[d]&\Theta_{B/A}\ar[r]^{\pi_{B/A}}\ar[d]&\Omega^1_A\ar[d]\\
  0\ar[r]&\Xi_{B'/A'}\ar[r]&\Theta_{B'/A'}\ar[r]^{\pi_{B'/A'}}&\Omega^1_{A'}
  }
  \]
  \item[(ii)] $A'$을 임의의 $A$-대수라 하고
  $B'=B\otimes_A A'$이라 놓자. 그러면
  $i'=i\otimes1:A'\to B'$과 $j'=j\otimes1:B'\to A'$은
  \sref{0.21.3.1.1-ko}과 \sref{0.21.3.1.2-ko}를 만족하고,
  \[
    \Omega^1_{B'/A'}=\Omega^1_{B/A}\otimes_A A'
    =\Omega^1_{B/A}\otimes_B B'
  \]
  이다 \sref{0.20.5.5-ko}. 이로부터 표준 $A'$-동형
  \begin{equation}
  \tag{21.3.6.1}\label{0.21.3.6.1-ko}
  \Theta_{B/A}\otimes_A A'\xrightarrow{\sim}\Theta_{B'/A'}
  \end{equation}
  과 표준 $B'$-준동형
  \begin{equation}
  \tag{21.3.6.2}\label{0.21.3.6.2-ko}
  (\Omega^1_{B/A})^{(p)}\otimes_B B'\xrightarrow{\sim}
  (\Omega^1_{B'/A'})^{(p)}
  \end{equation}
  도 얻는다.
\end{enumerate}
\end{remarks}

\subsection{체 확대의 경우}
\label{subsection:0.21.4-ko}

\begin{env}[21.4.0]
\label{0.21.4.0-ko}
$K$를 표수 $p>0$인 체, $k$를 $K$의 부분체라 하자. $k$의 모든
원소가 $K^p$ 위에서 대수적이므로 환 $k[K^p]$는 체 $k(K^p)$와
같다. 따라서 앞 번호들의 결과에서 $A,B,A[B^p]$를 각각
$k,K,k(K^p)$로 바꾸어 적용할 수 있다.
\end{env}

\begin{lemma}[21.4.1]
\label{0.21.4.1-ko}
$k$를 표수 $p>0$인 체, $K$를 $k$의 확대체라 하자. 원소
$x\in K$가 $k$ 위에서 $p$-자유이기 위한 필요충분조건은
$x\notin k(K^p)$인 것이다.
\end{lemma}

\begin{proof}
실제로 $x$는 $k(K^p)[X]$의 다항식 $X^p-x^p$의 근이다.
$x\notin k(K^p)$이면 이 다항식이 기약임이 알려져 있으므로
(Bourbaki, \emph{Alg.}, chap.~V, §~8, n$^\circ$~1, prop.~1), 원소
$1,x,\ldots,x^{p-1}$은 $k(K^p)$-가군 $k(K^p)(x)$의 기저를 이룬다.
\end{proof}

\begin{theorem}[21.4.2]
\label{0.21.4.2-ko}
$k$를 표수 $p>0$인 체, $K$를 $k$의 확대체,
$S$를 $k$ 위 $K$의 $p$-생성원계, $L\subset S$를 $k$ 위에서
$p$-자유인 부분집합이라 하자. 그러면
$L\subset B\subset S$인 $k$ 위 $K$의 $p$-기저 $B$가 존재한다.
특히 $k$의 모든 확대체는 $k$ 위에서 $p$-기저를 갖는다.
\end{theorem}
\begin{proof}
$K^p\subset k$인 경우로 한정할 수 있다. 초른의 보조정리에 의하여
$L\subset B\subset S$이고 $k$ 위에서 $p$-자유이며, 이 성질들을
갖는 $S$의 부분집합들 가운데 극대인 부분집합 $B$가 $K$ 안에
존재한다. $k$와 $B$가 생성하는 $K$의 부분체 $K'$이 $K$와 같음을
보이면 충분하다. 그렇지 않다면 $K'$에 속하지 않는 $x\in S$가
존재한다. $K^p\subset k$이므로 $K'=K'(K^p)$이고, 따라서 $x$는
$K'$ 위에서 $p$-자유이다 \sref{0.21.4.1-ko}. 그러면
$B\cup\{x\}$는 $k$ 위에서 $p$-자유이므로
\sref{0.21.1.10-ko}, $B$의 정의에 모순된다. \sref{0.21.4.2-ko}의
마지막 주장은 $L=\varnothing$, $S=K$로 택하면 얻는다.
\end{proof}

\begin{corollary}[21.4.3]
\label{0.21.4.3-ko}
$k$를 표수 $p>0$인 체, $K$를 $k$의 확대체라 하자. $K$의 원소족
$(x_\alpha)$가 $k$ 위에서 $p$-자유이기 위한 필요충분조건은, 모든
$\alpha$에 대하여 $x_\alpha$가 $k(K^p)$와 첨자
$\beta\neq\alpha$인 $x_\beta$들이 생성하는 체 $K_\alpha$에 속하지
않는 것이다.
\end{corollary}

\begin{proof}
조건의 필요성은 \sref{0.21.1.10-ko}에서 따른다. 거꾸로 이 조건이
성립한다고 하자. $(x_\alpha)$가 $K$의 $p$-생성원계인 경우로
한정할 수 있다. 그러면 $(x_\alpha)$의 어떤 부분족이 $K$의
$p$-기저이다 \sref{0.21.4.2-ko}. 그러나 이 부분족이
$(x_\alpha)$와 다를 수는 없다. 그렇지 않으면 가정에 의하여 그
부분족은 $K$의 $p$-생성원계가 아니기 때문이다.
\end{proof}

\oldpage[0\textsubscript{IV-1}]{163}
\begin{corollary}[21.4.4]
\label{0.21.4.4-ko}
$k$를 표수 $p>0$인 체, $K$를 $k$의 확대체라 하자.
$\Omega^1_{K/k}=0$이기 위한 필요충분조건은 $K=k(K^p)$인 것이다.
특히 $\Omega^1_K=0$이기 위한 필요충분조건은 $K$가 완전체인 것이다.
\end{corollary}

\begin{proof}
실제로 $(x_\alpha)$가 $k$ 위 $K$의 $p$-기저이면
$d_{K/k}(x_\alpha)$들은 $K$-벡터공간 $\Omega^1_{K/k}$의 기저를
이룬다 \sref{0.21.2.5-ko}.
\end{proof}

\begin{theorem}[21.4.5]
\label{0.21.4.5-ko}
$k$를 표수 $p>0$인 체, $K$를 $k$의 확대체,
$(x_\alpha)$를 $K$의 원소족이라 하자. $(x_\alpha)$가 $k$ 위에서
$p$-자유족(각각 $k$ 위 $K$의 $p$-생성원계, $p$-기저)이기 위한
필요충분조건은 $(d_{K/k}(x_\alpha))$가 $K$-벡터공간
$\Omega^1_{K/k}$에서 자유족(각각 생성원계, 기저)인 것이다.
\end{theorem}

\begin{proof}
$K'$을 $k(K^p)$와 $x_\alpha$들이 생성하는 $K$의 부분체(부분환과
같다)라 하자. \sref{0.20.4.7-ko}을 고려하면
$d_{K'/k}(x_\alpha)$들은
$\Omega^1_{K'/k(K^p)}=\Omega^1_{K'/k}$를 생성한다.
$d_{K/k}(x_\alpha)$들이
$\Omega^1_{K/k}=\Omega^1_{K/k(K^p)}$를 생성하면, 완전열
\sref{0.20.5.7.1-ko}
\[
\Omega^1_{K'/k(K^p)}\otimes_{K'}K\longrightarrow
\Omega^1_{K/k(K^p)}\longrightarrow\Omega^1_{K/K'}\longrightarrow0
\]
의 왼쪽 화살표는 전사이다. 따라서 $\Omega^1_{K/K'}=0$이고,
\sref{0.21.4.4-ko}에 의하여 $K'=K$이다.

이제 $(x_\alpha)$가 $p$-자유족이면 이는 $k$ 위 $K$의 어떤
$p$-기저의 부분족이다 \sref{0.21.4.2-ko}. 따라서
$d_{K/k}(x_\alpha)$들은 $K$-벡터공간 $\Omega^1_{K/k}$의 한 기저의
일부를 이루므로 \sref{0.21.2.5-ko}, $K$ 위에서 선형독립이다.
거꾸로 $d_{K/k}(x_\alpha)$들이 $K$ 위에서 선형독립이면
$(x_\alpha)$가 $k$ 위에서 $p$-자유족임을 보이자.

\oldpage[0\textsubscript{IV-1}]{164}
\sref{0.21.4.3-ko}을 고려하면, 모든 $\alpha$에 대하여
$x_\alpha\notin K_\alpha$임을 보이면 충분하다. 그런데 완전열
\[
\Omega^1_{K_\alpha/k(K^p)}\otimes_{K_\alpha}K\longrightarrow
\Omega^1_{K/k(K^p)}\longrightarrow\Omega^1_{K/K_\alpha}\longrightarrow0
\]
에서, 첨자 $\beta\neq\alpha$인 $d_{K_\alpha/k}(x_\beta)$들의 왼쪽
화살표에 의한 상은 $d_{K/k}(x_\beta)$들이다. 따라서 이 원소들은
$d_{K/k}(x_\alpha)$를 포함하지 않는 $\Omega^1_{K/k(K^p)}$의
벡터 부분공간을 생성한다. 이 부분공간은 오른쪽 화살표의 핵이므로
$d_{K/K_\alpha}(x_\alpha)\neq0$이고, 따라서
$x_\alpha\notin K_\alpha$이다.
\end{proof}

\begin{corollary}[21.4.6]
\label{0.21.4.6-ko}
$k$를 표수 $p>0$인 체, $K$를 $k$의 확대체, $x$를 $K$의 원소라
하자. 다음 세 조건은 동치이다.
\begin{enumerate}
  \item[a)] $x\notin k(K^p)$.
  \item[b)] $d_{K/k}(x)\neq0$.
  \item[c)] 원소 $x$는 $k$ 위에서 $p$-자유이다.
\end{enumerate}
\end{corollary}

\begin{proposition}[21.4.7]
\label{0.21.4.7-ko}
$L$을 표수 $p>0$인 체, $K$를 $L^p\subset K\subset L$을 만족하는
$L$의 부분체라 하자. 그러면 다음 $L$-벡터공간들의 열은
\begin{equation}
\tag{21.4.7.1}\label{0.21.4.7.1-ko}
0\to\Omega^1_{K/L^p}\otimes_KL\to\Omega^1_{L/L^p}
\to\Omega^1_{L/K}\to0
\end{equation}
완전하다. 달리 말하면, $\Upsilon_{L/K/L^p}=0$이다.
\end{proposition}

\begin{proof}
이는 \sref{0.21.4.2-ko}을 고려하면 \sref{0.21.2.4-ko}의 특수한
경우이다.
\end{proof}

\begin{proposition}[21.4.8]
\label{0.21.4.8-ko}
\sref{0.21.4.7-ko}의 가정 아래 표준 준동형
\begin{equation}
\tag{21.4.8.1}\label{0.21.4.8.1-ko}
(\Omega^1_{L/K})^{(p)}\longrightarrow\Upsilon_{L/K}
\end{equation}
은 전단사이다. $(x_\lambda)$가 $K$ 위 $L$의 $p$-기저이면
$d_K(x_\lambda^p)\otimes1$들은 $L$-벡터공간
$\Upsilon_{L/K}$의 기저를 이룬다.
\end{proposition}

\begin{proof}
이는 \sref{0.21.3.5-ko}의 특수한 경우이다.
\end{proof}

\subsection{응용: 분리성 판정법}
\label{0.21.5-ko}

이 번호와 다음 두 번호에서는 고려하는 환들이 표수 $p>0$이라는
가정을 더 이상 두지 않는다.

\begin{env}[21.5.1]
\label{0.21.5.1-ko}
먼저 판정법 \sref{0.21.2.7-ko}에 의해 분리 확대체에 관한 Cohen 정리
\sref{0.19.6.1-ko}의 일부를 증명할 수 있음을 주목하자. 즉 $k$가
표수 $p>0$이고 $K$가 $k$의 분리 확대체이면, $K$는 형식적으로
매끄러운 $k$-대수이다. 실제로 $K$는 $k$ 위에서 $p$-기저를
갖는다 \sref{0.21.4.2-ko}. 한편 MacLane 판정법
(Bourbaki, \emph{Alg.}, chap.~V, §~8, n$^\circ$~2, prop.~3)에
따르면, $K$의 한 대수적 폐포 안에서 $k^{p^{-1}}$과 $K$는 $k$ 위에서
선형적으로 서로소이다. 따라서 표준 준동형
$k^{p^{-1}}\otimes_k K\to k^{p^{-1}}(K)$는 전단사이다. 동형사상
$x\mapsto x^p$에 의해 구조를 옮기면, 이것이 바로
\sref{0.21.2.7-ko}의 조건 (i)이다.
\end{env}

\begin{proposition}[21.5.2]
\label{0.21.5.2-ko}
\textnormal{(MacLane).}
$A$, $B$를 두 완비 이산 값매김환, $\mathfrak m$, $\mathfrak n$을
각각의 극대 아이디얼, $K=A/\mathfrak m$, $L=B/\mathfrak n$을
각각의 잉여체라 하자. 또한 $u:A\to B$를
$Bu(\mathfrak m)=\mathfrak n$인 준동형, $u_0=u\otimes1:K\to L$을
잉여체에 유도되는 대응 준동형이라 하자. 다음 조건들을 고려한다.
\begin{enumerate}
  \item[\rm a)] $L$은 $u_0$에 관하여 $K$의 분리 확대체이다.
  \item[\rm b)] 극대 아이디얼이 $\mathfrak n'$이고 잉여체가 $L'$인
  임의의 완비 이산 값매김환 $B'$,
\oldpage[0\textsubscript{IV-1}]{165}
  $B'u'(\mathfrak m)=\mathfrak n'$인 임의의 준동형 $u':A\to B'$,
  그리고 $u_0$와 $u'_0=u'\otimes1:K\to L'$에 관한 임의의
  $K$-동형사상 $\sigma:L\to L'$에 대하여, $u'=w\circ u$이고
  $w_0=w\otimes1:L\to L'$이 $\sigma$와 같은 동형사상
  $w:B\to B'$이 존재한다.
  \item[\rm b$'$)] $Bu'(\mathfrak m)=\mathfrak n$이고
  $u'_0=u'\otimes1:K\to L$이 $u_0$와 같은 임의의 준동형
  $u':A\to B$에 대하여, $u'=w\circ u$이고
  $w_0=w\otimes1:L\to L$이 항등사상인 $B$의 자기동형사상 $w$가
  존재한다.
  \item[\rm c)] $u$에서 몫으로 옮겨 얻는 준동형을
  $u_1:A/\mathfrak m^2\to B/\mathfrak n^2$이라 하자. 그러면
  $\operatorname{gr}_0(u'_1)=\operatorname{gr}_0(u_1)(=u_0)$을 만족하는
  임의의 국소 준동형
  $u'_1:A/\mathfrak m^2\to B/\mathfrak n^2$에 대하여,
  $\operatorname{gr}_0(w_1)$이 항등사상이고
  $u'_1=w_1\circ u_1$인 $B/\mathfrak n^2$의 자기동형사상 $w_1$이
  존재한다.
  \item[\rm c$'$)]
  $\operatorname{gr}_0(u'_1)=\operatorname{gr}_0(u_1)$ 및
  $\operatorname{gr}_1(u'_1)=\operatorname{gr}_1(u_1)$을 만족하는
  임의의 국소 준동형
  $u'_1:A/\mathfrak m^2\to B/\mathfrak n^2$에 대하여
  ($\operatorname{gr}_1(u'_1)$과 $\operatorname{gr}_1(u_1)$은
  $\mathfrak m/\mathfrak m^2$에서 $\mathfrak n/\mathfrak n^2$으로 가는
  준동형이다), $\operatorname{gr}_0(w_1)$이 항등사상이고
  $u'_1=w_1\circ u_1$인 $B/\mathfrak n^2$의 자기동형사상 $w_1$이
  존재한다.
\end{enumerate}
그러면 $c)\Leftrightarrow c')\Leftrightarrow a)\Rightarrow b)
\Rightarrow b')$이다. 또한 $A$가 Cohen 환이면 앞의 다섯 조건은
서로 동치이다.

\end{proposition}

\begin{proof}
$b)\Rightarrow b')$와 $c)\Rightarrow c')$은 자명하다.
$a)$가 $b)$를 함의함을 보이자.
준동형 $u$는 가정에 따라 $A$의 균등화원을 $B$의 균등화원으로
보내므로, $B$는 꼬임 없는 $A$-가군이다. 따라서 $B$는 평탄
$A$-가군이다 ($0_{\rm I}$, 6.3.4). 그러므로 $B$는 Cohen
$A$-대수이다 \sref{0.19.8.1-ko}. 이제 $G=B'$,
$\mathfrak J=\mathfrak n'$로 놓고 \sref{0.19.8.2-ko}, (i)를 적용하면
$a)$가 $b)$를 함의한다. 여기서 $B'$은 $u'$에 의해 $A$-대수로
보고, 준동형 $B\to B'/\mathfrak n'$은 합성
$B\to B/\mathfrak n\to B'/\mathfrak n'$이다. 이는 가정
$u'_0=\sigma\circ u_0$에 의해 $A$-준동형이다. 같은 논증에서 이번에는
$G=B/\mathfrak n^2$, $\mathfrak J=\mathfrak n/\mathfrak n^2$로 놓고,
$G$를 $u'_1$에 의해 $A$-대수로 보면 $a)$가 $c)$를 함의함을 얻는다.

다음으로 $c')$가 $a)$를 함의함을 증명하자. 두 준동형 $u_1$과
$u'_1$에 대하여
$u'_1=u_1+D$이며, 여기서 $D$는 $A/\mathfrak m^2$에서
$\mathfrak n/\mathfrak n^2$으로 가는 미분이다
\sref{0.20.1.1-ko}. 더욱이 가정
$\operatorname{gr}_1(u'_1)=\operatorname{gr}_1(u_1)$은 $D$가
$\mathfrak m/\mathfrak m^2$에서 영임을 뜻한다. 따라서 $D$는
$K=A/\mathfrak m$에서 $\mathfrak n/\mathfrak n^2$으로 가는 미분으로
볼 수 있다. 한편 $B$의 균등화원을 하나 택하면
$\mathfrak n/\mathfrak n^2$은 $u_0$에 관한 $K$-가군 $L$과
동일시된다.

다른 한편 $w_1$에 부과된 조건들은
$\operatorname{gr}_1(w_1)$도 항등사상임을 함의한다. 실제로
$u_1(\mathfrak m/\mathfrak m^2)$이
$\mathfrak n/\mathfrak n^2$을 생성한다. 따라서 $w_1$은
$x\mapsto x+D'(x)$ 꼴이며, 여기서 $D'$은
$B/\mathfrak n^2$에서 $(B/\mathfrak n^2)$-가군
$\mathfrak n/\mathfrak n^2$으로 가는 미분이다. $w_1$은
$\mathfrak n/\mathfrak n^2$에서 항등사상이므로, 앞에서의
$\mathfrak n/\mathfrak n^2$과 $L$의 동일시에 의해 $D'$은 다시
$L$에서 $L$로 가는 미분으로 볼 수 있다. 끝으로 관계
$u'_1=w_1\circ u_1$은 $D'$이 $D$를 연장함을 뜻한다. 한편 $K$에서
$L$로 가는 모든 미분 $D$는 $c')$의 조건을 만족하는 준동형 $u'_1$에
대응한다 \sref{0.20.1.1-ko}. 그러므로 조건 $c')$은 $K$에서 $L$로
가는 모든 미분이 $L$에서 자기 자신으로 가는 미분으로 연장된다는
뜻이다. 즉 \sref{0.20.6.5-ko}에 따라 $L$은 $K$ 위에서 분리이다.

끝으로 $A$가 Cohen 환일 때 $b')$가 $c')$를 함의함을 증명하자.
먼저 $c')$의 가정 아래, $b')$의 가정을 만족하고 몫으로 옮기면
$u'_1:A/\mathfrak m^2\to B/\mathfrak n^2$을 주는 준동형
$u':A\to B$가 존재함을 보이자. 실제로 이는 합성 준동형
$v':A\to A/\mathfrak m^2\xrightarrow{u'_1}B/\mathfrak n^2$에
\sref{0.19.8.6-ko}, (i)를 적용하면 따른다.
\oldpage[0\textsubscript{IV-1}]{166}
따라서 이 합성은
$A\xrightarrow{u'}B\to B/\mathfrak n^2$으로 인수분해된다. 또한
$\operatorname{gr}_0(u'_1)$과 $\operatorname{gr}_1(u'_1)$에 관한
가정에서
$\operatorname{gr}_0(u')=\operatorname{gr}_0(u)=u_0$ 및
$\operatorname{gr}_1(u')=\operatorname{gr}_1(u)$가 따른다. 그러므로
$u'$에 의한 $A$의 균등화원의 상은 $B$의 균등화원이고, 따라서
$Bu'(\mathfrak m)=\mathfrak n$이다. 이제 필요한 자기동형사상 $w_1$을
얻기 위해서는 준동형 $u'$에 $b')$를 적용하여 얻는 $B$의
자기동형사상 $w$를 몫으로 옮겨 얻은 자기동형사상을 취하면 충분하다.
\end{proof}

\begin{remarks}[21.5.3]
\label{0.21.5.3-ko}
\begin{enumerate}
  \item[(i)] 체의 미분적 성질을 이용하면 완비 뇌터 국소환 $C$ 안의
  계수체의 \emph{유일성} 문제를 해결할 수 있다
  \sref{0.19.8.7-ko}, (ii). $\mathfrak J$를 $C$의 극대 아이디얼,
  $K=C/\mathfrak J$를 $C$의 잉여체라 하자. $\mathfrak J\neq0$인 경우만
  생각하면 충분하다. 증대 준동형 $C\to K$와 합성하면 항등사상이 되는
  준동형 $u:K\to C$가 존재한다고 하자. 이 준동형이 유일하기 위한
  필요충분조건은 $\Omega^1_K=0$인 것이다.

  실제로 $\Omega^1_K=0$이면 $K$는 그 소체 위에서 형식적으로
  비분기이다 \sref{0.20.7.4-ko}. 이 문제의 조건을 만족하는 두 번째
  준동형 $v\neq u$가 존재한다면, 모든 $x\in K$에 대한 $u(x)-v(x)$들의
  집합을 $\mathfrak J^n$이 포함하게 하는 가장 큰 정수 $n$이 존재한다.
  몫으로 옮기면 $u$와 $v$로부터 서로 다른 두 준동형
  $u_1,v_1:K\to C/\mathfrak J^{n+1}$이 얻어지며, 이들과
  $C/\mathfrak J^{n+1}\to C/\mathfrak J^n$의 합성은 서로 같다. 이는
  정의 \sref{0.19.10.2-ko}에 모순이다.

  역으로 $\Omega^1_K\neq0$이라고 하자. 그러면 $K$에서
  $\mathfrak J/\mathfrak J^2$으로 가는 영이 아닌 미분 $D$가 존재한다
  \sref{0.20.4.8-ko}. 따라서 $u$를 몫으로 옮겨 얻은 준동형을
  $u_1:K\to C/\mathfrak J^2$이라 할 때,
  $v_1=u_1+D$인 준동형 $v_1:K\to C/\mathfrak J^2$이 존재한다
  \sref{0.20.1.1-ko}. $k$를 $K$의 소체라 하면 $C$는 $k$-대수이고
  $K$는 형식적으로 매끄러운 $k$-대수이다 \sref{0.19.6.1-ko}.
  $u_1$과 $v_1$의 $k$에 대한 제한은 같으므로 $v_1$은
  $K\xrightarrow{v}C\to C/\mathfrak J^2$으로 인수분해되며,
  $v\neq u$이다.

  조건 $\Omega^1_K=0$은 $K$의 표수가 $0$이 아니면 $K$가
  완전체임을, 표수가 $0$이면 $K$가 $\mathbf Q$의 대수적 확대체임을
  뜻한다는 것을 상기하자
  (\sref{0.20.6.20-ko} 및 \sref{0.21.4.4-ko}).

  \item[(ii)] 마찬가지로 $W$를 Cohen 환, $C$를 완비 뇌터 국소환,
  $\mathfrak J\neq0$을 $C$의 극대 아이디얼에 포함되는 아이디얼이라
  하자. \sref{0.19.8.6-ko}, (i)의 인수분해
  $W\to C\to C/\mathfrak J$가 유일하기 위한 필요충분조건은 Cohen 환
  $W$의 잉여체 $K$가 $\Omega^1_K=0$을 만족하는 것이다.

  실제로 $\Omega^1_K\neq0$이고 $\mathfrak m$이 $W$의 극대
  아이디얼이면, $K$에서 $\mathfrak J/\mathfrak m\mathfrak J$으로 가는
  영이 아닌 미분 $D$를 증대 준동형 $W\to K$와 합성하여, $W$에서
  $\mathfrak J/\mathfrak m\mathfrak J$으로 가는 영이 아닌 미분 $D'$을
  얻으면 충분하다. 그러면 (i)에서와 같이 $D'$을 이용하여, $u$를
  몫으로 옮겨 얻은 준동형
  $u_1:W\to C/\mathfrak m\mathfrak J$와 서로 다른 준동형
  $v_1:W\to C/\mathfrak m\mathfrak J$을 만든다. 끝으로 이번에는
  \sref{0.19.8.6-ko}, (i)를 적용한다.

  반대로 $\Omega^1_K=0$이라고 하자. $W=K$가 표수 $0$인 체이면
  $v$의 유일성은 이미 (i)에서 따른다. 그 밖의 경우에는
  $\widehat{\Omega}^1_W=0$이다. 실제로 이때
  $\mathfrak m=pW$이고, $p$는 $K$의 표수이다
  \sref{0.19.8.5-ko}. 따라서 표준 준동형
  $\mathfrak m/\mathfrak m^2\to
  \Omega^1_W/\mathfrak m\Omega^1_W$
  \sref{0.20.5.11.2-ko}은 영이다. $W$와
  $K=W/\mathfrak m$에 완전열 \sref{0.20.5.12.1-ko}을 적용하면
  $\Omega^1_W=\mathfrak m\Omega^1_W$이므로 위의 주장을 얻는다.
  그런데 그러면 \sref{0.20.7.4-ko}에 의해 $W$는 그 아딕 위상에
  관하여 $\mathbf Z$ 위에서 형식적으로 비분기이고, $v$의 유일성은
  (i)에서와 같이 증명된다.
\end{enumerate}
\end{remarks}

\oldpage[0\textsubscript{IV-1}]{167}
\subsection{확대체에 대한 허용체}
\label{section:0.21.6-ko}

\begin{env}[21.6.1]
\label{0.21.6.1-ko}
네 체 $k_0\subset k\subset K\subset L$이 주어지면
\sref{0.20.6.16-ko}과 \sref{0.20.6.17-ko}에 의해 완전열
\begin{equation}
\tag{21.6.1.1}
\label{0.21.6.1.1-ko}
0\to\Upsilon_{K/k/k_0}\otimes_KL\xrightarrow{v}\Upsilon_{L/k/k_0}
\xrightarrow{u}\Upsilon_{L/K/k_0}\xrightarrow{s}\Upsilon_{L/K/k}\to0
\end{equation}
을 얻는다.

$k_0$, $K$, $L$을 고정하고 $k_0$와 $K$ 사이의 중간체 $k$를
``변화''시키면, $k=k_0$일 때에는 명백히
$\Upsilon_{L/K/k_0}=\Upsilon_{L/K/k}$이다. \sref{0.21.6.1.1-ko}의
표준 준동형 $s$가 여전히 전단사일 때, $k$를 $K$의 확대체 $L$에
대한 \emph{$k_0$-허용체}라고 한다. 어떤 조건 아래 이와 같은 체 $k$가
존재하면서도 $K$에 ``충분히 가까운'' 것, 예를 들어 $[K:k]$가 유한인
것이 중요한 까닭은 다음과 같다. 이 체들을 이용하면 몇몇 문제에서
미분 가군 $\Omega^1_{K/k_0}$와 $\Omega^1_{L/k_0}$를, 예를 들어
$k_0$가 소체일 때에는 ``너무 클'' 수 있는 이 가군들을, 더 다루기 쉬운
$\Omega^1_{K/k}$와 $\Omega^1_{L/k}$로 바꿀 수 있다.

$k_0$가 \emph{소체}이면 ``$k_0$-허용체'' 대신 ``허용체''라고 한다.

다음과 같이 놓는다.
\end{env}

\begin{env}[21.6.1.2]
\label{0.21.6.1.2-ko}
\[
\Delta(L/K,k/k_0)=\operatorname{Coker}
(\Upsilon_{K/k/k_0}\otimes_KL\to\Upsilon_{L/k/k_0})
\simeq\operatorname{Ker}(\Upsilon_{L/K/k_0}\to\Upsilon_{L/K/k})
\]
이는 $L$-벡터공간이다. 그 계수를 $d(L/K,k/k_0)$로 쓰고
\emph{$K$의 확대체 $L$에 대한 $k$의 $k_0$-허용성 결함}이라고 한다.
이는 $k$가 이 확대체에 대한 $k_0$-허용체일 필요충분조건이 영인
것임이 명백하다. $k_0$가 소체이면 $\Delta(L/K,k/k_0)$와
$d(L/K,k/k_0)$ 대신 $\Delta(L/K,k)$와 $d(L/K,k)$를 쓴다.
\end{env}

\begin{proposition}[21.6.2]
\label{0.21.6.2-ko}
$k_0\subset k\subset K\subset L$을 네 체라 하자.
\begin{enumerate}
  \item[\rm (i)] 다음 조건들은 동치이다.
  \begin{enumerate}
    \item[\rm a)] 체 $k$는 $K$의 확대체 $L$에 대한 $k_0$-허용체이다.
    달리 말하면, 준동형
    $\Upsilon_{L/K/k_0}\to\Upsilon_{L/K/k}$는 단사이고, 따라서
    전단사이다.
    \item[\rm b)] 표준 준동형
    $u:\Upsilon_{K/k/k_0}\to\Upsilon_{L/k/k_0}$는 영이다.
    \item[\rm c)] 표준 준동형
    $v:\Upsilon_{K/k/k_0}\otimes_KL\to\Upsilon_{L/k/k_0}$는 전사이고,
    따라서 전단사이다.
    \item[\rm d)] $d(L/K,k/k_0)=0$이다. 또는
    $\Delta(L/K,k/k_0)=0$이다.
  \end{enumerate}
  \item[\rm (ii)] 다음 경우에는 (i)의 동치 조건들이 성립한다.
  $\alpha)$ $L$은 $k$ 위에서 분리이다.
  $\beta)$ $L$은 $K$ 위에서 분리이다.
  $\gamma)$ $k\subset k_0(K^p)$이다. 여기서 $p$는 $k_0$의 표수
  지수이다.
\end{enumerate}
\end{proposition}

\begin{proof}
\begin{enumerate}
  \item[(i)] 각 주장은 완전열 \sref{0.21.6.1.1-ko}의 완전성에서
  자명하게 따른다.
  \item[(ii)] $L$이 $k$ 위에서 분리이면
  $\Upsilon_{L/k/k_0}=0$이므로 \sref{0.20.6.19-ko}, (i)의 조건 b)가
  성립한다. $L$이 $K$ 위에서 분리이면
  $\Upsilon_{L/K/k_0}=0$이므로 \sref{0.20.6.19-ko}, (i)의 조건 a)가
  성립한다. 끝으로 $k\subset k_0(K^p)$이면
  \sref{0.21.1.5.1-ko}에 의해
  $\Upsilon_{L/K/k_0}=\Upsilon_{L/K/k}$이므로 (i)의 조건 a)가 성립한다.
\end{enumerate}
\end{proof}

\oldpage[0\textsubscript{IV-1}]{168}
\begin{env}[21.6.3]
\label{0.21.6.3-ko}
체의 단사 준동형들로 이루어진 다음 가환 도식이 있다고 하자.
\[
\xymatrix{
 k'_0\ar[r]&k'\ar[r]&K'\ar[r]&L'\\
 k_0\ar[u]\ar[r]&k\ar[u]\ar[r]&K\ar[u]\ar[r]&L\ar[u]
}
\]
그러면 \sref{0.20.6.17-ko}, (ii)에 의해 표준 준동형
\begin{equation}
\tag{21.6.3.1}
\label{0.21.6.3.1-ko}
\Delta(L/K,k/k_0)\longrightarrow\Delta(L'/K',k'/k'_0)
\end{equation}
을 얻는다. 이는 명백한 추이성을 가지므로 $\Delta(L/K,k/k_0)$가
사중항 $(k_0,k,K,L)$의 \emph{함자}라고 말할 수 있다.
\end{env}

\begin{proposition}[21.6.4]
\label{0.21.6.4-ko}
\begin{enumerate}
  \item[\rm (i)] $k\subset k'\subset k''\subset K\subset L$을 다섯 체라
  하자. 표준 준동형의 완전열
\begin{equation}
\tag{21.6.4.1}
\label{0.21.6.4.1-ko}
0\to\Delta(L/K,k'/k)\to\Delta(L/K,k''/k)\to
\Delta(L/K,k''/k')\to0
\end{equation}
을 얻고, 따라서 등식
\begin{equation}
\tag{21.6.4.2}
\label{0.21.6.4.2-ko}
d(L/K,k''/k)=d(L/K,k''/k')+d(L/K,k'/k)
\end{equation}
을 얻는다.
  \item[\rm (ii)] $k_0\subset k\subset K\subset L\subset M$을 다섯 체라
  하자. 표준 준동형의 완전열
\begin{equation}
\tag{21.6.4.3}
\label{0.21.6.4.3-ko}
0\to\Delta(L/K,k/k_0)\otimes_LM\to\Delta(M/K,k/k_0)
\to\Delta(M/L,k/k_0)\to0
\end{equation}
을 얻고, 따라서 등식
\begin{equation}
\tag{21.6.4.4}
\label{0.21.6.4.4-ko}
d(M/K,k/k_0)=d(M/L,k/k_0)+d(L/K,k/k_0)
\end{equation}
을 얻는다.
\end{enumerate}
\end{proposition}

\begin{proof}
\begin{enumerate}
  \item[(i)] 가환 도식
\[
\xymatrix{
0\ar[r]&\Upsilon_{K/k'/k}\otimes_KL\ar[r]\ar[d]&
\Upsilon_{K/k''/k}\otimes_KL\ar[r]\ar[d]&
\Upsilon_{K/k''/k'}\otimes_KL\ar[r]\ar[d]&0\\
0\ar[r]&\Upsilon_{L/k'/k}\ar[r]\ar[d]&
\Upsilon_{L/k''/k}\ar[r]\ar[d]&
\Upsilon_{L/k''/k'}\ar[r]\ar[d]&0\\
0\ar[r]&\Delta(L/K,k'/k)\ar[r]&\Delta(L/K,k''/k)\ar[r]&
\Delta(L/K,k''/k')\ar[r]&0
}
\]
을 생각하자. 세 행을 각각 \emph{복합체} $T_1,T_2,T_3$로 볼 수 있다.
완전열 \sref{0.21.6.1.1-ko}과 $Delta$의 정의
\sref{0.21.6.1.2-ko}에 의해 복합체의 완전열
$0\to T_1\to T_2\to T_3\to0$을 얻는다. 여기에
\oldpage[0\textsubscript{IV-1}]{169}
코호몰로지 완전열을 적용하자. \sref{0.21.6.1.1-ko}의 완전성에 의해
$T_1$과 $T_2$의 코호몰로지는 오직 하나의 같은 차수를 제외하면 영이고,
그 차수에서 두 코호몰로지 가군은 모두
$\Upsilon_{k''/k'/k}\otimes_{k''}L$과 같다. 이처럼 $T_1$과 $T_2$의
코호몰로지가 같으므로 $T_3$의 코호몰로지는 반드시 영이다. 이로써
(i)가 증명되었다.
  \item[(ii)] 마찬가지로 가환 도식
\[
\xymatrix{
0\ar[r]&\Delta(L/K,k/k_0)\otimes_LM\ar[r]\ar[d]&
\Delta(M/K,k/k_0)\ar[r]\ar[d]&\Delta(M/L,k/k_0)\ar[r]&0\\
0\ar[r]&\Upsilon_{L/K/k_0}\otimes_LM\ar[r]\ar[d]&
\Upsilon_{M/K/k_0}\ar[r]\ar[d]&\Upsilon_{M/L/k_0}\ar[r]&0\\
0\ar[r]&\Upsilon_{L/K/k}\otimes_LM\ar[r]&
\Upsilon_{M/K/k}\ar[r]&\Upsilon_{M/L/k}\ar[r]&0
}
\]
을 생각하자. 여기서도 세 행을 \emph{복합체} $T'_1,T'_2,T'_3$로
본다. 완전열 \sref{0.21.6.1.1-ko}과 $Delta$의 정의
\sref{0.21.6.1.2-ko}에 의해 이제 복합체의 완전열
$0\to T'_1\to T'_2\to T'_3\to0$을 얻으며, 여기에 다시 코호몰로지
완전열을 적용한다. 이번에는 \sref{0.21.6.1.1-ko}의 완전성에 의해
$T'_2$와 $T'_3$의 코호몰로지가 오직 하나의 같은 차수를 제외하면
영이고, 그 차수에서 두 코호몰로지 가군은 모두
$\Upsilon_{M/L/k_0}$와 같다. 따라서 $T'_1$의 코호몰로지가 영임을
얻으며, 이로써 (ii)가 증명되었다.
\end{enumerate}
\end{proof}

\begin{corollary}[21.6.5]
\label{0.21.6.5-ko}
\begin{enumerate}
  \item[(i)] 다섯 체 $k\subset k'\subset k''\subset K\subset L$이
  주어졌다고 하자. $k''$가 $K$의 확대체 $L$에 대한 $k$-허용체일
  필요충분조건은, $k'$가 그 확대체에 대한 $k$-허용체이고 동시에
  $k''$가 그 확대체에 대한 $k'$-허용체인 것이다.
  \item[(ii)] 다섯 체 $k_0\subset k\subset K\subset L\subset M$이
  주어졌다고 하자. $k$가 $K$의 확대체 $M$에 대한 $k_0$-허용체일
  필요충분조건은, $k$가 $K$의 확대체 $L$과 $L$의 확대체 $M$ 각각에
  대한 $k_0$-허용체인 것이다.
\end{enumerate}

이는 $d$의 값들이 $\geqslant0$이므로 관계식
\sref{0.21.6.4.2-ko}와 \sref{0.21.6.4.4-ko}에서 곧바로 따른다.
\end{corollary}

\begin{corollary}[21.6.6]
\label{0.21.6.6-ko}
$k_0\subset k\subset K\subset L$을 네 체라 하고, $k$가 $K$의 확대체
$L$에 대한 $k_0$-허용체라고 하자. 네 체 $k'_0,k',K',L'$이
\[
k_0\subset k'_0\subset k'\subset k\subset K\subset K'\subset L'
\subset L
\]
을 만족하면, $k'$은 $K'$의 확대체 $L'$에 대한 $k'_0$-허용체이다.
\end{corollary}

\subsection{카르티에 등식}
\label{subsection:0.21.7-ko}

다음 결과는 미분에 관한 MacLane의 정리를 미분 가군의 말로 옮겨 쓴
것이다.

\oldpage[0\textsubscript{IV-1}]{170}
\begin{theorem}[21.7.1]
\label{0.21.7.1-ko}
\textnormal{(카르티에).}
$K$를 체, $L$을 $K$의 유한형 확대체라 하자. 그러면
$\Omega^1_{L/K}$와 $\Upsilon_{L/K}$는 유한 차원의 $L$-벡터공간이고,
다음 등식이 성립한다.
\begin{equation}
\tag{21.7.1.1}
\label{0.21.7.1.1-ko}
\operatorname{rg}\Omega^1_{L/K}-\operatorname{rg}\Upsilon_{L/K}
=\deg.\operatorname{tr}_K L.
\end{equation}
\end{theorem}

\begin{proof}
$K\subset L'\subset L$을 만족하는 체 $L'$에 대하여 완전열
\sref{0.20.6.15.1-ko}을 얻는다. 여기서는
$\Lambda=P$를 $K$의 소체로 두고 $A=K$, $B=L'$, $C=L$로 적용한다.
이 완전열은
\[
0\to\Upsilon_{L'/K}\otimes_{L'}L\to\Upsilon_{L/K}
\to\Upsilon_{L/L'}\to\Omega^1_{L'/K}\otimes_{L'}L
\to\Omega^1_{L/K}\to\Omega^1_{L/L'}\to0
\]
이다. 따라서 ($\mathbf 0_{\rm III}$, 11.10.2)
\[
\begin{split}
\operatorname{rg}_L\Omega^1_{L/K}-\operatorname{rg}_L\Upsilon_{L/K}
={}&(\operatorname{rg}_L\Omega^1_{L/L'}
-\operatorname{rg}_L\Upsilon_{L/L'})\\
&+(\operatorname{rg}_{L'}\Omega^1_{L'/K}
-\operatorname{rg}_{L'}\Upsilon_{L'/K}).
\end{split}
\]
한편
$\deg.\operatorname{tr}_K L=\deg.\operatorname{tr}_K L'
+\deg.\operatorname{tr}_{L'}L$이다. 따라서 확대체 $L$의 생성원 수에
관한 귀납법으로, $L=K(x)$일 때의 \sref{0.21.7.1.1-ko}을 증명하면
충분하다. 이제 세 경우로 나눈다.
\begin{enumerate}
  \item[\rm a)] $x$가 $K$ 위에서 초월적이면 $L$은 $K$ 위에서
  분리이므로 $\Upsilon_{L/K}=0$이다 \sref{0.20.6.3-ko}. 한편
  \sref{0.20.5.9-ko}와 \sref{0.20.4.13-ko}에 의해
  $\Omega^1_{L/K}$의 계수는 $1$이다. 따라서 이 경우
  \sref{0.21.7.1.1-ko}을 얻는다.
  \item[\rm b)] $L$이 $K$의 대수적 분리 확대체이면 다시
  $\Upsilon_{L/K}=0$이다 \sref{0.20.6.3-ko}. 한편
  \sref{0.20.6.20-ko}에 의해 $\Omega^1_{L/K}=0$이므로 역시
  \sref{0.21.7.1.1-ko}을 얻는다.
  \item[\rm c)] $L$이 $K$의 대수적 비분리 확대체이면, 처음의 논법에
  의해 $L^p\subset K$인 경우만 생각하면 충분하다. 그러므로
  \sref{0.21.4.8-ko}에 의해
  $\operatorname{rg}_L\Upsilon_{L/K}
  =\operatorname{rg}_L\Omega^1_{L/K}$이고, 역시
  \sref{0.21.7.1.1-ko}을 얻는다.
\end{enumerate}
\end{proof}

\begin{corollary}[21.7.2]
\label{0.21.7.2-ko}
$K$를 체, $L$을 $K$의 유한형 확대체, $k$를 $K$의 부분체라 하자.
그러면 $\Omega^1_{L/K}$와 $\Upsilon_{L/K/k}$는 유한 차원의
$L$-벡터공간이고, 다음 등식이 성립한다.
\begin{equation}
\tag{21.7.2.1}
\label{0.21.7.2.1-ko}
\operatorname{rg}\Omega^1_{L/K}-\operatorname{rg}\Upsilon_{L/K/k}
=\deg.\operatorname{tr}_K L+d(L/K,k).
\end{equation}
따라서 부등식
\begin{equation}
\tag{21.7.2.2}
\label{0.21.7.2.2-ko}
\operatorname{rg}\Omega^1_{L/K}-\operatorname{rg}\Upsilon_{L/K/k}
\geqslant\deg.\operatorname{tr}_K L
\end{equation}
이 성립한다. 또한 \sref{0.21.7.2.2-ko}의 양변이 같을 필요충분조건은
$k$가 $K$의 확대체 $L$에 대한 허용체인 것이다.
\end{corollary}

\begin{proof}
실제로 준동형 $\Upsilon_{L/K}\to\Upsilon_{L/K/k}$가 전사이므로,
정의 \sref{0.21.6.1.2-ko}에 의해
\[
\operatorname{rg}\Upsilon_{L/K}-\operatorname{rg}\Upsilon_{L/K/k}
=d(L/K,k).
\]
따라서 따름정리는 \sref{0.21.7.1-ko}과 \sref{0.21.6.2-ko}에서
곧바로 따른다.
\end{proof}

\begin{corollary}[21.7.3]
\label{0.21.7.3-ko}
$k\subset K\subset L$을 세 체라 하고, $L$이 $k$의 유한형 확대체라고
하자. 그러면
\begin{equation}
\tag{21.7.3.1}
\label{0.21.7.3.1-ko}
\operatorname{rg}_L\Omega^1_{L/k}-\operatorname{rg}_K\Omega^1_{K/k}
=\deg.\operatorname{tr}_K L+d(L/K,k).
\end{equation}
\oldpage[0\textsubscript{IV-1}]{171}
따라서 부등식
\begin{equation}
\tag{21.7.3.2}
\label{0.21.7.3.2-ko}
\operatorname{rg}_L\Omega^1_{L/k}-\operatorname{rg}_K\Omega^1_{K/k}
\geqslant\deg.\operatorname{tr}_K L
\end{equation}
이 성립한다. 여기서 등호가 성립할 필요충분조건은 $k$가 $K$의
확대체 $L$에 대한 허용체인 것이다.
\end{corollary}

\begin{proof}
완전열 \sref{0.20.6.1.1-ko}
\[
0\to\Upsilon_{L/K/k}\to\Omega^1_{K/k}\otimes_KL\to\Omega^1_{L/k}
\to\Omega^1_{L/K}\to0
\]
이 성립한다. 따라서 \sref{0.21.7.3.1-ko}의 좌변은
\[
\operatorname{rg}_L\Omega^1_{L/k}-\operatorname{rg}_L\Upsilon_{L/K/k}
\]
와 같으며, \sref{0.21.7.2-ko}을 적용하면 된다.
\end{proof}

마지막 따름정리의 의의는 불완전성 가군을 제외하고 오직 미분 가군만을
사용한다는 데 있다. 더욱이 아래의 \sref{0.21.8.6-ko}에서, $K$의 모든
유한형 확대체 $L$에 대하여 $[K:k]$가 유한이고 $K$의 확대체 $L$에
대한 허용체인 $K$의 부분체 $k$가 존재함을 볼 것이다. 따라서 이때
\sref{0.21.7.3.2-ko}의 양변은 같다.

\begin{corollary}[21.7.4]
\label{0.21.7.4-ko}
$K$가 체 $k$의 유한형 확대체라고 하자. 다음 조건들은 동치이다.
\begin{enumerate}
  \item[\rm a)] $K$는 $k$의 유한 분리 확대체이다.
  \item[\rm b)] $K$는 형식적으로 비분기인 $k$-대수이다
  \sref{0.19.10.2-ko}.
  \item[\rm c)] $K$는 형식적으로 에탈인 $k$-대수이다
  \sref{0.19.10.2-ko}.
  \item[\rm d)] $\Omega^1_{K/k}=0$이다.
\end{enumerate}
\end{corollary}

\begin{proof}
위상들이 이산일 때 조건 b)와 d)가 동치임은 이미 알고 있다
\sref{0.20.7.4-ko}. 또한 a)는 $K$가 형식적으로 매끄러운 $k$-대수임을
함의하므로 \sref{0.19.6.1-ko}, a)와 b)의 논리곱은 c)와 동치이다.
따라서 d)가 a)를 함의함만 보이면 된다. 그런데
\sref{0.21.7.1.1-ko}에 의해 $\Omega^1_{K/k}=0$이면 $K$는 $k$ 위에서
대수적이고, 따라서 유한이며, $\Upsilon_{K/k}=0$이다. 이는
\sref{0.20.6.3-ko}에 의해 $K$가 $k$ 위에서 분리임을 뜻한다.
\end{proof}

\begin{remark}[21.7.5]
\label{0.21.7.5-ko}
\sref{0.20.6.5.2-ko}과 ($\mathbf 0_{\rm III}$, 11.10.2)에 의해
\sref{0.21.7.1.1-ko}의 좌변은 $C=L$, $B=K$, $A=P$ ($K$의 소체)일 때
\sref{0.20.6.5-ko}에서 도입한 복합체 $\mathrm K_\bullet(C/B/A)$의
\emph{오일러--푸앵카레 특성}과 다름없다. 교차 이론과
리만--로흐 정리를 다루는 장에서는, \sref{0.20.6.22-ko}에서 생각한
복합체 $\mathrm F_\bullet(C/A)$를 일반화하는 복합체의 일반화된
오일러--푸앵카레 특성도 중요한 역할을 할 것이다. 이 특성들은
$\mathcal O_X$-가군의 동치류들로 이루어진 군에 값을 갖는다.
\end{remark}

\subsection{허용성 판정법}
\label{subsection:0.21.8-ko}

앞의 규약으로 돌아가, 이 절에서 생각하는 모든 체의 표수가 $p>0$이라고
가정한다.

\begin{lemma}[21.8.1]
\label{0.21.8.1-ko}
$K$를 체, $k$를 $K$의 부분체라 하고, $(k_\alpha)_{\alpha\in I}$를
$k=\bigcap_{\alpha\in I}k_\alpha$를 만족하는 $K$의 부분체들의 감소하는
여과족이라 하자. $V$를 $K$-벡터공간,
$(a_i)_{1\leqslant i\leqslant n}$을 $V$의 유한 벡터족이라 하자.
이 벡터족이 $k$ 위에서 선형 독립이면, 어떤 지표 $\gamma$에 대해서는
$k_\gamma$ 위에서도 선형 독립이다.
\end{lemma}

\oldpage[0\textsubscript{IV-1}]{172}
\begin{proof}
$(a_i)_{1\leqslant i\leqslant n}$의 $K$ 위 계수를 $r$이라 하고,
$n-r$에 관하여 귀납하자. $n=r$이면 명제는 자명하다. 실제로 이때
$(a_i)_{1\leqslant i\leqslant r}$은 $K$ 위에서 선형 독립이고, 따라서
$K$의 모든 부분체 위에서도 선형 독립이다. 예를 들어
$(a_i)_{1\leqslant i\leqslant r}$이 $K$ 위에서 선형 독립이라고 하고
\[
a_{r+1}=\sum_{i=1}^r\lambda_i a_i,\qquad \lambda_i\in K
\]
로 쓰자. $(a_i)_{1\leqslant i\leqslant r+1}$이 $k$ 위에서 선형
독립이므로 모든 $\lambda_i$가 $k$에 속할 수는 없다. 예를 들어
$\lambda_{i_0}\notin k$라고 하자. 그러면 어떤 지표 $\beta$에 대해서는
$\lambda_{i_0}\notin k_\beta$이다. 따라서
$(a_i)_{1\leqslant i\leqslant r+1}$은 $k_\beta$ 위에서 선형 독립이다.
실제로 $(a_i)_{1\leqslant i\leqslant r}$은 $K$의 모든 부분체 위에서
선형 독립이므로, 만일 앞의 $r+1$개 벡터가 $k_\beta$ 위에서 선형
독립이 아니라면 $a_{r+1}$은 $a_i$들의 $k_\beta$-계수 선형결합과
같다. 그 계수들은 반드시 $\lambda_i$이므로 모순이다. 이제 $K$를
$k_\beta$로 바꾸고, $(k_\alpha)_{\alpha\in I}$를 $k_\beta$에 포함되는
$k_\alpha$들로 이루어진 부분족으로 바꾸어 귀납가정을 적용하면 된다.
\end{proof}

\begin{lemma}[21.8.2]
\label{0.21.8.2-ko}
$K$를 체, $p$를 그 표수 지수, $k_0$를 $K$의 부분체라 하고,
$(k_\alpha)_{\alpha\in I}$를
\[
\bigcap_{\alpha\in I}k_\alpha(K^p)=k_0(K^p)
\]
를 만족하는 $K$의 부분체들의 감소하는 여과족이라 하자.
$(x_i)_{1\leqslant i\leqslant n}$이 $k_0$ 위에서 $p$-자유인
$K$의 원소들의 유한족이면, 어떤 지표 $\alpha$에 대해서는
$(x_i)$가 $k_\alpha$ 위에서도 $p$-자유이다.

실제로 $(x_i)$가 $K$의 부분체 $k$ 위에서 $p$-자유라는 것은
$0\leqslant m(i)<p$인 단항식들의 유한족
\[
\prod_{i=1}^n x_i^{m(i)}
\]
이 $k(K^p)$ 위에서 선형 독립이라는 뜻이다. 따라서 벡터공간
$V=K$ 안의 이 단항식족과 $K$의 부분체 $k_\alpha(K^p)$ 및
$k_0(K^p)$에 보조정리 \sref{0.21.8.1-ko}을 적용하면 충분하다.
\end{lemma}

\begin{theorem}[21.8.3]
\label{0.21.8.3-ko}
$K$를 표수가 $p>0$인 체, $k_0$를 $K$의 부분체,
$(k_\alpha)_{\alpha\in I}$를 $k_0$를 포함하는 $K$의 부분체들의
감소하는 여과족이라 하자. 다음 조건들은 동치이다.
\begin{enumerate}
  \item[\rm a)] $\displaystyle\bigcap_{\alpha\in I}k_\alpha(K^p)
  =k_0(K^p)$이다.
  \item[\rm b)] $\Upsilon_{L/K/k_0}$가 유한 차원의 $L$-벡터공간인
  $K$의 모든 확대체 $L$에 대하여, 어떤 $\alpha\in I$가 존재하여
  $k_\alpha$가 $K$의 확대체 $L$에 대한 $k_0$-허용체이다. 특히
  \sref{0.21.7.2-ko}에 의해 $L$이 유한형 확대체이면 이 유한 차원
  조건은 성립한다.
  \item[\rm b$'$)] $x^p\in K$인 $K$의 모든 확대체 $L=K(x)$에 대하여,
  어떤 $\alpha\in I$가 존재하여 $k_\alpha$가 $K$의 확대체 $L$에
  대한 $k_0$-허용체이다.
  \item[\rm c)] 표준 사상
\begin{equation}
\tag{21.8.3.1}
\label{0.21.8.3.1-ko}
\Omega^1_{K/k_0}\longrightarrow
\varprojlim_{\alpha}\Omega^1_{K/k_\alpha}
\end{equation}
  은 단사이다.
\end{enumerate}
\end{theorem}

\begin{proof}
표준 사상 \sref{0.21.8.3.1-ko}은 물론 준동형들의 사영계
\[
\Omega^1_{K/k_0}\longrightarrow\Omega^1_{K/k_\alpha}
\]
\sref{0.20.5.3.3-ko}에서 사영극한을 취하여 얻는다. 다음 논리적
도식에 따라 정리를 증명하겠다.
\[
c)\Rightarrow b)\Rightarrow b')\Rightarrow a)\Rightarrow c).
\]

$k_\alpha$가 $K$의 확대체 $L$에 대한 $k_0$-허용체라는 것은 표준
준동형
\oldpage[0\textsubscript{IV-1}]{173}
\[
\Upsilon_{L/K/k_0}\longrightarrow\Upsilon_{L/K/k_\alpha}
\]
이 단사라는 뜻이다. 이 준동형의 핵을 $N_\alpha$라 하면,
$N_\alpha$들은 $\Upsilon_{L/K/k_0}$의 벡터 부분공간들의 감소하는
여과계를 이룬다. 가정에 의해 $\Upsilon_{L/K/k_0}$의 차원이
유한이므로, 어떤 $N_\alpha$가 영이라는 것과 이들의 교집합이
영이라는 것은 동치이다. 그런데 이 교집합은 사영극한 준동형
\[
\Upsilon_{L/K/k_0}\longrightarrow
\varprojlim_\alpha\Upsilon_{L/K/k_\alpha}
\]
의 핵과 다름없다. 이제 가환도식
\[
\xymatrix{
\Upsilon_{L/K/k_0}\ar[r]\ar[d]&
\varprojlim\Upsilon_{L/K/k_\alpha}\ar[d]\\
\Omega^1_{K/k_0}\otimes_KL\ar[r]&
\varprojlim(\Omega^1_{K/k_\alpha}\otimes_KL)
}
\]
을 얻으며, 정의에 의해 왼쪽 세로 화살표는 단사이다. 따라서 $c)$가
$b)$를 함의함을 보이려면 모든 $K$-벡터공간 $V$에 대하여 표준 사상
\begin{equation}
\tag{21.8.3.2}
\label{0.21.8.3.2-ko}
\Omega^1_{K/k_0}\otimes_KV\longrightarrow
\varprojlim_{\alpha}(\Omega^1_{K/k_\alpha}\otimes_KV)
\end{equation}
이 단사임을 보이면 충분하다. 특히 $V=L$인 경우에도 그러하다.
$V=K^n$이면 모든 $K$-벡터공간 $W$에 대하여
$W\otimes_KV=W^n$이고 곱과 사영극한이 가환하므로 이는 명백하다.
한편 \sref{0.21.8.3.2-ko}의 왼쪽 항의 임의의 원소 $z$에 대하여,
계수가 유한인 $V$의 벡터 부분공간 $V'$가 존재하여
\[
z\in\Omega^1_{K/k_0}\otimes_KV'
\]
이다. 여기서 오른쪽 항을 $\Omega^1_{K/k_0}\otimes_KV$의 벡터
부분공간과 표준적으로 동일시한다. $z\ne0$이면
\[
\varprojlim_\alpha(\Omega^1_{K/k_\alpha}\otimes_KV')
\]
에서 그 상도 영이 아니다. 사영극한 함자는 왼쪽 완전이므로,
\sref{0.21.8.3.2-ko}의 오른쪽 항에서 $z$의 상도 영이 아니다.
이로써 $c)$가 $b)$를 함의함을 보였다.

$b)$가 $b')$을 함의함은 자명하다. 이제 $b')$가 $a)$를 함의함을
보이자. $b')$의 표기에서 $L\ne K$라고 가정할 수 있다. 그러면
\sref{0.21.4.7-ko}에 의해
\[
\operatorname{rg}\Upsilon_{L/K}\leqslant1
\]
이다. $\Upsilon_{L/K/k_\alpha}=0$이면 \sref{0.21.6.1.1-ko}에 의해
증명할 것이 없다. 그렇지 않으면 $\Upsilon_{L/K/k_\alpha}$는
$\Upsilon_{L/K}$와 표준적으로 동일시된다. $a=x^p$라 놓으면
$\Upsilon_{L/K}$는 하나의 원소
$d_K(a)\otimes1$을 기저로 갖는다 \sref{0.21.4.7-ko}. 따라서
$\Upsilon_{L/K/k_\alpha}$는 하나의 원소
$d_{K/k_\alpha}(a)\otimes1$을 기저로 갖는다. 그러므로 $k_\alpha$가
$K$의 확대체 $L$에 대한 $k_0$-허용체라는 것은
\[
d_{K/k_\alpha}(a)\ne0
\]
이라는 뜻이고, 이는 다시 \sref{0.21.4.5-ko}에 의해
$a\notin k_\alpha(K^p)$라는 뜻이다. 그런데 모든
$a\notin k_0(K^p)$에 대하여 $d_{K/k_0}(a)\ne0$이고
$K(a^{1/p})\ne K$이다. 따라서 $b')$을 적용하면
$a\notin k_\alpha(K^p)$인 $\alpha$가 존재한다. 다시 말해 $b')$는
$a)$를 함의한다.

이제 $a)$가 $c)$를 함의함을 보이면 된다. $(x_\mu)_{\mu\in M}$을
$k_0$ 위에서 $K$의 $p$-기저라 하자. 그러면
$d_{K/k_0}(x_\mu)$들은 $K$-벡터공간 $\Omega^1_{K/k_0}$의 기저를
이룬다 \sref{0.21.4.5-ko}. 조건 $c)$는 지표집합 $M$의 모든 유한
부분집합 $J$에 대하여 어떤 $\alpha\in I$가 존재하여,
$\mu\in J$인 $d_{K/k_\alpha}(x_\mu)$들이
$\Omega^1_{K/k_\alpha}$에서 선형 독립이라는 뜻이다. 그런데 이는
또한 \sref{0.21.4.5-ko}에 의해 $\mu\in J$인 $x_\mu$들이
$k_\alpha$ 위에서 $p$-자유족을 이룬다는 뜻이다. 이 성질을 갖는
$\alpha$의 존재는 보조정리 \sref{0.21.8.2-ko}과 가정 $a)$에서
따른다.
\end{proof}

\oldpage[0\textsubscript{IV-1}]{174}

\begin{corollary}[21.8.4]
\label{0.21.8.4-ko}
$K$를 표수가 $p>0$인 체, $k_0$를 $K$의 부분체,
$(k_\alpha)_{\alpha\in I}$를 $k_0$를 포함하는 $K$의 부분체들의
감소하는 여과족이라 하고,
\[
\bigcap_{\alpha\in I}k_\alpha(K^p)=k_0(K^p)
\]
라고 하자. $L$을 $K$의 확대체라 하고, $\Upsilon_{L/K/k_0}$가
유한 차원의 $L$-벡터공간이라고 하자. 그러면
$K\subset K'\subset L$인 모든 체 $K'$에 대하여, 어떤
$\alpha\in I$가 존재하여 $k_\alpha$는 $K'$의 확대체 $L$에 대한
$k_0$-허용체이다.
\end{corollary}

\begin{proof}
\sref{0.21.8.3-ko}과 \sref{0.21.6.6-ko}을 적용하면 충분하다.
\end{proof}

\begin{corollary}[21.8.5]
\label{0.21.8.5-ko}
$K$를 표수가 $p>0$인 체, $k_0$를 $K$의 부분체,
$(k_\alpha)_{\alpha\in I}$를 $k_0$를 포함하는 $K$의 부분체들의
감소하는 여과족이라 하고,
\[
\bigcap_{\alpha\in I}k_\alpha(K^p)=k_0(K^p)
\]
라고 하자. 그러면 $\Upsilon_{L/K/k_0}$가 유한 차원의
$L$-벡터공간인 $K$의 모든 확대체 $L$에 대하여
\[
\bigcap_{\alpha\in I}k_\alpha(L^p)=k_0(L^p)
\]
이다.
\end{corollary}

\begin{proof}
$M$을 $L$의 확대체라 하고, $\Upsilon_{M/L/k_0}$가 유한 차원의
$M$-벡터공간이라고 하자. 완전열 \sref{0.21.6.1.1-ko}
\[
0\longrightarrow\Upsilon_{L/K/k_0}\otimes_LM
\longrightarrow\Upsilon_{M/K/k_0}
\longrightarrow\Upsilon_{M/L/k_0}
\]
과 가정에 의해 $\Upsilon_{M/K/k_0}$도 유한 차원의
$M$-벡터공간이다. 따라서 어떤 지표 $\alpha$가 존재하여
$k_\alpha$는 $K$의 확대체 $M$에 대한 $k_0$-허용체이다
\sref{0.21.8.3-ko}. 그러므로 $k_\alpha$는 $L$의 확대체 $M$에
대해서도 $k_0$-허용체이다 \sref{0.21.6.6-ko}. 이 결론은
$\Upsilon_{M/L/k_0}$가 유한 차원인 $L$의 모든 확대체 $M$에 대하여
성립하므로, 따름정리는 \sref{0.21.8.3-ko}의 조건 a)와 b)의
동치에서 따른다.
\end{proof}

\begin{corollary}[21.8.6]
\label{0.21.8.6-ko}
$K$를 체, $p$를 그 표수 지수, $k_0$를 $K$의 부분체라 하자.
$L$이 $K$의 확대체이고 $\Upsilon_{L/K/k_0}$가 유한 차원의
$L$-벡터공간이면, $k_0(K^p)$를 포함하고 $[K:k]$가 유한이며
$K$의 확대체 $L$에 대한 $k_0$-허용체인 $K$의 부분체 $k$가
존재한다.
\end{corollary}

\begin{proof}
\sref{0.21.8.3-ko}에 의해, $K^p$와 $k_0$를 포함하고
$[K:k_\alpha]<+\infty$이며
\[
\bigcap_\alpha k_\alpha=k_0(K^p)
\]
인 $K$의 부분체들의 감소하는 여과족
$(k_\alpha)_{\alpha\in I}$를 구성하면 충분하다. 이를 위하여
$(x_\lambda)_{\lambda\in J}$를 $k_0$ 위에서 $K$의 $p$-기저라 하자.
$J$의 모든 유한 부분집합 $H$에 대하여, $k_0(K^p)$와
$\lambda\in J-H$인 $x_\lambda$들로 생성되는 $K$의 부분체를
$k_H$라 하자. 그러면 정의로부터
$(x_\lambda)_{\lambda\in H}$가 $k_H$ 위에서 $K$의 $p$-기저임을
얻으며, 따라서 $k_H$들은 곧바로 필요한 조건들을 만족한다.
\end{proof}

\begin{remarks}[21.8.7]
\label{0.21.8.7-ko}
\begin{enumerate}
  \item[(i)] $L$이 $K$의 \emph{유한형} 확대체이면 모든 $K$의
  부분체 $k_0$에 대하여 $\Upsilon_{L/K/k_0}$의 차원이 유한임은 이미
  \sref{0.21.7.2-ko}에서 보았다. $L$이 $K$의 \emph{분리}
  확대체인 경우에도 마찬가지이다. 실제로
  \sref{0.20.6.19-ko}에 의해 $\Upsilon_{L/K/k_0}=0$이다. 마지막으로,
  $L$이 $K$의 분리 확대체 $L_0$의 \emph{유한형 확대체}이면,
  \sref{0.21.8.5-ko}에서와 같은 논증으로
  $\Upsilon_{L/K/k_0}$의 차원이 여전히 유한임을 얻는다. 사실 이는
  $\Upsilon_{L/L_0/k_0}$의 벡터 부분공간과 동형이다.

  \oldpage[0\textsubscript{IV-1}]{175}
  \item[(ii)] \sref{0.21.8.5-ko}의 명제에서, 따라서
  \sref{0.21.8.3-ko}의 조건 b)에서도 $\Upsilon_{L/K/k_0}$가
  $L$ 위에서 유한 차원이라는 가정을 생략할 수 없다. $k_0$를 소체
  $\mathbf F_p$라 하고, $[K:K^p]$가 가산 무한인 체 $K$를 택하자.
  예를 들어 가산 무한 개의 부정원에 대한 유리 함수체
  $\mathbf F_p(X_1,\ldots,X_n,\ldots)$를 택할 수 있다.
  \sref{0.21.8.6-ko}의 구성법에 의해, $K_0=K$이고
  $\bigcap_nK_n=K^p$인 $K$의 부분체들의 무한한 강한 감소열
  $(K_n)$이 존재한다. 이제 $M$을 $M_n$들의 합집합이라 할 때
  $K$의 확대체 $L=M^{1/p}$가 \sref{0.21.8.5-ko}의 반례가 되도록,
  $K^p$의 유한 확대체 $M_n\subset K$들의 증가열을 구성하겠다.
  $x\in K-K^p$를 택하고, $M_0=K^p$라 놓으며, $n\geqslant1$에
  대하여
  \[
  M_n=K^p(a_1x,a_2x,\ldots,a_nx)
  \]
  라 하자. 여기서 $a_n$들은 모든 $n$에 대하여
  $a_n\in K_n$이고 $a_{n+1}\notin M_n(x)$가 되도록 귀납적으로
  구성한다. 이는 $M_n(x)$가 $K^p$ 위에서 유한 차수인 반면
  $K_n$은 그렇지 않으므로 가능하다. 그러면 곧바로
  $x\notin M=L^p$임을 얻는다. 그러나
  $x=(a_nx)a_n^{-1}$이므로 모든 $n$에 대하여
  $x\in K_n(M)=K_n(L^p)$이다. 따라서
  $x\in\bigcap_nK_n(L^p)$이고, 이것이 필요한 반례이다.
\end{enumerate}
\end{remarks}

\begin{proposition}[21.8.8]
\label{0.21.8.8-ko}
$k$를 표수가 $p>0$인 체라 하고,
\[
A=k[[T_1,\ldots,T_r]],\qquad K=k((T_1,\ldots,T_r))
\]
라 하자. 여기서 $A$는 $k$ 위의 $r$개 부정원에 대한 형식적
멱급수환이고, $K$는 그 분수체이다. 그러면 $A$의 뇌터 부분환들의
감소하는 여과족 $(A_\alpha)$가 존재하여, 모든 $\alpha$에 대하여
$A$는 유한형 자유 $A_\alpha$-가군이고, $K_\alpha$를
$A_\alpha$의 분수체라 할 때
\[
\bigcap_\alpha K_\alpha=K^p
\]
이다.
\end{proposition}

\begin{proof}
다음과 같이 쓸 수 있다.
\[
K^p=k^p((T_1^p,\ldots,T_r^p)).
\]
\sref{0.21.8.6-ko}의 증명에서, 모든 $\alpha$에 대하여
$[k:k_\alpha]$가 유한이고
$\bigcap_\alpha k_\alpha=k^p$인 $k$의 부분체들의 감소족
$(k_\alpha)$가 존재함을 보았다. 이제
\[
A_\alpha=k_\alpha[[T_1^p,\ldots,T_r^p]]
\]
라 놓으면 $A$가 유한형 자유 $A_\alpha$-가군임은 명백하다.
따라서 다음 관계를 증명하면 충분하다.
\begin{equation}
\tag{21.8.8.1}
\label{0.21.8.8.1-ko}
\bigcap_\alpha k_\alpha((T_1^p,\ldots,T_r^p))
=k^p((T_1^p,\ldots,T_r^p)).
\end{equation}
$\bigcap_\alpha k_\alpha=k^p$이므로, 이는 다음 두 보조정리에서
따를 것이다.
\end{proof}

\begin{lemma}[21.8.8.2]
\label{0.21.8.8.2-ko}
$k'$이 체 $k$의 확대체이면
\[
k'[[T_1,\ldots,T_r]]\cap k((T_1,\ldots,T_r))
=k[[T_1,\ldots,T_r]]
\]
이다.
\end{lemma}

\begin{proof}
$C=k[[T_1,\ldots,T_r]]$, $D=k'[[T_1,\ldots,T_r]]$라 놓자.
$k((T_1,\ldots,T_r))$는 $C$의 분수체이므로, $D$가 충실히 평탄한
$C$-가군임을 보이면 충분하다 (Bourbaki, \emph{Alg. comm.}, chap.~I,
\S~3, n\textsuperscript{o}~5, prop.~10). 그런데 $C$와 $D$는 뇌터
국소환이고, $\mathfrak m$을 $C$의 극대 아이디얼이라 하면
\[
D/\mathfrak m^jD=(C/\mathfrak m^j)\otimes_k k'
\]
이다. 따라서 $D/\mathfrak m^jD$는 평탄한
$(C/\mathfrak m^j)$-가군이다. 이제 ($\mathbf 0_{\mathrm{III}}$,
10.2.1)과 ($\mathbf 0_{\mathrm I}$, 6.6.2)을 적용하면 충분하다.
\end{proof}

\begin{lemma}[21.8.8.3]
\label{0.21.8.8.3-ko}
$k$를 체, $(k_\alpha)$를 $k$의 부분체들의 감소하는 여과족이라 하고
\[
k_0=\bigcap_\alpha k_\alpha
\]
라 놓자. $k$의 표수 지수의 어떤 거듭제곱 $q$가 존재하여
$k^q\subset k_0$라고 가정하자. 그러면
\begin{equation}
\tag{21.8.8.4}
\label{0.21.8.8.4-ko}
\bigcap_\alpha k_\alpha((T_1,\ldots,T_r))
=k_0((T_1,\ldots,T_r))
\end{equation}
이다.
\end{lemma}

\begin{proof}
\oldpage[0\textsubscript{IV-1}]{176}
\sref{0.21.8.8.4-ko}의 좌변의 영이 아닌 원소 $f$가 오른쪽 항에
속함을 보이면 충분하다. 지표 $\gamma$와
$g\in k_\gamma[[T_1,\ldots,T_r]]$를
\[
gf\in k_\gamma[[T_1,\ldots,T_r]]
\]
가 되도록 택하자. 필요하면 $g$를 $g^q$로 바꾸어
$g\in k_0[[T_1,\ldots,T_r]]$라고 가정할 수 있다. 그러면 모든
$\alpha\geqslant\gamma$에 대하여 보조정리
\sref{0.21.8.8.2-ko}에 의해
\[
gf\in k_\gamma[[T_1,\ldots,T_r]]
\cap k_\alpha((T_1,\ldots,T_r))
=k_\alpha[[T_1,\ldots,T_r]]
\]
이다. 한편 환 $k_\alpha[[T_1,\ldots,T_r]]$들의 교집합은 명백히
$k_0[[T_1,\ldots,T_r]]$이다. 따라서
$f\in k_0((T_1,\ldots,T_r))$이고, 증명이 끝난다.
\end{proof}

\subsection{형식적 멱급수환에서의 미분 가군의 완비화}
\label{subsection:0.21.9-ko}

이 번호에서는 체의 표수가 더 이상 반드시 양수라고 가정하지 않는다.

\begin{lemma}[21.9.1]
\label{0.21.9.1-ko}
$k$를 체, $A$를 완비 뇌터 국소환인 $k$-대수, $K$를 $A$의
잉여체라 하자. $K$가 $k$의 유한형 확대체이면 $A$-가군
$\widehat{\Omega}^1_{A/k}$는 유한형이다.
\end{lemma}

\begin{proof}
\sref{0.20.7.15-ko}에 의해 $\Omega^1_{K/k}$가 유한 차원의
$K$-벡터공간임을 보이면 충분하다. 그런데 가정에 의해 $K$는 어떤
유한 생성 $k$-대수 $B$의 분수체이다. 또한
$\Omega^1_{B/k}$는 \sref{0.20.4.7-ko}에 의해 유한형 $B$-가군이므로,
결론은 \sref{0.20.5.9-ko}에서 따른다.
\end{proof}

\begin{proposition}[21.9.2]
\label{0.21.9.2-ko}
$k_0$를 체, $A$를 형식적으로 매끄러운 $k_0$-대수인 완비 뇌터
국소환, $\mathfrak m$을 그 극대 아이디얼,
$K=A/\mathfrak m$을 그 잉여체라 하자. 따라서 환으로서 $A$는
형식적 멱급수환 $K[[T_1,\ldots,T_r]]$과 동형이다
\sref{0.19.6.5-ko}. $K$가 $k_0$의 유한형 확대체이면
$\widehat{\Omega}^1_{A/k_0}$은 계수가
\[
\dim(A)+\deg.\operatorname{tr}_{k_0}K
\]
인 자유 $A$-가군이다. 더욱이 $\Omega^1_{k_0/k}$가 유한 차원인
$k_0$의 모든 부분체 $k$에 대하여 $\widehat{\Omega}^1_{A/k}$는
계수가
\[
\dim(A)+\deg.\operatorname{tr}_{k_0}K
+\operatorname{rg}_{k_0}(\Omega^1_{k_0/k})
\]
인 자유 $A$-가군이다.
\end{proposition}

\begin{proof}
$P$를 $k_0$의 소체라 하면 $K$가 형식적으로 매끄러운
$P$-대수임은 명백하다 \sref{0.19.6.1-ko}. 한편
$\Upsilon^K_{A/k_0/P}=0$임을 지적하자. 실제로 이는 준동형
\[
  u_{A/k_0/P}\otimes1_K:
  \Omega^1_{k_0/P}\otimes_{k_0}K
  \longrightarrow \Omega^1_{A/P}\otimes_A K
\]
이 단사임을 보이는 것과 같다 \sref{0.20.6.14.5-ko}. 여기서 $u$는
구조 준동형 $k_0\to A$를 나타낸다. 그런데 가정에 의해 $A$는
형식적으로 매끄러운 $k_0$-대수이므로, \sref{0.20.7.2-ko}에 의해
$u_{A/k_0/P}$는 형식적으로 왼쪽 가역이다. 따라서
\sref{0.19.1.7-ko}에 의해 $u_{A/k_0/P}\otimes1_K$는 왼쪽 가역이고,
그러므로 특히 단사이다. 이로써 주장이 증명된다.

완전열 \sref{0.20.6.22.1-ko}에서 $\Lambda,A,B,C$를 각각
$P,k_0,A,K$로 바꾸어 적용하면 완전열
\[
  0\longrightarrow \Upsilon^K_{A/k_0}
  \longrightarrow \mathfrak m/\mathfrak m^2
  \longrightarrow \Omega^1_{A/k_0}\otimes_A K
  \longrightarrow \Omega^1_{K/k_0}\longrightarrow0
\]
을 얻는다. 따라서
\[
\begin{split}
  \operatorname{rg}_K(\Omega^1_{A/k_0}\otimes_AK)
  &={\operatorname{rg}}_K(\mathfrak m/\mathfrak m^2)
    +\bigl(\operatorname{rg}_K(\Omega^1_{K/k_0})
    -\operatorname{rg}_K(\Upsilon_{K/k_0})\bigr)\\
  &={\operatorname{rg}}_K(\mathfrak m/\mathfrak m^2)
    +\deg.\operatorname{tr}_{k_0}K
\end{split}
\]
이다 \sref{0.21.7.1-ko}. 한편 $\widehat\Omega^1_{A/k_0}$은 형식적
사영 $A$-가군이고 \sref{0.20.4.9-ko}, 그 위상은
$\mathfrak m$-준아딕 위상이므로 \sref{0.20.4.5-ko}, 모든 $j$에 대하여
\[
\widehat\Omega^1_{A/k_0}/
\mathfrak m^{j+1}\widehat\Omega^1_{A/k_0}
\]
은 사영 $(A/\mathfrak m^{j+1})$-가군이다 \sref{0.19.2.4-ko}. 따라서
이는 자유 가군이고 ($\mathbf 0_{\mathrm{III}}$, 10.1.2), 그 계수는

\oldpage[0\textsubscript{IV-1}]{177}
$n=\operatorname{rg}_K(\Omega^1_{A/k_0}\otimes_AK)$이다.
\sref{0.21.9.1-ko}에 의해 $\widehat\Omega^1_{A/k_0}$은 유한형
$A$-가군이다. 또한 이 $A$-가군은 ($\mathbf 0_{\mathrm{III}}$,
10.2.1)에 의해 평탄하므로, 계수 $n$인 자유 가군이다
($\mathbf 0_{\mathrm{III}}$, 10.1.3). 이로써 첫 번째 주장이 증명된다.

가정에 의해 $\Omega^1_{k_0/k}$는 $k_0$ 위에서 유한 차원이므로 완비
텐서곱 $\Omega^1_{k_0/k}\widehat\otimes_{k_0}A$는
$\Omega^1_{k_0/k}\otimes_{k_0}A$와 동일시된다. $A$가 형식적으로
매끄러운 $k_0$-대수이므로, \sref{0.20.7.18-ko}에 의해 준동형
\[
  \widehat u_{A/k_0/k}:
  \Omega^1_{k_0/k}\otimes_{k_0}A
  \longrightarrow\widehat\Omega^1_{A/k}
\]
은 연속 $A$-준동형인 왼쪽 역을 가진다. 특히 그 상은
$\widehat\Omega^1_{A/k}$에서 닫혀 있다. 따라서 완전열
\[
  0\longrightarrow \Omega^1_{k_0/k}\otimes_{k_0}A
  \xrightarrow{\widehat u_{A/k_0/k}}\widehat\Omega^1_{A/k}
  \longrightarrow\widehat\Omega^1_{A/k_0}\longrightarrow0
\]
은 \sref{0.20.7.17-ko}과 위에서 보인 것에 의해 분할된다. 이것으로
명제의 증명이 끝난다.
\end{proof}

\begin{corollary}[21.9.3]
\label{0.21.9.3-ko}
$k$를 체, $A$를 통상적인 $k$-대수 구조를 갖춘 형식적 멱급수환
$k[[T_1,\ldots,T_r]]$이라 하자. 그러면 $\widehat\Omega^1_{A/k}$는
계수가 $r=\dim(A)$인 자유 $A$-가군이다.
\end{corollary}

\begin{proof}
$A$가 형식적으로 매끄러운 $k$-대수임에 유의하고
\sref{0.19.3.4-ko}, \sref{0.21.9.2-ko}에서 $k_0=k=K$로 두면 충분하다.
\end{proof}

\begin{lemma}[21.9.4]
\label{0.21.9.4-ko}
$k$를 표수가 $p>0$인 체, $A$를 완비 뇌터 국소환인 $k$-대수,
$B$를 $k[[T_1,\ldots,T_r]]$ 꼴의 위상 $k$-대수와 동형인 $A$의
$k$-부분대수라 하자. 또한 $A$가 유한형 $B$-대수라고 가정하자.
$B_1$이 $B$의 부분대수 $k[[T_1^p,\ldots,T_r^p]]$이면,
$\widehat\Omega^1_{A/k}$는 $\Omega^1_{A/B_1}$과 표준적으로 동일시된다.
\end{lemma}

\begin{proof}
$A$-위상가군 $P$로 가는 $B_1$의 연속 미분 가운데 $A$에서 $P$로
가는 연속 $k$-미분의 제한인 것은 모두 \emph{영}이다. 실제로 그러한
미분은 다항식 부분환 $k[T_1^p,\ldots,T_r^p]$에서 영이고, 이 부분환은
$B_1$에서 조밀하다. 따라서 완전열 \sref{0.20.3.7.1-ko}에 의해 표준
준동형
\[
  \operatorname{D\acute{e}r.\,cont}_{B_1}(A,P)
  \longrightarrow \operatorname{D\acute{e}r.\,cont}_k(A,P)
\]
은 전단사이다. $\Omega^1_{A/B_1}$은 유한형 $A$-가군이고
\sref{0.20.4.7-ko}, $A$가 완비 뇌터 국소환이므로 이 가군은 분리이고
완비이다. 이제 결론은 \sref{0.20.7.14.4-ko}에서 따른다.
\end{proof}

\begin{proposition}[21.9.5]
\label{0.21.9.5-ko}
$A$를 완비 뇌터 국소 정역, $A_0$를 체 $k$ 위의 형식적 멱급수환
$k[[T_1,\ldots,T_r]]$과 동형인 $A$의 부분환이라 하자. $A$가 유한
$A_0$-대수이고, $A$의 분수체 $E$가 $A_0$의 분수체 $L_0$의 분리
확대체라고 가정하자. 그러면
\begin{equation}
\tag{21.9.5.1}
\label{0.21.9.5.1-ko}
  \operatorname{rg}_A\widehat\Omega^1_{A/k}
  =\operatorname{rg}_E(\widehat\Omega^1_{A/k}\otimes_AE)=\dim(A)
\end{equation}
이다.
\end{proposition}

\begin{proof}
$\mathfrak m_0$를 $A_0$의 극대 아이디얼이라 하자. $A$가 유한
$A_0$-대수이므로 $A$의 위상은 $\mathfrak m_0$-아딕 위상이다
(Bourbaki, \emph{Alg. comm.}, chap.~IV, \S~2,
n\textsuperscript{o}~5, prop.~9의 cor.~3). 또한 이 위상이 $A_0$에
유도하는 위상은 $\mathfrak m_0$-아딕 위상이다
(Bourbaki, \emph{Alg. comm.}, chap.~III, \S~3,
n\textsuperscript{o}~4, th.~3). \sref{0.21.9.1-ko}에 의해
$\widehat\Omega^1_{A/k}$는 유한형 $A$-가군이다. 한편

\oldpage[0\textsubscript{IV-1}]{178}
\sref{0.21.9.3-ko}에 의해 $\widehat\Omega^1_{A_0/k}$는 계수가
$r=\dim(A_0)=\dim(A)$인 자유 $A_0$-가군이다
\sref{0.16.1.5-ko}. 따라서 텐서곱
$\widehat\Omega^1_{A_0/k}\otimes_{A_0}A$는 계수가 $r$인 자유
$A$-가군이고 그 분리 완비화와 같다. 끝으로
$\widehat\Omega^1_{A/A_0}$은 유한형 $A$-가군이므로
\sref{0.20.4.7-ko}, 그 분리 완비화와 같다
($\mathbf 0_{\mathrm I}$, 7.3.6).

이제 준동형들의 열
\begin{equation}
\tag{21.9.5.2}
\label{0.21.9.5.2-ko}
  \widehat\Omega^1_{A_0/k}\otimes_{A_0}A
  \xrightarrow{u}\widehat\Omega^1_{A/k}
  \xrightarrow{v}\widehat\Omega^1_{A/A_0}\longrightarrow0
\end{equation}
에서 $v$는 전사이고 $\operatorname{Im}(u)$는
$\operatorname{Ker}(v)$에서 조밀하다 \sref{0.20.7.17-ko}. 그런데
$\widehat\Omega^1_{A/k}$의 모든 $A_0$-부분가군은
$\mathfrak m_0$-아딕 위상에 대하여 닫혀 있다
($\mathbf 0_{\mathrm I}$, 7.3.5). 따라서 열
\sref{0.21.9.5.2-ko}는 완전하다. 한편 $A$는 $A_0$ 위에서 정수적이고,
$A_0$는 정수적으로 닫혀 있으며, $E$는 $L_0$의 분리 확대체이므로
$\widehat\Omega^1_{A/A_0}$은 꼬임 $A$-가군이다
\sref{0.20.4.13-ko}, (iv). 완전열 \sref{0.21.9.5.2-ko}에 $E$를
텐서하면 완전열
\[
  \widehat\Omega^1_{A_0/k}\otimes_{A_0}E
  \longrightarrow \widehat\Omega^1_{A/k}\otimes_AE
  \longrightarrow0
\]
을 얻는다. 따라서
$\operatorname{rg}_E(\widehat\Omega^1_{A/k}\otimes_AE)\leqslant r$이다.

이제 $\widehat\Omega^1_{A/k}\otimes_AE$ 안에 $E$ 위에서 선형
독립인 원소 $r$개가 존재함을 보이면 된다. $D_i=\partial/\partial T_i$
($1\leqslant i\leqslant r$)를 $A_0$에서 자기 자신으로 가는 표준
미분이라 하자. $E$가 $L_0$의 유한 분리 확대체이므로 이 미분들은
$E$에서 자기 자신으로 가는 미분들로 유일하게 연장되며, 이 연장들도
$D_i$로 나타내자. 이를 $A$에 제한하면 $A$에서 $E$로 가는 미분들을
얻는다. 이 미분들의 값이 모두 $E$의 어떤 하나의 유한형
$A$-부분가군에 속함은 명백하다. 이 미분들은 $A_0$에서 연속이고
$A$의 위상은 $\mathfrak m_0$-아딕 위상이므로, 이로써 어떤
$A$-위상가군으로 가는 $A$의 연속 미분 $r$개를 얻었다. 또한
$D_i(T_j)=\delta_{ij}$이므로 이들은 명백히 선형 독립이다. 이것으로
명제의 증명이 끝난다.
\end{proof}

다음 명제는 \sref{0.21.7.3-ko}의 ``형식적'' 유사물이다.

\begin{proposition}[21.9.6]
\label{0.21.9.6-ko}
$k_0\subset k\subset K_0$를 표수가 $p>0$인 세 체라 하고
$[K_0:k]<+\infty$라고 가정하자. 다음과 같이 놓는다.
\[
  L_0=K_0((T_1,\ldots,T_r)),\qquad
  M_0=K_0((T_1^p,\ldots,T_r^p)),
\]
\[
  L=k((T_1,\ldots,T_r)),\qquad
  M=k((T_1^p,\ldots,T_r^p)),\qquad
  N=k((T_1^{p^2},\ldots,T_r^{p^2})).
\]
$E$를 $L_0$의 유한형 확대체라 하자.
\begin{enumerate}
  \item[\rm(i)] $\Upsilon_{K_0/k_0}$가 유한 차원이면
  \begin{equation}
  \tag{21.9.6.1}
  \label{0.21.9.6.1-ko}
    \operatorname{rg}_E\Omega^1_{E/M}
    -\operatorname{rg}_{K_0}\Omega^1_{K_0/k}
    =(r+\deg.\operatorname{tr}_{L_0}E)
      +d(E/K_0,k_0)+d(E/M_0,N/k_0)
  \end{equation}
  이다. 여기서는 \sref{0.21.6.1-ko}의 기호를 쓴다. 특히
  \begin{equation}
  \tag{21.9.6.2}
  \label{0.21.9.6.2-ko}
    \operatorname{rg}_E\Omega^1_{E/M}
    -\operatorname{rg}_{K_0}\Omega^1_{K_0/k}
    \geqslant r+\deg.\operatorname{tr}_{L_0}E+d(E/K_0,k_0)
  \end{equation}
  이다. 더욱이 \sref{0.21.9.6.2-ko}의 두 변이 같다면,
  $k_0\subset k'\subset k$인 체 $k'$로 $k$를 바꾸어도 두 변은
  여전히 같다.

  \item[\rm(ii)] 더 나아가 $[k_0:k_0^p]<+\infty$라고 가정하자.
  $(k_\alpha)$를 $k_0$를 포함하는 $K_0$의 부분체들의 감소하는
  여과족이라 하자. 모든 $\alpha$에 대하여
  $[K_0:k_\alpha]<+\infty$이고
  \[
    \bigcap_\alpha k_\alpha(K_0^p)=k_0(K_0^p)
  \]
  라고 가정하자. 그러면 $k=k_\alpha$로 둘 때
  \sref{0.21.9.6.2-ko}의 두 변이 같아지는 $\alpha$가 존재한다.
\end{enumerate}
\end{proposition}

\begin{proof}
\begin{enumerate}
  \oldpage[0\textsubscript{IV-1}]{179}
  \item[(i)] $M$을 $M(L_0^p)$로 바꾸어도
  $\operatorname{rg}_E\Omega^1_{E/M}$는 변하지 않는다
  \sref{0.21.1.5-ko}, (ii). 그런데
  $L_0^p=K_0^p((T_1^p,\ldots,T_r^p))$이고
  $[k(K_0^p):k]<+\infty$이므로
  \[
    M(L_0^p)=k(K_0^p)((T_1^p,\ldots,T_r^p))
  \]
  이다. 마찬가지로 $\Omega^1_{K_0/k}$에서 $k$를 $k(K_0^p)$로
  바꾸어도 이 가군은 변하지 않는다. 따라서 이후에는
  \emph{$K_0^p\subset k$라고 가정}할 수 있고, 계속 이 가정을 사용한다.
  또한 체
  \[
    k'_0=k_0(K_0^p),\qquad
    N'=k'_0((T_1^{p^2},\ldots,T_r^{p^2}))
  \]
  를 도입하면 다음 체들의 도식을 얻는다.
  \[
    \xymatrix{
      K_0\ar[r] & M_0\ar[r] & L_0\ar[r] & E\\
      k\ar[u]\ar[r] & N\ar[u]\ar[r] & M\ar[u]\ar[r] & L\ar[u]\\
      k_0\ar[u]\ar[r] & k'_0\ar[u]\ar[r] & N'\ar[u] &
    }
  \]

  \sref{0.21.9.6.2-ko}의 첫째 변에서 둘째 변을 뺀 차 $\delta$를
  계산하자. \sref{0.21.7.3-ko}에 의해
  \[
    \operatorname{rg}_E\Omega^1_{E/M}
    -\operatorname{rg}_{M_0}\Omega^1_{M_0/M}
    =\deg.\operatorname{tr}_{M_0}E+d(E/M_0,M)
  \]
  이다. 또한 $L_0$는 $M_0$의 유한 확대체이므로
  $\deg.\operatorname{tr}_{L_0}E=\deg.\operatorname{tr}_{M_0}E$이다.
  한편 $k=k(K_0^p)$ 위의 $K_0$의 $p$-기저는 $M$ 위의 $M_0$의
  $p$-기저이기도 하므로, \sref{0.21.4.5-ko}에 의해
  \[
    \operatorname{rg}_{M_0}\Omega^1_{M_0/M}
    =\operatorname{rg}_{K_0}\Omega^1_{K_0/k}
  \]
  이다. 따라서
  \begin{equation}
  \tag{21.9.6.3}
  \label{0.21.9.6.3-ko}
    \delta=d(E/M_0,M)-d(E/K_0,k_0)-r
  \end{equation}
  이다.

  이제 다음 고전적 보조정리를 기록해 두자.

  \begin{lemma}[21.9.6.4]
  \label{0.21.9.6.4-ko}
  임의의 체 $k$에 대하여 형식적 멱급수체
  $K=k((T_1,\ldots,T_r))$는 $k$의 분리 확대체이다.
  \end{lemma}

  \begin{proof}
  완결성을 위하여 이 보조정리의 증명을 간단히 상기하자. $k$의 모든
  유한 확대체 $k'$에 대하여 $K\otimes_k k'$에 멱영원이 없음을
  보이면 충분하다 (Bourbaki, \emph{Alg.}, chap.~VIII, \S~7,
  n\textsuperscript{o}~3, th.~1의 증명). 실제로
  $A=k[[T_1,\ldots,T_r]]$, $S=A-\{0\}$라 놓으면
  $K\otimes_k k'=S^{-1}(A\otimes_k k')$이다. 또한
  $A\otimes_k k'$은 정역 $A'=k'[[T_1,\ldots,T_r]]$과 표준적으로
  동일시되고 $A$는 $A'$의 부분환과 동일시된다. 이로부터 결론이
  따른다.
  \end{proof}

  이 보조정리와 \sref{0.21.6.2-ko}, (ii)를 사용하면
  $d(L_0/K_0,k_0)=0$이고, 따라서 공식 \sref{0.21.6.4.4-ko}에 의해
  \[
    d(E/K_0,k_0)=d(E/L_0,k_0)
  \]
  이다.

  \oldpage[0\textsubscript{IV-1}]{180}
  한편 공식 \sref{0.21.6.4.4-ko}과 \sref{0.21.6.4.2-ko}에 의해
  \[
    d(E/M_0,M)=d(E/L_0,M)+d(L_0/M_0,M),
  \]
  \[
    d(L_0/M_0,M)=d(L_0/M_0,M/N)+d(L_0/M_0,N)
  \]
  이므로
  \[
    \delta=d(E/L_0,M)+d(L_0/M_0,N)-d(E/L_0,k_0)
    +d(L_0/M_0,M/N)-r
  \]
  이다.

  이제 다음 관계를 보이자.
  \begin{equation}
  \tag{21.9.6.5}
  \label{0.21.9.6.5-ko}
    d(L_0/M_0,M/N)=r
  \end{equation}
  정의에 의해
  \[
    d(L_0/M_0,M/N)
    =\operatorname{rg}_{L_0}\Upsilon_{L_0/M/N}
    -\operatorname{rg}_{M_0}\Upsilon_{M_0/M/N}
  \]
  이다. 두 완전열
  \[
  \begin{aligned}
    0&\longrightarrow\Upsilon_{L_0/M/N}
      \longrightarrow\Omega^1_{M/N}\otimes_ML_0
      \longrightarrow\Omega^1_{L_0/N}
      \longrightarrow\Omega^1_{L_0/M}\longrightarrow0,\\
    0&\longrightarrow\Upsilon_{M_0/M/N}
      \longrightarrow\Omega^1_{M/N}\otimes_MM_0
      \longrightarrow\Omega^1_{M_0/N}
      \longrightarrow\Omega^1_{M_0/M}\longrightarrow0
  \end{aligned}
  \]
  을 고려하면
  \[
  \begin{split}
    d(L_0/M_0,M/N)={}&\operatorname{rg}_{L_0}\Omega^1_{L_0/M}
      -\operatorname{rg}_{L_0}\Omega^1_{L_0/N}\\
      &-\operatorname{rg}_{M_0}\Omega^1_{M_0/M}
      +\operatorname{rg}_{M_0}\Omega^1_{M_0/N}
  \end{split}
  \]
  을 얻는다. 그런데 $N(L_0^p)\supset M$이므로
  $\Omega^1_{L_0/N}=\Omega^1_{L_0/M}$이다
  \sref{0.21.1.5-ko}, (ii). 한편 $T_i^p$ ($1\leqslant i\leqslant r$)는
  $N$ 위의 $M$의 $p$-기저를 이루고, $k$ 위의 $K_0$의 $p$-기저는
  $M$ 위의 $M_0$의 $p$-기저이기도 하다. 따라서 $M_0^p\subset N$임을
  함께 고려하면 \sref{0.21.4.5-ko}와 \sref{0.21.1.10-ko}에 의해
  \[
    \operatorname{rg}_{M_0}\Omega^1_{M_0/N}
    =r+\operatorname{rg}_{K_0}\Omega^1_{K_0/k}
    =r+\operatorname{rg}_{M_0}\Omega^1_{M_0/M}
  \]
  이다. 이로써 \sref{0.21.9.6.5-ko}가 증명된다.

  다시 \sref{0.21.6.4.2-ko}에 의해
  \[
    d(E/L_0,M)=d(E/L_0,k_0)+d(E/L_0,M/k_0)
  \]
  이고
  \[
    d(L_0/M_0,N)=d(L_0/M_0,N/k)+d(L_0/M_0,k)
  \]
  이다. 그런데 $L_0$는 $K_0$의 분리 확대체이므로
  \sref{0.21.9.6.4-ko}, $d(L_0/K_0,k)=0$이고, 따라서
  $d(L_0/M_0,k)=0$이기도 하다
  (\sref{0.21.6.2-ko}와 \sref{0.21.6.4-ko}). 그러므로
  \[
    \delta=d(E/L_0,M/k_0)+d(L_0/M_0,N/k)
  \]
  이다. 이 식의 우변 첫째 항에 \sref{0.21.6.4.2-ko}를 다시 두 번
  적용하면
  \[
    d(E/L_0,M/k_0)
    =d(E/L_0,M/N)+d(E/L_0,N/k)+d(E/L_0,k/k_0)
  \]
  을 얻는다. $M\subset N(L_0^p)$이므로
  $d(E/L_0,M/N)=0$이다 \sref{0.21.6.2-ko}. 따라서
  \sref{0.21.6.4.4-ko}에 의해
  \[
    \delta=d(E/L_0,k/k_0)+d(E/M_0,N/k)
  \]
  이다. 그런데 $M_0$와 $L_0$는 $K_0$의 분리 확대체이므로
  \sref{0.21.9.6.4-ko},
  \[
    d(L_0/K_0,k/k_0)=d(M_0/K_0,k/k_0)=0
  \]
  이다 \sref{0.21.6.2-ko}. 다시 \sref{0.21.6.4.4-ko}를 두 번 적용하면

  \oldpage[0\textsubscript{IV-1}]{181}
  \[
    d(E/L_0,k/k_0)=d(E/K_0,k/k_0)=d(E/M_0,k/k_0)
  \]
  이고, 마지막으로 \sref{0.21.6.4.2-ko}를 적용하면
  \[
    \delta=d(E/M_0,N/k_0)
  \]
  이다. 이로써 \sref{0.21.9.6.1-ko}이 증명된다. 더욱이
  $k_0\subset k'\subset k$인 중간체 $k'$로 $k$를 바꾸면 $N$은 어떤
  부분확대체 $N'$으로 바뀐다. 따라서 \sref{0.21.6.4.2-ko}에 의해
  \[
    d(E/M_0,N'/k_0)\leqslant d(E/M_0,N/k_0)
  \]
  이고, 이것이 (i)의 마지막 주장을 증명한다.

  \item[(ii)] 모든 $\alpha$에 대하여
  \[
    N_\alpha=k_\alpha((T_1^{p^2},\ldots,T_r^{p^2}))
  \]
  라 놓자. \sref{0.21.9.6.1-ko}과 \sref{0.21.8.3-ko}에 의해,
  $K_0^p\subset k_\alpha$이고 $E$가 $M_0$의 유한형 확대체임을
  고려하면
  \[
    \bigcap_\alpha N_\alpha=k_0(M_0^p)
  \]
  임을 보이면 충분하다. 그런데 \sref{0.21.8.8.3-ko}과 관계
  $\bigcap_\alpha k_\alpha=k_0(K_0^p)$
  \sref{0.21.8.6-ko}에 의해
  \begin{equation}
  \tag{21.9.6.6}
  \label{0.21.9.6.6-ko}
    \bigcap_\alpha N_\alpha=N'
    =(k_0(K_0^p))((T_1^{p^2},\ldots,T_r^{p^2}))
  \end{equation}
  이다. 한편
  \[
    k_0(M_0^p)=k_0(K_0^p((T_1^{p^2},\ldots,T_r^{p^2})))
  \]
  이다. 따라서 \sref{0.21.9.6-ko}, (ii)의 증명을 끝내려면 다음
  보조정리만 증명하면 된다.

  \begin{lemma}[21.9.6.7]
  \label{0.21.9.6.7-ko}
  $k_0$를 $[k_0:k_0^p]<+\infty$인 표수 $p>0$의 체, $K_0$를
  $k_0$의 확대체라 하자. 그러면
  \begin{equation}
  \tag{21.9.6.8}
  \label{0.21.9.6.8-ko}
    (k_0(K_0^p))((T_1,\ldots,T_r))
    =k_0\bigl(K_0^p((T_1,\ldots,T_r))\bigr)
  \end{equation}
  이다.
  \end{lemma}

  \begin{proof}
  $k_0^p\subset K_0^p$이므로
  \[
    k_0^p\bigl(K_0^p((T_1,\ldots,T_r))\bigr)
    =K_0^p((T_1,\ldots,T_r))
  \]
  이다. $(c_i)_{1\leqslant i\leqslant n}$이 $k_0^p$ 위의 $k_0$의
  기저이면, $(c_i)$가 \sref{0.21.9.6.8-ko}의 양변 각각을
  $K_0^p((T_1,\ldots,T_r))$ 위에서 생성함은 명백하다.
  \end{proof}
\end{enumerate}
\end{proof}

\begin{remark}[21.9.7]
\label{0.21.9.7-ko}
$[k_0:k_0^p]$가 무한인 경우에는 \sref{0.21.9.6-ko}, (ii)의 주장이
\emph{확장되지 않는다}. 실제로 $k_0$를 무한히 많은 부정원에 관한
유리함수체 $\mathbf F_p(X_n)_{n\geqslant0}$라 하자. 그러면 $X_n$들은
$\mathbf F_p$ 위의 $k_0$의 $p$-기저를 이루므로
$[k_0:k_0^p]=+\infty$이다. \sref{0.21.9.6-ko}의 명제에서
$K_0=k_0$로 두자. 그러면 반드시 $k=k_0$이다. 또한
$L_0=k_0((T))$로 두고, $E$를 다항식
\[
  Y^p-\sum_{n=0}^{\infty}X_nT^{np}
\]
의 근 $y$를 첨가한 확대체 $L_0(y)$라 하자. 이때
\sref{0.21.9.6.2-ko}의 양변의 차 $\delta$는 \emph{영이 아니다}.
실제로 $E$는 $k_0$의 분리 확대체이다. 그렇지 않다면
$y^p\in k_0(L_0^p)$이어야 하는데, 대수적으로 독립인 $X_n$이 무한히
많다는 사실 때문에 이것은 즉시 불가능함을 알 수 있다. 따라서
$d(E/K_0,k_0)=0$이고, 공식 \sref{0.21.9.6.3-ko}은
\[
  \delta=d(E/M_0,M_0)-1
\]
로 환원된다. 또한
$d(E/M_0,M_0)=\operatorname{rg}_E\Upsilon_{E/M_0}$임은 명백하다.
따라서 이 계수가 $2$임을 확인하면 된다. 그런데 $E^p\subset M_0$이고,
$L_0$ 위의 $E$의 $p$-기저는 하나의 원소 $y$로 이루어지며, $M_0$
위의 $L_0$의 $p$-기저는 하나의 원소 $T$로 이루어진다. 따라서 $E$가
$M_0$의 대수적 확대체임을 함께 고려하면, 주장은
\sref{0.21.7.1-ko}과 \sref{0.21.4.5-ko}에서 따른다.
\end{remark}

\begin{proposition}[21.9.8]
\label{0.21.9.8-ko}
$k_0$를 임의의 표수의 체, $K_0$를 $k_0$의 분리 확대체, $A$를
$k_0$-대수인 완비 뇌터 국소환이라 하자. $A$의 잉여체가 $K_0$의

\oldpage[0\textsubscript{IV-1}]{182}
유한 확대체라고 가정하자. $A$의 각 소아이디얼 $\mathfrak p$에 대하여
$k(\mathfrak p)$를 $\mathfrak p$에서 $A$의 잉여체라 하자. 그러면
$k_0\subset k\subset K_0$이고 $[K_0:k]<+\infty$인 모든 체 $k$에
대하여
\begin{equation}
\tag{21.9.8.1}
\label{0.21.9.8.1-ko}
\begin{split}
  &\operatorname{rg}_{k(\mathfrak p)}
    \bigl(\widehat\Omega^1_{(A/\mathfrak p)/k}
      \otimes_{A/\mathfrak p}k(\mathfrak p)\bigr)
    -\operatorname{rg}_{K_0}\Omega^1_{K_0/k}\\
  &\hspace{35mm}\geqslant
    \dim(A/\mathfrak p)
    +\operatorname{rg}_{k(\mathfrak p)}\Upsilon_{k(\mathfrak p)/k_0}
\end{split}
\end{equation}
이다. 더 나아가 $[k_0:k_0^p]<+\infty$라고 가정하자. 여기서 $p$는
$k_0$의 표수 지수이다. 그러면
$k_0(K_0^p)\subset k\subset K_0$이고 $[K_0:k]<+\infty$인 체 $k$를
\sref{0.21.9.8.1-ko}의 두 변이 같아지도록 택할 수 있다.
\end{proposition}

\begin{proof}
$A/\mathfrak p$는 $A$와 같은 잉여체 $K$를 가지며 완비이므로, $A$가
정역이고 $\mathfrak p=0$인 경우로 환원할 수 있다. 이 경우
$E=k(\mathfrak p)$를 $A$의 분수체라 놓는다. $K_0$는 $k_0$ 위에서
형식적으로 매끄러우므로 \sref{0.19.6.1-ko}, $A$를 $K_0$-대수로
만드는 $k_0$-단사 준동형 $K_0\to A$가 존재한다. 특히 $A$ 안에는
$K_0[[T_1,\ldots,T_r]]$과 $K_0$-동형인 $K_0$-부분대수 $A_0$가
존재하고, $A$는 유한 $A_0$-대수이다 \sref{0.19.8.9-ko}. 따라서
$E$는 $A_0$의 분수체 $L_0=K_0((T_1,\ldots,T_r))$의 유한 확대체이다.
$K_0$가 $k_0$의 분리 확대체이므로 $\Upsilon_{K_0/k_0}=0$이고,
$\deg.\operatorname{tr}_{L_0}E=0$이다. 또한 $A$가 유한
$A_0$-대수이므로 $\dim(A)=\dim(A_0)$이다 \sref{0.16.1.5-ko}.

$p=1$이면 \sref{0.21.9.5-ko}와 \sref{0.17.1.4-ko}, (iii)에 의해
\[
  \operatorname{rg}_E(\widehat\Omega^1_{A/k}\otimes_AE)=r
\]
이다. 또한 $K_0$는 $k$의 유한 분리 확대체이므로
$\Omega^1_{K_0/k}=0$이다 \sref{0.21.7.4-ko}. 한편 $E$는 $k_0$의
분리 확대체이므로 $\Upsilon_{E/k_0}=0$이다. 따라서 이 경우
\sref{0.21.9.8.1-ko}의 두 변은 같다.

이제 $p>1$이라고 가정하자. $B=k[[T_1,\ldots,T_r]]$라 하면 $A$는
유한 $B$-대수이다. 따라서 $B_1=k[[T_1^p,\ldots,T_r^p]]$라 놓으면
$\widehat\Omega^1_{A/k}$는 $\Omega^1_{A/B_1}$과 동일시된다
\sref{0.21.9.4-ko}. 또한 $M$을 분수체
$k((T_1^p,\ldots,T_r^p))$라 하면, \sref{0.20.5.9-ko}에 의해
$\widehat\Omega^1_{A/k}\otimes_AE$는 $\Omega^1_{E/M}$과 동일시된다.
그러면 부등식 \sref{0.21.9.8.1-ko}은 바로
\sref{0.21.9.6.2-ko}이고, 명제의 마지막 주장은
\sref{0.21.9.6-ko}, (ii)에서 따른다.
\end{proof}

\section{형식적 매끄러움과 정칙성의 미분 판정법}
\label{section:0.22-ko}
\label{0.22-ko}

이 절의 주요 결과는 다음과 같다.
\begin{enumerate}
  \item[a)] 판정법 (22.2.2)와 (22.5.8)은 체 $k$를 포함하는 완비
  뇌터 국소환 $A$가 $k$ 위에서 형식적으로 매끄럽기 위한 필요충분조건을
  준다. 판정법 (22.5.8)은 미분 개념을 사용하지 않고 진술되며
  \sref{0.19.6.6-ko}을 보완한다. 조금 더 기술적인 성격의 명제
  (22.2.2)는 $k$를 고정했을 때 여기서 다루는 국소환 $A$를
  \emph{분류}하게 해 준다 (22.2.5). 이들은 차원 $n$과 잉여체 $K$로
  주어지며, 유일한 조건은
  $n\geqslant\operatorname{rg}_K\Upsilon_{K/k}$이다.

  \item[b)] 정리 (22.3.3)과 나가타의 야코비 판정법 (22.7.3)의 따름정리
  (22.7.6)은 체 $k$를 포함하는 완비 뇌터 국소환 $A$의 한
  \emph{국소화환}이 $k$ 위에서 기하학적으로 정칙임을 보장하는 조건을
  준다. 이 결과들은 제IV장 \textsection7에서 전개할 국소환의
  ``정밀한'' 성질을 연구하는 데 중요한 역할을 할 것이다. 현재로서는
  나가타의 야코비 판정법과 독립적인 (22.7.6)의 증명은 알려져 있지 않다.

  \item[c)] 자리스키의 야코비 판정법 (22.6.7)과 그 여러 변형은
  \sref{0.20.7.8-ko}의 쉬운 귀결이다.
\end{enumerate}

\oldpage[0\textsubscript{IV-1}]{183}

\subsection{형식적 매끄러움의 올림}
\label{subsection:0.22.1-ko}
\label{0.22.1-ko}

\begin{theorem}[22.1.1]
\label{0.22.1.1-ko}
$s:\Lambda\to A$, $u:A\to B$를 두 환 준동형, $\mathfrak J$를 $B$의
아이디얼, $C$를 몫환 $B/\mathfrak J$이라 하자. $C$가 이산 위상에
대하여 형식적으로 매끄러운 $\Lambda$-대수라고 가정하자.
\begin{enumerate}
  \item[\rm(i)] 다음 조건들이 성립한다고 가정하자.
  \begin{enumerate}
    \item[$\alpha$)] $\mathfrak J/\mathfrak J^2$은 사영 $C$-가군이다.
    \item[$\beta$)] 표준 준동형
    \[
      \mathbf S_C^*(\mathfrak J/\mathfrak J^2)
      \longrightarrow \operatorname{gr}_{\mathfrak J}^*(B)
    \]
    은 전단사이다 \sref{0.19.5.2-fr}.
    \item[$\gamma$)] 특성 준동형
    \[
      \chi:\Upsilon_{C/B/\Lambda}\longrightarrow
      \mathfrak J/\mathfrak J^2
    \]
    는 단사이다 \sref{0.20.6.20-ko}.
    \item[$\delta$)] $C$-가군 $\Omega^1_{B/A}\otimes_BC$는 사영이다.
  \end{enumerate}
  그러면 $\mathfrak J$-준아딕 위상을 준 $B$는 형식적으로 매끄러운
  $\Lambda$-대수이다.

  \item[\rm(ii)] 역으로 $B$가 $\mathfrak J$-준아딕 위상에 대하여
  형식적으로 매끄러운 $\Lambda$-대수이고, $A$가 이산 위상에 대하여
  형식적으로 매끄러운 $\Lambda$-대수라고 가정하자. 그러면 (i)의 조건
  $\alpha$), $\beta$), $\gamma$), $\delta$)가 성립한다.
\end{enumerate}
\end{theorem}

\begin{proof}
완전열 \sref{0.20.6.22.1-ko}을 적용할 수 있다. 실제로
$B/\mathfrak J^2$은 $\mathfrak J/\mathfrak J^2$에 의한 $C$의
$\Lambda$-자명 확대이고, $C$는 형식적으로 매끄러운 $\Lambda$-대수라고
가정했다 \sref{0.19.4.4-fr}. 따라서 조건 $\gamma$)는
$\Upsilon^C_{B/A/\Lambda}=0$과 동치이며, 다시 다음 조건과 동치이다.
\begin{enumerate}
  \item[$\gamma'$)] 준동형
  \[
    u_{B/A/\Lambda}\otimes 1_C:
    \Omega^1_{A/\Lambda}\otimes_A C
    \longrightarrow \Omega^1_{B/\Lambda}\otimes_B C
  \]
  는 단사이다.
\end{enumerate}

(i) 조건 $\alpha$)와 $\beta$)에 의해 $B$는 $\mathfrak J$-준아딕
위상에 대하여 형식적으로 매끄러운 $\Lambda$-대수이다
\sref{0.19.5.4-fr}. 따라서 $B$가 이 위상에 대하여 형식적으로
매끄러운 $A$-대수라고 하는 것과 $\Lambda$에 대하여 상대적으로
형식적으로 매끄러운 $A$-대수라고 하는 것은 같다. 그러므로 $B$가
형식적으로 매끄러운 $A$-대수임을 보이려면 \sref{0.20.7.2-ko}에 의해
위상적 $B$-가군의 연속 준동형
\[
  u_{B/A/\Lambda}:\Omega^1_{A/\Lambda}\otimes_A B
  \longrightarrow\Omega^1_{B/\Lambda}
\]
이 \emph{형식적으로 왼쪽 가역}임을 보이면 충분하다. 그런데 $B$는
형식적으로 매끄러운 $\Lambda$-대수이므로 위상적 $B$-가군
$\Omega^1_{B/\Lambda}$는 형식적으로 사영이고
\sref{0.20.4.9-ko}, 그 위상은 $B$의 위상에서 유도된다
\sref{0.20.4.5-ko}. 따라서 \sref{0.19.1.9-fr}에 의해
$u_{B/A/\Lambda}$가 형식적으로 왼쪽 가역임을 보이려면
$u_{B/A/\Lambda}\otimes1_C$가 왼쪽 가역임을 보이면 충분하다. 한편
$\gamma'$)에 의해 이 마지막 사상은 단사이므로 완전열
\sref{0.20.6.14.7-ko}
\[
  0\longrightarrow\Omega^1_{A/\Lambda}\otimes_A C
  \longrightarrow\Omega^1_{B/\Lambda}\otimes_B C
  \longrightarrow\Omega^1_{B/A}\otimes_B C\longrightarrow0
\]
을 얻는다. 끝으로 $\delta$)에 의해 $C$-가군
$\Omega^1_{B/A}\otimes_BC$는 사영이므로 이 완전열은 분할된다. 이로써
(i)의 증명이 끝난다.

\oldpage[0\textsubscript{IV-1}]{184}

(ii) $B$가 $\mathfrak J$-준아딕 위상에 대하여 형식적으로 매끄러운
$\Lambda$-대수이고 $A$가 이산 위상에 대하여 형식적으로 매끄러운
$\Lambda$-대수이면, $B$는 $\mathfrak J$-준아딕 위상에 대하여
형식적으로 매끄러운 $A$-대수이다 \sref{0.19.3.5-fr}, (ii). 따라서
조건 $\alpha$)와 $\beta$)는 \sref{0.19.5.6-fr}에서 따른다. 또한
\sref{0.20.7.2-ko}에 의해 $u_{B/A/\Lambda}$는 형식적으로 왼쪽
가역이므로 $u_{B/A/\Lambda}\otimes1_C$는 왼쪽 가역이고
\sref{0.19.1.7-fr}, 따라서 특히 단사이다. 끝으로
$\Omega^1_{B/A}$는 형식적으로 사영인 $B$-가군이므로
\sref{0.20.4.9-ko}, $\Omega^1_{B/A}\otimes_BC$는 사영 $C$-가군이다
\sref{0.19.2.4-fr}.
\end{proof}

\begin{corollary}[22.1.2]
\label{0.22.1.2-ko}
$s:\Lambda\to A$, $u:A\to B$를 두 환 준동형, $\mathfrak m$을 $B$의
극대 아이디얼, $K=B/\mathfrak m$을 잉여체라 하자. $\Lambda$-대수
$A$와 $K$가 이산 위상에 대하여 형식적으로 매끄럽다고 가정하자.
그러면 다음 조건들이 동치이다.
\begin{enumerate}
  \item[a)] $B$는 $\mathfrak m$-준아딕 위상에 대하여 형식적으로
  매끄러운 $A$-대수이다.
  \item[b)] 표준 준동형
  \begin{equation}
  \tag{22.1.2.1}
  \label{0.22.1.2.1-ko}
    \mathbf S_K^*(\mathfrak m/\mathfrak m^2)
    \longrightarrow\operatorname{gr}_{\mathfrak m}^*(B)
  \end{equation}
  은 전단사이고, 특성 준동형
  \begin{equation}
  \tag{22.1.2.2}
  \label{0.22.1.2.2-ko}
    \chi:\Upsilon_{K/B/\Lambda}\longrightarrow\mathfrak m/\mathfrak m^2
  \end{equation}
  은 단사이다.
\end{enumerate}
\end{corollary}

\begin{proof}
$K$가 체이므로 조건 $\alpha$)와 $\delta$)는 여기서 자명하게 성립한다.
따라서 주장은 \sref{0.22.1.1-ko}에서 곧바로 따른다.
\end{proof}

\begin{remark}[22.1.3]
\label{0.22.1.3-ko}
\sref{0.22.1.2-ko}의 가정들이 성립하고, 더 나아가
$\mathfrak m/\mathfrak m^2$이 \emph{유한} 계수의 $K$-벡터공간이라고
가정하자. $B$가 $\mathfrak m$-준아딕 위상에 대하여 형식적으로
매끄러운 $A$-대수이려면 다음 조건들이 성립하는 것으로 충분하다.

\begin{env}[22.1.3.1]
\label{0.22.1.3.1-ko}
다항식 대수 $E=K[T_1,\ldots,T_n]$, $T_i$들 $(1\leqslant i\leqslant n)$로
생성되는 $E$의 극대 아이디얼 $\mathfrak n$, 그리고 $E$의 아이디얼
$\mathfrak b\supset\mathfrak n^k$가 주어졌다고 하자. 몫으로 내려가면
항등사상 $K\to K$를 주는 모든 $\Lambda$-준동형
$v:B\to E/\mathfrak b$는
\[
  B\xrightarrow{w}E/\mathfrak n^k\longrightarrow E/\mathfrak b
\]
로 인수분해된다. 여기서 $w$는 $\Lambda$-준동형이다.
\end{env}

\begin{env}[22.1.3.2]
\label{0.22.1.3.2-ko}
$F$가 자명 확대 $D_K(K)$(``$K$ 위의 쌍대수''),
$\rho:A\to F$가 합성 $A\xrightarrow{\rho}F\to K$가 표준 준동형인 환
준동형이라고 하자. 몫으로 내려가면 항등사상 $K\to K$를 주는 모든
$\Lambda$-준동형 $v:B\to K$는
\[
  B\xrightarrow{w}F\longrightarrow K
\]
로 인수분해된다. 여기서 $w$는 $\Lambda$-준동형이다. 이때 $F$에는
$\rho$가 정하는 $A$-대수 구조를 준다.
\end{env}

실제로 \sref{0.19.5.5-fr}에서 조건 \sref{0.22.1.3.1-ko}이 표준
준동형 \sref{0.22.1.2.1-ko}의 전단사성을 함의함을 보았다.
\sref{0.22.1.2-ko}에 의해 이제 표준 준동형
\[
  u_{B/A/\Lambda}\otimes1_K:
  \Omega^1_{A/\Lambda}\otimes_AK
  \longrightarrow\Omega^1_{B/\Lambda}\otimes_BK
\]
이 단사임을 보이면 충분하다. 이는 두 항이 $K$-벡터공간이므로, 모든
$K$-벡터공간 $L$에 대하여 표준 준동형
\[
  \operatorname{D\acute{e}r}_\Lambda(B,L)
  \longrightarrow\operatorname{D\acute{e}r}_\Lambda(A,L)
\]
이 \emph{전사}인 것과 동치이다 \sref{0.20.4.8-ko}. 이를 위해서는 표준
준동형
\[
  \operatorname{D\acute{e}r}_\Lambda(B,K)
  \longrightarrow\operatorname{D\acute{e}r}_\Lambda(A,K)
\]
이 전사인 것이 필요충분함이 명백하며, 따라서 주장은
\sref{0.20.1.5-ko}에서 따른다.

두 조건 \sref{0.22.1.3.1-ko}과 \sref{0.22.1.3.2-ko}은 다음 조건에서
따름을 유의하자.

\begin{env}[22.1.3.3]
\label{0.22.1.3.3-ko}
잉여체가 $K$인 모든 아르틴 국소 $\Lambda$-대수 $E$와 $E$의 모든
멱영 아이디얼 $\mathfrak r$에 대하여, 몫으로 내려가면 항등사상
$K\to K$를 주는 모든 $\Lambda$-준동형 $B\to E/\mathfrak r$는
$B\xrightarrow{u}E\to E/\mathfrak r$로 인수분해된다. 여기서 $u$는
$\Lambda$-준동형이다.
\end{env}
\end{remark}

\oldpage[0\textsubscript{IV-1}]{185}

이 결과는 다음과 같이 일반화된다.

\begin{proposition}[22.1.4]
\label{0.22.1.4-ko}
$\varphi:A\to B$를 뇌터 국소환들의 국소 준동형, $k$와 $K$를 각각
$A$와 $B$의 잉여체라 하자. $B$가 준아딕 위상들에 대하여 형식적으로
매끄러운 $A$-대수일 필요충분조건은 다음과 같다. 잉여체가 $K$인 모든
아르틴 국소환 $C$, $C$를 위상적 $A$-대수로 만드는 모든 국소 준동형
$A\to C$, 그리고 $C$의 모든 멱영 아이디얼 $\mathfrak J$에 대하여,
몫으로 내려가서 얻는 준동형 $K\to K$가 항등사상인 모든 국소
$A$-준동형 $B\to C/\mathfrak J$는
\[
  B\xrightarrow{f}C\longrightarrow C/\mathfrak J
\]
로 인수분해된다. 여기서 $f$는 국소 $A$-준동형이다.
\end{proposition}

\begin{proof}
조건의 필요성은 정의 \sref{0.19.3.1-fr}에서 따르므로 충분성을 보이자.
이산 위상을 준 국소 $\widehat A$-대수로 가는 $\widehat B$의 국소
$\widehat A$-준동형들은 그 대수로 가는 $B$의 국소 $A$-준동형을
확장하여 얻어진다. 따라서 \sref{0.19.3.6-fr}에 의해 $A$와 $B$가
\emph{완비}인 경우로 한정할 수 있다. $A$가 \emph{체} $k$일 때에는,
$\Lambda$를 $k$의 소체로 둔 \sref{0.22.1.3.3-ko}과
\sref{0.19.6.1-ko}에서 명제가 따른다.

일반적인 경우에는 $B_0=B\otimes_Ak=B/\mathfrak mB$라 놓자. 여기서
$\mathfrak m$은 $A$의 극대 아이디얼이고, $B_0$는 완비이다. 잉여체가
$K$인 아르틴 국소 $k$-대수 $C_0$, $C_0$의 멱영 아이디얼
$\mathfrak J_0$, 그리고 몫으로 내려가면 $K$ 위의 항등사상을 주는
국소 $k$-준동형 $g_0:B_0\to C_0/\mathfrak J_0$가 주어졌다고 하자.
이를 합성하면 국소 $A$-준동형
\[
  g:B\longrightarrow B_0\xrightarrow{g_0}C_0/\mathfrak J_0
\]
을 얻는다. 가정에 따라 이는
$B\xrightarrow{f}C_0\to C_0/\mathfrak J_0$로 인수분해되며, 여기서
$f$는 국소 $A$-준동형이다. 그런데 $C_0$가 $k$-대수이므로 $f$는
$B\to B_0\xrightarrow{f_0}C_0$로 인수분해되며, 여기서 $f_0$는 국소
$k$-준동형이다. 따라서 명제의 가정들에서 $B$와 $A$를 각각 $B_0$와
$k$로 바꾸어도 여전히 성립한다. 그러므로 방금 본 바에 의해 $B_0$는
형식적으로 매끄러운 $k$-대수이다.

\sref{0.19.7.2-ko}와 \sref{0.19.7.1-ko}을 적용하면, 완비 뇌터 국소환
$B'$, $B'$을 평탄 $A$-가군이자 형식적으로 매끄러운 $A$-대수로 만드는
국소 준동형 $A\to B'$, 그리고 $k$-동형사상
\[
  u'_0:B'_0=B'\otimes_Ak\xrightarrow{\sim}B_0
\]
이 존재한다. $v'_0$를 $u'_0$의 역동형사상이라 하자. 이제 몫으로
내려가면 $v'_0$를 주는 국소 $A$-준동형 $v:B\to B'$이 존재함을 보이자.
이를 위해 $\mathfrak n$과 $\mathfrak n'$을 각각 $B$와 $B'$의 극대
아이디얼이라 하고, 각 $j$에 대하여
\[
  w_j:B/(\mathfrak n^{j+1}+\mathfrak mB)
  \longrightarrow B'/(\mathfrak n'^{j+1}+\mathfrak mB')
\]
를 $v'_0$에서 몫으로 내려가 얻는 준동형이라 하자. 각 $j$에 대하여
다음 두 도식이 가환하도록 국소 $A$-준동형
$v_j:B/\mathfrak n^{j+1}\to B'/\mathfrak n'^{j+1}$을 구성하면 충분하다.
\[
\xymatrix@C=4em@R=3em{
  B/\mathfrak n^{j+1} \ar[r]^{v_j} \ar[d] &
  B'/\mathfrak n'^{j+1} \ar[d] \\
  B/\mathfrak n^j \ar[r]_{v_{j-1}} & B'/\mathfrak n'^j
}
\qquad
\xymatrix@C=4em@R=3em{
  B/\mathfrak n^{j+1} \ar[r]^{v_j} \ar[d] &
  B'/\mathfrak n'^{j+1} \ar[d] \\
  B/(\mathfrak n^{j+1}+\mathfrak mB) \ar[r]_{w_j} &
  B'/(\mathfrak n'^{j+1}+\mathfrak mB')
}
\]
$B$와 $B'$이 완비이므로 $v=\varprojlim v_j$가 원하는 준동형이다.

\oldpage[0\textsubscript{IV-1}]{186}

$v_j$의 귀납적 구성은 $B$에 대한 가정과
\sref{0.19.3.10.1-fr}에서와 같은 논증으로부터 따른다. 여기서는 보조정리
\sref{0.19.3.10.2-fr}와 등식
\[
  \mathfrak n'^j+(\mathfrak n'^{j+1}+\mathfrak mB')
  =\mathfrak n'^j+\mathfrak mB'
\]
을 사용한다. 이제 \sref{0.19.7.1.4-ko}에 의해 $v$는
$A$-\emph{동형사상}이다. 실제로 $B'$은 평탄 $A$-가군이고, $B$는
$\mathfrak m$-준아딕 위상에 대하여 완비이며, 아이디얼
$\mathfrak m^iB$는 $\mathfrak n$-아딕 위상에 대하여 $B$에서 닫혀 있다
($\mathbf 0_{\mathrm I}$, 7.3.5).
\end{proof}

\subsection{체 위에서 형식적으로 매끄러운 국소 대수의 미분적 특징짓기}
\label{subsection:0.22.2-ko}
\label{0.22.2-ko}

\begin{env}[22.2.1]
\label{0.22.2.1-ko}
$k$를 체, $P$를 그 소체, $A$를 국소환인 $k$-대수, $\mathfrak m$을
그 극대 아이디얼, $K=A/\mathfrak m$을 그 잉여체라 하자. $K$는 $P$ 위에서
분리적이므로 이산 위상에 대하여 형식적으로 매끄러운 $P$-대수이다
\sref{0.19.6.1-ko}. 따라서 $\mathfrak m/\mathfrak m^2$에 의한 $K$의
$P$-확대 $A/\mathfrak m^2$은 $P$-\emph{자명}하고, 이에 따라 특성
준동형
\begin{equation}
\tag{22.2.1.1}
\label{0.22.2.1.1-ko}
  \chi_{A/k}:\Upsilon_{K/k}\longrightarrow\mathfrak m/\mathfrak m^2
\end{equation}
이 정의된다 \sref{0.20.6.23-ko}.
\end{env}

\begin{theorem}[22.2.2]
\label{0.22.2.2-ko}
\sref{0.22.2.1-ko}의 조건 아래, $A$가 $\mathfrak m$-준아딕 위상에
대하여 형식적으로 매끄러운 $k$-대수일 필요충분조건은 다음 조건들이
성립하는 것이다.
\begin{enumerate}
  \item[\rm(i)] 표준 준동형
  $\mathbf S_K^*(\mathfrak m/\mathfrak m^2)\to
  \operatorname{gr}_{\mathfrak m}^*(A)$는 전단사이다.
  \item[\rm(ii)] 특성 준동형
  $\chi_{A/k}:\Upsilon_{K/k}\to\mathfrak m/\mathfrak m^2$은 단사이다.
\end{enumerate}
\end{theorem}

\begin{proof}
이는 \sref{0.22.1.2-ko}에서 $\Lambda$를 $P$로, $A$를 $k$로, $B$를
$A$로 바꾸어 얻는 특별한 경우이다. 실제로 $K$와 $k$는 $P$ 위에서
분리적이므로 형식적으로 매끄러운 $P$-대수이다 \sref{0.19.6.1-ko}.
\end{proof}

\begin{remark}[22.2.3]
\label{0.22.2.3-ko}
\sref{0.22.2.2-ko}의 조건 (ii)는 다음 조건 가운데 어느 하나로
바꾸어도 된다.
\begin{enumerate}
  \item[a)] 표준 준동형
  $\Omega^1_k\otimes_kK\to\Omega^1_A\otimes_AK$는 단사이다.
  \item[b)] 모든 $n$에 대하여 표준 준동형
  \[
    \Omega^1_k\otimes_k(A/\mathfrak m^{n+1})
    \longrightarrow
    \Omega^1_A\otimes_A(A/\mathfrak m^{n+1})
  \]
  은 왼쪽 가역이다.
  \item[c)] 표준 준동형
  $\Omega^1_k\widehat\otimes_kA\to\widehat\Omega^1_A$는 왼쪽 가역이다.
  \item[d)] $A$는 $\mathfrak m$-준아딕 위상에 대하여 소체에
  \emph{상대적으로} 형식적으로 매끄러운 $k$-대수이다.
\end{enumerate}
실제로 \sref{0.22.1.1-ko}에서 $\chi_{A/k}$의 단사성이 a)와 동치이고,
a)와 \sref{0.22.2.2-ko}의 조건 (i)를 함께 요구하는 것이 d)와 (i)를
함께 요구하는 것과 동치임을 이미 보았다. 끝으로
\sref{0.22.1.1-ko}에서 b)와 d)가 동치임도 보았고, b)와 c)의 동치는
\sref{0.19.1.8-fr}에서 따른다.
\end{remark}

\begin{env}[22.2.4]
\label{0.22.2.4-ko}
\sref{0.22.2.2-ko}의 주된 응용은 \emph{완비} 뇌터 국소 $k$-대수에
관한 것이다. 이는 $\mathfrak m/\mathfrak m^2$이 유한 차원
$K$-벡터공간인 경우이다. 더 정확히 말해, $k$를 체, $K$를 $k$의 확대체,
$V$를 유한 차원 $K$-벡터공간이라 하자.

\oldpage[0\textsubscript{IV-1}]{187}

다음과 같은 삼쌍 $(A,\alpha,\beta)$들을 생각하자. $A$는 극대 아이디얼이
$\mathfrak m$인, 형식적으로 매끄러운 완비 뇌터 국소 $k$-대수이고,
$\alpha:A/\mathfrak m\xrightarrow{\sim}K$는 $k$-대수 동형사상이며,
$\beta:\mathfrak m/\mathfrak m^2\xrightarrow{\sim}V$는 $\alpha$에 관한
벡터공간의 디-동형사상이다. 이러한 두 삼쌍 $(A,\alpha,\beta)$와
$(A',\alpha',\beta')$이 주어졌을 때, 첫 번째 것에서 두 번째 것으로
가는 \emph{동치}란 다음 두 도식을 가환이 되게 하는 $k$-동형사상
$u:A\xrightarrow{\sim}A'$을 말한다. 여기서 $\mathfrak m'$은 $A'$의
극대 아이디얼이다.
\[
\xymatrix@C=4em@R=3em{
  A/\mathfrak m \ar[r]^{\operatorname{gr}^0(u)} \ar[d]_{\alpha} &
  A'/\mathfrak m' \ar[d]^{\alpha'} \\
  K \ar[r]_{1_K} & K
}
\qquad
\xymatrix@C=4em@R=3em{
  \mathfrak m/\mathfrak m^2 \ar[r]^{\operatorname{gr}^1(u)}
    \ar[d]_{\beta} &
  \mathfrak m'/\mathfrak m'^2 \ar[d]^{\beta'} \\
  V \ar[r]_{1_V} & V
}
\]
삼쌍 $(A,\alpha,\beta)$들의 \emph{동치류}들은 하나의 \emph{집합}을
이루며, 이를 $\operatorname{FL}(K/k,V)$로 나타내기로 하자. 이제 각 삼쌍
$(A,\alpha,\beta)$에는 벡터공간의 합성 $K$-준동형
\[
  \Upsilon_{K/k}\xrightarrow{\chi_{A/k}}
  \mathfrak m/\mathfrak m^2\xrightarrow{\beta}V
\]
이 대응한다. 여기서 $\chi_{A/k}$는 $\alpha^{-1}$에 관한 디-준동형으로
본다. 앞의 정의와 \sref{0.20.6.22.2-ko}에 따르면 이 $K$-준동형은
$\operatorname{FL}(K/k,V)$에서 $(A,\alpha,\beta)$의 \emph{동치류에만
의존한다}. 다시 말해, 표준 사상
\begin{equation}
\tag{22.2.4.1}
\label{0.22.2.4.1-ko}
  \operatorname{FL}(K/k,V)\longrightarrow
  \operatorname{Hom}_K(\Upsilon_{K/k},V)
\end{equation}
을 정의하였다.
\end{env}

\begin{proposition}[22.2.5]
\label{0.22.2.5-ko}
$k$를 체, $K$를 $k$의 확대체, $V$를 유한 차원 $K$-벡터공간이라 하자.
표준 사상 \sref{0.22.2.4.1-ko}은 $\operatorname{FL}(K/k,V)$에서
$\Upsilon_{K/k}$로부터 $V$로 가는 단사 $K$-준동형들의 집합으로 가는
전단사이다.
\end{proposition}

\begin{proof}
(i) \sref{0.22.2.4.1-ko}이 단사임을 보이려면 다음을 증명하면 된다.
$A$, $A'$을 형식적으로 매끄러운 두 완비 뇌터 국소 $k$-대수,
$\mathfrak m$, $\mathfrak m'$을 각각의 극대 아이디얼,
$\alpha:A/\mathfrak m\xrightarrow{\sim}K$와
$\alpha':A'/\mathfrak m'\xrightarrow{\sim}K$를 두 $k$-동형사상이라 하자.
또한 두 $k$-동형사상
\[
  v^0:A/\mathfrak m\xrightarrow{\sim}A'/\mathfrak m',\qquad
  v^1:\mathfrak m/\mathfrak m^2\xrightarrow{\sim}
      \mathfrak m'/\mathfrak m'^2
\]
이 주어지고, 다음 두 도식이 가환이라고 가정하자.
\begin{equation}
\tag{22.2.5.1}
\label{0.22.2.5.1-ko}
\xymatrix@C=4em@R=3em{
  A/\mathfrak m \ar[r]^{v^0} \ar[dr]_{\alpha} &
  A'/\mathfrak m' \ar[dl]^{\alpha'} \\
  & K
}
\qquad
\xymatrix@C=4em@R=3em{
  \mathfrak m/\mathfrak m^2 \ar[r]^{v^1} &
  \mathfrak m'/\mathfrak m'^2 \\
  \Upsilon_{K/k} \ar[u]^{\chi_A} \ar[ur]_{\chi_{A'}} &
}
\end{equation}
그러면 $\operatorname{gr}^0(u)=v^0$이고
$\operatorname{gr}^1(u)=v^1$인 $k$-동형사상
$u:A\xrightarrow{\sim}A'$이 존재함을 보여야 한다. 먼저 $K$가 체이므로
표준 준동형 \sref{0.20.6.7.1-ko}은 전단사이다
\sref{0.20.6.11-ko}. 따라서
$\operatorname{gr}^0(v)=v^0$이고 $\operatorname{gr}^1(v)=v^1$인
$k$-동형사상
$v:A/\mathfrak m^2\xrightarrow{\sim}A'/\mathfrak m'^2$이 이미 존재한다.
$A'$은 거리화 가능하고 완비이며 $A$는 형식적으로 매끄러운
$k$-대수이므로, 합성 준동형
$A\to A/\mathfrak m^2\xrightarrow{v}A'/\mathfrak m'^2$은
$A\xrightarrow{u}A'\to A'/\mathfrak m'^2$로 인수분해된다. 여기서 $u$는

\oldpage[0\textsubscript{IV-1}]{188}

$k$-준동형이다 \sref{0.19.3.11-fr}. 따라서 이미
$\operatorname{gr}^0(u)=v^0$이고 $\operatorname{gr}^1(u)=v^1$이다.
한편 다음 도식은 가환이다.
\[
\xymatrix@C=5em@R=3em{
  \mathbf S_K^*(\mathfrak m/\mathfrak m^2) \ar[r] \ar[d]_{w} &
  \operatorname{gr}_{\mathfrak m}^*(A)
    \ar[d]^{\operatorname{gr}^*(u)} \\
  \mathbf S_K^*(\mathfrak m'/\mathfrak m'^2) \ar[r] &
  \operatorname{gr}_{\mathfrak m'}^*(A')
}
\]
여기서 가로 화살표들은 표준 준동형들이고 $w$는 $v^0$과 $v^1$에서
유도된다. $v^0$과 $v^1$이 전단사이므로 $w$도 전단사이다. 또한 $A$와
$A'$은 형식적으로 매끄러운 $k$-대수이므로 \sref{0.22.2.2-ko}, (i)에
의해 $\operatorname{gr}(u)$는 전단사이다. 그런데 $A$와 $A'$은 분리적이고
완비이므로 이로부터 $u$ 자체가 \emph{전단사}임이 따른다
(Bourbaki, \emph{Alg. comm.}, chap.~III, \textsection~2,
n\textsuperscript{o}~8, cor.~3 du th.~1).

(ii) 이제 \sref{0.22.2.4.1-ko}의 상이 $\Upsilon_{K/k}$에서 $V$로 가는
\emph{단사} 준동형들의 집합임을 보이면 된다. \sref{0.22.2.2-ko}, (ii)에
의해 그 상이 이 집합에 포함됨은 이미 안다. 따라서 다음 보조정리를
보이면 충분하다.
\end{proof}

\begin{lemma}[22.2.5.2]
\label{0.22.2.5.2-ko}
모든 $K$-준동형 $h:\Upsilon_{K/k}\to V$에 대하여, 극대 아이디얼이
$\mathfrak m$이고 잉여체가 $K$인 완비 정칙 뇌터 국소환인 $k$-대수
$A$와 $K$-동형사상
$\beta:\mathfrak m/\mathfrak m^2\xrightarrow{\sim}V$가 존재하여
$h=\beta\circ\chi_{A/k}$이다.
\end{lemma}

\begin{proof}
실제로 이 보조정리가 증명되었다고 하자. $A$가 정칙이라는 사실에서
\sref{0.22.2.2-ko}의 조건 (i)가 성립한다
\sref{0.17.1.1-fr}. 더 나아가 $h$가 단사이면
\sref{0.22.2.2-ko}에 의해 $A$는 형식적으로 매끄러운 $k$-대수이고,
그 동치류는 \sref{0.22.2.4.1-ko}에 의하여 $h$에 대응한다. 반면 $h$가
단사가 아니면 \sref{0.22.2.5.2-ko}가 그 존재를 보장하는 $k$-대수
$A$는 형식적으로 매끄럽지 않다.

\sref{0.22.2.5.2-ko}을 증명하기 위해 $B$를 $K$-대수
$\mathbf S_K^*(V)$, $\mathfrak n$을 $B$의 증대 아이디얼, $A=\widehat B$를
$\mathfrak n$-준아딕 위상에 대한 $B$의 완비화라 하자. 따라서
$n=\operatorname{rg}_KV$이면 $A$는 형식적 멱급수환
$K[[T_1,\ldots,T_n]]$과 동형이다. \sref{0.20.6.11-ko}에 의해
$A/\mathfrak m^2$ 위에는, 여기서 $\mathfrak m=\mathfrak nA$라 할 때,
$h$가 그 특성 준동형이 되는 $V$에 의한 $K$의 $k$-확대 구조가 존재한다.
이 구조는 일반적으로 표준 단사 $k\to K\to A$가 정하는 $k$-자명 구조와
\emph{다르다}. $f:k\to A/\mathfrak m^2$을 구조 준동형이라 하자. $k$가
소체 위에서 분리적이므로 \sref{0.19.6.1-ko}에 의해 $f$는
\[
  k\xrightarrow{g}A\longrightarrow A/\mathfrak m^2
\]
로 인수분해된다. 이렇게 정의한 $k$-대수 $A$가 원하는 조건을 만족한다
\sref{0.20.6.22.2-ko}.
\end{proof}

\begin{theorem}[22.2.6]
\label{0.22.2.6-ko}
$k$를 체, $K$를 $k$의 확대체라 하자. 극대 아이디얼이 $\mathfrak m$이고
$A/\mathfrak m=K$이며 $\mathfrak m$-준아딕 위상에 대하여 형식적으로
매끄러운 완비 뇌터 국소 $k$-대수 $A$가 존재할 필요충분조건은
$\Upsilon_{K/k}$가 유한 차원 $K$-벡터공간인 것이다. 이때 $k$-대수
$A$는 그 차원에 의해 결정된다. 여기서 동형은 몫으로 내려가면 항등사상
$K\to K$를 주는 $k$-동형사상까지 생각하며, 차원에 부과되는 유일한
조건은
\[
  \dim(A)\geqslant\operatorname{rg}_K\Upsilon_{K/k}
\]
이다.
\end{theorem}

\begin{proof}
$\mathfrak m/\mathfrak m^2$은 유한 차원 $K$-벡터공간이므로 조건의
필요성은 \sref{0.22.2.2-ko}, (ii)에서 따른다. 역으로
\sref{0.22.2.5-ko}에 의해 이 조건은 충분하고,
$\mathfrak m/\mathfrak m^2$로 차원이 임의의
\[
  n\geqslant\operatorname{rg}_K(\Upsilon_{K/k})
\]
인 $K$-벡터공간을 택할 수 있다.
\end{proof}

\oldpage[0\textsubscript{IV-1}]{189}

특히 $\Upsilon_{K/k}$의 항등 준동형은 $K$에
$\Upsilon_{K/k}\simeq\mathfrak m/\mathfrak m^2$인 형식적으로 매끄러운
완비 뇌터 $k$-대수를 \emph{표준적으로} 대응시킨다.

\begin{remarks}[22.2.7]
\label{0.22.2.7-ko}
\begin{enumerate}
  \item[\rm(i)] $\Lambda$를 잉여체가 $k$인 뇌터 국소환이라 하자.
  \sref{0.19.7.1-ko}과 \sref{0.19.7.2-ko}에 따르면, 완비 뇌터
  국소환이고 형식적으로 매끄러운 $\Lambda$-대수 $A$들을 결정하는 문제는
  $\Lambda$를 $k$로 바꾸어 얻는 같은 문제와 \emph{동치}이다. 따라서
  원칙적으로 이 문제는 \sref{0.22.2.5-ko}에 의해 완전히 풀리며, 그
  명제는 찾는 $k$-대수들의 \emph{잉여체}를 $k$의 확대체로서 연구하는
  문제로 환원한다.

  \item[\rm(ii)] $\Upsilon_{K/k}$가 $K$ 위에서 유한 차원이라는 사실은
  $K$가 $k$ 위에서 ``유한 순수 비분리 중복도''를 가짐을
  \emph{함의하지 않는다} \sref{0.19.6.6-ko}. 예를 들어 $k_0$를
  표수가 $p>0$인 완전체, $k=k_0(T)$를 $k_0$ 위의 한 부정원에 대한
  유리함수체, $K=k^{p^{-\infty}}$를 $k$의 가장 작은 완전 확대체라 하자.
  그러면
  $\Omega^1_K=\Omega^1_{K/k}=0$이다 \sref{0.21.4.4-ko}. 따라서 정의
  \sref{0.20.6.1.1-ko}에 의해
  $\Upsilon_{K/k}=\Omega^1_k\otimes_kK$이다. 또한
  $k^p=k_0(T^p)$이므로 $T$는 $k^p$ 위에서 $k$의 $p$-기저이고, 이에 따라
  \sref{0.21.4.5-ko}에 의해 $\Omega^1_k$는 $k$ 위에서 차원 $1$이다.
  그러므로 $\Upsilon_{K/k}$는 $K$ 위에서 차원 $1$이다. 하지만 모든
  $n$에 대하여 국소환 $K\otimes_kk^{p^{-n}}$의 잉여체는 $K$와
  동형이고, 따라서 $k^{p^{-n}}$ 위에서 분리적이지 않다.
\end{enumerate}
\end{remarks}

정리 \sref{0.22.2.2-ko}은 다음과 같이 일반화된다.

\begin{proposition}[22.2.8]
\label{0.22.2.8-ko}
\sref{0.22.2.1-ko}의 조건 아래 $k$의 표수 $p$가 양수라고 가정하자.
$(k_\alpha)$를
\[
  \bigcap_\alpha k_\alpha(k^p)=k^p
\]
을 만족하는 $k$의 부분체들로 이루어진 감소하는 필터족이라 하자.
그러면 \sref{0.22.2.2-ko}의 조건 (ii)는 다음 조건 가운데 어느 하나로
바꾸어도 된다.
\begin{enumerate}
  \item[a)] 어떤 $\alpha_0$가 존재하여 모든
  $\alpha\geqslant\alpha_0$에 대하여 표준 준동형
  \[
    \Omega^1_{k/k_\alpha}\otimes_kK
    \longrightarrow\Omega^1_{A/k_\alpha}\otimes_AK
  \]
  는 단사이다.
  \item[b)] 어떤 $\alpha_0$가 존재하여 모든
  $\alpha\geqslant\alpha_0$에 대하여 $A$는 $\mathfrak m$-준아딕
  위상에 대하여 $k_\alpha$에 상대적으로 형식적으로 매끄러운
  $k$-대수이다.
  \item[c)] 모든 $\alpha$와 모든 정수 $n\geqslant0$에 대하여 표준
  준동형
  \[
    \Omega^1_{k/k_\alpha}\otimes_k(A/\mathfrak m^{n+1})
    \longrightarrow
    \Omega^1_{A/k_\alpha}\otimes_A(A/\mathfrak m^{n+1})
  \]
  은 왼쪽 가역이다.
\end{enumerate}
\end{proposition}

\begin{proof}
$A$가 형식적으로 매끄러운 $k$-대수이면 $k$의 모든 부분체에 대하여도
상대적으로 형식적으로 매끄럽다. 따라서 \sref{0.20.7.2-ko}와
\sref{0.19.1.7-fr}에 의해 c)가 성립한다. 또한
\sref{0.20.7.2-ko}에 의해 c)가 b)를 함의하고, b)가 a)를 함의함은
명백하다. 끝으로 \sref{0.22.2.3-ko}을 고려하면, a)가 표준 준동형
\[
  u:\Omega^1_k\otimes_kK\longrightarrow\Omega^1_A\otimes_AK
\]
의 단사성을 함의함을 보이면 충분하다. 모든 $\alpha$에 대하여 다음
도식은 가환이다.
\[
\xymatrix@C=5em@R=3em{
  \Omega^1_k\otimes_kK \ar[r]^{u} \ar[d]_{v_\alpha} &
  \Omega^1_A\otimes_AK \ar[d] \\
  \Omega^1_{k/k_\alpha}\otimes_kK \ar[r]_{u_\alpha} &
  \Omega^1_{A/k_\alpha}\otimes_AK
}
\]

\oldpage[0\textsubscript{IV-1}]{190}

$\alpha\geqslant\alpha_0$일 때 $u_\alpha$가 단사라는 가정으로부터
$\operatorname{Ker}(u)$는 $\operatorname{Ker}(v_\alpha)$들의
\emph{교집합}에 포함된다. 그런데
$\bigcap_\alpha k_\alpha(k^p)=k^p$라는 가정에 의해 이 교집합은
영임을 \sref{0.21.8.3-ko}에서 보았다. 이로써 증명이 끝난다.
\end{proof}

\begin{proposition}[22.2.9]
\label{0.22.2.9-ko}
\sref{0.22.2.1-ko}의 가정 아래 다음 조건들은 동치이다.
\begin{enumerate}
  \item[a)] $A/\mathfrak m^2$은 $\mathfrak m/\mathfrak m^2$에 의한
  $K$의 $k$-자명 확대이다.
  \item[b)] $\chi_{A/k}=0$이다.
  \item[c)] 표준 준동형 \sref{0.20.5.11.2-ko}
  \[
    \delta_{K/A/k}:\mathfrak m/\mathfrak m^2
    \longrightarrow\Omega^1_{A/k}\otimes_AK
  \]
  는 단사이다. 다시 말해, 완전열
  \begin{equation}
  \tag{22.2.9.1}
  \label{0.22.2.9.1-ko}
    0\longrightarrow\mathfrak m/\mathfrak m^2
    \longrightarrow\Omega^1_{A/k}\otimes_AK
    \longrightarrow\Omega^1_{K/k}\longrightarrow0
  \end{equation}
  을 이룬다.
\end{enumerate}
\end{proposition}

\begin{proof}
b)와 c)의 동치는 완전열 \sref{0.20.6.22.1-ko}에서 따르고, c)와 a)의
동치는 \sref{0.20.5.12-ko}에서 따른다. 실제로
\sref{0.22.2.9.1-ko}은 벡터공간들의 완전열이므로 완전하면 분할된다.

$K$가 $k$ 위에서 \emph{분리적}이면 \sref{0.22.2.9-ko}의 동치인 조건들이
성립함을 유의하자. 이는 $\chi_{A/k}$의 정의
\sref{0.22.2.1-ko}과 \sref{0.20.6.3-ko}에서 따른다.
\end{proof}

\begin{corollary}[22.2.10]
\label{0.22.2.10-ko}
\sref{0.22.2.1-ko}의 가정 아래 $k_0$를 $k$의 부분체라 하자. 다음
조건들은 동치이다.
\begin{enumerate}
  \item[a)] $A/\mathfrak m^2$은 $\mathfrak m/\mathfrak m^2$에 의한
  $K$의 $k_0$-자명 확대이다.
  \item[b)] 합성 준동형
  \[
    \Upsilon_{K/k_0}\longrightarrow\Upsilon_{K/k}
    \xrightarrow{\chi_{A/k}}\mathfrak m/\mathfrak m^2
  \]
  은 영이다.
  \item[c)] 특성 준동형
  $\chi_{A/k}:\Upsilon_{K/k}\to\mathfrak m/\mathfrak m^2$은
  \[
    \Upsilon_{K/k}\longrightarrow\Upsilon_{K/k/k_0}
    \xrightarrow{\chi'}\mathfrak m/\mathfrak m^2
  \]
  로 인수분해된다.
\end{enumerate}
더욱이 이 조건들이 성립하면 준동형 $\chi'$은 유일하게 결정되며,
$\mathfrak m/\mathfrak m^2$에 의한 $K$의 $k_0$-자명 확대
$A/\mathfrak m^2$의 특성 준동형과 일치한다.
\end{corollary}

\begin{proof}
b)의 합성은 바로 $\chi_{A/k_0}$이므로 a)와 b)의 동치는
\sref{0.22.2.9-ko}에서 따른다. 한편 b)와 c)는 완전열
\sref{0.20.6.18.1-ko}
\[
  \Upsilon_{K/k_0}\longrightarrow\Upsilon_{K/k}
  \longrightarrow\Upsilon_{K/k/k_0}\longrightarrow0
\]
에 의해 동치이다. 표준 준동형
$\Upsilon_{K/k}\to\Upsilon_{K/k/k_0}$이 전사이므로 $\chi'$의 유일성도
따른다. $\chi'$이 $k_0$-확대 $A/\mathfrak m^2$의 특성 준동형이라는
사실은 \sref{0.20.6.22.2-ko}에서 따른다.
\end{proof}

\begin{corollary}[22.2.11]
\label{0.22.2.11-ko}
\sref{0.22.2.1-ko}의 가정 아래 $p$를 $k$의 특성지수라 하고,
$(k_\alpha)$를
\[
  \bigcap_\alpha k_\alpha(k^p)=k^p
\]
을 만족하는 $k$의 부분체들로 이루어진 감소하는 필터족이라 하자.
$\Upsilon_{K/k}$의 $K$ 위 계수가 유한이면, 어떤 지표 $\alpha$가 존재하여
$A/\mathfrak m^2$은 $\mathfrak m/\mathfrak m^2$에 의한 $K$의
$k_\alpha$-자명 $k_\alpha$-확대이다.
\end{corollary}

\begin{proof}
실제로 \sref{0.21.8.3-ko}에 의해 표준 준동형
\[
  \Upsilon_{K/k}\longrightarrow\Upsilon_{K/k/k_\alpha}
\]
가 전단사인 지표 $\alpha$가 존재한다. 따라서 $k_0$를 $k_\alpha$로
바꾸면 \sref{0.22.2.10-ko}의 조건 b)가 성립한다.
\end{proof}

\oldpage[0\textsubscript{IV-1}]{191}

\subsection{몇몇 국소환과 그 완비화 사이의 관계에 대한 응용}
\label{subsection:0.22.3-ko}

\begin{lemma}[22.3.1]
\label{0.22.3.1-ko}
$A_0$를 뇌터 반국소환, $A$를 유한 $A_0$-대수라 하자. 그러면 $A$도
뇌터 반국소환이다 (cf.~Bourbaki, \emph{Alg. comm.}, chap.~IV, §~2,
n\textsuperscript{o}~5, prop.~9의 cor.~3). $\widehat A_0$와
$\widehat A$를 각각 $A_0$와 $A$의 각 준아딕 위상에 대한 완비화라 하자.
$A_0$, $A$, $\widehat A$에 이산 위상을 주면, $\widehat A$는
$A_0$에 대하여 상대적으로 형식적으로 매끄러운 $A$-대수이다
\sref{0.19.9.1-ko}.
\end{lemma}

\begin{proof}
실제로
\[
  \widehat A=A\otimes_{A_0}\widehat A_0
\]
임이 알려져 있다 (Bourbaki, \emph{Alg. comm.}, chap.~IV, §~2,
n\textsuperscript{o}~5, prop.~9의 cor.~3 및 chap.~III, §~3,
n\textsuperscript{o}~4, th.~1). 따라서 주장은 \sref{0.19.9.3-ko}에서
따른다.
\end{proof}

\begin{proposition}[22.3.2]
\label{0.22.3.2-ko}
$A$를 정칙 뇌터 국소환, $K$를 그 분수체, $p$를 $K$의 표수라 하자.
다음 가정 가운데 하나가 성립한다고 하자.
\begin{enumerate}
  \item[(i)] $p=0$.
  \item[(ii)] $p>0$이고, $A$의 뇌터 부분환들로 이루어진 감소하는 필터족
  $(A_\alpha)$가 존재하여, 모든 $\alpha$에 대하여 $A$는 유한
  $A_\alpha$-대수이고, 적당한 정수 $h(\alpha)$에 대하여
  $A^{p^{h(\alpha)}}\subset A_\alpha$이며, $K_\alpha$를
  $A_\alpha$의 분수체라 하면
  \[
    \bigcap_\alpha K_\alpha(K^p)=K^p
  \]
  이다.
\end{enumerate}
이제 $B$를 $A$를 포함하는 유한 $A$-대수인 정역,
$\mathfrak n$을 $B$의 소 아이디얼, $C$를 국소환 $B_\mathfrak n$이라
하자. 또한 $\mathfrak q$를
$\mathfrak q\cap C=0$인 완비화 $\widehat C$의 소 아이디얼이라 하자.
그러면 국소환 $\widehat C_\mathfrak q$는 $B$의 분수체 $L$ 위의
대수이다. 이때 $\widehat C_\mathfrak q$는 그 $\mathfrak q$-아딕 위상에
대하여 형식적으로 매끄러운 $L$-대수이고, 따라서 $L$ 위에서
기하학적으로 정칙인 환이다 \sref{0.19.6.5-ko}.
\end{proposition}

\begin{proof}
두 경우로 나눈다.

\emph{I)} $\mathfrak n$이 $A$의 극대 아이디얼 $\mathfrak m$을 포함한다고
하자. 그러면 $\mathfrak n\cap A=\mathfrak m$이다. 이때
$\mathfrak n$은 극대 아이디얼이다 (Bourbaki, \emph{Alg. comm.},
chap.~V, §~2, n\textsuperscript{o}~1, prop.~1). $\mathfrak r$을 반국소환
$B$의 근기라 하면, $\widehat C$는 $B$의 $\mathfrak n$-준아딕 위상에
대한 분리 완비화이며, $B$의 $\mathfrak r$-준아딕 위상에 대한 완비화인
반국소환 $\widehat B$의 한 성분이다 (Bourbaki, \emph{Alg. comm.},
chap.~III, §~3, n\textsuperscript{o}~4, prop.~8). 따라서
\[
  \widehat C_\mathfrak q=\widehat B_{\mathfrak q'}
\]
로 쓸 수 있는데, 여기서 $\mathfrak q'$은 $\mathfrak q$의
$\widehat B$에서의 역상인 소 아이디얼이다. 이제 $B$는 $\widehat B$의
부분환이고 $\mathfrak q\cap C=0$이라는 가정은
$\mathfrak q'\cap B=0$을 함의한다. 따라서
$\widehat B_{\mathfrak q'}$은
$\widehat B\otimes_B L$의 어느 소 아이디얼 $\mathfrak q''$에 대한
국소화이기도 하며, $\mathfrak q'$은 $\mathfrak q''$의 $\widehat B$에서의
역상이다. 또한
\[
  \widehat B=\widehat A\otimes_A B
\]
이다 (Bourbaki, \emph{Alg. comm.}, chap.~IV, §~2,
n\textsuperscript{o}~5, prop.~9의 cor. 및 chap.~III, §~3,
n\textsuperscript{o}~4, th.~1). 따라서
\[
  \widehat B\otimes_B L=\widehat A\otimes_A L.
\]
끝으로 $\mathfrak p$를 $\mathfrak q''$의 $\widehat A$에서의 역상인 소
아이디얼이라 하자. $A\to B$가 단사라는 사실과 다음 도식의 가환성으로부터
\[
\xymatrix{
  B \ar[r] & \widehat B\\
  A \ar[u] \ar[r] & \widehat A \ar[u]
}
\]
$\mathfrak p\cap A=0$을 얻는다. 그러므로
$(\widehat A\otimes_A L)_{\mathfrak q''}$은
\[
  \widehat A_\mathfrak p\otimes_A L
  =(\widehat A_\mathfrak p\otimes_A K)\otimes_K L
\]
의 어느 소 아이디얼에 대한 국소화이기도 하다.

\oldpage[0\textsubscript{IV-1}]{192}

이 국소환이 형식적으로 매끄러운 $L$-대수임을 보이려면
$\widehat A_\mathfrak p\otimes_A K$가 형식적으로 매끄러운
$K$-대수임을 보이면 충분하다 \sref{0.19.3.5-ko}, (iii), (iv). 또한
$\mathfrak p\cap A=0$이므로
\[
  \widehat A_\mathfrak p\otimes_A K=\widehat A_\mathfrak p.
\]
이제 판정법 \sref{0.22.2.2-ko}과 \sref{0.22.2.8-ko}을 적용한다. 먼저
$A$가 정칙이므로 $\widehat A$도 정칙이고 \sref{0.17.1.5-ko}, 따라서
$\widehat A_\mathfrak p$도 정칙이다 \sref{0.17.3.2-ko}. 그러므로
\sref{0.22.2.2-ko}의 조건 (i)가 성립한다 \sref{0.17.1.1-ko}.
따라서 $p>0$일 때에는 \sref{0.22.2.8-ko}의 조건 b)를 확인하면
충분하다. 그런데 \sref{0.22.3.1-ko}에 의해 모든 $\alpha$에 대하여
$\widehat A$는 이산 위상에서 $A_\alpha$에 대하여 상대적으로
형식적으로 매끄러운 $A$-대수이다. 따라서
\[
  \widehat A\otimes_{A_\alpha}K_\alpha
  =\widehat A\otimes_A(A\otimes_{A_\alpha}K_\alpha)
\]
는 이산 위상에서 $K_\alpha$에 대하여 상대적으로 형식적으로 매끄러운
$(A\otimes_{A_\alpha}K_\alpha)$-대수이다
\sref{0.19.9.2-ko}, (iii). 한편 $A$는 유한 $A_\alpha$-대수이므로
\[
  A\otimes_{A_\alpha}K_\alpha=K.
\]
$\widehat A_\mathfrak p$는 $\widehat A\otimes_A K$의 분수환이므로,
\sref{0.19.9.2-ko}, (iv)에 의해 $\widehat A_\mathfrak p$는 이산
위상에서 $K_\alpha$에 대하여 상대적으로 형식적으로 매끄러운
$K$-대수이고, 따라서 더욱이 그 $\mathfrak p$-아딕 위상에 대해서도
그러하다 \sref{0.19.3.8-ko}, \sref{0.19.9.5-ko}. 이것이 바로
\sref{0.22.2.8-ko}의 조건 b)이므로 $p>0$인 이 경우의 명제가
증명되었다. $p=0$이면 $\widehat A_\mathfrak p$의 잉여체는 $K$ 위에서
분리적이고, $\widehat A_\mathfrak p$는 정칙환이므로
\sref{0.19.6.4-ko}에 의해 형식적으로 매끄러운 $K$-대수이다.

\emph{II)} 일반적인 경우. $\mathfrak s=\mathfrak n\cap A$를 $A$의
소 아이디얼이라 하고 $S=A-\mathfrak s$라 놓자. 그러면
\[
  S^{-1}A=A_\mathfrak s,
  \qquad C=(S^{-1}B)_{S^{-1}\mathfrak n},
\]
이며, 이제 $S^{-1}\mathfrak n$은 $A_\mathfrak s$의 극대 아이디얼
$\mathfrak sA_\mathfrak s$를 포함한다. $A_\mathfrak s$는 정칙이고
\sref{0.17.3.2-ko}, $S^{-1}B$는 $A_\mathfrak s$를 포함하는 유한
$A_\mathfrak s$-대수인 정역이다. 따라서 $p>0$일 때
$A_\mathfrak s$에 대하여 조건 (ii)를 확인하면 충분하다.
$\mathfrak s_\alpha=\mathfrak s\cap A_\alpha$라 놓고 국소환
$(A_\alpha)_{\mathfrak s_\alpha}$를 생각하자. 모든
$t\notin\mathfrak s$에 대하여 $q=p^{h(\alpha)}$라 놓으면, 가정에 의해
$t^q\in A_\alpha$이고 $t^q\notin\mathfrak s$이므로
$t^q\notin\mathfrak s_\alpha$이다. 따라서
$S_\alpha=A_\alpha-\mathfrak s_\alpha$라 놓으면
\[
  A_\mathfrak s=S_\alpha^{-1}A.
\]
가정에 의해 $A_\mathfrak s$는 유한
$(A_\alpha)_{\mathfrak s_\alpha}$-대수이다. 또한 $K_\alpha$는
$(A_\alpha)_{\mathfrak s_\alpha}$의 분수체이기도 하므로 증명이 끝난다.
\end{proof}

\begin{theorem}[22.3.3]
\label{0.22.3.3-ko}
$A$를 완비 뇌터 국소환, $\mathfrak p$를 $A$의 소 아이디얼,
$B=A_\mathfrak p$를 국소화환, $\mathfrak q$를 완비화 $\widehat B$의
소 아이디얼, $\mathfrak r=\mathfrak q\cap B$라 하자. 그러면
$\widehat B$의 국소화환 $\widehat B_\mathfrak q$는 준아딕 위상들에
대하여 형식적으로 매끄러운 $B_\mathfrak r$-대수이다.
\end{theorem}

\begin{proof}
$B$의 소 아이디얼 $\mathfrak r$은
$\mathfrak n\subset\mathfrak p$인 $A$의 어떤 소 아이디얼
$\mathfrak n$에 대하여 $\mathfrak n_\mathfrak p$ 꼴이고,
\[
  B_\mathfrak r=A_\mathfrak n
\]
이다. $k$를 $B_\mathfrak r$의 잉여체라 하자. 이는 완비 국소 정역
$A'=A/\mathfrak n$의 분수체이기도 하다. $\widehat B$는 평탄한
$B$-가군이므로 $(\mathbf 0_\mathrm I,~7.3.5)$,
$\widehat B_\mathfrak q$는 평탄한 $B_\mathfrak r$-가군이다
$(\mathbf 0_\mathrm I,~6.3.2)$. 따라서
$\widehat B_\mathfrak q$가 형식적으로 매끄러운
$B_\mathfrak r$-대수임을 보이려면
\[
  \widehat B_\mathfrak q\otimes_{B_\mathfrak r}k
  =\widehat B_\mathfrak q/\mathfrak r\widehat B_\mathfrak q
\]
가 형식적으로 매끄러운 $k$-대수임을 보이면 충분하다
\sref{0.19.7.1-ko}. 그런데
\[
  \widehat B_\mathfrak q/\mathfrak r\widehat B_\mathfrak q
  =(\widehat B/\mathfrak r\widehat B)_{\mathfrak q'},
\]
여기서 $\mathfrak q'$은 $\mathfrak q$의
$\widehat B/\mathfrak r\widehat B$에서의 표준상이다. 정의에 의해
\[
  \mathfrak q'\cap(B/\mathfrak r)=0,
  \qquad B/\mathfrak r=A'_{\mathfrak p'},
  \qquad \mathfrak p'=\mathfrak p/\mathfrak n.
\]
또한 $\widehat B/\mathfrak r\widehat B$는 $B/\mathfrak r$의
완비화이다.

따라서 다음 따름정리를 증명하는 것으로 환원된다.
\end{proof}

\begin{corollary}[22.3.4]
\label{0.22.3.4-ko}
$A$를 완비 뇌터 국소 정역, $K$를 그 분수체, $\mathfrak p$를 $A$의
소 아이디얼, $B=A_\mathfrak p$를 국소화환이라 하자. 그러면 완비화
$\widehat B$의 모든 소 아이디얼 $\mathfrak q$에 대하여

\oldpage[0\textsubscript{IV-1}]{193}

$\mathfrak q\cap B=0$이면, 국소환 $\widehat B_\mathfrak q$는 그
$\mathfrak q$-아딕 위상에 대하여 형식적으로 매끄러운 $K$-대수이고,
따라서 $K$ 위에서 기하학적으로 정칙인 환이다.
\end{corollary}

\begin{proof}
실제로 \sref{0.19.8.9-ko}에 의해 $A$의 부분환 $A_0$가 존재하여,
$A_0$는 표수가 $p>0$인 체 위의 형식적 멱급수환이거나 코언 환 위의
형식적 멱급수환이고, $A$는 유한 $A_0$-대수이다. 코언 환의 분수체는
표수가 $0$임을 상기하자. 한편 $A_0$는 정칙이고
\sref{0.17.3.8-ko}, \sref{0.21.8.8-ko}에 의해
\sref{0.22.3.2-ko}의 가정 (i), (ii) 가운데 하나를 만족한다. 따라서
\sref{0.22.3.2-ko}에서 $B$를 $A$로, $A$를 $A_0$로 바꾸어 적용하면
충분하다.
\end{proof}

\subsection{극대 아이디얼의 제곱이 영인 국소환의 유한 확대에 관한 예비 결과}
\label{subsection:0.22.4-ko}
\label{0.22.4-ko}

\begin{proposition}[22.4.1]
\label{0.22.4.1-ko}
$A$를 극대 아이디얼 $\mathfrak m$의 제곱이 영인 국소환,
$K=A/\mathfrak m$을 그 잉여체라 하자. 아이디얼 $\mathfrak m$을
$K$-벡터공간으로 본 것을 $V$로 나타내자. $T_i$
$(1\leqslant i\leqslant r)$를 부정원들이라 하고, 각 $i$에 대하여
$F_i$를 차수가 $d_i$인 $A[T_i]$의 일계수 다항식이라 하자.
$r$개의 다항식 $F_i$가 생성하는 아이디얼 $\mathfrak b$에 의한
다항식환 $A[T_1,\ldots,T_r]$의 몫을 $A'$이라 하자. $A'$이 국소환이라
가정하고, 그 극대 아이디얼을 $\mathfrak m'$, 잉여체를
$K'=A'/\mathfrak m'$이라 하자. 그러면 다음이 성립한다.
\begin{enumerate}
  \item[(i)] $d=\prod_{i=1}^r d_i$라 놓으면 $A'$은 계수가 $d$인 자유
  $A$-가군이고, $[K':K]\leqslant d$이다.
  \item[(ii)] $[K':K]=d$일 필요충분조건은
  \[
    A'\otimes_AK=A'/\mathfrak mA'
  \]
  이 체인 것이다. 이 체는 필연적으로 $K'$과 동형이다. 다시 말해,
  필요충분조건은 $\mathfrak m'=\mathfrak mA'$이다. 이 경우
  $\mathfrak m'^2=0$이고, 아이디얼 $\mathfrak m'$을 $K'$-벡터공간으로
  본 것을 $V'$이라 하면 표준 준동형
  \[
    V\otimes_KK'\longrightarrow V'
  \]
  은 전단사이다. 따라서
  $\operatorname{rg}_{K'}V'=\operatorname{rg}_K V$이다.
\end{enumerate}
\end{proposition}

\begin{proof}
$A'$이 국소라는 가정 없이도, 유클리드 나눗셈과 $r$에 대한 귀납법에
의해 다음은 명백하다. $t_i$를 $T_i$의 $\mathfrak b$ 법류라 하면,
\[
  M_\nu=t_1^{\nu_1}t_2^{\nu_2}\cdots t_r^{\nu_r},
  \qquad 0\leqslant\nu_i\leqslant d_i-1,
\]
인 단항식들이 $A$-가군 $A'$의 기저를 이룬다. $A'$이 국소이면
$K'$은 $A'/\mathfrak mA'=A'\otimes_AK$의 몫이므로
\[
  [K':K]\leqslant[A'\otimes_AK:K]=d.
\]
등호가 성립하려면 반드시 $\mathfrak m'=\mathfrak mA'$이어야 한다.
이로써 (i)와 (ii)의 첫 번째 주장이 증명된다. 한편
$\mathfrak m'=\mathfrak mA'$이면 $\mathfrak m'^2=0$이고,
$A'$이 자유 $A$-가군이므로
\[
  V\otimes_KK'=\mathfrak m\otimes_AA'
  =\mathfrak mA'=V'.
\]
\end{proof}

\begin{remark}[22.4.2]
\label{0.22.4.2-ko}
$A$가 아르틴 환인 경우, 다시 말해 $\operatorname{rg}_K(V)$가 유한인
경우에는, $A'$이 국소라는 가정으로부터 언제나
\[
  \operatorname{rg}_{K'}(\mathfrak m'/\mathfrak m'^2)
  \geqslant\operatorname{rg}_K(V)
\]
가 따른다. 이는 뒤의 \sref{0.22.5.2-ko}에서 볼 것이다. 이 두 계수는
$\mathfrak m'=\mathfrak mA'$이 아니어도 같을 수 있음을 유의하자.
\end{remark}

\begin{proposition}[22.4.3]
\label{0.22.4.3-ko}
$A$에 관한 가정과 기호는 \sref{0.22.4.1-ko}과 같다고 하자. $F_i$들이
\begin{equation}
\label{0.22.4.3.1-ko}
  F_i(T_i)=T_i^{d_i}+\sum_{1\leqslant k\leqslant d_i}
  c_{ik}T_i^{d_i-k},
  \tag{22.4.3.1}
\end{equation}
꼴이고, 모든 $i$에 대하여 $d_i\geqslant2$이며
$c_{ik}\in\mathfrak m$이라고 하자 (``아이젠슈타인 다항식''). 그러면
다음이 성립한다.
\begin{enumerate}
  \item[(i)] $A'$은 국소환이고, 그 잉여체
  $K'=A'/\mathfrak m'$은 $K$와 동형이다.

  \oldpage[0\textsubscript{IV-1}]{194}

  \item[(ii)] 표준 준동형
  \[
    V\longrightarrow V'=\mathfrak m'/\mathfrak m'^2
  \]
  의 핵은 $\xi_i=c_{i,d_i}$, 즉 $F_i$의 ``상수항''들이 생성하는
  $A$의 아이디얼 $\mathfrak n$이다. 이 준동형의 여핵은 $T_i$들의
  상들을 $K'$ 위의 기저로 갖는다. 특히
  \begin{equation}
  \label{0.22.4.3.2-ko}
    \operatorname{rg}_{K'}(V')=\operatorname{rg}_K(V)+r-r',
    \tag{22.4.3.2}
  \end{equation}
  여기서 $r'=\operatorname{rg}_K(\mathfrak n)$이다.
  $\operatorname{rg}_K(V)$가 유한인 경우, 다시 말해 $A$가 아르틴
  환인 경우에는
  \[
    \operatorname{rg}_{K'}(V')=\operatorname{rg}_K(V)
  \]
  일 필요충분조건은 $\xi_i$들이 $K$ 위에서 선형 독립인 것이다.
\end{enumerate}
\end{proposition}

\begin{proof}
(i) $A'\otimes_AK=A'/\mathfrak mA'$은
\[
  K[T_1,\ldots,T_r]/\mathfrak b'
\]
과 동형이다. 여기서 $\mathfrak b'$은 $r$개의 다항식
$T_i^{d_i}$가 생성하는 아이디얼이다. 따라서 이는 잉여체가 $K$이고,
극대 아이디얼이 $T_i$들의 $\mathfrak b'$ 법류들로 생성되는 국소환이다.
이로부터 (i)가 따른다.

(ii) $t_i$ ($1\leqslant i\leqslant r$)를 $T_i$의 $\mathfrak b$ 법류라
하자. (i)에 의해 $A'$의 극대 아이디얼은
\[
  \mathfrak m'=\mathfrak mA'+\sum_{i=1}^r t_iA'
\]
이고, $\mathfrak m^2=0$임을 함께 쓰면
\[
  \mathfrak m'^2=\sum_{i,j}t_it_jA'+\sum_i t_i\mathfrak mA'
\]
을 얻는다. 따라서 $A'/\mathfrak m'^2$은
$A[T_1,\ldots,T_r]/\mathfrak c$와 동형이다. 여기서
$\mathfrak c$는 다음 원소들이 생성하는 아이디얼이다.
\[
  F_i\ (1\leqslant i\leqslant r),\qquad
  T_iT_j\ (1\leqslant i,j\leqslant r),\qquad
  xT_i\ (1\leqslant i\leqslant r,\ x\in\mathfrak m).
\]
$d_i\geqslant2$이고 $T_i^2\in\mathfrak c$이므로, 가정에 의해
$\mathfrak c$는 다음 원소들로도 생성된다.
\[
  \xi_i\ (1\leqslant i\leqslant r),\qquad
  T_iT_j\ (1\leqslant i,j\leqslant r),\qquad
  T_ix\ (1\leqslant i\leqslant r,\ x\in\mathfrak m).
\]
다음과 같이 놓자.
\[
  \mathfrak c_1=\sum_{i,j}T_iT_jA[T_1,\ldots,T_r],\qquad
  \mathfrak c_2=\mathfrak c_1+
  \sum_i\mathfrak mT_iA[T_1,\ldots,T_r].
\]
그러면
\[
  \mathfrak c=\mathfrak c_2+
  \sum_i\xi_iA[T_1,\ldots,T_r].
\]
$A$-가군 $A[T_1,\ldots,T_r]/\mathfrak c_1$은 $A$와
$At_i^{(1)}$들의 직합이다. 여기서 $t_i^{(1)}$은 $T_i$의
$\mathfrak c_1$ 법류이다. 따라서
$A[T_1,\ldots,T_r]/\mathfrak c_2$는 $A$와
$Kt_i^{(2)}$들의 직합이고, $t_i^{(2)}$는 $T_i$의
$\mathfrak c_2$ 법류이다. 끝으로
$A[T_1,\ldots,T_r]/\mathfrak c$는 $A/\mathfrak n$과
$Kt_i'$들의 직합이다. 여기서 $t_i'$은 $T_i$의 $\mathfrak c$ 법류이고,
$\mathfrak n$은 $\xi_i$들이 생성하는 $A$의 아이디얼이다.
$A'/\mathfrak m'^2$을 $A[T_1,\ldots,T_r]/\mathfrak c$와 동일시하면
$V$의 상은 $\mathfrak m/\mathfrak n$과 동일시된다. 이로써 (ii)가
증명된다.
\end{proof}

\begin{proposition}[22.4.4]
\label{0.22.4.4-ko}
$A$에 관한 가정과 기호는 \sref{0.22.4.1-ko}과 같다고 하자. $K$의
표수가 $p>0$이라고 가정한다. 따라서 $A$에서
$p\cdot1\in\mathfrak m$이다. 또한 $F_i$들이
\begin{equation}
\label{0.22.4.4.1-ko}
  F_i(T_i)=T_i^p+p\sum_{k=1}^{p-1}c_{ik}T_i^{p-k}-f_i,
  \tag{22.4.4.1}
\end{equation}
꼴이라고 하자. 여기서 모든 $i$와 $1\leqslant k\leqslant p-1$에
대하여 $c_{ik}\in A$이고 $f_i\in A$이다. 나아가 $a_i$를 $f_i$의
$\mathfrak m$ 법류라 할 때, 원소족
$(a_i)_{1\leqslant i\leqslant r}$가 $K^p$ 위에서 $p$-자유라고
가정하자 (``절대적으로 $p$-자유''). 그러면 다음이 성립한다.
\begin{enumerate}
  \item[(i)] $A'$은 극대 아이디얼이 $\mathfrak mA'$인 국소환이고,
  그 잉여체 $K'$은 $K$에 $a_i^{1/p}$
  $(1\leqslant i\leqslant r)$들을 첨가하여 얻는 확대체와 동형이다.

  \oldpage[0\textsubscript{IV-1}]{195}

  \item[(ii)] 원소들 $d_Af_i\otimes1_{K'}$은 표준 준동형
  \[
    \Omega^1_A\otimes_AK'\longrightarrow
    \Omega^1_{A'}\otimes_{A'}K'
  \]
  의 핵 $N=\Upsilon^{K'}_{A'/A}$의 기저를 이룬다.
\end{enumerate}
\end{proposition}

\begin{proof}
(i) 이번에는 $A'\otimes_AK=A'/\mathfrak mA'$이
$K[T_1,\ldots,T_r]/\mathfrak b'$과 동형이다. 여기서
$\mathfrak b'$은 다항식들 $T_i^p-a_i$가 생성하는 아이디얼이다.
$a_i$들에 대한 가정에 의해 이 몫환은 체이다
\sref{0.21.1.8-ko}. 이로부터 (i)의 주장들이 따른다.

(ii)를 증명하기 전에 다음 보조정리를 보이자.

\begin{lemma}[22.4.4.2]
\label{0.22.4.4.2-ko}
$\Lambda$를 환, $A$를 $\Lambda$-대수, $\mathfrak J$를
$C=A/\mathfrak J$의 표수가 $p>0$이 되게 하는 $A$의 아이디얼이라
하자. 이는 $A$에서 $p\cdot1\in\mathfrak J$라는 것과 같다.
$\pi$를 $p\cdot1$의 $\mathfrak J/\mathfrak J^2$에서의 표준상이라
하자. $C$가 이산 위상들에 대하여 형식적으로 매끄러운
$(\Lambda/p\Lambda)$-대수이면 다음 표준 완전열이 존재한다.
\begin{equation}
\label{0.22.4.4.3-ko}
  0\longrightarrow(\mathfrak J/\mathfrak J^2)/C\cdot\pi
  \longrightarrow\Omega^1_{A/\Lambda}\otimes_AC
  \longrightarrow\Omega^1_{C/\Lambda}\longrightarrow0.
  \tag{22.4.4.3}
\end{equation}
\end{lemma}

\begin{proof}[보조정리의 증명]
$A'=A/pA$, $\mathfrak J'=\mathfrak J/pA$라 놓으면
$C=A'/\mathfrak J'$이다. 또한
\[
  A'=A\otimes_\Lambda(\Lambda/p\Lambda)
\]
이므로, $\Lambda'=\Lambda/p\Lambda$라 쓰면 표준 동형 아래
\[
  \Omega^1_{A/\Lambda}\otimes_AA'
  =\Omega^1_{A'/\Lambda'}
\]
이다 \sref{0.20.5.5-ko}. 이제 완전열 \sref{0.20.5.14.1-ko}을
형식적으로 매끄러운 $\Lambda'$-대수 $C=A'/\mathfrak J'$에 적용하면
완전열
\[
  0\longrightarrow\mathfrak J'/\mathfrak J'^2
  \longrightarrow\Omega^1_{A'/\Lambda'}\otimes_{A'}C
  \longrightarrow\Omega^1_{C/\Lambda'}\longrightarrow0
\]
을 얻는다. 준동형 $\Lambda\to\Lambda'$은 전사이므로 표준 동형 아래
\[
  \Omega^1_{C/\Lambda'}=\Omega^1_{C/\Lambda}
\]
이다 \sref{0.20.5.15-ko}. 따라서 앞의 사실들로부터 완전열
\sref{0.22.4.4.3-ko}이 따른다.
\end{proof}

(ii) $\mathfrak m'^2=0$이므로, 아이디얼 $\mathfrak m'$을
$K'$-벡터공간으로 본 것도 계속 $V'$으로 나타내고 다음과 같이 놓는다.
\begin{equation}
\label{0.22.4.4.4-ko}
  V_0=V/K\cdot p=\mathfrak m/(\mathfrak m^2+pA),\qquad
  V_0'=V'/K'\cdot p=\mathfrak m'/(\mathfrak m'^2+pA').
  \tag{22.4.4.4}
\end{equation}
완전열 \sref{0.22.4.4.3-ko}에서 $\Lambda=\mathbf Z$로 놓고,
$\mathbf Z$-대수 $A$와 그 아이디얼 $\mathfrak m$에 적용한다. 마찬가지로
$A'$과 $\mathfrak m'$에도 적용한다. $K$와 $K'$은 각각 소체
$\mathbf Z/p\mathbf Z$ 위에서 분리적이므로, 이산 위상에서 형식적으로
매끄러운 $(\mathbf Z/p\mathbf Z)$-대수이다 \sref{0.19.6.1-ko}.
따라서 다음 열들은 완전하다.
\begin{equation}
\label{0.22.4.4.5-ko}
  0\longrightarrow V_0\longrightarrow\Omega^1_A\otimes_AK
  \longrightarrow\Omega^1_K\longrightarrow0,\qquad
  0\longrightarrow V_0'\longrightarrow\Omega^1_{A'}\otimes_{A'}K'
  \longrightarrow\Omega^1_{K'}\longrightarrow0.
  \tag{22.4.4.5}
\end{equation}

이제 도식 \sref{0.22.4.4.6-ko}을 생각하자. 그 가운데 두 열은 각각
\sref{0.22.4.4.5-ko}의 두 번째 완전열과 첫 번째 완전열을 $K'$과
텐서한 열이다. 첫 번째 완전열은 $K$-벡터공간들로 이루어져 있으므로,
도식 \sref{0.22.4.4.6-ko}의 가운데 두 열은 완전하다. 또한 이 두 열과
그들을 잇는 수평 화살표들로 이루어진 도식은
\sref{0.22.4.4.2-ko}의 증명과 도식 \sref{0.20.5.11.3-ko}의
가환성에 의해 가환이다.

도식 \sref{0.22.4.4.6-ko}의 마지막 행은
\sref{0.20.6.1.1-ko}에서 $A,B,C$를 각각
$\mathbf Z,K,K'$으로 바꾸어 얻는 완전열이다. 가운데 행은
\sref{0.20.6.1.1-ko}에서 $A,B,C$를 각각
$\mathbf Z,A,K$로 바꾸어 얻는 $K$-가군 완전열을 $K'$과 텐서하여
얻는다. 이 두 행과 그들을 잇는 수직 화살표들로 이루어진 도식은
\sref{0.20.6.4.3-ko}의 가환성에 의해 가환이다. 한편
\sref{0.22.4.1-ko}, (ii)를 적용할 조건이 성립하므로 도식
\sref{0.22.4.4.6-ko}의 맨 위 행은 완전하다. 따라서 도식의 양 끝 열은

\oldpage[0\textsubscript{IV-1}]{196}

\begin{equation}
\label{0.22.4.4.6-ko}
  \xymatrix{
    && 0 \ar[d] & 0 \ar[d] && \\
    & 0 \ar[r] & V_0\otimes_KK' \ar[r] \ar[d] & V_0' \ar[r] \ar[d] & 0 & \\
    0 \ar[r] & N \ar[r] \ar[d] & \Omega^1_A\otimes_AK' \ar[r] \ar[d] &
      \Omega^1_{A'}\otimes_{A'}K' \ar[r] \ar[d] &
      \Omega^1_{A'/A}\otimes_{A'}K' \ar[r] \ar[d] & 0 \\
    0 \ar[r] & \Upsilon_{K'/K} \ar[r] \ar[d] &
      \Omega^1_K\otimes_KK' \ar[r] \ar[d] & \Omega^1_{K'} \ar[r] \ar[d] &
      \Omega^1_{K'/K} \ar[r] \ar[d] & 0 \\
    & 0 & 0 & 0 & 0 &
  }
  \tag{22.4.4.6}
\end{equation}
각각 가운데 수평 화살표들의 핵과 여핵으로 이루어진다. 이로써 도식의
가환성에 대한 증명이 완성된다.
더욱이 다음이 성립한다.

\begin{lemma}[22.4.4.7]
\label{0.22.4.4.7-ko}
도식 \sref{0.22.4.4.6-ko}의 준동형
\[
  N\longrightarrow\Upsilon_{K'/K}\qquad\text{및}\qquad
  \Omega^1_{A'/A}\otimes_{A'}K'\longrightarrow\Omega^1_{K'/K}
\]
은 전단사이다.
\end{lemma}

\begin{proof}
실제로 \sref{0.22.4.4.6-ko}의 가운데 두 열에 뱀 보조정리의 도식을
적용하면 (Bourbaki, \emph{Alg. comm.}, chap.~I, §~1,
n\textsuperscript{o}~4, prop.~2) 완전열
\[
  0\longrightarrow N\longrightarrow\Upsilon_{K'/K}\longrightarrow0
  \longrightarrow\Omega^1_{A'/A}\otimes_{A'}K'
  \longrightarrow\Omega^1_{K'/K}\longrightarrow0
\]
을 얻는다.

$t_i$를 $T_i$의 $A'$에서의 표준상이라 하자. 관계 $F_i(t_i)=0$과
$A'$에서
\[
  p\cdot1\in\mathfrak mA'\subset\mathfrak m'
\]
이라는 사실로부터 곧
$d_{A'}f_i\in\mathfrak m'\Omega^1_{A'}$를 얻고, 따라서
$d_{A'}f_i\otimes1_{K'}=0$이다. 이는
$d_Af_i\otimes1_{K'}$들이 $N$에 속함을 보인다. 또한 이들의
$\Upsilon_{K'/K}$에서의 상은 $d_Ka_i\otimes1_{K'}$들이며, 이들은
$K'$ 위에서 $\Upsilon_{K'/K}$의 기저를 이룬다
\sref{0.21.4.7-ko}. \sref{0.22.4.4.7-ko}을 함께 쓰면
\sref{0.22.4.4-ko}의 증명이 끝난다.
\end{proof}
\end{proof}

\begin{env}[22.4.5]
\label{0.22.4.5-ko}
이제 표수가 $p>0$인 두 환 $A$, $B$와 두 준동형
\[
  A\xrightarrow{i}B\xrightarrow{j}A
\]

\oldpage[0\textsubscript{IV-1}]{197}

을 생각하자. 이들은 \sref{0.21.3.1-ko}의 조건을 만족한다고 하자.
또한 $\mathfrak m$을 $A$의 제곱영 아이디얼이라 하고 다음과 같이 놓는다.
\[
  K=A/\mathfrak m,\qquad
  B_{(K)}=B\otimes_AK=B/Bi(\mathfrak m).
\]
그리고
\[
  \varphi:A\longrightarrow K,\qquad
  \psi:B\longrightarrow B_{(K)}
\]
를 표준 준동형들이라 하자.

$\varphi$가 전사이므로 $\Omega^1_{B_{(K)}/A}$는
$\Omega^1_{B_{(K)}/K}$와 표준적으로 동일시된다
\sref{0.20.5.15-ko}. $p\geqslant2$이므로 임의의
$x\in\mathfrak m$에 대하여 $j(i(x))=x^p=0$, 즉
$j(i(\mathfrak m))=0$이다. 따라서 $j$를 몫으로 내려 다음 준동형을 얻는다.
\[
  j':B_{(K)}\longrightarrow A.
\]
$i'=\psi\circ i:A\longrightarrow B_{(K)}$라 하면
\[
  j'(i'(x))=x^p\qquad\text{이고}\qquad i'(j'(y))=y^p
  \qquad\text{이다}\quad(x\in A,\ y\in B_{(K)}).
\]
\sref{0.21.3.2-ko}의 표기에 따라
\[
  \Theta_{B_{(K)}/A}
  =\Omega^1_{B_{(K)}/A}\otimes_{B_{(K)}}A_{[j']}
  =\Omega^1_{B_{(K)}/K}\otimes_{B_{(K)}}A_{[j']}
\]
라 놓으면, \sref{0.21.3.2.2-ko}에 의해 표준 준동형
\begin{equation}
\label{0.22.4.5.1-ko}
  \pi_{B_{(K)}/A}:\Theta_{B_{(K)}/A}\longrightarrow\Omega^1_A
  \tag{22.4.5.1}
\end{equation}
을 얻는다.

한편 $j(i(\mathfrak m))=0$임을 고려하여 $i$와 $j$를 몫으로 내리면
준동형
\[
  K\xrightarrow{i_0}B_{(K)}\xrightarrow{j_0}K
\]
을 얻으며, 이에 따라 다음 도식은 가환이다.
\[
  \xymatrix{
    A \ar[r]^i \ar[d]_\varphi & B \ar[r]^j \ar[d]^\psi & A \ar[d]^\varphi \\
    K \ar[r]_{i_0} & B_{(K)} \ar[r]_{j_0} \ar[ur]_{j'} & K
  }
\]
또한 다음이 성립함은 분명하다.
\[
  j_0(i_0(s))=s^p\qquad\text{이고}\qquad i_0(j_0(t))=t^p
  \qquad\text{이다}\quad(s\in K,\ t\in B_{(K)}).
\]
그러므로 마찬가지로
\[
  \Theta_{B_{(K)}/K}
  =\Omega^1_{B_{(K)}/K}\otimes_{B_{(K)}}K_{[j_0]}
  =\Theta_{B_{(K)}/A}\otimes_AK
\]
를 정의할 수 있다. 따라서 \sref{0.22.4.5.1-ko}을 $K$와 텐서하면
표준 $K$-준동형
\begin{equation}
\label{0.22.4.5.2-ko}
  \pi'_{B_{(K)}/K}=\pi_{B_{(K)}/A}\otimes1_K:
  \Theta_{B_{(K)}/K}\longrightarrow\Omega^1_A\otimes_AK
  \tag{22.4.5.2}
\end{equation}
을 얻는다.

\oldpage[0\textsubscript{IV-1}]{198}

정의에서 곧 다음 도식이 가환임을 알 수 있다.
\[
  \xymatrix{
    \Theta_{B_{(K)}/K} \ar[r]^{\pi'_{B_{(K)}/K}} \ar[dr]_{\pi_{B_{(K)}/K}}
      & \Omega^1_A\otimes_AK \ar[d]^{\varphi_{K/A}} \\
    & \Omega^1_K
  }
\]
여기서 $\pi_{B_{(K)}/K}$는 $i_0$와 $j_0$에 대응하는 표준 준동형이다
\sref{0.21.3.2.2-ko}. 따라서 $\pi'_{B_{(K)}/K}$를
$\Xi_{B_{(K)}/K}=\operatorname{Ker}(\pi_{B_{(K)}/K})$에 제한하면
표준 준동형
\begin{equation}
\label{0.22.4.5.3-ko}
  \Xi_{B_{(K)}/K}\longrightarrow
  \operatorname{Ker}(\Omega^1_A\otimes_AK\longrightarrow\Omega^1_K)=\Upsilon_{K/A}
  \tag{22.4.5.3}
\end{equation}
을 얻는다.
\end{env}

\begin{env}[22.4.6]
\label{0.22.4.6-ko}
가정과 표기는 \sref{0.22.4.5-ko}와 같다고 하자. 더 나아가
$\mathfrak m$이 제곱영 극대 아이디얼이고, 따라서 $K$가 체라고 하자.
아이디얼 $\mathfrak m$을 $K$-벡터공간으로 본 것을 $V$라 쓰자.
$K$는 자신의 소체 $\mathbf F_p$ 위에서 형식적으로 매끄러운 대수이고
\sref{0.21.5.2-ko}, $A$는 $\mathbf F_p$-대수이므로
($\Omega^1_A=\Omega^1_{A/\mathbf F_p}$), \sref{0.20.5.14-ko}에서
다음 완전열을 얻는다.
\begin{equation}
\label{0.22.4.6.1-ko}
  0\longrightarrow V\longrightarrow\Omega^1_A\otimes_AK
  \longrightarrow\Omega^1_K\longrightarrow0
  \tag{22.4.6.1}
\end{equation}
따라서 \sref{0.22.4.5.3-ko}에 의해 다음과 같이 쓰는 준동형을 얻는다.
\begin{equation}
\label{0.22.4.6.2-ko}
  \chi'_{B/A}:\Xi_{B_{(K)}/K}\longrightarrow V
  \tag{22.4.6.2}
\end{equation}
이를 $A$-대수 $B$의 ($i$, $j$, $\mathfrak m$에 대한)
\emph{특성 준동형}이라 부르겠다. 정의에 의하면
$\chi'_{B/A}$가 단사이기 위한 필요충분조건은
$\pi'_{B_{(K)}/K}$가 단사인 것이다. 실제로 뒤 준동형의 핵은 분명히
$\Xi_{B_{(K)}/K}$에 포함된다.
\end{env}

\begin{proposition}[22.4.7]
\label{0.22.4.7-ko}
$A$를 극대 아이디얼 $\mathfrak m$이 제곱영인 아르틴 국소환,
$K=A/\mathfrak m$을 그 잉여체, $V$를 아이디얼 $\mathfrak m$을
$K$-벡터공간으로 본 것이라 하자. $A$의 표수는 $p>0$이라고 가정한다.
다음을 놓자.
\begin{equation}
\label{0.22.4.7.1-ko}
  F_i(T_i)=T_i^p-f_i\qquad\text{단 }f_i\in A\quad(1\leqslant i\leqslant r)
  \tag{22.4.7.1}
\end{equation}
$B$를 $A[T_1,\ldots,T_r]$를 $r$개의 다항식 $F_i$가 생성하는
아이디얼로 나눈 몫환이라 하자. 그러면 다음이 성립한다.
\begin{enumerate}
  \item[\rm(i)] $B$는 국소환이다.
  \item[\rm(ii)] $\mathfrak m'$을 $B$의 극대 아이디얼,
  $K'=B/\mathfrak m'$을 그 잉여체, $V'$을 $B$-가군
  $\mathfrak m'/\mathfrak m'^2$을 $K'$-벡터공간으로 본 것이라 하자.
  그러면 $\operatorname{rg}_K(V)=\operatorname{rg}_{K'}(V')$이기 위한
  필요충분조건은 $K$-벡터공간의 특성 준동형
  \begin{equation}
  \label{0.22.4.7.2-ko}
    \chi'_{B/A}:\Xi_{B_{(K)}/K}\longrightarrow V
    \tag{22.4.7.2}
  \end{equation}
  이 단사인 것이다 (\sref{0.22.4.6.2-ko} 참조).
\end{enumerate}
\end{proposition}

\begin{proof}
$a_i$ ($1\leqslant i\leqslant r$)로도 $f_i$의 $\mathfrak m$을 법으로 한
류를 나타내자. $a_i$들이 반드시 $K^p$ 위에서 $p$-독립인 것은 아니지만,
부분족

\oldpage[0\textsubscript{IV-1}]{199}

$(a_i)_{s+1\leqslant i\leqslant r}$이 $K^p$ 위에서 $p$-자유이고,
체 $K^p(a_1,\ldots,a_r)$의 $K^p$ 위의 $p$-기저를 이루도록 언제나
택할 수 있다. 따라서
\[
  a_i\in K^p(a_{s+1},\ldots,a_r)\qquad(1\leqslant i\leqslant s)
\]
이다. $A''$을 $A[T_{s+1},\ldots,T_r]$를 첨자 $i\geqslant s+1$인
$F_i$들이 생성하는 아이디얼로 나눈 몫환이라 하자. 그러면
\sref{0.22.4.4-ko}에 의해 $A''$은 극대 아이디얼이
$\mathfrak m''=\mathfrak mA''$인 국소환이다. 이 아이디얼은 제곱영이고,
잉여체는
\[
  K''=K(a_{s+1}^{1/p},\ldots,a_r^{1/p})
\]
이다. 따라서 $1\leqslant i\leqslant s$에 대하여 $a_i\in K''^p$이고,
$A''$에서 $f_i$는 다음과 같이 쓸 수 있다.
\begin{equation}
\label{0.22.4.7.3-ko}
  f_i=g_i^p+h_i
  \tag{22.4.7.3}
\end{equation}
여기서 $g_i\in A''$이고 $h_i\in\mathfrak m''$이다.
$\mathfrak m''^2=0$이고 $p\mathfrak m''=0$이므로, 이 조건들에 의해
$h_i$가 \emph{유일하게} 결정됨은 곧 알 수 있다. $T_i$를
$T_i+g_i$로 바꾸고
\[
  G_i(T_i)=T_i^p+\sum_{k=1}^{p-1}\binom pk g_i^{p-k}T_i^k-h_i
  \qquad(1\leqslant i\leqslant s)
\]
라 놓으면, $B$는 $A''[T_1,\ldots,T_s]$를 $G_i$들이 생성하는
아이디얼로 나눈 몫환이다. 그런데 $G_i$의 최고차항 계수를 제외한
모든 계수는 $\mathfrak m''$에 속하므로 \sref{0.22.4.3-ko}의 상황에
해당한다. 따라서 $B$는 국소환이고 그 잉여체 $K'$은 $K''$과
\emph{동형}이다. 이로써 (i)이 증명된다. 한편 (i)은
$B^p\subset A$이고, 따라서 $\mathfrak m$ 위에 놓이는 $B$의 아이디얼이
하나뿐이라는 사실에서도 직접 따른다.

이제 \sref{0.22.4.7.2-ko}이 단사이기 위한 필요충분조건은
\[
  \pi'_{B_{(K)}/K}:\Omega^1_{B_{(K)}/K}\otimes_{B_{(K)}}K
  \longrightarrow\Omega^1_A\otimes_AK
\]
가 단사인 것이다 (\sref{0.22.4.6-ko}). 그런데
$B_{(K)}=B/\mathfrak mB$는 $C=K[T_1,\ldots,T_r]$를 다항식
$\overline F_i=T_i^p-a_i$들이 생성하는 아이디얼로 나눈 몫대수이다.
$d_{C/K}\overline F_i=0$이므로 \sref{0.20.5.13-ko}에 의해
$\Omega^1_{B_{(K)}/K}$는 $B_{(K)}$-자유가군이고,
$d_{B_{(K)}/K}t_i$들이 그 기저를 이룬다. 여기서 $t_i$는 $T_i$의
표준상이다 ($1\leqslant i\leqslant r$). 한편 정의
\sref{0.22.4.5-ko}에 의하면 $x\in B$의 $\mathfrak mB$을 법으로 한
류의 $j':B_{(K)}\longrightarrow A$에 의한 표준상은 $x^p\in A$이다.
따라서 $t_i^p=f_i$이므로
$d_{B_{(K)}/K}t_i\otimes1$의 $\pi'_{B_{(K)}/K}$에 의한 상은
$\Omega^1_A\otimes_AK$ 안의 $d_Af_i\otimes1$이다. 그러므로
$\pi'_{B_{(K)}/K}$가 단사이기 위한 필요충분조건은
$d_Af_i\otimes1_K$들이 $K$ 위에서 선형 독립인 것이다. 이는 다시
그 상들 $d_Af_i\otimes1_{K'}$이
$\Omega^1_A\otimes_AK'$에서 $K'$ 위에 선형 독립인 것과 같다.

이제 $A$-대수 $A''$에 \sref{0.22.4.4-ko}를 적용하면, 표준 준동형
\[
  \Omega^1_A\otimes_AK'\longrightarrow
  \Omega^1_{A''}\otimes_{A''}K'
\]
의 핵은 $s+1\leqslant i\leqslant r$인
$d_Af_i\otimes1_{K'}$들을 기저로 갖는다. 따라서 앞의 조건은
$1\leqslant i\leqslant s$인 $d_{A''}f_i\otimes1_{K'}$들이
$\Omega^1_{A''}\otimes_{A''}K'$에서 선형 독립이라는 것과도 같다.
그런데 \sref{0.22.4.7.3-ko}에 의해
$d_{A''}f_i=d_{A''}h_i$이다. 한편 $V''$을 $K'$-벡터공간
$\mathfrak m''$이라 하면
\[
  \operatorname{rg}_{K'}(V'')=\operatorname{rg}_K(V)
\]
이다 \sref{0.22.4.1-ko}, (ii). 따라서
$\operatorname{rg}_{K'}(V')=\operatorname{rg}_K(V)$라는 조건은
$\operatorname{rg}_{K'}(V'')=\operatorname{rg}_{K'}(V')$이라는 조건과
같다. \sref{0.22.4.3-ko}에 의하면 이 마지막 등식은 $h_i$들의 류가

\oldpage[0\textsubscript{IV-1}]{200}

$V''$에서 $K'$ 위에 선형 독립이라는 것과 같다. 그런데
$\Omega^1_{A''}\otimes_{A''}K'$ 안의
$d_{A''}h_i\otimes1_{K'}$들은 단사 준동형
\[
  V''\longrightarrow\Omega^1_{A''}\otimes_{A''}K'
\]
에 의한 $h_i$들의 표준상들이다
\sref{0.20.5.11.2-ko} (\sref{0.22.4.6.1-ko} 참조). 이로써
\sref{0.22.4.7-ko}의 증명이 끝난다.
\end{proof}

더 나아가 다음도 증명하였다.

\begin{corollary}[22.4.7.4]
\label{0.22.4.7.4-ko}
$\operatorname{rg}_K(V)=\operatorname{rg}_{K'}(V')$이기 위한 필요충분조건은
원소 $d_Af_i\otimes1_K$ ($1\leqslant i\leqslant r$)들이
$K$-벡터공간 $\Omega^1_A\otimes_AK$에서 선형 독립인 것이다.
\end{corollary}

\begin{remark}[22.4.8]
\label{0.22.4.8-ko}
준동형 \sref{0.22.4.5.2-ko}과 \sref{0.22.4.6.2-ko}은 조금 다른
조건에서도 정의할 수 있다. $A$를 환, $\mathfrak m$을 $A$의 제곱영
아이디얼, $K=A/\mathfrak m$이라 하고, $p\cdot1\in\mathfrak m$인
소수 $p$를 택하자. 다시 말해 $K$의 표수는 $p$이지만 $A$의 표수는
반드시 $p$일 필요가 없다. 또한 $B$를 $A$-가군으로서 충실하게 평탄한
$A$-대수라 하자. 따라서 $A\subset B$이다. 모든 $y\in B$에 대하여
\begin{equation}
\label{0.22.4.8.1-ko}
  y^p\in A+pB\subset A+\mathfrak mB
  \tag{22.4.8.1}
\end{equation}
이라고 가정하자. $y^p=x+pz$이고 $x\in A$, $z\in B$이면
$x$의 $(A\cap pB)$을 법으로 한 \emph{류}는 잘 결정된다. $B$가
충실하게 평탄한 $A$-가군이므로 $A\cap pB=pA$이다. 따라서 표준 사상
\[
  j:B\longrightarrow A/pA
\]
를 정의할 수 있으며, $x\in A$에 대하여 $j(x)$는 $x^p$의 $pA$를
법으로 한 류이다. 더욱이 $y'\in B$이고
$y'^p=x'+pz'$ ($x'\in A$, $z'\in B$)이면, 이항계수
$\binom pi$가 $1\leqslant i\leqslant p-1$에 대하여 $p$의 배수라는
사실로부터
\[
  (y+y')^p-x-x'\in pB,\qquad (yy')^p-xx'\in pB
\]
임을 곧 알 수 있다. 따라서 $j$는 \emph{환 준동형}이다. 합성하여
$j$로부터 준동형
\[
  j':B\xrightarrow{j}A/pA\longrightarrow A/\mathfrak m=K
\]
를 얻고, 이어서 준동형 $j_0:B_{(K)}\longrightarrow K$를 얻는다.
이때 $j'$은 합성
\[
  B\longrightarrow B_{(K)}=B\otimes_AK\xrightarrow{j_0}K
\]
이다. 또한 $i_0:K\longrightarrow B_{(K)}$가 표준 사상이면,
$i_0$와 $j_0$는 표수가 $p$인 환 $K$와 $B_{(K)}$에 대하여
\sref{0.21.3.1-ko}의 조건을 만족한다.

이제 $K$-준동형
\[
  \pi'_{B_{(K)}/K}:\Omega^1_{B_{(K)}/K}\otimes_{B_{(K)}}K_{[j_0]}
  \longrightarrow\Omega^1_A\otimes_AK
\]
를 정의하려면 \sref{0.20.4.8-ko}에 의해 $K$-미분
\begin{equation}
\label{0.22.4.8.2-ko}
  D_0:B_{(K)}\longrightarrow\Omega^1_A\otimes_AK
  \tag{22.4.8.2}
\end{equation}
를 정의하면 충분하다. 여기서 오른쪽 항은 $j_0$를 통하여
$B_{(K)}$-가군으로 본다. 이를 위해서는 $A$-미분
\begin{equation}
\label{0.22.4.8.3-ko}
  D:B\longrightarrow\Omega^1_A\otimes_AK
  \tag{22.4.8.3}
\end{equation}
로서 $\mathfrak mB$에서 영이 되는 것을 정의하면 충분하다. 여기서
오른쪽 항은 $j'$을 통하여 $B$-가군으로 본다. 실제로
$y^p=x+pz$이고 $y\in B$, $x\in A$, $z\in B$이면,
$\Omega^1_A\otimes_AK$의 원소 $d_Ax\otimes1_K$는 $x$의 $pA$를
법으로 한 류 $j(y)$에만 의존하고, 따라서 $y$에만 의존한다. 왜냐하면
모든 $t\in A$에 대하여 $d_A(pt)=p\,d_A(t)$이므로,
$p\cdot1\in\mathfrak m$인 것에 의해
\[
  d_A(pt)\otimes1=0
\]
이 $\Omega^1_A\otimes_AK=\Omega^1_A/\mathfrak m\Omega^1_A$에서
성립하기 때문이다. 따라서
\[
  D(y)=d_Ax\otimes1_K
\]
로 놓을 수 있다. $D$가 미분이라는 것은 $j$가 환 준동형이라는
사실로부터 쉽게 따른다. 또한 $t\in A$에 대하여
\[
  d_A(t^px)=t^p d_A(x)+pt^{p-1}x d_A(t)
\]
이므로 $D(ty)=tD(y)$를 얻는다. 따라서 $D$는 $A$-미분이다. 끝으로
$t\in\mathfrak m$이기도 하면

\oldpage[0\textsubscript{IV-1}]{201}

$D(ty)=tD(y)=0$이다. 실제로 $\Omega^1_A\otimes_AK$는
$\mathfrak m$에 의해 소멸한다. 이로써 원하던 준동형
\sref{0.22.4.5.2-ko}을 정의하였다. 또한 합성 준동형
\[
  \Theta_{B_{(K)}/K}\xrightarrow{\pi'_{B_{(K)}/K}}
  \Omega^1_A\otimes_AK\longrightarrow\Omega^1_K
\]
이 여전히 $i_0$와 $j_0$에 대한 표준 준동형
$\pi_{B_{(K)}/K}$임은 곧 확인된다. 따라서
\sref{0.22.4.5.3-ko}과 같은 표준 준동형도 얻는다.

이제 $K$가 \emph{체}이면 $A/pA$는 $\mathbf F_p$-대수이고,
\sref{0.20.5.14-ko}에 의해 다음 완전열을 얻는다.
\[
  0\longrightarrow V/pA\longrightarrow
  \Omega^1_{A/pA}\otimes_{A/pA}K\longrightarrow\Omega^1_K\longrightarrow0
\]
한편 \sref{0.20.5.12.1-ko}의 완전열
\[
  pA\longrightarrow\Omega^1_A\otimes_A(A/pA)
  \longrightarrow\Omega^1_{A/pA}\longrightarrow0
\]
을 $K$와 텐서하면 $K$-벡터공간
$\Omega^1_A\otimes_AK$와
$\Omega^1_{A/pA}\otimes_{A/pA}K$가 동형임을 알 수 있다. 따라서 이
경우 \sref{0.22.4.6.2-ko}에 해당하는 것은 다음 준동형이다.
\begin{equation}
\label{0.22.4.8.4-ko}
  \Xi_{B_{(K)}/K}\longrightarrow V/pA=\mathfrak m/(\mathfrak m^2+pA)
  \tag{22.4.8.4}
\end{equation}

이제 \emph{$V$와 $V'$을 각각 $V/pA$와 $V'/pB$로 바꾼
\sref{0.22.4.7-ko}의 판정법은 $K$의 표수가 $p>0$이라고만 가정해도}
(즉 $A$에서 $p\cdot1\in\mathfrak m$이라고만 가정해도), 그리고
\emph{$F_i(T_i)$들이 더 일반적인 꼴 \sref{0.22.4.4.1-ko}이라고
가정해도 그대로 성립한다}. 실제로 $B$는 자유 $A$-가군이므로
충실하게 평탄하다. 따라서 \sref{0.22.4.7-ko}의 논증에서
\sref{0.22.4.6.2-ko}을 \sref{0.22.4.8.4-ko}로, $V''$을
$V''/pA''$로 바꾸어 되풀이하면 충분하다.
\end{remark}

\subsection{기하학적으로 정칙인 대수와 형식적으로 매끄러운 대수}
\label{subsection:0.22.5-ko}

\begin{proposition}[22.5.1]
\label{0.22.5.1-ko}
$A$를 정칙 국소환, $A'$을 $A$를 포함하는 국소환이자 유한
$A$-대수라 하자. $\mathfrak m$과 $\mathfrak m'$을 각각 $A$와 $A'$의
극대 아이디얼, $K=A/\mathfrak m$과 $K'=A'/\mathfrak m'$을 각각의
잉여체라 하자. $A'$이 정칙환이기 위한 필요충분조건은 다음 등식이다.
\begin{equation}
\label{0.22.5.1.1-ko}
  \operatorname{rg}_K(\mathfrak m/\mathfrak m^2)
  =\operatorname{rg}_{K'}(\mathfrak m'/\mathfrak m'^2)
  \tag{22.5.1.1}
\end{equation}
\end{proposition}

\begin{proof}
실제로 \sref{0.22.5.1.1-ko}의 왼쪽 항은 $\dim(A)$이다
\sref{0.17.1.1-ko}. 또한 $A\subset A'$이고 $A'$은 유한
$A$-대수이므로 $\dim(A')=\dim(A)$이다 \sref{0.16.1.5-ko}.
따라서 등식 \sref{0.22.5.1.1-ko}은 $A'$이 정칙이기 위한
필요충분조건이다 \sref{0.17.1.1-ko}.
\end{proof}

\begin{remark}[22.5.2]
\label{0.22.5.2-ko}
\begin{enumerate}
  \item[\rm(i)] 다음과 같이 놓자.
  \[
    A_1=A/\mathfrak m^2,\qquad
    A'_1=A'\otimes_AA_1=A'/\mathfrak m^2A'.
  \]
  $\mathfrak m^2A'\subset\mathfrak m'^2$이므로
  $\mathfrak m'_1=\mathfrak m'/\mathfrak m^2A'$은 $A'_1$의 극대
  아이디얼이고, $\mathfrak m'_1/\mathfrak m_1'^2$은
  $\mathfrak m'/\mathfrak m'^2$과 동형인 $K'$-벡터공간이다. 따라서
  $A'$의 정칙성 조건은 $A'_1$의 $A_1$-대수 구조에만 의존한다.
  이 사실로부터 아래에서 아르틴 환 위의 유한 대수에 관한
  \sref{0.22.4-ko}의 예비 결과들을 적용할 수 있다.
  \item[\rm(ii)] \sref{0.22.5.1-ko}의 증명과
  \sref{0.16.2.6.1-ko}에서 모든 경우에 다음 부등식을 얻는다.
  \begin{equation}
  \label{0.22.5.2.1-ko}
    \operatorname{rg}_K(\mathfrak m/\mathfrak m^2)
    \leqslant\operatorname{rg}_{K'}(\mathfrak m'/\mathfrak m'^2)
    \tag{22.5.2.1}
  \end{equation}

  \oldpage[0\textsubscript{IV-1}]{202}

  \item[\rm(iii)] $A$와 $A'$이 \sref{0.22.4.1-ko}의 일반 가정들을
  만족한다고 하자. \sref{0.19.8.10-ko}에 의해 극대 아이디얼이
  $\mathfrak n$인 정칙 국소환 $B$로서
  $A\simeq B/\mathfrak n^2$인 것이 존재한다. $G_i\in B[T_i]$를
  $A[T_i]$에서의 표준상이 $F_i$인 일계수 다항식이라 하자
  ($1\leqslant i\leqslant r$). $B'$을 $B[T_1,\ldots,T_r]$를
  $G_i$들이 생성하는 아이디얼로 나눈 몫환이라 하자. 그러면 $B'$은
  계수가 $d$인 자유 $B$-가군이고
  \[
    A'=B'\otimes_BA=B'/\mathfrak n^2B'
  \]
  임은 분명하다. $A'$이 국소환이라고 가정했으므로 $B'$도 국소환이다.
  $\mathfrak n'$을 $B'$의 극대 아이디얼이라 하면
  $B'/\mathfrak n'\simeq K'$이고,
  $\mathfrak n'/\mathfrak n'^2$은 $K'$-벡터공간으로서
  $\mathfrak m'/\mathfrak m'^2$과 동형이다. 따라서 부등식
  \sref{0.22.5.2.1-ko}에 의해 \sref{0.22.4.1-ko}의 가정 아래
  \begin{equation}
  \label{0.22.5.2.2-ko}
    \operatorname{rg}_K(V)
    \leqslant\operatorname{rg}_{K'}(\mathfrak m'/\mathfrak m'^2)
    \tag{22.5.2.2}
  \end{equation}
  이다.
\end{enumerate}
\end{remark}

\begin{corollary}[22.5.3]
\label{0.22.5.3-ko}
$A$를 정칙 국소환, $\mathfrak m$을 그 극대 아이디얼,
$T_i$ ($1\leqslant i\leqslant r$)를 부정원이라 하자. 각 $i$에 대하여
\begin{equation}
\label{0.22.5.3.1-ko}
  F_i(T_i)=T_i^{d_i}+\sum_{1\leqslant k\leqslant d_i}c_{ik}T_i^{d_i-k}
  \tag{22.5.3.1}
\end{equation}
을 $A[T_i]$의 일계수 다항식이라 하자. 모든 $i$에 대하여
$d_i\geqslant2$이고 $c_{ik}\in\mathfrak m$이라고 가정한다. $A'$을
다항식환 $A[T_1,\ldots,T_r]$을 $r$개의 다항식 $F_i$가 생성하는
아이디얼로 나눈 몫환이라 하자. 그러면 다음이 성립한다.
\begin{enumerate}
  \item[\rm(i)] $A'$은 국소환이고, 그 잉여체 $K'$은 $A$의 잉여체
  $K$와 동형이다.
  \item[\rm(ii)] $A'$이 정칙환이기 위한 필요충분조건은 원소
  $c_{i,d_i}$의 $\mathfrak m^2$을 법으로 한 류 $\xi_i$
  ($1\leqslant i\leqslant r$)들이 $K$ 위에서 선형 독립인 것이다.
\end{enumerate}
\end{corollary}

\begin{proof}
실제로 \sref{0.22.5.2-ko}, (i)의 표기를 사용하자. $\overline F_i$를
$F_i$의 계수들에 준동형 $A\longrightarrow A_1$을 적용하여 얻는
$A_1[T_i]$의 다항식이라 하면, $A'_1$은 $A_1$과 다항식
$\overline F_i$들로부터 \sref{0.22.4.1-ko}과 같이 정의된다. 따라서
결론은 \sref{0.22.5.1-ko}과 \sref{0.22.4.3-ko}에서 따른다.
\end{proof}

\begin{theorem}[22.5.4]
\label{0.22.5.4-ko}
$A$를 정칙 국소환, $K$를 그 잉여체라 하고, $K$의 표수가
$p>0$이라고 가정하자. 다음과 같이 놓자.
\begin{equation}
\label{0.22.5.4.1-ko}
  F_i(T_i)=T_i^p+p\sum_{k=1}^{p-1}c_{ik}T_i^{p-k}-f_i
  \tag{22.5.4.1}
\end{equation}
여기서 $c_{ik}\in A$, $f_i\in A$이다
($1\leqslant i\leqslant r$). $B$를 $A[T_1,\ldots,T_r]$을 $r$개의
다항식 $F_i$가 생성하는 아이디얼로 나눈 몫환이라 하자. 그러면
$B$는 국소환이고, 다음 조건들은 서로 동치이다.
\begin{enumerate}
  \item[a)] $B$는 정칙이다.
  \item[b)] 원소 $d_Af_i\otimes1_K$ ($1\leqslant i\leqslant r$)들이
  $K$-벡터공간 $\Omega^1_A\otimes_AK$에서 선형 독립이다.
  \item[c)] $B'=B/\mathfrak m^2B$라 하고
  $\Xi_{B_{(K)}/K}=\Xi_{B'_{(K)}/K}$라 놓으면
  \sref{0.22.4.5-ko}, 특성 준동형
  \[
    \Xi_{B_{(K)}/K}\longrightarrow\mathfrak m/\mathfrak m^2
  \]
  이 단사이다 \sref{0.22.4.6.2-ko}.
\end{enumerate}
\end{theorem}

\begin{proof}
실제로 $B'$도 $A_1=A/\mathfrak m^2$과 $F_i$의 계수들에 준동형
$A\longrightarrow A_1$을 적용하여 얻은 다항식 $\overline F_i$들로부터
정의된다. $\Omega^1_A\otimes_AK$가
$\Omega^1_{A_1}\otimes_{A_1}K$와 표준적으로 동형임을 고려하면
\sref{0.20.5.12-ko}, (ii), 정리는 \sref{0.22.5.1-ko}과
\sref{0.22.4.7.4-ko}에서 따른다. $A_1$의 표수가 $p$가 아닐 때에는
비고 \sref{0.22.4.8-ko}을 사용한다.
\end{proof}

\oldpage[0\textsubscript{IV-1}]{203}

\begin{env}[22.5.5]
\label{0.22.5.5-ko}
$k$를 표수가 $p>0$인 체, $A$를 국소환, $k'$을 $k'^p\subset k$인
$k$의 \emph{유한 순수 비분리 확대체}, $\mathfrak m$을 $A$의 극대
아이디얼, $K=A/\mathfrak m$을 그 잉여체라 하자. 다음 환
\[
  A'=A\otimes_kk'
\]
은 국소환이고, 그 잉여체는 국소환 $K\otimes_kk'$의 잉여체와 같음을
상기하자 (Bourbaki, \emph{Alg. comm.}, chap.~V, §~2,
n\textsuperscript{o}~3, lemme~4). 또한
\[
  A'_{(K)}=A'\otimes_AK=K\otimes_kk'
\]
이므로 \sref{0.21.3.6-ko}에 의해 표준 동형을 제외하면
\begin{equation}
\label{0.22.5.5.1-ko}
  \Theta_{A'_{(K)}/K}=\Theta_{k'/k}\otimes_kK
  \tag{22.5.5.1}
\end{equation}
이다. 더욱이 \sref{0.21.4.7-ko}에 의해 $\Xi_{k'/k}=0$이다. 다시 말해
표준 준동형 \sref{0.21.3.2.2-ko}
\[
  \pi_{k'/k}:\Theta_{k'/k}\longrightarrow\Omega^1_k
\]
는 단사이다. 따라서 다음 \emph{표준 단사 준동형}을 얻는다.
\begin{equation}
\label{0.22.5.5.2-ko}
  \Theta_{A'_{(K)}/K}\longrightarrow\Omega^1_k\otimes_kK
  \tag{22.5.5.2}
\end{equation}
이는 \sref{0.22.5.5.1-ko}을 고려하면 다음과 같이 구체적으로 주어진다.
모든 $x'\in k'$에 대하여
$\Omega^1_{A'_{(K)}/K}\otimes_{A'_{(K)}}K$의 원소
\[
  d_{A'_{(K)}/K}(x'\otimes1_K)\otimes1_K
\]
에 $\Omega^1_k\otimes_kK$의 원소 $d_k(x'^p)\otimes1_K$를 대응시킨다.
\end{env}

\begin{lemma}[22.5.6]
\label{0.22.5.6-ko}
\sref{0.22.5.5-ko}의 표기를 사용하면 다음이 성립한다.
\begin{enumerate}
  \item[\rm(i)] 표준 준동형 $\pi_{A'_{(K)}/K}$
  \sref{0.21.3.2.2-ko}의 핵은 \sref{0.22.5.5.2-ko}에 의해
  \[
    (\Theta_{k'/k}\otimes_kK)\cap\Upsilon_{K/k}
  \]
  와 동일시된다.
  \item[\rm(ii)] $A_1=A/\mathfrak m^2$,
  $A'_1=A_1\otimes_kk'=A'/\mathfrak m^2A'$라 하자. 그러면 특성 준동형
  $\chi'_{A'_1/A_1}$ \sref{0.22.4.6.2-ko}은 특성 준동형
  \[
    \chi_{A/k}:\Upsilon_{K/k}\longrightarrow V=\mathfrak m/\mathfrak m^2
  \]
  을 $(\Theta_{k'/k}\otimes_kK)\cap\Upsilon_{K/k}$에 제한한 것과
  동일시된다 \sref{0.20.6.24-ko}.
\end{enumerate}
\end{lemma}

\begin{proof}
(i) \sref{0.20.6.9-ko}에서 $A$를 소체로, $B$를 $k$로, $C$를
$K$로, $E$를 $V=\mathfrak m/\mathfrak m^2$에 의한 $K$의 확대
$A_1=A/\mathfrak m^2$로 바꾸어 적용한다. 그러면
\sref{0.22.5.5-ko}도 고려할 때 다음 사실만 확인하면 충분하다.
$x'\in k'$이고 $z\in A_1$의 $\mathfrak m/\mathfrak m^2$을 법으로 한
류가 $x'^p\in k$이면, 표준 준동형
\[
  \Omega^1_k\otimes_kK\longrightarrow
  \Omega^1_{A_1}\otimes_{A_1}K
\]
에 의한 $d_k(x'^p)\otimes1_K$의 상은 $d_{A_1}z\otimes1_K$이다.
이는 정의에서 곧 따른다.

(ii) $K=A_1/(\mathfrak m/\mathfrak m^2)$이므로
\[
  A'_{(K)}=K\otimes_kk'=(A'_1)_{(K)}.
\]
따라서 $\Theta_{A'_{(K)}/K}$는
$\Theta_{(A'_1)_{(K)}/K}$와 동일시된다. 이제 (i)과 같은 계산 및
$\pi'_{(A'_1)_{(K)}/K}$의 정의 \sref{0.22.4.5.2-ko}에서 결론이 따른다.
\end{proof}

\begin{proposition}[22.5.7]
\label{0.22.5.7-ko}
$k$를 표수가 $p>0$인 체, $A$를 정칙 국소환인 $k$-대수,
$\mathfrak m$을 그 극대 아이디얼, $K=A/\mathfrak m$을 그 잉여체라
하고, $V=\mathfrak m/\mathfrak m^2$을 $K$-벡터공간으로 보자.
$k'$을 $k'^p\subset k$인 $k$의 유한 확대체라 하고
$A'=A\otimes_kk'$라 하자. 국소환 $A'$이 정칙이기 위한 필요충분조건은
특성 준동형
\[
  \chi_{A/k}:\Upsilon_{K/k}\longrightarrow V
\]
\sref{0.20.6.24-ko}을
$(\Theta_{k'/k}\otimes_kK)\cap\Upsilon_{K/k}$에 제한한 것이 단사인
것이다.
\end{proposition}

\begin{proof}
$\mathfrak m'$을 $A'$의 극대 아이디얼, $K'=A'/\mathfrak m'$을 그
잉여체라 할 때, 명제의 조건이 \sref{0.22.5.1.1-ko}과 동치임을 보이면
된다. 비고 \sref{0.22.5.2-ko}, (i)의 표기로
$A'_1=A_1\otimes_kk'$이다. 이제 $(a_i)_{1\leqslant i\leqslant r}$를
$k'$의 $k$ 위의 $p$-기저라 하자. 그러면 $a_i^p=b_i\in k$이고,
$k'$은 다항식환

\oldpage[0\textsubscript{IV-1}]{204}

$k[T_1,\ldots,T_r]$을 다항식 $F_i=T_i^p-b_i$
($1\leqslant i\leqslant r$)가 생성하는 아이디얼로 나눈 몫과 동형이다.
따라서 $A'_1$은 $A_1[T_1,\ldots,T_r]$을 $F_i$들이 생성하는
아이디얼로 나눈 몫과 동형이다. 이제 판정법 \sref{0.22.4.7-ko}을
적용할 수 있고, 명제의 조건은 $\chi'_{A'_1/A_1}$이 단사인 것과
동치이다. 결론은 \sref{0.22.5.6-ko}에서 따른다.
\end{proof}

이제 \sref{0.19.6.5-ko}의 역을 증명할 수 있다.

\begin{theorem}[22.5.8]
\label{0.22.5.8-ko}
$k$를 체, $p$를 그 표수 지수, $A$를 국소 뇌터 $k$-대수라 하자.
다음 조건들은 서로 동치이다.
\begin{enumerate}
  \item[a)] $A$는 그 준아딕 위상에 대하여 형식적으로 매끄러운
  $k$-대수이다.
  \item[b)] $A$는 $k$ 위에서 기하학적으로 정칙이다.
  \item[b')] $k'^p\subset k$인 모든 $k$의 유한 확대체 $k'$에 대하여
  $A'=A\otimes_kk'$은 정칙환이다.
\end{enumerate}
\end{theorem}

\begin{proof}
\sref{0.19.6.6-ko}을 고려하면 $p>1$인 경우만 다루면 충분하다.

a)가 b)를 함의한다는 것은 바로 \sref{0.19.6.5-ko}이고, b)는
자명하게 b')을 함의한다. b')에서 \sref{0.22.2.2-ko}의 조건들이
성립함을 보이자. b')이 먼저 $A$가 정칙임을 함의하므로 첫 번째 조건은
\sref{0.17.1.1-ko}에 의해 분명하다. 한편
$(x_\alpha)_{\alpha\in I}$를 $k^{1/p}$의 $k$ 위의 $p$-기저라 하자.
원소 $d_k(x_\alpha^p)$들은 $k$-벡터공간 $\Omega^1_k$의 기저를
이룬다 \sref{0.21.4.5-ko}. 따라서
$d_k(x_\alpha^p)\otimes1_K$들은 $K$-벡터공간
$\Omega^1_k\otimes_kK$의 기저를 이룬다. 이로부터
\sref{0.22.5.5-ko}도 고려하면, $k'$이 $k^{1/p}$의 $k$ 위의 유한
부분확대체들을 달릴 때 부분공간 $\Theta_{k'/k}\otimes_kK$들의 족은
증가 여과족이고 그 합집합은 $\Omega^1_k\otimes_kK$임을 알 수 있다.
따라서 \sref{0.22.5.7-ko}에 의해 b')은 $\chi_{A/k}$가 단사임을
함의한다. 이것은 바로 \sref{0.22.2.2-ko}의 조건 (ii)이다.
\end{proof}

\begin{corollary}[22.5.9]
\label{0.22.5.9-ko}
$k$를 체, $A$를 형식적으로 매끄러운 국소 뇌터 $k$-대수라 하자.
그러면 $A$의 모든 소아이디얼 $\mathfrak p$에 대하여 국소환
$A_{\mathfrak p}$은 $\mathfrak p$-준아딕 위상에 관하여 형식적으로
매끄러운 $k$-대수이다.
\end{corollary}

\begin{proof}
실제로 \sref{0.22.5.8-ko}, b')의 표기 아래 국소환
$A_{\mathfrak p}\otimes_kk'$이 정칙임을 보이면 충분하다. 이 환은
$A\otimes_kk'$의 분수환이고, 뒤의 환은 가정에 의해 정칙이므로 앞의
환도 정칙이다 \sref{0.17.3.6-ko}.
\end{proof}

\begin{remark}[22.5.10]
\label{0.22.5.10-ko}
\begin{enumerate}
  \item[\rm(i)] 조건 b')에서 $k$의 \emph{단순 확대체}
  $k'=k(x)$ ($x^p\in k$)만 고려하면 \sref{0.22.5.8-ko}의 결론은
  성립하지 않는다. 실제로 모든 이러한 확대체 $k'$에 대하여
  \[
    \Upsilon_{K/k}\cap(\Theta_{k'/k}\otimes_kK)=0
  \]
  이지만 $\Upsilon_{K/k}\neq0$인 경우가 있을 수 있다. 이는
  $x\notin k$인 모든 $x\in k^{1/p}$에 대하여
  $d_k(x^p)\otimes1_K\notin\Upsilon_{K/k}$, 또는 동치로
  $d_K(x^p)\neq0$, 즉 $x^p\notin K^p$이어야 한다는 뜻이다. 다시 말해
  $K$가 $k$의 \emph{비분리} 확대체인데도 $K^p\cap k=k^p$이어야 한다.

  이러한 확대체의 예는 쉽게 구성할 수 있다. 표수가 $p>0$인 완전체
  $k_0$에서 시작하여 $X,Y,Z$를 세 부정원이라 하고
  \[
    k=k_0(X,Y),\qquad K=k(X^{1/p}+ZY^{1/p})
  \]
  라 놓자. $K$가 앞의 조건들을 만족하고 따라서
  $\Upsilon_{K/k}\neq0$임은 쉽게 확인된다. 이제
  \sref{0.22.2.5.2-ko}에서 $V=K$로 놓고 $h$를 영준동형으로 하여
  적용하자. 그러면 $A$는 잉여체가 $K$인 이산 값매김환이고
  $\chi_{A/k}=0$이다. 따라서 \sref{0.22.2.2-ko}에 의해 $A$는
  형식적으로 매끄러운 $k$-대수가 아니다. 그런데도 $k^{1/p}$ 안의
  모든 $k$의 단순 부분확대체 $k'$에 대하여 $A\otimes_kk'$은
  정칙환이다.

  \oldpage[0\textsubscript{IV-1}]{205}

  \item[\rm(ii)] 따름정리 \sref{0.22.5.9-ko}는 다음 문제를 생각하게
  한다. $A$가 뇌터 $k$-대수일 때, $A_{\mathfrak p}$이 형식적으로
  매끄러운 $k$-대수가 되는 소아이디얼
  $\mathfrak p\in\operatorname{Spec}(A)$들의 집합은 어느 조건에서
  \emph{열린집합}인가? 우리는 나중에 몇몇 특수한 경우를 다룰 것이다.
\end{enumerate}
\end{remark}

\subsection{자리스키의 야코비 판정법}
\label{subsection:0.22.6-ko}

\begin{theorem}[22.6.1]
\label{0.22.6.1-ko}
\textnormal{(형식적 매끄러움에 대한 야코비 판정법.)}
$A,B$를 두 위상환, $u:A\longrightarrow B$를 $B$가 형식적으로
매끄러운 $A$-대수가 되게 하는 연속 준동형, $\mathfrak J$를 $B$의
아이디얼(닫혀 있을 필요는 없다), $C=B/\mathfrak J$를 그 몫 위상
$A$-대수라 하자. $C$가 형식적으로 매끄러운 $A$-대수일 필요충분조건은
정준 준동형(\sref{0.20.5.11.2-ko} 참조)
\begin{equation}
\label{0.22.6.1.1-ko}
  \delta_{C/B/A}:\mathfrak J/\mathfrak J^2
  \longrightarrow\Omega^1_{B/A}\otimes_BC
  \tag{22.6.1.1}
\end{equation}
이 \emph{형식적으로 왼쪽 가역}인 것이다
(\sref{0.19.1.5-ko} 참조).
\end{theorem}

\begin{proof}
실제로 $B$ (각각 $C$)가 형식적으로 매끄러운 $A$-대수라는 것은
$B$ (각각 $C$)의 열린 아이디얼에 의해 소멸되는 모든 이산
$B$-가군 (각각 $C$-가군) $L$에 대하여
\[
  \operatorname{Exalcotop}_A(B,L)=0
  \quad\text{(각각 }\operatorname{Exalcotop}_A(C,L)=0\text{)}
\]
이라는 뜻이다 \sref{0.19.4.4-ko}. 가정에 의해 $C$의 열린 아이디얼에
의해 소멸되는 모든 이산 $C$-가군 $L$에 대하여
$\operatorname{Exalcotop}_A(B,L)=0$이므로, $C$가 형식적으로 매끄러운
$A$-대수라는 것은 정준 준동형
\[
  \operatorname{Exalcotop}_A(C,L)\longrightarrow
  \operatorname{Exalcotop}_A(B,L)
\]
이 단사라는 뜻이다. 따라서 정리는 \sref{0.20.7.8-ko}에서 따른다.
\end{proof}

$B$와 $C$가 형식적으로 매끄러운 $A$-대수일 때
$\mathfrak J/\mathfrak J^2$은 형식적으로 사영인 위상 $C$-가군임을
상기하자 \sref{0.19.5.3-ko}. 또한 $B$가 \emph{준허용환}이라고
가정하면 정준 준동형
\[
  \varphi_n:\mathbf S_C^n(\mathfrak J/\mathfrak J^2)
  \longrightarrow\mathfrak J^n/\mathfrak J^{n+1}
\]
들은 \emph{형식적 전단사 준동형}이다.

\begin{corollary}[22.6.2]
\label{0.22.6.2-ko}
\sref{0.22.6.1-ko}의 가정과 표기를 유지하자. 더 나아가 $B$의 모든
열린 아이디얼의 제곱이 열려 있고, $\mathfrak J$ 위의 부분공간 위상이
$B$의 위상에서 유도된 위상 \sref{0.19.0-ko}과 같다고 하자. (이 두
조건은 $B$가 뇌터이고 그 위상이 준아딕일 때
$(\mathbf 0_{\mathrm I},~7.3.2)$, 또는 $B$의 위상이
$\mathfrak J$-준아딕 위상일 때 성립한다.)
$(\mathfrak b_\lambda)$를 $B$의 열린 아이디얼 기본계라 하고 모든
$\lambda$에 대하여 $B_\lambda=B/\mathfrak b_\lambda$라 놓자. 그러면
다음이 성립한다.
\begin{enumerate}
  \item[\rm(i)] 다음 조건들은 서로 동치이다.
  \begin{enumerate}
    \item[a)] $C$는 형식적으로 매끄러운 $A$-대수이다.
    \item[b)] 모든 $\lambda$에 대하여 $\delta_{C/B/A}$에 텐서곱을
    취하여 얻은 $B_\lambda$-가군 준동형
    \begin{equation}
    \label{0.22.6.2.1-ko}
      (\mathfrak J/\mathfrak J^2)\otimes_BB_\lambda
      \longrightarrow\Omega^1_{B/A}\otimes_BB_\lambda
      \tag{22.6.2.1}
    \end{equation}
    은 왼쪽 가역이다. 다시 말해
    $\Omega^1_{B/A}\otimes_BB_\lambda$의 한 직합인자로 가는
    동형사상이다.
  \end{enumerate}
  \item[\rm(ii)] 더 나아가 위상환 $C$가 준허용환
  $(\mathbf 0_{\mathrm I},~7.1.2)$이고

  \oldpage[0\textsubscript{IV-1}]{206}

  $\mathfrak L$이 $C$의 정의 아이디얼이라고 하자. 그러면 (i)의 조건
  a), b)는 다음 조건과도 동치이다.
  \begin{enumerate}
    \item[c)] $(C/\mathfrak L)$-가군 준동형
    \[
      (\mathfrak J/\mathfrak J^2)\otimes_C(C/\mathfrak L)
      \longrightarrow\Omega^1_{B/A}\otimes_B(C/\mathfrak L)
    \]
    은 왼쪽 가역이다.
  \end{enumerate}
\end{enumerate}
\end{corollary}

\begin{proof}
가정에 의해 $\Omega^1_{B/A}$의 위상은 $B$의 위상에서 유도된 위상이다
\sref{0.20.4.5-ko}. 따라서 (i)는 \sref{0.22.6.1-ko}과
\sref{0.19.1.7-ko}에서 따른다. 한편 $B$가 형식적으로 매끄러운
$A$-대수이므로 $\Omega^1_{B/A}$는 형식적으로 사영인 $B$-가군임을
상기하자 \sref{0.20.4.9-ko}. 따라서 모든 $\lambda$에 대하여
$\Omega^1_{B/A}\otimes_BB_\lambda$는 사영 $B_\lambda$-가군이다
\sref{0.19.2.4-ko}. 이제 $C$가 준허용환일 때 c)가 a), b)와 동치라는
것은 \sref{0.19.1.9-ko}에서 따른다.
\end{proof}

특히 다음을 얻는다.

\begin{corollary}[22.6.3]
\label{0.22.6.3-ko}
$A$를 환, $B$를 $A$-대수, $\mathfrak J$를 $B$의 아이디얼,
$C=B/\mathfrak J$라 하자. $A,B,C$에는 이산 위상을 주고, $B$가
형식적으로 매끄러운 $A$-대수라고 가정한다. $C$가 형식적으로 매끄러운
$A$-대수일 필요충분조건은 정준 준동형 \sref{0.22.6.1.1-ko}이 왼쪽
가역인 것이다.
\end{corollary}

모든 $A$-대수 $C$는 다항식 대수 $A[T_\alpha]_{\alpha\in I}$의
몫이고, 이 다항식 대수는 형식적으로 매끄러우므로
\sref{0.19.3.3-ko}, 판정법 \sref{0.22.6.3-ko}을 이용하면 원칙적으로
$C$가 형식적으로 매끄러운 $A$-대수인지 판정할 수 있다.

\begin{proposition}[22.6.4]
\label{0.22.6.4-ko}
$A$를 환, $B$를 (이산 위상들에 관하여) 형식적으로 매끄러운
$A$-대수, $\mathfrak J$를 $B$의 아이디얼, $C=B/\mathfrak J$라 하자.
$\mathfrak J/\mathfrak J^2$이 유한형 $C$-가군이라고 가정한다.
$\mathfrak p$를 $C$의 소아이디얼, $k(\mathfrak p)$를
$C_{\mathfrak p}$의 잉여체, $\mathfrak p'$을 $B$에서
$\mathfrak p$의 역상, $\mathfrak q$를 $A$에서 $\mathfrak p'$의
역상인 소아이디얼이라 하자. 다음 조건들은 서로 동치이다.
\begin{enumerate}
  \item[a)] $C_{\mathfrak p}$은 (이산 위상들에 관하여) 형식적으로
  매끄러운 $A_{\mathfrak q}$-대수이다.
  \item[a')] $C_{\mathfrak p}$은 (이산 위상들에 관하여) 형식적으로
  매끄러운 $A$-대수이다.
  \item[b)] 정준 준동형
  \begin{equation}
  \label{0.22.6.4.1-ko}
    (\mathfrak J/\mathfrak J^2)\otimes_Ck(\mathfrak p)
    \longrightarrow\Omega^1_{B/A}\otimes_Bk(\mathfrak p)
    \tag{22.6.4.1}
  \end{equation}
  은 단사이다.
  \item[c)] $C_f$가 형식적으로 매끄러운 $A$-대수가 되게 하는
  $f\in C-\mathfrak p$가 존재한다.
\end{enumerate}
더 나아가 $B$가 뇌터이면, 앞의 조건들은 다음 조건과도 동치이다.
\begin{enumerate}
  \item[d)] $C_{\mathfrak p}$은 $C_{\mathfrak p}$의
  $\mathfrak p$-준아딕 위상과 $A_{\mathfrak q}$의 이산 위상(또는
  $\mathfrak q$-준아딕 위상)에 관하여 형식적으로 매끄러운
  $A_{\mathfrak q}$-대수이다.
\end{enumerate}
\end{proposition}

\begin{proof}
$A_{\mathfrak q}$가 형식적으로 매끄러운 $A$-대수라는 것은
\sref{0.19.3.5-ko}, (iv) 및 (ii)에서 따른다. 따라서 a)와 a')이
동치임은 이미 알려져 있다 \sref{0.19.3.5-ko}, (ii). 또한
\[
  C_{\mathfrak p}=B_{\mathfrak p'}/\mathfrak J B_{\mathfrak p'}
\]
이고, $B_{\mathfrak p'}$은 이산 위상에 관하여 형식적으로 매끄러운
$A$-대수이다 \sref{0.19.3.5-ko}, (iv). 더 나아가
\[
  \Omega^1_{B_{\mathfrak p'}/A}
  =(\Omega^1_{B/A})_{\mathfrak p'}
\]
이다 \sref{0.20.5.9-ko}. 이제 $B_{\mathfrak p'}$과
$C_{\mathfrak p}$에 \sref{0.22.6.3-ko}을 적용하고,
$\Omega^1_{B/A}$가 사영 $B$-가군이라는 사실
\sref{0.20.4.9-ko}, \sref{0.19.2.11-ko} 및
$\mathfrak J/\mathfrak J^2$이 유한형 $B$-가군이라는 사실을 고려하여
\sref{0.19.1.12-ko}을 쓰면 a')과 b)의 동치가 따른다.
\sref{0.19.1.12-ko}을 적용하면 b)와 c)의 동치도 얻는다.

끝으로 $B_{\mathfrak p'}$은 이산 위상들에 관하여 형식적으로 매끄러운
$A_{\mathfrak q}$-대수이므로 \sref{0.19.3.5-ko}, (iv),

\oldpage[0\textsubscript{IV-1}]{207}

$B_{\mathfrak p'}$에 $\mathfrak p'$-준아딕 위상을 주고
$A_{\mathfrak q}$에 이산 위상 또는 $\mathfrak q$-준아딕 위상을
주어도 형식적으로 매끄러운 $A_{\mathfrak q}$-대수이다
\sref{0.19.3.8-ko}. $B$가 뇌터일 때 b)와 d)의 동치를 보이기 위해,
이산 위상(또는 $\mathfrak q$-준아딕 위상)을 준 환
$A_{\mathfrak q}$와 $\mathfrak p'$-준아딕 위상을 준 환
$B_{\mathfrak p'}$에 \sref{0.22.6.2-ko}을 적용하자.
$C_{\mathfrak p}$은 준허용환이고
$\mathfrak pC_{\mathfrak p}$은 그 위상의 정의 아이디얼이므로,
\sref{0.22.6.2-ko}의 c)와 a)의 동치를 적용할 수 있다.
\end{proof}

\begin{corollary}[22.6.5]
\label{0.22.6.5-ko}
\sref{0.22.6.4-ko}의 일반 가정 아래, 이산 위상들에 관하여
$C_{\mathfrak p}$이 형식적으로 매끄러운 $A$-대수인 (또는
$\mathfrak q$가 $A$에서 $\mathfrak p$의 역상일 때 형식적으로
매끄러운 $A_{\mathfrak q}$-대수인)
$\mathfrak p\in\operatorname{Spec}(C)$들의 집합은
$\operatorname{Spec}(C)$에서 열린집합이다.
\end{corollary}

이는 \sref{0.22.6.4-ko}의 조건 c)에서 따른다. $B$가 뇌터이면 이산
위상들을 준아딕 위상들로 바꿀 수 있음에 유의하자.

\begin{corollary}[22.6.6]
\label{0.22.6.6-ko}
\sref{0.22.6.4-ko}의 일반 가정 아래 다음 조건들은 서로 동치이다.
\begin{enumerate}
  \item[a)] $C$는 (이산 위상들에 관하여) 형식적으로 매끄러운
  $A$-대수이다.
  \item[b)] 모든 $\mathfrak p\in\operatorname{Spec}(C)$에 대하여
  (또는 $C$의 모든 극대 아이디얼 $\mathfrak p$에 대해서만)
  $C_{\mathfrak p}$은 이산 위상들에 관하여 형식적으로 매끄러운
  $A$-대수이다. 동치로, $\mathfrak q$가 $A$에서 $\mathfrak p$의
  역상일 때 형식적으로 매끄러운 $A_{\mathfrak q}$-대수이다.
\end{enumerate}
더 나아가 $B$가 뇌터이면, b)에서 이산 위상들을
$A_{\mathfrak q}$와 $C_{\mathfrak p}$의 준아딕 위상들로 바꿀 수
있다.
\end{corollary}

\begin{proof}
마지막 주장은 \sref{0.22.6.4-ko}에서 따른다.
\sref{0.22.6.3-ko}에 의해 조건 a)는 준동형
\sref{0.22.6.1.1-ko}이 왼쪽 가역이라는 것과 동치이다. 따라서
$\Omega^1_{B/A}$가 사영 $B$-가군이고
$\mathfrak J/\mathfrak J^2$이 유한형 $C$-가군임을 고려하면,
a)와 b)의 동치는 \sref{0.19.1.14-ko}에서 따른다.
\end{proof}

\begin{proposition}[22.6.7]
\label{0.22.6.7-ko}
$k_0$를 체, $k$를 $k_0$의 분리 확대체, $C$를 유한형 $k$-대수라
하자.
\begin{enumerate}
  \item[\rm(i)] $\mathfrak p$를 $C$의 소아이디얼이라 하자. 다음
  조건들은 서로 동치이다.
  \begin{enumerate}
    \item[a)] $C_{\mathfrak p}$은 이산 위상에 관하여 형식적으로
    매끄러운 $k_0$-대수이다.
    \item[b)] $C_{\mathfrak p}$은 $\mathfrak p$-준아딕 위상에
    관하여 형식적으로 매끄러운 $k_0$-대수이다.
    \item[c)] $C_f$가 이산 위상에 관하여 형식적으로 매끄러운
    $k_0$-대수가 되게 하는 $f\in C-\mathfrak p$가 존재한다.
    \item[d)] $C_{\mathfrak p}$은 $k_0$ 위에서 기하학적으로 정칙인
    환이다 \sref{0.19.6.5-ko}.
  \end{enumerate}
  더 나아가 이 조건들 가운데 어느 하나를 만족하는
  $\mathfrak p\in\operatorname{Spec}(C)$들의 집합은
  $\operatorname{Spec}(C)$에서 열린집합이다.
  \item[\rm(ii)] 다음 조건들은 서로 동치이다.
  \begin{enumerate}
    \item[a)] $C$는 이산 위상에 관하여 형식적으로 매끄러운
    $k_0$-대수이다.
    \item[b)] 모든 $\mathfrak p\in\operatorname{Spec}(C)$가
    (i)의 동치인 조건들을 만족한다.
    \item[c)] $C$의 모든 극대 아이디얼 $\mathfrak m$이 (i)의
    동치인 조건들을 만족한다.
    \item[d)] $k_0$의 모든 확대체 $k'$에 대하여
    $C\otimes_{k_0}k'$은 정칙환이다 \sref{0.17.3.6-ko}. 이때 $C$를
    $k_0$ 위에서 기하학적으로 정칙인 환이라고도 한다.
  \end{enumerate}

  \oldpage[0\textsubscript{IV-1}]{208}

  \item[\rm(iii)] $B=k[T_1,\ldots,T_r]$를 다항식 대수,
  $\mathfrak J$를 $C$가 $B/\mathfrak J$와 동형이 되게 하는 $B$의
  아이디얼이라 하자. $\mathfrak q$를 $\mathfrak J$를 포함하는
  $B$의 소아이디얼이라 하고 $\mathfrak p=\mathfrak q/\mathfrak J$라
  하자. 그러면 (i)의 조건들은 다음 조건과도 동치이다
  (“자리스키의 야코비 판정법”).
  \begin{enumerate}
    \item[e)] $B_{\mathfrak q}$에서 아이디얼
    $\mathfrak J_{\mathfrak q}$을 생성하는 $\mathfrak J$의 원소계
    $(f_i)_{1\leq i\leq m}$와 $B$에서 자신으로 가는 $k_0$-미분
    $D_j$ $(1\leq j\leq m)$가 존재하여
    $\det(D_j(f_i))\notin\mathfrak q$이다.
  \end{enumerate}
\end{enumerate}
\end{proposition}

\begin{proof}
$C$는 항상 $B/\mathfrak J$ 꼴의 $k$-대수와 동형이다. $B$는
형식적으로 매끄러운 $k$-대수이고 \sref{0.19.3.3-ko}, $k$는
$k_0$ 위에서 분리적이므로 형식적으로 매끄러운 $k_0$-대수이다
\sref{0.19.6.1-ko}. 따라서 $B$는 형식적으로 매끄러운
$k_0$-대수이다 \sref{0.19.3.5-ko}, (ii). 그러므로 (i)의 조건 a),
b), c)의 동치는 \sref{0.22.6.4-ko}에서 따르고, (ii)의 a), b), c)의
동치는 \sref{0.22.6.6-ko}에서 따른다. (i)의 a)와 d)의 동치는
\sref{0.22.5.8-ko}에서 따른다. $C\otimes_{k_0}k'$의 모든 국소화환은
적당한 소아이디얼 $\mathfrak p\in\operatorname{Spec}(C)$에 대하여
$C_{\mathfrak p}\otimes_{k_0}k'$의 국소화환이기도 하므로, (ii)의
d)와 b)의 동치는 (i)의 d)와 a)의 동치에서 따른다. 끝으로
$\Omega^1_{B/k_0}$는 사영 $B$-가군이므로, 자리스키 판정법 e)와
(i)의 다른 조건들의 동치는 \sref{0.20.4.8-ko}을 고려한
\sref{0.19.1.12-ko}과 \sref{0.22.6.3-ko}에서 따른다.
\end{proof}

실제로 자리스키가 관심을 가진 것은 국소환 $C_{\mathfrak p}$의
정칙성에 대한 미분 판정법이었다. 체를 포함하는 뇌터 국소환이
정칙이라는 것과 그 소체 위의 대수로서 형식적으로 매끄럽다는 것은
동치이므로 \sref{0.19.6.4-ko}, \sref{0.22.6.7-ko}에서 $k_0$를
$k$의 소체로 바꾸면 그러한 판정법을 곧바로 얻는다. 특히 다음 결과를
얻는데, 이는 나중에 나가타의 더 일반적인 결과들의 특수한 경우로서
(IV, 6.12.5)에서 다시 만나게 될 것이다.

\begin{corollary}[22.6.8]
\label{0.22.6.8-ko}
\textnormal{(자리스키.)}
$C$를 체 위의 유한형 대수라 하자. 그러면
$C_{\mathfrak p}$이 정칙 국소환인
$\mathfrak p\in\operatorname{Spec}(C)$들의 집합은
$\operatorname{Spec}(C)$에서 열린집합이다.
\end{corollary}

\begin{remark}[22.6.9]
\label{0.22.6.9-ko}
\begin{enumerate}
  \item[\rm(i)] $k$가 $k_0$ 위에서 분리적이라고 가정하는 대신
  $k$가 \emph{유한 순수 비분리 중복도}를 갖는다고만 가정해도
  \sref{0.22.6.7-ko}의 명제들은 여전히 성립한다. 다시 말해
  \sref{0.19.6.6-ko}, $k_0$의 유한 순수 비분리 확대체 $k'_0$가
  존재하여 국소 아르틴 환 $k\otimes_{k_0}k'_0$의 잉여체가
  $k'_0$의 분리 확대체 $k'$이 된다고 하자. 그러면
  $k'$은 형식적으로 매끄러운 $k'_0$-대수이므로
  \sref{0.19.6.1-ko}, 표준 준동형
  $k\otimes_{k_0}k'_0\longrightarrow k'$과 합성했을 때 항등사상이
  되는 $k'_0$-단사 준동형
  \[
    k'\longrightarrow k\otimes_{k_0}k'_0
  \]
  이 존재한다. 따라서 $C\otimes_{k_0}k'_0$에는 유한형 $k'$-대수
  구조가 주어진다. $k'$은 $k'_0$ 위에서 분리적이므로 이
  $k'$-대수에 \sref{0.22.6.7-ko}을 적용할 수 있고, 이어서
  \sref{0.19.4.7-ko}을 적용하면 주장이 따른다. 이 일반화는 사용할
  일이 없으므로 세부사항은 독자에게 맡긴다. 반면 $k$의 $k_0$ 위의
  확대에 \emph{아무} 가정도 두지 않을 때
  \sref{0.22.6.7-ko}의 결과들이 성립하는지는 알지 못한다.
  \item[\rm(ii)] \sref{0.22.6.7-ko}의 일반 가정 아래,
  $(k_\alpha)$를 $k_0$를 포함하는 $k$의 부분체들의 감소 여과족으로
  두고
  \[
    \bigcap_\alpha k_\alpha(k^p)=k_0(k^p)
  \]
  라 하자. 여기서 $p$는 $k_0$의 표수 지수이다. 나가타의 차원 계산
  방법을 사용하면 야코비 판정법 \sref{0.22.6.7-ko}, (iii)에서
  적당한 $\alpha$에 대한 $B$의 $k_\alpha$-미분 $D_j$들만
  고려해도 됨을 보일 수 있다.

  \oldpage[0\textsubscript{IV-1}]{209}

  이 결과의 의의는 $[k:k_\alpha]$가 유한인 이러한 족
  $(k_\alpha)$가 언제나 존재한다는 데 있다 \sref{0.21.8.6-ko}.
  우리는 사용하지 않을 자리스키 판정법의 이 정밀화를 증명하지
  않는다. (22.7)에서는 완비 국소환에 대하여 본질적으로
  같은 방법으로 증명되는 자리스키 판정법의 변형을 제시할 것이다.
  이 변형도 나가타에 의한 것이며 기술적 어려움은 조금 더 크다.
\end{enumerate}
\end{remark}

\subsection{나가타의 야코비 판정법}
\label{subsection:0.22.7-ko}

\begin{env}[22.7.1]
\label{0.22.7.1-ko}
나가타의 야코비 판정법은 자리스키의 야코비 판정법에 대응하는
판정법이지만, 체 위의 \emph{형식적 멱급수}환의 몫환을 대상으로 한다.
나가타 [31]를 따라 두 가지 형태를 제시하며, 여기서는 형식적
매끄러움의 판정법으로 서술한다.
\end{env}

\begin{proposition}[22.7.2]
\label{0.22.7.2-ko}
$k$를 체, $A=k[[T_1,\ldots,T_r]]$를 통상적인 $k$-대수 구조를 준
$k$ 위의 형식적 멱급수환, $\mathfrak q$를 $A$의 아이디얼,
$B=A/\mathfrak q$라 하자. $\mathfrak p$를 $\mathfrak q$를 포함하는
$A$의 소아이디얼이라 하자. $C=A/\mathfrak p$의 부분 $k$-대수
$C'$이 존재하여, $C'$은 형식적 멱급수 대수
$k[[X_1,\ldots,X_s]]$와 동형이고, $C$는 $C'$ 위에서 유한이며,
$C$의 분수체 $L$은
\[
  L'=k((X_1,\ldots,X_s))
\]
의 분리 확대체라고 가정한다. (이 가정은 $k$의 표수가 $0$이면
언제나 성립한다.) 그러면 다음 조건들은 서로 동치이다.
\begin{enumerate}
  \item[a)] $B_{\mathfrak p}
  =A_{\mathfrak p}/\mathfrak qA_{\mathfrak p}$은
  $\mathfrak p$-준아딕 위상에 관하여 형식적으로 매끄러운
  $k$-대수이다.
  \item[b)] $A$에서 자신으로 가는 $k$-미분
  $D_i$ $(1\leq i\leq m)$와 $\mathfrak q$의 원소
  $f_i$ $(1\leq i\leq m)$가 존재하여, $f_i$들의
  $\mathfrak qA_{\mathfrak p}$ 안의 상들이 $A_{\mathfrak p}$의 이
  아이디얼을 생성하고
  \[
    \det(D_i(f_j))\notin\mathfrak p
  \]
  이다.
  \item[c)] $B_{\mathfrak p}$은 정칙환이다.
\end{enumerate}
\end{proposition}

\begin{proof}
$B_{\mathfrak p}$의 잉여체는 $k(\mathfrak p)$이다.
$k((X_1,\ldots,X_s))$는 $k$ 위에서 분리적이므로
\sref{0.21.9.6.4-ko}, 가정에 의해
$k(\mathfrak p)=L$은 $k$ 위에서 분리적이다. 따라서 이 조건 아래
a)와 c)가 동치임은 이미 알려져 있다 \sref{0.19.6.4-ko}.
한편 $A_{\mathfrak p}$은 정칙환이고
(\sref{0.17.1.4-ko} 및 \sref{0.17.3.2-ko}), 그 잉여체
$k(\mathfrak p)$가 $k$ 위에서 분리적이므로
$A_{\mathfrak p}$은 $\mathfrak p$-준아딕 위상에 관하여 형식적으로
매끄러운 $k$-대수이다 \sref{0.19.6.4-ko}. 그러므로
\sref{0.22.6.2-ko}, (ii)에 의해 조건 a)는 $L$-벡터공간의 준동형
\begin{equation}
\label{0.22.7.2.1-ko}
  j_0:(\mathfrak q/\mathfrak q^2)\otimes_AL
  \longrightarrow\Omega^1_{A_{\mathfrak p}/k}\otimes_{A_{\mathfrak p}}L
  =\widehat\Omega^1_{A/k}\otimes_AL
  \tag{22.7.2.1}
\end{equation}
이 단사라는 것과 동치이다.

한편 합성 준동형
\begin{equation}
\label{0.22.7.2.2-ko}
  j:(\mathfrak q/\mathfrak q^2)\otimes_AL
  \xrightarrow{j_0}\Omega^1_{A/k}\otimes_AL
  \longrightarrow\widehat\Omega^1_{A/k}\otimes_AL
  \tag{22.7.2.2}
\end{equation}
을 생각하고, 조건 b)가 $j$가 단사라는 것과 동치임을 보이자. 실제로
\[
  (\mathfrak q/\mathfrak q^2)\otimes_AL
  =(\mathfrak qA_{\mathfrak p}/\mathfrak q^2A_{\mathfrak p})
  \otimes_{A_{\mathfrak p}}L
\]
이다. $f'_i$가
$\mathfrak qA_{\mathfrak p}/\mathfrak q^2A_{\mathfrak p}$에서
$f_i$의 상이고 $F_i=j(f'_i\otimes1)$이라 하자. 조건 b)는 행렬
\[
  (\langle F_i,D_j\otimes1\rangle)
\]
이 가역이라는 뜻이다. 따라서 $F_i$들은 선형 독립이고, 더욱이
$f'_i\otimes1$들도 선형 독립이다. 후자들은
\[
  (\mathfrak qA_{\mathfrak p}/\mathfrak q^2A_{\mathfrak p})
  \otimes_{A_{\mathfrak p}}L
\]
을 생성하므로 $j$는 단사이다. 역으로 $j$가 단사라고 가정하고,
$f'_i\otimes1$들이

\oldpage[0\textsubscript{IV-1}]{210}

$(\mathfrak qA_{\mathfrak p}/\mathfrak q^2A_{\mathfrak p})
\otimes_{A_{\mathfrak p}}L$의 기저를 이루도록 $f_i$들을 택하자.
그러면 $F_i$들은 $j$의 상의 기저를 이룬다. 그런데
$\widehat\Omega^1_{A/k}$는 계수가 $r$인 자유 $A$-가군이고
\sref{0.21.9.3-ko}, $A$에서 자신으로 가는 $k$-미분들은 그 쌍대가군을
생성한다. 따라서 $F_i$들이 선형 독립이라는 사실로부터 행렬
$(\langle F_i,D_j\otimes1\rangle)$이 가역이 되게 하는, $A$에서
자신으로 가는 $k$-미분 $D_j$들이 존재한다. 이는 조건 b)이다.

이 논의로부터 이미 b)가 a)를 함의한다. 역으로 조건 a)가 $j$의
단사성을 함의함을 보이기 위해 가환 그림
\begin{equation}
\label{0.22.7.2.3-ko}
\xymatrix{
  (\mathfrak q/\mathfrak q^2)\otimes_AL
    \ar[r]^j \ar[d]_\alpha &
  \widehat\Omega^1_{A/k}\otimes_AL \ar@{=}[d] \\
  (\mathfrak p/\mathfrak p^2)\otimes_AL
    \ar[r]_i &
  \widehat\Omega^1_{A/k}\otimes_AL
}
\tag{22.7.2.3}
\end{equation}
을 생각하고 다음 보조정리를 기록하자.
\end{proof}

\begin{lemma}[22.7.2.4]
\label{0.22.7.2.4-ko}
$R$을 정칙 국소환, $\mathfrak m$을 그 극대 아이디얼,
$K=R/\mathfrak m$을 그 잉여체라 하자. $\mathfrak n$이 $R$의
아이디얼이고 $R/\mathfrak n$이 정칙이면, 정준 준동형
\begin{equation}
\label{0.22.7.2.5-ko}
  \alpha:(\mathfrak n/\mathfrak n^2)\otimes_RK
  \longrightarrow(\mathfrak m/\mathfrak m^2)\otimes_RK
  \tag{22.7.2.5}
\end{equation}
는 단사이다.
\end{lemma}

\begin{proof}
실제로 $\mathfrak n$은 $R$의 정칙 매개변수계
$(x_i)_{1\leq i\leq n}$의 일부를 이루는 원소계
$(x_i)_{1\leq i\leq t}$에 의해 생성된다 \sref{0.17.1.6-ko}.
$1\leq i\leq t$인 $x_i$들의 상은
$(\mathfrak n/\mathfrak n^2)\otimes_RK$의 기저를 이루고,
$\alpha$에 의한 그 상들은 $1\leq i\leq n$인 $x_i$들의 상으로
이루어진 $(\mathfrak m/\mathfrak m^2)\otimes_RK$의 기저의 일부를
이룬다. 따라서 보조정리가 따른다.
\end{proof}

\begin{proof}[명제~{\rm(22.7.2)}의 증명 계속]
이 보조정리와 그림 \sref{0.22.7.2.3-ko}에 의해, 가정 a) 아래에서는
$i$가 단사임을 보이면 충분하다. 열
\begin{equation}
\label{0.22.7.2.6-ko}
  \mathfrak p/\mathfrak p^2\xrightarrow{h}
  \widehat\Omega^1_{A/k}\widehat\otimes_A(A/\mathfrak p)
  \xrightarrow{g}\widehat\Omega^1_{(A/\mathfrak p)/k}\longrightarrow0
  \tag{22.7.2.6}
\end{equation}
을 생각하자. \sref{0.20.7.20-ko}에 의해 $g$는 전사이고 $h$의 상은
$\mathfrak m$-아딕 위상에서 $g$의 핵에 조밀하다. 여기서
$\mathfrak m$은 $A$의 극대 아이디얼이다. 한편
$\widehat\Omega^1_{A/k}$가 유한형 $A$-가군이므로
$(A/\mathfrak p)$-가군
\[
  \widehat\Omega^1_{A/k}\otimes_A(A/\mathfrak p)
  =\widehat\Omega^1_{A/k}\widehat\otimes_A(A/\mathfrak p)
\]
의 모든 부분가군은 $\mathfrak m$-아딕 위상에 관하여 닫혀 있다
$(\mathbf 0_{\mathrm I},~7.3.5)$. 따라서
\sref{0.22.7.2.6-ko}은 완전열이다. $L$을 텐서곱하면 완전열
\[
  (\mathfrak p/\mathfrak p^2)\otimes_AL\xrightarrow{i}
  \widehat\Omega^1_{A/k}\otimes_AL\longrightarrow
  \widehat\Omega^1_{(A/\mathfrak p)/k}\otimes_{A/\mathfrak p}L
  \longrightarrow0
\]
을 얻는다.

$L$에 대한 분리성 가정과 \sref{0.21.9.5-ko}에 의해 마지막 가군의
$L$ 위의 계수는 $s$이다. 가운데 가군의 $L$ 위의 계수는
$r$이다 \sref{0.21.9.3-ko}. 따라서
\[
  \operatorname{rg}_L(\operatorname{Im}i)=r-s.
\]
또한 $A/\mathfrak p$은 $k[[T_1,\ldots,T_s]]$와 동형인 부분대수
위에서 유한이므로
\[
  \dim(A/\mathfrak p)=s
\]
이다 (\sref{0.17.1.4-ko} 및 \sref{0.16.1.5-ko}). 그러므로
\[
  r-s=\dim(A)-\dim(A/\mathfrak p)=\dim(A_{\mathfrak p})
\]
이다 \sref{0.16.5.11-ko}.

\oldpage[0\textsubscript{IV-1}]{211}

끝으로 $A_{\mathfrak p}$은 정칙이므로 \sref{0.17.3.2-ko},
$\dim(A_{\mathfrak p})$은
$\mathfrak pA_{\mathfrak p}/\mathfrak p^2A_{\mathfrak p}$의 $L$ 위의
계수와 같고 \sref{0.17.1.1-ko}, 따라서
$(\mathfrak p/\mathfrak p^2)\otimes_AL$의 $L$ 위의 계수와 같다.
그러므로 $i$의 정의역과 상의 계수가 같아 $i$는 단사이다.
\end{proof}

\begin{theorem}[22.7.3]
\label{0.22.7.3-ko}
$k$를 체, $A$를 잉여체가 $K$인 완비 뇌터 국소환이라 하자. 다음을
가정한다.
\begin{enumerate}
  \item[$1^\circ$] $[k:k^p]<+\infty$이다. 여기서 $p$는 $k$의 표수
  지수이다.
  \item[$2^\circ$] $K$는 $k$의 분리 확대체 $K_0$의 유한 확대체이다.
  \item[$3^\circ$] $A$에는 (준아딕 위상에 관하여) 형식적으로
  매끄러운 $K_0$-대수 구조가 주어져 있다.
\end{enumerate}
$\mathfrak q$를 $A$의 아이디얼, $B=A/\mathfrak q$,
$\mathfrak p$를 $\mathfrak q$를 포함하는 $A$의 소아이디얼이라 하자.
다음 조건들은 서로 동치이다.
\begin{enumerate}
  \item[a)] $B_{\mathfrak p}$은 ($\mathfrak p$-준아딕 위상에 관하여)
  형식적으로 매끄러운 $k$-대수이다.
  \item[b)] $A$에서 자신으로 가는 $k$-미분
  $D_i$ $(1\leq i\leq m)$와 $\mathfrak q$의 원소
  $f_i$ $(1\leq i\leq m)$가 존재하여, $f_i$들의
  $\mathfrak qA_{\mathfrak p}$ 안의 상들이 $A_{\mathfrak p}$의 이
  아이디얼을 생성하고 $\det(D_i(f_j))\notin\mathfrak p$이다.
  \item[b')] $K_0^p$를 포함하고 $[K_0:k']<+\infty$인 $K_0$의
  부분확대체 $k'$, $A$에서 자신으로 가는 $k'$-미분
  $D_i$ $(1\leq i\leq m)$, 그리고 $\mathfrak q$의 원소
  $f_i$ $(1\leq i\leq m)$가 존재하여, $f_i$들의
  $\mathfrak qA_{\mathfrak p}$ 안의 상들이 $A_{\mathfrak p}$의 이
  아이디얼을 생성하고 $\det(D_i(f_j))\notin\mathfrak p$이다.
\end{enumerate}
\end{theorem}

\begin{proof}
$p=1$인 경우와 $p>1$인 경우를 나누어 생각하자.

\emph{A)} $k$의 \emph{표수가} $0$이라고 하자. 이때 $K$는 $K_0$의
분리 확대체이므로, \sref{0.19.6.4-ko}에 의해 $A$는 통상적인
$K$-대수 구조를 준 형식적 멱급수환
\[
  K[[T_1,\ldots,T_r]]
\]
과 $K_0$-동형이다. 이제 $A/\mathfrak p$에
\sref{0.19.8.8-ko}, (ii)를 적용하면, \sref{0.22.7.2-ko}의 일반
가정에서 $k$를 $K$로 바꾼 조건들이 성립함을 알 수 있다. 더 나아가
\sref{0.19.6.4-ko}에 의해 $B_{\mathfrak p}$이 형식적으로 매끄러운
$k$-대수라는 것과 형식적으로 매끄러운 $K$-대수라는 것은 서로
동치이며, 두 조건 모두 $B_{\mathfrak p}$이 정칙환이라는 것과
동치이다. 따라서 \sref{0.22.7.2-ko}의 결론을 적용할 수 있고, 이로써
a), b), b')의 동치가 곧바로 증명된다. 여기서 b')에서는
$k'=K_0$로 둔다.

\emph{B)} $k$의 \emph{표수가} $p>0$이라고 하자. $A$는 형식적으로
매끄러운 $K_0$-대수이고 $K_0$는 $k$ 위에서 분리적이므로,
\sref{0.19.3.5-ko}, (ii)와 \sref{0.19.6.1-ko}에 의해 $A$는
형식적으로 매끄러운 $k$-대수이다. \sref{0.22.5.9-ko}에 의해
$A_{\mathfrak p}$도 $\mathfrak p$-준아딕 위상에 관하여 형식적으로
매끄러운 $k$-대수이다. 그러므로 다시 \sref{0.22.6.2-ko}, (ii)에
의해 조건 a)는 \sref{0.22.7.2.1-ko}에서 정의한 준동형 $j_0$가
단사라는 것과 동치이다. 여기서 $L=k(\mathfrak p)$이다. 따라서
\sref{0.19.1.12-ko}, c)는 이미 b)가 a)를 함의함을 보인다. 한편
b')은 자명하게 b)를 함의하므로, $j_0$가 단사라는 가정이 b')을
함의함을 보이면 충분하다.

먼저 $k'$을 $[K_0:k']<+\infty$인 $K_0$의 임의의 부분확대체라 하자.
$\widehat\Omega^1_{A/k'}$은 유한형 자유 $A$-가군이다. 따라서
\sref{0.22.7.2-ko}의 논법에 의해, 적당한 $k'$을 택했을 때 합성
준동형
\begin{equation}
\label{0.22.7.3.1-ko}
  j':(\mathfrak q/\mathfrak q^2)\otimes_AL
  \xrightarrow{j_0}\Omega^1_{A/k}\otimes_AL
  \longrightarrow\widehat\Omega^1_{A/k'}\otimes_AL
  \tag{22.7.3.1}
\end{equation}

\oldpage[0\textsubscript{IV-1}]{212}

이 단사임을 보이면 충분하다. 다시 가환 그림
\begin{equation}
\label{0.22.7.3.2-ko}
\xymatrix{
  (\mathfrak q/\mathfrak q^2)\otimes_AL
    \ar[r]^{j_0} \ar[d]_\alpha &
  \Omega^1_{A/k}\otimes_AL \ar[r] \ar@{=}[d] &
  \widehat\Omega^1_{A/k'}\otimes_AL \ar@{=}[d] \\
  (\mathfrak p/\mathfrak p^2)\otimes_AL
    \ar[r]_{i_0} &
  \Omega^1_{A/k}\otimes_AL \ar[r] &
  \widehat\Omega^1_{A/k'}\otimes_AL
}
\tag{22.7.3.2}
\end{equation}
을 생각하자.

$j_0$가 단사이므로
$\operatorname{Ker}(i_0\circ\alpha)=0$이다. 한편
$A_{\mathfrak p}$은 정칙환이고 가정 a)에 의해
$B_{\mathfrak p}=A_{\mathfrak p}/\mathfrak qA_{\mathfrak p}$도
정칙환이므로 \sref{0.19.6.5-ko}, 보조정리
\sref{0.22.7.2.4-ko}에 의해 $\alpha$는 단사이다. 따라서
$\alpha^{-1}(\operatorname{Ker}i_0)=0$이다. $i'$을
\sref{0.22.7.3.2-ko}의 둘째 행에 있는 준동형들의 합성이라 하면,
마찬가지로
$\operatorname{Ker}(j')=\alpha^{-1}(\operatorname{Ker}(i'))$이다.
그러므로 적당한 $k'$을 택하여
\begin{equation}
\label{0.22.7.3.3-ko}
  \operatorname{Ker}(i')=\operatorname{Ker}(i)
  \tag{22.7.3.3}
\end{equation}
임을 보이면 충분하다.

그런데 \sref{0.22.7.2-ko}에서와 같은 논법으로 완전열
\[
  (\mathfrak p/\mathfrak p^2)\otimes_AL\xrightarrow{i'}
  \widehat\Omega^1_{A/k'}\otimes_AL\longrightarrow
  \widehat\Omega^1_{(A/\mathfrak p)/k'}\otimes_{A/\mathfrak p}L
  \longrightarrow0
\]
을 얻는다. 주어진 가정들 아래에서 \sref{0.21.9.8-ko}을 적용할 수
있으므로, $K_0^p$를 포함하고 $[K_0:k']<+\infty$인 $K_0$의
부분확대체 $k'$을 택하여
\[
  \operatorname{rg}_L\bigl(\widehat\Omega^1_{(A/\mathfrak p)/k'}
  \otimes_{A/\mathfrak p}L\bigr)
  =\dim(A/\mathfrak p)+\operatorname{rg}_L\Upsilon_{L/k}
  +\operatorname{rg}_{K_0}\Omega^1_{K_0/k'}
\]
이 되게 할 수 있다. 한편 \sref{0.16.5.11-ko}에 의해
\[
  \dim(A/\mathfrak p)=\dim(A)-\dim(A_{\mathfrak p}),
\]
$A_{\mathfrak p}$이 정칙이므로 \sref{0.17.1.1-ko}
\[
  \dim(A_{\mathfrak p})
  =\operatorname{rg}_L\bigl((\mathfrak p/\mathfrak p^2)\otimes_AL\bigr),
\]
그리고 \sref{0.21.9.2-ko}에 의해
\[
  \dim(A)+\operatorname{rg}_{K_0}\Omega^1_{K_0/k'}
  =\operatorname{rg}_L\bigl(\widehat\Omega^1_{A/k'}\otimes_AL\bigr)
\]
이다. 따라서
\begin{equation}
\label{0.22.7.3.4-ko}
  \operatorname{rg}_L(\operatorname{Ker}(i'))
  =\operatorname{rg}_L\Upsilon_{L/k}.
  \tag{22.7.3.4}
\end{equation}

다른 한편 $A_{\mathfrak p}$의 잉여체는 $k$의 소체 위에서
형식적으로 매끄러우므로 완전열 \sref{0.20.6.22.1-ko}
\[
  \Upsilon_{L/k}\xrightarrow{\chi_{A_{\mathfrak p}/k}}
  (\mathfrak pA_{\mathfrak p})/(\mathfrak pA_{\mathfrak p})^2
  \xrightarrow{i}\Omega^1_{A_{\mathfrak p}/k}
  \otimes_{A_{\mathfrak p}}L
\]
이 있다. 또한 $A_{\mathfrak p}$은 $\mathfrak p$-준아딕 위상에 관하여
형식적으로 매끄러운 $k$-대수이므로, \sref{0.22.2.2-ko}에 의해
$\chi_{A_{\mathfrak p}/k}$는 단사이다. 따라서
$\operatorname{Ker}(i)$는 $\Upsilon_{L/k}$와 동형이다.
이제 \sref{0.22.7.3.4-ko}와
$\operatorname{Ker}(i)\subset\operatorname{Ker}(i')$를 함께 쓰면
\sref{0.22.7.3.3-ko}이 증명되고, 따라서 정리가 증명된다.
\end{proof}

\oldpage[0\textsubscript{IV-1}]{213}

\begin{remark}[22.7.4]
\label{0.22.7.4-ko}
사실 \sref{0.21.9.6-ko}을 이용하면 다음도 증명된다. $k'$을
$K_0^p$를 포함하는 $K_0$의 부분체들의 감소하는 여과족
$(k_\alpha)$ 가운데 하나로 제한하고, 모든 $\alpha$에 대하여
$[K_0:k_\alpha]<+\infty$이며
$\bigcap_\alpha k_\alpha=K_0^p$라고 하더라도, 조건 b')은
\sref{0.22.7.3-ko}의 다른 조건들과 여전히 동치이다.
\end{remark}

\begin{corollary}[22.7.5]
\label{0.22.7.5-ko}
\sref{0.22.7.3-ko}의 일반 가정 아래에서, $B_{\mathfrak n}$이
($\mathfrak n$-준아딕 위상에 관하여) 형식적으로 매끄러운
$k$-대수가 되는 소아이디얼
$\mathfrak n\in\operatorname{Spec}(B)$들의 집합은
$\operatorname{Spec}(B)$에서 열린집합이다.
\end{corollary}

\begin{proof}
\sref{0.22.7.3-ko}의 표기를 쓰자. $B_{\mathfrak p}$이 $k$ 위에서
형식적으로 매끄러울 때, $\operatorname{Spec}(A)$에서
$\mathfrak p$에 충분히 가깝고 $\mathfrak p'\in V(\mathfrak q)$인
모든 $\mathfrak p'$에 대하여 $B_{\mathfrak p'}$도 그러함을 보이면
충분하다. 그런데 $k$-미분 $D_i$와 $\mathfrak q$의 원소 $f_i$가
\sref{0.22.7.3-ko}의 판정법 b)를 만족하면,
$f=\det(D_i(f_j))$라 둘 때 $\mathfrak p'\in D(f)$에 대하여
$f\notin\mathfrak p'$이다. 한편 $g\notin\mathfrak p$인 원소 $g$가
존재하여, $\mathfrak p'\in D(g)$일 때 $f_i$들의
$\mathfrak qA_{\mathfrak p'}$ 안의 상들이 이 $A_{\mathfrak p'}$의
아이디얼을 생성한다 (Bourbaki, \emph{Alg. comm.}, chap.~II, §~5,
n\textsuperscript{o}~1, prop.~2). 이로써 따름정리가 증명된다.
\end{proof}

\begin{corollary}[22.7.6]
\label{0.22.7.6-ko}
\textnormal{(나가타).} $B$를 체를 포함하는 완비 뇌터 국소환이라 하자.
$B_{\mathfrak n}$이 정칙환이 되는
$\mathfrak n\in\operatorname{Spec}(B)$들의 집합은
$\operatorname{Spec}(B)$에서 열린집합이다.
\end{corollary}

\begin{proof}
실제로 $B$는 어떤 소체 $k$ 위의 대수이므로, \sref{0.19.8.8-ko},
(ii)에 의해 $B$는 $A/\mathfrak q$ 꼴이다. 여기서
$A=K_0[[T_1,\ldots,T_n]]$은 형식적 멱급수환이고 $K_0$는 $k$의
확대체이다. 다른 한편 $B_{\mathfrak n}$이 정칙이라는 것은 이것이
$\mathfrak n$-준아딕 위상에 관하여 형식적으로 매끄러운
$k$-대수라는 것과 동치이다 \sref{0.19.6.4-ko}. 따라서
\sref{0.22.7.5-ko}을 적용하는 데 필요한 모든 조건이 성립한다.
\end{proof}

\begin{remarks}[22.7.7]
\label{0.22.7.7-ko}
\begin{enumerate}
  \item[(i)] 뒤에서 나가타와 함께 \sref{0.22.7.6-ko}을 이용하여,
  임의의 완비 뇌터 국소환으로 \sref{0.22.7.6-ko}의 결론을 확장할
  수 있음을 보일 것이다 (IV, 6.12.7).

  \item[(ii)] $[k:k^p]$가 유한이라고 가정하지 않으면 정리
  \sref{0.22.7.3-ko}의 결론은 더 이상 반드시 성립하지 않는다.
  예를 들어
  
  \[
    k=\mathbf F_p(X_n)_{n\geqslant0},\qquad
    A=k[[T,U]],
  \]
  
  로 두고, $\mathfrak q$를 원소
  \[
    f=U^p-\sum_{n=0}^{\infty}X_nT^{np}
  \]
  가 생성하는 주아이디얼 $Af$로 두자. $A$는 유일분해정역이고,
  그 안에서 $f$는 기약원이다. 실제로 $f$는 다항식환
  $k((T))[U]$에서 기약원이므로 형식적 멱급수환
  $k((T))[[U]]$에서도 기약원이며 (Bourbaki,
  \emph{Alg. comm.}, chap.~VII, §~3), 따라서 더욱 $A$에서도
  기약원이다. $\mathfrak p=\mathfrak q$로 두면
  $B_{\mathfrak p}$은 \sref{0.21.9.7-ko}에서 $E$로 표시한 체이다.
  이 체는 $k$ 위에서 분리적이므로 (같은 곳), 형식적으로 매끄러운
  $k$-대수이다. 그러나 여기서는 \sref{0.22.7.3-ko}의 표기로
  $k=K_0=K$이고, $A$에서 자신으로 가는 $k$-미분들은
  $\partial/\partial T$와 $\partial/\partial U$의 $A$-선형결합이다.
  그런데
  \[
    \partial f/\partial T=\partial f/\partial U=0
  \]
  이므로 \sref{0.22.7.3-ko}의 판정법 b)와 b')은 성립하지 않는다.
\end{enumerate}
\end{remarks}

\section{일본환}
\label{section:0.23-ko}

이 절의 결과들은 IV, 7.6 및 7.7에서 보충될 것이다.

\subsection{일본환}
\label{subsection:0.23.1-ko}

\begin{definition}[23.1.1]
\label{0.23.1.1-ko}
정역 $A$가 다음 조건을 만족할 때 $A$를 \emph{일본환}이라 한다.
$A$의 분수체 $K$의 임의의 유한 확대체 $K'$에 대하여, $K'$ 안에서의
$A$의 정수적 폐포 $A'$이

\oldpage[0\textsubscript{IV-1}]{214}

유한형 $A$-가군이다. 다시 말해 $A'$은 유한 $A$-대수이다. 환 $A$의
모든 유한형 정역 $A$-대수가 일본환일 때 $A$를 \emph{보편 일본환}이라
한다.
\end{definition}

$A$가 일본환이면 모든 분수환 $S^{-1}A$도 일본환임은 명백하다.
실제로 위 표기에서 $S^{-1}A'$은 $K'$ 안에서의 $S^{-1}A$의 정수적
폐포이고, 명백히 유한형 $S^{-1}A$-가군이다.

뇌터 정역은, 심지어 \emph{이산 값매김환}이라 하더라도, 반드시
일본환인 것은 아니다 \cite{IV-30}.

\begin{proposition}[23.1.2]
\label{0.23.1.2-ko}
$A$를 뇌터 정역, $K$를 그 분수체라 하자. $K$의 모든 유한 순수
비분리 확대체 $K'$에 대하여 $K'$ 안에서의 $A$의 정수적 폐포가
유한형 $A$-가군이면, $A$는 일본환이다.
\end{proposition}

\begin{proof}
$A$가 뇌터이므로, $A$가 일본환임을 확인하려면 $K$의 모든 유한
준갈루아 확대체 $L$에 대하여 $L$ 안에서의 $A$의 정수적 폐포 $B$가
유한형 $A$-가군임을 보이면 충분하다. 그런데 $K'$을 $L$에 포함되는
$K$의 최대 순수 비분리 확대체라 하면 $L$은 $K'$의 갈루아
확대체이다. 또한 $A'$을 $K'$ 안에서의 $A$의 정수적 폐포라 하면,
$B$는 $L$ 안에서의 $A'$의 정수적 폐포이다. 한편 $K'$의 분리 확대체
안에서 $A'$의 정수적 폐포는 유한형 $A'$-가군임이 알려져 있다
(Bourbaki, \emph{Alg. comm.}, chap.~V, §~1,
n\textsuperscript{o}~6, prop.~20의 cor.~1). 따라서 명제가 성립한다.
\end{proof}

\sref{0.23.1.2-ko}에 의해, $K$의 표수가 $0$일 때 $A$가 일본환이라는
것은 그 \emph{정수적 폐포} $A'$이 유한형 $A$-가군이라는 뜻이다.

\begin{theorem}[23.1.3]
\label{0.23.1.3-ko}
\textnormal{(테이트).} $A$를 뇌터 정역, $x\neq0$을 $A$의 원소라 하자.
다음을 가정한다.
\begin{enumerate}
  \item[(i)] $A$는 정수적으로 닫혀 있다.
  \item[(ii)] $\mathfrak p=xA$는 소아이디얼이고, $A$는
  $\mathfrak p$-준아딕 위상에 관하여 분리이고 완비이다.
  \item[(iii)] $A/xA$는 일본환이다.
\end{enumerate}
그러면 $A$는 일본환이다.
\end{theorem}

다음 보조정리를 사용할 것이다.

\begin{lemma}[23.1.3.1]
\label{0.23.1.3.1-ko}
$A$를 환, $x$를 $A$에서 영인자가 아닌 원소라 하고
$xA=\mathfrak p$가 소아이디얼이라고 하자. 그러면 모든 정수 $n>0$에
대하여 $A_{\mathfrak p}$에서 $x^nA_{\mathfrak p}$의 $A$ 안의
역상은 $x^nA$이다.
\end{lemma}

\begin{proof}
실제로 $a\in A$, $s\notin\mathfrak p$이고
$A_{\mathfrak p}$에서 $b/1=x^na/s$인 $b\in A$를 택하자. 그러면
$s'\notin\mathfrak p$인 $s'$이 존재하여
$s'sb=x^nas'$이다. 따라서 $s'sb\in\mathfrak p$이고
$s's\notin\mathfrak p$이므로 $b\in\mathfrak p$이다. 즉 어떤
$b'\in A$에 대하여 $b=xb'$이다. $x$가 영인자가 아니므로
$s'sb'=x^{n-1}as'$을 얻으며, 이제 $n$에 관한 귀납법을 적용하면
된다.
\end{proof}

\begin{proof}[정리~{\rm\sref{0.23.1.3-ko}}의 증명]
$A$가 정수적으로 닫혀 있으므로, 정리를 증명할 때 $A$의 분수체
$K$의 표수가 $p>0$인 경우만 생각하면 된다. $K'$을 $K$의 유한 순수
비분리 확대체라 하자. 그러면 어떤 $q=p^e$에 대하여
$K'^q\subset K$이다. $K'$을 더 큰 순수 비분리 확대체로 바꾸면,
$y^q=x$인 $y\in K'$이 존재한다고 가정할 수도 있다. $A'$을 $K'$
안에서의 $A$의 정수적 폐포라 하자. $A$가 정수적으로 닫혀 있으므로
$A'$은
\[
  A'=\{x'\in K'\mid x'^q\in A'\cap K=A\}
\]
이다.

$V=A_{\mathfrak p}$로 두자. $V$의 극대 아이디얼
$\mathfrak m=xV$가 주아이디얼이므로, 뇌터 국소환 $V$는 이산
값매김환이다 \sref{0.17.1.4-ko}. $V$도 정수적으로 닫혀 있으므로,
위와 같은 논법에 의해 $K'$ 안에서의 $V$의 정수적 폐포 $V'$은

\oldpage[0\textsubscript{IV-1}]{215}

$x'^q\in V$인 $x'\in K'$들의 집합이다. 또한 $V'$은 이산
값매김환이고 (Bourbaki, \emph{Alg. comm.}, chap.~VI, §~8,
n\textsuperscript{o}~6, prop.~6 및 chap.~V, §~2,
n\textsuperscript{o}~3, lemme~4), 그 극대 아이디얼 $\mathfrak m'$은
$x'^q\in\mathfrak m$인 $x'\in K'$들의 집합이다. 더 나아가 잉여체
$V'/\mathfrak m'$은 $V/\mathfrak m$의 유한 확대체이다 (같은 곳,
chap.~VI, §~8, n\textsuperscript{o}~1, lemme~2). 모든 정수 $n>0$에
대하여
\begin{equation}
\label{0.23.1.3.2-ko}
  \mathfrak m'^n\cap A'=y^nA'
  \tag{23.1.3.2}
\end{equation}
임을 보이자.

실제로 $y^q=x$이므로 $y\in\mathfrak m'$이고, 따라서
$y^nA'\subset\mathfrak m'^n\cap A'$이다. 역으로
$x'\in\mathfrak m'^n\cap A'$이라 하고 $x'=z'y^n$인 $z'\in K'$을
택하자. 앞에서 본 바와 같이 $x'^q\in A'\cap K=A$이다. 한편 $x'$은
$t'_i\in\mathfrak m'$인 곱 $t'_1\cdots t'_n$들의 합이므로
$t_i'^q\in\mathfrak m$이다. 따라서
$x'^q\in\mathfrak m^n\cap A$이다. 보조정리
\sref{0.23.1.3.1-ko}에 의해
$\mathfrak m^n\cap A=x^nA$이므로
$z'^qx^n\in x^nA$이다. $x\neq0$이므로 $z'^q\in A$이고, 이로써
\sref{0.23.1.3.2-ko}이 증명된다.

둘째로 $A'$ 위의 $xA'$-준아딕 위상이 \emph{분리}임을 보이자. 이
위상은 $A'$의 $yA'$-준아딕 위상과 같으며,
\sref{0.23.1.3.2-ko}에 의해 이 위상은 $V'$의
$\mathfrak m'$-준아딕 위상이 유도하는 위상이다. $V'$은 이산
값매김환이므로 후자의 위상은 분리이다.

이제 $A'/xA'$이 유한형 $A$-가군임을 보이자. $x=y^q$이고
$y^kA'/y^{k+1}A'$은 $A'/yA'$과 동형이므로, $A'/yA'$이 유한형
$A$-가군임을 보이면 충분하다. 그런데
\sref{0.23.1.3.2-ko}에서 $n=1$로 두면, $A'/yA'$은 $V'$의 잉여체
$V'/\mathfrak m'$의 부분환이다. 이 잉여체는 $V$의 잉여체
$V/\mathfrak m$의 유한 확대체이다. 한편 $V/\mathfrak m$은
$A/\mathfrak p$의 분수체이고, $A'$이 $A$ 위에서 정수적이므로
$A'/yA'$은 $A/\mathfrak p$ 위에서 정수적이다. 따라서
$A'/yA'$은 $V'/\mathfrak m'$ 안에서의 $A/\mathfrak p$의 정수적
폐포에 포함된다. 가정에 의해 $A/\mathfrak p$는 뇌터 일본환이므로,
$A'/yA'$은 유한형 $(A/\mathfrak p)$-가군이고, 따라서 더욱 유한형
$A$-가군이다.

이제 $A'$의 $xA'$-준아딕 위상이 분리이므로, $A'$은 이 위상에 관한
완비화 $\widehat{A'}$의 부분환이다. 그런데 $A'/xA'$은 유한형
$A$-가군이고 $A$는 $xA$-준아딕 위상에 관하여 분리이고 완비이므로,
$(\mathbf 0_{\mathrm I},~7.2.9)$에 의해 $\widehat{A'}$은 유한형
$A$-가군이다. $A$가 뇌터이므로 그 부분가군 $A'$도 유한형이다.
\end{proof}

\begin{corollary}[23.1.4]
\label{0.23.1.4-ko}
$A$를 뇌터이고 정수적으로 닫힌 일본환이라 하자. 그러면 모든 형식적
멱급수환 $A[[T_1,\ldots,T_r]]$은 일본환이다.
\end{corollary}

\begin{proof}
뇌터이고 정수적으로 닫힌 임의의 환 $B$에 대하여 형식적 멱급수환
$B[[T]]$도 뇌터이고 정수적으로 닫혀 있음이 알려져 있다
(Bourbaki, \emph{Alg. comm.}, chap.~V, §~1,
n\textsuperscript{o}~4, prop.~14). 따라서 변수의 수 $r$에 관하여
귀납하면, $A[[T]]$이 일본환임을 보이는 경우로 환원된다. 그런데
$x=T$는 \sref{0.23.1.3-ko}의 모든 조건을 만족하므로 결론이 따른다.
\end{proof}

\begin{theorem}[23.1.5]
\label{0.23.1.5-ko}
\textnormal{(나가타).} 모든 완비 뇌터 국소 정역 $A$는 일본환이다.
\end{theorem}

\begin{proof}
\sref{0.19.8.8-ko}, (ii)에 의해, $A$의 부분환 $B$로서 완비 정칙
국소환이고 $A$가 유한 $B$-대수가 되는 것이 존재한다. $B$가
일본환임을 보이면 충분하므로,

\oldpage[0\textsubscript{IV-1}]{216}

$A$가 정칙이라고 더 가정하는 경우로 환원된다. $n=\dim(A)$에 관하여
귀납하자. $n=0$이면 정리는 명백하다. 이제 $n>0$이라 하고,
$\mathfrak m$을 $A$의 극대 아이디얼이라 하며
$x\in\mathfrak m\setminus\mathfrak m^2$을 택하자.
\sref{0.17.1.8-ko}에 의해 $A/xA$는 정칙환이고,
\sref{0.17.1.7-ko} 및 \sref{0.16.3.4-ko}에 의해
\[
  \dim(A/xA)=n-1.
\]
또한 아이디얼 $xA$는 $A$에서 닫혀 있으므로
$(\mathbf 0_{\mathrm I},~7.3.5)$, $A/xA$는 완비이다. 귀납가정에 의해
$A/xA$는 일본환이고, 이제 \sref{0.23.1.3-ko}에 의해 $A$도
일본환이다.
\end{proof}

\begin{corollary}[23.1.6]
\label{0.23.1.6-ko}
$A$를 완비 뇌터 국소 정역, $K$를 그 분수체, $K'$을 $K$의 유한
확대체라 하자. 그러면 $K'$ 안에서의 $A$의 정수적 폐포 $A'$은 완비
국소환이다.
\end{corollary}

\begin{proof}
\sref{0.23.1.5-ko}에 의해 이미 $A'$은 유한 $A$-대수이다. 따라서
$A'$은 완비 뇌터 반국소환이고, 그러므로 국소환들의 \emph{직접곱}이다.
그런데 $A'$은 정역이므로 반드시 국소환이다.
\end{proof}

\begin{proposition}[23.1.7]
\label{0.23.1.7-ko}
$A$를 뇌터 국소 정역, $K$를 그 분수체, $K'$을 $K$의 유한 확대체,
$A'$을 $K'$ 안에서의 $A$의 정수적 폐포라 하자. 완비화
$\widehat A$가 축소환이라고 가정하고, 그 전분수환을 $R$이라 하자.
\begin{enumerate}
  \item[(i)] 환 $R\otimes_KK'$이 축소환이면 $A'$은 유한형
  $A$-가군이다. 특히 $K'$이 $K$의 분리 확대체이면 이 가정은
  성립한다. 실제로 이때 $R$은 $K$의 확대체들인 체들의 직접곱이다
  (Bourbaki, \emph{Alg.}, chap.~VIII, §~7,
  n\textsuperscript{o}~3, th.~1의 cor.~1).
  \item[(ii)] 특히 $R$이 분리 $K$-대수이면 $A$는 일본환이다.
\end{enumerate}
\end{proposition}

\begin{proof}
\emph{(i)} 간단히
\[
  A_1=\widehat A,\qquad A'_1=A'\otimes_AA_1,\qquad
  K'_1=K'\otimes_AA_1
\]
로 두자. $A_1$은 평탄 $A$-가군이므로
$(\mathbf 0_{\mathrm I},~7.3.5)$, $A'_1$은 $K'_1$의 부분환과
동일시되며 명백히 $A_1$ 위에서 정수적이다. 다른 한편
$K_1=K\otimes_AA_1$로 두면
\[
  K'_1=K'\otimes_KK_1
\]
로 쓸 수 있다. $A$가 정역이고 $A_1$이 평탄 $A$-가군이므로 $A$의
모든 영 아닌 원소는 $A_1$-정칙이다
$(\mathbf 0_{\mathrm I},~6.3.4)$. 따라서
$S=A\setminus\{0\}$라 하면
\[
  K_1=S^{-1}A_1
\]
이고, $K_1$은 $A_1$의 전분수환 $R$의 부분환과 동일시된다. 모든
$K$-가군이 평탄하므로 $K'_1$은 $K'\otimes_KR$의 부분환과
동일시된다. 가정에 의해 후자는 축소환이고 유한 $R$-대수이다.

$A_1$의 극소 소아이디얼들을
$\mathfrak q_i$ $(1\leq i\leq n)$라 하고,
$B_i=A_1/\mathfrak q_i$를 정역, $L_i$를 그 분수체라 하자. 그러면
$R$은 $L_i$들의 직접곱이고, $K'\otimes_KR$은
$K'\otimes_KL_i$들의 직접곱이다. 뒤의 각 $K$-대수는 축소환이므로
$L_i$의 유한 확대체들인 체들의 직접곱이다. $B_i$는 완비 국소
정역이므로, 나가타 정리 \sref{0.23.1.5-ko}에 의해
$K'\otimes_KL_i$ 안에서의 $B_i$의 정수적 폐포는 유한형
$B_i$-가군이고, 따라서 더욱 유한형 $A_1$-가군이다. 또한
$K'\otimes_KL_i$의 원소가 $A_1$ 위에서 정수적이면
$\mathfrak q_i$에 의해 소멸하므로 $B_i$ 위에서도 정수적이다.
따라서 $K'\otimes_KR$ 안에서의 $A_1$의 정수적 폐포는 유한형
$A_1$-가군이다. $A'_1$은 이 정수적 폐포에 포함되고 $A_1$은
뇌터이므로, $A'_1$도 유한형 $A_1$-가군이다. 끝으로 $A_1$은 충실
평탄 $A$-가군이므로 $(\mathbf 0_{\mathrm I},~7.3.5)$, $A'$은 유한형
$A$-가군이다 (Bourbaki, \emph{Alg. comm.}, chap.~I, §~3,
n\textsuperscript{o}~6, prop.~11).

\oldpage[0\textsubscript{IV-1}]{217}

\emph{(ii)} 가정에 의해 $L_i$들은 $K$의 분리 확대체이다. 따라서
$K'\otimes_KL_i$는 축소환이다 (Bourbaki, \emph{Alg.}, chap.~VIII,
§~7, n\textsuperscript{o}~3, th.~1). 그러므로 $K$의 임의의 유한
확대체 $K'$에 (i)를 적용할 수 있고, 이로써 주장이 증명된다.
\end{proof}

\subsection{뇌터 국소 정역의 정수적 폐포}
\label{subsection:0.23.2-ko}

\begin{env}[23.2.1]
\label{0.23.2.1-ko}
이하에서는 임의의 환 $A$에 대하여, 간단히 $A$를 그 멱영근기로
나눈 몫을 $A_{\mathrm{red}}$로 나타낸다. 따라서
$X=\operatorname{Spec}(A)$이면
\[
  X_{\mathrm{red}}=\operatorname{Spec}(A_{\mathrm{red}})
\]
이다.

\emph{국소}환 $A$에 대하여 $A_{\mathrm{red}}$가 정역이고
$A_{\mathrm{red}}$의 정수적 폐포가 \emph{국소환}이면, $A$를
\emph{단일분지}라 한다. 이는 (III, 4.3.6)의 정의를 일반화한다.
$A$가 단일분지이고, $A_{\mathrm{red}}$의 정수적 폐포인 국소환의
잉여체가 $A$의 잉여체의 \emph{순수 비분리 확대체}이면, $A$를
\emph{기하학적 단일분지}라 한다. 정수적으로 닫힌 국소환은
기하학적 단일분지임이 명백하다. \sref{0.23.1.6-ko}에 의해 완비
뇌터 국소 정역은 단일분지이다.
\end{env}

\begin{lemma}[23.2.2]
\label{0.23.2.2-ko}
\textnormal{(*)}\footnote{이 번호의 이하에서는 장 IV, §§~2와 5에서
설명한 개념과 결과를 사용한다. 이 번호의 결과는 장 IV, §~6 이전에는
사용되지 않으므로 논리적 순환은 생기지 않는다.}
$A$를 정역, $A'$을 그 정수적 폐포라 하자. 표준 사상
\[
  f:\operatorname{Spec}(A')\longrightarrow\operatorname{Spec}(A)
\]
이 라디시엘 사상일 필요충분조건은 모든
$\mathfrak p\in\operatorname{Spec}(A)$에 대하여
$A_{\mathfrak p}$가 기하학적 단일분지인 것이다. 이 조건이 성립하면
$f$는 위상동형이다.
\end{lemma}

\begin{proof}
실제로 모든 $\mathfrak p\in\operatorname{Spec}(A)$에 대하여
$A_{\mathfrak p}$의 정수적 폐포는 $A'_{\mathfrak p}$이다
(Bourbaki, \emph{Alg. comm.}, chap.~V, §~1,
n\textsuperscript{o}~5, prop.~16). 또한
$\mathfrak pA_{\mathfrak p}$ 위에 놓이는
$A'_{\mathfrak p}$의 모든 소아이디얼은 극대이다
(\emph{loc. cit.}, §~2, n\textsuperscript{o}~1, prop.~1).
따라서 $f$가 단사라는 것은 모든
$\mathfrak p\in\operatorname{Spec}(A)$에 대하여
$A'_{\mathfrak p}$가 \emph{국소환}이라는 것, 곧
$A_{\mathfrak p}$들이 단일분지라는 것과 같다. 이제 모든
$A_{\mathfrak p}$가 기하학적 단일분지라는 것은 (I, 3.5.8)에 의해
$f$가 라디시엘 사상이라는 것과 같다. 이 조건이 성립하면 $f$는
전사이고 닫힌 사상이므로 (II, 6.1.10), 위상동형이다.
\end{proof}

\begin{lemma}[23.2.3]
\label{0.23.2.3-ko}
$A$를 정역, $K$를 그 분수체, $B$를 뇌터환,
$\varphi:A\to B$를 $B$가 평탄 $A$-가군이 되게 하는 환 준동형이라
하자. $K'$을 $K$의 확대체라 하고, $B_1=B_{\mathrm{red}}$로 두며,
$R$을 $B_1$의 전분수환이라 하자. 그러면 $K'$의 모든
$A$-부분대수 $A'$에 대하여 표준 준동형
\begin{equation}
\label{0.23.2.3.1-ko}
  (A'\otimes_AB)_{\mathrm{red}}
  \longrightarrow(K'\otimes_AR)_{\mathrm{red}}
  \tag{23.2.3.1}
\end{equation}
은 단사이다.
\end{lemma}

\begin{proof}
표준 준동형
$h:A'\otimes_AB\to K'\otimes_AR$를 다음 준동형들의 합성으로 볼 수
있다.
\[
  A'\otimes_AB\xrightarrow{u}K'\otimes_AB
  \xrightarrow{v}K'\otimes_AB_1\xrightarrow{w}K'\otimes_AR.
\]
$B$가 평탄 $A$-가군이므로 $u$는 단사이다. 마찬가지로 $K'$은
평탄 $K$-가군이고 $K$는 평탄 $A$-가군이므로 $K'$은 평탄
$A$-가군이다. 또 $B_1$은 축소환이므로 $B_1\to R$은 단사이고,
따라서 $w$도 단사이다. 끝으로 같은 이유로, $\mathfrak N$을

\oldpage[0\textsubscript{IV-1}]{218}

$B$의 멱영근기라 하면 $v$의 핵은 $K'\otimes_A\mathfrak N$이고,
이는 $K'\otimes_AB$의 멱영근기에 포함된다. 그러므로
$x\in A'\otimes_AB$의 상
$h(x)=(w\circ v\circ u)(x)$가 멱영이면 $x$도 멱영이다. 이로써
보조정리가 증명된다.
\end{proof}

\begin{proposition}[23.2.4]
\label{0.23.2.4-ko}
$A$를 정역, $K$를 그 분수체, $B$를 뇌터환,
$\mathfrak q_i$ $(1\leq i\leq n)$를 $B$의 극소 소아이디얼,
$\varphi:A\to B$를 환 준동형이라 하자. 각
$B/\mathfrak q_i$가 일본환이고 $B$가 평탄 $A$-가군이라고
가정한다. $K'$을 $K$의 유한 확대체라 하고, $(C_\lambda)$를
$K'$을 분수체로 가지는 유한 $A$-대수인 $K'$의 부분환들의 증가
여과족이라 하자. 그러면 $C_\lambda$들의 합집합 $A'$은 $K'$ 안에서의
$A$의 정수적 폐포이다. 이때 다음이 성립한다.
\begin{enumerate}
  \item[(i)] 어떤 지표 $\alpha$가 존재하여, $\lambda\geq\alpha$이면
  표준 준동형
  \begin{equation}
  \label{0.23.2.4.1-ko}
    (C_\alpha\otimes_AB)_{\mathrm{red}}
    \longrightarrow(C_\lambda\otimes_AB)_{\mathrm{red}}
    \tag{23.2.4.1}
  \end{equation}
  은 전단사이다.
  \item[(ii)] 더 나아가 $B$가 충실 평탄 $A$-가군이면 사상
  \[
    \operatorname{Spec}(A')\longrightarrow
    \operatorname{Spec}(C_\alpha)
  \]
  는 라디시엘 사상이다.
\end{enumerate}
\end{proposition}

\begin{proof}
\emph{(i)} $B_1=B_{\mathrm{red}}$로 두고, $R$을 $B_1$의 전분수환이라
하자. $A$가 정역이고 $B$가 평탄 $A$-가군이므로, $A$의 모든 영
아닌 원소는 $B$-정칙이다 $(\mathbf 0_{\mathrm I},~6.3.4)$.
따라서 합성 준동형 $A\xrightarrow{\varphi}B\to B_1$은
$K\to R$인 준동형으로 연장된다. 그러므로
\[
  K'\otimes_AR=(K'\otimes_AK)\otimes_KR
\]
로 쓸 수 있고, $K'\otimes_AK$는 $K'$과 동일시되므로
\[
  K'\otimes_AR=K'\otimes_KR.
\]
$R$은 $B/\mathfrak q_i$의 분수체 $L_i$들의 직접곱이므로
$K'\otimes_AR$은 $K'\otimes_KL_i$들의 직접곱이고, 따라서
$(K'\otimes_AR)_{\mathrm{red}}$는 $L_i$들의 유한 확대체 유한 개의
직접곱이다. 가정에 의해 $L_i$의 유한 확대체 안에서의
$B/\mathfrak q_i$의 정수적 폐포는 유한형
$(B/\mathfrak q_i)$-가군이고, 따라서 유한형 $B$-가군이다.
그러므로 $(K'\otimes_AR)_{\mathrm{red}}$ 안에서의 $B$의 정수적
폐포 $B'$은 유한형 $B$-가군이다.

\sref{0.23.2.3-ko}에 의해
$(C_\lambda\otimes_AB)_{\mathrm{red}}$들은
$(K'\otimes_AR)_{\mathrm{red}}$의 부분환들과 표준적으로 동일시된다.
이들은 유한 $B$-대수이므로 $B'$에 포함된다. $B$가 뇌터환이고
$B'$이 유한형 $B$-가군이므로,
$(C_\lambda\otimes_AB)_{\mathrm{red}}$들의 여과족은
\emph{최대 원소} $(C_\alpha\otimes_AB)_{\mathrm{red}}$를 갖는다.
이로써 (i)이 증명된다.

\emph{(ii)} $B$가 충실 평탄 $A$-가군이므로 (IV, 2.6.1, (v)),
사상
\[
  \operatorname{Spec}(A'\otimes_AB)\longrightarrow
  \operatorname{Spec}(C_\alpha\otimes_AB)
\]
이 라디시엘 사상임을 보이면 충분하다. 이는 (I, 5.1.6)에 의해 사상
\[
  \operatorname{Spec}\bigl((A'\otimes_AB)_{\mathrm{red}}\bigr)
  \longrightarrow
  \operatorname{Spec}\bigl((C_\alpha\otimes_AB)_{\mathrm{red}}\bigr)
\]
이 라디시엘 사상이라는 것과 같다. 그런데
\[
  A'\otimes_AB=\varinjlim_\lambda(C_\lambda\otimes_AB)
\]
이므로 (IV, 5.13.2),
\[
  (A'\otimes_AB)_{\mathrm{red}}
  =\varinjlim_\lambda(C_\lambda\otimes_AB)_{\mathrm{red}}.
\]
따라서 결론은 (i)에서 따른다.
\end{proof}

\begin{corollary}[23.2.5]
\label{0.23.2.5-ko}
$A$를 뇌터 국소 정역, $K$를 그 분수체, $K'$을 $K$의 유한
확대체라 하자. $(C_\lambda)$를 $K'$을 분수체로 가지는 유한
$A$-대수인 $K'$의 부분환들의 증가 여과족이라 하자. 이 환들은 뇌터
반국소환이다 (Bourbaki, \emph{Alg. comm.}, chap.~IV, §~2,
n\textsuperscript{o}~5, prop.~9의 cor.~3). 그러면 어떤 $\alpha$가
존재하여 $\lambda\geq\alpha$이면 준동형
\[
  (\widehat C_\alpha)_{\mathrm{red}}
  \longrightarrow(\widehat C_\lambda)_{\mathrm{red}}
\]
은 동형이다. 또한 $A'$을 $K'$ 안에서의 $A$의 정수적 폐포라 하면,
사상
\[
  \operatorname{Spec}(A')\longrightarrow
  \operatorname{Spec}(C_\alpha)
\]
은 라디시엘 사상이다 (cf.~\sref{0.23.2.2-ko}).
\end{corollary}

\oldpage[0\textsubscript{IV-1}]{219}

\begin{proof}
\sref{0.23.2.4-ko}에서 $B=\widehat A$로 두어 적용한다.
$B/\mathfrak q_i$들은 완비 뇌터 국소환이므로 일본환이고
\sref{0.23.1.5-ko}, $B$는 충실 평탄 $A$-가군이다
$(\mathbf 0_{\mathrm I},~7.3.5)$. 또한
\[
  C_\lambda\otimes_AB=\widehat C_\lambda
\]
이다 (Bourbaki, \emph{Alg. comm.}, chap.~III, §~3,
n\textsuperscript{o}~4, th.~3 및 chap.~IV, §~2,
n\textsuperscript{o}~5, prop.~9의 cor.~3).
\end{proof}

\begin{corollary}[23.2.6]
\label{0.23.2.6-ko}
\sref{0.23.2.5-ko}의 가정 아래에서, $K'$ 안에서의 $A$의 정수적 폐포
$A'$은 반국소환이다. $\mathfrak m$을 $A$의 극대 아이디얼이라 하면,
$A'$의 모든 극대 아이디얼 $\mathfrak m'$에 대하여
$A'/\mathfrak m'$은 $A/\mathfrak m$의 유한 확대체이다.
\end{corollary}

\begin{proof}
$\lambda\geq\alpha$에 대하여 준동형
\[
  (\widehat C_\alpha)_{\mathrm{red}}
  \longrightarrow(\widehat C_\lambda)_{\mathrm{red}}
\]
이 전단사라는 사실에서, $\lambda\geq\alpha$일 때
$C_\lambda$의 극대 아이디얼 개수는 일정함이 따른다. 또한
$\mathfrak m_\alpha$가 $C_\alpha$의 극대 아이디얼이고
$\mathfrak m_\lambda$가 $\mathfrak m_\alpha$ 위에 놓이는
$C_\lambda$의 유일한 극대 아이디얼이면, 체
$C_\alpha/\mathfrak m_\alpha$와
$C_\lambda/\mathfrak m_\lambda$는 표준적으로 동형이다. 이제
\[
  \mathfrak m'=\varinjlim\mathfrak m_\lambda,\qquad
  A'/\mathfrak m'
  =\varinjlim(C_\lambda/\mathfrak m_\lambda)
\]
이므로 $(\mathbf 0_{\mathrm{III}},~10.3.1.1.3)$ 결론이 따른다.
\end{proof}

\begin{proposition}[23.2.7]
\label{0.23.2.7-ko}
\textnormal{(Y.~Mori).}
$A$를 뇌터 국소 정역, $K$를 그 분수체, $A'$을 $A$의 정수적
폐포라 하자. 그러면 $A'$은 반국소 크룰환이다
(Bourbaki, \emph{Alg. comm.}, chap.~VII, §~1). 다시 말해, $K$를
분수체로 가지는 이산 값매김환들의 족 $(V_\lambda)$가 존재하여
\[
  A'=\bigcap_\lambda V_\lambda
\]
이고, 모든 $x\in K$는 유한 개를 제외한 모든 $V_\lambda$에 속한다.
더 나아가, $K$의 부분환 $C$로서 유한 $A$-대수인 것이 존재하여
$V_\lambda$들은 환 $C_{\mathfrak p_\lambda}$의 정수적 폐포가 된다.
여기서 $(\mathfrak p_\lambda)$는 국소환 $C$의 높이 $1$인
소아이디얼들의 족이다.
\end{proposition}

먼저 다음 보조정리를 증명한다.

\begin{lemma}[23.2.7.1]
\label{0.23.2.7.1-ko}
$A$, $B$를 두 환, $\varphi:A\to B$를 $B$가 충실 평탄
$A$-가군이 되게 하는 준동형이라 하자. $A$가 정역이라고 가정한다.
그러면 $C=B_{\mathrm{red}}$라 할 때 합성 준동형
\[
  \psi:A\longrightarrow B\longrightarrow B_{\mathrm{red}}=C
\]
는 단사이고, $A$의 분수체 $K$에서 $C$의 전분수환 $R$로 가는
단사 준동형으로 연장된다. 또한 $C'$을 $R$ 안에서의 $C$의 정수적
폐포라 하면, $C'\cap K$는 $A$의 정수적 폐포이다.
\end{lemma}

\begin{proof}
$A$가 정역이고 $\varphi$가 단사이므로, $\varphi(A)$와 $B$의
멱영근기 $\mathfrak N$의 교집합은 $0$이다. $C=B/\mathfrak N$이므로
$\psi$는 단사이다. $A$의 모든 $a\neq0$은 $A$에서 영인자가 아니고,
평탄성에 의해 $B$에서도 영인자가 아니다
$(\mathbf 0_{\mathrm I},~6.3.4)$. 따라서 $a$는 $C$에서도 영인자가
아니다. 실제로 $B$의 어떤 $x\notin\mathfrak N$에 대하여
$ax\in\mathfrak N$이라면, 어떤 정수 $n$에 대하여
$a^nx^n=0$이 되어 $x^n\neq0$이라는 사실과 모순된다. 그러므로
$\psi$는 $K$에서 $R$로 가는 단사 준동형으로 연장된다.

마지막 주장을 증명하자. $A$의 정수적 폐포 $A'$이
$C'\cap K$에 포함됨은 명백하다. 역으로 $x\in C'\cap K$라 하자.
그러면 $x$는 $C$ 위에서 정수적이고, 따라서 더욱 $B$ 위에서도
정수적이다. 즉 $B[x]$는 유한 $B$-대수이다. 한편 $K$는
$B\otimes_AK$의 부분환과 동일시되고, 평탄성에 의해
$B\otimes_AK$의 부분환 $B[x]$은 $B\otimes_AA[x]$와 동일시된다.
따라서 $A[x]$은 유한형 $A$-가군이다
(Bourbaki, \emph{Alg. comm.}, chap.~I\textsuperscript{er}, §~3,
n\textsuperscript{o}~6, prop.~11). 그러므로 $x\in A'$이다.
\end{proof}

\begin{proof}[명제~\sref{0.23.2.7-ko}의 증명]
이 보조정리를 적용하여, $A$에서 그 완비화 $\widehat A$로 가는 표준
단사 준동형을 생각하자. $\widehat A$는 충실 평탄 $A$-가군이다.
$\mathfrak N$을 $\widehat A$의 멱영근기라 하면
\[
  B=(\widehat A)_{\mathrm{red}}=\widehat A/\mathfrak N
\]
은 축소 완비 뇌터 국소환이다. 따라서 그 전분수환 $R$은 유한 개의
체 $L_i$의 직접곱이다.

\oldpage[0\textsubscript{IV-1}]{220}

$B$의 $L_i$ 안에서의 표준상 $B_i$는 $L_i$를 분수체로 가지는
완비 뇌터 국소 정역이고, $R$ 안에서의 $B$의 정수적 폐포 $B'$은
$B_i$의 정수적 폐포 $B'_i$들의 \emph{직접곱}이다. 그런데
\sref{0.23.1.5-ko}에 의해 $B'_i$는 유한 $B_i$-대수이므로 정수적으로
닫힌 뇌터 국소환이다. 또한 Bourbaki, \emph{Alg. comm.}, chap.~VII,
§~1, n\textsuperscript{o}~3, th.~2의 cor.에 의해 $B'_i$는
크룰환이다. 각 $i$에 대하여 단사 준동형
\[
  \varphi_i:K\longrightarrow L_i
\]
가 존재하므로, $\varphi_i^{-1}(B'_i)$는 $K$ 안의 크룰환이다. 한편
보조정리에 의해 $A'=B'\cap K$이고
\[
  A'=\bigcap_i\varphi_i^{-1}(B'_i)
\]
이므로, $A'$은 유한 개의 크룰환들의 교집합이고 따라서 크룰환이다.

명제의 마지막 주장을 증명하자. 유한 $A$-대수 $C\subset K$가
존재하여 사상
\[
  \operatorname{Spec}(A')\longrightarrow\operatorname{Spec}(C)
\]
이 라디시엘 사상이면서 위상동형이다 (\sref{0.23.2.4-ko}). 따라서
모든 소아이디얼 $\mathfrak p'\in\operatorname{Spec}(A')$에 대하여
$\mathfrak p'$은 $\mathfrak p=\mathfrak p'\cap C$ 위에 놓이는
$A'$의 유일한 소아이디얼이고, $A'_{\mathfrak p'}$은
$C_{\mathfrak p}$의 정수적 폐포이다. 다른 한편 위 사상이
위상동형이라는 사실에서
\[
  \mathfrak p'\longmapsto\mathfrak p'\cap C
\]
가 $A'$의 높이 $1$인 소아이디얼들의 집합에서 $C$의 높이 $1$인
소아이디얼들의 집합으로 가는 전단사임이 따른다. 이제 Bourbaki,
\emph{Alg. comm.}, chap.~VII, §~1, n\textsuperscript{o}~6,
th.~4를 적용하면 결론을 얻는다.
\end{proof}

\begin{remarks}[23.2.8]
\label{0.23.2.8-ko}
\begin{enumerate}
  \item[(i)] \sref{0.23.2.6-ko}과 \sref{0.23.2.7-ko}을 적용할 때,
  $A'$이 \emph{반드시 뇌터환인 것은 아니다}라는 점에 유의해야 한다.
  $\dim(A)=3$인 나가타의 예가 이를 보여 준다 [30].

  \item[(ii)] $A'$이 $K$의 유한 확대체 $K'$ 안에서의 $A$의 정수적
  폐포인 경우에도 \sref{0.23.2.7-ko}의 결론들은 성립한다. 실제로
  $K'$을 분수체로 가지는 유한 $A$-대수 $B$를 택하면 충분하다.
  $B$는 뇌터 반국소환이므로 유한 개의 뇌터 국소환
  $B_{\mathfrak m_i}$의 교집합이다. 여기서 $\mathfrak m_i$들은
  $B$의 극대 아이디얼들이다. $B$의 정수적 폐포는 $A'$과 같고,
  \[
    A'=\bigcap_i A'_{\mathfrak m_i}
  \]
  이다 (Bourbaki, \emph{Alg. comm.}, chap.~II, §~3,
  n\textsuperscript{o}~3, th.~1의 cor.~4). 각
  $A'_{\mathfrak m_i}$는 $B_{\mathfrak m_i}$의 정수적 폐포이므로
  \sref{0.23.2.7-ko}에 의해 크룰환이다. 따라서 $A'$도 크룰환이다.

  \item[(iii)] 모든 뇌터 정역 $A$에 대하여, $A$가 국소환이라는
  가정 없이도 $A$의 정수적 폐포가 크룰환임을 보일 수 있다
  ([30], 3.3.10).
\end{enumerate}
\end{remarks}

\begin{corollary}[23.2.9]
\label{0.23.2.9-ko}
$A$를 뇌터 정역, $A'$을 그 정수적 폐포라 하자. $A\subset C\subset A'$이고
$C$가 유한 $A$-대수인 모든 환 $C$와, $C$의 높이 $1$인 모든
소아이디얼 $\mathfrak q$에 대하여 $\mathfrak q\cap A$가 $A$의
높이 $1$인 소아이디얼이라고 가정한다. $P$를 $A$의 높이 $1$인
소아이디얼들의 집합이라 하면
\[
  A'=\bigcap_{\mathfrak p\in P}A'_{\mathfrak p}.
\]
\end{corollary}

\begin{proof}
$A'$은 꼬임 없는 $A$-가군이므로
\[
  A'=\bigcap_{\mathfrak m}A'_{\mathfrak m},
\]
이다. 여기서 $\mathfrak m$은 $A$의 극대 아이디얼 전체를 달린다
(Bourbaki, \emph{Alg. comm.}, chap.~II, §~3,
n\textsuperscript{o}~3, th.~1의 cor.~4).

\oldpage[0\textsubscript{IV-1}]{221}

따라서 $A'_{\mathfrak m}$이 $\mathfrak p\subset\mathfrak m$인 높이
$1$의 소아이디얼 $\mathfrak p$에 대한 $A'_{\mathfrak p}$들의
교집합임을 보이면 충분하다. $A'_{\mathfrak m}$은
$A_{\mathfrak m}$의 정수적 폐포이므로, 따져야 할 명제를
$A_{\mathfrak m}$에 대해서만 증명하면 된다. 그런데
\sref{0.23.2.7-ko}에 의해 $A'_{\mathfrak m}$에 포함되는 유한
$A_{\mathfrak m}$-대수 $B$가 존재하여, $A'_{\mathfrak m}$은 $B$의
높이 $1$인 소아이디얼 $\mathfrak q$ 전체에 대한
$B_{\mathfrak q}$의 정수적 폐포들의 교집합이 된다. 한편 $B$는 어떤
유한 $A$-대수 $C$에 대하여 $C_{\mathfrak m}$ 꼴이다. 따라서 $B$의
높이 $1$인 모든 소아이디얼 $\mathfrak q$에 대하여, 가정에 의해
$\mathfrak q\cap A$는 $A$의 높이 $1$인 소아이디얼이다.
$B_{\mathfrak q}$의 정수적 폐포는
$A_{\mathfrak q\cap A}$의 정수적 폐포를 포함하고, 뒤의 환은
$A'_{\mathfrak m}$을 포함한다. 그러므로 $A'_{\mathfrak m}$은
$\mathfrak p\in P$이고 $\mathfrak p\subset\mathfrak m$인
$\mathfrak p$에 대한 $A'_{\mathfrak p}$들의 교집합이다. 이로써
증명이 끝난다.
\end{proof}

\begin{remarks}[23.2.10]
\label{0.23.2.10-ko}
\begin{enumerate}
  \item[(i)] $A$가 \emph{보편 카테너리 환}이면
  \sref{0.23.2.9-ko}의 가정이 성립함을 뒤에서 볼 것이다
  (IV, 5.6.3 및 5.6.10). 특히 $A$가 \emph{정칙환의 몫}이면 그
  가정이 성립한다 (IV, 5.10.17 및 5.11.2). 끝으로 사상
  \[
    \operatorname{Spec}(A')\longrightarrow\operatorname{Spec}(A)
  \]
  이 위상동형이면 가정은 명백히 성립하며, 특히 이 사상이 라디시엘
  사상이면 그러하다 (\sref{0.23.2.2-ko}).

  \item[(ii)] $A'$의 유한 부분 $A$-대수들에 관한 가정을 생략하면
  \sref{0.23.2.9-ko}의 결론은 더 이상 반드시 성립하지 않는다.
  실제로 다음 성질을 갖는 차원 $2$의 뇌터 국소 정역 $A$를 정의할
  수 있다. 그 정수적 폐포 $A'$은 유한 $A$-대수이지만, $A$의 높이
  $2$인 극대 아이디얼 $\mathfrak m$ 위에 놓이는 $A'$의 높이 $1$인
  소아이디얼 $\mathfrak p'$이 존재한다. 또한 $A$의 높이 $1$인 모든
  소아이디얼 $\mathfrak p$에 대하여 $A_{\mathfrak p}$는 정수적으로
  닫혀 있다 (IV, 5.6.11). 이때 이 환들 $A_{\mathfrak p}$의
  교집합은 $A'$과 같을 수 없다. 실제로 $A'$은 크룰환이다
  (Bourbaki, \emph{Alg. comm.}, chap.~VII, §~1,
  n\textsuperscript{o}~5, prop.~9).
\end{enumerate}
\end{remarks}

\begin{flushright}
\emph{(계속.)}
\end{flushright}

\cleardoublepage
\markboth{제IV장}{제IV장}
\oldpage[IV-1]{222}
\thispagestyle{plain}

\begin{center}
{\Large 제IV장}

\vspace{1.5cm}
{\LARGE\bfseries 스킴과 스킴 사상의 국소적 연구}

\vspace{1.5cm}
{\large\bfseries 목차}
\end{center}

\medskip
\begin{tabular}{@{}r p{0.84\textwidth}}
§ 1. & 상대적 유한성 조건. 준스킴의 구성 가능 집합.\\
§ 2. & 밑변경과 평탄성.\\
§ 3. & 연관 소사이클과 일차 분해.\\
§ 4. & 기초체의 변경.\\
§ 5. & 준스킴의 차원과 깊이.\\
§ 6. & 국소 뇌터 준스킴의 평탄 사상.\\
§ 7. & 뇌터 국소환과 그 완비화 사이의 관계에 대한 응용. 엑설런트 환.\\
§ 8. & 준스킴의 사영 극한.\\
§ 9. & 구성 가능 성질.\\
§ 10. & 야콥슨 준스킴.\\
§ 11.\footnotemark[1] & 유한 표시 평탄 사상의 위상적 성질. 평탄성의 국소 판정법.\\
§ 12. & 유한 표시 평탄 사상의 올에 대한 연구.\\
§ 13. & 등차원 사상.\\
§ 14. & 보편 열린 사상.\\
§ 15. & 보편 열린 사상의 올에 대한 연구.\\
§ 16. & 미분 불변량. 미분적으로 매끄러운 사상.\\
§ 17. & 매끄러운 사상, 비분기 사상, 에탈 사상.\\
§ 18. & 에탈 사상에 관한 보충. 헨젤 국소환.\\
§ 19. & 정칙 몰입과 횡단 정칙 몰입.\\
§ 20. & 초평면 단면; 일반 사영.\\
§ 21. & 무한소 확장.\\
\end{tabular}
\footnotetext[1]{§§~11부터 21까지의 순서와 내용은 참고용으로만
제시한 것이며, 출판 전에 변경될 수 있다.}

\clearpage
\oldpage[IV-1]{223}

이 장에서 다루는 주제에 관하여 다음과 같은 점들을 말해 두자.

\medskip
\noindent
\emph{a)} 이 주제들의 공통점은 준스킴 또는 사상의 \emph{국소적}
성질, 즉 한 점에서, 한 올의 점들에서, 또는 한 점이나 한 올의
(지정되지 않은) 근방에서 고찰하는 성질을 다룬다는 데 있다. 이러한
성질은 일반적으로 \emph{위상적}, \emph{미분적} 또는
\emph{차원적} 성격을 지니며, 마지막 것은 \emph{차원}과
\emph{깊이}의 개념을 포함한다. 이 성질들은 해당 점에서의
\emph{국소환}의 성질과 연결된다. 전형적인 문제 하나는 주어진 사상
$f:X\to Y$와 점 $x\in X$에 대하여, $x$에서의 $X$의 성질을
$y=f(x)$에서의 $Y$의 성질 및 $x$에서의 올
$X_y=f^{-1}(y)$의 성질과 관련짓는 것이다. 또 다른 문제는 $X$가
어떤 성질을 갖거나 $x$를 지나는 올 $X_{f(x)}$가 $x$에서 어떤
성질을 갖는 점 $x\in X$ 전체의 위상적 성격, 예를 들어 구성
가능성이나 열리거나 닫혀 있다는 성질을 판정하는 것이다. 마찬가지로
$X$가 올 $X_y$의 모든 점에서 어떤 성질을 갖거나 그 올 자체가 어떤
성질을 갖는 점 $y\in Y$ 전체의 위상적 성격도 연구한다.

\medskip
\noindent
\emph{b)} 뒤의 장들에서 가장 중요한 개념은 \emph{유한 표시 평탄
사상}, 그리고 그 특수한 경우인 \emph{매끄러운 사상}과
\emph{에탈 사상}이다. 이 개념들과 관련 문제들에 대한 상세한 연구는
본격적으로 §~11에서 시작한다.

\medskip
\noindent
\emph{c)} §§~1부터 10까지는 예비적인 부분으로 볼 수 있으며, 이
장의 나머지 절들뿐 아니라 뒤의 장들에서도 사용하는 세 종류의
기법을 전개한다.

\medskip
\noindent
\emph{c 1)} §§~1부터 4까지에서는 특히 \emph{유한성}이나
\emph{평탄성} 조건과 관련하여 \emph{밑변경} 개념의 여러 측면을
고찰한다. 여기서는 \emph{강하} 기법의 가장 기초적인 측면을
시작한다. 이 기법에 관련된 ``유효성'' 문제는 제V장에서 연구할
것이다.

\medskip
\noindent
\emph{c 2)} §§~5부터 7까지는 \emph{뇌터 기법}이라 부를 수 있는
것을 중심으로 한다. 여기서 고찰하는 준스킴은 언제나 국소 뇌터라고
가정하지만, 반대로 \emph{사상}에는 일반적으로 아무 유한성 조건도
부과하지 않는다. 이는 본질적으로 차원과 깊이의 개념이 뇌터
국소환에 대해서만 다루기 쉽기 때문이다. §~7은 뇌터 국소환의
``정밀한'' 이론을 이루며, 이 장의 뒤쪽에서는 비교적 적게 쓰인다는
점을 상기하자.

\medskip
\noindent
\emph{c 3)} §§~8부터 10까지에서는 특히, 연구하는 준스킴에 대한
\emph{뇌터 가정을 제거}하고 대신 고찰하는 \emph{사상}에 적절한
\emph{유한성} 가정, 곧 ``유한 표시''를 두는 방법을 제시한다. 이
대체의 장점은 뒤의 가정들이 밑변경에 대하여 \emph{안정적}인 반면,
준스킴에 대한 뇌터 가정은 그렇지 않다는 데 있다. 이 이행을 가능하게
하는 기법은 한편으로 준스킴의 \emph{사영극한} 개념을 사용하여 한
문제를 뇌터 가정 아래의 같은 문제로 환원하는 데 의존하고, 다른
한편으로 \emph{구성 가능 집합}을 체계적으로 사용하는 데 의존한다.
구성 가능 집합은 임의의 사상에 의한 역상과 유한 표시 사상에 의한
직상 아래에서 보존된다는 점, 그리고 국소 뇌터 준스킴에서 다루기
쉬운 위상적 성질을 갖는다는 점에서 이중의 이점이 있다.

\oldpage[IV-1]{224}

같은 기법으로 흔히 더 특수한 뇌터환, 예를 들어 유한형
$\mathbf Z$-대수로까지 환원할 수 있다. 바로 여기서 §~7에서 연구할
``엑설런트'' 환의 성질이 결정적인 구실을 한다. 뇌터 가정을 제거하는
문제와 별개로도, §§~8부터 10까지의 매우 기초적인 기법은 거의 모든
응용에서 끊임없이 사용된다.

\setcounter{section}{0}

\section{상대적 유한성 조건. 준스킴의 구성 가능 집합}
\label{section:IV.1-fr}

이 절에서는 (I, 6.3 및 6.6)에서 제시한 준스킴의 사상
$f:X\to Y$에 대한 ``유한성 조건''을 보충하여 다시 설명한다. $X$에
대하여 \emph{전역적} 성격을 갖는 ``유한성'' 개념은 본질적으로 둘,
곧 \emph{준콤팩트 사상} (I, 6.6.1에서 정의)과 \emph{준분리 사상}의
개념이다. 다른 한편 $X$에 대하여 \emph{국소적} 성격을 갖는
``유한성'' 개념도 둘, 곧 \emph{국소 유한형 사상} (I, 6.6.2에서
정의)과 \emph{국소 유한 표시 사상}의 개념이다. 이 국소 개념들을
앞의 전역 개념들과 결합하면 \emph{유한형 사상} (I, 6.3.1에서
정의)과 \emph{유한 표시 사상}의 개념을 얻는다. 독자의 편의를 위해
(I, 6.3 및 6.6)에서 진술한 주요 성질들을 이 절에서 다시 제시하되,
그 증명은 물론 제I장의 해당 번호들을 참조한다.

n\textsuperscript{o}~(1.8)과 (1.9)에서는 앞의 유한성 개념을
사용하여 $(0\textsubscript{III},~§~9)$의 구성 가능 집합에 관한
결과를 준스킴의 틀 안에서 보충한다.

\subsection{준콤팩트 사상}
\label{subsection:IV.1.1-fr}

\begin{definition}[1.1.1]
\label{IV.1.1.1-fr}
준스킴의 사상 $f:X\to Y$에 대하여, 위상공간 $X$에서 위상공간
$Y$로 가는 연속사상 $f$가 준콤팩트이면 $f$를 \emph{준콤팩트}라고
한다 $(0\textsubscript{III},~9.1.1)$. 다시 말해, $Y$의 모든
준콤팩트 열린집합 $U$의 역상 $f^{-1}(U)$가 준콤팩트인 경우이다
(cf.~(I, 6.6.1)).
\end{definition}

$\mathfrak B$가 아핀 열린집합들로 이루어진 $Y$의 위상의 기저라
하자. $f$가 준콤팩트일 필요충분조건은 모든 $V\in\mathfrak B$에
대하여 $f^{-1}(V)$가 \emph{아핀 열린집합들의 유한 합집합}인 것이다.
예를 들어 $Y$가 아핀이고 $X$가 준콤팩트이면, \emph{모든} 사상
$f:X\to Y$는 준콤팩트이다 (I, 6.6.1).

$f:X\to Y$가 준콤팩트 사상이면, $Y$의 모든 열린집합 $V$에 대하여
$f$를 $f^{-1}(V)$에 제한한 사상
\[
  f^{-1}(V)\longrightarrow V
\]
도 준콤팩트임은 명백하다. 역으로 $(U_\alpha)$가 $Y$의 열린
덮개이고, 사상 $f:X\to Y$의 모든 제한
$f^{-1}(U_\alpha)\to U_\alpha$가 준콤팩트이면 $f$도 준콤팩트이다.

\oldpage[IV-1]{225}

따라서 $f:X\to Y$가 $S$-준스킴들의 $S$-사상이고, $S$의 열린 덮개
$(S_\lambda)$가 존재하여 $f$의 제한
\[
  g^{-1}(S_\lambda)\longrightarrow h^{-1}(S_\lambda)
\]
가 모두 준콤팩트이면 $f$는 준콤팩트이다. 여기서 $g$와 $h$는 구조
사상이다.

\begin{proposition}[1.1.2]
\label{IV.1.1.2-fr}
\begin{enumerate}
  \item[(i)] 몰입 $X\to Y$가 닫힌 몰입이거나, $Y$의 바탕공간이
  국소 뇌터이거나, 또는 $X$의 바탕공간이 뇌터이면 이 몰입은
  준콤팩트이다.
  \item[(ii)] 두 준콤팩트 사상의 합성은 준콤팩트이다.
  \item[(iii)] $f:X\to Y$가 준콤팩트 $S$-사상이면, 기초 준스킴의
  모든 확장 $g:S'\to S$에 대하여
  \[
    f_{(S')}:X_{(S')}\longrightarrow Y_{(S')}
  \]
  도 준콤팩트이다.
  \item[(iv)] $f:X\to X'$과 $g:Y\to Y'$이 두 준콤팩트
  $S$-사상이면
  \[
    f\times_S g:X\times_S Y\longrightarrow X'\times_S Y'
  \]
  는 준콤팩트이다.
  \item[(v)] 두 사상 $f:X\to Y$, $g:Y\to Z$의 합성 $g\circ f$가
  준콤팩트이고, $g$가 분리 사상이거나 $X$의 바탕공간이 국소
  뇌터이면 $f$는 준콤팩트이다.
  \item[(vi)] 사상 $f$가 준콤팩트일 필요충분조건은
  $f_{\mathrm{red}}$가 준콤팩트인 것이다.
\end{enumerate}
\end{proposition}

증명은 (I, 6.6.4)를 보라. (vi)는 다음의 더 일반적인 명제에서도
따른다는 점에 유의하자.

\begin{proposition}[1.1.3]
\label{IV.1.1.3-fr}
$f:X\to Y$, $g:Y\to Z$를 두 사상이라 하자. $g\circ f$가
준콤팩트이고 $f$가 전사이면 $g$는 준콤팩트이다.
\end{proposition}

\begin{proof}
실제로 $V$가 $Z$의 준콤팩트 열린집합이면, 가정에 의해
$f^{-1}(g^{-1}(V))$는 준콤팩트이다. $f$가 전사이므로
\[
  g^{-1}(V)=f\bigl(f^{-1}(g^{-1}(V))\bigr).
\]
따라서 $g^{-1}(V)$는 준콤팩트이다.
\end{proof}

\begin{corollary}[1.1.4]
\label{IV.1.1.4-fr}
$f:X\to Y$, $g:Y'\to Y$를 두 사상이라 하고,
$X'=X\times_YY'$, $f'=f_{(Y')}:X'\to Y'$라 하자. $g$가 준콤팩트
전사라고 가정한다 (각각, 전사라고 가정한다). 그러면 $f'$이
준콤팩트이면 $f$도 준콤팩트이다 (각각, $f'$이 우세이면 $f$도
우세이다).
\end{corollary}

\begin{proof}
$g':X'\to X$을 표준 사영이라 하면 $g'$은 전사이고 (I, 3.5.2),
$f\circ g'=g\circ f'$이다. $f'$이 준콤팩트이면 (각각, 우세이면),
$g$가 준콤팩트이므로 (1.1.2) (각각, 전사이므로) $g\circ f'$도
준콤팩트이다 (각각, 우세이다). $f\circ g'$이 준콤팩트이고
(각각, 우세이고) $g'$이 전사이므로, $f$는 준콤팩트이다 (1.1.3)
(각각, 우세이다).
\end{proof}

\oldpage[IV-1]{226}

간단히 하기 위해, 준스킴 $X$의 기약 성분의 모든 일반점을 $X$의
\emph{극대점}이라 부르겠다. $X=\operatorname{Spec}(A)$가
아핀이면 $X$의 극대점들은 $A$의 \emph{극소 소아이디얼}들이다.

\begin{proposition}[1.1.5]
\label{IV.1.1.5-fr}
$f:X\to Y$를 준콤팩트 사상이라 하자. 다음 조건들은 서로 동치이다.
\begin{enumerate}
  \item[a)] $f$는 우세이다.
  \item[b)] $Y$의 모든 극대점 $y$에 대하여
  $f^{-1}(y)\neq\varnothing$이다.
  \item[c)] $Y$의 모든 극대점 $y$에 대하여 $f^{-1}(y)$는 $X$의
  극대점을 하나 포함한다.
\end{enumerate}
\end{proposition}

\emph{a)}와 \emph{c)}의 동치는 (I, 6.6.5)에서 증명하였다.
\emph{c)}가 \emph{b)}를 함의함은 명백하다. 다른 한편 $X'$이
$f^{-1}(y)$와 만나는 $X$의 한 기약 성분이면, $f(X')$은 기약이고
$y$를 포함하므로 $Y$의 기약 성분
$Y'=\overline{\{y\}}$에 포함된다 $(0\textsubscript{I},~2.1.6)$.
$x$를 $X'$의 일반점이라 하면 반드시 $f(x)=y$이다
$(0\textsubscript{I},~2.1.5)$. 따라서 \emph{b)}는 \emph{c)}를
함의한다.

\begin{proposition}[1.1.6]
\label{IV.1.1.6-fr}
$f':X'\to Y$, $f'':X''\to Y$를 두 사상이라 하고, $X$를 합 준스킴
$X'\coprod X''$, $f:X\to Y$를 $X'$에서는 $f'$과 같고 $X''$에서는
$f''$과 같은 사상이라 하자. $f$가 준콤팩트일 필요충분조건은
$f'$과 $f''$이 모두 준콤팩트인 것이다.
\end{proposition}

이는 정의에서 곧바로 따른다.

\subsection{준분리 사상}
\label{subsection:IV.1.2-fr}

\begin{definition}[1.2.1]
\label{IV.1.2.1-fr}
$X$, $Y$를 두 준스킴이라 하자. 대각사상
\[
  \Delta_f=(1_X,1_X)_Y:X\longrightarrow X\times_YX
\]
이 준콤팩트이면 사상 $f:X\to Y$를 \emph{준분리}라고 한다. 이때
$X$가 $Y$ 위에서 \emph{준분리}라고도 한다. 준스킴 $X$가
$\operatorname{Spec}(\mathbf Z)$ 위에서 준분리이면 $X$를
\emph{준분리 준스킴}이라고 한다.
\end{definition}

정의상 모든 분리 사상은 준분리이다. 실제로 $\Delta_f$는 닫힌
몰입이다 (1.1.2, (i)). \emph{스킴} $X$는
$\operatorname{Spec}(\mathbf Z)$ 위에서 분리이므로 (I, 5.4.1)
준분리 준스킴이다.

\begin{proposition}[1.2.2]
\label{IV.1.2.2-fr}
\begin{enumerate}
  \item[(i)] 준스킴의 모든 단사사상 $f:X\to Y$, 특히 모든 몰입은
  준분리이다.
  \item[(ii)] $f:X\to Y$와 $g:Y\to Z$가 두 준분리 사상이면
  $g\circ f:X\to Z$는 준분리이다.
  \item[(iii)] $X$, $Y$를 두 $S$-준스킴, $f:X\to Y$를 준분리
  $S$-사상이라 하자. 모든 밑변경 $g:S'\to S$에 대하여 사상
  \[
    f_{(S')}:X_{(S')}\longrightarrow Y_{(S')}
  \]
  은 준분리이다.
  \item[(iv)] $f:X\to Y$, $f':X'\to Y'$이 두 준분리 $S$-사상이면
  \[
    f\times_Sf':X\times_SX'\longrightarrow Y\times_SY'
  \]
  은 준분리이다.
  \item[(v)] $f:X\to Y$, $g:Y\to Z$가 두 사상이고 $g\circ f$가
  준분리이면 $f$는 준분리이다.
  \item[(vi)] $f:X\to Y$가 준분리 사상이면
  $f_{\mathrm{red}}:X_{\mathrm{red}}\to Y_{\mathrm{red}}$도
  준분리이다.
\end{enumerate}
\end{proposition}

\begin{proof}
(I, 5.5.12)에 의해 (i), (ii), (iii)만 증명하면 충분하다.

모든 단사사상은 분리이므로 (I, 5.5.1), (i)은 자명하다. (iii)을
증명할 때는 곧바로 $Y=S$인 경우로 환원할 수 있고 (I, 3.3.11),
이때 주장은
\[
  \Delta_{f_{(S')}}=(\Delta_f)_{(S')}
\]
(I, 5.3.4)와 (1.1.2, (iii))에서 따른다. (ii)를 증명하기 위해
$X\times_YX$에서 $X$로 가는 두 사영을 $p$, $q$라 하자.
$i=(p,q)_Z$라 두면 다음 도표는 가환이다.
\[
\xymatrix{
  X\times_Y X \ar[r]^-{i} \ar[d]_-{\pi} & X\times_Z X
    \ar[d]^-{f\times_Z f} \\
  Y \ar[r]_-{\Delta_g} & Y\times_Z Y
}
\]

\oldpage[IV-1]{227}

여기서 $\pi$는 구조사상이고, 이 도표는 $X\times_YX$를
$(Y\times_ZY)$-준스킴 $Y$와 $X\times_ZX$의 곱과 동일시한다
(I, 5.3.5). $g$가 준분리이면 $\Delta_g$가 준콤팩트이므로
$i$도 준콤팩트이다 (1.1.2, (iii)). 더 나아가 $f$가 준분리이면
$\Delta_f$가 준콤팩트이므로 $i\circ\Delta_f$도 준콤팩트이다
(1.1.2, (ii)). 그런데 이 합성은 $\Delta_{g\circ f}$와 같다.
\end{proof}

\begin{corollary}[1.2.3]
\label{IV.1.2.3-fr}
\begin{enumerate}
  \item[(i)] $X$가 준분리 준스킴이면 모든 사상 $f:X\to Y$는
  준분리이다.
  \item[(ii)] $Y$가 준분리 준스킴이면, 사상 $f:X\to Y$가
  준분리일 필요충분조건은 준스킴 $X$가 준분리인 것이다.
\end{enumerate}
\end{corollary}

이는 $X$, $Y$, $\operatorname{Spec}(\mathbf Z)$에 (1.2.2, (ii),
(v))를 적용하면 곧바로 따른다.

\begin{proposition}[1.2.4]
\label{IV.1.2.4-fr}
$f:X\to Y$를 한 사상, $g:Y\to Z$를 준분리 사상이라 하자.
$g\circ f$가 준콤팩트이면 $f$도 준콤팩트이다.
\end{proposition}

\begin{proof}
$f$는 합성사상
\[
  X\xrightarrow{\Gamma_f}X\times_ZY
  \xrightarrow{\operatorname{pr}_2}Y
\]
이다 (I, 5.3.13). 한편 $\operatorname{pr}_2$는
$(g\circ f)\times_Z1_Y$와 동일시된다 (I, 3.3.4). 따라서
$g\circ f$가 준콤팩트이면 $\operatorname{pr}_2$도 준콤팩트이다
(1.1.2, (iv)). 끝으로 다음 가환도표를 생각하자.
\[
\xymatrix{
  X \ar[r]^-{\Gamma_f} \ar[d]_-{f} & X\times_ZY
    \ar[d]^-{f\times 1_Y} \\
  Y \ar[r]_-{\Delta_g} & Y\times_ZY
}
\tag{1.2.4.1}
\label{IV.1.2.4.1-fr}
\]
이 도표는 $X$를 $(Y\times_ZY)$-준스킴 $Y$와 $X\times_ZY$의 곱과
동일시한다 (I, 5.3.7). 가정에 의해 $\Delta_g$가 준콤팩트이므로
$\Gamma_f$도 준콤팩트이다 (1.1.2, (iii)). 따라서 결론은
(1.1.2, (ii))에서 따른다.
\end{proof}

\begin{proposition}[1.2.5]
\label{IV.1.2.5-fr}
$f:X\to Y$, $g:Y'\to Y$를 두 사상이라 하고,
$X'=X\times_Y Y'$, $f'=f_{(Y')}:X'\to Y'$이라 하자. $g$가
준콤팩트이고 전사이며 $f'$이 준분리이면, $f$는 준분리이다.
\end{proposition}

\begin{proof}
실제로 $(X'\times_{Y'}X')=(X\times_YX)_{(Y')}$이고
$\Delta_{f'}=(\Delta_f)_{(Y')}$이다 (I, 5.3.4). 사영
$X'\times_{Y'}X'\to X\times_YX$는 준콤팩트이고 (1.1.2, (iii))
전사이므로 (I, 3.5.2), (I, 3.3.11)에 따라 (1.1.4)를 적용할 수
있다. 따라서 $\Delta_{f'}$가 준콤팩트 사상이면 $\Delta_f$도
준콤팩트 사상이다.
\end{proof}

\begin{proposition}[1.2.6]
\label{IV.1.2.6-fr}
$f:X\to Y$를 한 사상, $(U_\alpha)$를 각 유도 준스킴
$U_\alpha$가 준분리인 열린집합들로 이루어진 $Y$의 덮개라 하자.
$f$가 준분리일 필요충분조건은 $X$가 각 열린집합
$f^{-1}(U_\alpha)$ 위에 유도하는 준스킴이 준분리인 것이다.
\end{proposition}

\begin{proof}
$X\times_YX$에서 $U_\alpha$의 역상은
$X_\alpha=f^{-1}(U_\alpha)$라 둘 때
$X_\alpha\times_{U_\alpha}X_\alpha$이고, $\Delta_f$의 제한
$X_\alpha\to X_\alpha\times_{U_\alpha}X_\alpha$는 $f$의 제한
$f_\alpha:X_\alpha\to U_\alpha$의 대각사상 $\Delta_{f_\alpha}$와
다르지 않다. 정의 (1.2.1)와 준콤팩트 사상이라는 개념의 국소성
(1.1.1)에 따라, $f$가 준분리일 필요충분조건은 각 사상
$f_\alpha$가 준분리인 것이다. 그런데 가정에 의해 사상
$U_\alpha\to\operatorname{Spec}(\mathbf Z)$는 준분리이므로,
$f_\alpha$가 준분리라는 것과 합성사상
$X_\alpha\xrightarrow{f_\alpha}U_\alpha\to\operatorname{Spec}(\mathbf Z)$
이 준분리라는 것은 같은 말이다 (1.2.2, (ii), (v)). 이로부터 명제가
따른다.
\end{proof}

\oldpage[IV-1]{228}

그러므로 어떤 사상이 준분리인지 확인하는 문제는 어떤 준스킴이
준분리인지 판정하는 다음 기준들을 제시하는 문제로 환원된다.

\begin{proposition}[1.2.7]
\label{IV.1.2.7-fr}
$X$를 한 준스킴, $(U_\alpha)$를 준콤팩트 열린집합들로 이루어진
$X$의 덮개라 하자. 다음 조건들은 서로 동치이다.
\begin{enumerate}
  \item[a)] $X$는 준분리이다.
  \item[b)] $X$의 모든 준콤팩트 열린집합 $U$에 대하여 정준 몰입
  $U\to X$는 준콤팩트이다 (다시 말해, $U$는 $X$ 안에서
  역콤팩트이다).
  \item[b')] $X$의 두 준콤팩트 열린집합의 교집합은 준콤팩트이다.
  \item[c)] 모든 지표쌍 $(\alpha,\beta)$에 대하여 교집합
  $U_\alpha\cap U_\beta$는 준콤팩트이다.
\end{enumerate}
\end{proposition}

\begin{proof}
준콤팩트 사상의 정의에 의해 조건 \emph{b)}와 \emph{b')}은 자명하게
동치이다. 준콤팩트 열린집합은 유한 개의 아핀 열린집합의 합집합이므로,
$X$의 두 준콤팩트 열린집합 $U$, $V$에 대하여
$U\times_{\mathbf Z}V$는 $X\times_{\mathbf Z}X$의 준콤팩트
열린집합이고 (I, 3.2.7), $\Delta_X$에 의한 그 역상은 $U\cap V$이다.
따라서 \emph{a)}는 \emph{b')}을 함의한다. \emph{b')}이
\emph{c)}를 함의함은 자명하다. 끝으로 \emph{c)}가 성립하면
$U_\alpha\times_{\mathbf Z}U_\beta$들은 준콤팩트 열린집합들로
$X\times_{\mathbf Z}X$를 덮는다. 사상
$\Delta_X:X\to X\times_{\mathbf Z}X$가 준콤팩트이기 위해서는
$\Delta_X$에 의한 $U_\alpha\times_{\mathbf Z}U_\beta$들의 역상
$U_\alpha\cap U_\beta$가 준콤팩트이면 충분하다 (1.1.1).
따라서 \emph{c)}는 \emph{a)}를 함의한다.
\end{proof}

\begin{corollary}[1.2.8]
\label{IV.1.2.8-fr}
바탕공간이 국소 뇌터인 모든 준스킴 $X$ (예를 들어 국소 뇌터
준스킴)는 준분리이다. 따라서 모든 사상 $f:X\to Y$는 준분리이다.
\end{corollary}

\begin{proposition}[1.2.9]
\label{IV.1.2.9-fr}
$X'$, $X''$을 두 준스킴, $X=X'\coprod X''$을 그 합,
$f':X'\to Y$, $f'':X''\to Y$를 두 사상, $f:X\to Y$를
$X'$에서는 $f'$과 같고 $X''$에서는 $f''$과 같은 사상이라 하자.
$f$가 준분리일 필요충분조건은 $f'$과 $f''$이 모두 준분리인 것이다.
\end{proposition}

\begin{proof}
실제로 $X\times_YX$는 네 준스킴
$X'\times_YX'$, $X'\times_YX''$, $X''\times_YX'$,
$X''\times_YX''$의 합이고, $\Delta_f$는 $X'$에서는
$\Delta_{f'}$과 같고 $X''$에서는 $\Delta_{f''}$과 같은 사상이다.
따라서 명제는 정의에서 곧바로 따른다.
\end{proof}

\subsection{국소 유한형 사상}
\label{subsection:IV.1.3-fr}

\begin{env}[1.3.1]
\label{IV.1.3.1-fr}
$A$가 한 환일 때, (가환) $A$-대수 $B$가 $B$의 유한 개 원소로
$A$-대수로서 생성되면 $B$를 \emph{유한 생성 $A$-대수}라고 부른다는 것을
상기하자. 동치인 표현으로, $B$가 다항식 대수
$A[T_1,\ldots,T_n]$를 그 아이디얼 하나로 나눈 몫과 동형이면 된다.
그러면 모든 (가환) $A$-대수 $A'$에 대하여 $B\otimes_AA'$은
유한 생성 $A'$-대수임이 명백하다. $B$가 유한 생성 $A$-대수이고 $C$가
유한 생성 $B$-대수이면 $C$는 유한 생성 $A$-대수이다. 실제로 $C$가
$B[T_1,\ldots,T_n]$의 몫이고 $B$가 $A[S_1,\ldots,S_m]$의
몫이면,
$B[T_1,\ldots,T_n]=B\otimes_AA[T_1,\ldots,T_n]$은
$A[S_1,\ldots,S_m,T_1,\ldots,T_n]$의 몫이므로 $C$도 그러하다.
\end{env}

\begin{definition}[1.3.2]
\label{IV.1.3.2-fr}
$f:X\to Y$를 준스킴의 사상, $x$를 $X$의 한 점,
$y=f(x)$라 하자. $y$의 아핀 열린근방 $V$와 $x$의 아핀 열린근방
$U$가 존재하여 $f(U)\subset V$이고 환 $A(U)$가 유한 생성
$A(V)$-대수이면, $f$가 \emph{$x$에서 유한형}이라고 한다. $f$가
$X$의 모든 점에서 유한형이면 $f$를 \emph{국소 유한형}이라고 한다.
\end{definition}

\oldpage[IV-1]{229}

\begin{proposition}[1.3.3]
\label{IV.1.3.3-fr}
$Y$가 국소 뇌터 준스킴이고 $f:X\to Y$가 국소 유한형 사상이면,
$X$는 국소 뇌터이다.
\end{proposition}

증명은 (I, 6.3.7)을 보라.

\begin{proposition}[1.3.4]
\label{IV.1.3.4-fr}
\begin{enumerate}
  \item[(i)] 모든 국소 몰입 사상은 국소 유한형이다.
  \item[(ii)] 두 사상 $f:X\to Y$, $g:Y\to Z$가 국소 유한형이면
  $g\circ f$도 국소 유한형이다.
  \item[(iii)] $f:X\to Y$가 국소 유한형 $S$-사상이면, 밑준스킴을
  바꾸는 모든 사상 $S'\to S$에 대하여
  $f_{(S')}:X_{(S')}\to Y_{(S')}$은 국소 유한형이다.
  \item[(iv)] $f:X\to X'$과 $g:Y\to Y'$이 국소 유한형
  $S$-사상이면 $f\times_Sg$도 국소 유한형이다.
  \item[(v)] 두 사상의 합성 $g\circ f$가 국소 유한형이면 $f$는
  국소 유한형이다.
  \item[(vi)] 사상 $f$가 국소 유한형이면 $f_{\mathrm{red}}$도
  국소 유한형이다.
\end{enumerate}
\end{proposition}

증명은 (I, 6.6.6)을 보라.

\begin{corollary}[1.3.5]
\label{IV.1.3.5-fr}
$f:X\to Y$를 국소 유한형 사상이라 하자. $Y'$이 국소 뇌터가 되는
모든 사상 $Y'\to Y$에 대하여 $X\times_YY'$은 국소 뇌터이다.
\end{corollary}

\begin{proof}
실제로 (1.3.4, (iii))에 의해
$f_{(Y')}:X\times_YY'\to Y'$은 국소 유한형이므로, (1.3.3)을
적용하면 된다.
\end{proof}

\begin{proposition}[1.3.6]
\label{IV.1.3.6-fr}
$\varphi:A\to B$를 환 준동형이라 하자. 이에 대응하는 사상
$f:\operatorname{Spec}(B)\to\operatorname{Spec}(A)$가 국소
유한형일 필요충분조건은 $B$가 유한 생성 $A$-대수인 것이다.
\end{proposition}

$f$가 준콤팩트라는 사실과 (I, 6.6.3)을 고려하여, 증명은
(I, 6.3.3)을 보라.

\begin{proposition}[1.3.7]
\label{IV.1.3.7-fr}
국소 유한형 사상 $f:X\to Y$가 전사일 필요충분조건은 모든 대수적으로
닫힌 체 $\Omega$에 대하여 $f$에 대응하는 사상
$X(\Omega)\to Y(\Omega)$ (I, 3.4.1)가 전사인 것이다.
\end{proposition}

증명은 (I, 6.3.10)을 보라.

\begin{env}[1.3.8]
\label{IV.1.3.8-fr}
$A$를 한 환이라 하자. $A$-대수 $B$가 어떤 유한 생성 $A$-대수 $C$와
$C$의 곱셈적 부분집합 $S$에 대하여 $S^{-1}C$ 꼴의 $A$-대수와
$A$-동형이면, $B$를 \emph{본질적으로 유한 생성된 $A$-대수}라고 한다.
\end{env}

\par\medskip\noindent
\begin{proposition}[1.3.9]
\label{IV.1.3.9-fr}
\begin{enumerate}
  \item[(i)] $B$가 본질적으로 유한 생성된 $A$-대수이고 $C$가
  본질적으로 유한 생성된 $B$-대수이면, $C$는 본질적으로 유한
  생성된 $A$-대수이다.
  \item[(ii)] $B$를 본질적으로 유한 생성된 $A$-대수, $A'$을
  $A$-대수라 하자. 그러면 $B'=B\otimes_AA'$은 본질적으로 유한
  생성된 $A'$-대수이다.
\end{enumerate}
\end{proposition}

\begin{proof}
\begin{enumerate}
  \item[(i)] $B=S'^{-1}B'$, $C=T'^{-1}C'$라 하자. 여기서 $B'$은
  유한 생성 $A$-대수이고 $C'$은 유한 생성 $B$-대수이며, $S'$과
  $T'$은 각각 $B'$과 $C'$의 곱셈적 부분집합이다. 그러면 $C'$은
  $S'^{-1}C''$ 꼴이고, 여기서 $C''$은 유한 생성 $B'$-대수이다.
  따라서 $C$도 $C''$의 분수환이다. $C''$은 유한 생성 $A$-대수이므로
  이것이 주장을 증명한다.
  \item[(ii)] $B$가 유한 생성 $A$-대수 $C$의 분수환이면, $B'$은
  유한 생성 $A'$-대수 $C\otimes_AA'$의 분수환이다. 이로부터
  (ii)가 따른다.
\end{enumerate}
\end{proof}

\oldpage[IV-1]{230}

\begin{env}[1.3.10]
\label{IV.1.3.10-fr}
$B$가 본질적으로 유한 생성된 \emph{국소} $A$-대수이면,
$B=C_{\mathfrak q}$ 꼴이다. 여기서 $C$는 유한 생성 $A$-대수이고
$\mathfrak q$는 $C$의 소아이디얼이다 (Bourbaki,
\emph{Alg. comm.}, chap. II, \S~5, prop. 11). $\mathfrak p$를
$A$에서 $\mathfrak q$의 역상이라 하고 $S=A-\mathfrak p$라 두자.
$C_{\mathfrak q}$는 $S^{-1}C$의 한 소아이디얼에서 취한 국소환이기도
하다. $S^{-1}C$는 $A_{\mathfrak p}=S^{-1}A$ 위의 유한 생성
대수이므로, $B$는 본질적으로 유한 생성된 $A_{\mathfrak p}$-대수이며,
더 나아가 준동형 $A_{\mathfrak p}\to B$는 \emph{국소 준동형}이다.
\end{env}

\begin{proposition}[1.3.11]
\label{IV.1.3.11-fr}
$B$가 본질적으로 유한 생성된 국소 $A$-대수이면, $B$는
$C_{\mathfrak q}$ 꼴의 $A$-대수를 나눈 몫 $A$-대수와 $A$-동형이다.
여기서 $C$는 다항식 대수 $A[T_1,\ldots,T_n]$이고
$\mathfrak q$는 $C$의 소아이디얼이다.
\end{proposition}

\begin{proof}
실제로 정의에 의해 $B$는 $C'_{\mathfrak q'}$과 동형이다. 여기서
$C'$은 유한 생성 $A$-대수이고 $\mathfrak q'$은 $C'$의 소아이디얼이다.
그런데 $C'=C/\mathfrak b$이고, 여기서
$C=A[T_1,\ldots,T_n]$이며 $\mathfrak b$는 $C$의 아이디얼이다.
따라서 $\mathfrak q'=\mathfrak q/\mathfrak b$이고, 여기서
$\mathfrak q$는 $C$의 소아이디얼이며, $C'_{\mathfrak q'}$은
$C_{\mathfrak q}/\mathfrak bC_{\mathfrak q}$과 동형이다.
\end{proof}

\subsection{국소 유한 표시 사상}
\label{subsection:IV.1.4-fr}

\begin{env}[1.4.1]
\label{IV.1.4.1-fr}
환 $A$가 주어졌다고 하자. (가환) $A$-대수 $B$가 $A$ 위의 다항식
대수를 그 유한 생성 아이디얼 하나로 나눈 몫
$A[T_1,\ldots,T_n]/\mathfrak a$와 동형이면, $B$를
\emph{유한 표시 $A$-대수}라고 하겠다. 모든 (가환) $A$-대수
$A'$에 대하여 $B\otimes_AA'$은 유한 표시 $A'$-대수이다. 실제로
이는 $A'[T_1,\ldots,T_n]$을 $\mathfrak a\otimes_AA'$의 이 환
안에서의 정준상으로 나눈 몫과 동형이고, 이 정준상은 명백히 유한 생성
$A'[T_1,\ldots,T_n]$-가군이다. $B$가 유한 표시 $A$-대수이고
$C$가 유한 표시 $B$-대수이면, $C$는 유한 표시 $A$-대수이다.
실제로
\[
  B=A[S_1,\ldots,S_m]/\mathfrak a,\qquad
  C=B[T_1,\ldots,T_n]/\mathfrak b
\]
라 하자. 여기서 $\mathfrak a$와 $\mathfrak b$는 각각
$A[S_1,\ldots,S_m]$와 $B[T_1,\ldots,T_n]$의 유한 생성
아이디얼이다. 환 $B[T_1,\ldots,T_n]$은
\[
  A[S_1,\ldots,S_m,T_1,\ldots,T_n]/\mathfrak a'
\]
과 동형이고, 여기서 $\mathfrak a'$은
$\mathfrak a\otimes_AA[T_1,\ldots,T_n]$의 정준상이므로
$A[S_1,\ldots,S_m,T_1,\ldots,T_n]$의 유한 생성 아이디얼이다.
$B[T_1,\ldots,T_n]$의 아이디얼 $\mathfrak b$는
$\mathfrak b'/\mathfrak a'$ 꼴이고, 여기서 $\mathfrak b'$은
$A[S_1,\ldots,S_m,T_1,\ldots,T_n]$의 아이디얼이다.
$\mathfrak a'$과 $\mathfrak b'/\mathfrak a'$이
$A[S_1,\ldots,S_m,T_1,\ldots,T_n]$ 위의 유한 생성 가군이므로
$\mathfrak b'$도 그러하다. 따라서
$A[S_1,\ldots,S_m,T_1,\ldots,T_n]/\mathfrak b'$과 동형인
$C$는 실제로 유한 표시 $A$-대수이다. 끝으로 $A$가
\emph{뇌터}이면 $A[T_1,\ldots,T_n]$도 뇌터이므로, 모든 유한 생성
$A$-대수는 유한 표시를 갖는다는 점에 유의하자.
\end{env}

\begin{definition}[1.4.2]
\label{IV.1.4.2-fr}
$f:X\to Y$를 준스킴의 사상, $x$를 $X$의 한 점,
$y=f(x)$라 하자. $y$의 아핀 열린근방 $V$와 $x$의 아핀 열린근방
$U$가 존재하여 $f(U)\subset V$이고 환 $A(U)$가 유한 표시
$A(V)$-대수이면, $f$가 \emph{$x$에서 유한 표시}라고 한다. $f$가
$X$의 모든 점에서 유한 표시이면 $f$를 \emph{국소 유한 표시}라고 한다.
\end{definition}

$Y$가 국소 뇌터이면, $f$가 국소 유한형이라고 하는 것과 국소 유한
표시라고 하는 것은 같은 말이다.

\begin{proposition}[1.4.3]
\label{IV.1.4.3-fr}
\begin{enumerate}
  \item[(i)] 모든 국소 동형사상은 국소 유한 표시이다.
  \item[(ii)] 두 사상 $f:X\to Y$, $g:Y\to Z$가 국소 유한 표시이면
  $g\circ f$도 국소 유한 표시이다.

\oldpage[IV-1]{231}

  \item[(iii)] $f:X\to Y$가 국소 유한 표시 $S$-사상이면, 밑준스킴을
  바꾸는 모든 사상 $S'\to S$에 대하여
  $f_{(S')}:X_{(S')}\to Y_{(S')}$은 국소 유한 표시이다.
  \item[(iv)] $f:X\to X'$과 $g:Y\to Y'$이 국소 유한 표시
  $S$-사상이면 $f\times_Sg$도 국소 유한 표시이다.
  \item[(v)] 두 사상의 합성 $g\circ f$가 국소 유한 표시이고 $g$가
  국소 유한형이면, $f$는 국소 유한 표시이다.
\end{enumerate}
\end{proposition}

\begin{proof}
주장 (i)은 자명하다. (iii)을 증명할 때는 $S=Y$인 경우로 환원할 수
있다 (I, 3.3.11). $x'$을 $X_{(S')}$의 한 점, $s'$과 $x$를 각각
$S'$과 $X$ 위로 내린 그 사영, $s$를 $S$ 위에서 $s'$과 $x$의 공통
사영이라 하자. 가정에 의해 $s$의 아핀 열린근방 $V$와 $x$의 아핀
열린근방 $U$가 존재하여 $U$의 상은 $V$에 포함되고 $A(U)$는 유한
표시 $A(V)$-대수이다. $s'$의 아핀 열린근방 $V'$을 택하되 $S$에서의
그 상이 $V$에 포함되게 하자. 그러면 $U\times_VV'$은 $x'$의 아핀
열린근방이고, $S'$에서의 그 상은 $V'$에 포함되며, 그 환
$A(U)\otimes_{A(V)}A(V')$은 유한 표시 $A(V')$-대수이다 (1.4.1).

(ii)를 증명하기 위해 $x\in X$를 한 점이라 하고, $y=f(x)$,
$z=g(f(x))$라 두자. 가정에 의해 $z$의 아핀 열린근방
$W=\operatorname{Spec}(C)$와 $y$의 아핀 열린근방
$V=\operatorname{Spec}(B)$가 존재하여 $g(V)\subset W$이고 $B$는
유한 표시 $C$-대수이다. 다른 한편 $y$의 아핀 열린근방
$V'=\operatorname{Spec}(B')$와 $x$의 아핀 열린근방
$U=\operatorname{Spec}(A)$가 존재하여 $f(U)\subset V'$이고 $A$는
유한 표시 $B'$-대수이다. $V$ 안의 열린집합
$D(s)=\operatorname{Spec}(B_s)$가 $V'$에 포함된 $y$의 근방이 되게
하는 $s\in B$를 택하고, $U'=f^{-1}(D(s))\subset U$라 두자. 그러면
$U'=\operatorname{Spec}(A\otimes_{B'}B_s)$이고,
$A\otimes_{B'}B_s$는 유한 표시 $B_s$-대수이다 (1.4.1). 다른 한편
$B_s=B[T]/(1-sT)$는 정의상 유한 표시 $B$-대수이다. 따라서
(1.4.1)에 의해 $A\otimes_{B'}B_s$는 유한 표시 $C$-대수이고,
이것이 (ii)를 증명한다.

(iv)는 (i), (ii), (iii)에서 따른다는 것을 알고 있다 (I, 3.5.1).
(v)를 증명하기 위해 다음 가환도표를 생각하자.
\[
\xymatrix{
X \ar[r]^{\Gamma_f} \ar[d] & X\times_ZY \ar[d]^{f\times 1} \\
Y \ar[r]_{\Delta_g} & Y\times_ZY
}
\]
이 도표는 $X$를 $(Y\times_ZY)$-준스킴 $Y$와 $X\times_ZY$의 곱과
동일시한다 (I, 5.3.7; 여기서 $\Delta_g$는 대각사상이다).
$\Delta_g$가 국소 유한 표시임을 알면 (iii)에 의해 $\Gamma_f$도
그러하다. 다른 한편 $f$는
\[
  X\xrightarrow{\Gamma_f}X\times_ZY\xrightarrow{p_2}Y
\]
로 인수분해되고 (I, 5.3.13), 여기서 사영 $p_2$는
$(g\circ f)\times_Z1_Y$와 같다. 가정에 의해 $g\circ f$가 국소
유한 표시이므로 (iv)에 의해 $p_2$도 그러하며, 끝으로 (ii)에 의해
$f$도 국소 유한 표시임이 따른다. 따라서 다음 따름정리를 증명하는
것으로 환원된다.
\end{proof}

\begin{corollary}[1.4.3.1]
\label{IV.1.4.3.1-fr}
$g:Y\to Z$를 국소 유한형 사상이라 하자. 그러면 대각사상
$\Delta_g:Y\to Y\times_ZY$는 국소 유한 표시이다.
\end{corollary}

\oldpage[IV-1]{232}

\begin{proof}
$Z=\operatorname{Spec}(A)$, $Y=\operatorname{Spec}(B)$이고 $B$가
유한 생성 $A$-대수인 경우로 환원할 수 있다. 이때 사상 $\Delta_g$에
대응하는 것은 $\delta(x\otimes y)=xy$로 주어지는
\emph{증대 준동형} $\delta:B\otimes_AB\to B$이다. 정의 (1.4.2)에
비추어 보면, $\delta$의 핵 $\mathfrak J$가 \emph{유한 생성}
아이디얼임을 보이면 충분하고, 이는 (0, 20.4.4)에서 따른다.
\end{proof}

\begin{proposition}[1.4.4]
\label{IV.1.4.4-fr}
$A$를 한 환, $B$를 $A$-대수, $B'$을 다항식 대수
$A[T_1,\ldots,T_n]$, $u:B'\to B$를 전사 $A$-대수 준동형이라
하자. $B$가 유한 표시 $A$-대수일 필요충분조건은 $u$의 핵
$\mathfrak J$가 $B'$의 유한 생성 아이디얼인 것이다.
\end{proposition}

\begin{proof}
이 조건이 충분함은 정의 (1.4.1)에서 따른다. 필요함을 보이기 위해
사상 $g:\operatorname{Spec}(B')\to\operatorname{Spec}(A)$가 국소
유한형이라는 점에 유의하자. $B$가 유한 표시 $A$-대수이면 $u$에
대응하는 사상 $f:\operatorname{Spec}(B)\to\operatorname{Spec}(B')$에
대하여 $g\circ f:\operatorname{Spec}(B)\to\operatorname{Spec}(A)$가
국소 유한 표시이므로, (1.4.3, (v))에 의해 $f$도 국소 유한 표시이다.
그런데 아이디얼 $\mathfrak J$가 유한 생성임을 증명하려면
$\widetilde{\mathfrak J}$가 유한 생성 $\widetilde{B'}$-가군임을
증명하면 충분하다 (II, 6.1.4.1). 정의 (1.4.2)로 돌아가면 끝으로
다음 보조정리를 증명하는 것으로 환원된다.
\end{proof}

\begin{lemma}[1.4.4.1]
\label{IV.1.4.4.1-fr}
$\mathfrak a$를 환 $C$의 한 아이디얼, $s$를 $C$의 한 원소라 하고,
$C_s/\mathfrak a_s$가 유한 표시 $C$-대수라고 하자. 그러면
$\mathfrak a_s$는 환 $C_s$의 유한 생성 아이디얼이다.
\end{lemma}

\begin{proof}
가정에 의해 $C$-다항식 대수 $C'=C[T_1,\ldots,T_n]$과 전사
$C$-준동형 $v:C'\to C_s/\mathfrak a_s$가 존재하며, 그 핵
$\mathfrak b$는 유한 생성이다. $p:C_s\to C_s/\mathfrak a_s$를
정준 준동형이라 하자. 모든 $i$에 대하여
$p(t_i/s^k)=v(T_i)$가 되는 $t_i\in C$가 있다 (지수 $k$는 모든
지표 $i$에 대하여 같다고 해도 된다). 이제 $C_s$-대수
$D=C_s[T_1,\ldots,T_n]$과 $w(T_i)=t_i/s^k$로 주어지는
$C_s$-대수 준동형 $w:D\to C_s$를 생각하자. $w$는 명백히 전사이고,
따라서 $v'=p\circ w:D\to C_s/\mathfrak a_s$도 전사이다. 다른 한편
$D$의 모든 다항식은 $P\in C'=C[T_1,\ldots,T_n]$인 $P/s^m$ 꼴로
쓸 수 있고,
\[
  v'(P/s^m)=(1/s^m)v(P)
\]
이다. $1/s$가 $C_s$에서 가역이므로, $C_s/\mathfrak a_s$에서
$v'(P/s^m)=0$이라는 관계는 $v(P)=0$과 동치이다. 따라서 $v'$의 핵
$\mathfrak b'$은 $D$ 안에서 $\mathfrak b$의 정준상으로 생성되며,
특히 환 $D$에서 유한 생성이다. 그런데
\[
  \mathfrak b'=v'^{-1}(0)=w^{-1}(p^{-1}(0))
  =w^{-1}(\mathfrak a_s),
\]
이고 $w$가 전사이므로 $\mathfrak a_s=w(\mathfrak b')$이다. 이는
$\mathfrak a_s$가 $C_s$의 유한 생성 아이디얼임을 증명한다.
\end{proof}

\begin{corollary}[1.4.5]
\label{IV.1.4.5-fr}
$X$, $Y$를 두 준스킴, $j:X\to Y$를 한 몰입, $U$를 $j(X)$가
그 안에서 닫힌 $Y$의 열린집합이라 하고, $\mathcal J$를 $j$에
대응하는 $U$의 닫힌 부분준스킴을 정의하는 $\mathcal O_U$의
준연접 아이디얼이라 하자. $j$가 국소 유한 표시일 필요충분조건은
$\mathcal J$가 유한 생성 $\mathcal O_U$-가군인 것이다.
\end{corollary}

\par\medskip\noindent
\begin{proposition}[1.4.6]
\label{IV.1.4.6-fr}
$\varphi:A\to B$를 환 준동형이라 하자. 이에 대응하는 사상
$f:\operatorname{Spec}(B)\to\operatorname{Spec}(A)$가 국소 유한
표시일 필요충분조건은 $B$가 유한 표시 $A$-대수인 것이다.
\end{proposition}

\begin{proof}
이 조건이 충분함은 자명하므로 필요함만 증명하면 된다. $f$가 국소
유한 표시이면 (1.3.6)에 의해 이미 $B$는 유한 생성 $A$-대수이다.
따라서 전사 $A$-대수 준동형
$u:B'=A[T_1,\ldots,T_n]\to B$가 존재한다. (1.4.3, (v))에 의해
몰입 $j:\operatorname{Spec}(B)\to\operatorname{Spec}(B')$은 국소
유한 표시이다.

\oldpage[IV-1]{233}

그러므로 $\mathfrak J$를 $u$의 핵이라 하면, (1.4.5)에 의해
$\widetilde{\mathfrak J}$는 유한 생성 $\widetilde{B'}$-가군이고,
따라서 (II, 6.1.4.1)에 의해 $\mathfrak J$는 유한 생성 아이디얼이다.
\end{proof}

\begin{proposition}[1.4.7]
\label{IV.1.4.7-fr}
$\rho:A\to B$를 $B$를 유한 $A$-대수로 만드는 환 준동형이라 하자.
$B$가 유한 표시 $A$-대수일 필요충분조건은 $B$가 유한 표시
$A$-가군인 것이다.
\end{proposition}

다음 보조정리를 사용한다.

\begin{lemma}[1.4.7.1]
\label{IV.1.4.7.1-fr}
$B$가 유한 $A$-대수이면, 전사 $A$-대수 준동형 $u:B'\to B$가
존재하며, 여기서 $B'$은 유한 $A$-대수이자 유한 표시 $A$-대수이고
자유 $A$-가군이다.
\end{lemma}

\begin{proof}
실제로 $A$-대수 $B$의 유한 생성계 $(a_i)_{1\leq i\leq m}$을 각
$a_i$가 관계 $F_i(a_i)=0$을 만족하도록 택할 수 있다. 여기서
$F_i\in A[X]$는 차수가 $>0$인 일계수 다항식이다. 따라서 몫
$A$-대수 $B'_i=A[X]/(F_i)$는 $A$ 위의 유한 계수 자유 가군이다.
$B'_i$에서 $X$의 유를 $c_i$라 하자. $u_i(c_i)=a_i$를 만족하는
$A$-대수 준동형 $u_i:B'_i\to B$가 존재한다. 이제
$B'$을 텐서곱
$B'_1\otimes_AB'_2\otimes_A\cdots\otimes_AB'_m$으로 잡고
$u_1\otimes u_2\otimes\cdots\otimes u_m:B'\to B$를 생각하면 충분하며,
이 준동형은 구성상 전사이다. 또 $B''=A[T_1,\ldots,T_m]$라 하고
$b'_i$를 $B'$에서 $c_i$의 표준상이라 하면, $v(T_i)=b'_i$
$(1\leq i\leq m)$로 주어지는 $A$-대수 준동형 $v:B''\to B'$은
전사이다. 그 핵 $\mathfrak b'$은 $F_i(T_i)$들로 생성되는 $B''$의
아이디얼이므로 유한 생성이다. 이는 $B'$이 유한 표시
$A$-대수임을 증명한다.
\end{proof}

\begin{proof}
이 보조정리의 표기를 그대로 쓰고 $w=u\circ v:B''\to B$라 하자.
$\mathfrak b$와 $\mathfrak a$가 각각 $w$와 $u$의 핵이면, $v$가
전사이므로 $\mathfrak a=v(\mathfrak b)$이다. 이제 (1.4.4)에 의해,
$B$가 유한 표시 $A$-대수이면 $\mathfrak b$는 $B''$의 유한 생성
아이디얼이고, 따라서 $\mathfrak a$는 $B'$의 유한 생성 아이디얼이다.
$B'$이 유한 $A$-대수이므로 $\mathfrak a$는 유한 생성 $A$-가군이고,
$B'$이 자유 $A$-가군이므로 $B$는 유한 표시 $A$-가군이다. 역으로
$B$가 유한 표시 $A$-가군이면, $\mathfrak a$는 유한 생성 $A$-가군이고
(Bourbaki, \emph{Alg. comm.}, chap. I, \S~2,
n\textsuperscript{o}~8, lemme 9), 따라서 더욱이 $B'$의 유한 생성
아이디얼이다. 그러므로 정의상 $B$는 유한 표시 $B'$-대수이고,
$B'$이 유한 표시 $A$-대수이므로 $B$는 유한 표시 $A$-대수이다.
\end{proof}

\subsection{유한형 사상}
\label{subsection:IV.1.5-fr}

\begin{proposition}[1.5.1]
\label{IV.1.5.1-fr}
$f:X\to Y$를 한 사상, $(U_\alpha)$를 아핀 열린집합들로 이루어진
$Y$의 한 덮개라 하자. 다음 조건들은 서로 동치이다.
\begin{enumerate}
  \item[a)] $f$는 국소 유한형이고 준콤팩트이다.
  \item[b)] 모든 $\alpha$에 대하여 $f^{-1}(U_\alpha)$는 유한 개의
  아핀 열린집합 $V_{\alpha i}$들의 합집합이고, 각 환
  $A(V_{\alpha i})$는 유한 생성 $A(U_\alpha)$-대수이다.
  \item[c)] $Y$의 모든 아핀 열린집합 $U$에 대하여 $f^{-1}(U)$는
  유한 개의 아핀 열린집합 $V_j$들의 합집합이고, 환 $A(V_j)$들은
  유한 생성 $A(U)$-대수이다.
\end{enumerate}
\end{proposition}

증명은 (I, 6.3.2 및 6.6.3)을 보라.

\begin{definition}[1.5.2]
\label{IV.1.5.2-fr}
사상 $f$가 명제 (1.5.1)의 동치인 조건들을 만족하면 $f$를
\emph{유한형}이라고 한다.
\end{definition}

\oldpage[IV-1]{234}

\begin{proposition}[1.5.3]
\label{IV.1.5.3-fr}
$f:X\to Y$가 유한형 사상이고 $Y$가 뇌터이면 $X$도 뇌터이다.
\end{proposition}

증명은 (I, 6.3.7)을 보라.

\begin{proposition}[1.5.4]
\label{IV.1.5.4-fr}
\begin{enumerate}
  \item[(i)] $j:X\to Y$를 몰입 사상이라 하자. $j$가 준콤팩트이면
  $j$는 유한형이다. 특히 $j$가 닫힌 몰입 사상이거나, $X$의
  바탕공간이 뇌터이거나, 또는 $Y$의 바탕공간이 국소 뇌터이면
  $j$는 준콤팩트이다.
  \item[(ii)] 두 유한형 사상의 합성은 유한형이다.
  \item[(iii)] $f:X\to Y$가 유한형 $S$-사상이면, 밑준스킴을 바꾸는
  모든 사상 $S'\to S$에 대하여
  $f_{(S')}:X_{(S')}\to Y_{(S')}$은 유한형이다.
  \item[(iv)] $f:X\to X'$과 $g:Y\to Y'$이 유한형 $S$-사상이면
  $f\times_Sg$는 유한형이다.
  \item[(v)] 두 사상의 합성 $g\circ f$가 유한형이고, $g$가
  준분리이거나 $X$가 뇌터이면 $f$는 유한형이다.
  \item[(vi)] 사상 $f$가 유한형이면 $f_{\mathrm{red}}$도 유한형이다.
\end{enumerate}
\end{proposition}

이는 정의와 (1.1.2), (1.2.4), (1.3.4)에서 곧바로 따른다. 다만
$X$가 뇌터인 경우의 (v)는 (I, 6.3.6)에서 증명하였다.

\begin{corollary}[1.5.5]
\label{IV.1.5.5-fr}
$f:X\to Y$를 유한형 사상이라 하자. $Y'$이 뇌터인 모든 사상
$Y'\to Y$에 대하여 $X\times_YY'$은 뇌터이다.
\end{corollary}

이는 (1.5.4, (iii))과 (1.5.3)에서 따른다.

\begin{corollary}[1.5.6]
\label{IV.1.5.6-fr}
$X$를 국소 뇌터 준스킴 $S$ 위의 유한형 준스킴이라 하자. 그러면
모든 $S$-사상 $f:X\to Y$는 유한형이다.
\end{corollary}

증명은 (I, 6.3.9)를 보라.

\begin{proposition}[1.5.7]
\label{IV.1.5.7-fr}
$\varphi:A\to B$를 환 준동형이라 하자. 이에 대응하는 사상
$f:\operatorname{Spec}(B)\to\operatorname{Spec}(A)$가 유한형일
필요충분조건은 $\varphi$가 $B$를 유한 생성 $A$-대수로 만드는 것이다.
\end{proposition}

증명은 (I, 6.3.3)을 보라.

\subsection{유한 표시 사상}
\label{subsection:IV.1.6-fr}

\begin{definition}[1.6.1]
\label{IV.1.6.1-fr}
$X$, $Y$를 두 준스킴이라 하자. 사상 $f:X\to Y$가 다음 세 조건을
만족하면 $f$를 \emph{유한 표시}라고 한다.
\begin{enumerate}
  \item[(i)] $f$는 국소 유한 표시이다.
  \item[(ii)] $f$는 준콤팩트이다. (이는 (i)가 성립할 때 $f$가
  유한형이라는 것과 동치이다 (1.5.2).)
  \item[(iii)] $f$는 준분리이다.
\end{enumerate}
이때 $X$가 \emph{$Y$ 위에서 유한 표시}라고도 하고,
\emph{유한 표시 $Y$-준스킴}이라고도 한다.
\end{definition}

$f$가 분리 사상이거나 $X$가 국소 뇌터이면 조건 (iii)은 자동으로
성립한다 (1.2.8). $Y$가 국소 뇌터이면 $f$가 \emph{유한 표시}라는
것과 $f$가 \emph{유한형}이라는 것은 같은 말이다. 실제로 뒤의
조건에서 $X$가 국소 뇌터임이 따른다 (1.3.3).

\oldpage[IV-1]{235}

\begin{proposition}[1.6.2]
\label{IV.1.6.2-fr}
\begin{enumerate}
  \item[(i)] 준콤팩트이고 국소 유한 표시인 모든 몰입 사상은 유한
  표시이다. 특히 모든 준콤팩트 열린 몰입 사상은 유한 표시이다.
  \item[(ii)] 두 유한 표시 사상의 합성은 유한 표시이다.
  \item[(iii)] $X$, $Y$를 두 $S$-준스킴, $f:X\to Y$를 유한 표시
  $S$-사상이라 하자. 그러면 모든 사상 $S'\to S$에 대하여
  $f_{(S')}:X_{(S')}\to Y_{(S')}$은 유한 표시이다.
  \item[(iv)] $f:X\to Y$, $f':X'\to Y'$을 두 유한 표시 $S$-사상이라
  하자. 그러면
  $f\times_Sf':X\times_SX'\to Y\times_SY'$은 유한 표시이다.
  \item[(v)] 두 사상의 합성 $g\circ f$가 유한 표시이고 $g$가
  준분리이면서 국소 유한 표시이면, $f$는 유한 표시이다.
\end{enumerate}
\end{proposition}

이는 (1.1.2), (1.2.2), (1.2.4), (1.4.3)에서 곧바로 따른다.

특히 (1.6.2, (iii))에 의해, $f$가 유한 표시 사상이고 $U$가 $Y$의
열린집합이면 $f$의 제한 $f^{-1}(U)\to U$는 유한 표시이다. 역으로
$(U_\alpha)$를 아핀 열린집합들로 이루어진 $Y$의 덮개라 하고,
$f$의 제한 $f^{-1}(U_\alpha)\to U_\alpha$가 모두 유한 표시라고
가정하자. 그러면 $f$는 명백히 국소 유한 표시이고, (1.1.1)에 의해
준콤팩트이며, (1.2.6)에 의해 준분리이다. 따라서 $f$는 유한
표시이다. 그러므로 사상 $f:X\to Y$가 유한 표시라는 성질은
\emph{$Y$ 위에서 국소적}이라고 할 수 있다.

$X$가 준분리 준스킴이면 모든 사상 $f:X\to Y$는 준분리이다
(1.2.3). 따라서 $f$가 준콤팩트이고 국소 유한 표시이면 $f$는 유한
표시이다.

특히 다음을 얻는다.

\begin{corollary}[1.6.3]
\label{IV.1.6.3-fr}
$\varphi:A\to B$를 환 준동형이라 하자. 이에 대응하는 사상
$f:\operatorname{Spec}(B)\to\operatorname{Spec}(A)$가 유한 표시일
필요충분조건은 $\varphi$가 $B$를 유한 표시 $A$-대수로 만드는 것이다.
\end{corollary}

실제로 이 조건의 필요성은 (1.4.6)에서 따른다.

\begin{remark}[1.6.4]
\label{IV.1.6.4-fr}
정의 (1.6.1)에서 조건 (iii)은 다른 두 조건의 귀결이 아니다. 예를
들어 바탕공간이 뇌터가 아닌 아핀 스킴 $Y$와 $Y$의 준콤팩트가 아닌
열린집합 $U$를 택하자. $Y$와 동형인 두 준스킴 $Y_1$, $Y_2$를 $U$에
대응하는 열린집합 $U_1$, $U_2$를 따라 붙여서 얻는 준스킴을 $X$라
하자 (0$_{\mathrm I}$, 4.1.7). 그러면 $X$는 각각 $Y_1$, $Y_2$와
동형인 두 열린집합의 합집합이며, 이들을 $Y_1$, $Y_2$와 동일시하면
$Y_1\cap Y_2=U$이다. 또한 각 $Y_i$에서 주어진 동형사상
$Y_i\to Y$와 일치하는 사상 $f:X\to Y$가 존재한다 ($i=1,2$).
$f$는 국소 동형사상이므로 (1.6.1)의 조건 (i)를 만족한다. $Y$의
준콤팩트 열린집합의 역상은 $Y_1$과 $Y_2$에 있는 그 두 사본의
합집합이므로 (ii)도 만족한다. 그러나 $Y_1\cap Y_2$가 준콤팩트가
아니므로 $f$는 준분리가 아니다 (1.2.7).
\end{remark}

\begin{proposition}[1.6.5]
\label{IV.1.6.5-fr}
$X'$, $X''$을 두 준스킴, $X=X'\amalg X''$을 그 합, $f':X'\to Y$와
$f'':X''\to Y$를 두 사상, $f:X\to Y$를 $X'$에서 $f'$과 일치하고
$X''$에서 $f''$과 일치하는 사상이라 하자. $f$가 유한 표시일
필요충분조건은 $f'$과 $f''$이 모두 유한 표시인 것이다.
\end{proposition}

\par\medskip\noindent
\begin{proof}
$f$가 정의 (1.6.1)의 세 성질 가운데 어느 하나를 가질 필요충분조건이
$f'$과 $f''$이 모두 그 성질을 갖는 것임을 보이면 충분하다.

\oldpage[IV-1]{236}

국소 유한 표시라는 성질은 $X$ 위에서 국소적이므로 이는 그 성질에
대해서는 명백하다. 준콤팩트라는 성질에 대해서는 (1.1.6)에서,
준분리라는 성질에 대해서는 (1.2.9)에서 보았다.
\end{proof}

\subsection{앞선 결과들의 개선}
\label{subsection:IV.1.7-fr}

이 절에서는 앞 장들에서 증명한 명제들 가운데, 위에서 도입한 새로운
유한성 조건을 사용하여 명제를 더 강하게 만들 수 있는 것들을
열거한다.

\begin{env}[1.7.1]
\label{IV.1.7.1-fr}
(I, 6.4.2, 6.4.3, 6.4.9)의 명제들에서는 $X$와 $Y$가 체 $K$ 위에서
유한형이라고 가정하는 대신, 이들이 $K$ 위에서 \emph{국소
유한형}이라고만 가정할 수 있다. (I, 6.4.2)에 대해서는 준스킴 $X$의
한 점이 $X$의 모든 아핀 열린집합에서 닫혀 있으면 $X$에서도 닫혀
있다는 것과, $K$ 위에서 국소 유한형인 아핀 스킴은 그 자체로
유한형이라는 것에 유의하면 충분하다 (1.5.1). (I, 6.4.9)에 대해서는
(1.3.7)과 (1.3.4, (v))를 사용한다.
\end{env}

\begin{env}[1.7.2]
\label{IV.1.7.2-fr}
(I, 6.5.1, (ii))에서는 $S$가 국소 뇌터이고 $Y$가 $S$ 위에서
유한형이라고 가정하는 대신, $Y$가 $S$ 위에서 \emph{국소 유한
표시}라고만 가정할 수 있다. 이는 정의 (1.4.2)를 적용하면 증명에서
곧바로 드러난다. 마찬가지로 (I, 6.5.4, (ii))와 (I, 6.5.5, (ii))에서는
$Y$가 $S$ 위에서 국소 유한 표시라는 조건 아래 $f$가 $S$-사상이라고만
가정하면 충분하다.
\end{env}

\begin{env}[1.7.3]
\label{IV.1.7.3-fr}
(I, 7.1.11)의 두 번째 주장에서는 $S$가 국소 뇌터이고 $Y$가 $S$
위에서 유한형이라고 가정하는 대신, $Y$가 $S$ 위에서 국소 유한
표시라고만 가정할 수 있다. 실제로 그 증명은 (I, 6.5.1, (ii))에
의존한다. (I, 7.1.12에서 7.1.15까지)에서도 마찬가지이다.
\end{env}

\begin{env}[1.7.4]
\label{IV.1.7.4-fr}
(I, 9.2.2)에서는 세 가정을, $f$가 \emph{준콤팩트이고 준분리}라는
하나의 가정으로 바꿀 수 있다. 이 조건은 각 원래 가정의 귀결이며,
이는 (1.2.6)과 (1.2.7)에서 따른다. (I, 9.2.1)의 가정은 정확히
$f$가 준콤팩트이고 준분리라는 뜻임에 유의하자.
\end{env}

\begin{env}[1.7.5]
\label{IV.1.7.5-fr}
(I, 9.3.1, 9.3.2, 9.3.3)에서는 $X$에 대한 두 가정을, $X$가
\emph{준콤팩트이고 준분리}인 준스킴이라는 더 약한 가정으로 바꿀
수 있다. 실제로 그 증명은 바로 이 두 성질만을 사용한다 (1.2.7).
\end{env}

\begin{env}[1.7.6]
\label{IV.1.7.6-fr}
(I, 9.3.4, 9.3.5)에서는 $X$가 준콤팩트이고 $\mathcal F$와
$\mathcal J$가 준연접이면서 유한 생성이라고 가정하면 충분하다.
\end{env}

\begin{env}[1.7.7]
\label{IV.1.7.7-fr}
(I, 9.4.7)에서는 $X$가 \emph{준분리}라고만 가정하여 조건 b)를
약화할 수 있다. 실제로 표준 몰입 사상 $U\to X$가 준콤팩트이기만
하면 충분하다 (1.2.7). (I, 9.4.8, 9.4.9, 9.4.10)에서도 같은 변경을
할 수 있다.
\end{env}

\begin{env}[1.7.8]
\label{IV.1.7.8-fr}
(I, 9.5.1)에서는 명제의 괄호 안에 적힌 조건들을 $f$가 준콤팩트이고
준분리라는 조건으로 바꿀 수 있다.
\end{env}

\begin{env}[1.7.9]
\label{IV.1.7.9-fr}
(I, 9.6.5)에서는 $X$에 대한 가정들을 $X$가 준콤팩트이고 준분리라는
더 약한 가정으로 바꿀 수 있다. 실제로 그 증명은 (I, 9.4.7)만을
사용한다.

\oldpage[IV-1]{237}

(I, 9.6.6)에서는 ``$X$가 준콤팩트 스킴이다''라는 가정을 ``$X$가
준콤팩트이고 준분리인 준스킴이다''로 바꿀 수 있다.
\end{env}

\begin{env}[1.7.10]
\label{IV.1.7.10-fr}
(II, 1.7.15)에서는 $Y$가 준콤팩트이고 준분리라고만 가정할 수 있다.
실제로 그 증명은 (I, 9.6.5)만을 사용한다.
\end{env}

\begin{env}[1.7.11]
\label{IV.1.7.11-fr}
(II, 3.4.5)의 두 번째 주장에서는 $Y$에 대한 두 가정을, $Y$가
준콤팩트이고 준분리라는 더 약한 가정으로 바꿀 수 있다. (II,
3.4.8)에서도 마찬가지이다.
\end{env}

\begin{env}[1.7.12]
\label{IV.1.7.12-fr}
(II, 3.8.5)에서도 마찬가지로 $Y$가 준콤팩트이고 준분리라고만
가정할 수 있다. 이 가정은 (I, 9.4.7)을 통해 쓰인다.
\end{env}

\begin{env}[1.7.13]
\label{IV.1.7.13-fr}
(II, 4.4.1, 4.4.6, 4.4.7)에서는 $Y$가 \emph{준콤팩트이고
준분리}라고 가정하면 충분하다. (II, 4.4.7)의 증명에서는 (1.2.3)도
고려한다. 마찬가지로 (II, 4.4.10.1)에서는 $Z$가 준콤팩트이고
준분리라고 가정하면 충분하다.
\end{env}

\begin{env}[1.7.14]
\label{IV.1.7.14-fr}
(II, 4.5.2)와 (II, 4.5.5)에서는 $X$가 \emph{준콤팩트이고
준분리}라고 가정하면 충분하다. 이 가정은 (I, 9.3.1, 9.3.2)를 통해
쓰인다.
\end{env}

\begin{env}[1.7.15]
\label{IV.1.7.15-fr}
(II, 4.6.8)에서도 $X$가 \emph{준콤팩트이고 준분리}라고 가정하면
충분하다. 이 가정은 (I, 9.4.7)을 통해 쓰인다.
\end{env}

\begin{env}[1.7.16]
\label{IV.1.7.16-fr}
(II, 5.1.2)에서는 $X$가 준콤팩트이고 준분리라고 가정하면 충분하다.
이 가정은 (II, 4.5.2, 4.5.5)를 통해 쓰인다. (I, 9.6.5)를 사용하는
(II, 5.1.9)에서도 마찬가지이다.
\end{env}

\begin{env}[1.7.17]
\label{IV.1.7.17-fr}
(II, 5.2.1)에서도 여전히 $X$가 준콤팩트이고 준분리라고 가정하면
충분하다. 이 가정은 (II, 4.5.2)를 통해 쓰인다.
\end{env}

\begin{env}[1.7.18]
\label{IV.1.7.18-fr}
(II, 5.2.2)에서는 $f$가 준콤팩트이고 $X$가 준분리라고 가정해야
한다. 실제로 현재의 가정만으로는 기존 증명이 충분하지 않다. 모든
준연접 $\mathcal O_X$-가군 $\mathcal F$에 대하여
$\mathrm R^1f_*(\mathcal F)=0$이더라도, $Y$의 아핀 열린집합 $U$에
대하여 같은 성질이 준연접 $(\mathcal O_X|f^{-1}(U))$-가군들에
성립한다고 반드시 결론지을 수는 없다. 그러려면 그런 가군이 준연접
$\mathcal O_X$-가군의 \emph{제한}임을 알아야 한다. 같은 지적은 조건
b)와 c')에도 적용된다. $f$가 준콤팩트이고 $X$가 준분리이면, 위에서
수정한 (I, 9.4.7), 즉 (1.7.7)에 의해 이러한 연장이 가능하다.
\end{env}

\begin{env}[1.7.19]
\label{IV.1.7.19-fr}
(II, 5.3.2)에서는 $Y$가 \emph{준콤팩트이고 준분리}라고 가정하면
충분하다. 이는 (II, 4.4.6, 4.4.7)을 통해 쓰인다. (II, 5.5.3,
(ii))에서도 마찬가지이다.
\end{env}

\begin{env}[1.7.20]
\label{IV.1.7.20-fr}
(III, 1.2.2, 1.2.3, 1.2.4)에서는 $X$에 대한 두 가정을, $X$가
\emph{준콤팩트이고 준분리}인 준스킴이라는 더 약한 가정으로 바꿀
수 있다. 실제로 그 증명은 (I, 9.3.3)만을 사용한다.
\end{env}

\begin{env}[1.7.21]
\label{IV.1.7.21-fr}
(III, 1.4.10에서 1.4.14까지)에서는 ``분리''를 ``준분리''로 바꿀
수 있다. 이 가정은 오직 (I, 9.2.2)를 적용하기 위해서만 쓰이기
때문이다 (위의 (1.7.4)를 보라). 마찬가지로 (III, 1.4.15)에서는
사상 $u$가 \emph{준분리이고 준콤팩트}라고 가정하면 충분하다.
\end{env}

\begin{env}[1.7.22]
\label{IV.1.7.22-fr}
(III, 2.1.3)에서도 가정들을 $X$가 준콤팩트이고 준분리인 준스킴이라는
더 약한 가정으로 바꿀 수 있다. 논법은 (III, 1.2.2)에서와 같다.
\end{env}

\subsection{유한 표시 사상과 구성 가능 집합}
\label{subsection:IV.1.8-fr}

\oldpage[IV-1]{238}

\begin{env}[1.8.1]
\label{IV.1.8.1-fr}
$X$를 한 준스킴이라 하자. $X$의 모든 준콤팩트 열린집합은 유한 개의
아핀 열린집합의 합집합이고, 그 역도 성립한다. 따라서 $X$의 열린집합
$U$가 $X$ 안에서 \emph{역콤팩트}라는 것은
$(0\textsubscript{III}, 9.1.1)$ $X$의 모든 \emph{아핀 열린집합}
$V$에 대하여 $U\cap V$가 준콤팩트라는 것과 동치이다. $X$가
\emph{준분리}이면
(1.2.1), $X$의 모든 준콤팩트 열린집합은 $X$ 안에서 역콤팩트이고
(1.2.7), 따라서 $X$의 모든 국소 구성 가능 부분집합은 $X$ 안에서
역콤팩트이다 $(0\textsubscript{III}, 9.1.13)$. \emph{더욱이} $X$가
\emph{준콤팩트}이면, $X$의 구성 가능 부분집합과 국소 구성 가능
부분집합은 서로 같고 $(0\textsubscript{III}, 9.1.12)$, 구성 가능한
열린 부분집합과 준콤팩트 열린 부분집합도 서로 같다
$(0\textsubscript{III}, 9.1.5)$.
\end{env}

\begin{proposition}[1.8.2]
\label{IV.1.8.2-fr}
$X$, $Y$를 두 준스킴, $f:X\to Y$를 한 사상이라 하자. $Y$의 모든
구성 가능 부분집합 (각각 국소 구성 가능 부분집합) $Z$에 대하여
$f^{-1}(Z)$는 $X$의 구성 가능 부분집합 (각각 국소 구성 가능
부분집합)이다.
\end{proposition}

\begin{proof}
$Z$가 국소 구성 가능하다고 하자. 모든 $x\in X$에 대하여 $Z\cap V$가
$V$에서 구성 가능이 되는 $f(x)$의 아핀 열린근방 $V$가 존재한다.
구성 가능 부분집합에 대하여 명제가 증명되었다면
$f^{-1}(Z)\cap f^{-1}(V)$는 $f^{-1}(V)$에서 구성 가능이고, 따라서
$f^{-1}(Z)$는 국소 구성 가능이다. 그러므로 $Z$가 구성 가능한 경우만
생각하면 충분하다. $(0\textsubscript{III}, 9.1.3)$을 고려하면, $Z$가
$Y$ 안에서 열린 역콤팩트 부분집합인 경우, 다시 말해 표준 포함 사상
$Z\to Y$가 준콤팩트 사상인 경우로 환원된다. 이때 $f^{-1}(Z)$가
$X$ 안에서 역콤팩트임, 다시 말해 표준 포함 사상
$f^{-1}(Z)\to X$가 준콤팩트임을 보이면 충분하다. 그런데
$f^{-1}(Z)=Z\times_YX$이므로 (I, 4.4.1), 이는 (1.1.2, (iii))에서
따른다.
\end{proof}

\begin{lemma}[1.8.3]
\label{IV.1.8.3-fr}
$X$를 준콤팩트이고 준분리인 준스킴, $Z$를 $X$의 구성 가능
부분집합이라 하자. 그러면 아핀 스킴 $X'$과 유한 표시 사상
$f:X'\to X$가 존재하여 $f(X')=Z$이다.
\end{lemma}

\begin{proof}
$X$가 준콤팩트이므로 유한 개의 아핀 열린집합 $X_i$ ($i\in I$)로
덮인다. $X$가 준분리이므로 (1.2.7)에 의해 각 표준 몰입 사상
$g_i:X_i\to X$는 유한 표시이다. 따라서 $f_i:X'_i\to X_i$가 유한
표시이면 $g_i\circ f_i:X'_i\to X$도 유한 표시이다 (1.6.2). 또
$X'$을 $X'_i$들의 합인 준스킴이라 하면, 각 $X'_i$에서
$g_i\circ f_i$와 일치하는 사상 $h:X'\to X$도 유한 표시이다
(1.6.5). $Z\cap X_i$는 $X_i$에서 구성 가능이므로
$(0\textsubscript{III}, 9.1.8)$, 보조정리를 $X_i$가 환 $A$의
\emph{아핀} 스킴인 경우에 증명하는 것으로 환원된다.
$Z$는 $U$, $V$가 준콤팩트 열린집합인 $U\cap\complement V$ 꼴의
집합들의 유한 합집합이므로 $(0\textsubscript{III}, 9.1.3)$, 다시
적당한 준스킴의 합을 생각하면 $Z=U\cap\complement V$인 경우로
환원된다. 또 $U$는 유한 개의 아핀 열린집합의 합집합이므로 $U$가
\emph{아핀}인 경우만 생각해도 된다. $V$가 준콤팩트이므로
$f_i\in A$인 유한 개의 열린집합 $X_{f_i}$의 합집합이다.
$\mathfrak J$를 $f_i$들이 생성하는 $A$의 아이디얼, $Z'$을
$\mathfrak J$가 정의하는 $X$의 아핀 닫힌 부분스킴이라 하자. $Z'$은
구성상 $X$ 위에서 유한 표시이다. 정의에 의해 $V=\complement Z'$이고,
따라서 $U\cap\complement V=U\cap Z'$이다. 아핀 스킴
$X'=Z'\times_XU$를 생각하자.

\oldpage[IV-1]{239}

$f:X'\to X$를 구조사상이라 하자. 방금 본 대로 표준 몰입 사상
$Z'\to X$는 유한 표시이다. 열린 몰입 사상 $U\to X$도 유한
표시인데, 이는 $U$와 $X$가 아핀이므로 이 사상이 준콤팩트이기
때문이다. 따라서 (1.6.2, (iv))에 의해 $f$는 유한 표시이고,
$f(X')=Z'\cap U$이다 (I, 3.4.8). 이로써 증명이 끝난다.
\end{proof}

\begin{theorem}[1.8.4]
\label{IV.1.8.4-fr}
\textup{(슈발레).} --- $f:X\to Y$를 유한 표시 사상이라 하자. $X$의
모든 국소 구성 가능 부분집합 $Z$에 대하여 $f(Z)$는 $Y$에서 국소
구성 가능이다.
\end{theorem}

\begin{proof}
$y\in Y$이고 $V$가 $y$의 아핀 열린근방이라고 하자. $f$가
준콤팩트이고 준분리이므로 그 제한 $f^{-1}(V)\to V$도 그러하고,
따라서 $f^{-1}(V)$는 준콤팩트이고 준분리인 준스킴이다 (1.2.3).
$Z\cap f^{-1}(V)$는 $f^{-1}(V)$에서 구성 가능이므로 (1.8.1), $Y$가
아핀이고 $Z$가 구성 가능한 경우만 생각하면 충분하다. 이때 $X$는
준콤팩트이고 준분리이므로 (1.8.3)에 의해 $Z=g(X')$가 되는 유한 표시
사상 $g:X'\to X$가 존재한다. $f\circ g$가 유한 표시이므로 $Z=X$인
경우로 환원된다. 다시 말해 다음 보조정리를 증명하면 충분하다.
\end{proof}

\begin{lemma}[1.8.4.1]
\label{IV.1.8.4.1-fr}
$Y$를 아핀 스킴, $f:X\to Y$를 준콤팩트이고 국소 유한 표시인
사상이라 하자. 그러면 $f(X)$는 $Y$의 구성 가능 부분집합이다.
\end{lemma}

\begin{proof}
$X$가 준콤팩트이므로 유한 개의 아핀 열린집합으로 덮인다. 따라서
$Y=\operatorname{Spec}(A)$, $X=\operatorname{Spec}(B)$이고 $B$가
유한 표시 $A$-대수인 경우로 환원된다 (1.4.6). 이제 다음 보조정리를
사용한다.
\end{proof}

\begin{lemma}[1.8.4.2]
\label{IV.1.8.4.2-fr}
$A_0$를 한 환, $(A_\lambda)_{\lambda\in L}$를 $A_0$-대수들의
귀납계, $A=\varinjlim A_\lambda$라 하자. $B$가 유한 표시
$A$-대수이면, 어떤 $\lambda$와 유한 표시 $A_\lambda$-대수
$B_\lambda$가 존재하여
$B\simeq B_\lambda\otimes_{A_\lambda}A$이다.
\end{lemma}

\begin{proof}
가정에 의해 $B$는
$A[T_1,\ldots,T_n]/\mathfrak J$ 꼴의 대수와 동형이다. 여기서
$T_i$들은 부정원이고 $\mathfrak J$는 유한 생성 아이디얼이다.
$(F_j)$ ($1\leq j\leq m$)를 $\mathfrak J$의 한 생성계라 하고,
$\varphi_\lambda:A_\lambda\to A$를 표준 준동형이라 하자. $L$이
여과되어 있으므로, 각 다항식 $F_j$의 모든 계수가
$\varphi_\lambda$에 의한 $A_\lambda$의 원소의 상이 되는 $\lambda$가
존재한다. 다시 말해
$F_j=\varphi_\lambda(F_{j\lambda})$ ($1\leq j\leq m$)를 만족하는
$A_\lambda[T_1,\ldots,T_n]$의 다항식계
$(F_{j\lambda})_{1\leq j\leq m}$가 존재한다. $\mathfrak J_\lambda$를
$F_{j\lambda}$들이 생성하는 $A_\lambda[T_1,\ldots,T_n]$의
아이디얼이라 하자. 그러면 $\mathfrak J$는
\[
\mathfrak J_\lambda\otimes_{A_\lambda}A
\]
의 $A[T_1,\ldots,T_n]
=A_\lambda[T_1,\ldots,T_n]\otimes_{A_\lambda}A$ 안에서의 상이다.
따라서 환
$B_\lambda=A_\lambda[T_1,\ldots,T_n]/\mathfrak J_\lambda$가
요구를 만족한다.
\end{proof}

여기서는 $A$를 그 유한 생성 부분 $\mathbf Z$-대수들의 귀납극한으로
보고 이 보조정리를 적용한다. 따라서 $A$의 그러한 부분
$\mathbf Z$-대수 $A_0$와 유한 표시 $A_0$-대수 $B_0$가 존재하여
$B\simeq B_0\otimes_{A_0}A$이다. $X_0=\operatorname{Spec}(B_0)$,
$Y_0=\operatorname{Spec}(A_0)$라 두면
$X=X_0\times_{Y_0}Y$이고 $f$는 사영 $X\to Y$이다. 구조사상들을
$f_0:X_0\to Y_0$, $g_0:Y\to Y_0$라 하자. (I, 3.4.8)에 의해
\[
  f(X)=g_0^{-1}(f_0(X_0)).
\]
따라서 (1.8.2)를 고려하면 $f_0(X_0)$가 구성 가능임을 보이면
충분하다. 다시 말해 끝으로 다음 따름정리를 증명하는 것으로
환원된다.

\begin{corollary}[1.8.5]
\label{IV.1.8.5-fr}
$Y$를 뇌터 준스킴, $f:X\to Y$를 유한형 사상이라 하자. $X$의 모든
구성 가능 부분집합 $Z$에 대하여 $f(Z)$는 $Y$의 구성 가능
부분집합이다.
\end{corollary}

\oldpage[IV-1]{240}

\begin{proof}
위와 같이 $f(X)$가 구성 가능임을 증명하는 것으로 환원된다. 판정법
$(0\textsubscript{III}, 9.2.3)$을 적용하자. $Y$의 모든 기약 닫힌
부분집합 $T$에 대하여
\[
  T\cap f(X)=f(f^{-1}(T))
\]
가 $T$에서 어디에서도 조밀하지 않거나 $T$의 공집합이 아닌
열린집합을 포함함을 증명해야 한다. 바탕공간이 $T$인 $Y$의
부분준스킴을 택하고 (I, 5.2.1) 그 $f$에 의한 역상을 생각하면,
(1.5.4, (iii))을 고려하여 결국 $Y$가 \emph{기약이고 $f$가
우세}라고 가정할 때 $f(X)$가 $Y$의 공집합이 아닌 열린집합을
포함함을 증명하는 것으로 환원된다. 더 나아가 $Y$가 아핀이라고
가정할 수 있다. 이때 $X$는 유한 개의 아핀 열린집합 $V_i$의
합집합이고, $f(X)$가 $Y$에서 조밀하므로 $f(V_i)$ 가운데 적어도
하나는 $Y$에서 조밀하다. 따라서 $X$도 아핀이라고 가정할 수 있다.
$(X_j)$를 $X$의 유한한 기약 성분족이라 하면, 위와 같이
$f(X_j)$ 가운데 적어도 하나는 $Y$에서 조밀하다. 따라서 $X$도
기약이라고 가정할 수 있다. 끝으로 $f$를 $f_{\mathrm{red}}$로
바꾸어 (1.5.4, (vi)) $X$와 $Y$가 \emph{정역 준스킴}이라고 가정할
수 있다. 그러면 $X=\operatorname{Spec}(B)$,
$Y=\operatorname{Spec}(A)$이고, $A$는 뇌터 정역이며 $B$는 $A$를
포함하는 유한 생성 $A$-대수인 정역이다 (I, 1.2.7). 어떤 $g\in A$가
존재하여 모든 $y\in D(g)$, 즉 $g\notin\mathfrak p_y$인 모든 $y$에
대하여 $\mathfrak p_y$가 $B$의 어떤 소 아이디얼과 $A$의 교집합이
됨을 보이면 충분하다. 그러면 $D(g)\subset f(X)$이기 때문이다.
끝으로 다음 보조정리를 증명하면 충분하다.
\end{proof}

\begin{lemma}[1.8.5.1]
\label{IV.1.8.5.1-fr}
$A$를 정역, $B$를 유한 생성 $A$-대수인 정역이라 하자. 어떤
$g\in A$가 존재하여, $g$에서 0이 아닌 모든 준동형
$A\to\Omega$ ($\Omega$는 대수적으로 닫힌 체)는 $B\to\Omega$인
준동형으로 연장된다.
\end{lemma}

\begin{proof}
이는 가환대수학의 고전적 결과이다 (Bourbaki, \emph{Alg. comm.},
chap. V, \S~3, n\textsuperscript{o}~1, th. 1의 cor. 3).
\end{proof}

\begin{corollary}[1.8.6]
\label{IV.1.8.6-fr}
$Y$를 기약 준스킴, $\eta$를 그 일반점, $f:X\to Y$를 국소 유한형
사상이라 하자. $f^{-1}(\eta)$가 공집합이 아니면, 모든 $y\in U$에
대하여 $f^{-1}(y)$가 공집합이 아닌 $Y$에서의 $\eta$의 열린근방
$U$가 존재한다. $\mathscr F$가 유한 생성 준연접
$\mathcal O_X$-가군이고
$\mathscr F_\eta=\mathscr F\otimes_{\mathcal O_Y}k(\eta)$가 0이
아니면, 모든 $y\in U'$에 대하여
$\mathscr F_y=\mathscr F\otimes_{\mathcal O_Y}k(y)$가 0이 아닌
$Y$에서의 $\eta$의 열린근방 $U'$가 존재한다.
\end{corollary}

\begin{proof}
$p_y$를 올
$f^{-1}(y)=X\times_Y\operatorname{Spec}(k(y))$에서 $X$로 가는 표준
사영이라 하면 $\mathscr F_y=p_y^*(\mathscr F)$이고, 따라서
\[
\operatorname{Supp}(\mathscr F_y)
=p_y^{-1}(\operatorname{Supp}(\mathscr F))
=\operatorname{Supp}(\mathscr F)\cap f^{-1}(y)
\]
이다. 이는 $\mathscr F$가 유한 생성이므로 (I, 9.1.13.1)과
(I, 3.6.1)에서 따른다. 또한 $\operatorname{Supp}(\mathscr F)$는
$X$에서 닫혀 있다 $(0\textsubscript{I}, 5.2.2)$. 바탕공간이
$\operatorname{Supp}(\mathscr F)$인 $X$의 닫힌 부분준스킴을 $Z$라
하고 (I, 5.2.1), $j:Z\to X$를 표준 몰입 사상이라 하면 $f\circ j$는
국소 유한형이다 (1.3.4). 이는 첫 번째 주장에서 두 번째 주장이
따름을 증명한다.
첫 번째 주장을 증명하기 위해, 먼저 $X$와 $Y$가 축소 준스킴이라고
가정할 수 있음에 주의하자 (1.3.4, (vi)). 특히 $Y$는 기약이므로
정역 준스킴이라고 가정할 수 있다. $\xi$를 $f^{-1}(\eta)$의 한 점이라
하자. $X$와 $Y$를 각각 $\xi$와 $\eta$의 아핀 열린근방으로 바꿀 수
있으므로
\[
  X=\operatorname{Spec}(B),\qquad Y=\operatorname{Spec}(A)
\]
이고 $B$가 유한 생성 $A$-대수라고 가정할 수 있다. 더 나아가 $Z$를
바탕집합이 $\overline{\{\xi\}}$인 $X$의 축소 닫힌 부분준스킴이라
하면 (I, 5.2.1), 위와 같이 $X$를 $Z$로 바꿀 수 있다. 달리 말해
$B$가 \emph{정역}이라고 가정할 수 있다 (I, 5.1.4). 이때 사상 $f$가
우세이므로 대응하는 준동형 $\varphi:A\to B$는 단사이다 (I, 1.2.7).
따라서 따름정리는 끝으로 보조정리 (1.8.5.1)에서 따른다.
\end{proof}

\oldpage[IV-1]{241}

\begin{proposition}[1.8.7]
\label{IV.1.8.7-fr}
$S$를 준스킴, $g:X\to S$, $h:Y\to S$를 두 유한 표시 사상,
$f:X\to Y$를 $S$-사상이라 하자. 각 $s\in S$에 대하여
\[
  X_s=g^{-1}(s)=X\times_S\operatorname{Spec}(k(s)),\qquad
  Y_s=h^{-1}(s)=Y\times_S\operatorname{Spec}(k(s)),
\]
$f_s=f\times 1:X_s\to Y_s$라 두자. 그러면 $f_s$가 전사인 (각각
라디시엘 사상인) $s\in S$들의 집합은 국소 구성 가능이다.
\end{proposition}

\begin{proof}
$E$를 $f_s$가 전사인 $s\in S$들의 집합이라 하자. 집합 $S-E$는
$h(Y-f(X))$와 같다. 그런데 $f$는 유한 표시 사상이므로 (1.6.2, (v))
$f(X)$는 $Y$에서 국소 구성 가능이고 (1.8.4), $h$도 유한 표시
사상이므로 $h(Y-f(X))$는 $S$에서 국소 구성 가능이다. 따라서 $E$도
국소 구성 가능이다.

$f_s$가 라디시엘 사상인 $s\in S$들의 집합 $E'$이 국소 구성 가능임을
보이기 위해 다음 보조정리를 사용한다.
\end{proof}

\begin{lemma}[1.8.7.1]
\label{IV.1.8.7.1-fr}
$U$, $V$를 두 준스킴이라 하자. 사상 $f:U\to V$가 라디시엘
사상일 필요충분조건은 대각사상
$\Delta_f:U\to U\times_VU$가 전사인 것이다. 따라서 모든 라디시엘
사상은 분리 사상이다.
\end{lemma}

\begin{proof}
실제로 $\Delta_f$는 몰입 사상이므로 (I, 5.3.9) 국소 유한형 사상이다
(1.3.4). 따라서 $\Delta_f$가 전사라는 것은 모든 대수적으로 닫힌 체
$K$에 대하여 대응하는 사상
\[
  U(K)\longrightarrow (U\times_VU)(K)
  =U(K)\times_{V(K)}U(K)
\]
이 전사라는 것과 같다 (1.3.7 및 I, 3.4.2.1). 그러나 집합의
섬유곱의 정의에 의해 이는 $f$에 대응하는 사상 $U(K)\to V(K)$가
단사라는 뜻이고, 이는 $f$가 라디시엘 사상이라는 것과 같다
(I, 3.5.5).

이제 $\Delta_{f_s}$는 대각사상
$\Delta_f:X\to X\times_YX$에 밑변환
$\operatorname{Spec}(k(s))\to S$를 하여 얻는다 (I, 5.3.4). $f$가
유한 표시 사상이므로 구조사상 $X\times_YX\to Y$도 유한 표시
사상이고 (1.6.2, (iv)), 따라서 $\Delta_f$도 유한 표시 사상이다
(1.6.2, (v)). 그러므로 보조정리 (1.8.7.1)을 사용하여 이 명제의
첫 번째 부분을 $\Delta_f$에 적용하면 충분하다.
\end{proof}

\subsection{프로-구성 가능 집합과 인드-구성 가능 집합}
\label{subsection:IV.1.9-fr}

\begin{lemma}[1.9.1]
\label{IV.1.9.1-fr}
$(A_\alpha)_{\alpha\in I}$를 지표집합 $I$가 오른쪽 여과 집합인 환의
귀납계라 하고 $A=\varinjlim A_\alpha$라 하자. $A=0$일 필요충분조건은
어떤 지표 $\gamma$에 대하여 $A_\gamma=0$인 것이다. 이때 모든
$\beta\geq\gamma$에 대하여 $A_\beta=0$이다.
\end{lemma}

\begin{proof}
실제로 각 $\alpha\in I$에 대하여 표준 준동형
$\varphi_\alpha:A_\alpha\to A$는 단위원을 단위원으로 보낸다.
$A=0$이라는 것은 $\varphi_\alpha(1)=0$, 즉
$\varphi_\alpha(1)=\varphi_\alpha(0)$이라는 뜻이다. 이는 어떤
$\gamma\geq\alpha$가 존재하여 준동형
$\varphi_{\gamma\alpha}:A_\alpha\to A_\gamma$가
$\varphi_{\gamma\alpha}(1)=\varphi_{\gamma\alpha}(0)$을 만족한다는
것과 같다고 알려져 있다. 마지막 관계는 $A_\gamma=0$과 같으므로
보조정리를 얻는다.
\end{proof}

\begin{proposition}[1.9.2]
\label{IV.1.9.2-fr}
$B$를 환, $(A_\alpha)_{\alpha\in I}$를 지표집합 $I$가 오른쪽 여과
집합인 $B$-대수의 귀납계라 하고 $A=\varinjlim A_\alpha$라 하자.
\[
  Y=\operatorname{Spec}(B),\qquad
  X_\alpha=\operatorname{Spec}(A_\alpha),\qquad
  X=\operatorname{Spec}(A)
\]
라 두고, $u:X\to Y$, $u_\alpha:X_\alpha\to Y$를 각각 구조 준동형
$B\to A$, $B\to A_\alpha$에 대응하는 사상이라 하자. 그러면 다음이
성립한다.
\end{proposition}

\begin{env}[1.9.2.i]
\label{IV.1.9.2.i-fr}
\emph{$X=\varnothing$일 필요충분조건은 어떤 $\gamma\in I$에 대하여
$X_\gamma=\varnothing$인 것이다. 이때 모든 $\alpha\geq\gamma$에
대하여 $X_\alpha=\varnothing$이다.}
\end{env}

\oldpage[IV-1]{242}

\begin{env}[1.9.2.ii]
\label{IV.1.9.2.ii-fr}
\emph{다음 등식이 성립한다.}
\[
  u(X)=\bigcap_{\alpha\in I}u_\alpha(X_\alpha). \tag{1.9.2.1}
  \label{IV.1.9.2.1-fr}
\]
\end{env}

주장 (i)는 (1.9.1)을 옮겨 쓴 것에 지나지 않는다. (ii)를 증명하자.
$u$가
\[
  X\longrightarrow X_\alpha\xrightarrow{u_\alpha}Y
\]
로 분해되므로 (1.9.2.1)의 왼쪽은 오른쪽에 포함된다. 역으로
$y\in Y-u(X)$라 하자. 그러면 $u^{-1}(y)=\varnothing$이고,
$u^{-1}(y)$는 $k(y)$-준스킴
$\operatorname{Spec}(A\otimes_Bk(y))$의 바탕공간이다. $B$-가군의
범주에서 함자 $\varinjlim$은 텐서곱과 가환하므로
$A\otimes_Bk(y)$는 환의 귀납계
$A_\alpha\otimes_Bk(y)$의 귀납극한이다. 보조정리 (1.9.1)에 의해
$A_\alpha\otimes_Bk(y)=0$인 $\alpha$가 존재한다. 이는
$u_\alpha^{-1}(y)=\varnothing$이라는 뜻이므로
$y\notin u_\alpha(X_\alpha)$이다. 증명 끝.

\begin{proposition}[1.9.3]
\label{IV.1.9.3-fr}
$X$를 준콤팩트이고 준분리인 준스킴, $E$를 $X$의 부분집합,
$(X_\alpha)_{\alpha\in I}$를 $X$의 아핀 열린덮개라 하자. 각
$\alpha\in I$에 대하여 $E_\alpha=E\cap X_\alpha$라 두면 다음 조건들은
서로 동치이다.
\begin{enumerate}
  \item[a)] 어떤 준콤팩트 사상 $f:X'\to X$가 존재하여
  $E=f(X')$이다.
  \item[a$'$)] 각 $\alpha$에 대하여 어떤 준콤팩트 사상
  $f_\alpha:X'_\alpha\to X_\alpha$가 존재하여
  $E_\alpha=f_\alpha(X'_\alpha)$이다.
  \item[a$''$)] 각 $\alpha$에 대하여 어떤 아핀 스킴 $X''_\alpha$와
  사상 $f'_\alpha:X''_\alpha\to X_\alpha$가 존재하여
  $f'_\alpha(X''_\alpha)=E_\alpha$이다.
  \item[b)] $E$는 $X$의 구성 가능 부분집합들의 교집합이다.
  \item[b$'$)] $F=X-E$는 $X$의 구성 가능 부분집합들의 합집합이다.
\end{enumerate}
\end{proposition}

\begin{proof}
$b)$와 $b')$가 동치임은 명백하다. 조건 $a)$는 $X'$의 열린집합
$f^{-1}(X_\alpha)$ 위의 유도된 준스킴을 $X'_\alpha$로, $f$의 제한을
$f_\alpha$로 택하면 $a')$를 함의한다. $a')$가 $a'')$를 함의함을
보자. $f_\alpha$가 준콤팩트이고 $X_\alpha$가 아핀이므로
$X'_\alpha$는 준콤팩트이다. 따라서 유한 개의 아핀 열린집합
$X''_{\alpha\lambda}$ ($\lambda\in L$)로 덮을 수 있다. 이들의 합
준스킴을 $X''_\alpha$라 하면 이는 아핀 스킴이다. 각
$X''_{\alpha\lambda}$에서 항등사상과 같은 사상을
$g_\alpha:X''_\alpha\to X'_\alpha$라 하자. 그러면 $g_\alpha$가
전사이므로 $f'_\alpha=f_\alpha\circ g_\alpha$라 둘 때
$E_\alpha=f'_\alpha(X''_\alpha)$임이 명백하다.

$a'')$가 $a)$를 함의함을 보자. $\alpha\in J$인 $X_\alpha$들이
$X$를 덮는 $I$의 유한 부분집합 $J$가 존재한다. $X'$을
$\alpha\in J$인 $X''_\alpha$들의 합 준스킴이라 하고,
$j_\alpha:X_\alpha\to X$를 표준 포함 사상이라 하자. 각
$\alpha\in J$에서 $j_\alpha\circ f'_\alpha$와 같은 사상을
$f:X'\to X$라 하자. $E$는 $\alpha\in J$인
$E\cap X_\alpha$들의 합집합이므로 실제로 $E=f(X')$이다. 이제
$f$가 준콤팩트 사상임만 보이면 된다. $X$가 준분리라는 가정에 의해
$j_\alpha$는 준콤팩트이고 (1.2.7), 따라서
$j_\alpha\circ f'_\alpha$도 준콤팩트이며 (1.1.1 및 1.1.2), 끝으로
$f$도 준콤팩트이다 (1.1.6).

따라서 $a'')$와 $b)$의 동치만 증명하면 된다. 먼저 $a'')$가 $b)$를
함의함을 보이자. $\alpha\in J$인 $X_\alpha$들이 $X$를 덮는 $I$의
유한 부분집합 $J$를 택하자. 각 $\alpha\in J$에 대하여 $E_\alpha$가
$X_\alpha$의 구성 가능 부분집합 $F_{\alpha\lambda}$
($\lambda\in L_\alpha$)들의 교집합임을 보이면 충분하다. 실제로
$X$가 준분리라는 가정에 의해 각 $X_\alpha$는 $X$ 안에서
역콤팩트이고 (1.2.7), 따라서 각 $F_{\alpha\lambda}$는
$(0\textsubscript{III}, 9.1.8, (ii))에 의해 $X$에서도 구성 가능이다.
$E$는 $\alpha\in J$인 $E_\alpha$들의 합집합이므로, 모든
$\lambda(\alpha)\in L_\alpha$의 선택에 대한 유한 합집합
\[
  \bigcup_{\alpha\in J}F_{\alpha,\lambda(\alpha)}
\]
들의 교집합이다. 이 합집합들은 $X$에서 구성 가능이므로 주장이
따른다. 따라서 다음 경우로 환원할 수 있다.

\oldpage[IV-1]{243}

$X=\operatorname{Spec}(A)$이고 $E=f(X')$라 하자. 여기서
$X'=\operatorname{Spec}(A')$는 아핀이고 $f:X'\to X$는 환 준동형
$A\to A'$에 대응하는 사상이다. 이제 $A'$은 포함관계로 순서가 주어진
그 유한 생성 부분 $A$-대수들의 족 $(A'_\lambda)$의 귀납극한이다.
$X'_\lambda=\operatorname{Spec}(A'_\lambda)$라 하면 $f$는
\[
  X'\longrightarrow X'_\lambda\xrightarrow{f_\lambda}X
\]
로 분해되고, (1.9.2, (ii))에 의해
\[
  f(X')=\bigcap_\lambda f_\lambda(X'_\lambda)
\]
이다. 각 $f_\lambda(X'_\lambda)$가 $X$의 구성 가능 부분집합들의
교집합임을 보이면 $f(X')$도 그러하다. 따라서 $A'$이 유한 생성
$A$-대수인 경우로 환원할 수 있다. 이제 다음 보조정리를 사용한다.
\end{proof}

\begin{lemma}[1.9.3.1]
\label{IV.1.9.3.1-fr}
$A$를 환, $A'$을 유한 생성 $A$-대수라 하자. 그러면 $A'$은 유한 표시
$A$-대수들 $A''_\mu$의 여과 귀납극한
$\varinjlim A''_\mu$와 $A$-동형이다 (1.4.1).
\end{lemma}

\begin{proof}
실제로 $B=A[T_1,\ldots,T_n]$이고 $\mathfrak J$가 $B$의 아이디얼일 때
$A'=B/\mathfrak J$로 쓸 수 있다. 그런데 $\mathfrak J$는
$\mathfrak J$에 포함되고 포함관계로 순서가 주어진 $B$의 유한 생성
아이디얼들 $\mathfrak J'_\mu$의 귀납극한이다. $B$-가군의 범주에서
함자 $\varinjlim$이 완전하므로, $A$-동형을 통해
\[
  A'=\varinjlim B/\mathfrak J'_\mu
\]
를 얻는다.
\end{proof}

위와 같은 논증으로 $A'$이 유한 표시 $A$-대수인 경우로 환원할 수
있다. 이때 Chevalley의 정리 (1.8.5)에 의해 $f(X')$는 $X$에서 구성
가능이므로 $b)$가 증명된다.

끝으로 $b)$가 $a'')$를 함의함을 보이자. $E$가 $X$의 구성 가능
부분집합족 $(G_\lambda)_{\lambda\in L}$의 교집합이면, 각 $E_\alpha$는
$X_\alpha$에서 구성 가능한 $G_\lambda\cap X_\alpha$들의 교집합이다
$(0\textsubscript{III}, 9.1.8, (i))$. 따라서
$X=\operatorname{Spec}(A)$가 아핀인 경우로 환원된다.

이때 (1.8.3)에 의해 각 $\lambda\in L$에 대하여
$X'_\lambda=\operatorname{Spec}(A'_\lambda)$가 아핀이고
$f_\lambda(X'_\lambda)=G_\lambda$인 사상
$f_\lambda:X'_\lambda\to X$가 존재한다. 따라서 다음 보조정리만
증명하면 충분하다.

\begin{lemma}[1.9.3.2]
\label{IV.1.9.3.2-fr}
$(A'_\lambda)_{\lambda\in L}$를 $A$-대수의 족,
$X=\operatorname{Spec}(A)$,
$X'_\lambda=\operatorname{Spec}(A'_\lambda)$라 하고,
$f_\lambda:X'_\lambda\to X$를 구조사상이라 하자. $L$의 모든 유한
부분집합 $J$에 대하여
\[
  A'_J=\bigotimes_{\lambda\in J}A'_\lambda
\]
를 $A$-대수의 텐서곱으로 정의하고,
$X'_J=\operatorname{Spec}(A'_J)$라 하며 $f_J:X'_J\to X$를
구조사상이라 하자. 그러면
\[
  f_J(X'_J)=\bigcap_{\lambda\in J}f_\lambda(X'_\lambda).
\]
유한 부분집합 $J$들의 여과 집합을 따라
$A'=\varinjlim A'_J$라 하고, $X'=\operatorname{Spec}(A')$,
$f:X'\to X$를 구조사상이라 하면
\[
  f(X')=\bigcap_{\lambda\in L}f_\lambda(X'_\lambda).
\]
\end{lemma}

\begin{proof}
첫 번째 주장은 (I, 3.4.7)에서 따르고, 두 번째 주장은
\[
  f(X')=\bigcap_Jf_J(X'_J)
\]
를 주는 (1.9.2, (ii))에서 따른다.
\end{proof}

\begin{definition}[1.9.4]
\label{IV.1.9.4-fr}
$X$를 위상공간이라 하자. $X$의 부분집합 $E$가 다음 조건을 만족하면
$X$에서 \emph{프로-구성 가능} (각각 \emph{인드-구성 가능})이라고
한다. 모든 $x\in X$에 대하여 $X$에서 $x$의 어떤 열린근방 $U$가
존재하여 $E\cap U$가 $U$의 국소 구성 가능 부분집합들의 교집합
(각각 합집합)이다.
\end{definition}

(1.2.7)과 $(0\textsubscript{III}, 9.1.8 및 9.1.12)$을 고려하면,
명제 (1.9.3)의 동치조건들은 $E$가 $X$에서 \emph{프로-구성 가능}이라고
말하는 것과 같다. 실제로 (1.9.3)의 가정 아래에서는 $X$의 국소 구성
가능 집합과 구성 가능 집합이 서로 같다.

\oldpage[IV-1]{244}

\begin{proposition}[1.9.5]
\label{IV.1.9.5-fr}
$X$를 준스킴이라 하자.
\begin{enumerate}
  \item[(i)] $X$의 부분집합 $E$가 프로-구성 가능일 필요충분조건은
  $X-E$가 인드-구성 가능인 것이다.
  \item[(ii)] $X$의 프로-구성 가능 부분집합들의 모든 유한 합집합과
  모든 교집합은 프로-구성 가능이다. 인드-구성 가능 부분집합들의
  모든 유한 교집합과 모든 합집합은 인드-구성 가능이다.
  \item[(iii)] $X$의 국소 구성 가능 부분집합들의 모든 교집합 (각각
  합집합)은 프로-구성 가능 (각각 인드-구성 가능)이다. 역으로 $X$가
  준콤팩트이고 준분리이면, $X$의 모든 프로-구성 가능 (각각
  인드-구성 가능) 부분집합은 구성 가능 부분집합들의 교집합 (각각
  합집합)이다.
  \item[(iv)] $(U_\alpha)$를 $X$의 열린덮개라 하자. $X$의 부분집합
  $E$가 $X$에서 프로-구성 가능 (각각 인드-구성 가능)일
  필요충분조건은 모든 $\alpha$에 대하여 $E\cap U_\alpha$가
  $U_\alpha$에서 프로-구성 가능 (각각 인드-구성 가능)인 것이다.
  \item[(v)] $X$의 모든 프로-구성 가능 부분집합 $E$는 $X$ 안에서
  역콤팩트이다. $X$의 국소 닫힌 부분집합 $E$가 $X$ 안에서
  역콤팩트일 필요충분조건은 $X$에서 프로-구성 가능인 것이다.
  \item[(vi)] $f:X'\to X$를 준스킴의 사상이라 하자. $X$의 모든
  프로-구성 가능 (각각 인드-구성 가능) 부분집합 $E$에 대하여
  $f^{-1}(E)$는 $X'$에서 프로-구성 가능 (각각 인드-구성 가능)이다.
  \item[(vii)] $f:X'\to X$를 준콤팩트 사상이라 하자. $X'$의 모든
  프로-구성 가능 부분집합 $E'$에 대하여 $f(E')$은 $X$에서
  프로-구성 가능이다. 특히 $f(X')$은 $X$에서 프로-구성 가능이다.
  \item[(viii)] $f:X'\to X$를 국소 유한 표시 사상이라 하자. $X'$의
  모든 인드-구성 가능 부분집합 $E'$에 대하여 $f(E')$은 $X$에서
  인드-구성 가능이다. 특히 $f(X')$은 $X$에서 인드-구성 가능이다.
  \item[(ix)] $X$가 준콤팩트이고 준분리라고 하자. $X$의 부분집합
  $E$가 프로-구성 가능일 필요충분조건은 어떤 아핀 스킴 $X'$과
  사상 $f:X'\to X$가 존재하여 $E=f(X')$인 것이다. 이 $f$는
  반드시 준콤팩트이다.
\end{enumerate}
\end{proposition}

\begin{proof}
(i), (ii), (iv) 및 (iii)의 첫 번째 주장은 정의 (1.9.4)에서 따르며,
임의의 위상공간에서 성립한다. 여기서는 교집합과 합집합의 상호
분배법칙 및 $T$가 $X$에서 국소 구성 가능이면 $X$의 모든 열린집합
$U$에 대하여 $T\cap U$가 $U$에서 국소 구성 가능이라는 사실을
사용한다.

$X$가 준콤팩트이고 준분리이며 $E$가 $X$에서 프로-구성 가능이라고
하자. $X$의 유한 아핀 열린덮개 $(U_i)$를 택하여 각 $E\cap U_i$가
$U_i$의 구성 가능 부분집합 $M_{i\lambda}$
($\lambda\in L_i$)들의 교집합이 되게 할 수 있다. (1.2.7)과
$(0\textsubscript{III}, 9.1.8, (ii))$에 의해 $M_{i\lambda}$들은
$X$에서도 구성 가능이고, $E$는 모든 $i$에 대한
$\lambda(i)\in L_i$의 각 선택에 대응하는 유한 합집합
\[
  \bigcup_i M_{i,\lambda(i)}
\]
들의 교집합이다. 이 유한 합집합들은 $X$의 구성 가능 부분집합이므로
(iii)의 두 번째 주장이 증명된다. (ix)는 (iii)과 (1.9.3)에서 따른다.

(v)의 첫 번째 주장을 증명하기 위해 $X$가 아핀인 경우로 환원하고,
$E$가 준콤팩트임을 보이면 충분하다 $(0\textsubscript{III}, 9.1.1)$.
이때 $X$는 준분리이므로 (ix)에 의해 $E=f(X')$인 준콤팩트 사상
$f:X'\to X$가 존재한다. $X'$가 준콤팩트이므로 연속사상에 의한 그
상 $E$도 준콤팩트이다.

(vi)을 증명할 때에는 (iv)에 의해 $X$가 아핀인 경우로 환원할 수
있고, 결론은 (iii)과 (1.8.2)에서 따른다.

\oldpage[IV-1]{245}

마찬가지로 (vii)을 증명할 때에도 $X$가 아핀인 경우로 환원할 수
있다. 그러면 $X'$은 유한 개의 아핀 열린집합 $X'_i$들의 합집합이고
$f(E')$은 $f(E'\cap X'_i)$들의 합집합이므로, (ii)에 의해 $X'$도
아핀이라고 가정할 수 있다. 이때 (ix)에 의해 어떤 준콤팩트 사상
$g:X''\to X'$에 대하여 $E'=g(X'')$이다. 따라서
$f(E')=f(g(X''))$이고, 합성사상 $f\circ g$가 준콤팩트이므로 $f(E')$은
프로-구성 가능이다.

(v)의 두 번째 주장도 (vii)에서 얻는다. 실제로 $E$를 $X$의 국소
닫힌 역콤팩트 부분집합이라 하고, 바탕공간이 $E$인 $X$의 부분준스킴을
생각하자 (I, 5.2.1). 표준 포함 사상 $j:E\to X$는 가정에 의해
준콤팩트이므로, (vii)에 의해 $E=j(E)$는 $X$에서 프로-구성 가능이다.

끝으로 (viii)을 증명할 때에도 $X$가 아핀이라고 가정할 수 있다.
또 $(X'_\alpha)$를 $X'$의 아핀 열린덮개로 택하여 각
$E'\cap X'_\alpha$가 $X'_\alpha$의 구성 가능 부분집합들의 합집합이
되게 하자 ((iii) 및 (iv)). $f(E')$은 $f(E'\cap X'_\alpha)$들의
합집합이고, Chevalley의 정리 (1.8.4)에 의해 각
$f(E'\cap X'_\alpha)$는 $X$의 구성 가능 부분집합들의 합집합이다.
따라서 결론을 얻는다.
\end{proof}

\begin{env}[1.9.6]
\label{IV.1.9.6-fr}
\textbf{예} (1.9.6). --- 준스킴 $X$의 모든 \emph{유한} 부분집합은
프로-구성 가능이다. 실제로 (1.9.5, (ii))에 의해 한 점으로 이루어진
부분집합 $\{x\}$만 생각하면 충분하다. $\{x\}$는 준콤팩트인 표준 사상
\[
  \operatorname{Spec}(k(x))\longrightarrow X
\]
에 의한 $\operatorname{Spec}(k(x))$의 상이므로, (1.9.5, (vii))에서
결론이 따른다. 반면 한 점으로 이루어진 부분집합 $\{x\}$가 반드시
인드-구성 가능한 것은 아니다. 예를 들어 $A$를 극대 아이디얼이
무한히 많은 뇌터 정역이라 하고 $x$를
$X=\operatorname{Spec}(A)$의 일반점이라 하자. $\{x\}$가 $X$에서
인드-구성 가능이라면 (1.9.5, (iii))에 의해 구성 가능이고, 따라서
$X$의 공집합이 아닌 열린집합을 포함한다
$(0\textsubscript{III}, 9.2.2)$. 그러나 Artin--Tate의 정리
$(0, 16.3.3)$에 의해 이는 $A$의 극대 아이디얼이 무한히 많다는
가정과 모순된다.

준스킴 $X$의 모든 \emph{닫힌} 부분집합 $Y$는 (1.9.5, (v))에 의해
프로-구성 가능이다. 따라서 $X$의 모든 열린 부분집합은 인드-구성
가능이지만, $X$의 열린 부분집합 $U$가 프로-구성 가능이면
(1.9.5, (v))에 의해 반드시 \emph{역콤팩트}이다.

끝으로 $X$가 \emph{뇌터} 준스킴이면 준분리이고 (1.2.8),
(1.9.5, (iii))과 $(0\textsubscript{III}, 9.1.7)$에 의해 $X$의
인드-구성 가능 부분집합들은 \emph{$X$의 국소 닫힌 부분집합들의
합집합}이다.
\end{env}

\begin{theorem}[1.9.7]
\label{IV.1.9.7-fr}
$X$를 준콤팩트 준스킴, $E$를 $X$의 인드-구성 가능 부분집합,
$(F_\lambda)_{\lambda\in L}$를
\[
  \bigcap_{\lambda\in L}F_\lambda\subset E
\]
를 만족하는 $X$의 프로-구성 가능 부분집합족이라 하자. 그러면
\[
  \bigcap_{\lambda\in J}F_\lambda\subset E
\]
를 만족하는 $L$의 유한 부분집합 $J$가 존재한다.
\end{theorem}

\begin{proof}
집합 $F=X-E$와 $F\cap F_\lambda$들은 프로-구성 가능이므로 다음
따름정리로 환원된다.
\end{proof}

\begin{corollary}[1.9.8]
\label{IV.1.9.8-fr}
$X$를 준콤팩트 준스킴,
$(F_\lambda)_{\lambda\in L}$를 교집합이 공집합인 $X$의 프로-구성
가능 부분집합족이라 하자. 그러면 교집합이 공집합인 유한 부분족
$(F_\lambda)_{\lambda\in J}$가 존재한다.
\end{corollary}

\begin{proof}
$X$를 유한 개의 아핀 열린집합으로 덮고 (1.9.5, (iv))를 사용하면
$X$가 아핀인 경우로 환원된다. (1.9.5, (ix))에 의해
$X'_\lambda=\operatorname{Spec}(A'_\lambda)$가 아핀이고
$f_\lambda:X'_\lambda\to X$인 사상에 대하여
$F_\lambda=f_\lambda(X'_\lambda)$로 쓸 수 있다.

\oldpage[IV-1]{246}

(1.9.3.2)의 표기를 사용하자. 가정에 의해 $f(X')=\varnothing$이므로
$X'=\varnothing$이다. 그러면 (1.9.2, (i))에 의해 $X'_J=\varnothing$인
$L$의 유한 부분집합 $J$가 존재한다. 따라서
\[
  f_J(X'_J)=\bigcap_{\lambda\in J}F_\lambda=\varnothing.
\]
\end{proof}

\begin{corollary}[1.9.9]
\label{IV.1.9.9-fr}
$X$를 준콤팩트 준스킴, $F$를 $X$의 프로-구성 가능 부분집합,
$(E_\lambda)_{\lambda\in L}$를
\[
  F\subset\bigcup_{\lambda\in L}E_\lambda
\]
를 만족하는 $X$의 인드-구성 가능 부분집합족이라 하자. 그러면
\[
  F\subset\bigcup_{\lambda\in J}E_\lambda
\]
를 만족하는 $L$의 유한 부분집합 $J$가 존재한다. 특히 $X$의 인드-구성
가능 부분집합들로 이루어진 임의의 덮개에서 유한 덮개를 추출할 수 있다.
\end{corollary}

\begin{proof}
여집합을 취하여 (1.9.7)을 적용하면 된다.
\end{proof}

\begin{proposition}[1.9.10]
\label{IV.1.9.10-fr}
$X$를 준스킴이라 하자. $X$의 부분집합 $E$가 인드-구성 가능이기
위해서는, 모든 $x\in E$에 대하여
$E\cap\overline{\{x\}}$가 $\overline{\{x\}}$에서 $x$의 근방이어야
한다. $X$가 국소 뇌터이면 이 조건은 충분하다.
\end{proposition}

\begin{proof}
$E$가 인드-구성 가능이라고 가정하고 $x\in E$를 고정하자. $x$를
포함하는 $X$의 아핀 열린집합과 $\overline{\{x\}}$의 교집합을 $Y$라
하자. 그러면 바탕공간이 $Y$인 $X$의 부분준스킴이 존재한다
(I, 5.2.1). 표준 포함 사상을 $j:Y\to X$라 하면
\[
  E\cap Y=j^{-1}(E)
\]
는 $Y$에서 인드-구성 가능이다 (1.9.5, (vi)). $Y$는 준콤팩트이고
분리이므로, $E\cap Y$는 $Y$의 구성 가능 부분집합들의 합집합이다
(1.9.5, (iii)). 따라서 $Y$의 두 열린집합 $U,V$가 존재하여
\[
  x\in U\cap\complement V\subset E\cap Y
\]
가 된다 $(0\textsubscript{III}, 9.1.3)$. 그런데 $x$는 기약공간 $Y$의
일반점이므로 $V$는 반드시 공집합이고, 따라서 $U\subset E\cap Y$이다.

이제 $X$가 국소 뇌터이고 $E$가 명제의 조건을 만족한다고 가정하자.
정의 (1.9.4)에 의해 $X$가 뇌터인 경우만 생각하면 충분하다. 그러면
모든 $x\in E$에 대하여 $X$의 어떤 열린집합 $U$가 존재하여
\[
  x\in U\cap\overline{\{x\}}\subset E\cap\overline{\{x\}}
\]
가 된다. 따라서 $E$는 $X$의 국소 닫힌
부분집합들의 합집합이고, 따라서 (1.9.6)에 의해 인드-구성 가능이다.
\end{proof}

\begin{proposition}[1.9.11]
\label{IV.1.9.11-fr}
$X$를 준스킴, $E$를 $X$의 부분집합이라 하자. 다음 성질들은
동치이다.
\begin{enumerate}
  \item[a)] $E$는 국소 구성 가능이다.
  \item[b)] $E$는 인드-구성 가능이고 프로-구성 가능이다.
  \item[c)] $E$와 $X-E$는 프로-구성 가능이다.
  \item[d)] $E$와 $X-E$는 인드-구성 가능이다.
\end{enumerate}
\end{proposition}

\begin{proof}
d)가 a)를 함의함을 보이면 충분하며, $X$가 아핀인 경우만 생각하면
된다. (1.9.5, (iii))에 의해 $E$와 $X-E$를 각각 $X$의 구성 가능
부분집합족 $(E_\lambda)$와 $(E'_\mu)$의 합집합으로 쓸 수 있다.
$E_\lambda$들과 $E'_\mu$들이 $X$를 덮으므로, (1.9.9)에 의해 유한
개의 지표 $\lambda_i,\mu_j$를 택하여 $E_{\lambda_i}$들과
$E'_{\mu_j}$들이 $X$를 덮게 할 수 있다. 따라서 $E$는 유한 개의
$E_{\lambda_i}$들의 합집합이고, 그러므로 구성 가능이다.
\end{proof}

\begin{corollary}[1.9.12]
\label{IV.1.9.12-fr}
$f:X\to Y$를 전사 사상이라 하고, $f$가 준콤팩트이거나 국소 유한
표시라고 하자. $Y$의 부분집합 $E$가 국소 구성 가능 (각각 프로-구성
가능, 인드-구성 가능)이기 위한 필요충분조건은 $f^{-1}(E)$가 $X$에서
같은 성질을 갖는 것이다.
\end{corollary}

\oldpage[IV-1]{247}

\begin{proof}
조건의 필요성은 이미 알려져 있다 ((1.8.2) 및 (1.9.5, (vi))). 충분성을
보이기 위해 (1.9.5, (iv))에 의해 $X$가 아핀인 경우로 환원한다. 또한
(1.9.11)에 의해 $f^{-1}(E)$가 프로-구성 가능한 경우와 인드-구성 가능한
경우를 각각 생각하면 충분하다. $f$가 전사이고 준콤팩트이며
$f^{-1}(E)$가 프로-구성 가능이면, (1.9.5, (vii))에 의해
\[
  E=f(f^{-1}(E))
\]
는 프로-구성 가능이다. $f$가 전사이고 국소 유한 표시이며
$f^{-1}(E)$가 인드-구성 가능이면, (1.9.5, (viii))에 의해
$E=f(f^{-1}(E))$는 인드-구성 가능이다. 이로써 증명이 끝난다.
\end{proof}

\begin{env}[1.9.13]
\label{IV.1.9.13-fr}
$X$를 준스킴, $\mathscr T$를 그 위상이라 하자. (1.9.5, (i), (ii))에
의해 $X$의 인드-구성 가능 부분집합들 (각각 프로-구성 가능 부분집합들)은
$X$ 위의 어떤 위상의 열린집합들 (각각 닫힌집합들)을 이룬다. 이 위상을
\emph{구성 가능 위상}이라 하고 $\mathscr T^{\mathrm{cons}}$로 나타내자.
집합 $X$에 위상 $\mathscr T^{\mathrm{cons}}$를 준 공간을
$X^{\mathrm{cons}}$로 나타낸다.
\end{env}

\begin{proposition}[1.9.14]
\label{IV.1.9.14-fr}
$X$를 준스킴, $\mathscr T$를 그 위상,
$\mathscr T^{\mathrm{cons}}$를 $X$의 구성 가능 위상이라 하자.
\begin{enumerate}
  \item[(i)] 위상 $\mathscr T^{\mathrm{cons}}$는 $\mathscr T$보다
  세밀하다.
  \item[(ii)] $X$의 국소 구성 가능 부분집합들은
  $X^{\mathrm{cons}}$에서 동시에 열리고 닫힌 부분집합들과 일치한다.
  \item[(iii)] 모든 사상 $f:X\to Y$에 대하여 바탕사상
  $X^{\mathrm{cons}}\to Y^{\mathrm{cons}}$는 연속이다. 이를
  $f^{\mathrm{cons}}$로 나타낸다.
  \item[(iv)] 사상 $f:X\to Y$가 준콤팩트이면
  $f^{\mathrm{cons}}$는 닫힌 사상이다. 특히 $f$가 준콤팩트이고
  전단사이면 $f^{\mathrm{cons}}$는 위상동형사상이다.
  \item[(v)] 사상 $f:X\to Y$가 국소 유한 표시이면
  $f^{\mathrm{cons}}$는 열린 사상이다. 특히 $f$가 전단사이고 국소 유한
  표시이면 $f^{\mathrm{cons}}$는 위상동형사상이다.
  \item[(vi)] $X$의 모든 열린집합 $U$에 대하여
  $\mathscr T^{\mathrm{cons}}$가 $U$에 유도하는 위상은
  $U^{\mathrm{cons}}$의 위상과 일치한다.
\end{enumerate}
\end{proposition}

\begin{proof}
(i)은 $\mathscr T$의 모든 열린집합이 $\mathscr T^{\mathrm{cons}}$에서도
열린집합이라는 사실에서 따르고 (1.9.6), (ii)는 (1.9.11)을 다시 쓴
것이다. (iii), (iv), (v)는 각각 (1.9.5)의 (vi), (vii), (viii)을 다시
쓴 것이다. 끝으로 (vi)을 보이려면 표준 포함 사상 $j:U\to X$가 국소
유한 표시이고 (1.4.3), 따라서
\[
  j^{\mathrm{cons}}:U^{\mathrm{cons}}\longrightarrow X^{\mathrm{cons}}
\]
가 연속인 열린 사상임을 주의하면 충분하다.
\end{proof}

\begin{proposition}[1.9.15]
\label{IV.1.9.15-fr}
$X$를 준스킴이라 하자.
\begin{enumerate}
  \item[(i)] $X^{\mathrm{cons}}$가 준콤팩트일 필요충분조건은 $X$가
  준콤팩트인 것이다.
  \item[(ii)] $X^{\mathrm{cons}}$가 하우스도르프 공간일 필요충분조건은
  $X$가 준분리인 것이다. 이때 $X^{\mathrm{cons}}$는 국소 콤팩트이고
  완전 비연결이다.
  \item[(iii)] $X^{\mathrm{cons}}$가 콤팩트일 필요충분조건은 $X$가
  준콤팩트이고 준분리인 것이다.
  \item[(iv)] $X$의 모든 점은 위상 $\mathscr T^{\mathrm{cons}}$에
  관하여 열리고 콤팩트인 근방을 갖는다.
  \item[(v)] 사상 $f:X\to Y$가 준콤팩트일 필요충분조건은 연속사상
  \[
    f^{\mathrm{cons}}:X^{\mathrm{cons}}\longrightarrow Y^{\mathrm{cons}}
  \]
  이 고유사상인 것이다.
\end{enumerate}
\end{proposition}

\begin{proof}
(i) $\mathscr T^{\mathrm{cons}}$가 $\mathscr T$보다 세밀하므로,
$X^{\mathrm{cons}}$가 준콤팩트이면 $X$도 준콤팩트임은 명백하다. 역은
(1.9.9)에서 따른다.

(ii) $X$가 준분리라고 가정하고 $X^{\mathrm{cons}}$가 완전 비연결임을
보이자. 실제로 $x,y$가 $X$의 서로 다른 두 점이고, 예를 들어
$x\notin\overline{\{y\}}$라고 하자.

\oldpage[IV-1]{248}

그러면 $y$를 포함하지 않는 $x$의 아핀 열린근방 $U$가 존재한다.
$X$가 준분리이므로 $U$는 $X$에서 역콤팩트이다 (1.2.7). 따라서
$U$와 $X-U$는 $X$에서 국소 구성 가능이고, (1.9.11)에 의해
$X^{\mathrm{cons}}$에서 열린집합이다. 이는 두 점을 서로 여집합인
열린닫힌집합들로 분리하므로 주장을 얻는다. 또한 $X$의 모든 아핀
열린집합 $U$ 위에 $\mathscr T^{\mathrm{cons}}$가 유도하는 위상은
(1.9.14, (vi))에 의해 $U^{\mathrm{cons}}$의 위상이다. 앞에서 보인
것으로부터 $U^{\mathrm{cons}}$는 콤팩트이고 $U$는
$X^{\mathrm{cons}}$에서 열린집합이므로, $X^{\mathrm{cons}}$는 국소
콤팩트이다.
$X^{\mathrm{cons}}$가 하우스도르프이면 $X$가 준분리임을 증명하는
일만 남는다. 실제로 $X$의 아핀 열린집합 $U$를 생각하자. 이때
$\mathscr T^{\mathrm{cons}}$가 $U$에 유도하는 위상은
$U^{\mathrm{cons}}$의 위상이다 (1.9.14, (vi)). 따라서 앞 논증의 첫
부분에 의해 $U^{\mathrm{cons}}$가 콤팩트이므로, $U$는 이 유도위상에
관하여 콤팩트이다. $V$가 $X$의 또 하나의 아핀 열린집합이면,
$U\cap V$는 하우스도르프 공간 $X^{\mathrm{cons}}$의 두 콤팩트
부분집합의 교집합이므로 그 공간의 콤팩트 부분집합이다. 한편
$\mathscr T^{\mathrm{cons}}$가 $U\cap V$에 유도하는 위상은
$(U\cap V)^{\mathrm{cons}}$의 위상이고 (1.9.14, (vi)), 항등사상
\[
  (U\cap V)^{\mathrm{cons}}\longrightarrow U\cap V
\]
은 연속이다. 따라서 $U\cap V$는 $\mathscr T$가 유도하는 위상에
관하여 준콤팩트이다. 그러므로 (1.2.7)에 의해 $X$는 준분리이다.

(iii)은 (i)와 (ii)의 직접적인 따름정리이다.

(iv) 임의의 $x\in X$에 대하여, $\mathscr T$에 관한 $x$의 아핀
열린근방 $U$는 $\mathscr T^{\mathrm{cons}}$에 관해서도 $x$의
근방이다. 또한 $\mathscr T^{\mathrm{cons}}$가 $U$에 유도하는 위상은
$U^{\mathrm{cons}}$의 위상과 같으므로, (iii)과 (1.9.14, (vi))에
의해 $U$는 콤팩트이다.

(v) $f$가 준콤팩트라고 가정하자. 그러면 이미 (1.9.14, (iv))에 의해
$f^{\mathrm{cons}}$는 닫힌사상이다. 한편 $y$를 $Y$의 한 점이라 하고
\[
  Z=f^{-1}(y)=X\times_Y\operatorname{Spec}(k(y))
\]
라 놓자. $Z$는 준콤팩트이고 (1.1.2, (iii)), 표준사상
$p:Z\to X$는 단사이다. 따라서 (i)와 $p^{\mathrm{cons}}$의 연속성으로부터,
$X^{\mathrm{cons}}$의 위상이 $f^{-1}(y)$에 유도하는 위상 아래에서
$f^{-1}(y)$는 준콤팩트 공간이다. 이는 $f^{\mathrm{cons}}$가
고유사상임을 증명한다 (Bourbaki, \emph{Top. gén.}, chap.
I\textsuperscript{er}, 3\textsuperscript{e} éd., \S~10,
n\textsuperscript{o}~2, th.~1).

역으로 연속사상 $f^{\mathrm{cons}}$가 고유사상이라고 가정하고,
$V$를 $Y$의 준콤팩트 열린집합이라 하자. $h:V\to Y$가 표준 단사이면,
$h^{\mathrm{cons}}:V^{\mathrm{cons}}\to Y^{\mathrm{cons}}$는 연속이고
단사이며, (i)에 의해 $V^{\mathrm{cons}}$는 준콤팩트이다. 따라서
$Y^{\mathrm{cons}}$의 위상이 $V$에 유도하는 위상 아래에서 $V$는
준콤팩트 공간이다. 이제 $f^{\mathrm{cons}}$가 고유사상이라는 가정에
의해, $X^{\mathrm{cons}}$의 위상이 $f^{-1}(V)$에 유도하는 위상
아래에서 $f^{-1}(V)$는 준콤팩트 공간이다 (\emph{loc. cit.}, prop.~6).
그러므로 $f^{-1}(V)$는 $X$의 준콤팩트 부분공간이기도 하며, 이는
사상 $f$가 준콤팩트임을 보인다.
\end{proof}

\begin{env}[1.9.16]
\label{IV.1.9.16-ko}
뒤에서 임의의 준스킴 $X$에 대하여 공간 $X^{\mathrm{cons}}$에 환의
층을 주어, 모든 국소환이 체이고 $X$의 각 점에서의 잉여체와 같은
준스킴으로 만드는 방법을 보일 것이다. 이러한 준스킴은 예를 들어
Néron [32]의 아벨 다양체 축소 이론에 나오는 구성을 스킴의 언어로
옮길 때 자연스럽게 나타난다.
\end{env}

\oldpage[IV-1]{249}

\subsection{열린 사상에의 응용}
\label{subsection:IV.1.10-ko}

\begin{theorem}[1.10.1]
\label{IV.1.10.1-ko}
$X$를 준스킴, $E$를 $X$의 인드-구성 가능 부분집합 (1.9.4), $x$를
$X$의 한 점이라 하자. $x$가 $E$의 내부점일 필요충분조건은 $x$의
모든 일반화 ($0_{\mathrm I}$, 2.1.2) $x'$이 $E$에 속하는 것이다.
다시 말해 (I, 2.4.2),
\[
  \operatorname{Spec}(\mathcal O_{X,x})\subset E
\]
가 성립하는 것이다.
\end{theorem}

\begin{proof}
$x$의 모든 근방은 $x$의 일반화들을 포함하므로 이 조건이 필요함은
명백하다. 충분성을 보이기 위해서는 $X=\operatorname{Spec}(A)$가
아핀이고 $x$가 $A$의 소 아이디얼 $\mathfrak p$인 경우로 한정해도
된다. 이때 (1.9.5, (ix))에 의해 아핀 스킴
$X'=\operatorname{Spec}(A')$과 사상 $f:X'\to X$가 존재하여
$f(X')=X-E$이다. $Y=\operatorname{Spec}(\mathcal O_{X,x})$와
$Y'=X'\times_XY$라 놓자. 가정은 (I, 3.6.5)에 의해
$Y'=\varnothing$, 즉
\[
  A'\otimes_AA_{\mathfrak p}=0
\]
임을 뜻한다. $t$가 $A-\mathfrak p$를 달릴 때
$A_{\mathfrak p}=\varinjlim A_t$이므로 ($0_{\mathrm I}$, 1.4.5),
귀납극한이 텐서곱과 가환한다는 사실로부터
$\varinjlim A_t'=0$을 얻는다. 따라서 (1.9.1)에 의해 어떤
$t\in A-\mathfrak p$에 대하여 $A_t'=0$이고, 이때 $\mathrm D(t)$는
$E$에 포함되는 $X$에서의 $x$의 열린근방이다.
\end{proof}
\begin{corollary}[1.10.2]
\label{IV.1.10.2-ko}
$f:X\to Y$를 국소 유한 표시 사상이라 하고 $y\in Y$라 하자. $y$가
$f(X)$의 내부점일 필요충분조건은 $y$의 모든 일반화 $y'$이 $f(X)$에
속하는 것이다.
\end{corollary}

\begin{proof}
실제로 $f(X)$가 $Y$에서 인드-구성 가능임을 알고 있으므로 (1.9.5,
(viii)), (1.10.1)을 적용하면 충분하다.
\end{proof}

연속사상 $f:X\to Y$가 다음 조건을 만족하면 $x\in X$에서
\emph{열린다}고 하자. $X$에서 $x$의 \emph{모든} 근방의 $f$에 의한
상은 $Y$에서 $f(x)$의 근방이다.

\begin{proposition}[1.10.3]
\label{IV.1.10.3-ko}
$f:X\to Y$를 사상, $x$를 $X$의 한 점이라 하고 $y=f(x)$라 하자.
다음 조건들을 생각하자.
\begin{enumerate}
  \item[a)] $f$는 점 $x$에서 열린다.
  \item[b)] $y$의 모든 일반화 $y'$에 대하여, $f(x')=y'$을 만족하는
  $x$의 일반화 $x'$이 존재한다.
  \item[b')] $\operatorname{Spec}(\mathcal O_{X,x})$의 $f$에 의한 상은
  $\operatorname{Spec}(\mathcal O_{Y,y})$이다.
  \item[c)] $y$를 포함하는 $Y$의 모든 기약 닫힌 부분집합 $Y'$에
  대하여, $X'=f^{-1}(Y')$의 어떤 기약성분은 $x$를 포함하고 그
  제한사상은 $Y'$로 가는 우세 사상이다.
\end{enumerate}
그러면
$\mathrm{a)}\Rightarrow\mathrm{b)}\Leftrightarrow\mathrm{b')}
\Leftrightarrow\mathrm{c)}$가 성립한다. 또한 $f$가 국소 유한 표시이면
네 조건은 서로 동치이다.
\end{proposition}

\begin{proof}
일반화의 정의에 의해 $\operatorname{Spec}(\mathcal O_{X,x})$의
$f$에 의한 상은 언제나 $\operatorname{Spec}(\mathcal O_{Y,y})$에
포함되므로, b)와 b')은 동치이다. c)가 b)를 함의함은 직접적이다.
실제로 $y'$이 $y$의 일반화이면
$Y'=\overline{\{y'\}}$은 $y$를 포함하는 $Y$의 기약 닫힌 부분집합이다.
$x$를 포함하고 $Y'$로 가는 우세 사상을 주는 $f^{-1}(Y')$의 한
기약성분의 일반점을 $x'$이라 하면, $f(x')=y'$이고
($0_{\mathrm I}$, 2.1.5) $x'$은 $x$의 일반화이다.

역으로 b)는 c)를 함의한다. 실제로 $y'$을 $Y'$의 일반점이라 하고,
$f(x')=y'$을 만족하는 $x$의 일반화 $x'$을 택하자. $x'$을 포함하는
$f^{-1}(Y')$의 한 기약성분을 $Z'$이라 하자. 그러면 $Z'$은 $x$도
포함한다. 이미 $f(x')$이 $Y'$에서 조밀하므로 $f(Z')$은 더욱이
$Y'$에서 조밀하다.

a)가 b')을 함의함을 보이기 위해서는 $X=\operatorname{Spec}(B)$와
$Y=\operatorname{Spec}(A)$가 아핀이고, $x$가 $B$의 소 아이디얼
$\mathfrak q$인 경우로 한정해도 된다. 그러면
\[
  \mathcal O_{X,x}=B_{\mathfrak q}=\varinjlim B_t,
\]
여기서 $t$는 $B-\mathfrak q$를 달린다.

\oldpage[IV-1]{250}

따라서 (1.9.2)에 의해
\[
  f(\operatorname{Spec}(\mathcal O_{X,x}))
  =\bigcap_t f(\operatorname{Spec}(B_t))=\bigcap_t f(\mathrm D(t))
\]
이다. 가정에 의해 모든 $t\in B-\mathfrak q$에 대하여 $y$는
$f(\mathrm D(t))$의 내부점이다. 따라서 $f(\mathrm D(t))$는
$\operatorname{Spec}(\mathcal O_{Y,y})$를 포함하며, 이로부터
\[
  \operatorname{Spec}(\mathcal O_{Y,y})
  \subset f(\operatorname{Spec}(\mathcal O_{X,x}))
\]
을 얻어 주장이 증명된다. 마지막으로 $f$가 국소 유한 표시이면,
$x$의 임의의 열린근방 $U$에 대한 $f$의 제한 $U\to Y$에 (1.10.2)를
적용하여 b)가 a)를 함의함을 얻는다.
\end{proof}

\begin{corollary}[1.10.4]
\label{IV.1.10.4-ko}
$f:X\to Y$를 사상이라 하자. 다음 조건들을 생각하자.
\begin{enumerate}
  \item[a)] $f$는 열린사상이다.
  \item[b)] 모든 $x\in X$와 $y=f(x)$의 모든 일반화 $y'$에 대하여,
  $f(x')=y'$을 만족하는 $x$의 일반화 $x'$이 존재한다.
  \item[b')] 모든 $x\in X$에 대하여,
  $\operatorname{Spec}(\mathcal O_{X,x})$의 $f$에 의한 상은
  $\operatorname{Spec}(\mathcal O_{Y,f(x)})$이다.
  \item[c)] $Y$의 모든 기약 닫힌 부분집합 $Y'$에 대하여,
  $f^{-1}(Y')$의 모든 기약성분에서 $Y'$로 가는 제한사상은
  우세 사상이다.
\end{enumerate}
그러면
$\mathrm{a)}\Rightarrow\mathrm{b)}\Leftrightarrow\mathrm{b')}
\Leftrightarrow\mathrm{c)}$가 성립한다. 또한 $f$가 국소 유한 표시이면
네 조건은 서로 동치이다.
\end{corollary}
\begin{proof}
$f$가 열린사상이라는 것은 모든 점 $x\in X$에서 열린다는 뜻이다.
따라서 $\mathrm{a)}\Rightarrow\mathrm{b)}\Leftrightarrow
\mathrm{b')}$와, $f$가 국소 유한 표시일 때의
$\mathrm{b)}\Rightarrow\mathrm{a)}$는 (1.10.3)의 대응하는
함의들에서 따른다. $\mathrm{c)}\Rightarrow\mathrm{b)}$도 (1.10.3)의
대응하는 함의에서 따른다. 마지막으로 b)가 c)를 함의함을 보이자.
실제로 $x'$을 $f^{-1}(Y')$의 한 기약성분의 일반점이라 하고,
$y'=f(x')$이 $Y'$의 일반점임을 보이자. $y''$이 $y'$의 일반화이면,
가정에 의해 $f(x'')=y''$을 만족하는 $x'$의 일반화 $x''$이 존재한다.
$x''\in f^{-1}(Y')$이고 $x'$이 그 기약성분의 일반점이므로 반드시
$x''=x'$이다. 따라서 $y''=y'$이고, 이로써 증명이 끝난다.
\end{proof}

\begin{center}
\emph{(계속.)}
\end{center}

\clearpage
\oldpage[IV-1]{251}
\phantomsection
\addcontentsline{toc}{section}{참고문헌 (계속)}
\section*{참고문헌 (\emph{계속})}

\begingroup
\small
\setlength{\parindent}{0pt}
\setlength{\parskip}{3pt}
[30] M.~Nagata, \emph{Local rings}, \emph{Interscience Tracts}, 13 (1962),
Interscience, New York.

[31] M.~Nagata, \emph{A Jacobian criterion of simple points},
\emph{Ill. Journ. of Math.}, 1 (1957), p.~427--432.

[32] A.~Néron, \emph{Modèles minimaux des variétés abéliennes sur les corps
locaux et globaux}, \emph{Publ. Inst. Hautes Études Sci.}, n\textsuperscript{o}
21 (1964).

\clearpage
\oldpage[IV-1]{252}
\phantomsection
\addcontentsline{toc}{section}{기호 색인}
\section*{기호 색인}

$\dim_c(X)$, $\dim(X)$ ($X$는 위상공간) : 0, 14.1.2.\par
$\dim_x(X)$ ($X$는 위상공간, $x$는 $X$의 점) : 0, 14.1.2.\par
$\operatorname{codim}(Y,X)$ ($X$는 위상공간, $Y$는 $X$의 닫힌 부분집합) : 0, 14.2.1.\par
$\operatorname{codim}_x(Y,X)$ ($X$는 위상공간, $Y$는 $X$의 닫힌 부분집합, $x$는 $X$의 점) : 0, 14.2.4.\par
$N[T_1,\ldots,T_n]$ ($N$은 $A$-가군) : 0, 15.1.1.\par
$\dim(A)$ ($A$는 환) : 0, 16.1.1.\par
$\operatorname{ht}(\mathfrak a)$ ($\mathfrak a$는 아이디얼) : 0, 16.1.3.\par
$\dim_A(M)$, $\dim(M)$ ($M$은 $A$-가군) : 0, 16.1.7.\par
$P_{\mathfrak q}(M,n)$ ($M$은 뇌터 반국소환 $A$ 위의 유한형 가군,
$\mathfrak q$는 $A$의 정의 아이디얼) : 0, 16.2.1.\par
$\operatorname{prof}_A(M)$, $\operatorname{prof}(M)$ ($M$은 $A$-가군) : 0, 16.4.5.\par
$\operatorname{coprof}_A(M)$, $\operatorname{coprof}(M)$ ($M$은 $A$-가군) : 0, 16.4.9.\par
$\operatorname{dim.proj}_A(M)$, $\operatorname{dim.proj}(M)$,
$\operatorname{dim.inj}_A(M)$, $\operatorname{dim.inj}(M)$ ($M$은 $A$-가군) : 0, 17.2.1.\par
$\operatorname{dim.coh}(A)$ ($A$는 환) : 0, 17.2.8.\par
$\operatorname{dim.proj}(\mathscr F)$, $\operatorname{dim.inj}(\mathscr F)$
($\mathscr F$는 환 달린 공간 위의 가군) : 0, 17.2.14.\par
$\operatorname{dim.coh}(X)$ ($X$는 환 달린 공간) : 0, 17.2.14.\par
$E\times_B F$ ($B$, $E$, $F$는 $A$-환) : 0, 18.1.2.\par
$D_B(L)$ ($B$는 환, $L$은 $B$-쌍가군) : 0, 18.2.3.\par
$Y\oplus_X Z$ ($X$, $Y$, $Z$는 $A$-쌍가군) : 0, 18.2.7.\par
$\operatorname{Exan}_A(B,L)$ ($B$는 $A$-환, $L$은 $B$-쌍가군) : 0, 18.3.4.\par
$\operatorname{Exan}_{A/A'}(B,L)$ : 0, 18.3.7.\par
$\operatorname{Exal}_A(B,L)$, $\operatorname{Exal}_{A/A'}(B,L)$ ($A$는
가환환, $B$는 $A$-대수, $L$은 $B$-쌍가군) : 0, 18.4.1.\par
$\operatorname{Exalcom}_A(B,L)$, $\operatorname{Exalcom}_{A/A'}(B,L)$
($A$는 가환환, $B$는 가환 $A$-대수, $L$은 $B$-가군) : 0, 18.4.2.\par
$\mathrm H_A^2(B,L)$, $\mathrm H_A^2(B,L)^s$ : 0, 18.4.3.\par
$\operatorname{Exantop}_A(B,L)$, $\operatorname{Exaltop}_A(B,L)$,
$\operatorname{Exalcotop}_A(B,L)$ : 0, 18.5.1.\par
$\operatorname{Exantop}_{A/A'}(B,L)$, $\operatorname{Exaltop}_{A/A'}(B,L)$,
$\operatorname{Exalcotop}_{A/A'}(B,L)$ : 0, 18.5.2.\par
$\operatorname{Hom.\ cont}_C(M,N)$ ($M$, $N$은 위상 $C$-쌍가군) : 0, 18.5.3.\par
$\operatorname{Dér}_A(B,L)$ ($B$는 $A$-환, $L$은 $B$-쌍가군) : 0, 20.1.2.\par
$\operatorname{Dér.\ cont}_A(B,L)$ : 0, 20.3.1.\par
$p_{B/A}$, $\mathfrak J_{B/A}$ ($B$는 가환 위상 $A$-대수) : 0, 20.4.1.\par
$P^1_{B/A}$, $\varepsilon_{B/A}$ ($B$는 가환 위상 $A$-대수) : 0, 20.4.2.\par
$\Omega^1_{B/A}$, $\widehat{\Omega}^1_{B/A}$ ($B$는 가환 위상 $A$-대수) : 0, 20.4.3.\par
$\Omega^1_B$ ($B$는 위상환) : 0, 20.4.3.\par
$d_{B/A}$, $d(x)$, $dx$, $d_B$ : 0, 20.4.6.\par
$\Upsilon_{C/B/A}$, $\Upsilon_{C/B}$ ($A$, $B$, $C$는 환) : 0, 20.6.1.\par
$\chi_E$ ($E$는 $A$-자명한, $L$에 의한 $C$의 $B$-확장) : 0, 20.6.8.\par

\clearpage
\oldpage[IV-1]{253}
$\Upsilon^C_{B/A/\Lambda}$, $\Upsilon^C_{B/A}$ ($\Lambda$, $A$, $B$, $C$는 환) : 0, 20.6.14.\par
$\chi_{B/A}$, $\chi_B$ ($B$는 $A$-대수) : 0, 20.6.24.\par
$\chi_E$ ($E$는 위상 $B$-대수) : 0, 20.7.10.\par
$F_A$, $A^p$, $A^{(p)}$, $E^{(p)}$ ($A$는 표수 $p>0$인 환,
$E$는 $A$-가군) : 0, 21.1.4.\par
$\pi_{B/A}$, $\Theta_{B/A}$, $\Xi_{B/A}$ ($A$, $B$는 표수 $p>0$인 환) : 0, 21.3.2.\par
$\Delta(L/K,k/k_0)$, $d(L/K,k/k_0)$, $\Delta(L/K,k)$, $d(L/K,k)$
($k_0$, $k$, $K$, $L$은 체) : 0, 21.6.1.\par
$\chi'_{B/A}$ ($A$, $B$는 표수 $p>0$인 환) : 0, 22.4.6.\par
$\mathscr T^{\mathrm{cons}}$, $X^{\mathrm{cons}}$ ($X$는 준스킴,
$\mathscr T$는 $X$의 위상) : IV, 1.9.13.\par
$f^{\mathrm{cons}}$ ($f$는 준스킴의 사상) : IV, 1.9.14.\par

\clearpage
\oldpage[IV-1]{254}
\phantomsection
\addcontentsline{toc}{section}{용어 색인}
\section*{용어 색인}

1차 주요 부분의 증대 대수 : 0, 20.4.2.\par
Cohen $A$-대수 : 0, 19.8.1.\par
유한 표시 대수 : IV, 1.4.1.\par
본질적 유한형 $A$-대수 : IV, 1.3.8.\par
$A$-환 : 0, 18.1.1.\par
$B$ 위에 증대된 $A$-환 : 0, 18.1.4.\par
쌍등차원 환 : 0, 16.1.4.\par
카테너리 환 : 0, 16.1.4.\par
표수 $p$인 환 : 0, 21.1.1.\par
Cohen 환 : 0, 19.8.4.\par
Cohen--Macaulay 환 : 0, 16.5.1 및 16.5.13.\par
등여차원 환, 등차원 환 : 0, 16.1.4.\par
기하학적 단일분지 국소환 : 0, 23.1.7.\par
일본환 : 0, 23.1.1.\par
소 국소환, 완비 소 국소환 : 0, 19.8.3.\par
정칙환 : 0, 17.1.1 및 17.3.6.\par
단일분지 국소환 : 0, 23.1.7.\par
보편 일본환 : 0, 23.1.1.\par
한 점에서 열린 사상 : IV, 1.10.2.\par
증대 : 0, 18.1.4.\par

형식적 전단사 준동형 : 0, 19.1.2.\par

순서집합의 사슬, 위상공간의 기약 닫힌 부분집합들의 사슬 : 0, 14.1.1.\par
포화 사슬 : 0, 14.3.1.\par
조합적 여차원, 위상공간의 닫힌 부분집합의 여차원 : 0, 14.2.1.\par
위상공간의 닫힌 부분집합의 한 점에서의 여차원 : 0, 14.2.4.\par
사슬 조건 : 0, 14.3.2.\par
뇌터 국소환 위 가군의 여깊이 : 0, 16.4.9.\par
$k_0$-허용체, 확대체에 대한 허용체 : 0, 21.6.1.\par
부분체 위에서 유한 순수 비분리 중복도를 갖는 체 : 0, 19.6.6.\par

확대체의 $k_0$-허용성 결함 : 0, 21.6.1.\par
미분, $A$-미분 : 0, 20.1.2.\par
미분 가군, 1차 미분 가군, 절대 미분 가군 : 0, 20.4.3.\par
외미분 : 0, 20.4.6.\par
조합적 차원, 위상공간의 차원 : 0, 14.1.2.\par
위상공간의 한 점에서의 차원 : 0, 14.1.2.\par
크룰 차원, 환의 차원 : 0, 16.1.1.\par
가군의 차원 : 0, 16.1.7.\par
환의 코호몰로지 차원, 전역 코호몰로지 차원 : 0, 17.2.8.\par
한 점에서의 환 달린 공간의 코호몰로지 차원, 환 달린 공간의
코호몰로지 차원 : 0, 17.2.14.\par

\clearpage
\oldpage[IV-1]{255}
가군의 단사 차원, 사영 차원 : 0, 17.2.1.\par
한 점에서의 단사 차원, 층의 단사 차원 : 0, 17.2.14.\par
한 점에서의 사영 차원, 층의 사영 차원 : 0, 17.2.14.\par

$A$-가군에서 영인자가 아닌 $A$의 원소 : 0, 15.1.1.\par
$M$-정칙 원소, $M$-준정칙 원소 : 0, 15.1.4.\par
인드-구성 가능 집합, 프로-구성 가능 집합 : IV, 1.9.4.\par
형식적 전사 준동형 : 0, 19.1.2.\par
$A$-확장의 $A$-동치 : 0, 18.2.2.\par
쌍등차원 공간 : 0, 14.3.3.\par
카테너리 공간 : 0, 14.3.2.\par
등여차원 공간 : 0, 14.2.1.\par
등차원 공간 : 0, 14.1.3.\par
$B$-쌍가군에 의한 $A$-환 $B$의 $A$-확장 : 0, 18.2.2.\par
쌍가군 준동형으로부터 다른 확장에서 유도된 확장 : 0, 18.2.8.\par
호흐실트 확장 : 0, 18.4.3.\par
$A$-자명한 $B$-확장 : 0, 18.3.7.\par
표준 자명 확장 : 0, 18.2.3.\par
$A$-동치인 $A$-확장들 : 0, 18.2.2.\par

$p$-자유족, 절대 $p$-자유족 : 0, 21.1.9.\par
형식적으로 에탈인 대수 : 0, 19.10.2.\par
형식적으로 오른쪽 가역인 준동형 : 0, 19.1.15.\par
형식적으로 왼쪽 가역인 준동형 : 0, 19.1.5.\par
형식적으로 비분기인 대수 : 0, 19.10.2.\par
형식적으로 사영인 가군 : 0, 19.2.1.\par
형식적으로 매끄러운 대수 : 0, 19.3.1.\par
환에 대하여 상대적으로 형식적으로 매끄러운 대수 : 0, 19.9.1.\par

체 위에서 기하학적으로 정칙인 국소환 : 0, 19.6.5.\par
$A$-확장류의 군 : 0, 18.3.4.\par

아이디얼의 높이 : 0, 16.1.3.\par
$A$-환의 $A$-준동형 : 0, 18.1.1.\par
$A$-자명한 $B$-확장의 특성 준동형 : 0, 20.6.8.\par
환과 아이디얼에 대한 $A$-대수의 특성 준동형 : 0, 20.6.24.\par
표수 $p>0$인 대수의 특성 준동형 : 0, 22.4.6.\par

증대 아이디얼 : 0, 18.1.4.\par
증대환의 역상 : 0, 18.1.5.\par
$A$-확장의 역상 : 0, 18.2.5.\par

사슬의 길이 : 0, 14.1.1.\par

쌍등차원, 카테너리, 등여차원, 등차원 가군 : 0, 16.1.7.\par
Cohen--Macaulay 가군 : 0, 16.5.1 및 16.5.13.\par
불완전성 가군 : 0, 20.6.1.\par
형식적 단사 준동형 : 0, 19.1.2.\par
$A$-확장의 사상 : 0, 18.2.4.\par
한 점에서 유한 표시인 사상, 국소 유한 표시 사상 : IV, 1.4.2.\par
유한 표시 사상 : IV, 1.6.1.\par
준분리 사상 : IV, 1.2.1.\par

Cohen의 $p$-환 : 0, 19.8.4.\par
$p$-기저, 절대 $p$-기저 : 0, 21.1.9.\par
준스킴의 극대점 : IV, 1.1.4.\par

\clearpage
\oldpage[IV-1]{256}
힐베르트--사무엘 다항식 : 0, 16.2.1.\par
다른 준스킴 위에서 유한 표시인 준스킴 : IV, 1.6.1.\par
준분리 준스킴 : IV, 1.2.1.\par
$A$-환의 섬유곱 : 0, 18.1.2.\par
가군의 깊이 : 0, 16.4.5.\par

쌍가군의 밂 : 0, 18.2.7.\par
엄밀하게 형식적으로 사영인 가군 : 0, 19.2.3.\par
$M$-정칙열, $M$-준정칙열 : 0, 15.1.7.\par
$\mathscr F$-정칙열, $\mathscr F$-준정칙열 : 0, 15.2.2.\par
뇌터 반국소환 위 가군의 매개변수계 : 0, 16.3.6.\par
$p$-생성원계, 절대 $p$-생성원계 : 0, 21.1.9.\par
열린 아이디얼들, 열린 부분가군들의 기본계 : 0, 19.0.2.\par
정칙 국소환의 정칙 매개변수계 : 0, 17.1.6.\par

준스킴의 구성 가능 위상 : IV, 1.9.13.\par
$A$의 위상으로부터 유도된 $A$-가군의 위상 : 0, 19.0.2.\par
자명한 증대 $A$-환 : 0, 18.1.4.\par
$A$-자명한 확장 : 0, 18.2.3 및 18.3.7.\par

\clearpage
\oldpage[IV-1]{257}
\phantomsection
\addcontentsline{toc}{section}{목차}
\section*{목차}

\newcommand{\ivonecontentsline}[3]{\noindent\makebox[17mm][l]{#1}\parbox[t]{0.77\linewidth}{#2\dotfill}\makebox[12mm][r]{#3}\par}
\ivonecontentsline{제0장.}{\textbf{예비사항} (\emph{계속})}{5}
\ivonecontentsline{\S~14.}{위상공간의 조합적 차원}{6}
\ivonecontentsline{14.1.}{위상공간의 조합적 차원}{6}
\ivonecontentsline{14.2.}{닫힌 부분집합의 여차원}{8}
\ivonecontentsline{14.3.}{사슬 조건}{10}
\ivonecontentsline{\S~15.}{$M$-정칙열과 $\mathscr F$-정칙열}{12}
\ivonecontentsline{15.1.}{$M$-정칙열과 $M$-준정칙열}{12}
\ivonecontentsline{15.2.}{$\mathscr F$-정칙열}{20}
\ivonecontentsline{\S~16.}{뇌터 국소환의 차원과 깊이}{22}
\ivonecontentsline{16.1.}{환의 차원}{22}
\ivonecontentsline{16.2.}{뇌터 반국소환의 차원}{25}
\ivonecontentsline{16.3.}{뇌터 국소환의 매개변수계}{28}
\ivonecontentsline{16.4.}{깊이와 여깊이}{32}
\ivonecontentsline{16.5.}{Cohen--Macaulay 가군}{36}
\ivonecontentsline{\S~17.}{정칙환}{39}
\ivonecontentsline{17.1.}{정칙환의 정의}{39}
\ivonecontentsline{17.2.}{가군의 사영 차원과 단사 차원에 관한 복습}{42}
\ivonecontentsline{17.3.}{정칙환의 코호몰로지 이론}{46}
\ivonecontentsline{\S~18.}{대수의 확장에 관한 보충}{51}
\ivonecontentsline{18.1.}{증대환의 역상}{51}
\ivonecontentsline{18.2.}{쌍가군에 의한 환의 확장}{54}
\ivonecontentsline{18.3.}{$A$-확장류의 군}{59}
\ivonecontentsline{18.4.}{대수의 확장}{64}
\ivonecontentsline{18.5.}{위상환의 경우}{66}
\ivonecontentsline{\S~19.}{형식적으로 매끄러운 대수와 Cohen 환}{69}
\ivonecontentsline{19.0.}{서론}{69}
\ivonecontentsline{19.1.}{형식적 전사 준동형과 형식적 단사 준동형}{71}
\ivonecontentsline{19.2.}{형식적으로 사영인 가군}{78}

\clearpage
\oldpage[IV-1]{258}
\ivonecontentsline{19.3.}{형식적으로 매끄러운 대수}{79}
\ivonecontentsline{19.4.}{형식적 매끄러움의 첫 판정법들}{86}
\ivonecontentsline{19.5.}{형식적 매끄러움과 연관 등급환}{90}
\ivonecontentsline{19.6.}{체 위의 대수의 경우}{100}
\ivonecontentsline{19.7.}{국소 준동형의 경우: 존재성과 유일성 정리}{104}
\ivonecontentsline{19.8.}{Cohen 대수와 Cohen의 $p$-환; 완비 국소환의 구조에 대한 응용}{109}
\ivonecontentsline{19.9.}{상대적으로 형식적으로 매끄러운 대수}{114}
\ivonecontentsline{19.10.}{형식적으로 비분기인 대수와 형식적으로 에탈인 대수}{115}
\ivonecontentsline{\S~20.}{미분과 미분 가군}{116}
\ivonecontentsline{20.1.}{미분과 대수의 확대}{117}
\ivonecontentsline{20.2.}{미분의 함자적 성질}{119}
\ivonecontentsline{20.3.}{위상환의 연속 미분}{121}
\ivonecontentsline{20.4.}{주요 부분과 미분}{123}
\ivonecontentsline{20.5.}{$\Omega^1_{B/A}$의 기본 함자적 성질}{128}
\ivonecontentsline{20.6.}{불완전성 가군과 특성 준동형}{136}
\ivonecontentsline{20.7.}{위상환으로의 일반화}{147}
\ivonecontentsline{\S~21.}{표수 $p$인 환에서의 미분}{153}
\ivonecontentsline{21.1.}{$p$-생성원계와 $p$-기저}{154}
\ivonecontentsline{21.2.}{$p$-기저와 형식적 매끄러움}{157}
\ivonecontentsline{21.3.}{$p$-기저와 불완전성 가군}{160}
\ivonecontentsline{21.4.}{체 확대의 경우}{162}
\ivonecontentsline{21.5.}{응용: 분리성 판정법}{164}
\ivonecontentsline{21.6.}{확대체에 대한 허용체}{167}
\ivonecontentsline{21.7.}{카르티에 등식}{169}
\ivonecontentsline{21.8.}{허용성 판정법}{171}
\ivonecontentsline{21.9.}{형식적 멱급수환에서의 미분 가군의 완비화}{176}
\ivonecontentsline{\S~22.}{형식적 매끄러움과 정칙성의 미분 판정법}{182}
\ivonecontentsline{22.1.}{형식적 매끄러움의 올림}{183}
\ivonecontentsline{22.2.}{체 위에서 형식적으로 매끄러운 국소 대수의 미분적 특징짓기}{186}
\ivonecontentsline{22.3.}{몇몇 국소환과 그 완비화 사이의 관계에 대한 응용}{191}
\ivonecontentsline{22.4.}{극대 아이디얼의 제곱이 영인 국소환의 유한 확대에 관한 예비 결과}{193}

\clearpage
\oldpage[IV-1]{259}
\ivonecontentsline{22.5.}{기하학적으로 정칙인 대수와 형식적으로 매끄러운 대수}{201}
\ivonecontentsline{22.6.}{자리스키의 야코비 판정법}{205}
\ivonecontentsline{22.7.}{나가타의 야코비 판정법}{209}
\ivonecontentsline{\S~23.}{일본환}{213}
\ivonecontentsline{23.1.}{일본환}{213}
\ivonecontentsline{23.2.}{뇌터 국소 정역의 정수적 폐포}{217}
\ivonecontentsline{}{\textsc{제IV장.} --- \textbf{스킴과 스킴 사상의 국소적 연구}}{222}
\ivonecontentsline{\S~1.}{상대적 유한성 조건. 준스킴의 구성 가능 집합}{224}
\ivonecontentsline{1.1.}{준콤팩트 사상}{224}
\ivonecontentsline{1.2.}{준분리 사상}{226}
\ivonecontentsline{1.3.}{국소 유한형 사상}{228}
\ivonecontentsline{1.4.}{국소 유한 표시 사상}{230}
\ivonecontentsline{1.5.}{유한형 사상}{233}
\ivonecontentsline{1.6.}{유한 표시 사상}{234}
\ivonecontentsline{1.7.}{앞선 결과들의 개선}{236}
\ivonecontentsline{1.8.}{유한 표시 사상과 구성 가능 집합}{238}
\ivonecontentsline{1.9.}{프로-구성 가능 집합과 인드-구성 가능 집합}{241}
\ivonecontentsline{1.10.}{열린 사상에의 응용}{249}
\ivonecontentsline{}{참고문헌 (\emph{계속})}{251}
\ivonecontentsline{}{기호 색인}{252}
\ivonecontentsline{}{용어 색인}{254}

\vspace{2em}
\begin{flushright}
\emph{1963년 3월 15일 접수.}
\end{flushright}
\let\ivonecontentsline\relax
\endgroup
